\section*{第 1 章 集合与逻辑(必修 1)}

\section*{知识梳理*}

\section*{一、集合及其运算}

1.集合的特性:集合中的元素具有确定性、无序性、互异性;

2. 列举法:将集合中的元素一一列举出来写在大括号内表示集合的方法;

3. 描述法:将集合中元素的通性描述出来写在大括号内表示集合的方法；通式: $\{ x \mid  x$ 满足 $P\}$ ；读法:满足条件 $P$ 的 $x$ 的集合；

4. 空集(《》)是不含任何元素的集合；它是任何集合的子集，是任何非空集合的真子集；

5. “ $\in  , \notin$ ”表示元素与集合间的从属关系; “ $\subset  , \subseteq$ ”表示集合与集合间的包含关系.

6. 给出下列条件:

(1)集合 $A$ 中任何一个元素都是集合 $B$ 中的元素;

(2)集合 $B$ 至少存在一个元素不在集合 $A$ 中；

(3)集合 $B$ 中任何一个元素都是集合 $A$ 中的元素.

如果集合 $A$ 、 $B$ 满足 \ding{172} ，则 $A$ 是 $B$ 的子集( $A \subseteq  B$ )；

如果集合 $A$ 、 $B$ 满足 \ding{172}  \ding{173} ，则 $A$ 是 $B$ 的真子集( $A \subset  B$ )；

如果集合 $A$ 、 $B$ 满足 \ding{172}  \ding{174} ，则 $A$ 与 $B$ 是相等的集合 $\left( {A = B}\right)$ .

注意: \ding{172}  $A \subseteq  B$ ，讨论时别忘 $A = \varnothing$ 的情况；考察集合的关系借助韦恩图；

 \ding{173} 几个特殊数集之间的包含关系为: ${N} \subset  \mathbf{Z} \subset  \mathbf{Q} \subset  \mathbf{R} \subset  \mathbf{C}$ (本书中 ${\mathbf{N}}^{ * }$ 表示正整数集);

 \ding{174} 正确区分 $\varnothing \text{ 、 }\{ 0\} \text{ 、 }\{ \varnothing \}  : \varnothing$ 是不含任何元素的集合，即空集. \{0\}是含有一个元素 0 的集合，它不是空集， 因为它有一个元素,这个元素是 $0.\{ \varnothing \}$ 是含有一个元素 $\varnothing$ 的集合,其关系为: $\varnothing  \subset  \{ 0\} ,\varnothing  \subset  \{ \varnothing \} ,\varnothing  \in  \{ \varnothing \}$ .

7. 集合的含义:

(1) $A = \{ x \mid  y = f\left( x\right) \}$ 表示函数的定义域；

(2) $B = \{ y \mid  y = f\left( x\right) \}$ 表示函数的值域；

(3) $C = \{ \left( {x,y}\right)  \mid  y = f\left( x\right) \}$ 表示方程 $y = f\left( x\right)$ 的解的集合，或表示曲线上的点的集合……

8. 集合的运算: $A \cap  B = \{ x \mid  x \in  A$ 且 $x \in  B\}$ ； $A \cup  B = \{ x \mid  x \in  A$ 或 $x \in  B\}$ ； $\overline{A} = \{ x \mid  x \in  U$ ，且 $x \notin  A\}$ (其中 $U$ 表示全集).

9. 运算性质:

$A \subseteq  B$ ,且 $B \subseteq  C \Rightarrow  A \subseteq  C$ ;

$\varnothing  \subseteq  A \cap  B = B \cap  A \subseteq  A$ (或 $B$ ) $\subseteq  A \cup  B = B \cup  A$ ;

$\bar{A} \cap  \bar{B} = \overline{A \cup  B};\bar{A} \cup  \bar{B} = \overline{A \cap  B};\overline{\varnothing } = U;\bar{U} = \varnothing$ .

10.几个结论:

$B \cup  \bar{A} = U \Leftrightarrow  A \subseteq  B;\bar{B} \subseteq  \bar{A} \Leftrightarrow  A \cap  \bar{B} = \varnothing$ ;

$A \subseteq  B \Leftrightarrow  A \cap  B = A \Leftrightarrow  A \cup  B = B \Leftrightarrow  A \cap  B \supseteq  A \Leftrightarrow  A \cup  B \subseteq  B$ .

11. 正整数 (整数) 分类: 被 2 整除与否可分为 ${2k} - 1,{2k}\left( {k \in  \mathbf{N}\text{ (或 }\mathbf{Z}}\right) ,k \geq  1$ );

被 3 整除与否可分为 ${3k} - 2,{3k} - 1,{3k}\left( {k \in  \mathbf{N}\text{ (或 }\mathbf{Z}}\right) ,k \geq  1$ );

被 4 整除与否可分为 ${4k} - 3,{4k} - 2,{4k} - 1,{4k}\left( {k \in  \mathbf{N}\text{ (或 }\mathbf{Z}}\right) ,k \geq  1$ ); 其余依此类推.

\HRule

* 注 本章节所属知识模块是“预备知识”.

\HRule

12. $n$ 个元素的子集有 ${{C}}_{n}^{0} + {{C}}_{n}^{1} + {{C}}_{n}^{2} + \cdots  + {{C}}_{n}^{n} = {2}^{n}$ 个;

真子集有 ${{C}}_{n}^{0} + {{C}}_{n}^{1} + {{C}}_{n}^{2} + \cdots  + {{C}}_{n}^{n - 1} = {2}^{n} - 1$ 个;

非空子集有 ${{C}}_{n}^{1} + {{C}}_{n}^{2} + {{C}}_{n}^{3} + \cdots  + {{C}}_{n}^{n} = {2}^{n} - 1$ 个;

非空真子集有 ${{C}}_{n}^{1} + {{C}}_{n}^{2} + {{C}}_{n}^{3} + \cdots  + {{C}}_{n}^{n - 1} = {2}^{n} - 2$ 个.

13.* 容斥原理:记集合 $M$ 中的元素个数为 $\operatorname{Card}\left( M\right)$ ，则

$\operatorname{Card}\left( {A \cup  B}\right)  = \operatorname{Card}\left( A\right)  + \operatorname{Card}\left( B\right)  - \operatorname{Card}\left( {A \cap  B}\right)$ (可推广到有限个集合).

\section*{二、命题与充要条件}

1. 可以判断真假的语句叫命题. 其一般形式是:“如果……，那么 ······”(或“若 $\alpha$ ，则 $\beta$ ”)，命题有真命题和假命题两类.

要判断命题“若 $\alpha$ ，则 $\beta$ ”是假命题，只要存在一个满足条件 $\alpha$ 但不满足结论 $\beta$ 的对象就可以了. 要判断命题“若 $\alpha$ ,则 $\beta$ ”是真命题,就需要证明所有满足条件 $\alpha$ 的对象都满足结论 $\beta$ ; 若要否定命题,切记举反例.

2. 证明命题成立的方法:

综合法:从已知的条件出发，根据已知的定义、公理、定理和推论进行严密的逻辑推理，逐步得到结论的方法. 即 “由因导果”. 连接词: “ “ …… ， ， …… ... ”。

分析法:从需要的结论出发，一步步寻找结论成立的充分条件，直到找到的条件是已知的定义、公理、定理和推论或题设条件的方法. 即“执果索因”. 连接词:“要证……，只要证……”.

反证法:当证明“若 $\alpha$ ，则 $\beta$ ”感到困难时，可从它的反面入手证明.

证题步骤: \ding{172} 假设命题的结论 $\beta$ 不成立； \ding{173} 从这个假设出发，经过正确的推理论证，得出矛盾； \ding{174} 由矛盾判断假设不成立，从而肯定结论正确.

矛盾的来源: \ding{172}  与原命题的条件； \ding{173}  与假设； \ding{174}  恒真命题.

适用于待证命题的结论涉及“不可能”“不是”“至少”“至多”“唯一”等否定词.

\begin{center}
\begin{tabular}{|c|c|c|c|c|c|c|c|c|}
\hline
正面词语 & 等于 & \multicolumn{2}{|c|}{大于} & \multicolumn{2}{|c|}{小于} & \multicolumn{2}{|c|}{是} & 都是 \\
\cline{1-9}
否定 & 不等于 & \multicolumn{2}{|c|}{小于或等于} & \multicolumn{2}{|c|}{大于或等于} & \multicolumn{2}{|c|}{不是} & 不都是 \\
\cline{1-9}
正面词语 & \multicolumn{2}{|c|}{至少一个} & \multicolumn{2}{|c|}{任意的} & \multicolumn{2}{|c|}{所有的} & \multicolumn{2}{|c|}{至多有 $n$ 个} \\
\cline{1-9}
否定 & \multicolumn{2}{|c|}{一个都没有} & \multicolumn{2}{|c|}{存在一个} & \multicolumn{2}{|c|}{存在一个} & \multicolumn{2}{|c|}{至少 $n + 1$ 个} \\
\cline{1-9}
\hline
\end{tabular}
\end{center}

3. 充要条件及其判定 $\left( {A\text{ 、 }B\text{ 分别可以是命题，也可以是一个“开语句”): }}\right)$

若 $A \Leftrightarrow  B$ (或 $\bar{B} \Rightarrow  \bar{A}$ ),则 $A$ 是 $B$ 的充分条件;

若 $A \Leftarrow  B$ (或 $\overline{\overline{A \Rightarrow  B}}$ ),则 $A$ 是 $B$ 的必要条件;

若 $A \Leftrightarrow  B$ (或 $\bar{A} \Leftrightarrow  \bar{B}$ ),则 $A$ 是 $B$ 的充分必要条件,简称充要条件.

4. 充要条件的判断中,分清“条件”和“结论”是关键,其次是看推出符号的箭头指向.

5. 若命题 $\alpha$ 成立的范围是集合 $A$ ，命题 $\beta$ 成立的范围是集合 $B\texttt{ ( }A\text{ 、 }B\texttt{ 非空 }\texttt{ ),则 }{}^{\alpha }A\text{ 、 }B$ 之间的包含关系” 与“充要条件”之间有如下的关系:

若 $A \subseteq  B$ ,则 $\alpha$ 是 $\beta$ 的充分条件;

若 $A \subset  B$ ,则 $\left\{  \begin{array}{l} \alpha  \Rightarrow  \beta , \\  \beta  \nRightarrow  \alpha , \end{array}\right.$ 即 $\alpha$ 是 $\beta$ 的充分非必要条件;

若 $B \subseteq  A$ ,则 $\alpha$ 是 $\beta$ 的必要条件;

\begin{center}
\includegraphics[width=0.2\textwidth]{images/bo_d7tk7oc91nqc73875fu0_1_1166_1790_304_248_0.jpg}
\end{center}
 

若 $B \subset  A$ ,则 $\left\{  \begin{array}{l} \alpha  \Leftarrow  \beta , \\  \alpha  \nRightarrow  \beta , \end{array}\right.$ 即 $\alpha$ 是 $\beta$ 的必要非充分条件;

若 $A = B$ ,则 $\alpha$ 是 $\beta$ 的充要条件.

6. 一般地, 充要条件的范围具有如右图的包含关系.

7. 集合的运算是“交、并、补”; 命题的运算是“或、且、非”，两者都是逻辑运算.

\section*{$\text{ § }{1.1}$ 集合及其运算}

\section*{典型例题精讲}

【例 1】设集合 $A = \{ x \mid   - 2 \leq  x \leq  5\} ,B = \{ x \mid  m + 1 \leq  x \leq  {2m} - 1\}$ .

(1)若 $A \cap  B = B$ ，求实数 $m$ 的取值范围；

(2)当 $x \in  \mathbf{R}$ 时，不存在 $x$ 使得 $x \in  A$ 与 $x \in  B$ 同时成立，求实数 $m$ 的取值范围.

【例 2】已知 $a\text{ 、 }m \in  \mathbf{R}$ ,集合 $A = \left\{  {x \mid  {x}^{2} - {3x} + 2 = 0}\right\}  ,B\left\{  {x \mid  {x}^{2} + 2\left( {m + 1}\right) x + {m}^{2} - 5 = 0,x \in  \mathbf{R}}\right\}  ,C =  \left\{  {x \mid  {x}^{2} - {ax} + a - 1 = 0,x \in  \mathbf{R}}\right\}  .$

(1)若 $A \cup  C = A$ ，求实数 $a$ 的值；

(2)若 $A \cap  B = B$ ，求实数 $m$ 的取值范围.

【例 3】设 $f\left( x\right)  = {x}^{2} + {ax} + b\left( {a\text{ 、 }b \in  \mathbf{R}}\right) ,A = \{ x \mid  f\left( x\right)  = x\} ,B = \left\{  {x\left| {f\left\lbrack  {f\left( x\right) }\right\rbrack   = x}\right. }\right\}$ .

(1)证明 $A \subseteq  B$ ；

(2)若 $A = \{  - 1,3\}$ ，求集合 $B$ .

【例 4】已知实数集 $A = \left\{  {{a}_{1},{a}_{2},\cdots ,{a}_{n}}\right\}  \left( {n \in  \mathbf{N},n \geq  3}\right)$ ，定义 $\varphi \left( A\right)  = \left\{  {{a}_{i}{a}_{j} \mid  {a}_{i}\text{ 、 }{a}_{j} \in  A\left( {i \neq  j}\right) }\right.$ .

(1)若 $A = \{  - 2,0,1,2\}$ ，求 $\varphi \left( A\right)$ ；

(2)若 $\varphi \left( A\right)  = \{ 0, - 6, - 8, - {12},{12},{18},{24}\}$ ，求集合 $A$ .

【例 5】已知函数 $f\left( x\right)  = {x}^{2} - {2x} - {a}^{2} + {2a}\left( {a \in  \mathbf{R}}\right)$ ,集合 $A = \{ x \mid  f\left( x\right)  \leq  0\}$ .

(1)若集合 $A$ 中有且仅有 3 个整数，求实数 $a$ 的取值范围；

(2)集合 $B = \{ x \mid  f\left( {f\left( x\right)  + b}\right)  \leq  0\}$ ，若存在实数 $a \leq  1$ ，使得 $A \subseteq  B$ ，求实数 $b$ 的取值范围.

\section*{四基达标精练}

\section*{一、填空题}

1. 已知集合 $A = \{ 1,2,4\}$ ，集合 $B = \{ 2,4,5\}$ ，则 $A \cap  B =$ \_\_\_\_\_.

2. 集合 $A = \{ \left( {x,y}\right)  \mid  x + y = 4$ 且 $x - y = 2\}$ ，用列举法表示为\_\_\_\_\_.

3. 已知集合 $A = \{ x \mid  x \leq  1\} ,B = \{ x \mid  x \geq  a\}$ ，且 $A \cup  B = \mathbf{R}$ ，则实数 $a$ 的取值范围是\_\_\_\_\_.

\begin{center}
\includegraphics[width=0.2\textwidth]{images/bo_d7tk7oc91nqc73875fu0_4_1168_465_271_175_0.jpg}
\end{center}
 

(第 5 题图)

4. 全集 $U = \{ 2,3,5\} ,A = \{ \left| {a - 5}\right| ,2\} ,\overline{A} = \{ 5\}$ ,则 $a =$ \_\_\_\_\_.

5. 已知全集 $U = \mathbf{R}$ ,集合 $A = \{ x \mid  x <  - 1$ 或 $x > 4),B = \{ x \mid   - 2 \leq  x \leq  3)$ ,那么阴影部分表示的集合为\_\_\_\_\_.

6. 已知 $M = \{ \left( {x,y}\right)  \mid  {2x} + {3y} = {4320},x\text{ 、 }y \in  \mathbf{N}\} ,P = \{ \left( {x,y}\right)  \mid  {4x} - {3y} = 1$ , $x$ 、 $y \in  \mathbf{N}\}$ ，则其中的有限集是\_\_\_\_\_.

7. 已知集合 $A = \left\{  {x\left| {a{x}^{2} - {3x} + 2 = 0}\right. }\right\}$ 中有且仅有一个元素，则 $a$ 的取值组成的集合为\_\_\_\_\_.

8. 集合 $A = \{ \left( {x,y}\right) \left| {y = a}\right| x \mid  \} ,B = \{ \left( {x,y}\right)  \mid  y = x + 1\}$ ,若 $A \cap  B$ 为单元素集，实数 $a$ 的取值范围为\_\_\_\_\_.

9. 设数集 $M = \left\{  {x\left| {m \leq  x \leq  m + \frac{3}{4}}\right. }\right\}  ,N = \left\{  {x\left| {n - \frac{1}{3} \leq  x \leq  n}\right. }\right\}$ ,且 $M\text{ 、 }N$ 都是集合 $\{ x \mid  0 \leq  x \leq  1\}$ 的子集,如果把 $b - a$ 叫做集合 $\{ x \mid  a \leq  x \leq  b\}$ 的“长度”,那么集合 $M \cap  N$ 的长度的最小值是\_\_\_\_\_.

\section*{二、选择题}

10. 设集合 $P = \{ x \mid  0 \leq  x \leq  \sqrt{2}\} ,m = \sqrt{3}$ ，则下列关系中正确的是 ( )

A. $m \subseteq  P$ . B. $m \text{ ⊄ } P$ . C. $m \in  P$ . D. $m \notin  P$ .

11. 由 ${a}^{2} + 1\text{ 、 }a + 3\text{ 、 }a$ 组成的集合含有元素 2,则实数 $a$ 的可能取值的个数是 ( )

A. 1 个. B. 2 个. C. 3 个. D. 4 个.

12. 设 $a$ 、 $b$ 、 $c \in  \mathbf{R}$ ,函数 $f\left( x\right)  = \left( {x + a}\right) \left( {{x}^{2} + {bx} + c}\right) ,g\left( x\right)  = \left( {{ax} + 1}\right) \left( {c{x}^{2} + {bx} + 1}\right)$ . 记集合 $S = \{ x \mid  f\left( x\right)  = 0$ , $x \in  \mathbf{R}\} ,T = \{ x \mid  g\left( x\right)  = 0,x \in  \mathbf{R}\}$ . 若 $\left| S\right| \text{ 、 }\left| T\right|$ 分别为集合 $S\text{ 、 }T$ 的元素个数,则下列结论不可能的是 ( )

A. $\left| S\right|  = 1,\left| T\right|  = 0$ . B. $\left| S\right|  = 1,\left| T\right|  = 1$ . C. $\left| S\right|  = 2,\left| T\right|  = 2$ . D. $\left| S\right|  = 2,\left| T\right|  = 3$ .

\section*{三、解答题}

13. 已知集合 $A = \{ x \mid  x \geq  3\} ,B = \{ x \mid  1 \leq  x \leq  7\} ,C = \{ x \mid  x \geq  a - 1\}$ .

(1)求 $A \cap  B$ 和 $A \cup  B$ ；

(2)若 $C \cup  A = A$ ，求实数 $a$ 的取值范围.

14. 已知 $a\text{ 、 }x \in  \mathbf{R}$ ,集合 $A = \{ 2,4,{x}^{2} - {5x} + 9\} ,B = \left\{  {3,{x}^{2} + {ax} + a}\right\}  ,C = \left\{  {{x}^{2} + \left( {a + 1}\right) x - 3,1}\right\}$ .

(1)若 $2 \in  B$ ，且 $B \subseteq  A$ ，求 $a$ 和 $x$ 的值；

(2)若集合 $B = C$ ，求 $a$ 和 $x$ 的值.

15. 已知全集是 $\mathbf{R}$ ,函数 $f\left( x\right)  = \sqrt{x - 3} - \frac{1}{\sqrt{7 - x}}$ 的定义域为集合 $A,B = \{ x \mid  2 < x < {10},x \in  \mathbf{Z}\} ,C = \{ x \mid  x < a$ 或 $x > a + 1\}$ .

(1)求 $A$ 和 $\overline{A} \cap  B$ ；

(2)若 $A \cup  C = \mathbf{R}$ ，求实数 $a$ 的取值范围.

16. 设 $A$ 为集合 $M$ 的子集，且 $A = \left\{  {{a}_{1},{a}_{2},\cdots ,{a}_{n}}\right\}  \left( {n \in  \mathbf{N},n \geq  2}\right)$ ，若 ${a}_{1} + {a}_{2} + \cdots  + {a}_{n} = {a}_{1} \cdot  {a}_{2} \cdot  \cdots  \cdot  {a}_{n}$ ，则， 称 $A$ 为集合 $M$ 的 $n$ 元“好集”.

(1)写出实数集 $\mathbf{R}$ 的一个二元“好集”；

(2)求出正整数集的所有三元“好集”.

\section*{$\text{ § }{1.2}$ 充要条件与反证法}

\section*{典型例题精讲}

【例 1】已知集合 $A = \left\{  {y\left| {y = {x}^{2} - \frac{3}{2}x + 1}\right. ,x \in  \left\lbrack  {\frac{3}{4},2}\right\rbrack  }\right\}  ,B = \left\{  {x\left| {x + {m}^{2} \geq  1}\right. }\right\}$ .

(1)求集合 $A$ ；

(2)若 $p : x \in  A;q : x \in  B$ 且 $p$ 是 $q$ 的充分条件，求实数 $m$ 的取值范围.

【例 2】设 $p$ : 实数 $x$ 满足 ${x}^{2} - {4ax} + 3{a}^{2} \leq  0\left( {a > 0}\right) ,q$ : 实数 $x$ 满足 $\frac{x - 3}{x - 2} < 0$ .

(1)若 $a = 1,p$ 且 $q$ 是真命题，求实数 $x$ 的取值范围；

(2)若 $p$ 是 $q$ 的必要非充分条件，求实数 $a$ 的取值范围.

【例 3】已知函数 $f\left( x\right)  =  - {x}^{2} + {bx} + c$ ,则 “ $f\left( {f\left( \frac{b}{2}\right) }\right)  > 0$ ” 是 “方程 $f\left( x\right)  = 0$ 有两个不同实数解且方程 $f\left( {f\left( x\right) }\right)  = 0$ 恰有两个不同实数解”的 ( )

A. 充要条件. B. 充分非必要条件. C. 必要不充分条件. D. 既非充分也非必要条件.

【例 4】已知定义域为 $D$ 的函数 $f\left( x\right)$ ,若存在实数 $a$ ,使得对任意 ${x}_{1} \in  D$ ,都存在 ${x}_{2} \in  D$ 满足 $\frac{{x}_{1} + f\left( {x}_{2}\right) }{2} = a$ ,则称函数 $f\left( x\right)$ 具有性质 $P\left( a\right)$ .

(1)判断下列函数 $f\left( x\right)  = {2}^{x}$ 和 $g\left( x\right)  = {\log }_{2}x\left( {0 < x < 1}\right)$ 是否具有性质 $P\left( 0\right)$ ，说明理由；

(2)若函数 $f\left( x\right)$ 的定义域为 $D$ ，且具有性质 $P\left( 1\right)$ ，则 “ $f\left( x\right)$ 存在零点” 是 “ $2 \in  D$ ” 的什么条件，说明理由.

【例 5】证明下列命题是真命题:

(1)已知 $n \in  \mathbf{Z}$ . 用反证法证明:若 ${n}^{3}$ 是奇数，则 $n$ 是奇数；

(2)已知 $a$ 、 $b \in  \left( {0, + \infty }\right)$ ， $a + b > 2$ . 证明: $\frac{1 + a}{b}$ 与 $\frac{1 + b}{a}$ 中至少有一个小于 2 .

\section*{四基达标精练}

\section*{一、填空题}

1. 命题:“若 $a < b$ ，则 $a{c}^{2} < b{c}^{2}$ ”是\_\_\_\_\_命题(填“真”或“假”).

2. 设 $m\text{ 、 }n$ 是整数，则 “ $m\text{ 、 }n$ 均为偶数” 的否定是\_\_\_\_\_.

3. 已知 “ $p : x > 1$ ” 是 “ $q : {x}^{2} > x$ ”，则 $p\text{ 、 }q$ 的推出关系用 “ $\Rightarrow$ 、 $\Leftarrow$ 、 $\Leftrightarrow$ ”中的一个表示为\_\_\_\_\_.

4. 一元二次方程 $a{x}^{2} + {bx} + c = 0$ 有一正一负根的充要条件是\_\_\_\_\_.

5. 否定命题 “若 $a > b$ ，则 $\frac{1}{a} < \frac{1}{b}$ ” 的一个反例可以是\_\_\_\_\_.

6. 若 “ $a$ 满足 $\frac{{2a} - 1}{a + 1} \leq  1$ ” 为真命题，则实数 $a$ 的取值范围是\_\_\_\_\_.

7. 设条件 $p : \left| x\right|  \leq  m\left( {m > 0}\right) ,q :  - 1 \leq  x \leq  4$ ，若 $p$ 是 $q$ 的充分条件，则 $m$ 的最大值为\_\_\_\_\_.

8. 已知 $A = \{ x \mid  \left| {x - a}\right|  < 4\} ,B = \{ x \mid  \left( {x - 1}\right) \left( {2 - x}\right)  > 0\}$ ,若 $x \in  A$ 是 $x \in  B$ 的必要非充分条件,则实数 $a$ 的取值范围是\_\_\_\_\_.

9. 已知命题 $p :$ 对任意 $x \in  \left\lbrack  {1,2}\right\rbrack$ ， ${x}^{2} - a \geq  0$ 都成立；命题 $q :$ 存在 $x \in  \mathbf{R}$ ，使得成立 ${x}^{2} + {2ax} + 2 - a = 0$ ， 若命题“ $p$ 且 $q$ ”是真命题，则实数 $a$ 的取值范围为\_\_\_\_\_.

\section*{二、选择题}

10. 设 $a\text{ 、 }b\text{ 、 }c \in  {\mathbf{R}}^{n}{b}^{2} - {4ac} > 0$ ” 是 “方程 $a{x}^{2} + {bx} + c = 0$ 有实根” 的 ( )

A. 充分非必要条件. B. 必要非充分条件. C. 充要条件. D. 既非充分也非必要条件.

11. " $x + y \neq  3$ " 是 " $x \neq  1$ 或 $y \neq  2$ "的 ( )

A. 充分非必要条件. B. 必要非充分条件. C. 充要条件. D. 既非充分也非必要条件.

12. ” ${\sin }^{2}\alpha  + {\cos }^{2}\beta  = 1$ ” 是 “ $\alpha  = \beta$ ” 的充要条件,则 $\alpha$ 、 $\beta$ 满足的条件是 ( )

A. $\alpha  = \beta$ . B. $\alpha  - \beta  = {2k\pi }\left( {k \in  \mathbf{Z}}\right)$ . C. $\alpha  - \beta  = {k\pi }\left( {k \in  \mathbf{Z}}\right)$ . D. $\alpha  - \beta  = \frac{k\pi }{2}\left( {k \in  \mathbf{Z}}\right)$ .

\section*{三、解答题}

13. 已知命题 $p$ : 关于 $x$ 的方程 ${x}^{2} + \left( {m - 2}\right) x + 1 = 0$ 无实根; 命题 $q$ : 关于 $x$ 的方程 ${x}^{2} + {mx} + 1 = 0$ 有两个不等的负实根.

(1)若命题 $p$ 为假命题，求实数 $m$ 的取值范围；

(2)若命题 $p$ 、 $q$ 中有且仅有一个是真命题，求实数 $m$ 的取值范围.

14. 证明下列命题是真命题:

(1)已知 $x$ 为正数， $a =  - x + \frac{10}{x}$ ， $b = {5x} - \frac{1}{x}$ ，则 $a$ 、 $b$ 中至少有一个不小于 6；

(2)已知实数 $a$ 、 $b$ 满足 ${a}^{3} + {b}^{3} = 2$ ，则 $a + b \leq  2$ .

15. 设集合 $A = \left\{  {x \mid  {x}^{2} + \left( {p + 2}\right) x + 1 = 0,x \in  \mathbf{R}}\right\}$ .

(1)求 $A = \varnothing$ 的充要条件；

(2)求 $A \cap  {\mathbf{R}}^{ + } = \varnothing$ 的充要条件，其中 ${\mathbf{R}}^{ + }$ 表示正实数集.

16. 设 $a$ 、 $b$ 、 $c$ 为 $\bigtriangleup {ABC}$ 的三条边长,方程 ${x}^{2} + {2ax} + {b}^{2} = 0$ 与 ${x}^{2} + {2cx} - {b}^{2} = 0$ .

(1)求方程 ${x}^{2} + {2ax} + {b}^{2} = 0$ 有实数根的充要条件；

(2)求这两个方程有公共根的充要条件.

\section*{第 2 章 等式与不等式(必修 1)}

\section*{知识梳理*}

1. 不等式的依据:实数的有序性，即 $a > b \Leftrightarrow  a - b > 0;a < b \Leftrightarrow  a - b < 0;a = b \Leftrightarrow  a - b = 0$ ；

2. 等式与不等式的性质:

\begin{center}
\begin{tabular}{|c|c|c|}
\hline
性质 & 等式 & 不等式 \\
\cline{1-3}
意义 & $a = b \Leftrightarrow  a - b = 0$ & $a > b \Leftrightarrow  a - b > 0;a < b \Leftrightarrow  a - b < 0$ \\
\cline{1-3}
传递性 & $a = b,b = c \Leftrightarrow  a = c$ & $a > b,b > c \Rightarrow  a > c$ \\
\cline{1-3}
加法性质 & $a = b \Leftrightarrow  a + c = b + c$ & $a > b \Leftrightarrow  a + c > b + c$ \\
\cline{1-3}
推论 & $a = b,c = d \Rightarrow  a + c = b + d$ & $a > b,c > d \Rightarrow  a + c > b + d$ \\
\cline{1-3}
推广 & ${a}_{i} = {b}_{i}\left( {i \in  \mathbf{N},i \geq  2}\right)  \Rightarrow  \mathop{\sum }\limits_{{i = 1}}^{n}{a}_{i} = \mathop{\sum }\limits_{{i = 1}}^{n}{b}_{i}$ & ${a}_{i} > {b}_{i}\left( {i \in  \mathbf{N},i \geq  2}\right)  \Rightarrow  \mathop{\sum }\limits_{{i = 1}}^{n}{a}_{i} > \mathop{\sum }\limits_{{i = 1}}^{n}{b}_{i}$ \\
\cline{1-3}
乘法性质 & $a = b \Rightarrow  {ac} = {bc}$ & $a > b,c > 0 \Rightarrow  {ac} > {bc};a > b,c < 0 \Rightarrow  {ac} < {bc}$ \\
\cline{1-3}
推论 & $a = b,c = d \Rightarrow  {ac} = {bd}$ & $a > b > 0,c > d > 0 \Rightarrow  {ac} > {bd}$ \\
\cline{1-3}
推广 & \makecell{ ${a}_{i} = {b}_{i}\left( {i \in  \mathbf{N},i \geq  2}\right)$ \\ $\Rightarrow  {a}_{1}{a}_{2}\cdots {a}_{n} = {b}_{1}{b}_{2}\cdots {b}_{n}$ } & \makecell{ ${a}_{i} > {b}_{i} > 0\left( {i \in  \mathbf{N},i \geq  2}\right)$ \\ $\Rightarrow  {a}_{1}{a}_{2}\cdots {a}_{n} > {b}_{1}{b}_{2}\cdots {b}_{n}$ } \\
\cline{1-3}
乘方性质 & $a = b,n \in  \mathbf{N},n \geq  1 \Rightarrow  {a}^{n} = {b}^{n}$ & $a > b > 0,n \in  \mathbf{N},n \geq  1 \Rightarrow  {a}^{n} > {b}^{n} > 0$ \\
\cline{1-3}
开方性质 & $a = b > 0,n \in  \mathbf{N},n \geq  1 \Rightarrow  \sqrt[n]{a} = \sqrt[n]{b} > 0$ & $a > b > 0,n \in  \mathbf{N},n \geq  2 \Rightarrow  \sqrt[n]{a} > \sqrt[n]{b} > 0$ \\
\cline{1-3}
倒数性质 & $a = b \neq  0 \Rightarrow  \frac{1}{a} = \frac{1}{b}$ & $a > b,{ab} > 0 \Rightarrow  \frac{1}{a} < \frac{1}{b}$ \\
\cline{1-3}
\hline
\end{tabular}
\end{center}

3. 方程 (一元二次方程) 与方程的解集:

 \ding{172}  方程的解与数域有关；

 \ding{173} 一元二次方程的解集与韦达定理:

若一元二次方程 $a{x}^{2} + {bx} + c = 0\left( {a \neq  0}\right)$ 的解为 ${x}_{1}\text{ 、 }{x}_{2}$ ,则 ${x}_{1}{x}_{2} =  - \frac{b}{a},{x}_{1}{x}_{2} = \frac{c}{a}$ .

4. 不等式的解法:高次不等式低次化；分式不等式整式化；绝对值不等式常规化；无理不等式有理化；含参数的不等式分类讨论; 数形结合、换元法、根据函数的单调性等是转化的重要途径.

5. 均值不等式:

 \ding{172} 常用不等式: $a$ 、 $b \in  \mathbf{R} \Rightarrow  {a}^{2} + {b}^{2} \geq  {2ab}$ (当且仅当 $a = b$ 时，等号成立)；

 \ding{173} 平均值不等式: $a \geq  0,b \geq  0 \Rightarrow  \frac{a + b}{2} \geq  \sqrt{ab}$ (当且仅当 $a = b$ 时，等号成立)；

\HRule

* 注:本章节所属知识模块是“预备知识”.

\HRule

 \ding{174}  $a\text{ 、 }b \in  \mathbf{R} \Rightarrow  {\left( \frac{a + b}{2}\right) }^{2} \geq  {ab}$ (当且仅当 $a = b$ 时，等号成立)；

 \ding{175}  $a\text{ 、 }b \in  \mathbf{R} \Rightarrow  \frac{{a}^{2} + {b}^{2}}{2} \geq  {\left( \frac{a + b}{2}\right) }^{2} \geq  {ab}$ (当且仅当 $a = b$ 时，等号成立)；

 \circled{5}  ${ab} > 0 \Rightarrow  \frac{b}{a} + \frac{a}{b} \geq  2$ (当且仅当 $a = b$ 时，等号成立)； ${ab} < 0 \Rightarrow  \frac{b}{a} + \frac{a}{b} \leq   - 2$ (当且仅当 $a =  - b$ 时，等号成立).

注意: ${a}^{2} + {b}^{2}\text{ 、 }a + b\text{ 、 }a - b\text{ 、 }{ab}$ 之间的关系.

重要结论: 设 $x\text{ 、 }y$ 为正数,

 \ding{172} 如果积 ${xy}$ 为定值 $P$ ，那么当 $x = y$ 时，和 $x + y$ 有最小值 $2\sqrt{P}$ ；

 \ding{173} 如果和 $x + y$ 为定值 $S$ ，那么当 $x = y$ 时，积 $x \cdot  y$ 有最大值 $\frac{1}{4}{S}^{2}$ .

注意:利用平均值不等式求最值时要注意“一正、二定、三取等”.

6. 三角不等式:

 \ding{172}  $\left| {a + b}\right|  \leq  \left| a\right|  + \left| b\right|$ (当且仅当 ${ab} \geq  0$ 时，等号成立)

 \ding{173}  $\left| {a - b}\right|  \leq  \left| a\right|  + \left| b\right|$ (当且仅当 ${ab} \leq  0$ 时，等号成立)

 \ding{174}  $\left| \right| a\left| -\right| b\left| \right|  \leq  \left| {a + b}\right|$ (当且仅当 ${ab} \leq  0$ 时，等号成立)

 \ding{175}  $\left| \right| a\left| -\right| b\left| \right|  \leq  \left| {a - b}\right|$ (当且仅当 ${ab} \geq  0$ 时，等号成立)

即 $a\text{ 、 }b \in  \mathbf{R} \Rightarrow  \left| \right| a\left| -\right| b\left| \right|  \leq  \left| {a \pm  b}\right|  \leq  \left| a\right|  + \left| b\right|$ .

7. 不等式的证明:

有比较法(比差、比商)、综合法、反证法、分析法等.

(1)比较法:作差 $\rightarrow$ 变形 (配方或分解因式) $\rightarrow$ 判断符号.

两边均正时也可比商，判断与 1 的大小.

(2)综合法:从已知出发，依据性质和基本不等式，直接推出所需的不等式(连接词:: ···，，....).

(3)反证法:首先假设结论 $\beta$ 不成立( $\beta$ 为假)，然后经过正确的逻辑推理得出与已知条件或(已学)定理等相矛盾的结论,从而说明“ $\beta$ 为假”是不可能发生的,即结论 $\beta$ 是正确的. 这样的证明方法叫反证法.

(4)分析法:从欲证的不等式出发，寻找不等式成立的充分条件 (连接词:要证...，只要证...).

除此以外, 还有放缩法、数学归纳法等.

8. 不等式的恒成立和有解问题:

(1) $a{x}^{2} + {bx} + c > \left( { \geq  0}\right)$ 0 恒成立 $\Leftrightarrow  \left\{  \begin{array}{l} a > 0, \\  \Delta  < \left(  \leq  \right) 0 \end{array}\right.$ 或 $\left\{  \begin{array}{l} a = 0,b = 0, \\  c > \left(  \geq  \right) 0; \end{array}\right.$

(2)若 $m \leq  f\left( x\right)  \leq  M$ ，则 $a > f\left( x\right)$ 恒成立 $\Leftrightarrow  a > M;a > f\left( x\right)$ 有解 $\Leftrightarrow  a > m;a \geq  f\left( x\right)$ 恒成立 $\Leftrightarrow  a \geq  M$ ； $a \geq  f\left( x\right)$ 有解 $\Leftrightarrow  a \geq  m$ .

※若 $m < f\left( x\right)  < M$ ( $m$ 、 $M$ 分别是 $f\left( x\right)$ 的最大下界和最小上界),则 $a > f\left( x\right)$ 恒成立 $\Leftrightarrow  a \geq  M;a >  f\left( x\right)$ 有解 $\Leftrightarrow  a > m;a \geq  f\left( x\right)$ 恒成立 $\Leftrightarrow  a \geq  M;a \geq  f\left( x\right)$ 有解 $\Leftrightarrow  a > m$ . 对于 $a < f\left( x\right)$ 和 $a \leq  f\left( x\right)$ 类似考虑.

注意区别 “恒成立” 和 “有解”: “恒成立” 中的 $x$ 具有任意性,“有解” 中的 $x$ 只要具有存在性.

9. 常用解题策略:认清变量和参数，恰当分离参数，灵活转化问题是解决不等式问题的重要保证.

\section*{$\text{ § }{2.1}$ 等式与方程}

\section*{典型例题精讲}

【例 1】已知全集 $U = \{ 1,2,3,4,5\}$ ,集合 $A = \left\{  {x \mid  {x}^{2} - {5x} + m = 0}\right\}  ,B = \left\{  {x \mid  {x}^{2} + {nx} + {12} = 0}\right\}$ ,且 $\bar{A} \cup  B  = \{ 1,3,4,5\}$ ,求 $m + n$ 的值.

【例 2】已知一元二次函数 $f\left( x\right)  = a{x}^{2} + {bx} + c\left( {a > 0,c > 0}\right)$ 的图像与 $x$ 轴有两个不同的公共点,其中一个公共点的坐标为 $\left( {c,0}\right)$ ,且当 $0 < x < c$ 时,恒有 $f\left( x\right)  > 0$ .

(1)当 $a = 1\text{ 、 }c = \frac{1}{2}$ 时，求出不等式 $f\left( x\right)  < 0$ 的解;

(2)求不等式 $f\left( x\right)  < 0$ 的解 (用 $a$ 、 $c$ 表示)；

(3)若不等式 ${m}^{2} - {2km} + 1 + b + {ac} \geq  0$ 对所有 $k \in  \left\lbrack  {-1,1}\right\rbrack$ 恒成立，求实数 $m$ 的取值范围.

【例 3】已知二次函数 $f\left( x\right)  = a{x}^{2} + {bx}$ ( $a$ 、 $b$ 为常数，且 $a \neq  0$ )满足条件: $f\left( {x - 1}\right)  = f\left( {3 - x}\right)$ ，且方程 $f\left( x\right)  = {2x}$ 有两等根.

(1)求 $f\left( x\right)$ 的解析式；

(2)是否存在实数 $m$ 、 $n$ ( $m < n$ )，使 $f\left( x\right)$ 的定义域和值域分别为 $\left\lbrack  {m,n}\right\rbrack$ 和 $\left\lbrack  {{4m},{4n}}\right\rbrack$ ？如果存在，求出 $m\text{ 、 }n$ 的值；如果不存在，说明理由.

【例 4】函数 $y = f\left( x\right)$ 的定义域为 $D$ ,若存在 ${x}_{0} \in  D$ ,使得 $f\left( {x}_{0}\right)  = {x}_{0}$ 成立,则称 ${x}_{0}$ 为函数 $y = f\left( x\right)$ 的“不动点”.

(1)若 $f\left( x\right)  = a{x}^{2} + {bx} + c\left( {a \neq  0}\right)$ 有两个不动点 -1 和 3，求 ${bc}$ 的最小值；

(2)若 $a > 0$ ，且 $f\left( x\right)  = a{x}^{2} + {bx} + c$ 有两个不动点 $m$ 、 $n$ ，且满足: $0 < m < n < \frac{1}{a}$ ，求证:当 $x \in  \left( {0,m}\right)$ 时， $f\left( x\right)  < m$ .

【例 5】已知 $a > b > c,a + b + c = 0$ ,方程 $a{x}^{2} + {bx} + c = 0$ 的两个实根为 ${x}_{1}\text{ 、 }{x}_{2}$ .

(1)证明: $- \frac{1}{2} < \frac{b}{a} < 1$ ；

( 2 )若 ${x}_{1}^{2} + {x}_{1}{x}_{2} + {x}_{2}^{2} = 1$ ，求 ${x}_{1}^{2} - {x}_{1}{x}_{2} + {x}_{2}^{2}$ 的值；

(3)求 $\left| {{x}_{1}^{2} - {x}_{2}^{2}}\right|$ 的取值范围.

\section*{四基达标精练}

\section*{一、填空题}

1. 已知一元二次方程 $a{x}^{2} + {bx} + 2 = 0$ 的解是 -2 $\frac{1}{2}$ 和 $\frac{1}{3}$ ，则 $a + b =$ \_\_\_\_\_.

2. 若关于 $x$ 的不等式 $a{x}^{2} + {6x} - {a}^{2} < 0$ 的解集是 $\left( {-\infty ,1}\right)  \cup  \left( {m, + \infty }\right)$ ，则实数 $m =$ \_\_\_\_\_.

3. 已知实数 $a$ 、 $b$ 满足 ${\left( a + 2\right) }^{2} + {\left( {b}^{2} - 2b - 3\right) }^{2} = 0$ ，则 $a + b$ 的值为\_\_\_\_\_.

4. 已知方程 ${x}^{2} + {2mx} + m + 2 = 0$ 有两个不等正根，则实数 $m$ 的取值范围是\_\_\_\_\_.

5. 关于 $x$ 的方程 ${x}^{2} - {2ax} - 8{a}^{2} = 0$ 的解为 ${x}_{1}$ 和 ${x}_{2}$ ，且 ${x}_{1}^{2} - {x}_{2}^{2} = {15}$ ，则实数 $a =$ \_\_\_\_\_.

6. 二次函数 $y = a{x}^{2} + {bx} + c\left( {x \in  \mathbf{R}}\right)$ 的部分对应值如表:

\begin{center}
\begin{tabular}{|c|c|c|c|c|c|c|c|c|}
\hline
$x$ & -3 & -2 & -1 & 0 & 1 & 2 & 3 & 4 \\
\cline{1-9}
$y$ & 6 & 0 & -4 & -6 & -6 & -4 & 0 & 6 \\
\cline{1-9}
\hline
\end{tabular}
\end{center}

则不等式 $a{x}^{2} + {bx} + c > 0$ 的解集是\_\_\_\_\_.

7. 设 ${x}_{1}$ 、 ${x}_{2}$ 是关于 $x$ 的方程 $2{x}^{2} - {4mx} + 2{m}^{2} + {3m} - 2 = 0$ 的两个实数根，则 ${x}_{1}^{2} + {x}_{2}^{2}$ 的最小值为 .

8. 在 $\mathbf{R}$ 上定义运算 $\odot   : a \odot  b = {ab} + {2a} + b$ ，则不等式 $x \odot  \left( {x - 2}\right)  < 0$ 的解集是\_\_\_\_\_.

9. 已知 $\alpha$ 、 $\beta \left( {\alpha  < \beta }\right)$ 是方程 $\left( {x - a}\right) \left( {x - b}\right)  =  - 2\left( {a < b}\right)$ 的两个根，则 $\alpha$ 、 $\beta$ 、 $a$ 、 $b$ 的大小关系是\_\_\_\_\_.

\section*{二、选择题}

10. “ $- 2 < a < 2$ ” 是 “不等式 $\left( {a - 2}\right) {x}^{2} + 2\left( {a - 2}\right) x - 4 < 0$ 的解集为 $\mathbf{R}$ ” 的 ( )

A. 充分非必要条件. B. 必要非充分条件. C. 充要条件. D. 既非充分也非必要条件.


\begin{center}
\includegraphics[width=0.2\textwidth]{images/bo_d7tk7oc91nqc73875fu0_15_1268_157_163_345_0.jpg}
\end{center}
 

(第 12 题图)

11. 已知关于 $x$ 的不等式 ${x}^{2} - \left( {a + 1}\right) x + a < 0$ 的解集中恰有两个正整数,则实数 $a$ 的取值范围是 ( )

A. $\lbrack 2,4)$ . B. $\left\lbrack  {3,4}\right\rbrack$ . C. $(3,4\rbrack$ . D. $\left( {3,4}\right)$ .

12. 古希腊时期, 人们认为最美人体的头顶至肚脐的长度与肚脐至足底的长度之比是 $\frac{\sqrt{5} - 1}{2}\left( {\frac{\sqrt{5} - 1}{2} \approx  {0.618}}\right.$ ,称为黄金分割比例 $)$ ,著名的“断臂维纳斯”便是如此. 此外,最美人体的头顶至咽喉的长度与咽喉至肚脐的长度之比也是 $\frac{\sqrt{5} - 1}{2}$ . 若某人满足上述两个黄金分割比例, 且腿长为 ${105}{{cm}}$ ，头顶至脖子下端的长度为 ${26}{{cm}}$ ，则其身高可能是 ( )

A. ${165}{{cm}}$ . B. ${175}{{cm}}$ . C. ${185}{{cm}}$ . D. ${190}{{cm}}$ .

\section*{三、解答题}

13. 已知 $f\left( x\right)  = {x}^{2} + {mx} - 6\left( {m > 0}\right)$ ，实数 ${x}_{1}$ 、 ${x}_{2}$ 满足 $f\left( {x}_{1}\right)  = f\left( {x}_{2}\right)  = 0$ ，且 ${x}_{2} - {x}_{1} = 5$ .

(1)求函数 $y = f\left( x\right)$ 的解析式；

(2)解关于 $x$ 的不等式 $y < 4 - {2x}$ .

14. 已知集合 $\{ x \mid  {x}^{2} - \left( {k + 2}\right) x - {3k} + 1 \geq  0\}  = \left( {-\infty , - 1\rbrack \cup \lbrack 5, + \infty }\right)$ .

(1)求实数 $k$ 的值；

(2)已知 $t \in  \left( {-\infty ,2}\right)$ ，若不等式 ${x}^{2} - \left( {k + 2}\right) x - {3k} - {m}^{2} + {4m} + {15} \geq  0$ 在 $t \leq  x \leq  4$ 上恒成立，求实数 $m$ 的取值范围.

15. 已知函数 $f\left( x\right)  = {x}^{2} + {2x} + 1$ .

(1)求关于 $x$ 的不等式 $f\left( x\right)  \geq  {b}^{2}\left( {b \geq  0}\right)$ 的解集；

(2)若不等式 ${\left\lbrack  f\left( x\right) \right\rbrack  }^{2} - {2mf}\left( x\right)  + {m}^{2} - 1 \geq  0$ 对于任意 $x \in  \left\lbrack  {-3,2}\right\rbrack$ 都成立，求 $m$ 的取值范围.

16. 已知关于 $x$ 的一元二次方程 ${x}^{2} - {6x} + {2m} - 1 = 0$ 有两实数根 ${x}_{1}$ 、 ${x}_{2}$ .

(1)若 ${x}_{1} = 1$ ，求 ${x}_{2}$ 和 $m$ 的值；

(2)是否存在实数 $m$ ，满足 $\left( {{x}_{1} - 1}\right) \left( {{x}_{2} - 1}\right)  = \frac{6}{m - 5}$ ？若存在，求出实数 $m$ 的值；若不存在，请说明理由.

\section*{$\text{ § }{2.2}$ 不等式的性质与解法}

\section*{典型例题精讲}

【例 1】若不等式 $f\left( x\right)  \geq  0$ 的解集是 $\left\lbrack  {-1,2}\right\rbrack$ ,不等式 $g\left( x\right)  \geq  0$ 的解集为空集,且函数 $f\left( x\right) \text{ 、 }g\left( x\right)$ 的定义域都是 $\mathbf{R}$ ，则不等式 $\frac{f\left( x\right) }{g\left( x\right) } > 0$ 的解集为\_\_\_\_\_.

【例 2】解不等式组 $\left\{  \begin{array}{l} a\left( {x - 2}\right)  + 1 > 0, \\  {\left( x - 1\right) }^{2} > a\left( {x - 2}\right)  + 1, \end{array}\right.$ 其中 $a > 1$ .

【例 3】已知关于 $x$ 的不等式组 $\left\{  \begin{array}{l} {x}^{2} - x - 2 > 0, \\  2{x}^{2} + \left( {{2k} + 5}\right) x + {5k} < 0 \end{array}\right.$ 的解集为 $A,\mathbf{Z}$ 为整数集.

(1)若 $A \cap  \mathbf{Z} = \{  - 2\}$ ，求实数 $k$ 的取值范围；

(2)若 $A \cap  \mathbf{Z}$ 中恰有 $n$ ( $n$ 为正整数)个整数，求实数 $k$ 的取值范围(用 $n$ 表示).

【例 4】已知不等式: ${0x}^{2} - {4x} + 3 < 0,{0x}^{2} - {6x} + 8 < 0,{0x}^{3}{x}^{2} - {9x} + m < 0$ .

(1)若同时满足 \ding{172}  \ding{173} 的 $x$ 也满足 \ding{174} ，求实数 $m$ 的取值范围；

(2)若 \ding{174} 的解集非空，且满足  \ding{174}  的 $x$ 至少满足  \ding{172}  \ding{173}  中的一个，求实数 $m$ 的取值范围.

【例 5】已知关于 $x$ 的不等式 $\left( {{kx} - {k}^{2} - 4}\right) \left( {x - 4}\right)  > 0$ ，其中 $k \in  \mathbf{R}$ .

(1)试求不等式的解集 $A$ ；

(2)对于不等式的解集 $A$ ,若满足 $A \cap  \mathbf{Z} = B$ (其中 $\mathbf{Z}$ 为整数集). 试探究集合 $B$ 能否为有限集? 若能,求出使得集合 $B$ 中元素个数最少的 $k$ 的所有取值,并用列举法表示集合 $B$ ; 若不能,请说明理由.

\section*{四基达标精练}

\section*{一、填空题}

1. 不等式 $- \frac{1}{2} < \frac{1}{x} \leq  \frac{1}{3}$ 的解集是\_\_\_\_\_.

2. 已知 $x \in  \mathbf{R},M = 2{x}^{2} + {5x} + 3,N = {x}^{2} + {4x} + 2$ ，则 $M$ 、 $N$ 的大小关系是\_\_\_\_\_.

3. 已知 $- 1 < a < 2, - 1 < b < 2$ ,则 $a - b$ 的范围是\_\_\_\_\_.

4. 不等式 $- {x}^{2} + x + 6 < 0$ 的解集是\_\_\_\_\_.

5. 已知函数 $f\left( x\right)  = \left\{  \begin{array}{l} x + 2,x \leq  0, \\   - x + 2,x > 0, \end{array}\right.$ 则不等式 $f\left( x\right)  \geq  {x}^{2}$ 的解集为\_\_\_\_\_.

6. 若关于 $x$ 的不等式 ${ax} - b > 0$ 的解集是 $\left( {1, + \infty }\right)$ ,则关于 $x$ 的不等式 $\frac{{ax} + b}{x - 2} > 0$ 的解集是\_\_\_\_\_.

7. 已知不等式 $a{x}^{2} + {bx} + c > 0$ 的解集为 $\left( {-\infty ,\frac{1}{4}}\right)  \cup  \left( {\frac{1}{2}, + \infty }\right)$ ,则不等式 $c{x}^{2} + {bx} + a < 0$ 的解集为\_\_\_\_\_.

8. 若关于 $x$ 的不等式 $1 \leq  k{x}^{2} + {2x} + k \leq  2$ 有唯一实数解,则实数 $k =$ \_\_\_\_\_.

9. 对于问题:“已知关于 $x$ 的不等式 $a{x}^{2} + {bx} + c > 0$ 的解集为 $\left( {-1,2}\right)$ ，解关于 $x$ 的不等式 $a{x}^{2} - {bx} + c > 0$ ”， 给出如下一种解法:

解:由 $a{x}^{2} + {bx} + c > 0$ 的解集为 $\left( {-1,2}\right)$ ,得 $a{\left( -x\right) }^{2} + b\left( {-x}\right)  + c > 0$ 的解集为 $\left( {-2,1}\right)$ ,即关于 $x$ 的不等式 ${ax} - {bx} + c > 0$ 的解集为 $\left( {-2,1}\right)$ .

参考上述解法,若关于 $x$ 的不等式 $\frac{k}{x + a} + \frac{x + b}{x + c} < 0$ 的解集为 $\left( {-1, - \frac{1}{3}}\right)  \cup  \left( {\frac{1}{2},1}\right)$ ,则关于 $x$ 的不等式 $\frac{kx}{{ax} + 1} + \frac{{bx} + 1}{{cx} + 1} < 0$ 的解集为\_\_\_\_\_.

\section*{二、选择题}

10. 已知 $x > 0,P = \sqrt{x + 1} - \sqrt{x},Q = \sqrt{x + 3} - \sqrt{x + 2}$ ,则 $P$ 与 $Q$ 的大小关系为 ( )

A. $P > Q$ . B. $P < Q$ . C. $P = Q$ . D. 不确定.

11.16 世纪中叶, 英国数学家雷科德在《砺智石》一书中首先把 “ = ”作为等号使用，后来英国数学家哈利奥特首次使用“<”和“>”符号，并逐步被数学界接受，不等号的引入对不等式的发展影响深远. 已知 $b < a < 0$ ,则下列选项正确的是 ( )

A. ${a}^{2} > {b}^{2}$ . B. $a + b < {ab}$ . C. $\left| b\right|  < \left| a\right|$ . D. ${ab} > {b}^{2}$ .

12. 下列各组中, $M$ 是 $N$ 的充要条件的是 ( )

A. $M : \left| x\right|  + \left| y\right|  \leq  1,N : {x}^{2} + {y}^{2} \leq  1$ .

B. $M$ : 实数 $a$ 、 $b$ 满足 $a + b > 2$ 且 ${ab} > 1,N : a > 1$ 且 $b > 1$ .

C. 集合 $E\text{ 、 }F\text{ 、 }P,M : P \subset  E$ 且 $P \subset  F,N : P \subset  E \cap  F$ .

D. $M : \left| a\right|  \geq  1,N$ : 曲线 $y = \left| x\right|$ 与直线 $y = {ax} + 2$ 有且只有一个交点.

\section*{三、解答题}

13. 设 $x\text{ 、 }y\text{ 、 }a\text{ 、 }b \in  \mathbf{R}$ ,且 $a\text{ 、 }b > 0,a \neq  b$ ,分别比较下列两个数的大小.

(1) $m = {x}^{2} + {y}^{2} - {2x} + {2y},n =  - 3$ ，比较 $m$ 、 $n$ 的大小；

(2) $S = {a}^{3} + {b}^{3},T = {a}^{2}b + a{b}^{2}$ ，比较 $S$ 、 $T$ 的大小.

14. 解下列不等式:

(1) $3{x}^{2} - {2x} - 1 \geq  0$ ； (2) $\frac{x + 2}{x - 1} \leq  0$ ； (3) $\left| {x + 3}\right|  - \left| {{2x} - 1}\right|  < \frac{x}{2} + 1$ .

15. 已知 $m \in  \mathbf{R},f\left( x\right)  = \frac{x - m}{x + 1}$ .

(1)若 $m = 2$ ，解不等式 $f\left( x\right)  > 0$ ；

(2)讨论 $f\left( x\right)  > 0$ 的解集.

16. 如图,在梯形 ${PMNQ}$ 中, ${PQ}//{MN}$ ,对角线 ${PN}$ 和 ${MQ}$ 相交于点 $O$ ,并把梯形分成四部分,记这四部分的面积分别为 ${S}_{1}$ 、 ${S}_{2}$ 、 ${S}_{3}$ 、 ${S}_{4}$ . 试判断 ${S}_{1} + {S}_{2}$ 和 ${S}_{3} + {S}_{4}$ 的大小关系,并证明你的结论.

\begin{center}
\includegraphics[width=0.2\textwidth]{images/bo_d7tk7oc91nqc73875fu0_20_1107_1371_323_248_0.jpg}
\end{center}
 

(第 16 题图)

§ 2.3 基本不等式

\section*{典型例题精讲}

【例 1】解答下列问题:

(1)若 $x$ 、 $y > 0$ ，且 $x + {2y} = 1$ ，求 $\frac{3}{x} + \frac{2}{y}$ 的最小值；

(2)若 $x$ 、 $y > 0$ ，且 $\frac{1}{x} + \frac{2}{y} = 1$ ，求 ${3x} + {2y}$ 的最小值；

(3)若 $0 < x < 1$ ，求函数 $f\left( x\right)  = \frac{3}{x} + \frac{4}{1 - x}$ 的最小值.

【例 2】已知函数 $f\left( x\right)  = \left| {x + 1}\right|  + \left| x\right|$ .

(1)求不等式 $f\left( x\right)  \leq  3$ 的解集；

(2)若 $x \in  \left\lbrack  {a - 2,a + 2}\right\rbrack$ 时， $f\left( x\right)  < {3a} - 2$ ，求实数 $a$ 的取值范围.

【例 3】设函数 $f\left( x\right)  = \left| {x - 1}\right|  + \left| {{2x} + a}\right| \left( {a \in  \mathbf{R}}\right)$ .

(1)若 $a = 2$ ，求 $f\left( x\right)  \leq  8$ 的解集；

(2)若 $f\left( x\right)  \geq  3 - \left| {x - 1}\right|$ 对任意 $x \in  \mathbf{R}$ 恒成立，求实数 $a$ 的取值范围.

【例 4】已知 $a$ 、 $b$ 、 $c > 0$ ，且 $a + b + c = 1$ . 证明下列不等式:

(1) $\frac{1}{a} + \frac{1}{b + c} \geq  4$ ；

(2) $\frac{1}{a + {2b}} + \frac{1}{b + {2c}} + \frac{1}{c + {2a}} \geq  3$ .

【例 5】已知正实数 $a\text{ 、 }b$ 满足 ${\left( a + 1\right) }^{3} + {\left( b + 1\right) }^{3} = {16}$ ,设 $a + b$ 的最大值为 $m$ .

(1)求 $m$ 的值；

(2)若 $\left| {{2x} - a}\right|  \leq  m,\left| {{2y} - a}\right|  \leq  m$ ，求证: $\left| {{2x} - {4y} + a}\right|  \leq  6$ .

\section*{四基达标精练}

\section*{一、填空题}

1.若正数 $x\text{ 、 }y$ 满足 ${xy} = 1$ ，则 $x + {2y}$ 的最小值为\_\_\_\_\_.

2. 已知 $m > 0,n > 0,m + n = 1$ ,则 $\frac{1}{m} + \frac{2}{n}$ 的最小值为\_\_\_\_\_.

3. 若 $\left| {a - 1}\right|  - \left| a\right|  = 1$ ，则实数 $a$ 的取值范围是\_\_\_\_\_.

4. 若 $x$ 、 $y > 0$ ，且 $\frac{1}{x} + {2y} = 3$ ，则 $\frac{y}{x}$ 的最大值为\_\_\_\_\_.

5. 已知 $a$ 、 $b \in  \mathbf{R}$ ，且 $a - {3b} + 6 = 0$ ，则 ${2}^{a} + \frac{1}{{8}^{b}}$ 的最小值为\_\_\_\_\_.

6. 若不等式 $\left| {x - 4}\right|  - \left| {x - 3}\right|  \leq  a$ 对一切 $x \in  \mathbf{R}$ 恒成立，则实数 $a$ 的取值范围是\_\_\_\_\_.

7. 已知 $x > y > 0$ ，则 $\frac{{x}^{2} + {y}^{2}}{{xy} - {y}^{2}}$ 的最小值是\_\_\_\_\_.

8．正项等比数列 $\left\{  {a}_{n}\right\}$ 中， ${a}_{2023} = {a}_{2022} + 2{a}_{2021}$ ，若 ${a}_{m}{a}_{n} = {16}{a}_{1}^{2}$ ，则 $\frac{1}{m} + \frac{4}{n}$ 的最小值等于\_\_\_\_\_.

9. 已知定义在 $\mathbf{R}$ 上的偶函数 $f\left( x\right)  = \left| {x - m + 1}\right|  - 2$ ,若正实数 $a\text{ 、 }b$ 满足 $f\left( a\right)  + f\left( {2b}\right)  = m$ ,则 $\frac{1}{a} + \frac{2}{b}$ 的最小值为\_\_\_\_\_.

\section*{二、选择题}

10. 下列不等式恒成立的是 ( )

A. ${a}^{2} + {b}^{2} \leq  {2ab}$ . B. ${a}^{2} + {b}^{2} \geq   - {2ab}$ . C. $a + b \geq  2\sqrt{\left| ab\right| }$ . D. $a + b \geq   - 2\sqrt{\left| ab\right| }$ .

11. 等式 $\left| {a - b}\right|  = \left| {a - 1}\right|  + \left| {b - 1}\right|$ 成立的条件是 ( )

A. $\left( {a - 1}\right) \left( {b - 1}\right)  < 0$ . B. $\left( {a - 1}\right) \left( {b - 1}\right)  \leq  0$ . C. $\left( {a - 1}\right) \left( {b - 1}\right)  > 0$ . D. $\left( {a - 1}\right) \left( {b - 1}\right)  \geq  0$ .

12. 已知 ${2}^{m} = {3}^{n} = 6$ ,则 $m\text{ 、 }n$ 不可能满足的关系是 ( )

A. $m + n > 4$ . B. ${mn} > 4$ . C. ${m}^{2} + {n}^{2} < 8$ . D. ${\left( m - 1\right) }^{2} + {\left( n - 1\right) }^{2} > 2$ .

\section*{三、解答题}

13. 某化工厂生产某种产品,当年产量在 150 吨至 250 吨 (包括 150 吨与 250 吨) 时,每年的生产成本 $y$ 万元与年产量 $x$ 吨之间的关系可近似地表示为 $y = \frac{1}{10}{x}^{2} - {30x} + {4000}$ .

(1)若每年的生产总成本不超过 2000 万元，求年产量 $x$ 的取值范围；

(2)求年产量为多少吨时, 每吨的平均成本最低, 并求每吨的最低成本.

14. 已知 $a\text{ 、 }b$ 是正常数， $a \neq  b,x\text{ 、 }y \in  \left( {0, + \infty }\right)$ .

(1)求证: $\frac{{a}^{2}}{x} + \frac{{b}^{2}}{y} \geq  \frac{{\left( a + b\right) }^{2}}{x + y}$ ，指出等号成立的条件；

(2)利用(1)的结论求函数 $f\left( x\right)  = \frac{2}{x} + \frac{9}{1 - {2x}}$ ( $0 < x < \frac{1}{2}$ )的最小值，并指出取最小值时 $x$ 的值.

15. 设 $a > 0$ ，实数 $x$ 、 $y$ 满足 $\left| {x - 1}\right|  < \frac{a}{3}$ ， $\left| {y - 2}\right|  < \frac{a}{3}$ .

(1)若 $y = {3x} - 2$ ，且 $\left\{  {x\left| \right| x - 1\left| { < \frac{a}{3}\} \cap \{ y}\right| \left| {y - 2}\right|  < \frac{a}{3}}\right\}   \neq  \varnothing$ ，求 $a$ 的取值范围；

(2)求证: $\left| {{{2x} + y} - 4}\right|  < a$ .

16. 已知函数 $f\left( x\right)  = \left| {x - 2}\right| ,g\left( x\right)  =  - \left| {x + 3}\right|  + m\left( {m \in  \mathbf{R}}\right)$ .

(1)解关于 $x$ 的不等式 $f\left( x\right)  + a - 2 > 0\left( {a \in  \mathbf{R}}\right)$ ；

(2)若函数 $f\left( x\right)$ 的图像恒在函数 $g\left( x\right)$ 图像的上方，求 $m$ 的取值范围.

\section*{§ 2.4 不等式综合}

\section*{典型例题精讲}

【例 1】已知函数 $f\left( x\right)  = {ax} - 1,g\left( x\right)  = {x}^{2} - a$ .

(1)若不等式 $\frac{f\left( x\right) }{g\left( x\right) } < 0$ 的解集为 $S$ ,且 $3 \in  S,4 \notin  S$ ,求实数 $a$ 的取值范围；

(2)若不等式 $f\left( x\right)  \leq  g\left( x\right)$ 对任意 $x \in  \left\lbrack  {0,1}\right\rbrack$ 恒成立，求实数 $a$ 的取值范围.

【例 2】已知 $a\text{ 、 }b \in  \left( {0, + \infty }\right)$ ,证明下列结论:

(1) $\frac{a}{\sqrt{b}} + \frac{b}{\sqrt{a}} \geq  \sqrt{a} + \sqrt{b}$ ；

(2)若 $a + b = {10}$ ，求证: $\sqrt{1 + {3a}} + \sqrt{1 + {3b}} \leq  8$ .

【例 3】已知 $m + {2n} = 2$ ，且 $m >  - 1$ ， $n > 0$ .

(1)求 $\frac{1}{m + 1} + \frac{2}{n}$ 的最小值；

(2)求 $\frac{{m}^{2}}{{2n} + 2} + \frac{4{n}^{2}}{m + 1}$ 的最小值.

【例 4】已知正数 $a$ 、 $b$ 和实数 $t$ 满足 ${a}^{2} + {tab} + {b}^{2} = 1$ .

(1)若 $t = 1$ ，求 $a + b$ 的最大值；

(2)若 $a + b$ 存在最大值，求 $t$ 的取值范围.

【例 5】我们知道 $\left| {AB}\right|  = \left| {{x}_{A} - {x}_{B}}\right|$ 表示数轴上 $A$ 、 $B$ 两点的距离，它也可以看作满足一定条件的一种运算. 这样,我们也可以将满足下列三个条件的一个 $x$ 与 $y$ 间的运算 $p\left( {x,y}\right)$ 叫做 $x\text{ 、 }y$ 之间的距离: 条件 \ding{172} : 非负性 $p\left( {x,y}\right)  \geq  0$ ,当且仅当 $x = y$ 等号成立; 条件 \ding{173} : 交换律 $p\left( {x,y}\right)  = p\left( {y,x}\right)$ ; 条件 \ding{174} : $p\left( {x,z}\right)  \leq  p\left( {x,y}\right)  + p\left( {y,z}\right)$ 成立. 试判断运算 $s\left( {x,y}\right)  = \frac{\left| x - y\right| }{1 + \left| {x - y}\right| }$ 是否为一个距离? 如果是,请加以证明; 如果不是, 请举出反例.

\section*{四基达标精练}

\section*{一、填空题}

1. 若函数 $f\left( x\right)  = x + \frac{1}{x - 2}\left( {x > 2}\right)$ 在 $x = a$ 处取得最小值，则 $a =$ \_\_\_\_\_.

2. 已知 ${\log }_{5}x + {\log }_{5}y = 2$ ,则 $x + {4y}$ 的最小值为\_\_\_\_\_.

3. 若 $x = k$ 时，函数 $f\left( x\right)  = \left| {x - 2}\right|  + \left| {x + 1}\right|$ 取得最小值时，则实数 $k$ 的取值范围是\_\_\_\_\_.

4. 若不等式 $\left| {x - 1}\right|  - \left| {x + 1}\right|  > {2a}$ 的解集非空，则实数 $a$ 的取值范围是\_\_\_\_\_.

5. 已知 $a\text{ 、 }b > 0,a \geq  \frac{1}{a} + \frac{2}{b},b \geq  \frac{1}{b} + \frac{2}{a}$ ,则 $a + b$ 的最小值为\_\_\_\_\_.

6. 关于 $x$ 的不等式 $a{x}^{2} - \left| x\right|  + {2a} \geq  0$ 的解集是 $\mathbf{R}$ ，则实数 $a$ 的取值范围为\_\_\_\_\_.

7. 已知 $x + \frac{4}{x} - a - \frac{4}{a} \geq  0$ 对任意 $x > 0$ 恒成立，则实数 $a$ 的取值范围是\_\_\_\_\_. 8. 已知 $x$ 、 $y > 0$ ， $x + y + \sqrt{xy} = 2$ ，则 $x + y$ 的取值范围是\_\_\_\_\_.

9. 设 $A$ 、 $B$ 是非空数集，若对任意 $x \in  A$ 、 $y \in  B$ ，都有唯一确定的 $f\left( {x,y}\right)$ 与之对应，则称 $f\left( {x,y}\right)$ 为关于 $x\text{ 、 }y$ 的二元函数,现定义满足下列性质的二元函数 $f\left( {x,y}\right)$ 为关于实数 $x\text{ 、 }y$ 的广义“距离”.

(1)非负性: $f\left( {x,y}\right)  \geq  0$ ，当且仅当 $x = y$ 时取等号；

(2)对称性: $f\left( {x,y}\right)  = f\left( {y,x}\right)$ ；

(3)三角形不等式: $f\left( {x,y}\right)  \leq  f\left( {x,z}\right)  + f\left( {y,z}\right)$ 对任意的实数 $z$ 均成立.

给出下列三个二元函数: \ding{172}  $f\left( {x,y}\right)  = \sqrt{x - y}$ ;  \ding{173}  $f\left( {x,y}\right)  = {\left( x - y\right) }^{2}$ ;  \ding{174}  $f\left( {x,y}\right)  = \left| {x - y}\right|$ .

其中能够成为关于 $x$ 、 $y$ 的广义“距离”的二元函数的序号是\_\_\_\_\_.

\section*{二、选择题}

10. 已知 $x \in  \mathbf{R}$ ,下列不等式中正确的是 ( )

A. $\frac{1}{{2}^{x}} > \frac{1}{{3}^{3}}$ . B. $\frac{1}{{x}^{2} - x + 1} > \frac{1}{{x}^{2} + x + 1}$ .

C. $\frac{1}{2\left| x\right| } > \frac{1}{{x}^{2} + 1}$ . D. $\frac{1}{{x}^{2} + 1} > \frac{1}{{x}^{2} + 2}$ .

11. 在 $\mathbf{R}$ 上定义运算 $\otimes   : x \otimes  y = \frac{x}{2 - y}$ ,若关于 $x$ 的不等式 $\left( {x - a}\right)  \otimes  \left( {x - 1 - a}\right)  \geq  0$ 的解集是集合 $\{ x \mid   - 2 < x \leq  4\}$ 的子集,则实数 $a$ 的取值范围为 ( )

A. $- 2 < a < 1$ . B. $- 2 \leq  a < 1$ . C. $- 2 < a \leq  1$ . D. $- 2 \leq  a \leq  1$ .

12. 已知 $a\text{ 、 }b > 0$ ,有下列不等式:  \ding{172}  $\left| {a - 1}\right|  + \left| a\right|  \geq  1$ ;  \ding{173}  ${ab} + \frac{2}{ab} > 2$ ;  \ding{174}  ${a}^{2} + {b}^{2} > {4ab} - 4{b}^{2}$ ;  \ding{175}  $\frac{b}{a} > \frac{b + 1}{a + 1}$ , 其中恒成立的有 ( )

A.  \ding{172}  \ding{173} . B.  \ding{173}  \ding{175} . C.  \ding{172}  \ding{173}  \ding{174} . D.  \ding{172}  \ding{175} .

\section*{三、解答题}

13. 已知 $x > 0$ ,求 $x - 1 + \frac{2}{x}$ 的最小值. 某同学的解答如下:

因为 $x > 0$ ,所以 $x - 1 + \frac{2}{x} \geq  2\sqrt{\left( {x - 1}\right)  \cdot  \frac{2}{x}}$ ,当且仅当 $x - 1 = \frac{2}{x}$ 上式中等号成立,

即 ${x}^{2} - x - 2 = 0$ ,解得 ${x}_{1} = 2,{x}_{2} =  - 1$ (舍)

当 $x = 2$ 时, $2\sqrt{\left( {x - 1}\right)  \cdot  \frac{2}{x}} = 2$ ,所以当 $x = 2$ 时, $x - 1 + \frac{2}{x}$ 的最小值为 2 .

请判断该同学的结论的正误，如果你认为正确，则解答完毕；如果你认为错误，请指出所有的错误，并给出正确的解答.

14.某商品计划两次提价,有甲、乙、丙三种方案,其中 $p > q > 0$ .

\begin{center}
\begin{tabular}{|c|c|c|}
\hline
\makecell{ 次 \\ 方案 } & 第一次提价 & 第二次提价 \\
\cline{1-3}
甲 & p\% & q\% \\
\cline{1-3}
乙 & q\% & p\% \\
\cline{1-3}
丙 & $\frac{p + q}{2}\%$ & $\frac{p + q}{2}\%$ \\
\cline{1-3}
\hline
\end{tabular}
\end{center}

经两次提价后, 哪一种方案的提价幅度最大? 为什么?

15. 已知函数 $f\left( x\right)  = \left| {x - 1}\right|  + \left| {x + 2}\right| ,a > 0,b > 0$ .

(1)求不等式 $f\left( x\right)  \leq  7$ 的解集；

(2)若 $f\left( x\right)$ 的最小值为 $a + {2b}$ ，且 $\frac{1}{a} + \frac{2}{b} \geq  {m}^{2}$ 恒成立，求实数 $m$ 的取值范围.

16. 已知函数 $f\left( x\right)  = {kx} + b$ 的图像与 $x$ 、 $y$ 轴分别相交于点 $A$ 、 $B$ ， $\overrightarrow{AB} = 2\overrightarrow{i} + 2\overrightarrow{j}$ ( $\overrightarrow{i}$ 、 $\overrightarrow{j}$ 分别是与 $x$ 、 $y$ 轴正半轴同方向的单位向量),函数 $g\left( x\right)  = {x}^{2} - x - 6$ .

(1)求 $k$ 、 $b$ 的值；

(2)当 $x$ 满足 $f\left( x\right)  > g\left( x\right)$ 时,求函数 $y = \frac{g\left( x\right)  + 1}{f\left( x\right) } + \left| {x - 4}\right|  + \left| {x + 2}\right|$ 的最小值.

\section*{第 3 章 函数的概念与性质(必修 1)}

\section*{知识梳理*}

\section*{一、函数的概念}

1. 函数的两要素:定义域，对应法则；注意定义域优先的原则.

2. 函数定义域的求法:列不等式组，解不等式组.

原则: \ding{172}  分式分母不为零； \ding{173}  偶次方根被开方数非负； \ding{174}  零次方底数非零； \ding{175}  对数的底数为正且不等于 1，真数大于零； \circled{5} 正切函数的角的终边不能落在 $y$ 轴上； \circled{6} 实际问题的定义域要根据实际意义来确定； \circled{7} 复合函数定义域,若已知 $f\left( x\right)$ 的定义域是 $\left\lbrack  {a,b}\right\rbrack$ ,则 $f\left\lbrack  {g\left( x\right) }\right\rbrack$ 的定义域由 $a \leq  g\left( x\right)  \leq  b$ 的解集，且 $g\left( x\right)$ 有意义的 $x$ 的范围确定.

3. 注意函数的定义域与不等式恒成立之间的关系.

\section*{二、函数的基本性质:}

1. 函数的奇偶性 (函数定义域内的整体性质、定义域关于原点对称是必要条件):

(1)若 $f\left( x\right)$ 是偶函数，那么 $f\left( x\right)  = f\left( {-x}\right)  = f\left( \left| x\right| \right)$ (图像关于 $y$ 轴对称)；

(2)若 $f\left( x\right)$ 是奇函数，那么 $f\left( x\right)  =  - f\left( {-x}\right)$ (图像关于原点中心对称)；

(3)判断函数奇偶性可用定义的等价形式: $f\left( x\right)  \pm  f\left( {-x}\right)  = 0$ 或 $\frac{f\left( {-x}\right) }{f\left( x\right) } =  \pm  1\left( {f\left( x\right)  \neq  0}\right)$ ;

(4)判断函数奇偶性的实质是判断两个互为相反数的自变量的函数值相等或互为相反数;

(5)确定一个函数既不是奇函数也不是偶函数,通过举反例:如 $f\left( {-1}\right)  \neq  f\left( 1\right)$ 且 $f\left( {-1}\right)  \neq   - f\left( 1\right)$ 等, 注意此时不能用连续的不等号写成 $f\left( {-1}\right)  \neq  f\left( 1\right)  \neq   - f\left( 1\right)$ 的形式;

(6)分段函数的奇偶性分段求值和判断，复杂形式先化简后判断；

(7)原点处有意义的奇函数图像必过原点 (可用于求参数，但要给予一般性证明)；但 $f\left( 0\right)  = 0$ 是一个函数为奇函数的既不充分也不必要条件, 因为有定义域的限制;

(8)奇(偶)函数在关于原点对称的区间内单调性相同(反);

(9)奇(偶)函数在关于原点对称的区间内符号相反(同).

2. 函数图像 (或方程曲线) 的对称性 (中心对称关键:中点坐标公式；轴对称关键:垂直、平分):

(1) $f\left( x\right)$ 满足 $f\left( {a + x}\right)  = f\left( {a - x}\right)$ (或 $f\left( x\right)  = f\left( {{2a} - x}\right)$ )，则 $f\left( x\right)$ 图像关于 $x = a$ 对称；

(2) $f\left( x\right)$ 满足 $f\left( {a + x}\right)  =  - f\left( {a - x}\right)$ (或 $f\left( x\right)  =  - f\left( {{2a} - x}\right)$ )，则 $f\left( x\right)$ 图像关于(a，0)对称；

(3) $f\left( x\right)$ 满足 $f\left( {a + x}\right)  = {2b} - f\left( {a - x}\right)$ (或 $f\left( x\right)  = {2b} - f\left( {{2a} - x}\right)$ )，则 $f\left( x\right)$ 图像关于 $\left( {a,b}\right)$ 对称；

(4)特别的, $f\left( {x,y}\right)  = 0$ 关于 $y = x + a\left( {y =  - x + a}\right)$ 对称的曲线方程为: $f\left( {y - a,x + a}\right)  = 0(f(a - y$ , $a - x) = 0)$ ;

(5) 曲线 $f\left( {x,y}\right)  = 0$ 关于点 $\left( {a,b}\right)$ 对称的曲线方程为: $f\left( {{2a} - x,{2b} - y}\right)  = 0$ ;

(6) 函数 $y = f\left( {x - a}\right)$ 与 $y = f\left( {b - x}\right)$ 的图像关于直线 $x = \frac{a + b}{2}$ 对称.

说明:

 \ding{172} 证明函数图像的对称性，即证图像上任意点关于对称中心(对称轴)的对称点仍在图像上； \ding{173} 证明图像 ${C}_{1}\text{ 、 }{C}_{2}$ 的对称性,即证 ${C}_{1}$ 上任意点关于中心(轴)的对称点仍在 ${C}_{2}$ 上,反之亦然.

\HRule

* 注:本章节所属知识模块是 “函数与分析”.

\HRule

3. 函数的单调性 (函数定义域内的局部性质,单调性相同但间断的区间一般不能并在一起书写):

(1)单调性的定义:设 ${x}_{1}$ 、 ${x}_{2}$ 是函数 $y = f\left( x\right)$ 定义域内某区间 $I$ 上的任意两个自变量，

若 ${x}_{1} < {x}_{2}$ ,都有 $f\left( {x}_{1}\right)  \leq  f\left( {x}_{2}\right) \left( {f\left( {x}_{1}\right)  \geq  f\left( {x}_{2}\right) }\right)$ ,则 $y = f\left( x\right)$ 在区间 $I$ 上是增 (减) 函数;

若 ${x}_{1} < {x}_{2}$ ,都有 $f\left( {x}_{1}\right)  < f\left( {x}_{2}\right) \left( {f\left( {x}_{1}\right)  > f\left( {x}_{2}\right) }\right)$ ,则 $y = f\left( x\right)$ 在区间 $I$ 上是严格增 (减) 函数;

以上 “若……，都有……” 可改写为: “若 ${x}_{1} \neq  {x}_{2}$ ，都有 $\frac{f\left( {x}_{1}\right)  - f\left( {x}_{2}\right) }{{x}_{1} - {x}_{2}} > 0\left( {\frac{f\left( {x}_{1}\right)  - f\left( {x}_{2}\right) }{{x}_{1} - {x}_{2}} < 0}\right)$ ”，或改写为:“若 ${x}_{1} \neq  {x}_{2}$ ，都有 $\left( {{x}_{1} - {x}_{2}}\right) \left\lbrack  {f\left( {x}_{1}\right)  - f\left( {x}_{2}\right) }\right\rbrack   > 0$ ( $\left( {{x}_{1} - {x}_{2}}\right) \left\lbrack  {f\left( {x}_{1}\right)  - f\left( {x}_{2}\right) }\right\rbrack   < 0$ )”；

(2)规范证明步骤:任意取 ${x}_{1}\text{ 、 }{x}_{2} \in  I \subseteq  D$ 且 ${x}_{1} < {x}_{2}$ ，则 $f\left( {x}_{1}\right)  - f\left( {x}_{2}\right)  = \cdots \cdots ;\because {x}_{1}\text{ 、 }{x}_{2} \in  I \subseteq  D$ 且 ${x}_{1} < {x}_{2},\therefore \cdots \cdots ,\therefore f\left( {x}_{1}\right)  > f\left( {x}_{2}\right)$ (或 $f\left( {x}_{1}\right)  < f\left( {x}_{2}\right) ),\therefore f\left( x\right)$ 在 $I$ 上是减 (增函数);

(3)求单调区间的方法:定义法、图像法；(注意:极限思想的应用、单调区间一定是定义域的子集.

(4)复合函数 $y = f\left\lbrack  {g\left( x\right) }\right\rbrack$ 的单调性(同增异减): $y = f\left\lbrack  {g\left( x\right) }\right\rbrack$ 由 $y = f\left( u\right)$ 和 $u = g\left( x\right)$ 复合而成，若 $x \in  \left( {a,b}\right)$ 时, $u \in  \left( {c,d}\right)$ ,当 $u = g\left( x\right)$ 在 $x \in  \left( {a,b}\right)$ 上的单调性与 $y = f\left( u\right)$ 在 $u \in  \left( {c,d}\right)$ 上的单调性相同 (反) 时 $y = f\left\lbrack  {g\left( x\right) }\right\rbrack$ 在 $x \in  \langle a,b\rangle$ 上单调递增 (减); (其中 “ ( ) ” 随题目的变化确定区间的开闭)

(5)函数单调性的简要运算性质:在公共定义域内，增十增是增函数；减十减是减函数；增一减是增函数; 减一增是减函数.

注意:(4)(5)可作为判断单调性的依据，但不可作为解答题的证明过程.

4. 函数的最值(函数定义域内的整体性质):

(1)最值的定义:

对于函数 $y = f\left( x\right) \left( {x \in  D}\right)$ ,若存在 ${x}_{0} \in  D$ ,使得 $f\left( x\right)  \geq  \left(  \leq  \right) f\left( {x}_{0}\right)$ 对任意 $x \in  D$ 恒成立,则 $f\left( {x}_{0}\right)$ 称为函数 $f\left( x\right)$ 的最小 (大) 值. (本书用 $f{\left( x\right) }_{\max }$ 表示 $f\left( x\right)$ 的最大值， $f{\left( x\right) }_{\min }$ 表示最小值)

(2)函数值域(和最值)的求法:

 \ding{172} 配方法:适用于二次函数(或可转化为二次函数)及区间上的二次函数问题;

 \ding{173} 不等式法:适用于“耐克函数”；注意“一正二定三取等”；

 \ding{174} 换元法:通常用于对函数式进行变形以获得更为熟悉的函数；

 \ding{175} 有界性法:利用已知函数的有界性，借助解不等式求解；

 \circled{5} 单调性法:当所产生的函数单调性较明显或其他方法不太奏效时，考察函数的单调性有时很有效；

 \circled{6} 三角函数的特殊方法:辅助角公式；

 \circled{7} 数列的最大项和最小项的特殊性和特殊方法:最常用的是单调性.

注意:

 \ding{172} 数形结合是解决区间上的函数的最值和值域问题的重要手段，是求解这类问题时避免发生错误的根本保证.  \ding{173}  注意“最值”与“极值”的区别:最值是函数定义域内的整体性质，而极值是局部的.

5. 函数的零点 (函数定义域内的整体性质):

(1)零点的意义:

对于函数 $y = f\left( x\right) \left( {x \in  D}\right)$ ,若存在 ${x}_{0} \in  D$ ,使得 $f\left( {x}_{0}\right)  = 0$ 成立,则 ${x}_{0}$ 称为函数 $f\left( x\right)$ 的零点; 注意函数的零点是函数定义域内的一个实数，而不是有序实数对；

(2)零点存在定理:如果在区间 $\left\lbrack  {a,b}\right\rbrack$ 上，函数 $y = f\left( x\right)$ 的图像满足一段连续曲线，并且 $f\left( a\right)  \cdot  f\left( b\right)  < 0$ ，那么 $y = f\left( x\right)$ 在区间 $\left( {a,b}\right)$ 上一定有零点;

(3)二分法求零点:在确定了区间 $\left( {a,b}\right)$ 上一定有零点的前提下，将区间一分为二，这两部分中总有一个含有零点，而含有零点的区间的长度变为原先的一半，反复执行这种“一分为二”的操作，就能将零点限制在一个足够小的区间中，从而容易求得其近似值了.

6. 函数的周期性(函数定义域内的整体性质):

(1)对于函数 $y = f\left( x\right) \left( {x \in  D}\right)$ ,存在常数 $T \neq  0$ ,使 $f\left( {x + T}\right)  = f\left( x\right)$ 对任意 $x \in  D$ 恒成立,则 $T$ 为 $f\left( x\right)$ 的一个周期;

(2)若偶函数 $f\left( x\right)$ 的图像关于直线 $x = a\left( {a \neq  0}\right)$ 对称，则 $f\left( x\right)$ 是周期函数， $2\left| a\right|$ 是其一个周期；

二若函数 $f\left( x\right)$ 的图像关于直线 $x = a$ 和 $x = b\left( {a \neq  b}\right)$ 对称,则 $f\left( x\right)$ 是周期函数， $2\left| {a - b}\right|$ 是其周期； (3)若奇函数 $f\left( x\right)$ 的图像关于直线 $x = a\left( {a \neq  0}\right)$ 对称，则 $f\left( x\right)$ 是周期函数， $4\left| a\right|$ 是其一个周期；

* 若函数 $f\left( x\right)$ 的图像关于 $\left( {a,0}\right)$ 和 $x = b\left( {a \neq  b}\right)$ 对称,则 $f\left( x\right)$ 是周期函数， $4\left| {a - b}\right|$ 是其一个周期； (4)若奇函数 $f\left( x\right)$ 的图像关于点 $\left( {b,0}\right) \left( {b \neq  0}\right)$ 对称，则 $f\left( x\right)$ 是周期函数， $2\left| b\right|$ 是其一个周期；

二若函数 $f\left( x\right)$ 的图像关于 $\left( {a,0}\right)$ 和 $\left( {b,0}\right) \left( {a \neq  b}\right)$ 对称,则 $f\left( x\right)$ 是周期函数, $2\left| {a - b}\right|$ 是其一个周期.

7. 处理二次函数的问题勿忘数形结合；二次函数在闭区间上必有最值；求解最值问题用“两看”:一看开口方向；二看对称轴与所给区间的相对位置关系.

8.二次函数涉及到方程根的分布问题看:开口、对称轴、端点的函数值与零的关系，特别提醒:端点特别考虑.

\section*{三、函数图像的变换:}

1. 函数图像平移变换 (不改变形状、只改变位置):

(1)将 $y = f\left( x\right)$ 图像沿 $x$ 轴向左 $\left( {a > 0}\right)$ (或右 $\left( {a < 0}\right)$ )平移 $\left| a\right|$ 个单位得到函数 $y = f\left( {x + a}\right)$ 的图像；

(2)将 $y = f\left( x\right)$ 图像沿 $y$ 轴向上 $\left( {b > 0}\right)$ (或下 $\left( {b < 0}\right)$ 〕平移 $\left| b\right|$ 个单位得到函数 $y = f\left( x\right)  + b$ 的图像.

2. 几种特殊的对称变换:

(1)函数 $y = f\left( x\right)$ 与函数 $y = f\left( {-x}\right)$ 的图像关于 $y$ 轴 ( $x = 0$ ) 对称;

(2)函数 $y = f\left( x\right)$ 与函数 $y =  - f\left( x\right)$ 的图像关于 $x$ 轴( $y = 0$ )对称；

(3) 函数 $y = f\left( x\right)$ 与函数 $y =  - f\left( {-x}\right)$ 的图像关于坐标原点对称.

3. 伸缩变换:

(1)将 $y = f\left( x\right)$ 图像上各点的纵坐标伸长 $\left( {a > 1}\right)$ 或缩短 $\left( {0 < a < 1}\right)$ 到原来的 $a$ 倍(横坐标不变)，可得 $y  = {af}\left( x\right) \left( {a > 0}\right)$ 图像;

(2) 将 $y = f\left( x\right)$ 的图像上各点的横坐标伸长 $\left( {0 < a < 1}\right)$ 或缩短 $\left( {a > 1}\right)$ 到原来的 $\frac{1}{a}$ 倍 (纵坐标不变),可得 $y = f\left( {ax}\right) \left( {a > 0}\right)$ 的图像.

4. 翻折变换:

(1) $y = f\left( \left| x\right| \right)$ 的图像作法:保持 $y = f\left( x\right)$ 在 $y$ 轴右侧的图像不变，同时翻折到 $y$ 轴左侧；

(2) $y = \left| {f\left( x\right) }\right|$ 的图像作法:保持 $y = f\left( x\right)$ 在 $x$ 轴上方的图像不变，同时将下方的图像翻折到 $x$ 轴上方.

5. 对于函数 $y = f\left( x\right) \left( {x \in  \mathbf{R}}\right) ,f\left( {x + a}\right)  = f\left( {b - x}\right)$ 恒成立，则函数 $f\left( x\right)$ 的对称轴是 $x = \frac{a + b}{2}$ ；两个函数 $y = f\left( {x + a}\right)$ 与 $y = f\left( {b - x}\right)$ 的图像关于直线 $x = \frac{b - a}{2}$ 对称.

\section*{$\text{ § }{3.1}$ 函数的概念及其表示法}

\section*{典型例题精讲}

【例 1】已知函数 $f\left( x\right)  = \left\{  \begin{array}{l} {x}^{2},x < 0, \\  4 - {2x},x > 0. \end{array}\right.$

\begin{center}
\includegraphics[width=0.2\textwidth]{images/bo_d7tk7oc91nqc73875fu0_32_1141_538_336_339_0.jpg}
\end{center}
 

(例 1 图)

(1)画出函数 $f\left( x\right)$ 的图像；

(2)若 $x \in  \left\lbrack  {-2,0)\cup (0,3}\right\rbrack$ ，求其值域；

(3)当 $f\left( x\right)  \geq  2$ 时,求实数 $x$ 的取值范围.

【例 2】已知函数 $f\left( {{\log }_{3}x}\right)  = \frac{{2x} - 3}{x - 1}$ .

(1)求函数 $f\left( x\right)$ 的解析式及定义域；

(2)求函数 $f\left( x\right)$ 在 $x \in  \left( {-\infty ,0}\right)$ ∪ $\left( {0,2}\right)$ 时的值域.

【例 3】( 1 )已知函数 $f\left( x\right)$ 的定义域是 $\lbrack 1,7)$ ，求函数 $f\left( {{2x} - 3}\right)$ 的定义域；

(2)若函数 $f\left( \left| {x - 2}\right| \right)$ 的定义域是 $\lbrack 1,7)$ ，求 $f\left( x\right)$ 的定义域；

(3)若函数 $f\left( {{3x} - 2}\right)$ 的定义域是 $\lbrack 1,7)$ ，求 $f\left( {{2x} - 1}\right)$ 的定义域.

【例 4】已知函数 $g\left( x\right)  = \sqrt{x} + 1,h\left( x\right)  = \frac{1}{x + 3}\left( {-3 < x \leq  a}\right)$ ,其中 $a$ 为常数,且 $a > 0$ ,令函数 $f\left( x\right)  =  g\left( x\right) h\left( x\right)$ .

(1)求函数 $f\left( x\right)$ 的解析式，指出其定义域，并求当 $a = \frac{1}{4}$ 时，函数 $f\left( x\right)$ 的值域；

(2)是否存在正整数 $a$ ，使得函数 $f\left( x\right)$ 的值域恰为 $\left\lbrack  {\frac{1}{3},\frac{1}{2}}\right\rbrack$ ? 若存在，试写出所有满足条件的正整数 $a$ 所构成的集合; 若不存在,试说明理由.

【例 5】如图, $\bigtriangleup {OAB}$ 是边长为 2 的正三角形,记 $\bigtriangleup {OAB}$ 位于直线 $x = t\left( {t > 0}\right)$ 左侧的图形的面积为 $f\left( t\right)$ .

(1)求函数 $f\left( t\right)$ 解析式；

(2)当函数 $g\left( t\right)  = f\left( t\right)  - {at}$ 有且只有一个零点时，求实数 $a$ 的值.

\begin{center}
\includegraphics[width=0.2\textwidth]{images/bo_d7tk7oc91nqc73875fu0_33_1139_1055_298_267_0.jpg}
\end{center}
 

(例 5 图)

\section*{四基达标精练}

\section*{一、填空题}

1. 函数 $f\left( x\right)  = \frac{1}{x - 1} + \ln x$ 的定义域是\_\_\_\_\_.

2. 已知函数 $f\left( x\right)  = \left\{  \begin{array}{l} {\log }_{2}x,x > 0, \\  {3}^{x},x \leq  0, \end{array}\right.$ 则 $f\left\lbrack  {f\left( \frac{1}{4}\right) }\right\rbrack   =$ \_\_\_\_\_.

3. 若函数 $f\left( {{2x} + 1}\right)  = {4}^{x} - 2$ ，则函数 $f\left( x\right)$ 的解析式为\_\_\_\_\_.

4. 已知 $f\left( x\right)  = \left\{  {\begin{array}{l} {x}^{2},x \geq  0, \\   - x,x < 0, \end{array}g\left( x\right)  = f\left( {-x}\right) }\right.$ ，则 $g\left( x\right)$ 的解析式是\_\_\_\_\_.

5. 已知 $f\left( {x + 1}\right)$ 的定义域为 $\left\lbrack  {2,3}\right\rbrack$ ，则 $f\left( {{2x} - 1}\right)$ 的定义域为\_\_\_\_\_.313 (2).

6. 设 $f\left( x\right)  = \left\{  \begin{array}{l} {\left( \frac{1}{2}\right) }^{x} - 3,x \leq  0, \\  {x}^{\frac{1}{2}},x > 0, \end{array}\right.$ 若 $f\left( a\right)  < 1$ ，则实数 $a$ 的范围为\_\_\_\_\_.

7. 若 $y = \sqrt{\left| {x + 2}\right|  - \left| {x - m}\right|  + 1}$ 的定义域为 $\mathbf{R}$ ，则实数 $m$ 的取值范围是\_\_\_\_\_.

8. 用长为 $l$ 的铁丝弯成下部为矩形,上部为半圆形的框架,若矩形底边长为 ${2x}$ ,则此框架围成的面积 $y$ 关于 $x$ 的函数解析式为\_\_\_\_\_.

9. 已知 $a \in  \mathbf{R}$ ，若存在定义域为 $\mathbf{R}$ 的函数 $f\left( x\right)$ 同时满足下列两个条件: \ding{172} 对任意 ${x}_{0} \in  \mathbf{R}$ ， $f\left( {x}_{0}\right)$ 的值为 ${x}_{0}$ 或 ${x}_{0}^{2}$ ;  \ding{173} 关于 $x$ 的方程 $f\left( x\right)  = a$ 无实数解; 则 $a$ 的取值范围是\_\_\_\_\_.

\section*{二、选择题}

10. 某校有一班级,设变量 $x$ 是该班同学的姓名,变量 $y$ 是该班同学的学号,变量 $z$ 是该班同学的身高, 变量 $w$ 是该班同学的数学考试成绩,则下列选项中正确的是 ( )

A. $y$ 是 $x$ 的函数. B. $w$ 是 $y$ 的函数. C. $w$ 是 $z$ 的函数. D. $w$ 是 $x$ 的函数.

11. 存在函数 $y = f\left( x\right)$ ,满足对任意 $x \in  \mathbf{R}$ 都有 ( )

A. $f\left( {\left| x\right|  - 1}\right)  = \left| {x + 1}\right|$ . B. $f\left( \left| {x - 1}\right| \right)  = \left| {x + 1}\right|$ .

C. $f\left( {{x}^{2} + {2x}}\right)  = \left| {x + 1}\right|$ . D. $f\left( {{x}^{2} - 1}\right)  = \left| {x + 1}\right|$ .

12. 四参数方程的拟合函数表达式为 $y = \frac{a - d}{1 + {\left( \frac{x}{c}\right) }^{b}} + d\left( {x > 0}\right)$ ,常用于竞争系统和免疫检测,它的图像是一个递增 (或递减) 的类似指数曲线或对数曲线或双曲线 (如 $y = {x}^{-1}$ ),还可以是一条 ${S}$ 形曲线. 当 $a = 4$ , $b =  - 1,c = 1,d = 1$ 时,该拟合函数图像是 ( )

A. 类似递增的双曲线. B. 类似递增的对数曲线.

C. 类似递减的指数曲线. D. 是一条 S 形曲线.

\section*{三、解答题}

13. 已知函数 $f\left( x\right)  = \left\{  \begin{array}{l} \frac{2}{x},x < 0, \\   - x,0 \leq  x \leq  2, \\  \frac{1}{2}x - 3,x > 2. \end{array}\right.$

(1)求 $f\left( 0\right)$ 和 $f\left\lbrack  {f\left( 1\right) }\right\rbrack$ 的值；

\begin{center}
\includegraphics[width=0.3\textwidth]{images/bo_d7tk7oc91nqc73875fu0_34_1055_1435_422_378_0.jpg}
\end{center}
 

(第 13 题图)

(2)若 $f\left( m\right)  =  - 1$ ，求 $m$ 的值；

(3)在给定的坐标系中，作出函数 $y = f\left( x\right)$ 的图像.

14. 函数 $f\left( x\right)  = \sqrt{2 - \frac{x + 3}{x + 1}}$ 的定义域为 $A,g\left( x\right)  = \lg \left\lbrack  {\left( {x - a - 1}\right) \left( {{2a} - x}\right) }\right\rbrack  \left( {a < 1}\right)$ 的定义域为 $B$ .

(1)求 $A$ ；

(2)若 $B \subseteq  A$ ，求实数 $a$ 的取值范围.

15. 已知函数 $f\left( {1 - {2x}}\right)$ 的定义域为 $A = \left\lbrack  {\frac{1}{2},1}\right\rbrack$ .

(1)求 $f\left( x\right)$ 的定义域 $B$ ；

(2)对于(1)中的集合 $B$ ，若存在 $x \in  B$ ，使得 $a > {x}^{2} - x + 1$ 成立，求实数 $a$ 的取值范围.

16. 已知函数 $f\left( x\right)  = \frac{{x}^{2} - {3x}}{\sqrt{{2x} + 1}}$ 和 $g\left( x\right)  = \frac{\sqrt{{2x} + 1}}{x - 3}$ ,设 $h\left( x\right)  = f\left( x\right)  \cdot  g\left( x\right)$ .

(1)若函数 $H\left( x\right)  = x$ ，试判断 $y = h\left( x\right)$ 与 $y = H\left( x\right)$ 是否为同一函数，并说明理由；

(2)求 $G\left( x\right)  = h\left( x\right)  - \sqrt{{2h}\left( x\right)  + 1}$ 的值域.

\section*{$\text{ § }{3.2}$ 函数的奇偶性}

\section*{典型例题精讲}

【例 1】已知函数 $f\left( x\right)  = a{x}^{2} + \frac{b}{x}\left( {x \neq  0}\right.$ ,常数 $\left. {a\text{ 、 }b \in  \mathbf{R}}\right)$ .

(1)当 $a = 0,b = 1$ 时，解关于 $x$ 的不等式: $f\left( x\right)  > k\left( {k \in  \mathbf{R}}\right)$ ；

(2)讨论函数 $f\left( x\right)$ 的奇偶性,并说明理由.

【例 2】已知函数 $f\left( x\right)  = {x}^{2} + 2\left| {x - a}\right| \left( {a \in  \mathbf{R}}\right)$ .

( 1 )若函数 $f\left( x\right)$ 为偶函数，求 $a$ 的值；

( 2 )若函数 $g\left( x\right)  = {af}\left( x\right)  + 2$ 的最小值为 8，求 $a$ 的值.

【例 3】已知函数 $f\left( x\right)  = \left\{  {\begin{array}{l} {x}^{2} + {ax},x \geq  0, \\  {x}^{2} + {2x},x < 0 \end{array}\left( {a \in  \mathbf{R}}\right) }\right.$ .

(1)若 $f\left( x\right)$ 是偶函数，求 $a$ 的值；

(2)若 $f\left( x\right)$ 在 $x =  - 1$ 时取到最小值，求 $a$ 的取值范围.

【例 4】已知 $f\left( x\right)$ 是定义在 $\mathbf{R}$ 上的函数,对任意 $x$ 、 $y \in  \mathbf{R}$ ，都有 $f\left( {x + y}\right)  = f\left( x\right)  + f\left( y\right)$ ，且当 $x > 0$ 时, $f\left( x\right)  < 0,f\left( {-1}\right)  = 2$ .

(1)求证: $f\left( x\right)$ 为奇函数；

(2)求证: $f\left( x\right)$ 是 $\mathbf{R}$ 上的减函数；

(3)求 $f\left( x\right)$ 在 $\left\lbrack  {-2,4}\right\rbrack$ 上的最值.

【例 5】对于函数 $f\left( x\right)$ ,若在定义域内存在实数 ${x}_{0}$ ,满足 $f\left( {-{x}_{0}}\right)  =  - f\left( {x}_{0}\right)$ ,则称 $f\left( x\right)$ 为“ $M$ 类函数”. (1)设 $f\left( x\right)  = {2}^{x} + m$ 是定义在 $\left\lbrack  {-1,1}\right\rbrack$ 上的“ $M$ 类函数”，求实数 $m$ 的最小值；

(2)若 $f\left( x\right)  = \left\{  \begin{array}{l} {\log }_{2}\left( {{x}^{2} - {2mx}}\right) ,x \geq  2, \\   - 3,x < 2 \end{array}\right.$ 为定义域上的“ $M$ 类函数”，求实数 $m$ 的取值范围.

\section*{四基达标精练}

\section*{一、填空题}

1. 已知 $y = f\left( x\right)$ 是奇函数，当 $x \geq  0$ 时， $f\left( x\right)  =  - {x}^{2}$ ，则 $f\left( {-1}\right)$ 的值是\_\_\_\_\_.

2. 若定义在 $\left( {-4,{\log }_{2}a}\right)$ 上的函数 $y = f\left( x\right)$ 是偶函数，则实数 $a =$ \_\_\_\_\_.

3. 若函数 $y = {x}^{2} + \left( {a - 3}\right) x + 4$ 为偶函数，则 $a =$ \_\_\_\_\_.

4. 已知函数 $f\left( x\right)$ 是定义在 $\mathbf{R}$ 上的偶函数. 当 $x \in  \left( {-\infty ,0}\right)$ 时, $f\left( x\right)  = {2}^{x} - {x}^{2}$ ,则当 $x \in  \left( {0, + \infty }\right)$ 时, $f\left( x\right)  =$

5. 函数 $f\left( x\right)  = a{x}^{3} - {bx} - \tan x + 2$ ，若 $f\left( m\right)  = 1$ ，则 $f\left( {-m}\right)  =$ \_\_\_\_\_.

6. 设 $y = f\left( x\right)$ 为定义在 $\mathbf{R}$ 上的奇函数，当 $x \geq  0$ 时， $f\left( x\right)  = {2}^{x} + {2x} + b$ ( $b$ 为常数),则 $f\left( {-1}\right)  =$

7. 设 $f\left( x\right)$ 是定义域为 $\mathbf{R}$ 的奇函数, $g\left( x\right)$ 是定义域为 $\mathbf{R}$ 的偶函数,若函数 $f\left( x\right)  + g\left( x\right)$ 的值域为 $\lbrack 1,3)$ , 则函数 $f\left( x\right)  - g\left( x\right)$ 的值域为\_\_\_\_\_.

8. 函数 $f\left( x\right)  = \frac{\sqrt{a - {x}^{2}}}{\left| {x + 1}\right|  - 1}$ 是奇函数，则实数 $a$ 的取值范围是\_\_\_\_\_.

9. 已知函数 $f\left( x\right)  = {x}^{2} - \cos x\left( {-\frac{\pi }{2} \leq  x \leq  \frac{\pi }{2}}\right)$ ,如果 $f\left( x\right)  < f\left( \frac{\pi }{3}\right)$ ,那么实数 $x$ 的取值范围是\_\_\_\_\_.

\section*{二、选择题}

10. 函数 $f\left( x\right)  =  - x + \frac{1}{2x}$ 的图像关于 ( )

A. $y$ 轴对称. B. 直线 $y =  - x$ 对称. C. 坐标原点对称. D. 直线 $y = x$ 对称.

11. 若 $f\left( x\right)$ 为奇函数，且 ${x}_{0}$ 是 $y = f\left( x\right)  - {{e}}^{x}$ 的一个零点，则 $- {x}_{0}$ 一定是下列哪个函数的零点( )

D. $y = f\left( {-x}\right) {{e}}^{x} + 1$ .


\begin{center}
\includegraphics[width=0.2\textwidth]{images/bo_d7tk7oc91nqc73875fu0_39_1176_155_266_243_0.jpg}
\end{center}
 

(第 12 题图)

12. 如图，在平面直角坐标系内有一个边长为 $a$ 、中心在原点 $O$ 的正六边形 ${ABCDEF},{AB}//{Ox}$ . 直线 $l : y = {kx} + t\left( {k\text{ 为常数 }}\right)$ 与正六边形的边交于 $M\text{ 、 }N$ 两点, 记 $\bigtriangleup {OMN}$ 的面积为 $S$ ,则函数 $S = f\left( t\right)$ 的奇偶性为 ( )

A. 偶函数.

B. 奇函数.

C. 不是奇函数，也不是偶函数.

D. 奇偶性与 $k$ 有关.

\section*{三、解答题}

13. 已知 $y = f\left( x\right)$ 是定义在 $\mathbf{R}$ 上的奇函数,当 $x \geq  0$ 时, $f\left( x\right)  =  - {x}^{2} + {2x}$ .

(1)求函数 $f\left( x\right)$ 的解析式；

(2)若函数 $f\left( x\right)$ 在 $\left\lbrack  {-1,a - 2}\right\rbrack$ 上单调递增，求实数 $a$ 的取值范围.

14. 已知函数 $f\left( x\right)  = \ln \left( {3 + x}\right)  + \ln \left( {3 - x}\right)$ .

(1)判断函数 $f\left( x\right)$ 的奇偶性,并证明你的结论;

(2)若 $f\left( {{2m} - 1}\right)  < f\left( m\right)$ ，求 $m$ 的取值范围.

15. 已知函数 $f\left( x\right)  = \frac{\left( {x + 1}\right) \left( {x + a}\right) }{{x}^{2}}$ 为偶函数.

(1)求实数 $a$ 的值；

(2)当 $x \in  \left\lbrack  {\frac{1}{m},\frac{1}{n}}\right\rbrack  \left( {m\text{ 、 }n > 0}\right)$ 时,若函数 $f\left( x\right)$ 的值域为 $\left\lbrack  {2 - {3m},2 - {3n}}\right\rbrack$ ，求 $m\text{ 、 }n$ 的值.

16. 已知函数 $f\left( x\right)  = \frac{-{2}^{x} + 1}{{2}^{x + 1} + 2}$ .

(1)判断函数 $f\left( x\right)$ 的奇偶性,并证明;

(2)若不等式 $f\left( x\right)  > {\log }_{9}\left( {{2c} - 1}\right)$ 有解，求 $c$ 的取值范围.

\section*{§ 3.3 函数的单调性}

\section*{典型例题精讲}

【例 1】已知函数 $f\left( x\right)  = \frac{2x}{x - 1}$ .

(1)求 $f\left( x\right)$ 的定义域、值域并写出其单调区间及单调性(不要求证明)；

(2)判断并用定义证明函数 $g\left( x\right)  = {xf}\left( x\right)$ 在区间 $\left( {0,1}\right)$ 上的单调性.

【例 2】设 $a$ 为实数，函数 $f\left( x\right)  = {\left( x - a\right) }^{2} + \left| {x - a}\right|  - a\left( {a - 1}\right)$ .

(1)若 $f\left( 0\right)  \leq  1$ ，求 $a$ 的取值范围；

(2)讨论 $f\left( x\right)$ 的单调性.

【例 3】已知函数 $f\left( x\right)  = \left| {x - a}\right|  - \frac{a}{2}\ln x\left( {a \in  \mathbf{R},x > 0}\right) ,g\left( x\right)  = \left| {x - 2}\right|$ .

(1)若 $a = 0$ ，求函数 $F\left( x\right)  = f\left( x\right)  - g\left( x\right)$ 的值域；

(2)求函数 $f\left( x\right)$ 的单调区间.

【例 4】已知函数 $f\left( x\right)  = a{x}^{3} + x - a\left( {a \in  \mathbf{R}}\right) ,g\left( x\right)  = \frac{x}{1 - {x}^{3}}$ .

(1)判断 $g\left( x\right)$ 在 $\left( {-1, - \frac{\sqrt[3]{4}}{2}}\right\rbrack$ 和 $\left\lbrack  {-\frac{\sqrt[3]{4}}{2},1}\right)$ 的单调性,并说明理由;

(2)证明:函数 $f\left( x\right)$ 存在零点 $q$ ，使得 $a = q + {q}^{4} + {q}^{7} + \cdots  + {q}^{{3n} - 2} + \cdots$ 成立的充要条件是 $a \geq  \frac{-\sqrt[3]{4}}{3}$ .

【例 5】已知函数 $f\left( x\right)  = a\left( {x + \frac{1}{x}}\right)  - \left| {x - \frac{1}{x}}\right| \left( {x > 0}\right) ,a \in  \mathbf{R}$ .

(1)若 $a = \frac{1}{2}$ ，求 $y = f\left( x\right)$ 的单调区间；

(2)若关于 $x$ 的方程 $f\left( x\right)  = t$ 有四个不同的解,求实数 $a\text{ 、 }t$ 应满足的条件.

\section*{四基达标精练}

\section*{一、填空题}

1. 函数 $f\left( x\right)  = {x}^{2} - {3x} - 1$ 的严格递增区间是\_\_\_\_\_.

2. 函数 $y = \sqrt{-{x}^{2} - {2x} + 8}$ 的单调递减区间为\_\_\_\_\_.

3. 若函数 $f\left( x\right)  = \frac{5}{{3x} - 2}$ 的严格递减区间是\_\_\_\_\_.

4. 若函数 $f\left( x\right)  = {x}^{2} + 2\left( {a - 1}\right) x + 2$ 在 $( - \infty ,1\rbrack$ 上递减，则实数 $a$ 的取值范围是\_\_\_\_\_.

5. 若函数 $f\left( x\right)  = {x}^{2} + 2\left( {a - 1}\right) x + 2$ 在 $\left\lbrack  {-1,1}\right\rbrack$ 上是增函数，则实数 $a$ 的取值范围是\_\_\_\_\_.

6. 能说明“若 $f\left( x\right)  > f\left( 0\right)$ 对任意的 $x \in  (0,2\rbrack$ 都成立，则 $f\left( x\right)$ 在 $\left\lbrack  {0,2}\right\rbrack$ 上是增函数”为假命题的一个函数是\_\_\_\_\_.

7. 设 $a$ 为常数，函数 $f\left( x\right)  = {x}^{2} - {4x} + 3$ ，若 $f\left( {x + a}\right)$ 在 $\lbrack 0, + \infty )$ 上是增函数，则 $a$ 的取值范围是\_\_\_\_\_.

8. 若函数 $f\left( x\right)  = {x}^{2} + \left( {{2a} + 1}\right) \left| x\right|  + 1$ 的定义域被分成四个不同的单调区间，则实数 $a$ 的取值范围是\_\_\_\_\_.

9. 设函数 $y = f\left( x\right)$ 定义在 $\mathbf{R}$ 上,对于给定的正数 $K$ ,定义函数 ${f}_{K}\left( x\right)  = \left\{  \begin{array}{l} f\left( x\right) ,f\left( x\right)  \leq  K, \\  K,f\left( x\right)  > K. \end{array}\right.$ 若函数 $f\left( x\right)  = {2}^{-\left| x\right| }$ ,当 $K = \frac{1}{2}$ 时,函数 ${f}_{K}\left( x\right)$ 的严格递增区间为\_\_\_\_\_.

\section*{二、选择题}

10. 已知函数 $f\left( x\right)  = \frac{{ax} + 1}{x + 2}\left( {a \in  \mathbf{R}}\right)$ ,则 “ $f\left( 2\right)  < f\left( 3\right)$ ” 是 “ $f\left( x\right)$ 在区间 $\left( {-2, + \infty }\right)$ 上单调递增” 的

( )

B. 必要非充分条件. C. 充要条件. D. 既非充分也非必要条件.

11. 若函数 $y = f\left( x\right)$ 在区间 $I$ 上是增函数,且函数 $y = \frac{f\left( x\right) }{x}$ 在区间 $I$ 上是减函数,则称函数 $f\left( x\right)$ 是区间 $I$ 上的“ $H$ 函数”,对于命题: \ding{172} 函数 $f\left( x\right)  =  - x + 2\sqrt{x}$ 是 $\left( {0,1}\right)$ 上的“ $H$ 函数”;  \ding{173} 函数 $g\left( x\right)  = \frac{2x}{1 - {x}^{2}}$ 是 $\left( {0,1}\right)$ 上的“ $H$ 函数”; 下列判断正确的是 ( )

A.  \ding{172} 和 \ding{173} 均为真命题. B.  \ding{172} 为真命题， \ding{173} 为假命题.

C.  \ding{172} 为假命题， \ding{173} 为真命题. D.  \ding{172} 和 \ding{173} 均为假命题.

12. 命题 $p$ : 存在 $a \in  \mathbf{R}$ 且 $a \neq  0$ ,对于任意的 $x \in  \mathbf{R}$ ,使得 $f\left( {x + a}\right)  < f\left( x\right)  + f\left( a\right)$ ; 命题 ${q}_{1} : f\left( x\right)$ 单调递减且 $f\left( x\right)  > 0$ 恒成立; 命题 ${q}_{2} : f\left( x\right)$ 单调递增,存在 ${x}_{0} < 0$ 使得 $f\left( {x}_{0}\right)  = 0$ ,则下列说法正确的是 ( )

A. 只有 ${q}_{1}$ 是 $p$ 的充分条件. B. 只有 ${q}_{2}$ 是 $p$ 的充分条件.

C. ${q}_{1}\text{ 、 }{q}_{2}$ 都是 $p$ 的充分条件. D. ${q}_{1}\text{ 、 }{q}_{2}$ 都不是 $p$ 的充分条件.

\section*{三、解答题}

13. 已知函数 $f\left( x\right)  = 2 + \frac{1}{a} - \frac{1}{{a}^{2}x}\left( {a \in  \mathbf{R},a \neq  0}\right)$ .

(1)设 ${mn} > 0$ ，判断函数 $f\left( x\right)$ 在 $\left\lbrack  {m,n}\right\rbrack$ 上的单调性，并加以证明；

(2)若不等式 ${a}^{2}f\left( x\right)  \leq  {2x}$ 对 $x \geq  1$ 恒成立，求实数 $a$ 的取值范围.

14. 已知函数 $f\left( x\right)  = \frac{1 + {x}^{2}}{x}$ .

(1)求函数的定义域和值域；

(2)求函数的单调区间.


15. 某群体的人均通勤时间，是指单日内该群体中成员从居住地到工作地的平均用时. 某地上班族 $S$ 群体中的成员仅以自驾或公交方式通勤. 分析显示: 当 $S$ 中 $x\% \left( {0 < x < {100}}\right)$ 的成员自驾时,自驾群体的人均通勤时间为 $f\left( x\right)  = \left\{  \begin{array}{l} {30},0 < x \leq  {30}, \\  {2x} + \frac{1800}{x} - {90},{30} < x < {100} \end{array}\right.$ (单位: 分钟),而公交群体的人均通勤时间不受 $x$ 影响,恒为 40 分钟. 试根据上述分析结果回答下列问题:

(1)当 $x$ 在什么范围内时,公交群体的人均通勤时间少于自驾群体的通勤时间？

(2)求该地上班族 $S$ 群体的人均通勤时间 $g\left( x\right)$ 的表达式:讨论 $g\left( x\right)$ 的单调性，并说明其实际意义.

16. 已知函数 $f\left( x\right)  = \frac{{2}^{x} + a}{{2}^{x} + 1},x \in  \mathbf{R}$ .

(1)证明:当 $a > 1$ 时，函数 $y = f\left( x\right)$ 在 $\mathbf{R}$ 上是严格减函数；

(2)求函数 $y = f\left( x\right)$ 的值域.

\section*{§ 3.4 函数的最值和值域}

\section*{典型例题精讲}

【例 1】已知函数 $f\left( x\right)  = x + b$ 在区间 $\left\lbrack  {1,2}\right\rbrack$ 上的最大值为 2,且函数 $g\left( x\right)  = \frac{x + b}{{x}^{2} + 1}$ .

(1)求出 $b$ 的值，判断 $g\left( x\right)$ 的奇偶性，并证明你的结论；

(2)设 $h\left( x\right)  =  - x - {2c}$ ，对任意的 ${x}_{1} \in  \mathbf{R}$ ，总存在 ${x}_{2} \in  \left\lbrack  {-2,2}\right\rbrack$ ，使得 $g\left( {x}_{1}\right)  = h\left( {x}_{2}\right)$ 成立，求实数 $c$ 的取值范围.

【例 2】已知函数 $f\left( x\right)$ 的定义域为 $D$ ,对于给定的正整数 $k$ ,若存在 $\left\lbrack  {a,b}\right\rbrack   \subseteq  D$ ,使得函数 $f\left( x\right)$ 满足: 函数 $f\left( x\right)$ 在 $\left\lbrack  {a,b}\right\rbrack$ 上是单调函数,且 $f\left( x\right)$ 的最小值为 ${ka}$ ,最大值为 ${kb}$ ,则称函数 $f\left( x\right)$ 是 “倍缩函数”,区间 $\left\lbrack  {a,b}\right\rbrack$ 是函数 $f\left( x\right)$ 的“ $k$ 倍值区间”.

(1)判断函数 $f\left( x\right)  = {x}^{3}$ 是否是“倍缩函数”? 请说明理由;

(2)设函数 $g\left( x\right)  = \frac{8x}{4{x}^{2} + 1},x \in  \left\lbrack  {0,\frac{1}{2}}\right\rbrack$ ，若函数 $g\left( x\right)$ 存在 “ $k$ 倍值区间”，求 $k$ 的值.


【例 3】已知 $f\left( x\right)  = {\log }_{a}\left( {{ax} - 2}\right)$ (其中 $a > 0$ 且 $a \neq  1$ ).

(1)若 $a = 2$ ， $f\left( x\right)  < 2$ ，求实数 $x$ 的取值范围；

(2)若 $x \in  \left\lbrack  {4,6}\right\rbrack$ ， $f\left( x\right)$ 的最大值大于 1，求 $a$ 的取值范围.

【例 4】已知函数 $f\left( x\right)  = \sqrt{\left| {x + 2}\right|  + \left| {x - 4}\right|  - m}$ 的定义域为 $\mathbf{R}$ .

(1)求实数 $m$ 的范围；

(2)若 $m$ 的最大值为 $n$ ，当正数 $a$ 、 $b$ 满足 $\frac{4}{a + {5b}} + \frac{1}{{3a} + {2b}} = n$ 时，求 ${4a} + {7b}$ 的最小值.

【例 5】已知 $a \in  \mathbf{R}$ ,函数 $f\left( x\right)  = x\left| {x - a}\right|$ .

(1)判断函数 $f\left( x\right)$ 的奇偶性,请说明理由;

(2)设 $a > 0$ ，求函数 $f\left( x\right)$ 在区间 $\left\lbrack  {1,3}\right\rbrack$ 上的最小值.

\section*{四基达标精练}

\section*{一、填空题}

1. 函数 $f\left( x\right)  = \left| {x - 3}\right|  + \left| {x + 1}\right|$ 的最小值是\_\_\_\_\_.

2. 函数 $y = \frac{{x}^{2} + 5}{\sqrt{{x}^{2} + 4}}$ 的最小值为\_\_\_\_\_.

3. 函数 $y = \sqrt{x\left( {3 - {2x}}\right) }$ 的值域是\_\_\_\_\_.

4. 若函数 $f\left( x\right)  = {\left( \frac{1}{3}\right) }^{{ax} + {a}^{2}}$ 在 $\lbrack 1, + \infty )$ 上有最大值 $\frac{1}{9}$ ，则实数 $a =$ \_\_\_\_\_.

5. 定义 $\min \{ a,b\}  = \left\{  \begin{array}{l} a,a \leq  b, \\  b,a > b, \end{array}\right.$ 则函数 $y = \min \{ {3x} - 1, - x + 2\}$ 的最大值为\_\_\_\_\_.

6. 函数 $y = \sqrt{{2x} + 1} - \sqrt{2 - x}$ 的值域是\_\_\_\_\_.

7. 若函数 $f\left( x\right)  = x + \frac{m}{x}\left( {1 < x < 2}\right)$ 存在最大值或最小值，则实数 $m$ 的取值范围是\_\_\_\_\_.

8. 函数 $f\left( x\right)  = \left| {x - 1}\right|  - \ln x$ 的最小值为\_\_\_\_\_.

9. 设 ${M}_{I}$ 表示函数 $f\left( x\right)  = \left| {{x}^{2} - {4x} + 2}\right|$ 在闭区间 $I$ 上的最大值. 若 $a > 0$ ，且 ${M}_{\left\lbrack  0,a\right\rbrack  } \geq  2{M}_{\left\lbrack  a,2a\right\rbrack  }$ ，则正实数 $a$ 的取值范围是\_\_\_\_\_.

\section*{二、选择题}

10. 已知 $f\left( x\right)$ 是定义在 $\left\lbrack  {0,1}\right\rbrack$ 上的函数，那么 “函数 $f\left( x\right)$ 在 $\left\lbrack  {0,1}\right\rbrack$ 上单调递增” 是 “函数 $f\left( x\right)$ 在 $\left\lbrack  {0,1}\right\rbrack$ 上的最大值为 $f\left( 1\right)$ ”的 ( )

A. 充分非必要条件. B. 必要非充分条件. C. 充要条件. D. 既非充分也非必要条件.

11. 关于 “函数 $f\left( x\right)  = \frac{x - \sqrt{13}}{x - \sqrt{14}},x \in  \left( {-\infty ,\sqrt{14}}\right)  \cup  \left( {\sqrt{14}, + \infty }\right)$ 的最大、最小值与函数 $g\left( x\right)  =  \frac{x - \sqrt{13}}{x - \sqrt{14}},x \in  \mathbf{Z}$ 的最大、最小值”,下列说法中正确的是 ( )

A. $f\left( x\right)$ 有最大、最小值, $g\left( x\right)$ 有最大、最小值. B. $f\left( x\right)$ 有最大、最小值, $g\left( x\right)$ 无最大、最小值.

C. $f\left( x\right)$ 无最大、最小值, $g\left( x\right)$ 有最大、最小值. D. $f\left( x\right)$ 无最大、最小值, $g\left( x\right)$ 无最大、最小值.


12. 若函数 $f\left( x\right)  = \left| {x - a}\right|  + b\left( {-1 \leq  x \leq  1}\right)$ 的最大值是 $M$ ,最小值是 $m$ ,则 $m - n$ ( )

A. 与 $a$ 有关,且与 $b$ 有关. B. 与 $a$ 有关,但与 $b$ 无关.

C. 与 $a$ 无关，且与 $b$ 无关. D. 与 $a$ 无关，但与 $b$ 有关.

\section*{三、解答题}

13. 设函数 $f\left( x\right)  = a{x}^{2} + \left( {b - 2}\right) x + 3\left( {a \neq  0}\right)$ .

(1)若关于 $x$ 的不等式 $f\left( x\right)  > 0$ 的解集为 $\left( {-1,1}\right)$ ，求实数 $a$ 和 $b$ 的值；

(2)若 $f\left( 1\right)  = 3$ ，求函数 $f\left( x\right)$ 在区间 $\left\lbrack  {1,3}\right\rbrack$ 上的最小值.

14. 已知二次函数 $f\left( x\right)$ 满足 $f\left( x\right)  = f\left( {4 - x}\right) ,f\left( 2\right)  = f\left( 1\right)  - 1$ ,若不等式 $f\left( x\right)  \leq   - {2x} + 2$ 有唯一实数解.

(1)求函数 $f\left( x\right)$ 的解析式；

(2)若函数 $f\left( x\right)$ 在 $\left\lbrack  {t,t + 1}\right\rbrack$ 上的最小值为 $g\left( t\right)$ . 求 $g\left( t\right)$ 的解析式.

15. 已知函数 $f\left( x\right)  = \left| {a{x}^{2} - b}\right|  + a{x}^{2} + {bx}\left( {a > 0}\right)$ .

(1)若 $a = b = 1$ ，求函数 $f\left( x\right)$ 的最小值；

(2)若 $b \leq  0$ ，函数 $f\left( x\right)$ 存在两个不同的零点 ${x}_{1}$ 和 ${x}_{2}$ ，求 $\frac{{x}_{2}}{{x}_{1}} + \frac{{x}_{1}}{{x}_{2}}$ 的取值范围.

16. 已知函数 $f\left( x\right)  = \left| {x - 1}\right|  + {\left( x - 1\right) }^{2} + \frac{\sqrt{3}}{3}$ .

(1)求 $g\left( x\right)  = f\left( x\right)  + \left| {f\left( x\right)  - \frac{\sqrt{3}}{3}}\right|$ 的最小值；

(2)若存在实数 $x$ ，使得 $\frac{\sqrt{3}}{{4f}\left( x\right) } + \frac{{4f}\left( x\right) }{\sqrt{3}} \leq  \left| {{m}^{2} - {2m} - 1}\right|$ 成立，求 $m$ 的取值范围.

\section*{$\text{ § }{3.5}$ 函数的零点}

\section*{典型例题精讲}

【例 1】已知函数 $f\left( x\right)  = \frac{a}{x - 1},a$ 为常数.

(1)当 $a > 0$ 时，根据定义证明函数 $f\left( x\right)$ 在区间 $\left( {1, + \infty }\right)$ 上单调递减；

(2)设 $g\left( x\right)  = f\left( x\right)  + {ax} - 3$ ,若 $g\left( x\right)$ 在 $\left( {1, + \infty }\right)$ 上存在两个零点,求实数 $a$ 的取值范围.

【例 2】已知函数 $f\left( x\right)  = {x}^{3} + {3a}{x}^{2} + {bx} + {a}^{2}$ 在 $x =  - 1$ 时有极值 0 .

(1)求函数 $f\left( x\right)$ 的解析式；

(2)记 $h\left( x\right)  = f\left( x\right)  - m + 1$ ，若函数 $h\left( x\right)$ 有三个零点，求实数 $m$ 的取值范围.

【例 3】设函数 $f\left( x\right)  = \left( {x - a}\right) \left( {x - b}\right) \left( {x - c}\right) ,a\text{ 、 }b\text{ 、 }c \in  \mathbf{R},{f}^{\prime }\left( x\right)$ 为 $f\left( x\right)$ 的导函数.

(1)若 $a = b = c$ ， $f\left( 4\right)  = 8$ ，求 $a$ 的值；

(2)若 $a \neq  b$ ， $b = c$ ，且 $f\left( x\right)$ 和 ${f}^{\prime }\left( x\right)$ 的零点均在集合 $\{  - 3,1,3\}$ 中，求 $f\left( x\right)$ 的极小值.

【例 4】已知函数 $f\left( x\right)  = \frac{{{e}}^{x} - {{e}}^{-x}}{2},g\left( x\right)  = \frac{{{e}}^{x} + {{e}}^{-x}}{2}$ ，其中 ${e}$ 是自然对数的底数.

(1)求证: $g\left( {2x}\right)  = {\left\lbrack  f\left( x\right) \right\rbrack  }^{2} + {\left\lbrack  g\left( x\right) \right\rbrack  }^{2}$ ；

(2)求函数 $h\left( x\right)  = g\left( {2x}\right)  - \frac{7}{2}g\left( x\right)$ 的零点.

【例 5】已知函数 $f\left( x\right)  = {\log }_{3}\left( {{9}^{x} + 1}\right)  - {kx}$ (其中 $k \in  \mathbf{R}$ ).

(1)若 $k = 2$ 且方程 $f\left( x\right)  - a + 1 = 0$ 有解,求实数 $a$ 的取值范围；

(2)若 $f\left( x\right)$ 是偶函数，讨论函数 $g\left( x\right)  = {\left( 3k\right) }^{f\left( x\right) } - \left( {m \cdot  {3}^{x} - \frac{4}{3}m}\right) \left( {m > 0}\right)$ 的零点情况.

\section*{四基达标精练}

\section*{一、填空题}

1. 函数 $f\left( x\right)  = {{e}}^{x} + \ln x$ 的零点个数为\_\_\_\_\_.

2. 已知函数 $f\left( x\right)  = x + {2}^{x}$ 的零点在区间 $\left( {n,n + 1}\right)$ 内, $n \in  \mathbf{Z}$ ,则 $n =$ \_\_\_\_\_.

3. 若函数 $f\left( x\right)  = {2x} + m$ 在 $\left( {1,2}\right)$ 上存在一个零点，则实数 $m$ 的取值范围是\_\_\_\_\_.

4. 已知函数 $h\left( x\right)$ 是奇函数，且 $f\left( x\right)  = h\left( x\right)  + 2$ ，若 $x = 2$ 是函数 $y = f\left( x\right)$ 的一个零点，则 $f\left( {-2}\right)  =$

5. 已知二次函数 $y =  - {x}^{2} + {bx} + c$ 的零点为 -2 和 1，则关于 $x$ 的不等式 ${x}^{2} + {bx} - c > 0$ 的解集为\_\_\_\_\_.

6. 已知函数 $f\left( x\right)  = {x}^{3} - 7{x}^{2} + {ax} - 8$ 有 3 个不同的零点分别为 ${x}_{1}$ 、 ${x}_{2}$ 、 ${x}_{3}$ ，且 ${x}_{1}$ ， ${x}_{2}$ ， ${x}_{3}$ 依次成等比数列，则实数 $a$ 的值为\_\_\_\_\_.

7. 若函数 $f\left( x\right)  = x - a\ln x$ 在定义域上恰有两个零点，则实数 $a$ 的取值范围是\_\_\_\_\_.

8. 已知定义域为 $\mathbf{R}$ 的函数 $f\left( x\right)  = \left\{  \begin{array}{l} {x}^{2} - {2x} + 2,0 < x \leq  2, \\  {\log }_{\frac{1}{2}}\left( {{2x} - \frac{15}{4}}\right) ,x > 2, \end{array}\right.$ 且满足 $f\left( {-x}\right)  =  - f\left( x\right) \left( {x \in  \mathbf{R}}\right)$ ,若方程 $f\left( x\right)  = m\left( {m > 0}\right)$ 的解为 ${x}_{1}\text{ 、 }{x}_{2}\text{ 、 }{x}_{3}\text{ 、 }{x}_{4}$ ,则 ${x}_{1} + {x}_{2} + {x}_{3} + {x}_{4}$ 的取值范围为\_\_\_\_\_.

9. 设 $a > 0$ ，函数 $f\left( x\right)  = \sqrt{\left| x - a\right| } + \sqrt{\left| x + a\right| }$ ，若 $g\left( x\right)  = f\left( x\right)  - b$ 恰有三个不同的零点，且 $b$ 是其中的一个零点，则实数 $b$ 的值为\_\_\_\_\_.

\section*{二、选择题}

10. 已知 $f\left( x\right)  = {x}^{2} - {ax} + a$ ，则 “ $a > 4$ ” 是 “ $f\left( x\right)$ 有两个不同的零点” 的 ( )

A. 充分非必要条件. B. 必要非充分条件. C. 充要条件. D. 既非充分也非必要条件.

11. 下列函数中不能用二分法求零点的是 ( )

A. $y = {3x} - 1$ . B. $y = {x}^{3}$ . C. $y = \left| x\right|$ . D. $y = \ln x$ .

12. 关于函数 $f\left( x\right)  = \frac{1}{4}\left( {x - \frac{1}{x}}\right)  - \ln x$ 的零点,下列说法正确的是 ( )

A. 函数 $f\left( x\right)$ 有两个零点 ${x}_{1}\text{ 、 }{x}_{2}$ ,且 ${x}_{1}{x}_{2} = 1$ .

B. 函数 $f\left( x\right)$ 有两个零点 ${x}_{1}\text{ 、 }{x}_{2}$ ,且 ${x}_{1}{x}_{2} \neq  1$ .

C. 函数 $f\left( x\right)$ 有三个零点 ${x}_{1}\text{ 、 }{x}_{2}\text{ 、 }{x}_{3}$ ,且 ${x}_{1}{x}_{2}{x}_{3} = 1$ .

D. 函数 $f\left( x\right)$ 有三个零点 ${x}_{1}\text{ 、 }{x}_{2}\text{ 、 }{x}_{3}$ ,且 ${x}_{1}{x}_{2}{x}_{3} \neq  1$ .

\section*{三、解答题}

13. 已知 $p :$ 存在 $x \in  \left\lbrack  {0,2}\right\rbrack$ ,使得 ${{e}}^{x} - 1 \leq  m;q :$ 函数 $y = {x}^{2} - {2mx} + 1$ 有两个零点.

(1)若 $p$ 或 $q$ 为假命题，求实数 $m$ 的取值范围；

(2)若 $p$ 或 $q$ 为真命题，而 $p$ 且 $q$ 为假命题，求实数 $m$ 的取值范围.

14. 已知奇函数 $f\left( x\right)  = {x}^{3} + b{x}^{2} + {cx} + d$ 的三个零点 ${x}_{1}$ 、 ${x}_{2}$ 、 ${x}_{3}$ 满足 ${x}_{1}{x}_{2} + {x}_{2}{x}_{3} + {x}_{3}{x}_{1} =  - 2$ .

(1)求函数 $f\left( x\right)$ 的解析式;

(2)求出所有零点.


15. 设函数 $f\left( x\right)  = {x}^{3} - {3a}{x}^{2} + {3bx}$ 的图像与直线 ${12x} + y - 1 = 0$ 相切于点 $\left( {1, - {11}}\right)$ .

(1)求a、b的值；

(2)若函数 $g\left( x\right)  = f\left( x\right)  - c$ 有 3 个零点，求实数 $c$ 的取值范围.

16. 已知函数 $f\left( x\right)  = \frac{1 - {x}^{2}}{1 + {x}^{2}}$ .

(1)求函数 $f\left( x\right)$ 的零点；

(2)设函数 $f\left( x\right)$ 的值域为 $A$ ,若至少有两个不同的 $x \in  \left\lbrack  {\frac{\pi }{2},\pi }\right\rbrack$ ，使得 $\sin \left( {{\omega x} + \frac{\pi }{6}}\right)  \notin  A$ ，求正数 $\omega$ 的取值范围.

\section*{$\text{ § }{3.6}$ 函数的周期性与对称性}

\section*{典型例题精讲}

【例 1】我们知道,函数 $y = f\left( x\right)$ 的图像关于坐标原点成中心对称图形的充要条件是函数 $y = f\left( x\right)$ 为奇函数,有同学发现可以将其推广为: 函数 $y = f\left( x\right)$ 的图像关于点 $P\left( {a,b}\right)$ 成中心对称图形的充要条件是函数 $y = f\left( {x + a}\right)  - b$ 为奇函数.

(1)利用上述材料，求函数 $f\left( x\right)  = {x}^{3} - 3{x}^{2} + {6x} - 2$ 图像的对称中心；

(2)利用函数单调性的定义，证明函数 $f\left( x\right)  = {x}^{3} - 3{x}^{2} + {6x} - 2$ 在 $\mathbf{R}$ 上是严格增函数.

【例 2】定义在 $\mathbf{R}$ 上的奇函数 $f\left( x\right)$ 有最小正周期 4,且 $x \in  \left( {0,2}\right)$ 时, $f\left( x\right)  = \frac{{2}^{x}}{{4}^{x} + 1}$ .

(1)判断 $f\left( x\right)$ 在 $\left( {0,2}\right)$ 上的单调性,并求 $f\left( x\right)$ 在 $\left\lbrack  {-2,2}\right\rbrack$ 上的解析式;

(2)当λ为何值时，关于 $x$ 的方程 $f\left( x\right)  =$ λ在 $\left\lbrack  {2,6}\right\rbrack$ 上有实数解?

【例 3】记函数 $f\left( x\right)$ 的定义域为 $D$ . 如果存在实数 $a\text{ 、 }b$ 使得 $f\left( {a - x}\right)  + f\left( {a + x}\right)  = b$ 对任意满足 $a - x \in  D$ 且 $a + x \in  D$ 的 $x$ 恒成立,则称 $f\left( x\right)$ 为 $\Psi$ 函数.

(1)设函数 $f\left( x\right)  = \frac{1}{x} - 1$ ，试判断 $f\left( x\right)$ 是否为 $\Psi$ 函数，并说明理由；

(2)若 $h\left( x\right)$ 是定义在 $\mathbf{R}$ 上的 $\Psi$ 函数，且函数 $h\left( x\right)$ 的图像关于直线 $x = m$ ( $m$ 为常数，且 $m \neq  a$ )对称，试判断 $h\left( x\right)$ 是否为周期函数? 并证明你的结论.

【例 4】对于定义域为 $\mathbf{R}$ 的函数 $g\left( x\right)$ ,若存在正常数 $T$ ,使得 $\cos \left\lbrack  {g\left( x\right) }\right\rbrack$ 是以 $T$ 为周期的函数,则称 $g\left( x\right)$ 为余弦周期函数,且称 $T$ 为其余弦周期. 已知 $f\left( x\right)$ 是以 $T$ 为余弦周期的余弦周期函数.

(1)证明: $h\left( x\right)  = x + \sin \frac{x}{3}$ 是以 ${6\pi }$ 为周期的余弦周期函数;

(2)证明: “ ${u}_{0}$ 为方程 $\cos \left\lbrack  {f\left( x\right) }\right\rbrack   = 1$ 在 $\left\lbrack  {0,T}\right\rbrack$ 上的解”的充要条件是 “ ${u}_{0} + T$ 为方程 $\cos \left\lbrack  {f\left( x\right) }\right\rbrack   = 1$ 在 $\left\lbrack  {T,{2T}}\right\rbrack$ 上有解”.

【例 5】已知函数 $f\left( x\right)$ 的定义域为 $\mathbf{R}$ ,若存在常数 $T \neq  0$ ,使得 $f\left( x\right)  = {Tf}\left( {x + T}\right)$ 对任意的 $x \in  \mathbf{R}$ 成立,则称函数 $f\left( x\right)$ 是 $\Omega$ 函数.

(1)判断函数 $F\left( x\right)  = x,h\left( x\right)  = \sin {\pi x}$ 是否是 $\Omega$ 函数，请说明理由;

(2)若函数 $f\left( x\right)$ 是 $\Omega$ 函数，且 $f\left( x\right)$ 是偶函数，求证:函数 $f\left( x\right)$ 是周期函数；

(3)若函数 $f\left( x\right)  = \sin {kx}$ 是 $\Omega$ 函数. 求实数 $k$ 的取值范围.

\section*{四基达标精练}

\section*{一、填空题}

1. 函数 $f\left( x\right)  = {2}^{\left| x - 1\right| }$ 的对称轴是直线\_\_\_\_\_.

2. 函数 $f\left( x\right)  = \frac{1}{{3x} + 1} + 1$ 的对称中心是\_\_\_\_\_.

3. 函数 $f\left( x\right)  = \left| {x - 3}\right|  - \left| {x + 1}\right|$ 的对称中心是\_\_\_\_\_.

4.若函数 $f\left( x\right)$ 满足 $f\left( {4 - x}\right)  = f\left( x\right)$ 对任意 $x \in  \mathbf{R}$ 恒成立,当 $x \geq  2$ 时, $f\left( x\right)  = {2x} - 1$ ,则 $f\left( 0\right)  =$

5. 若 $f\left( x\right)$ 是定义在 $\mathbf{R}$ 上的奇函数，且 $f\left( {x + 1}\right)  =  - f\left( x\right)$ ，则 $f\left( {10}\right)  =$ \_\_\_\_\_.

6. 若函数 $y = f\left( x\right)$ 是偶函数，且 $f\left( {x + 3}\right)  = f\left( x\right)$ 对任意 $x \in  \mathbf{R}$ 恒成立，则除 $x = 0$ 外，函数必有一条对称轴是\_\_\_\_\_(写出一条即可).

7. 已知函数 $f\left( x\right)  = \ln \frac{ex}{2 - x}$ ,则 $f\left( x\right)  + f\left( {2 - x}\right)  =$ \_\_\_\_\_；进一步有 $\mathop{\sum }\limits_{{k = 1}}^{{{2n} - 1}}f\left( \frac{k}{n}\right)  =$ \_\_\_\_\_.

8. 已知 $y = f\left( x\right)$ 是 $\mathbf{R}$ 上的奇函数，且 $f\left( x\right)  + f\left( {2 - x}\right)  = 0$ ，当 $- 1 < x < 0$ 时， $f\left( x\right)  = \sqrt{-x}$ ，则 $f\left( {2023}\right)  + f\left( \frac{2023}{2}\right)  =$ \_\_\_\_\_.

9. 已知 $f\left( x\right)$ 是定义在 $\mathbf{R}$ 上的函数,且满足 $f\left( {x + 2}\right)  + f\left( {x + 2}\right) f\left( x\right)  + f\left( x\right)  = 1,f\left( 1\right)  = \frac{1}{2}$ ,则 $f\left( 1\right)  +  f\left( 3\right)  + f\left( 5\right)  + \cdots  + f\left( {2023}\right)  =$

\section*{二、选择题}

10. 已知 $a$ 、 $b$ 、 $c$ 是常数，则函数 $f\left( x\right)  = \frac{{2x} + b}{x + a}$ 和 $g\left( x\right)  = \frac{{2x} + c}{x + a}\left( {c \neq  b}\right)$ 有相同的 ( )

A. 对应法则. B. 奇偶性. C. 对称中心. D. 单调区间和单调性.

11. 已知 $f\left( x\right)$ 是定义在 $\mathbf{R}$ 上的偶函数，且以 2 为周期，则 “ $f\left( x\right)$ 在 $\left\lbrack  {0,1}\right\rbrack$ 上是增函数” 是 “ $f\left( x\right)$ 在 $\left\lbrack  {3,4}\right\rbrack$ 上是减函数”的 ( )

A. 充分非必要条件. B. 必要非充分条件. C. 充要条件. D. 既非充分也非必要条件.

12. 定义在 $\mathbf{R}$ 上的函数 $f\left( x\right)$ 满足 $f\left( {x + 3}\right)  + f\left( x\right)  = 0$ ,且函数 $f\left( {x - \frac{3}{2}}\right)$ 为奇函数. 给出以下 3 个命题  \ding{172} 函数 $f\left( x\right)$ 的周期是 6;  \ding{173} 函数 $f\left( x\right)$ 的图像关于点 $\left( {-\frac{3}{2},0}\right)$ 对称;  \ding{174} 函数 $f\left( x\right)$ 的图像关于 $y$ 轴对称,则其中真命题的个数是 ( )

A.3. B.2. C.1. D.0.

\section*{三、解答题}

13. 已知函数 $f\left( x\right)  = \left\{  \begin{array}{l}  - {2x}, - 1 \leq  x < 0, \\  {x}^{2} - 1,0 \leq  x \leq  1. \end{array}\right.$

(1)求函数 $f\left( x\right)$ 的值域；

(2)函数 $f\left( x\right)$ 的图像上是否存在关于坐标原点对称的点？若存在，求出这些点的坐标；若不存在，说明理由.

14. 对于函数 $f\left( x\right)$ ,若在定义域 $D$ 内存在实数 ${x}_{0}$ 满足 $f\left( {2 - {x}_{0}}\right)  =  - f\left( {x}_{0}\right)$ ,则称函数 $y = f\left( x\right)$ 为 “类对称函数”.

(1)判断函数 $g\left( x\right)  = {x}^{2} - {2x} + 1$ 是否为 “类对称函数” ？若是，求出所有满足条件的 ${x}_{0}$ 的值；若不是， 请说明理由;

(2)若函数 $h\left( x\right)  = {3}^{x} + t$ 为定义在 $\lbrack  - 1,3)$ 上的“类对称函数”,求实数 $t$ 的取值范围.

15. 已知真命题: “函数 $y = f\left( x\right)$ 的图像关于点 $P\left( {a\text{ 、 }b}\right)$ 成中心对称图形” 的充要条件为 “函数 $y = f(x +  a) - b$ 是奇函数”.

(1)将函数 $f\left( x\right)  = {x}^{3} - 3{x}^{2}$ 图像向左平移 1 个单位，再向上平移 2 个单位，所得图像对应的函数是 $y =  g\left( x\right)$ ,请利用题设中的真命题求函数 $f\left( x\right)$ 图像对称中心的坐标;

(2)求函数 $h\left( x\right)  = {\log }_{2}\frac{2x}{4 - x}$ 图像对称中心的坐标.

16. 已知函数 $f\left( x\right)$ 定义在 $\mathbf{R}$ 上，且对任意 $x \in  \mathbf{R}$ ，都有 $f\left( {x + 1}\right)  = f\left( {x - 1}\right) ,f\left( {1 - x}\right)  = f\left( {1 + x}\right)$ ，当 $x \in  \left\lbrack  {0,1}\right\rbrack$ 时， $f\left( x\right)  = {3}^{x} - {3}^{-x}$ .

(1)请指出 $f\left( x\right)$ 在区间 $\left\lbrack  {-1,1}\right\rbrack$ 上的奇偶性、单调区间、最大(小)值和零点;

(2)试证明 $f\left( x\right)$ 是周期函数，并求其在区间 $\left\lbrack  {{2k} - 1,{2k}}\right\rbrack  \left( {k \in  \mathbf{Z}}\right)$ 上的解析式.

\section*{第 4 章 幂函数、指数函数和对数函数(必修 1)}

\section*{知识梳理*}

\section*{一、幂、指数与对数}

1. 指数幂的拓展:正整数指数幂、整数指数幂、有理数指数幂、实数指数幂.

2. 幂的运算性质:对任意给定的正数 $a$ 、 $b$ 及实数 $s$ 、 $t$ ，

(1) ${a}^{s}{a}^{t} = {a}^{s + t};\left( 2\right) {\left( {a}^{s}\right) }^{t} = {a}^{st};\left( 3\right) {\left( ab\right) }^{t} = {a}^{t}{b}^{t}$ ;

${\left( 4\right) }^{ * }{\left( \sqrt[n]{a}\right) }^{n} = a\left( {n \in  \mathbf{N},n \geq  2}\right) ;{\left( 5\right) }^{ * }\sqrt[n]{{a}^{n}} = \left\{  {\begin{array}{l} a,n\text{ 为奇数, } \\  \left| a\right| ,n\text{ 为偶数. } \end{array}\left( {n \in  \mathbf{N},n \geq  2,a \in  \mathbf{R}}\right) }\right.$

3. 指数基本不等式:当 $a > 1$ 时，若 $s > 0$ ，则 ${a}^{s} > 1$ ；若 $s < 0$ ，则 $0 < {a}^{s} < 1$ .

(当 $0 < a < 1$ 时，若 $s > 0$ ，则 $0 < {a}^{s} < 1$ ；若 $s < 0$ ，则 ${a}^{s} > 1$ )

4. 对数的定义: 若 ${a}^{x} = N$ ,则 $x = {\log }_{a}N$ ,其中 $a > 0$ ,且 $a \neq  1,N > 0$ .

5. 对数的基本性质: 设 $a > 0$ ,且 $a \neq  1,M\text{ 、 }N$ 是任意给定的正数, $c$ 是任意给定的实数.

(1) ${\log }_{a}1 = 0$ ;

(2) ${\log }_{a}a = 1$ ；

(3) ${\log }_{a}\left( {MN}\right)  = {\log }_{a}M + {\log }_{a}N$ ；

(4) ${\log }_{a}\frac{M}{N} = {\log }_{a}M - {\log }_{a}N$ ；

(5) ${\log }_{a}{N}^{c} = c{\log }_{a}N$ ;

(6) (换底公式) 如果 $b > 0$ ,且 $b \neq  1$ ,那么 ${\log }_{b}N = \frac{{\log }_{a}N}{{\log }_{a}b}$ .

6. 重要结论: ${a}^{{\log }_{a}N} = N\left( {N > 0}\right) ;{\log }_{a}{a}^{N} = N\left( {N \in  \mathbf{R}}\right) ;{\log }_{{a}^{n}}{N}^{m} = \frac{m}{n}{\log }_{a}N$ .

7. 常用对数和自然对数: 以 10 为底的对数叫常用对数,记作 $\lg N = {\log }_{10}N$ ; 以 ${e}$ 为底的对数叫自然对数,记作 $\ln N = {\log }_{{e}}N$ ,其中 ${e} = {2.71828}\cdots$ 是欧拉常数, ${e}$ 是一个无理数.

\section*{二、几个重要的初等函数的图像和性质}

1. 幂函数、指数函数、对数函数(属于基本初等函数，概念和性质任何时候可以直接应用)的图像与性质:

\begin{center}
\begin{tabular}{|c|c|c|c|c|}
\hline
幂函数 $y = {x}^{a}$ & \multicolumn{2}{|c|}{指数函数 $y = {a}^{x}\left( {a > 0,a \neq  1}\right)$} & \multicolumn{2}{|c|}{对数函数 $y = {\log }_{a}x\left( {a > 0,a \neq  1}\right)$} \\
\cline{1-5}
第一象限内图像 & $a > 1$ & $0 < a < 1$ & $a > 1$ & $0 < a < 1$ \\
\cline{1-5}
\includegraphics[width=0.2\textwidth]{images/bo_d7tk7oc91nqc73875fu0_61_185_1598_201_186_0.jpg} & \includegraphics[width=0.2\textwidth]{images/bo_d7tk7oc91nqc73875fu0_61_446_1594_187_184_0.jpg} & \includegraphics[width=0.2\textwidth]{images/bo_d7tk7oc91nqc73875fu0_61_672_1593_191_185_0.jpg} & \includegraphics[width=0.2\textwidth]{images/bo_d7tk7oc91nqc73875fu0_61_891_1596_220_180_0.jpg} & \includegraphics[width=0.2\textwidth]{images/bo_d7tk7oc91nqc73875fu0_61_1141_1594_223_180_0.jpg} \\
\cline{1-5}
定点 $\left( {1,1}\right)$ & 定点 $\left( {0,1}\right)$ & 定点 $\left( {0,1}\right)$ & 定点 $\left( {1,0}\right)$ & 定点 $\left( {1,0}\right)$ \\
\cline{1-5}
\multirow{2}{*}{\makecell{ $a > 0$ 时 $\lbrack 0, + \infty )$ \textbackslash \\ $a < 0$ 时 $\left( {0, + \infty }\right)  \downarrow$ }} & $\mathbf{R} \uparrow$ & R ↓ & $\left( {0, + \infty }\right)  \uparrow$ & $\left( {0, + \infty }\right)  \downarrow$ \\
\cline{2-5}
 & $x > 0 \Rightarrow  y > 1  x < 0 \Rightarrow  0 < y < 1$ & \makecell{ $x > 0 \Rightarrow  0 < y < 1$ \\ $x < 0 \Rightarrow  y > 1$ } & \makecell{ $x > 1 \Rightarrow  y > 0$ \\ $0 < x < 1 \Rightarrow  y < 0$ } & \makecell{ $x > 1 \Rightarrow  y < 0$ \\ $0 < x < 1 \Rightarrow  y > 0$ } \\
\cline{1-5}
\hline
\end{tabular}
\end{center}

\HRule

*注:本章节所属知识模块是“函数与分析”.

\HRule

2. “耐克 (双勾) 函数”:形如的 $f\left( x\right)  = {ax} + \frac{b}{x}\left( {{ab} \neq  0}\right)$ 函数，不属于基本初等函数，所以单调性用于解答题时应予以证明.

“耐克(双勾)函数”在各种情况下的图像(双曲线):

\begin{center}
\includegraphics[width=0.9\textwidth]{images/bo_d7tk7oc91nqc73875fu0_62_300_343_1081_302_0.jpg}
\end{center}
 

渐近线 $y$ 轴和 $y = {ax}\left( {a \neq  0}\right)$

性质:

\begin{center}
\begin{tabular}{|c|c|c|c|}
\hline
函数 & 参数范围 & 单调递增区间 & 单调递减区间 \\
\cline{1-4}
\multirow{4}{*}{$y = {ax} + \frac{b}{x}$} & $a > 0,b > 0$ & $\left( {-\infty , - \frac{\sqrt{ab}}{a}}\right\rbrack$ 和 $\left\lbrack  {\frac{\sqrt{ab}}{a}, + \infty }\right)$ & $\left\lbrack  {-\frac{\sqrt{ab}}{a},0)\text{ 和 }(0,\frac{\sqrt{ab}}{a}}\right\rbrack$ \\
\cline{2-4}
 & $a \geq  0,b < 0$ & $\left( {-\infty ,0}\right)$ 和 $\left( {0, + \infty }\right)$ & 无 \\
\cline{2-4}
 & $a < 0,b < 0$ & $\left\lbrack  {\frac{\sqrt{ab}}{a},0}\right)$ 和 $\left( {0, - \frac{\sqrt{ab}}{a}}\right\rbrack$ & $\left( {-\infty ,\frac{\sqrt{ab}}{a}}\right\rbrack$ 和 $\left\lbrack  {-\frac{\sqrt{ab}}{a}, + \infty }\right)$ \\
\cline{2-4}
 & $a \leq  0,b > 0$ & 无 & $\left( {-\infty ,0}\right)$ 和 $\left( {0, + \infty }\right)$ \\
\cline{1-4}
\hline
\end{tabular}
\end{center}

3. 由具体函数抽象出来的函数性质的特征表达式:

(1)正比例函数: $f\left( {x + y}\right)  = f\left( x\right)  + f\left( y\right)$ 或 $f\left( {x - y}\right)  = f\left( x\right)  - f\left( y\right)$ 或 $f\left( {nx}\right)  = {nf}\left( x\right) \left( {n \in  {\mathbf{N}}^{ * }}\right)$ ；

(2)幂函数: $f\left( {x \cdot  y}\right)  = f\left( x\right)  \cdot  f\left( y\right)$ 或 $f\left( \frac{x}{y}\right)  = \frac{f\left( x\right) }{f\left( y\right) }$ 或 $f\left( {x}^{n}\right)  = {f}^{n}\left( x\right) \left( {n \in  {\mathbf{N}}^{ * }}\right)$ ；

(3)指数函数: $f\left( {x + y}\right)  = f\left( x\right)  \cdot  f\left( y\right)$ 或 $f\left( {x - y}\right)  = \frac{f\left( x\right) }{f\left( y\right) }$ 或 $f\left( {nx}\right)  = {f}^{n}\left( x\right) \left( {n \in  {\mathbf{N}}^{ * }}\right)$ ；

(4)对数函数: $f\left( {x \cdot  y}\right)  = f\left( x\right)  + f\left( y\right)$ 或 $f\left( \frac{x}{y}\right)  = f\left( x\right)  - f\left( y\right)$ 或 $f\left( {x}^{n}\right)  = {nf}\left( x\right) \left( {n \in  {\mathbf{N}}^{ * }}\right)$ .

\section*{§ 4.1 幂函数、二次函数与“对勾”函数}

\section*{典型例题精讲}

【例 1】已知幂函数 $f\left( x\right)  = {x}^{\frac{1}{{m}^{2} + m}}\left( {m \in  {\mathbf{N}}^{ * }}\right)$ .

(1)试确定该函数的定义域，并指明该函数在其定义域上的单调性；

(2)若函数 $f\left( x\right)$ 的图像经过点 $\left( {2,\sqrt{2}}\right)$ ，试确定 $m$ 的值，并求满足条件 $f\left( {2 - a}\right)  > f\left( {a - 1}\right)$ 的实数 $a$ 的取值范围.

【例 2】已知幂函数 $f\left( x\right)  = {\left( m - 1\right) }^{2}{x}^{m{}^{2} - {4m} + 3}$ 在 $\left( {0, + \infty }\right)$ 上是严格增函数.

(1)求实数 $m$ 的值；

(2)若函数 $h\left( x\right)  =  - \sqrt[3]{{\left\lbrack  f\left( x\right) \right\rbrack  }^{2}} + {2bx} + 1 - b$ 在 $\left\lbrack  {0,2}\right\rbrack$ 上的最大值为 3，求实数 $b$ 的值.

【例 3】因函数 $y = x + \frac{t}{x}\left( {t > 0}\right)$ 的图像形状像对勾,我们称形如 “ $y = x + \frac{t}{x}\left( {t > 0}\right)$ ” 的函数为 “对勾函数”.

(1)已知 $f\left( x\right)  = {2x} + \frac{4}{{2x} - 1} - 5,x \in  \left\lbrack  {1,3}\right\rbrack$ ，利用对勾函数性质，求函数 $f\left( x\right)$ 的单调区间和值域；

(2)对于(1)中的函数 $f\left( x\right)$ 和函数 $g\left( x\right)  = {x}^{2} - {mx} + 4$ ，若对任意 ${x}_{1} \in  \left\lbrack  {1,3}\right\rbrack$ ，总存在 ${x}_{2} \in  \left\lbrack  {1,3}\right\rbrack$ ，使得 $g\left( {x}_{2}\right)  < f\left( {x}_{1}\right)$ 成立，求实数 $m$ 的取值范围.

【例 4】已知函数 $f\left( x\right)  = m\left( {x + \frac{1}{x}}\right)  + n$ 的图像与 $h\left( x\right)  = \frac{1}{4}\left( {x + \frac{1}{x}}\right)$ 的图像关于 $\left( {0, - 1}\right)$ 对称.

(1)求实数 $m$ 、 $n$ 的值；

(2)若 $g\left( x\right)  = f\left( x\right)  + \frac{a}{4x}$ ，求 $g\left( x\right)$ 在区间 $\left\lbrack  {1,2}\right\rbrack$ 上的最小值.

【例 5】设函数 $y = f\left( x\right)$ 的定义域为 $D$ ,若存在 $\left\lbrack  {a,b}\right\rbrack   \subseteq  D$ ,使得 $y = f\left( x\right)$ 在 $\left\lbrack  {a,b}\right\rbrack$ 上的值域也为 $\left\lbrack  {a,b}\right\rbrack$ ,则称函数 $y = f\left( x\right)$ 具有性质 $P$ ,已知幂函数 $f\left( x\right)  = \left( {{p}^{2} + p - 1}\right) {x}^{p - \frac{1}{2}}$ 在 $\left( {0, + \infty }\right)$ 上是严格增函数.

(1) $g\left( x\right)  = f\left( x\right)  - \frac{2}{9}$ 是否具有性质 $P$ ? 若是，请指出区间 $\left\lbrack  {a,b}\right\rbrack$ ；若不是，请说明理由；

(2)若函数 $h\left( x\right)  = n - f\left( {x + 1}\right)$ ，且 $h\left( x\right)$ 具有性质 $P$ ，求实数 $n$ 的取值范围.

\section*{四基达标精练}

\section*{一、填空题}

1. 已知幂函数 $y = f\left( x\right)$ 的图像过点 $P\left( {2,8}\right)$ ，则函数 $y = f\left( x\right)  - x$ 的零点为\_\_\_\_\_.

2. 若幂函数 $y = \left( {{m}^{2} - m - 1}\right) {x}^{{m}^{2} - {2m} - 1}$ 在 $\left( {0, + \infty }\right)$ 上是增函数，则实数 $m =$ \_\_\_\_\_.

3. 若 ${\left( \frac{1}{2}\right) }^{-\frac{1}{2}} > {a}^{-\frac{1}{2}}$ ，则实数 $a$ 的取值范围是\_\_\_\_\_.

4. 一元二次方程 ${x}^{2} + \left( {{a}^{2} - 1}\right) x + \left( {a - 2}\right)  = 0$ 的一根比 1 大，另一根比 -1 小，则实数 $a$ 的取值范围是\_\_\_\_\_.

5. 设定点 $A\left( {0,1}\right)$ ，若动点 $P$ 在函数 $y = \frac{x + 2}{x}\left( {x > 0}\right)$ 图像上，则 $\left| {PA}\right|$ 的最小值为\_\_\_\_\_.

6. 已知函数 $y = {x}^{2} + x + \frac{1}{2}\left( {n \leq  x \leq  n + 1}\right)$ ，其中 $n$ 为正整数的值域中共含有 ${a}_{n}$ 个整数，则 ${a}_{n} =$ \_\_\_\_\_.

7. 已知四个函数: \ding{172}  $y =  - x$ ;  \ding{173}  $y =  - \frac{1}{x}$ ;  \ding{174}  $y = {x}^{3}$ ； \ding{175}  $y = {x}^{\frac{1}{2}}$ . 从中任选 2 个，则事件“所选 2 个函数的图像有且仅有一个公共点”的概率为\_\_\_\_\_.

8. 已知函数 $f\left( x\right)  = \left| {x + \frac{9}{x} - a}\right|  + a$ 在区间 $\left\lbrack  {1,9}\right\rbrack$ 上的最大值是 10,则实数 $a$ 的取值范围是\_\_\_\_\_.

9. 设函数 $F\left( x\right)$ 和 $f\left( x\right)$ 都在区间 $D$ 上有定义,若对 $D$ 的任意子区间 $\left\lbrack  {u,v}\right\rbrack$ ,总有 $\left\lbrack  {u,v}\right\rbrack$ 上的实数 $p\text{ 、 }q$ , 使得不等式 $f\left( p\right)  \leq  \frac{F\left( u\right)  - F\left( v\right) }{u - v} \leq  f\left( q\right)$ 成立,则称 $F\left( x\right)$ 是 $f\left( x\right)$ 在区间 $D$ 上的甲函数, $f\left( x\right)$ 是 $F\left( x\right)$ 在区间 $D$ 上的乙函数. 已知 $F\left( x\right)  = {x}^{2} - {3x}\left( {x \in  \mathbf{R}}\right)$ ，那么 $F\left( x\right)$ 的乙函数 $f\left( x\right)  =$ \_\_\_\_\_.

\section*{二、选择题}

10. 设 $f\left( x\right)  = {x}^{2} - x + a$ ,若 $f\left( {-m}\right)  < 0$ ,则 $f\left( {m + 1}\right)$ 的值是 ( )

A. 正数. B. 负数. C. 非负数. D. 正负与 $m$ 和 $a$ 有关.

11. 已知幂函数的图像过 $\left( {\frac{1}{2},\frac{\sqrt{2}}{4}}\right) ,P\left( {{x}_{1},{y}_{1}}\right) \text{ 、 }Q\left( {{x}_{2},{y}_{2}}\right) \left( {{x}_{1} < {x}_{2}}\right)$ 是函数图像上的任意不同两点,则下列结论中正确的是 ( )

A. ${x}_{1}f\left( {x}_{1}\right)  > {x}_{2}f\left( {x}_{2}\right)$ . B. ${x}_{1}f\left( {x}_{2}\right)  < {x}_{2}f\left( {x}_{1}\right)$ . C. $\frac{f\left( {x}_{1}\right) }{{x}_{2}} > \frac{f\left( {x}_{2}\right) }{{x}_{1}}$ . D. $\frac{f\left( {x}_{1}\right) }{{x}_{1}} < \frac{f\left( {x}_{2}\right) }{{x}_{2}}$ .

12. 函数 $f\left( x\right)  = a{x}^{2} + {bx} + c\left( {a \neq  0}\right)$ 的图像关于直线 $x =  - \frac{b}{2a}$ 对称. 据此可推测,对任意的非零实数 $a$ 、 $b\text{ 、 }c\text{ 、 }m\text{ 、 }n\text{ 、 }p$ ,关于 $x$ 的方程 $m{\left\lbrack  f\left( x\right) \right\rbrack  }^{2} + {nf}\left( x\right)  + p = 0$ 的解集都不可能是 ( )

A. $\{ 1,2\}$ . B. $\{ 1,3,5\}$ . C. $\{ 1,4,7,{10}\}$ . D. $\{ 1,4,{16},{64}\}$ .

\section*{三、解答题}

13. 已知幂函数 $f\left( x\right)  = {x}^{9 - {3m}}\left( {m \in  {\mathbf{N}}^{ * }}\right)$ 的图像关于原点对称,且在 $\mathbf{R}$ 上严格递增.

(1)求 $f\left( x\right)$ 表达式；

(2)求满足 $f\left( {a + 1}\right)  + f\left( {{3a} - 4}\right)  < 0$ 实数 a 的取值范围.

14. 已知幂函数 $f\left( x\right)  = \left( {{m}^{2} - {5m} + 7}\right) {x}^{m - 1}$ 为偶函数.

(1)求 $f\left( x\right)$ 的解析式;

(2)若 $g\left( x\right)  = f\left( x\right)  - {ax} - 3$ 在 $\left\lbrack  {1,3}\right\rbrack$ 上不是单调函数，求实数 $a$ 的取值范围.

15. 已知二次函数 $f\left( x\right)  = a{x}^{2} + {bx}$ 对任意 $x \in  \mathbf{R}$ 均有 $f\left( {x - 1}\right)  = f\left( {-1 - x}\right)$ 成立,且函数的图像过点 $A\left( {1,\frac{3}{2}}\right)$ .

(1)求函数 $y = f\left( x\right)$ 的解析式;

(2)若不等式 $f\left( {x - t}\right)  \leq  x$ 的解集为 $\left\lbrack  {4,m}\right\rbrack$ ，求实数 $t$ 、 $m$ 的值.

16. 某通信公司需要在三角形地带 ${OAC}$ 区域内建造甲、乙两种通信信号加强中转站,甲中转站建在区域 ${BOC}$ 内,乙中转站建在区域 ${AOB}$ 内. 分界线 ${OB}$ 固定,且 ${OB} = \left( {1 + \sqrt{3}}\right)$ 百米,边界线 ${AC}$ 始终过点 $B$ ,边界线 ${OA}\text{ 、 }{OC}$ 满足 $\angle {AOC} = {75}^{ \circ  },\angle {AOB} = {30}^{ \circ  }$ ,设 ${OA} = x\left( {3 \leq  x \leq  6}\right)$ 百米, ${OC} = y$ 百米.

(1)试将 $y$ 表示成 $x$ 的函数，并求出函数 $y$ 的解析式；

(2)当 $x$ 取何值时， $\bigtriangleup {OAC}$ 的面积 ${S}_{\bigtriangleup {OAC}}$ 最小？求出其最小值.

\begin{center}
\includegraphics[width=0.2\textwidth]{images/bo_d7tk7oc91nqc73875fu0_67_1167_1110_268_278_0.jpg}
\end{center}
 

(第 16 题图)

\section*{$\text{ § }{4.2}$ 指数运算与指数函数}

\section*{典型例题精讲}

【例 1】求下列不等式的解集:

(1) ${3}^{x - 1} > {9}^{x}$ ；

(2) ${a}^{-{5x}} > {a}^{x + 7}\left( {a > 0,a \neq  1}\right)$ .

【例 2】已知函数 $g\left( x\right)  = {\left( a + 1\right) }^{x - 2} + 1\left( {a > 0}\right)$ 的图像恒过定点 $A$ ,且点 $A$ 又在函数 $f\left( x\right)  = {\log }_{\sqrt{3}}\left( {x + a}\right)$ 的图像上.

(1)求实数 $a$ 的值；

(2)当 $\left| {g\left( {x + 2}\right)  - 2}\right|  = {2b}$ 有两个不等实根时,求实数 $b$ 的取值范围.

【例 3】已知奇函数 $f\left( x\right)$ 与偶函数 $g\left( x\right)$ 的定义域都是 $\mathbf{R}$ ,且满足 $f\left( x\right)  + g\left( x\right)  = {2}^{x}$ .

(1)求 $f\left( x\right)$ 的解析式；

(2)设函数 $h\left( x\right)  = f\left( x\right)  + x$ .

 \ding{172} 判断 $h\left( x\right)$ 的单调性，并用定义法证明；

 \ding{173}  若 $f\left( {{\log }_{2}m}\right)  + f\left( {2{\log }_{2}m - 1}\right)  \leq  1 - 3{\log }_{2}m$ ，求实数 $m$ 的取值范围.

【例 4】已知函数 $f\left( x\right)  = {a}^{x} + b\left( {a > 0,a \neq  1}\right)$ .

(1)若 $f\left( x\right)$ 的图像如图(1)所示，求 $a$ 、 $b$ 的值；

(2)若 $f\left( x\right)$ 的图像如图(2)所示，求 $a$ 、 $b$ 的取值范围；

(3)在(1)中，若 $\left| {f\left( x\right) }\right|  = m$ 有且仅有一个实数解，求出 $m$ 的范围.

\begin{center}
\includegraphics[width=0.3\textwidth]{images/bo_d7tk7oc91nqc73875fu0_69_1003_1211_447_239_0.jpg}
\end{center}
 

(例 4 图)

【例 5】已知函数 $f\left( x\right)  = {4}^{x} - a \cdot  {2}^{x + 1} - 2$ .

(1)若 $f\left( x\right)$ 的最小值为 -3，求实数 $a$ 的值；

(2)若 $f\left( x\right)  < 0$ 对任意的 $x \in  \lbrack 0,1)$ ，恒成立，求实数 $a$ 的取值范围.

\section*{四基达标精练}

\section*{一、填空题}

1．化简 ${\left( {2}^{n + 1}\right) }^{2} \cdot  {2}^{-{2n} - 1} \div  {4}^{n}$ 的结果是\_\_\_\_\_.

2. 化简 $\frac{\sqrt{-{x}^{3}}}{x}$ 的结果是\_\_\_\_\_.

3. 用分数指数幂形式表示: $\sqrt{a\sqrt{a\sqrt{a\sqrt{a}}}} =$ \_\_\_\_\_.

4. 若指数函数 $y = {\left( 2a - 1\right) }^{x}$ 在 $\mathbf{R}$ 上严格递减，则实数 $a$ 的取值范围是\_\_\_\_\_.

5. 已知函数 $f\left( x\right)  = {a}^{x - 3} + 2$ 的图像恒过定点 $A$ ，则 $A$ 的坐标为\_\_\_\_\_.

6. 若 ${3}^{2x} + 9 = {10} \cdot  {3}^{x}$ ，则 ${x}^{2} + 1 =$ \_\_\_\_\_.

7. 函数 $y = {\left( \frac{1}{2}\right) }^{{x}^{2} + {2x} - 1}$ 的值域是\_\_\_\_\_.

8. 方程 $\left( {{4}^{x} + {4}^{-x}}\right)  - 2\left( {{2}^{x} + {2}^{-x}}\right)  + 2 = 0$ 的解集是\_\_\_\_\_.

9. 已知 $f\left( x\right)  + g\left( x\right)  = {a}^{x} - {a}^{-x} + 2\left( {a > 0,a \neq  1}\right)$ ,其中 $f\left( x\right)$ 是定义在 $\mathbf{R}$ 上的奇函数, $g\left( x\right)$ 是定义在 $\mathbf{R}$ 上的偶函数. 若 $g\left( a\right)  = a$ ,则 $f\left( a\right)  =$ \_\_\_\_\_.

\section*{二、选择题}

10. 已知函数 $g\left( x\right)  = {3}^{x} + t$ 的图像不经过第二象限,则 $t$ 的取值范围是 ( )

A. $t \leq   - 1$ . B. $t <  - 1$ . C. $t \leq   - 3$ . D. $t \geq   - 3$ .

11. 若 $0 < b < a < 1$ ,则在 ${a}^{b}$ 、 ${b}^{a}$ 、 ${a}^{a}$ 、 ${b}^{b}$ 中的最大值是 ( )

A. ${b}^{a}$ . B. ${a}^{a}$ . C. ${a}^{b}$ . D. ${b}^{b}$ .

12. 已知 $n \in  {\mathbf{N}}^{ * }$ ,则函数 $y = {x}^{n}\left( {x \in  \mathbf{R}}\right)$ 与 $y = {n}^{x}\left( {x \in  \mathbf{R}}\right)$ 图像的交点不可能 ( )

A. 只有 $\left( {n,{n}^{n}}\right)$ . B. 在直线 $y = {nx}$ 上. C. 多于三个. D. 在第二象限.

\section*{三、解答题}

13. 已知函数 $f\left( x\right)  = {a}^{x - 1}\left( {x \geq  0}\right)$ 的图像经过点 $\left( {2,\frac{1}{2}}\right)$ ，其中 $a > 0$ 且 $a \neq  1$ .

(1)求 $a$ 的值；

(2)求函数 $y = f\left( x\right)$ 的值域.

14. 解下列方程:

(1) ${3}^{{2x} + 5} = 5 \cdot  {3}^{x + 2} + 2$ ； (2) ${\left( \sqrt{3 + 2\sqrt{2}}\right) }^{x} + {\left( \sqrt{3 - 2\sqrt{2}}\right) }^{x} = 6$ .

15. 已知 $f\left( x\right)  = {{e}}^{x} - {{e}}^{-x},g\left( x\right)  = {{e}}^{x} + {{e}}^{-x}$ ,其中 ${e} = {2.71828}\cdots$ .

(1)求 ${\left\lbrack  f\left( x\right) \right\rbrack  }^{2} - {\left\lbrack  g\left( x\right) \right\rbrack  }^{2}$ 的值；

(2)设 $f\left( x\right)$ 、 $f\left( y\right)  = 4,g\left( x\right)  \cdot  g\left( y\right)  = 8$ ，求 $\frac{g\left( {x + y}\right) }{g\left( {x - y}\right) }$ 的值.

16. 已知 $f\left( x\right)  = \frac{{4}^{x} + 1}{{4}^{x} + a}$ 是奇函数,其中 $a$ 是常数.

(1)求 $f\left( x\right)$ 的定义域和值域；

(2)讨论 $f\left( x\right)$ 的单调性并证明.

\section*{§ 4.3 对数运算与对数函数}

\section*{典型例题精讲}

【例 1】已知函数 $f\left( x\right)  = {\log }_{a}x + b$ (其中 $a\text{ 、 }b$ 均为常数, $a > 0$ 且 $a \neq  1$ ) 的图像经过点 $\left( {2,5}\right)$ 和点 $\left( {8,7}\right)$ .

(1)求 $a$ 、 $b$ 的值；

(2)设函数 $g\left( x\right)  = {b}^{x} - {a}^{x + 2}$ ，若对任意的 ${x}_{1} \in  \left\lbrack  {1,4}\right\rbrack$ ，总存在 ${x}_{2} \in  \left\lbrack  {0,{\log }_{2}5}\right\rbrack$ ，使得 $f\left( {x}_{1}\right)  = g\left( {x}_{2}\right)  + m$ 成立,求实数 $m$ 的取值范围.

【例 2】如图,已知函数 $y = {kx}\left( {k > 0}\right)$ 的图像与函数 $y = {\log }_{2}x$ 的图像交于 $A\text{ 、 }B$ 两点.过 $A\text{ 、 }B$ 分别作 $x$ 轴的垂线,垂足分别为点 ${A}^{\prime }$ 、点 ${B}^{\prime }$ ，并且 $A{A}^{\prime }$ 、 $B{B}^{\prime }$ 分别交函数 $y = {\log }_{4}x$ 的图像于 $M$ 、 $N$ 两点.

(1)探究线段 ${AM}$ 与 $M{A}^{\prime }$ 长度的大小关系；

(2)若 ${AN}$ 平行于 $x$ 轴,求四边形 ${AMNB}$ 的面积.

\begin{center}
\includegraphics[width=0.2\textwidth]{images/bo_d7tk7oc91nqc73875fu0_73_1149_1317_288_216_0.jpg}
\end{center}
 

(例 2 图)

【例 3】( 1 )若方程 ${\log }_{2}\left( {x + 3}\right)  - {\log }_{4}{x}^{2} = a$ 的根在( 3,4 )内，求实数 $a$ 的取值范围；

(2)已知 $a > 0,a \neq  1$ ，试求使方程 ${\log }_{a}\left( {x - {ak}}\right)  = {\log }_{{a}^{2}}\left( {{x}^{2} - {a}^{2}}\right)$ 有解的实数 $k$ 的取值范围.

【例 4】已知函数 $f\left( x\right)  = {\log }_{\frac{1}{2}}\frac{1 - {ax}}{x - 1}$ 的图像关于原点对称,其中 $a$ 为常数.

(1)求 $a$ 的值；

(2)当 $x \in  \left( {1, + \infty }\right)$ 时， $f\left( x\right)  + {\log }_{\frac{1}{2}}\left( {x - 1}\right)  < m$ 恒成立，求实数 $m$ 的取值范围；

(3)若关于 $x$ 的方程 $f\left( x\right)  = {\log }_{\frac{1}{2}}\left( {x + k}\right)$ 在 $\left\lbrack  {2,3}\right\rbrack$ 上有解,求 $k$ 的取值范围.

【例 5】已知函数 $f\left( x\right)  = {\log }_{a}x\left( {a > 0\text{ 且 }a \neq  1}\right)$ .

(1)若 $f\left( {a + 4}\right)  \leq  f\left( {3a}\right)$ ，求实数 $a$ 的取值范围；

(2)设 $a = 2$ ，函数 $g\left( x\right)  =  - {f}^{2}\left( x\right)  + \left( {3 - {2m}}\right) f\left( x\right)  + m + 2\left( {0 < m \leq  1}\right)$ . 若 $x \in  \left\lbrack  {1,{2}^{m}}\right\rbrack$ ，证明: $g\left( x\right)  \leq  \frac{10}{3}$ .

\section*{四基达标精练}

\section*{一、填空题}

1. 若 $a > 1,b > 1,p = \frac{{\log }_{b}\left( {{\log }_{b}a}\right) }{{\log }_{b}a}$ ，则 ${a}^{p} =$ \_\_\_\_\_(用 $a$ 、 $b$ 表示).

2. 已知 $2\lg \left( {x - {2y}}\right)  = \lg x + \lg y$ ，则 $\frac{x}{y}$ 的值为\_\_\_\_\_.

3. 设 ${2}^{m} = {5}^{n} = {10}$ ，则 $\frac{1}{m} + \frac{1}{n} =$ \_\_\_\_\_.

4. 已知 ${\log }_{6}{27} = a$ ，试用 $a$ 表示 ${\log }_{18}{16} =$ \_\_\_\_\_.

5. 已知对数函数 $y = {\log }_{a}x$ 的图像过点 $P\left( {8,3}\right)$ ，则 $f\left( \frac{1}{32}\right)  =$ \_\_\_\_\_.

6．如果方程 ${\lg }^{2}x + \left( {\lg 7 + \lg 5}\right) \lg x + \lg 7 \cdot  \lg 5 = 0$ 的两根为 $\alpha$ 、 $\beta$ ，则 $\alpha  \cdot  \beta  =$ \_\_\_\_\_.

7. 方程 ${x}^{\lg x} = {10}$ 的所有实数根之积是\_\_\_\_\_.

8. 已知 ${\log }_{2}x = {\log }_{3}y = {\log }_{5}z =  - 2$ ，则 ${x}^{\frac{1}{2}}$ 、 ${y}^{\frac{1}{3}}$ 、 ${z}^{\frac{1}{5}}$ 由小到大的排列顺序是\_\_\_\_\_.

9. 已知函数 $f\left( x\right)$ 满足:对于任意 ${x}_{1}$ 、 ${x}_{2} \in  \mathbf{R}$ ，都有 $f\left( {{x}_{1} + {x}_{2}}\right)  = f\left( {x}_{1}\right)  + f\left( {x}_{2}\right)  - {2016}$ ，且当 $x > 0$ 时， $f\left( x\right)  > {2016}$ ，若 $f\left( x\right)$ 在区间 $\left\lbrack  {-{2023},{2023}}\right\rbrack$ 上的最大值、最小值分别为 $M$ 、 $N$ ，则 $M + N =$ \_\_\_\_\_.

\section*{二、选择题}

10. 设 $a$ 、 $b$ 、 $c$ 都是正数，且 ${3}^{a} = {4}^{b} = {6}^{c}$ ，则 ( )

A. $\frac{1}{c} = \frac{1}{a} + \frac{1}{b}$ . B. $\frac{2}{c} = \frac{2}{a} + \frac{1}{b}$ . C. $\frac{1}{c} = \frac{2}{a} + \frac{2}{b}$ . D. $\frac{2}{c} = \frac{1}{a} + \frac{2}{b}$ .

11. 已知函数 $f\left( x\right)  = \left| {{\log }_{2}x}\right|$ ，正实数 $m$ 、 $n\left( {m < n}\right)$ 满足 $f\left( m\right)  = f\left( n\right)$ ，若 $f\left( x\right)$ 在区间 $\left\lbrack  {{m}^{2},n}\right\rbrack$ 上的最大值为 2,则 $m\text{ 、 }n$ 的值分别为 ( )

A. $\frac{\sqrt{2}}{2},\sqrt{2}$ . B. $\frac{1}{4}\text{ 、 }2$ . C. $\frac{1}{2}\text{ 、 }2$ . D. $\frac{1}{4}$ 、4.

12. 已知函数 $f\left( x\right)$ 满足 $f\left( {x + 2}\right)  = f\left( x\right)$ 对任意 $x \in  \mathbf{R}$ 恒成立，且当 $x \in  \left\lbrack  {1,3}\right\rbrack$ 时， $f\left( x\right)  = \left| {x - 2}\right|  - 1$ ，若函数 $y  = f\left( x\right)  - {\log }_{a}\left( {x + 1}\right)$ 至少有三个零点，则实数 $a$ 的取值范围为 ( )

A. $\left( {0,\frac{\sqrt{3}}{3}}\right)$ . B. $\left( {0,\frac{1}{3}}\right)$ . C. $\left( {0,\frac{1}{2}}\right)$ . D. $\left( {\frac{1}{3},1}\right)$ .

\section*{三、解答题}

13. 已知函数 $f\left( x\right)  = \lg \left( {x - 1}\right) ,g\left( x\right)  = \lg \left( {4 - x}\right)$ .

(1)求函数 $h\left( x\right)  = f\left( x\right)  - g\left( x\right)$ 的定义域；

(2)当不等式 $f\left( x\right)  > g\left( x\right)$ 成立时,求实数 $x$ 的取值范围.

14. 解关于 $x$ 的方程.

(1) ${\log }_{4}\left( {3 - x}\right)  + {\log }_{0.25}\left( {3 + x}\right)  = {\log }_{4}\left( {1 - x}\right)  + {\log }_{0.25}\left( {{2x} + 1}\right)$ ;

(2) ${\log }_{\left( x + a\right) }{2x} = 2$ .

15. 函数 $f\left( x\right)  = b + {\log }_{a}x\left( {a > 0\text{ 且 }a \neq  1}\right)$ 的图像经过点 $\left( {8,2}\right)$ 和 $\left( {1, - 1}\right)$ .

(1)求函数 $f\left( x\right)$ 的解析式；

(2)函数 $g\left( x\right)  = {f}^{2}\left( x\right)  - f\left( x\right)$ ，求函数 $g\left( x\right)$ 的最小值.

16. 已知函数 $f\left( x\right)  = {a}^{x - 1}\left( {a > 0\text{ 且 }a \neq  1}\right)$ .

(1)若函数 $y = f\left( x\right)$ 的图像经过 $P\left( {3,4}\right)$ 点，求 $a$ 的值；

(2)若 $f\left( {\lg a}\right)  = {100}$ ，求 $a$ 的值.

\section*{第 5 章 导数及其应用(选择性必修 2)}

\section*{知识梳理*}

\section*{一、平均变化率与导数的概念:}

对于函数 $y = f\left( x\right)$ ,将 $\frac{f\left( {{x}_{0} + h}\right)  - f\left( {x}_{0}\right) }{h}$ 称为函数在以 ${x}_{0}$ 和 ${x}_{0} + h$ 为端点的区间上的平均变化率, 如果当 $h$ 趋近于 0 时, $\frac{f\left( {{x}_{0} + h}\right)  - f\left( {x}_{0}\right) }{h}$ 趋近于一个稳定值,则将这一稳定值称为 $h$ 趋于 0 时, $\frac{f\left( {{x}_{0} + h}\right)  - f\left( {x}_{0}\right) }{h}$ 的极限,记作 $\mathop{\lim }\limits_{{h \rightarrow  0}}\frac{f\left( {{x}_{0} + h}\right)  - f\left( {x}_{0}\right) }{h},\mathop{\lim }\limits_{{h \rightarrow  0}}\frac{f\left( {{x}_{0} + h}\right)  - f\left( {x}_{0}\right) }{h}$ 是函数在 $x = {x}_{0}$ 处的瞬时变化率,也称为函数 $y = f\left( x\right)$ 在 $x = {x}_{0}$ 处的导数,记作 ${f}^{\prime }\left( {x}_{0}\right)$ .

即 ${f}^{\prime }\left( {x}_{0}\right)  = \mathop{\lim }\limits_{{h \rightarrow  0}}\frac{f\left( {{x}_{0} + h}\right)  - f\left( {x}_{0}\right) }{h}$ ; 也就是说,导数值 $=$ 瞬时变化率 $=$ 平均变化率的极限.

\section*{二、平均变化率和导数值的几何意义}

定点 $P\left( {{x}_{0},f\left( {x}_{0}\right) }\right)$ 和动点 $Q\left( {{x}_{0} + h,f\left( {{x}_{0} + h}\right) }\right)$ 在函数 $y = f\left( x\right)$ 上,如果 ${f}^{\prime }\left( {x}_{0}\right)$ 存在,则 $\frac{f\left( {{x}_{0} + h}\right)  - f\left( {x}_{0}\right) }{h} =$ 割线 ${PQ}$ 的斜率,当 $h$ 趋于 0 时,割线 ${PQ}$ 逐渐趋近于 $y = f\left( x\right)$ 在点 $P$ 处的切线,因此 ${f}^{\prime }\left( {x}_{0}\right)$ 是 $y = f\left( x\right)$ 在点 $P$ 处切线的斜率.

\section*{三、导函数}

1. 对于函数 $y = f\left( x\right)$ ，在导数值存在的前提下，对不同的 ${x}_{0}$ ，总有一个确定的导数值 $y = f\left( {x}_{0}\right)$ 对应，换句话说，如果用 $x$ 表示自变量，那么 ${y}^{\prime } = {f}^{\prime }\left( x\right)$ 也是关于 $x$ 的一个函数，称为 $y = f\left( x\right)$ 的导函数.

注意:导数值和导函数的区别可类比函数值和函数的区别.

2. 驻点:将导数为零的点称为函数的驻点.

\section*{四、导数的运算}

1.基本初等函数的导数公式和运算法则

\begin{center}
\begin{tabular}{|c|c|}
\hline
基本初等函数 & 导函数 \\
\cline{1-2}
$f\left( x\right)  = C\left( {C\text{ 为常数 }}\right)$ & ${f}^{\prime }\left( x\right)  = 0$ \\
\cline{1-2}
$f\left( x\right)  = {x}^{a}\left( {a \in  Q,a \neq  0}\right)$ & ${f}^{\prime }\left( x\right)  = a{x}^{a - 1}$ \\
\cline{1-2}
$f\left( x\right)  = \sin x$ & ${f}^{\prime }\left( x\right)  = \cos x$ \\
\cline{1-2}
$f\left( x\right)  = \cos x$ & ${f}^{\prime }\left( x\right)  =  - \sin x$ \\
\cline{1-2}
$f\left( x\right)  = {e}^{x}$ & ${f}^{\prime }\left( x\right)  = {e}^{x}$ \\
\cline{1-2}
$f\left( x\right)  = {a}^{x}\left( {a > 0,a \neq  1}\right)$ & $f\left( x\right)  = {a}^{x}\ln a$ \\
\cline{1-2}
$f\left( x\right)  = \ln x$ & $f\left( x\right)  = \frac{1}{x}$ \\
\cline{1-2}
$f\left( x\right)  = {\log }_{a}x\left( {a > 0,a \neq  1}\right)$ & $f\left( x\right)  = \frac{1}{x\ln a}$ \\
\cline{1-2}
\hline
\end{tabular}
\end{center}

\begin{center}
\begin{tabular}{|c|}
\hline
运算法则 \\
\cline{1-1}
设 ${f}^{\prime }\left( x\right)$ 和 ${g}^{\prime }\left( x\right)$ 都存在,则 \\
\cline{1-1}
\makecell{  \ding{172}  ${\left\lbrack  f\left( x\right)  + g\left( x\right) \right\rbrack  }^{\prime } = {f}^{\prime }\left( x\right)  + {g}^{\prime }\left( x\right)$ \\  \ding{173}  ${\left\lbrack  f\left( x\right)  - g\left( x\right) \right\rbrack  }^{\prime } = {f}^{\prime }\left( x\right)  - {g}^{\prime }\left( x\right)$ \\  \ding{174}  ${\left\lbrack  f\left( x\right)  \cdot  g\left( x\right) \right\rbrack  }^{\prime } = {f}^{\prime }\left( x\right)  \cdot  g\left( x\right)  + f\left( x\right)  \cdot  {g}^{\prime }\left( x\right)$ \\  \ding{175}  ${\left\lbrack  \frac{f\left( x\right) }{g\left( x\right) }\right\rbrack  }^{\prime } = \frac{{f}^{\prime }\left( x\right)  \cdot  g\left( x\right)  - f\left( x\right)  \cdot  {g}^{\prime }\left( x\right) }{{g}^{2}\left( x\right) }$ \\ 其中 $g\left( x\right)  \neq  0$ } \\
\cline{1-1}
\hline
\end{tabular}
\end{center}

\HRule

* 注:本章节所属知识模块是“函数与分析”.

\HRule

2. 简单复合函数的求导法则:

简单复合函数 $y = f\left( {{ax} + b}\right)$ 的导数和函数 $y = f\left( u\right)$ 的关系为 ${\left\lbrack  f\left( ax + b\right) \right\rbrack  }^{\prime } = {f}^{\prime }\left( {{ax} + b}\right)  \cdot  a$ ,其中 ${\left\lbrack  f\left( ax + b\right) \right\rbrack  }^{\prime }$ 是指复合函数 $f\left( {{ax} + b}\right)$ 关于 $x$ 的导数, ${f}^{\prime }\left( {{ax} + b}\right)$ 是指将 ${f}^{\prime }\left( u\right)$ 中的 $u$ 用 ${ax} + b$ 代替.

注: \ding{172}  求函数的导数要准确地把函数拆分成基本初等函数的和、差、积、商，再利用运算法则求导.

 \ding{173} 抽象函数求导，恰当赋值是关键，然后活用方程思想求解.

 \ding{174} 复合函数求导，应由外到内逐层求导，必要时要进行换元.

\section*{五、切线方程}

函数 $y = f\left( x\right)$ 在 $x = {x}_{0}$ 处的导数的几何意义就是曲线 $y = f\left( x\right)$ 在点 $P\left( {{x}_{0},f\left( {x}_{0}\right) }\right)$ 处的切线的斜率, 相应的切线方程为 $y - f\left( {x}_{0}\right)  = {f}^{\prime }\left( {x}_{0}\right) \left( {x - {x}_{0}}\right)$ .

1．求曲线的切线方程要分清“在点处”与“过点处”的切线方程的不同.

求曲线在点 $P\left( {{x}_{0},{y}_{0}}\right)$ 处的切线,则表明点 $P$ 是切点,可直接求出函数在 $P$ 处的导数,然后利用点斜式写出切线方程,若在该点 $P$ 处的导数不存在,则切线垂直于 $x$ 轴,切线方程为 $x = {x}_{0}$ .

求过点 $P\left( {{x}_{0},{y}_{0}}\right)$ 的切线,点 $P$ 不一定是切点,要设出切点坐标,根据斜率相等建立方程 (组) 求解,求出切点坐标是解题的关键.

2. 处理与切线有关的参数问题,通常利用曲线、切线、切点的三个关系列出参数的方程 (组) 并解出参数:

(1)切点处的导数是切线的斜率；

(2)切点在切线上，故满足切线方程；

(3)切点在曲线上，故满足曲线方程。

3. 公切线问题: 两条曲线的公切线, 一般是把两条曲线分开考虑: 分别求两条曲线的切线, 化为斜截式方程，利用方程相等，列式求参数. 若其中一条曲线涉及二次曲线，可先求另一条曲线的切线，与二次曲线联立、 消元, 用判别式法.

4. 利用导数的几何意义求参数问题时, 注意利用数形结合, 化归与转化的思想方法.

六、函数的单调性与导数.

1. 函数的单调性与导数的关系.

\begin{center}
\begin{tabular}{|c|c|c|}
\hline
条件 & 恒有 & 结论 \\
\cline{1-3}
\multirow{3}{*}{函数 $y = f\left( x\right)$ 在区间 $\left( {a,b}\right)$ 上可导} & ${f}^{\prime }\left( x\right)  > 0$ & $f\left( x\right)$ 在 $\left( {a,b}\right)$ 上单调增 \\
\cline{2-3}
 & ${f}^{\prime }\left( x\right)  < 0$ & $f\left( x\right)$ 在 $\left( {a,b}\right)$ 上单调减 \\
\cline{2-3}
 & ${f}^{\prime }\left( x\right)  = 0$ & $f\left( x\right)$ 在 $\left( {a,b}\right)$ 上是常数函数 \\
\cline{1-3}
\hline
\end{tabular}
\end{center}

2. 利用导数判断函数单调性的步骤:

第 1 步，确定函数的定义域；

第 2 步，求出导函数 ${f}^{\prime }\left( x\right)$ 的零点；

第 3 步,用 ${f}^{\prime }\left( x\right)$ 的零点将 $f\left( x\right)$ 的定义域划分为若干个区间,列表给出 ${f}^{\prime }\left( x\right)$ 在各区间上的正负,由此得出函数 $y = f\left( x\right)$ 在定义域内的单调性.

(1)若函数 $f\left( x\right)$ 在区间 $\left( {a,b}\right)$ 上递增，则 ${f}^{\prime }\left( x\right)  \geq  0$ ，所以 “ ${f}^{\prime }\left( x\right)  > 0$ 在 $\left( {a,b}\right)$ 上成立” 是 “ $f\left( x\right)$ 在 $\left( {a,b}\right)$ 上单调增”的充分不必要条件.

(2)研究含参数的函数的单调性，要依据参数对不等式解集的影响进行分类讨论.

若导函数为二次函数式，首先看能否因式分解，再讨论二次项系数的正负及两根的大小；若不能因式分解,则需讨论判别式 $\Delta$ 的正负,二次项系数的正负,两根的大小及根是否在定义域内.

(3)个别导数为 0 的点不影响所在区间的单调性，如 $f\left( x\right)  = {x}^{3}$ ， ${f}^{\prime }\left( x\right)  = 3{x}^{2} \geq  0$ ( ${f}^{\prime }\left( x\right)  = 0$ 在 $x = 0$ 时取到), $f\left( x\right)$ 在 $\mathbf{R}$ 上是严格增函数.

3. 根据函数单调性求参数的一般思路:

(1)利用集合间的包含关系处理:

$y = f\left( x\right)$ 在 $\left( {a,b}\right)$ 上单调,则区间 $\left( {a,b}\right)$ 是相应单调区间的子集.

(2) $f\left( x\right)$ 为增 (减) 函数的充要条件是对任意的 $x \in  \left( {a,b}\right)$ 都有 ${f}^{\prime }\left( x\right)  \geq  0\left( {{f}^{\prime }\left( x\right)  \leq  0}\right)$ ，且在 $\left( {a,b}\right)$ 内的任一非空子区间上， ${f}^{\prime }\left( x\right)$ 不恒为零，应注意此时式子中的等号不能省略，否则会漏解.

(3)函数在某个区间上存在单调区间可转化为不等式有解问题.

\section*{七、函数的极值与最值}

1. 函数的极值

(1)对于函数 $y = f\left( x\right)$ ，如果在 $x = {x}_{1}$ 附近存在一个小区间，使得在该区间内， $f\left( {x}_{1}\right)$ 是最大值，那么， 就称 $y = f\left( x\right)$ 有极大值 $f\left( {x}_{1}\right) ,x = {x}_{1}$ 是极大值点. 从图像上说，极大值点附近是 “山峰”，极大值实际上是 “局部的最大值”，但不一定是整体的最大值.

类似地,如果在 $x = {x}_{2}$ 附近存在一个小区间,使得在该区间内, $f\left( {x}_{2}\right)$ 是最小值,那么,就称 $y = f\left( x\right)$ 有极小值 $f\left( {x}_{2}\right) ,x = {x}_{2}$ 是极小值点. 从图像上说,极小值点附近是 “山谷”,极小值实际上是 “局部的最小值”, 但不一定是整体的最小值.

(2)极大值和极小值统称为极值，极大值点和极小值点统称为极值点.

(3)在导数存在的前提下，“ $x = {x}_{0}$ 是函数的驻点”是 “ $x = {x}_{0}$ 是函数的极值点”的必要非充分条件.

2. 运用导数求函数 $f\left( x\right)$ 极值的一般步骤:

(1)确定函数 $f\left( x\right)$ 的定义域;

(2)求导数 ${f}^{\prime }\left( x\right)$ ；

(3)解方程 ${f}^{\prime }\left( x\right)  = 0$ ，求出函数定义域内的所有根；

(4)列表检验 ${f}^{\prime }\left( x\right)$ 在 ${f}^{\prime }\left( x\right)  = 0$ 的根 ${x}_{0}$ 左右两侧值的符号;

(5)求出极值.

注: \ding{172} 已知函数极值，确定函数解析式中的参数时，要注意:根据极值点的导数为 0 和极值这两个条件列方程组，利用待定系数法求解.  \ding{173} 导数值为 0 不是此点为极值点的充要条件，所以用待定系数法求解后必须检验.

3. 函数的最大 (小) 值

函数 $f\left( x\right)$ 在区间 $\left\lbrack  {a,b}\right\rbrack$ 上有最值的条件:

如果在区间 $\left\lbrack  {a,b}\right\rbrack$ 上函数 $y = f\left( x\right)$ 的图像是一条连续不断的曲线,那么它必有最大值和最小值.

4. 利用导数求函数 $f\left( x\right)$ 在 $\left\lbrack  {a,b}\right\rbrack$ 上的最值的一般步骤:

(1)求函数在 $\left( {a,b}\right)$ 内的极值.

(2)求函数在区间端点处的函数值 $f\left( a\right) ,f\left( b\right)$ .

(3)将函数 $f\left( x\right)$ 的各极值与 $f\left( a\right) ,f\left( b\right)$ 比较,其中最大的一个为最大值,最小的一个为最小值.

\section*{八、不等式证明的常用思路}

1. 移项构造函数法: 证明不等式 $f\left( x\right)  > g\left( x\right)$ (或 $f\left( x\right)  < g\left( x\right)$ ) 转化为证明 $f\left( x\right)  - g\left( x\right)  > 0$ (或 $f\left( x\right)  - g\left( x\right)  < 0)$ ,进而构造辅助函数 $h\left( x\right)  = f\left( x\right)  - g\left( x\right)$ ;

2. 适当放缩构造法:一是根据已知条件适当放缩；二是利用常见放缩结论；

3. 构造“形似”函数，稍作变形再构造，对原不等式同解变形，根据相似结构构造辅助函数；

4. 双变量不等式的处理策略:含两个变量的不等式，基本的思路是将之转化为一元的不等式，具体转化方法主要有三种:整体代换，分离变量，选取主元.

\section*{九、不等式成立问题常用的转化规则(注意等号的相对性)}

单变量不等式成立问题:一般利用参变分离法求解函数不等式恒 (能) 成立.

1. 对任意 $x \in  D,m \leq  f\left( x\right)  \Leftrightarrow  m \leq  f{\left( x\right) }_{\min }$ ;

2. 对任意 $x \in  D,m \geq  f\left( x\right)  \Leftrightarrow  m \geq  f{\left( x\right) }_{\min }$ ;

3. 存在 $x \in  D,m \leq  f\left( x\right)  \Leftrightarrow  m \leq  f{\left( x\right) }_{\min }$ ;

4. 存在 $x \in  D,m \geq  f\left( x\right)  \Leftrightarrow  m \geq  f{\left( x\right) }_{\min }$ .

\section*{$\text{ § }{5.1}$ 导数的概念}

\section*{典型例题精讲}

【例 1】求下列函数的导数.

(1) $y = {x}^{2}\sin x$ ; (2) $y = \ln x + \frac{1}{x}$ ;

(3) $y = \sin \left( {{2x} + \frac{\pi }{3}}\right)$ ； (4) $y = \ln \left( {{2x} - 5}\right)$ .

【例 2】已知某质点的运动方程为 $s\left( t\right)  = 3{t}^{2} + {2t} + 1$ ( $s$ 的单位为 ${m},t$ 的单位为 ${s}$ ).

(1)求从 $t = {2{s}}$ 到 $t = \left( {2 + {\Delta t}}\right) {s}$ 的平均速度；

(2)求当 $t = {2{s}}$ 时的瞬时速度.

【例 3】设 $f\left( 5\right)  = 5,{f}^{\prime }\left( 5\right)  = 3,g\left( 5\right)  = 4,{g}^{\prime }\left( 5\right)  = 1,h\left( x\right)  = \frac{f\left( x\right)  + 2}{g\left( x\right) }$ .

(1)求 $h\left( 5\right)$ 及 ${h}^{\prime }\left( 5\right)$ ；

(2)求曲线 $y = h\left( x\right)  + \sin \frac{\pi }{6}$ 在 $x = 5$ 处的切线方程.

【例 4】已知曲线 $S : y = {2x} - {x}^{3}$ .

(1)求曲线 $S$ 在点 $A\left( {1,1}\right)$ 处的切线方程；

(2)求过点 $B\left( {2,0}\right)$ 并与曲线 $S$ 相切的直线方程.


【例 5】记 ${f}^{\prime }\left( x\right) \text{ 、 }{g}^{\prime }\left( x\right)$ 分别为函数 $f\left( x\right) \text{ 、 }g\left( x\right)$ 的导函数. 若存在 ${x}_{0} \in  \mathbf{R}$ ,满足 $f\left( {x}_{0}\right)  = g\left( {x}_{0}\right)$ 且 ${f}^{\prime }\left( {x}_{0}\right)  = {g}^{\prime }\left( {x}_{0}\right)$ ,则称 ${x}_{0}$ 为函数 $f\left( x\right)$ 与 $g\left( x\right)$ 的一个“ $S$ 点”.

(1)证明:函数 $f\left( x\right)  = x$ 与 $g\left( x\right)  = {x}^{2} + {2x} - 2$ 不存在 “ $S$ 点”；

(2)若函数 $f\left( x\right)  = a{x}^{2} - 1$ 与 $g\left( x\right)  = \ln x$ 存在 “ $S$ 点”，求实数 $a$ 的值.

\section*{四基达标精练}

\section*{一、填空题}

1. 已知 $f\left( x\right)  = {2}^{x}$ ,则 ${f}^{\prime }\left( 1\right)  =$ \_\_\_\_\_.

2. 已知函数 $y = f\left( x\right)$ ，满足 $f\left( x\right)  = \ln \left( {2 - {3x}}\right)$ ，则它的导函数 ${y}^{\prime } =$ \_\_\_\_\_.

3. 函数 $f\left( x\right)  = x + \frac{1}{{x}^{2}}$ 的驻点为 $x =$ \_\_\_\_\_.

4. 已知直线 $y = x + 2$ 是曲线 $y = f\left( x\right)$ 在点 $\left( {0,2}\right)$ 处的切线,则 $\mathop{\lim }\limits_{{h \rightarrow  0}}\frac{f\left( {0 + h}\right)  - f\left( 0\right) }{h} =$ \_\_\_\_\_.

5. 已知函数 $f\left( x\right)$ 为定义在 $\mathbf{R}$ 上的偶函数,且当 $x \geq  0$ 时, $f\left( x\right)  = x + 2{e}^{x}$ ,则函数 $f\left( x\right)$ 在 $x =  - 1$ 处的切线斜率为\_\_\_\_\_.

6. 已知函数 $f\left( x\right)  = a{x}^{3} + {bx} + 1$ 的图像在点 $\left( {1,f\left( 1\right) }\right)$ 处的切线方程为 $y = {4x} - 1$ ，则 ${ab} =$ \_\_\_\_\_.

7. 若 $f\left( x\right)  = \ln x$ 与 $g\left( x\right)  = {x}^{2} + {3x} + a$ 两个函数的图像有一条与直线 $y = x$ 平行的公共切线，则实数 $a =$

\begin{center}
\includegraphics[width=0.2\textwidth]{images/bo_d7tk7oc91nqc73875fu0_83_1213_1667_257_298_0.jpg}
\end{center}
 

(第 8 题图)

8. 某高台跳水运动员在运动过程中的重心相对于水面的高度 $h$ (单位: ${m}$ ) 与跳起后的时间 $t$ (单位: ${s}$ ) 存在函数关系 $h\left( t\right)  =  - {4.9}{t}^{2} + {4.8t} + {11},h\left( t\right)$ 的图像如图所示,已知曲线 $h\left( t\right)$ 在 $t = {t}_{0}$ 处的切线 ${l}_{0}$ 平行于 $t$ 轴,根据图像,给出下列四个结论:

 \ding{172}  在 $t = {t}_{0}$ 时高度 $h$ 关于 $t$ 时间的瞬时变化率为 0 ;

 \ding{173} 曲线 $h\left( t\right)$ 在 $t = {t}_{2}$ 附近比在 $t = {t}_{1}$ 附近下降得慢;

 \ding{174} 曲线 $h\left( t\right)$ 在 $t = {t}_{3}$ 附近比在 $t = {t}_{4}$ 附近上升得快;

 \ding{175} 设在 $t = {t}_{2}$ 和 $t = {t}_{4}$ 时该运动员的瞬时速度分别为 ${v}_{2}$ 、 ${v}_{4}\left( {{m}/{s}}\right)$ ，

则 $\left| {v}_{2}\right|  > \left| {v}_{4}\right|$ . 其中所有正确结论的序号是\_\_\_\_\_.

\begin{center}
\includegraphics[width=0.2\textwidth]{images/bo_d7tk7oc91nqc73875fu0_84_1132_176_300_312_0.jpg}
\end{center}
 

(第 9 题图)

9. 如图,在墙角有一根长 $1{m}$ 的直木棒 ${AB}$ 紧贴墙面,墙面与底面垂直. 在 $t =  0{s}$ 时,木棒的端点 $A$ 从离地 $1{m}$ 的高度以 $1{m}/{s}$ 的速度竖直向下匀速运动,端点 $B$ 向右沿直线运动，则端点 $B$ 在 $t = \frac{\sqrt{2}}{2}{s}$ 这一时刻的瞬时速度为\_\_\_\_\_ ${m}/{s}$ .

\section*{二、选择题}

10. 下列函数的求导正确的是 ( )

A. ${\left( \ln {e}\right) }^{\prime } = \frac{1}{{e}}$ . B. ${\left( x\cos x\right) }^{\prime } = \cos x - x\sin x$ .

C. ${\left( {x}^{-2}\right) }^{\prime } =  - {2x}$ . D. ${\left( 2{{e}}^{x}\right) }^{\prime } = {{e}}^{x}$ .

11. 若函数 $y = f\left( x\right)$ 在 $x = {x}_{0}$ 处导数为 ${f}^{\prime }\left( {x}_{0}\right)$ ,则 $\mathop{\lim }\limits_{{h \rightarrow  0}}\frac{f\left( {{x}_{0} - {3h}}\right)  - f\left( {{x}_{0} + h}\right) }{h}$ 等于 ( )

A. $- {f}^{\prime }\left( {x}_{0}\right)$ . B. $3{f}^{\prime }\left( {x}_{0}\right)$ . C. $- 3{f}^{\prime }\left( {x}_{0}\right)$ . D. $- 4{f}^{\prime }\left( {x}_{0}\right)$ .

12. 已知函数 $y = f\left( x\right)$ ,其导函数 $y = {f}^{\prime }\left( x\right)$ 的图像如图所示. 如图四个选项中,可能表示函数 $y = f\left( x\right)$ 图像的是 ( )

yng

$y = {f}^{\prime }\left( x\right)$

-1 0 1 $x$

\begin{center}
\includegraphics[width=0.2\textwidth]{images/bo_d7tk7oc91nqc73875fu0_84_150_818_217_203_0.jpg}
\end{center}
 

A.

y

-1 0 1 $x$

B.

y

-1 0 1 $x$

C.

(第 12 题图)

$y$

-1 o $x$

D.

\section*{三、解答题}

13. 下列函数的导数.

(1) $y = {\left( 2x - 1\right) }^{3}$ ; (2) $y = x - \sin x\cos x$ ； (3) $y = \frac{\cos x}{{{e}}^{x}}$ .

14. 求下列各式的值.

(1) $f\left( x\right)  = {x}^{2} + 2\cos x$ ，求 ${f}^{\prime }\left( \frac{\pi }{2}\right)$ ；

(2)已知函数 $f\left( x\right)  = {{e}}^{x} + {2x}$ ，求 $\mathop{\lim }\limits_{{h \rightarrow  0}}\frac{{{e}}^{h} + {2h} - 1}{h}$ 的值.

15. 已知函数 $f\left( x\right)  = {x}^{3} - {x}^{2} - \frac{2}{3}{f}^{\prime }\left( {-1}\right) x$ ,其中 ${f}^{\prime }\left( x\right)$ 是 $f\left( x\right)$ 的导函数.

(1)求 ${f}^{\prime }\left( {-1}\right)$ ；

(2)求过原点与曲线 $y = f\left( x\right)$ 相切的切线方程.

16. 如图，酒杯的形状为倒立的圆锥，杯深 $8{{cm}}$ ，上口宽 $6{{cm}}$ ，水以 ${20}{{{cm}}}^{3}/{s}$ 的流量倒入杯中.

(1)给出水深 $h\left( {{cm}}\right)$ 和时间 $t\left( {s}\right)$ 的函数关系式;

(2)当水深为 $4{{cm}}$ 时,求水面升高的瞬时变化率.

\begin{center}
\includegraphics[width=0.2\textwidth]{images/bo_d7tk7oc91nqc73875fu0_85_1293_1099_182_239_0.jpg}
\end{center}
 

(第 16 题图)

\section*{§ 5.2 导数的应用一}

\section*{典型例题精讲}

【例 1】已知函数 $f\left( x\right)  = x - \ln x$ .

(1)求定义域及单调区间；

(2)求 $f\left( x\right)$ 的极值.

【例 2】已知函数 $f\left( x\right)  = x{\left( x - c\right) }^{2}$ 在 $x =  - 1$ 处取得极小值.

(1)求 $c$ 的值；

(2)求 $f\left( x\right)$ 在区间 $\left\lbrack  {-4,0}\right\rbrack$ 上的最值.

【例 3】设函数 $f\left( x\right)  = x - {x}^{3}{{e}}^{{ax} + b}$ ,曲线 $y = f\left( x\right)$ 在点 $\left( {1,f\left( 1\right) }\right)$ 处的切线方程为 $y =  - x + 1$ .

(1)设 $g\left( x\right)  = {f}^{\prime }\left( x\right)$ ，求 $g\left( x\right)$ 的单调区间；

(2)求 $f\left( x\right)$ 的极值点的个数.

【例 4】已知函数 $f\left( x\right)  = \ln x - a{x}^{2} + 1\left( {a \in  \mathbf{R}}\right)$ .

(1)求 $f\left( x\right)$ 的单调区间；

(2)若 $f\left( x\right)$ 有两个零点，求 $a$ 的取值范围.

【例 5】已知函数 $f\left( x\right)  = \left( {x + a}\right)  \cdot  {{e}}^{x}$ .

(1)当 $a = 1$ 时，求函数 $f\left( x\right)$ 的单调区间；

(2)过点 $O\left( {0,0}\right)$ 可作曲线 $y = f\left( x\right)$ 的两条切线,求实数 $a$ 的取值范围.

\section*{四基达标精练}

\section*{一、填空题}

1. 函数 $y = \frac{1}{x} - \ln x$ 的单调递减区间为\_\_\_\_\_.

2. 已函数 $f\left( x\right)  = \ln x - x + 1$ 的最大值为\_\_\_\_\_.

3. 已知函数 $f\left( x\right)  = {ax} - a - \ln x$ ，且 $f\left( x\right)$ 的最小值为 0，则实数 $a =$ \_\_\_\_\_.

4. 函数 $f\left( x\right)  = x - \ln x$ 在区间 $\left( {0,{e}}\right)$ 上的极小值为\_\_\_\_\_.

5. 函数 $f\left( x\right)  = \ln x - \left| {x - 2}\right|$ 的最大值为\_\_\_\_\_.

6. 已知函数 $f\left( x\right)  = \frac{1}{2}{x}^{2} - \left( {a + 2}\right) x + {2a}\ln x + 1$ 在 $\left( {4,6}\right)$ 上存在极值点,则实数 $a$ 的取值范围是\_\_\_\_\_.

\begin{center}
\includegraphics[width=0.2\textwidth]{images/bo_d7tk7oc91nqc73875fu0_88_1176_1667_287_222_0.jpg}
\end{center}
 

(第 8 题图)

7. 函数 $f\left( x\right)  = \frac{1}{3}{x}^{3} + a{x}^{2} + \left( {a + 2}\right) x - 1$ 有极大值又有极小值,则实数 $a$ 的范围是\_\_\_\_\_.

8. 已知直线 ${AB}$ 是函数 $y = f\left( x\right)$ 图像的一条割线 (如图所示), ${f}^{\prime }\left( x\right)$ 是函数 $f\left( x\right)$ 的导函数,若 $a = 2{f}^{\prime }\left( 3\right) \text{ 、 }b = 2{f}^{\prime }\left( 5\right) \text{ 、 }c = f\left( 5\right)  - f\left( 3\right)$ ,则关于 $a\text{ 、 }b\text{ 、 }c$ 从小到大的顺序是\_\_\_\_\_.

9. 已知函数 $f\left( x\right)  = \sin x - 1,g\left( x\right)  = a\ln x - x$ ,若对任意 ${x}_{1} \in  \mathbf{R}$ 都存在 ${x}_{2} \in  \left( {1,{e}}\right)$ ,使 $f\left( {x}_{1}\right)  < g\left( {x}_{2}\right)$ 成立，则实数 $a$ 的取值范围是\_\_\_\_\_.

\section*{二、选择题}

10. 函数 $f\left( x\right)  = \frac{x - 2}{{{e}}^{x}}$ 的单调递增区间为 ( )

A. $\left( {-\infty ,3}\right)$ . B. $\left( {0,3}\right)$ . C. $\left( {3, + \infty }\right)$ . D. $\left( {-\infty ,2}\right)$ .

11. 函数 $y = \frac{x}{{{e}}^{x}}$ 的极值点是 ( )

A.0. B.1. C. $\left( {0,1}\right)$ . D. $\left( {1,0}\right)$ .

12. 函数 $f\left( x\right)  = a{x}^{3} - {bx} - \ln x\left( {a > 0,b \in  \mathbf{R}}\right)$ 的一个极值点是 1,则 $\ln a - b + 1$ 的值 ( )

A. 恒大于 0 . C. 恒等于 0 . D. 不确定.

\section*{三、解答题}

13. 给出函数 $f\left( x\right)  = {x}^{3} - 4{x}^{2} + {4x}$ .

(1)求 $f\left( x\right)$ 的单调区间；

(2)求 $f\left( x\right)$ 在 $\left\lbrack  {0,3}\right\rbrack$ 上的最值.

14. 已知函数 $f\left( x\right)  = a{x}^{2} + b\ln x$ 在 $x = 1$ 处有极值 $\frac{1}{2}$ .

(1)求a、b的值；

(2)求函数 $y = f\left( x\right)$ 的单调性.

15. 已知函数 $f\left( x\right)  = \ln x - {ax} + 1\left( {a \in  \mathbf{R}}\right)$ .

(1)讨论 $f\left( x\right)$ 的单调性;

(2)若 $f\left( x\right)  \leq  0$ 恒成立，求 $a$ 的取值范围.

16. 设 $a > 0$ ,已知函数 $f\left( x\right)  = {\left( x - 2\right) }^{3} - {ax}$ .

(1)求函数 $y = f\left( x\right)$ 的单调区间；

(2)对于函数 $y = f\left( x\right)$ 的极值点 ${x}_{0}$ ，存在 ${x}_{1}\left( {{x}_{1} \neq  {x}_{0}}\right)$ ，使得 $f\left( {x}_{1}\right)  = f\left( {x}_{0}\right)$ ，试问对任意的正数 $a$ ， ${x}_{1} + 2{x}_{0}$ 是否为定值？若是，求出这个定值；若不是，请说明理由.

§ 5.3 导数的应用二

\section*{典型例题精讲}

【例 1】已知函数 $f\left( x\right)  = \left( {x + a}\right) {{e}}^{x}$ .

(1)若 $f\left( x\right)$ 在 $x = 1$ 处取得极小值，求实数 $a$ 的值；

(2)若 $f\left( x\right)$ 在 $\left( {-1,1}\right)$ 上单调递增，求实数 $a$ 的取值范围.

【例 2】已知函数 $f\left( x\right)  = \frac{1}{3}{x}^{3} - {x}^{2} + a$ .

(1)当 $a = 0$ 时，求函数 $f\left( x\right)$ 的极大值与极小值；

(2)若函数 $f\left( x\right)$ 在 $\left\lbrack  {1,3}\right\rbrack$ 上的最大值是最小值的 3 倍，求 $a$ 的值.

【例 3】已知 $x = 3$ 是函数 $f\left( x\right)  = a\ln x + {x}^{2} - {10x}$ 的一个极值点.

(1)求函数 $f\left( x\right)$ 的单调区间；

(2)若直线 $y = b$ 与 $y = f\left( x\right)$ 图像有 3 个交点，求 $b$ 的取值范围.

【例 4】已知函数 $f\left( x\right)  = {{e}}^{x} + m{{e}}^{-x} + \left( {m - 1}\right) x\left( {m \leq  0}\right)$ .

(1)讨论函数 $f\left( x\right)$ 的单调性;

(2)若对任意的 $x \in  \left( {0, + \infty }\right)$ ， $f\left( x\right)  \geq   - 2$ 恒成立，求 $m$ 的取值范围.


【例 5】已知函数 $f\left( x\right)  = \ln x + \left( {a - 1}\right) x + a + 1\left( {a \in  \mathbf{R}}\right)$ .

(1)讨论函数 $f\left( x\right)$ 的极值情况；

(2)证明:当 $a \leq  1$ 时， ${{e}}^{x} - f\left( x\right)  > 0$ .

\section*{四基达标精练}

\section*{一、填空题}

1. 已知函数 $f\left( x\right)  = \ln x + k$ 在 $\left\lbrack  {1,{e}}\right\rbrack$ 上的最大值为 2，则 $f\left( k\right)  =$ \_\_\_\_\_.

2. 已知函数 $f\left( x\right)  = {x}^{2} + \frac{a}{x}$ 在区间 $\left\lbrack  {1,3}\right\rbrack$ 上单调增，则实数 $a$ 的取值范围是\_\_\_\_\_.

3. 若函数 $f\left( x\right)  = 2{x}^{3} - 6{x}^{2} + m$ 在区间 $\left\lbrack  {-2,2}\right\rbrack$ 上有最大值 3,则 $f\left( x\right)$ 在区间 $\left\lbrack  {-2,2}\right\rbrack$ 上的最小值为\_\_\_\_\_.

4. 若函数 $f\left( x\right)  = {ax} - \ln x,x \in  \left( {1, + \infty }\right)$ 是严格增函数，则实数 $a$ 的最小值是\_\_\_\_\_.

5. 已知函数 $f\left( x\right)  = {x}^{3} - x$ 有 2 个极值点 ${x}_{1}\text{ 、 }{x}_{2}$ ,则 ${x}_{1} + {x}_{2} + f\left( {x}_{1}\right)  + f\left( {x}_{2}\right)  =$ \_\_\_\_\_.

6. 已知 $f\left( x\right)  = \ln x + {x}^{2} - {bx} + b$ 在区间 $\left( {1,2}\right)$ 上不单调，则实数 $b$ 的取值范围是\_\_\_\_\_.

7. 已知函数 $f\left( x\right)  = \ln x - \frac{1}{2}{x}^{2} - a$ ,如果存在 $x > 0$ ,使得 $f\left( x\right)  \geq  0$ ,则实数 $a$ 的取值范围

是

8. 已知函数 $f\left( x\right)  = {{e}}^{x} + m\ln x$ ( $m \in  \mathbf{R}$ , e 为自然对数的底数),若对任意正数 ${x}_{1}\text{ 、 }{x}_{2}$ ,当 ${x}_{1} > {x}_{2}$ 时都有 $f  \left( {x}_{1}\right)  - f\left( {x}_{2}\right)  > {x}_{1} - {x}_{2}$ 成立，则实数 $m$ 的取值范围是\_\_\_\_\_.

9. 已知点 $P$ 是曲线 $y = {x}^{2} - \ln x$ 上一点,求点 $P$ 到直线 $y = x - 3$ 的最小距离为\_\_\_\_\_.

\section*{二、选择题}

10. 设 $a \in  \mathbf{R}$ ,若函数 $f\left( x\right)  = {{e}}^{x} + {ax},x \in  \mathbf{R}$ 有大于零的极值点,则 $a$ 的取值范围为 ( )

A. $a <  - 1$ . B. $a <  - \frac{1}{{e}}$ C. $a >  - \frac{1}{{e}}$ . D. $a >  - 1$ .

11. 已知函数 $f\left( x\right)  = \left( {x - 1}\right) \ln x$ ,下列说法中正确的是 ( )

A. 1 既是 $y = f\left( x\right)$ 的一个零点,又是 $y = f\left( x\right)$ 的一个极小值点.

B. 1 既是 $y = f\left( x\right)$ 的一个零点,又是 $y = f\left( x\right)$ 的一个极大值点.

C. 1 是 $y = f\left( x\right)$ 的一个零点,不是 $y = f\left( x\right)$ 的极值点.

D. 1既不是 $y = f\left( x\right)$ 的一个零点,也不是 $y = f\left( x\right)$ 的极值点.

12. 函数 $f\left( x\right)$ 定义域为 $\left( {0, + \infty }\right)$ ,其导函数为 ${f}^{\prime }\left( x\right)$ ,若 $\forall x \in  \left( {0, + \infty }\right) ,x{f}^{\prime }\left( x\right)  + f\left( x\right)  < 1$ ,且 $f\left( 1\right)  =$ 2,则不等式 ${xf}\left( x\right)  - x < 1$ 的解集为 ( )

A. $\left( {1,2}\right)$ . B. $\left( {2, + \infty }\right)$ . C. $\left( {0,1}\right)$ . D. $\left( {1, + \infty }\right)$ .

\section*{三、解答题}

13. 已知函数 $y = f\left( x\right)  = a{x}^{3} - {x}^{2} + x - 5$ .

(1)若 $a =  - 1$ ，求 $y = f\left( x\right)$ 的极小值；

(2)若 $y = f\left( x\right)$ 在 $\mathbf{R}$ 上无极值，求 $a$ 的取值范围.

14. 设 $f\left( x\right)  =  - \frac{1}{3}{x}^{3} + \frac{1}{2}{x}^{2} + {ax}$ .

(1)若 $f\left( x\right)$ 在 $\left( {\frac{2}{3}, + \infty }\right)$ 上存在单调增区间，求 $a$ 的取值范围；

(2)若 $f\left( x\right)$ 在 $\left( {\frac{2}{3}, + \infty }\right)$ 上单调减，求 $a$ 的取值范围.

15. 已知函数 $f\left( x\right)  = \frac{1}{3}{x}^{3} - b{x}^{2} + {2x} - a,x = 2$ 是 $f\left( x\right)$ 的一个极值点.

(1)求 $f\left( x\right)$ 的单调增区间；

(2)当 $a > 0$ 时,求方程 $f\left( x\right)  = 0$ 的解的个数.

16. 已知函数 $f\left( x\right)  = \frac{1}{2}{x}^{2} + a\ln x$ .

(1)讨论 $f\left( x\right)$ 的单调性；

(2)若 $a = 1$ ，函数 $g\left( x\right)  = f\left( x\right)  - {mx}$ 是其定义域上的增函数，求实数 $m$ 的取值范围.

\section*{第 6 章 三角与三角函数(必修 2)}

\section*{知识梳理，}

\section*{一、任意角及其三角比}

1. 与 $\alpha$ 角终边相同、关于 $x$ 轴、 $y$ 轴、原点对称的角的集合分别是 $\{ \beta  \mid  \beta  = {2k\pi } + \alpha ,k \in  \mathbf{Z}\}$ 、 $\left\{  {\beta  \mid  \beta  = {2k\pi } - \alpha }\right\}$ ， $k \in  \mathbf{Z}\} \text{ 、 }\{ \beta  \mid  \beta  = {2k\pi } + \pi  - \alpha ,k \in  \mathbf{Z}\} \text{ 、 }\{ \beta  \mid  \beta  = {2k\pi } + \pi  + \alpha ,k \in  \mathbf{Z}\} .$

2. 弧度制: $\left| \alpha \right|  = \frac{l}{r}$ ( $\left| \alpha \right|$ 是将 $\left| \alpha \right|$ 看成圆心角时角 $\left| \alpha \right|$ 的弧度数， $r$ 是半径， $l$ 是弧长).

3. 半径为 $r$ 、圆心角为 $\alpha$ (弧度) (或 $\left. {n}^{ \circ  }\right)$ 的扇形弧长 $l = {r\alpha } = \frac{n\pi r}{180}$ ，面积 ${S}_{\text{ 扇 }} = \frac{1}{2}{lr} = \frac{1}{2}{r}^{2}\alpha  = \frac{{n\pi }{r}^{2}}{360}$ .

4. 设 $P\left( {x,y}\right)$ 是角 $\alpha$ 终边上任意一点 (异于原点) $\left| {OP}\right|  = \sqrt{{x}^{2} + {y}^{2}}$ ,则 $\sin \alpha  = \frac{y}{r};\cos \alpha  = \frac{x}{r};\tan \alpha  = \frac{y}{x}$ ; $\cot \alpha  = \frac{x}{y}$ ; 各三角比的符号由点 $P$ 所在的象限确定.

\section*{二、三角恒等变形}

1. 同角三角比基本关系: ${\sin }^{2}\theta  + {\cos }^{2}\theta  = 1,1 + {\tan }^{2}\theta  = \frac{1}{{\cos }^{2}\theta },1 + {\cot }^{2}\theta  = \frac{1}{{\sin }^{2}\theta },\sin \theta  = \tan \theta \cos \theta ,\cos \theta  =  \cot \theta \sin \theta ,\tan \theta  \cdot  \cot \theta  = 1$ .

2. 诱导公式: $n \cdot  \frac{\pi }{2} + \alpha \left( {n \in  \mathbf{Z}}\right)$ 变 $\alpha$ 的三角比遵循 “ $\left( n\right)$ 奇变偶不变，符号看象限(原象限原三角比)”.

3. 和角与差角三角比公式: $\sin \left( {\alpha  \pm  \beta }\right)  = \sin \alpha \cos \beta  \pm  \cos \alpha \sin \beta$ ;

$\cos \left( {\alpha  \pm  \beta }\right)  = \cos \alpha \cos \beta  \mp  \sin \alpha \sin \beta ,\tan \left( {\alpha  \pm  \beta }\right)  = \frac{\tan \alpha  \pm  \tan \beta }{1 \mp  \tan \alpha \tan \beta }.$

4. 辅助角公式 $a\sin \alpha  + b\cos \alpha  = \sqrt{{a}^{2} + {b}^{2}}\sin \left( {\alpha  + \varphi }\right)$ (其中 $\cos \varphi  = \frac{a}{\sqrt{{a}^{2} + {b}^{2}}},\sin \varphi  = \frac{b}{\sqrt{{a}^{2} + {b}^{2}}}$ ).

5. 二倍角公式: $\sin {2\alpha } = 2\sin \alpha \cos \alpha ;\cos {2\alpha } = {\cos }^{2}\alpha  - {\sin }^{2}\alpha  = 2{\cos }^{2}\alpha  - 1 = 1 - 2{\sin }^{2}\alpha$ .

6. 降幂公式: ${\sin }^{2}\alpha  = \frac{1 - \cos {2\alpha }}{2}$ ; ${\cos }^{2}\alpha  = \frac{1 + \cos {2\alpha }}{2}$ .

7. 万能公式: $\sin {2\alpha } = \frac{2\tan \alpha }{1 + {\tan }^{2}\alpha };\cos {2\alpha } = \frac{1 - {\tan }^{2}\alpha }{1 + {\tan }^{2}\alpha };\tan {2\alpha } = \frac{2\tan \alpha }{1 - {\tan }^{2}\alpha }$ .

\section*{三、解三角形}

1. $\bigtriangleup {ABC}$ 中,正弦定理: $\frac{a}{\sin A} = \frac{b}{\sin B} = \frac{c}{\sin C} = {2R}$ . 其中 $R$ 为 $\bigtriangleup {ABC}$ 外接圆半径;

2. 余弦定理: ${a}^{2} = {b}^{2} + {c}^{2} - {2bc}\cos A;{b}^{2} = {c}^{2} + {a}^{2} - {2ca}\cos B;{c}^{2} = {a}^{2} + {b}^{2} - {2ab}\cos C$ .

$\cos A = \frac{{b}^{2} + {c}^{2} - {a}^{2}}{2bc};\cos B = \frac{{c}^{2} + {a}^{2} - {b}^{2}}{2ca};\cos C = \frac{{a}^{2} + {b}^{2} - {c}^{2}}{2ab}$ .

3. 三角形的面积公式: \ding{172}  $S = \frac{1}{2}a{h}_{a} = \frac{1}{2}b{h}_{b} = \frac{1}{2}c{h}_{c}$ ( ${h}_{a}$ 、 ${h}_{b}$ 、 ${h}_{c}$ 分别表示 $a$ 、 $b$ 、 $c$ 边上的高)；

\HRule

* 注 本章节所属知识模块是 “函数与分析”.

\HRule

 \ding{173}  $S = \frac{1}{2}{ab}\sin C = \frac{1}{2}{bc}\sin A = \frac{1}{2}{ca}\sin B$ ； \ding{174}  ${S}_{\bigtriangleup {OAB}} = \frac{1}{2}\sqrt{{\left( \left| \overrightarrow{OA}\right|  \cdot  \left| \overrightarrow{OB}\right| \right) }^{2} - {\left( \overrightarrow{OA} \cdot  \overrightarrow{OB}\right) }^{2}}$ ；\_\_\_\_\_

公式的变式: $S = 2{R}^{2}\sin A\sin B\sin C = \frac{abc}{4R}(R$ 是 $\bigtriangleup {ABC}$ 的外接圆半径);

$S = \frac{1}{2}\left( {a + b + c}\right) r$ ( $r$ 是 $\bigtriangleup {ABC}$ 的内切圆半径).

\section*{四、三角函数的图像和性质}

\begin{center}
\begin{tabular}{|c|c|c|c|}
\hline
\phantom{X} & 正弦函数 & 余弦函数 & 正切函数 \\
\cline{1-4}
记 号 & $y = \sin x$ & $y = \cos x$ & $y = \tan x$ \\
\cline{1-4}
图像 & \includegraphics[width=0.2\textwidth]{images/bo_d7tk7oc91nqc73875fu0_97_289_557_326_267_0.jpg} & \includegraphics[width=0.2\textwidth]{images/bo_d7tk7oc91nqc73875fu0_97_674_557_328_266_0.jpg} & \includegraphics[width=0.2\textwidth]{images/bo_d7tk7oc91nqc73875fu0_97_1092_548_261_274_0.jpg} \\
\cline{1-4}
定义域 & R & $\mathbf{R}$ & $\{ x \mid  x \neq  {k\pi } + \frac{\pi }{2},k \in  \mathbf{Z}\}$ \\
\cline{1-4}
值域(最值) & $\left\lbrack  {-1,1}\right\rbrack  \left( {k \in  \mathbf{Z}}\right)  x = {2k\pi } - \frac{\pi }{2}$ 时, ${y}_{\min } =  - 1  x = {2k\pi } + \frac{\pi }{2}$ 时, ${y}_{\max } = 1$ & $\left\lbrack  {-1,1}\right\rbrack  \left( {k \in  \mathbf{Z}}\right)  x = {2k\pi } + \pi ,{y}_{\min } =  - 1  x = {2k\pi }$ 时， ${y}_{\max } = 1$ & R(无最值) \\
\cline{1-4}
单调性 & $\left\lbrack  {{2k\pi } - \frac{\pi }{2},{2k\pi } + \frac{\pi }{2}}\right\rbrack   \uparrow  \left( {k \in  \mathbf{Z}}\right)  \left\lbrack  {{2k\pi } + \frac{\pi }{2},{2k\pi } + \frac{3\pi }{2}}\right\rbrack   \downarrow$ & $\left\lbrack  {{2k\pi } - \pi ,{2k\pi }}\right\rbrack   \uparrow  \left( {k \in  \mathbf{Z}}\right)  \left\lbrack  {{2k\pi },{2k\pi } + \pi }\right\rbrack   \downarrow$ & $\left( {{k\pi } - \frac{\pi }{2},{k\pi } + \frac{\pi }{2}}\right)  \uparrow  \left( {k \in  \mathbf{Z}}\right)$ \\
\cline{1-4}
奇偶性 & 奇函数 & 偶函数 & 奇函数 \\
\cline{1-4}
对称轴 & $x = {k\pi } + \frac{\pi }{2}\left( {k \in  \mathbf{Z}}\right)$ & $x = {k\pi }\left( {k \in  \mathbf{Z}}\right)$ & 无 \\
\cline{1-4}
对称中心 & $\left( {{k\pi },0}\right) \left( {k \in  \mathbf{Z}}\right)$ & $\left( {{k\pi } + \frac{\pi }{2},0}\right) \left( {k \in  \mathbf{Z}}\right)$ & $\left( {\frac{k\pi }{2},0}\right) \left( {k \in  \mathbf{Z}}\right)$ \\
\cline{1-4}
周期性 & ${2\pi }$ & ${2\pi }$ & $\pi$ \\
\cline{1-4}
\hline
\end{tabular}
\end{center}

\section*{五、三角函数图像变换}

1. 变换法作函数 $y = A\sin \left( {{\omega x} + \varphi }\right) \left( {\omega  > 0,A > 0}\right)$ 的步骤:

(1)作函数 $y = \sin x$ 的图像；

(2)将 $y = \sin x$ 的图像右移 $\left( {\varphi  < 0}\right)$ 或左移 $\left( {\varphi  > 0}\right)$ 或 $\left| \varphi \right|$ 个单位得 $y = \sin \left( {x + \varphi }\right)$ 的图像；

(3)将 $y = \sin \left( {x + \varphi }\right)$ 的图像纵标不变，横标伸长(0 $< \omega  < 1$ )或缩短( $\omega  > 1$ )到原来的 $\frac{1}{\omega }$ 得到 $y = \sin \left( {{\omega x} + \varphi }\right)$ 的图像;

(4)将 $y = \sin \left( {{\omega x} + \varphi }\right)$ 的图像横标不变，纵标伸长 (A > 1) 或缩短 (0 $< A < 1$ ) 到原来的 $A$ 倍得 $y =  A\sin \left( {{\omega x} + \varphi }\right)$ 图像;

注意: 由 $y = \sin {\omega x}$ 的图像得到 $y = \sin \left( {{\omega x} + \varphi }\right)$ 的图像时是将 $y = \sin {\omega x}$ 的图像左 (或右) 移 $\left| \frac{\varphi }{\omega }\right|$ 个单

位. (要注意反向平移问题)

2. 已知图像求 $y = A\sin \left( {{\omega x} + \varphi }\right) \left( {A > 0,\omega  > 0}\right)$ 解析式的问题:

(1)根据图像确定 $\omega$ :

方法一:相邻最大值和最小值对应的横坐标之差的绝对值是周期的一半；

方法二:相邻两个零点对应的横坐标之差的绝对值是周期的一半；

方法三:最值与相邻的零点的横坐标之差的绝对值是周期的四分之一.

(2)根据图像确定 $A$ :最大值与最小值之差的一半的绝对值是 $\left| A\right|$ .

(3)根据图像确定 $\varphi$ :递增区间中零点的横坐标 ${x}_{0}$ 对应于 $\omega {x}_{0} + \varphi  = {2k\pi }\left( {k \in  \mathbf{Z}}\right)$ ；递减区间中零点的横坐标 ${x}_{0}$ 对应于 $\omega {x}_{0} + \varphi  = {2k\pi } + \pi \left( {k \in  \mathbf{Z}}\right)$ .

\section*{六、反三角和三角方程}

反三角函数的本质:

$\arcsin x\left( {x \in  \left\lbrack  {0,1}\right\rbrack  }\right)$ 表示 $\left\lbrack  {0,\frac{\pi }{2}}\right\rbrack$ 内正弦值为 $x$ 的角;

$\arccos x\left( {x \in  \left\lbrack  {0,1}\right\rbrack  }\right)$ 表示 $\left\lbrack  {0,\frac{\pi }{2}}\right\rbrack$ 内余弦值为 $x$ 的角;

$\arctan x\left( {x \geq  0}\right)$ 表示 $\left\lbrack  {0,\frac{\pi }{2}}\right)$ 内正切值为 $x$ 的角.

\section*{$\text{ § }{6.1}$ 三角及其恒等变形}

\section*{典型例题精讲}

【例 1】已知角 $\alpha$ 的终边过点 $P\left( {{7m}, - {3m}}\right) \left( {m \neq  0}\right)$ .

(1) 求 $\frac{\cos \left( {\frac{\pi }{2} + \alpha }\right) \sin \left( {-\pi  - \alpha }\right) }{\cos \left( {\frac{11\pi }{2} - \alpha }\right) \sin \left( {\frac{9\pi }{2} + \alpha }\right) }$ 的值；

(2) 求 $\sin \left( {\pi  - \alpha }\right) \cos \left( {{2\pi } - \alpha }\right)  - {\sin }^{2}\left( {\frac{\pi }{2} + \alpha }\right)  + 2$ 的值.

\section*{高分钥匙高考数学复习一轮通}

【例 2】已知是关于 $x$ 的方程 $k{x}^{2} - {kx} + k - 3 = 0$ 的两个实根.

(1)若 ${x}_{1} = \sin \alpha$ ， ${x}_{2} = \cos \alpha \left( {\alpha  \in  \mathbf{R}}\right)$ ，求实数 $k$ 的值；

(2)若 ${x}_{1} = \sin \alpha$ ， ${x}_{2} = \cos \beta \left( {\alpha \text{ 、 }\beta  \in  \mathbf{R}}\right)$ ，求实数 $k$ 的取值范围.

【例 3】已知角 $\theta$ 的顶点为坐标原点 $O$ ,始边为 $x$ 轴的非负半轴,终边与单位圆相交于点 $P\left( {x,y}\right)$ ,若点 $P$ 位于 $x$ 轴上方且 $x + y = \frac{1}{2}$ .

(1)求 $\sin \theta  - \cos \theta$ 的值；

(2)求 ${\sin }^{4}\theta  + {\cos }^{4}\theta$ 的值.

【例 4】已知 $0 < \alpha  < \frac{\pi }{3}, - \frac{\pi }{2} < \beta  <  - \frac{\pi }{3},\cos \left( {\alpha  + \frac{\pi }{6}}\right)  = \frac{\sqrt{3}}{3},\sin \left( {\alpha  - \beta  + \frac{\pi }{6}}\right)  = \frac{2}{3}$ .

(1)求 $\sin \alpha$ 的值；

(2)求 $\cos \left( {{3\alpha } - \beta }\right)$ 的值.

【例 5】已知向量 $\overrightarrow{a} = \left( {\cos \alpha ,\sin \alpha }\right) ,\overrightarrow{b} = \left( {\cos \beta , - \sin \beta }\right) ,\overrightarrow{c} = \left( {2,1}\right) ,\overrightarrow{a}//\overrightarrow{c},\overrightarrow{a} \cdot  \overrightarrow{b} = \frac{\sqrt{5}}{5}$ .

(1)求 $\frac{{\sin }^{2}\alpha  + 2\cos {2\alpha }}{1 + \tan \alpha }$ 的值；

(2)若 $\alpha$ 、 $\beta$ 均为锐角,求 $\tan \left( {\alpha  - \beta }\right)$ 的值.

\section*{四基达标精练}

\section*{一、填空题}

1. 角 $\alpha$ 的终边过点 $P\left( {-{6a}, - {8a}}\right) \left( {a > 0}\right)$ ，则 $\sin \alpha  - \cos \alpha  =$ \_\_\_\_\_.

2. 关于 $x$ 的不等式 $\sin x \leq   - \frac{\sqrt{3}}{2}$ 的解集为\_\_\_\_\_.

3. 若一个扇形的半径是 2,圆心角 $\alpha$ 所对的弦长为 $2\sqrt{3}$ ,则 $\alpha  =$ \_\_\_\_\_.

4. 若 $\cos \alpha  = \frac{2}{3}$ ，且 $\alpha$ 在第四象限，则 $\sin \left( {\alpha  - {2\pi }}\right)  + \sin \left( {-\alpha  - {3\pi }}\right) \cos \left( {\alpha  - {3\pi }}\right)  =$ \_\_\_\_\_.

5. 已知 $\tan \theta  = \frac{1}{3}$ ,则 $\frac{\tan \left( {{2\pi } - \theta }\right) }{\cos \left( {\frac{\pi }{2} - \theta }\right) \sin \left( {\frac{3\pi }{2} + \theta }\right) } =$ \_\_\_\_\_.

6. 已知点 $A$ 的坐标为 $\left( {-3,4}\right)$ ，将 ${OA}$ 绕坐标原点 $O$ 顺时针旋转 $\frac{\pi }{3}$ 至 ${OB}$ . 则点 $B$ 的坐标为\_\_\_\_\_.

7. 已知 $A\left( {{x}_{1},{y}_{1}}\right)$ 、 $B\left( {{x}_{2},{y}_{2}}\right)$ 是角 $\alpha$ 、 $\beta$ 终边与单位圆的两个不同交点，且 ${x}_{1}{y}_{2} = {x}_{2}{y}_{1}$ ，则 $2{x}_{1} - {x}_{2} +  2{y}_{1} - {y}_{2}$ 的最大值为\_\_\_\_\_.

8. 已知长方形的四个顶点: $A\left( {0,0}\right) \text{ 、 }B\left( {2,0}\right) \text{ 、 }C\left( {2,1}\right) \text{ 、 }D\left( {0,1}\right)$ ,一质点从点 $A$ 出发,沿与 ${AB}$ 夹角为 $\theta$ 的方向射到 ${BC}$ 上的点 ${P}_{1}$ 后,依次反射到 ${CD}\text{ 、 }{DA}$ 和 ${AB}$ 上的点 ${P}_{2}\text{ 、 }{P}_{3}\text{ 、 }{P}_{4}$ (入射角等于反射角). 设 ${P}_{4}$ 的坐标为 $\left( {{x}_{4},0}\right)$ ,若 $1 < {x}_{4} < 2$ ,则 $\tan \theta$ 的范围是\_\_\_\_\_.

9. 已知 $a$ 、 $b$ 、 $x \in  \mathbf{R},{ab} \neq  0$ ，且 $\frac{{\sin }^{4}x}{{a}^{2}} + \frac{{\cos }^{4}x}{{b}^{2}} = \frac{1}{{a}^{2} + {b}^{2}}$ ，则 $\frac{{\sin }^{2022}x}{{a}^{2020}} + \frac{{\cos }^{2022}x}{{b}^{2020}} =$ \_\_\_\_\_.

\section*{二、选择题}

10. 在平面直角坐标系 ${xOy}$ 中,角 $\alpha  = {120}^{ \circ  }$ ,其顶点为坐标原点 $O$ ,始边与 $x$ 轴的非负半轴重合,终边过点 $P\left( {-4,m}\right)$ ,则实数 $m$ 的值是 ( )

A.2. $\sqrt{3}$ B. $- 2\sqrt{3}$ C. $4\sqrt{3}$ D. $- 4\sqrt{3}$

11. 在直角坐标平面上,角 $\sin \frac{\pi }{6}$ 所在的象限是 ( )

A. 一 B. 二 C. 三 D. 四

12. 已知 $\sin \frac{q + p}{2} < 0,\tan \frac{q + p}{2} < 0,\cos \frac{p - q}{2} > 0$ ,则下列关系中必定不成立的是 ( )

A. $\sin q < 0,\cos q > 0$ B. $\sin q > 0,\cos q < 0$ C. $\sin q > 0,\cos q > 0$ D. $\sin q < 0,\cos q < 0$

\section*{三、解答题}

13. 已知 $\sin x + \cos x = \frac{1}{5},x \in  \left( {0,\pi }\right)$ .

(1) 求 $\sin x - \cos x$ 的值；

(2)求 $\cos {2x}$ 的值.

14. 已知关于 $x$ 的方程 ${13}{x}^{2} - {7x} + m = 0$ 的两根为 $\sin \theta \text{ 、 }\cos \theta ,\theta  \in  \left( {0,\pi }\right)$ .

(1)求 $m$ 的值；

(2)求 $\tan \theta$ 的值.

15. 已知 $\overrightarrow{a} = \left( {\cos \alpha ,\sin \alpha }\right) ,\overrightarrow{b} = \left( {\cos \beta ,\sin \beta }\right) ,0 < \beta  < \alpha  < \pi$ .

(1)若 $\left| {\overrightarrow{a} - \overrightarrow{b}}\right|  = \sqrt{2}$ ，求证: $\overrightarrow{a}\bot \overrightarrow{b}$ ；

(2)设 $\overrightarrow{c} = \left( {0,1}\right)$ ，若 $\overrightarrow{a} + \overrightarrow{b} = \overrightarrow{c}$ ，求 $\cos \left( {\alpha  - \beta }\right)$ 的值.

16. 如图，在平面直角坐标系中，角 $\alpha$ 和 $\beta$ 的终边与单位圆分别交于 $P$ 、 $Q$ 两点.

(1)若 $\overrightarrow{OP} + \overrightarrow{OQ} = \left( {\frac{1}{2},\frac{3}{2}}\right)$ ，求 $\cos \left( {\alpha  - \beta }\right)$ 的值；

(2)若 $\alpha  = \frac{\pi }{6},\left| {\overrightarrow{OP} - \overrightarrow{OQ}}\right|  = \frac{\sqrt{2}}{3}$ ，求 $\cos \left( {{2\beta } + \frac{2}{3}\pi }\right)$ 的值.

\begin{center}
\includegraphics[width=0.2\textwidth]{images/bo_d7tk7oc91nqc73875fu0_102_1159_1180_281_245_0.jpg}
\end{center}
 

(第 16 题图)

\section*{§ 6.2 解三角形}

\section*{典型例题精讲}

【例 1】在 $\bigtriangleup {ABC}$ 中,角 $A\text{ 、 }B\text{ 、 }C$ 所对的边分别为 $a\text{ 、 }b\text{ 、 }c$ ,已知 ${c}^{2} = {a}^{2} + {b}^{2} + {ab}$ .

(1)求角 $C$ 的大小；

(2)若 $c = 4\sqrt{7},\bigtriangleup {ABC}$ 的周长为 ${12} + 4\sqrt{7}$ ，求 $\sin A + \sin B$ .

【例 2】如图所示,某人为了测量出到河对岸铁塔的距离与铁塔的高,选与塔底 $B$ 同在水平面内的两个观测点 $C$ 与 $D$ ,在 $C$ 点测得塔底 $B$ 在北偏东 ${45}^{ \circ  }$ 方向,然后向正东方向前进 ${10}{m}$ 到达 $D$ ,测得此时塔底 $B$ 在北偏东 ${15}^{ \circ  }$ 方向.

(1)求点 $D$ 到塔底 $B$ 的距离 ${BD}$ ；

(2)若在点 $C$ 测得塔顶 $A$ 的仰角为 ${60}^{ \circ  }$ ，求铁塔 ${AB}$ 的高.

(注:结果保留根号)

\begin{center}
\includegraphics[width=0.2\textwidth]{images/bo_d7tk7oc91nqc73875fu0_103_1184_1447_268_377_0.jpg}
\end{center}
 

(例 2 图)

【例 3】已知 $\bigtriangleup  {ABC}$ 的内角和边满足 $a\cos C + \sqrt{3}a\sin C - b - c = 0$ ，其中角 $A$ 、 $B$ 、 $C$ 所对的边分别是 $a$ 、 $b$ 、 $c$ . (1)求角 $A$ 的大小；

(2)若 $a = \sqrt{3}$ ，求 $\bigtriangleup  {ABC}$ 面积的最大值.

【例 4】如图,已知点 $D$ 在 $\bigtriangleup {ABC}$ 的边 ${BC}$ 上,且 $\angle {DAC} = {90}^{ \circ  },\cos \angle {DAB} = \frac{2\sqrt{2}}{3},{AB} = 6$ .

(1)若 $\sin C = \frac{\sqrt{3}}{3}$ ，求线段 ${BC}$ 的长；

(2)若点 $E$ 是 ${BC}$ 的中点， ${AE} = \sqrt{17}$ ，求线段 ${AC}$ 的长.

\begin{center}
\includegraphics[width=0.3\textwidth]{images/bo_d7tk7oc91nqc73875fu0_104_1076_1215_363_180_0.jpg}
\end{center}
 

(例 4 图)

【例 5】已知 $\bigtriangleup {ABC}$ 的三个内角 $A\text{ 、 }B\text{ 、 }C$ 所对的边分别是 $a\text{ 、 }b\text{ 、 }c$ ,且满足 $3\overrightarrow{AB} \cdot  \overrightarrow{AC} + 4\overrightarrow{CA} \cdot  \overrightarrow{CB} =  \overrightarrow{BA} \cdot  \overrightarrow{BC}$ .

(1)求 $\frac{c}{b}$ 的值；

(2)若 $\overrightarrow{AD} = \frac{2}{3}\overrightarrow{AB} + \frac{1}{3}\overrightarrow{AC}$ ， $\left| \overrightarrow{AD}\right|  = 2$ ，求 $\bigtriangleup {ABC}$ 面积的最大值.

\section*{四基达标精练}

\section*{一、填空题}

1. 在 $\bigtriangleup {ABC}$ 中， ${a}^{2} + {bc} = {b}^{2} + {c}^{2}$ ，则 $A =$ \_\_\_\_\_.

2. 已知 $\bigtriangleup  {ABC}$ 的内角 $A$ 、 $B$ 、 $C$ 所对的边分别是 $a$ 、 $b$ 、 $c$ ，若 $A = {30}^{ \circ  }$ ， $B = {105}^{ \circ  }$ ， $a = 4$ ，则 $c =$ \_\_\_\_\_.

3. 在 $\bigtriangleup  {ABC}$ 中， $a = 2\sqrt{6}$ ， $b = {2c}$ ， $\cos A =  - \frac{1}{4}$ ，则 $\bigtriangleup  {ABC}$ 的面积 ${S}_{\bigtriangleup {ABC}} =$ \_\_\_\_\_.

4. 边长分别为 ${10}\text{ 、 }{14}\text{ 、 }{16}$ 的三角形中,最大角与最小角的和为\_\_\_\_\_.

5. 在 $\bigtriangleup  {ABC}$ 中， ${BC} = 6$ ， ${AC} = 8$ ， $\angle A = {40}^{ \circ  }$ ，则 $\angle B$ 的解的个数是\_\_\_\_\_.

6. 在 $\bigtriangleup {ABC}$ 中,角 $A\text{ 、 }B\text{ 、 }C$ 所对的边分别为 $a\text{ 、 }b\text{ 、 }c,\bigtriangleup {ABC}$ 的面积 $S\bigtriangleup {ABC} = \frac{\sqrt{3}}{4}\left( {{a}^{2} + {b}^{2} - {c}^{2}}\right)$ ,则 $\frac{c}{a + b}$ 的取值范围是\_\_\_\_\_.

7. 在 $\bigtriangleup  {ABC}$ 中，若 $\left( {\sin A + \sin B}\right) \left( {\sin A - \sin B}\right)  = {\sin }^{2}C$ ，则 $A =$ \_\_\_\_\_.

8. 在 $\bigtriangleup  {ABC}$ 中，角 $A$ 、 $B$ 、 $C$ 所对的边分别是 $a$ 、 $b$ 、 $c$ ，若 $\frac{a}{\sin B} = \frac{b}{\sin C} = \frac{c}{\sin A}$ ，则 $\bigtriangleup  {ABC}$ 的形状是

9. 已知 $\bigtriangleup {ABC}$ 的三个内角 $A$ 、 $B$ 、 $C$ 成等差数列，且 ${AB} = 1$ ， ${BC} = 4$ ，则边 ${AC}$ 上的中线 ${BD}$ 的长为\_\_\_\_\_.

\section*{二、选择题}

10. 在 $\bigtriangleup  {ABC}$ ，角 $A$ 、 $B$ 、 $C$ 所对的边分别为 $a$ 、 $b$ 、 $c$ ，若 $A : B : C = 3 : 2 : 1$ ，则 $a$ : $b : c$ 等于\_\_\_\_\_。

A. $1 : 2 : 3$ . B. $3 : 2 : 1$ . C. $1 : \sqrt{3} : 2$ . D. $2 : \sqrt{3} : 1$ .

11. 在 $\bigtriangleup  {ABC}$ 中，角 $A$ 、 $B$ 、 $C$ 所对的边分别为 $a$ 、 $b$ 、 $c$ ，若 $c = 4$ ， $a = 4\sqrt{2}$ ， $A = \frac{\pi }{4}$ ，则角 $C$ 的大小为( )

A. $\frac{\pi }{6}$ . B. $\frac{\pi }{6}$ 或 $\frac{5\pi }{6}$ . C. $\frac{\pi }{3}$ . D. $\frac{\pi }{3}$ 或 $\frac{2\pi }{3}$ .

12. 在 $\bigtriangleup  {ABC}$ 中，角 $A$ 、 $B$ 、 $C$ 所对应的边分别为 $a$ 、 $b$ 、 $c$ ，设 $\bigtriangleup  {ABC}$ 的面积为 $S$ ，则 $\frac{S}{{a}^{2} + {4bc}}$ 的最大值为 ()

A. $\frac{\sqrt{2}}{16}$ . B. $\frac{\sqrt{3}}{12}$ . C. $\frac{\sqrt{3}}{16}$ . D. $\frac{\sqrt{2}}{18}$ .

\section*{三、解答题}

13. 在 $\bigtriangleup {ABC}$ 中,已知 $a = 2,b = \sqrt{2},c = \sqrt{3} + 1$ .

(1)求角 $A$ 、 $B$ 、 $C$ 的大小；

(2)求 $\bigtriangleup  {ABC}$ 面积.

14. 在 $\bigtriangleup  {ABC}$ 中，内角 $A$ 、 $B$ 、 $C$ 所对的边分别为 $a$ 、 $b$ 、 $c$ ，已知 $a\sin C = \sqrt{3}c\cos A$ .

(1)求角 $A$ 的大小；

(2)若 $b = 2$ ，且 $\frac{\pi }{4} \leq  B \leq  \frac{\pi }{3}$ ，求边 $c$ 的取值范围.

15. 已知 $\bigtriangleup  {ABC}$ 的内角 $A$ 、 $B$ 、 $C$ 的对边分别记为 $a$ 、 $b$ 、 $c$ ，且 ${2c}\cos A = b - c$ .

(1)证明: $A = {2C}$ ；

(2)若 $\bigtriangleup  {ABC}$ 为锐角三角形，求 $\frac{a + b}{c}$ 的取值范围.

16. 海岸的 $A$ 哨所在凌晨 1 点 15 分发现哨所北偏东 ${30}^{ \circ  }$ 方向 20 海里处的点 $D$ 出现可疑船只，因天气恶劣能见度低,无法对船只进行识别,所以将该船雷达特征信号进行标记并上报周围哨所. 早上 5 点 15 分位于 $A$ 哨所正西方向 20 海里的 $B$ 哨所发现了该可疑船只位于 $B$ 哨所北偏西 30 \textdegree 方向 60 海里处的点 $E$ ，并识别出其为走私船,立刻命令位于 $B$ 哨所正西方向 30 海里处点 $C$ 的我方缉私船前往拦截，已知缉私船速度大小为 30 海里/时. (假设所有船只均保持匀速直线航行)

(1)求走私船的速度大小；

(2)缉私船沿什么方向行驶才能最快截获走私船，并求出截获走私船的具体时间.

\begin{center}
\includegraphics[width=0.2\textwidth]{images/bo_d7tk7oc91nqc73875fu0_107_1137_1144_322_261_0.jpg}
\end{center}
 

(第 16 题图)

\section*{§ 6.3 三角函数及其性质一}

\section*{典型例题精讲}

【例 1】已知函数 $f\left( x\right)  = \sqrt{3}\sin x\cos x - {\cos }^{2}x$ .

(1)求函数 $f\left( x\right)$ 的最小正周期和单调区间；

(2)求函数 $f\left( x\right)$ 的零点.

【例 2】已知 $f\left( x\right)  = \cos {\omega x}\left( {\omega  > 0}\right)$ .

(1) $f\left( x\right)$ 的周期是 $\pi$ ,求当 $x \in  \left\lbrack  {0,{2\pi }}\right\rbrack$ ，方程 $f\left( {x + \frac{\pi }{6}}\right)  = \frac{\sqrt{3}}{2}$ 的解集；

(2)已知 $\omega  = 1$ ， $g\left( x\right)  = {f}^{2}\left( x\right)  + \sqrt{3}f\left( {-x}\right)$ ， $f\left( {\frac{\pi }{2} - x}\right)$ ， $x \in  \left\lbrack  {0,\frac{\pi }{4}}\right\rbrack$ ，求 $g\left( x\right)$ 的值域.

【例 3】已知函数 $f\left( x\right)  = 2\cos \left( {{\omega x} + \varphi }\right) \left( {\omega  \in  \mathbf{N},\left| \varphi \right|  < \frac{\pi }{2}}\right)$ 在区间 $\left( {0,\frac{5\pi }{18}}\right)$ 上是单调函数，且 $f\left( \frac{\pi }{18}\right)  +  f\left( \frac{\pi }{6}\right)  = 0.$

(1)求 $f\left( x\right)$ 图像的一个对称中心；

(2)若 $f\left( \frac{\pi }{6}\right)  =  - 1$ ，求 $f\left( x\right)$ 的解析式.

【例 4】函数 $f\left( x\right)  = \cos \left( {{\omega x} + \varphi }\right) \left( {\omega  > 0,\left| \varphi \right|  < \frac{\pi }{2}}\right)$ 的部分图像如图所示.

(1)求 $f\left( x\right)$ 的解析式;

(2)若对任意 $x \in  \left\lbrack  {-\frac{\pi }{4},\frac{\pi }{4}}\right\rbrack$ ，都有 ${\left\lbrack  f\left( x\right) \right\rbrack  }^{2} - {mf}\left( x\right)  - 1 \leq  0$ ，求实数 $m$ 的取值范围.

\begin{center}
\includegraphics[width=0.2\textwidth]{images/bo_d7tk7oc91nqc73875fu0_109_1097_1328_322_208_0.jpg}
\end{center}
 

(例 4 图)

【例 5】已知 $f\left( x\right)  = 2{x}^{2} - {3x} + 1,g\left( x\right)  = k \cdot  \sin \left( {x - \frac{\pi }{6}}\right) \left( {k \neq  0}\right)$ .

(1)设 $f\left( x\right)$ 的定义域为 $\left\lbrack  {0,3}\right\rbrack$ ，值域为 $A$ ； $g\left( x\right)$ 的定义域为 $\left\lbrack  {0,3}\right\rbrack$ ，值域为 $B$ ，且 $A \subseteq  B$ ，求实数 $k$ 的取值范围.

(2)若方程 $f\left( {\sin x}\right)  + \sin x - a = 0$ 在 $\lbrack 0,{2\pi })$ 上恰有两个解，求实数 $a$ 的取值范围.

\section*{四基达标精练}

\section*{一、填空题}

1. $f\left( x\right)  = \left| {\sin x}\right|$ 的值域是\_\_\_\_\_.

2. 函数 $y = 2\cos x$ 的严格减区间为\_\_\_\_\_.

3. 函数 $y = \tan \left( {\frac{\pi }{6}x + \frac{\pi }{3}}\right)$ 的定义域是\_\_\_\_\_.

4. 方程 $\cos {2x} - \sin x = 0$ 在区间 $\left\lbrack  {0,{2\pi }}\right\rbrack$ 上的所有解的和为\_\_\_\_\_.

5. 若函数 $f\left( x\right)  = \cos {\omega x}\left( {\omega  > 0}\right)$ 相邻两条对称轴之间的距离为 2，则 $f\left( 2\right)  =$ \_\_\_\_\_.

6. 若函数 $y = \sin \left( {{\pi x} - \frac{\pi }{6}}\right)$ 在 $\left\lbrack  {0,m}\right\rbrack$ 上单调增,则 $m$ 的最大值为\_\_\_\_\_.

7. 若方程 $\cos x = \frac{a - 1}{2}$ 在 $x \in  \left\lbrack  {-\frac{\pi }{3},\pi }\right\rbrack$ 上有两个不同的实数根，则实数 $a$ 的取值范围为\_\_\_\_\_.

8. 已知函数 $f\left( x\right)  = a\sin x + {bx} + 1$ ，若 $f\left( {-1}\right)  = 2$ ，则 $f\left( 1\right)  =$ \_\_\_\_\_.

9. 已知 $f\left( x\right)  = {x}^{2} - \cos x$ ,对 $\left\lbrack  {-\frac{\pi }{2},\frac{\pi }{2}}\right\rbrack$ 上的任意 ${x}_{1}\text{ 、 }{x}_{2}$ ,有如下条件: $\text{  \ding{172}  }{x}_{1} > {x}_{2}$ ;  \ding{173}  ${x}_{1}^{2} > {x}_{2}^{2}$ ； \ding{174}  $\left| {x}_{1}\right|  > {x}_{2}$ . 其中能使 $f\left( {x}_{1}\right)  > f\left( {x}_{2}\right)$ 恒成立的条件序号是\_\_\_\_\_.

\section*{二、选择题}

10. 下列四个函数中,同时具有性质:  \ding{172}  最小正周期为 ${2\pi }$ ;  \ding{173}  图像关于直线 $x = \frac{\pi }{3}$ 对称的一个函数是 ( )

A. $y = \sin \left( {x - \frac{\pi }{6}}\right)$ . B. $y = \sin \left( {x + \frac{\pi }{6}}\right)$ . C. $y = \sin \left( {x + \frac{\pi }{3}}\right)$ . D. $y = \sin \left( {{2x} - \frac{\pi }{3}}\right)$ .

11. 函数 $y = x + \sin \left| x\right| \left( {x \in  \left\lbrack  {-\pi ,\pi }\right\rbrack  }\right)$ 的大致图像是 ( )

\begin{center}
\includegraphics[width=0.2\textwidth]{images/bo_d7tk7oc91nqc73875fu0_111_161_206_256_210_0.jpg}
\end{center}
 

A.

\begin{center}
\includegraphics[width=0.4\textwidth]{images/bo_d7tk7oc91nqc73875fu0_111_463_205_576_243_0.jpg}
\end{center}
 

(第 11 题图)

\begin{center}
\includegraphics[width=0.2\textwidth]{images/bo_d7tk7oc91nqc73875fu0_111_1117_199_254_210_0.jpg}
\end{center}
 

D.

12. 函数 $f\left( x\right)  = \cos \left( {{\omega x} + \varphi }\right) \left( {0 < \varphi  < \pi }\right)$ 周期为 $T = \frac{2\pi }{3}$ ，其图像关于 $\left( {\frac{\pi }{18},0}\right)$ 对称，且当 $x \in  \left\lbrack  {\frac{\pi }{6},m}\right\rbrack$ 时， $f\left( x\right)$ 的值域是 $\left\lbrack  {-1, - \frac{\sqrt{3}}{2}}\right\rbrack$ ,则 $m$ 的取值范围是 ( )

A. $\left\lbrack  {\frac{\pi }{9},\frac{7\pi }{18}}\right\rbrack$ . B. $\left\lbrack  {\frac{2\pi }{9},\frac{7\pi }{18}}\right\rbrack$ . C. $\left\lbrack  {\frac{\pi }{9},\frac{5\pi }{18}}\right\rbrack$ . D. $\left\lbrack  {\frac{2\pi }{9},\frac{5\pi }{18}}\right\rbrack$ .

\section*{三、解答题}

13. 已知函数 $f\left( x\right)  = 2\sin \left( {{2x} - \frac{\pi }{6}}\right)$ .

(1)求 $f\left( x\right)$ 的最小正周期及单调递减区间；

(2)求 $f\left( x\right)$ 在区间 $\left\lbrack  {\frac{\pi }{4},\frac{\pi }{2}}\right\rbrack$ 上的最大值与最小值.

14. 已知函数 $f\left( x\right)  = 2\sin \left( {{2\omega x} - \frac{\pi }{6}}\right) ,{x}_{1}\text{ 、 }{x}_{2}$ 是函数 $y = f\left( x\right)$ 与直线 $y = 2$ 的两个交点的横坐标，且 $\left| {{x}_{1} - {x}_{2}}\right|$ 的最小值为 $\pi$ .

(1)求 $f\left( x\right)$ 的解析式;

(2)若 $x \in  \left\lbrack  {\frac{\pi }{2},m}\right\rbrack$ ， $f\left( x\right)$ 的值域是 $\left\lbrack  {-2,1}\right\rbrack$ ，求 $m$ 的取值范围.

15. 已知函数 $f\left( x\right)  = A\sin \left( {{\omega x} + \varphi }\right) ,\left( {\omega  > 0,\left| \varphi \right|  < \frac{\pi }{2}}\right)$ 的部分图像如图所示.

(1)求 $f\left( x\right)$ 的解析式；

(2)若 $x \in  \left( {-\frac{\pi }{18},\frac{5\pi }{18}}\right)$ ，求 $f\left( x\right)$ 的取值范围.

\begin{center}
\includegraphics[width=0.2\textwidth]{images/bo_d7tk7oc91nqc73875fu0_112_1116_362_320_205_0.jpg}
\end{center}
 

(第 15 题图)

16. 如图，已知 ${OPQ}$ 是半径为 $\sqrt{7}$ ，圆心角为 $\frac{\pi }{3}$ 的扇形， $C$ 是圆弧 $\overset{\text{ ⏜ }}{PQ}$ 上的一点， ${ABCD}$ 是扇形的内接矩形,记 $\angle {BOC} = \alpha$ .

(1)若 $\alpha  = \frac{\pi }{6}$ ，求矩形 ${ABCD}$ 的周长和面积；

(2)若 ${{Rt}} \bigtriangleup  {CBO}$ 的周长为 $\frac{\sqrt{7}\left( {2\sqrt{10} + 5}\right) }{5}$ ，求 $\tan \alpha$ 的值.

\begin{center}
\includegraphics[width=0.2\textwidth]{images/bo_d7tk7oc91nqc73875fu0_112_1187_1213_253_253_0.jpg}
\end{center}
 

(第 16 题图)

\section*{§ 6.4 三角函数及其性质二}

\section*{典型例题精讲}

【例 1】把正弦函数 $f\left( x\right)  = \cos x$ 的图像向左平移 $\frac{\pi }{6}$ 个单位后,再将得到的图像上各点横坐标伸长到原来的 2 倍,纵坐标不变,得到函数 $y = g\left( x\right)$ 的图像.

(1)求函数 $g\left( x\right)$ 单调递减区间；

(2)若 $x \in  \lbrack 0,{4\pi })$ ，且 $g\left( {x}_{0}\right)  =  - \frac{\sqrt{3}}{2}$ ，求 ${x}_{0}$ 的值.

【例 2】已知函数 $f\left( x\right)  = \sin \left( {{\omega x} + \varphi }\right) \left( {\omega  > 0,\left| \varphi \right| \frac{\pi }{2}}\right)$ 的图像关于 $\left( {-\frac{\pi }{6},0}\right)$ 中心对称，且图像上相邻两个对称轴的距离为 $\frac{\pi }{2}$ .

(1)求函数 $f\left( x\right)$ 的解析式；

(2)设 ${x}_{1}$ 、 ${x}_{2} \in  \left( {-\frac{\pi }{6},\frac{\pi }{3}}\right)$ ，且 ${x}_{1} \neq  {x}_{2}$ ，若 $f\left( {x}_{1}\right)  = f\left( {x}_{2}\right)$ ，求 $f\left( {{x}_{1} + {x}_{2}}\right)$ 的值.

【例 3】若点 $\left( {{x}_{0},{y}_{0}}\right)$ 在函数 $f\left( x\right)$ 的图像上,且满足 ${y}_{0} \cdot  f\left( {y}_{0}\right)  \geq  0$ ,则称 ${x}_{0}$ 是 $f\left( x\right)$ 的 $\xi$ 点. 函数 $f\left( x\right)$ 的所有 $\xi$ 点构成的集合称为 $f\left( x\right)$ 的 $\xi$ 集.

(1)判断 $\frac{2\pi }{3}$ 是否是函数 $f\left( x\right)  = \tan x$ 的 $\xi$ 点,并说明理由;

(2)若函数 $f\left( x\right)  = \sin \left( {\pi x}\right)$ ，求 $f\left( x\right)$ 的 $\xi$ 集.

【例 4】某旅游区每年各个月份接待游客的人数近似地满足周期性规律,因而第 $n$ 个月从事旅游服务工作的人数 $f\left( n\right)$ 可近似地用函数 $f\left( n\right)  = A\cos \left( {{\omega n} + \theta }\right)  + k$ 来刻画,其中正整数 $n$ 表示月份且 $n \in  \left\lbrack  {1,{12}}\right\rbrack$ , 例如 $n = 1$ 表示 1 月份, $A$ 和 $k$ 是正整数, $\omega  > 0,\theta  \in  \left( {0,\pi }\right)$ . 统计发现,该地区每年各个月份从事旅游服务工作的人数有以下规律:

 \ding{172} 每年相同的月份，该地区从事旅游服务工作的人数基本相同；

 \ding{173} 该地区从事旅游服务工作的人数最多的 8 月份和最少的 2 月份相差 400 人；

 \ding{174} 2 月份该地区从事旅游服务工作的人数为 100 人，随后逐月递增直到 8 月份达到最多.

(1)试根据已知信息，求 $f\left( n\right)$ 的表达式；

(2)一般地，当该地区从事旅游服务工作的人数在 400 或 400 以上时，该地区也进入了一年中的旅游“旺季”，那么，一年中的哪几个月是该地区的旅游“旺季”？请说明理由.

【例 5】如图所示是某斜拉式大桥图片, 为了解桥的一些结构情况, 学校数学兴趣小组将大桥的结构进行了简化，取其部分可抽象成右图所示的模型，其中桥塔 ${AB}$ 、 ${CD}$ 与桥面 ${AC}$ 垂直，通过测量得知 ${AB} =  {50}{m},{AC} = {50}{m}$ ,当 $P$ 为 ${AC}$ 中点时, $\angle {PBD} = {45}^{ \circ  }$ .

(1)求 ${CD}$ 的长；

(2)设 ${AP} = x$ ，写出 $\tan \angle {PBD}$ 与 $x$ 的函数关系式；

(3)已知命题:函数 $y = \tan x$ 在 $x \in  \left( {0,\frac{\pi }{2}}\right)$ 内为严格增函数；求证该命题为真命题，并用该命题求解 $P$ 在线段 ${AC}$ 的何处时， $\angle {PBD}$ 达到最大，最大值为多少？

\begin{center}
\includegraphics[width=0.6\textwidth]{images/bo_d7tk7oc91nqc73875fu0_115_691_520_765_271_0.jpg}
\end{center}
 

(例 5 图)

\section*{四基达标精练}

\section*{一、填空题}

1. 将曲线 $y = \sin x$ 上所有点向左平移 $\varphi \left( {\varphi  > 0}\right)$ 个单位，得到函数 $y =  - \sin x$ 的图像，则 $\varphi$ 的最小值为\_\_\_\_\_.

2. 若函数 $f\left( x\right)  = \cos \left( {{2x} + \varphi }\right) \left( {-\pi  \leq  \varphi  < \pi }\right)$ 为奇函数，则 $\varphi  =$ \_\_\_\_\_.

3. 函数 $y = \tan x + x$ 在坐标原点处的切线的斜率为\_\_\_\_\_.

4. 已知 $\cos x = \frac{3}{5}$ ，则 $x$ 取值集合为\_\_\_\_\_. (答案用反余弦表示)

5. 已知函数 $f\left( x\right)  = \sqrt{3}\sin x - \cos x$ ，若将 $f\left( x\right)$ 的图像向左平行移动 $\frac{\pi }{6}$ 个单位长度后得到 $g\left( x\right)$ 的图像， 则 $g\left( x\right)$ 的对称中心为\_\_\_\_\_.

6. 如果函数 $y = \frac{1}{2}\sin {\omega x}$ 在区间 $\left\lbrack  {-\frac{\pi }{8},\frac{\pi }{12}}\right\rbrack$ 上单调递减，那么 $\omega$ 的取值范围是\_\_\_\_\_.

7. 函数 $f\left( x\right)  = 3\cos \left( {{2\omega x} + \frac{\pi }{3}}\right) \left( {\omega  > 0}\right)$ 在 $\left\lbrack  {-\frac{\pi }{3},\frac{5\pi }{6}}\right\rbrack$ 上是减函数,且在 $\left\lbrack  {0,{2\pi }}\right\rbrack$ 上恰好取得一次最小值 -3，则 $\omega$ 的取值范围是\_\_\_\_\_.

8. 方程 $\sin {2x} = \frac{1}{2}\tan x$ 在区间 $\left\lbrack  {0,\pi }\right\rbrack$ 上的所有解的和为\_\_\_\_\_.

9. 直线 $y = a$ 与函数 $f\left( x\right)  = \sin \left( {{\omega x} + \frac{\pi }{3}}\right)  + 1\left( {\omega  > 0}\right)$ 的图像相交,若自左至右的三个相邻交点依次为 $A\text{ 、 }B\text{ 、 }C$ ，且满足 $2\left| {AB}\right|  = \left| {BC}\right|$ ，则实数 $a =$ \_\_\_\_\_.

\section*{二、选择题}

10. 把函数 $y = f\left( x\right)$ 图像上所有点的横坐标缩短到原来的 $\frac{1}{3}$ ，纵坐标不变，再把所得曲线向左平移 $\frac{\pi }{4}$ 个单位长度,得到函数 $y = \sin \left( {x + \frac{\pi }{3}}\right)$ 的图像,则 ( )

A. $f\left( x\right)  = \sin \left( {\frac{x}{3} + \frac{\pi }{12}}\right)$ . B. $f\left( x\right)  = \sin \left( {\frac{x}{3} + \frac{7\pi }{12}}\right)$ . C. $f\left( x\right)  = \sin \left( {{3x} + \frac{\pi }{12}}\right)$ . D. $f\left( x\right)  = \sin \left( {{3x} + \frac{7\pi }{12}}\right)$ .

11. 关于函数 $f\left( x\right)  = \frac{1 + \tan x}{1 - \tan x}$ ,下列说法错误的是 ( )

A. 函数的最小正周期为 $\pi$ . B. 函数图像关于点 $\left( {\frac{\pi }{4},0}\right)$ 对称.

C. 函数图像关于直线 $x = \frac{\pi }{2}$ 对称. D. 函数的定义域为 $\{ x \mid  x \neq  \frac{\pi }{4} + {k\pi },k \in  \mathbf{Z}\}$ .

12. 已知函数 $f\left( x\right)  = \sin \left( {{\omega x} - \frac{\pi }{3}}\right) \left( {\omega  > 0}\right)$ 在 $\left( {\frac{\pi }{2},\frac{3\pi }{2}}\right)$ 上无零点,则 $\omega$ 的取值范围是 ( )

A. $\left( {0,\frac{2}{9}}\right\rbrack   \cup  \left\lbrack  {\frac{8}{9}, + \infty }\right)$ . B. $\left( {0,\frac{2}{9}}\right\rbrack   \cup  \left\lbrack  {\frac{2}{3},\frac{8}{9}}\right\rbrack$ . C. $\left( {0,\frac{2}{9}}\right\rbrack   \cup  \left\lbrack  {\frac{8}{9},1}\right\rbrack$ . D. $\left( {\frac{2}{9},\frac{8}{9}}\right\rbrack   \cup  \lbrack 1, + \infty )$ .

\section*{三、解答题}

13. 已知函数 $f\left( x\right)  = 2\sqrt{3}\sin x\cos x + 2{\cos }^{2}x$ .

(1)求函数 $f\left( x\right)$ 的最小正周期；

(2)现将函数 $f\left( x\right)$ 图像上所有的点的横坐标伸长到原来的 2 倍 (纵坐标不变)，得到函数 $g\left( x\right)$ 的图像， 求 $g\left( x\right)$ 在区间 $\left\lbrack  {0,\frac{\pi }{2}}\right\rbrack$ 上的值域.

14. 已知函数 $f\left( x\right)  = \sin x\cos x - \sqrt{3}{\sin }^{2}x + \frac{\sqrt{3}}{2}$ .

(1)若存在 $x \in  \left\lbrack  {-\frac{\pi }{3},\frac{\pi }{6}}\right\rbrack$ ，使得 $f\left( x\right)  \geq  a$ 成立，则求 $a$ 的取值范围；

(2)将函数 $f\left( x\right)$ 的图像上每个点纵坐标不变，横坐标缩短到原来的 $\frac{1}{2}$ ，得到函数 $g\left( x\right)$ 的图.

15. 某用电器电流 $I\left( {{mA}}\right)$ 随时间 $t\left( {s}\right)$ 的变化的部分图像如图所示， $I$ 和 $t$ 满足下列关系式:

$I\left( t\right)  = A\sin \left( {{\omega t} + \varphi }\right) \left( {A > 0,\omega  > 0,\left| \varphi \right|  < \frac{\pi }{2}}\right) .$

(1)求 $I\left( t\right)  = A\sin \left( {{\omega t} + \varphi }\right)$ 的解析式；

(2)若该用电器核心部件有效工作的电流 $\left| I\right|$ 必须大于 ${150}{{mA}}$ ，则在一个周期内，该用电器核心部件的有效工作时间是多少? (电流的正负表示电流的正反方向)

\begin{center}
\includegraphics[width=0.2\textwidth]{images/bo_d7tk7oc91nqc73875fu0_117_1142_1010_317_233_0.jpg}
\end{center}
 

(第 15 题图)

16. 已知函数 $f\left( x\right)  = 2\sqrt{3}\sin \left( {\frac{\pi }{4} + \frac{x}{2}}\right)  \cdot  \cos \left( {\frac{\pi }{4} + \frac{x}{2}}\right)  - \sin \left( {\pi  + x}\right)$ ,且函数 $y = g\left( x\right)$ 的图像与函数 $y =  f\left( x\right)$ 的图像关于直线 $x = \frac{\pi }{4}$ 对称.

(1)求函数 $g\left( x\right)$ 的解析式；

(2)若存在 $x \in  \left\lbrack  {0,\frac{\pi }{2}}\right)$ ，使等式 ${\left\lbrack  g\left( x\right) \right\rbrack  }^{2} - {mg}\left( x\right)  + 2 = 0$ 成立，求实数 $m$ 的最大值和最小值.

\section*{第 7 章 数列(选择性必修 1)}

\section*{知识梳理}

\section*{一、数列的有关概念}

1.数列的概念:数列是有序排列的一列数. 数列是定义在自然数集或其有限子集上的函数,当自变量由小到大依次取值时对应的一列函数值.

 \ding{172} 数列是特殊的函数:数列的学习是函数学习的延续; 可以借助函数的知识和方法研究数列;

 \ding{173} 数列是离散型函数:数列的图像是一系列孤立的点；

 \ding{174} 数列与函数比较更具有特殊性，数列的学习更多关注的是其特殊性.

2. 数列的单调性:

数列 $\left\{  {a}_{n}\right\}$ 是递增数列: 对任意 $n \in  {\mathbf{N}}^{ * },{a}_{n + 1} > {a}_{n}$ 成立;

数列 $\left\{  {a}_{n}\right\}$ 是递减数列: 对任意 $n \in  {\mathbf{N}}^{ * },{a}_{n + 1} < {a}_{n}$ 成立.

 \ding{172} 数列的单调性定义是判断数列单调性的基本依据.

 \ding{173} 数列的单调性与函数的单调性不同.

一般地，若函数 $f\left( x\right) \left( {x > 0}\right)$ 单调增 (减)，则数列 $\left\{  {f\left( n\right) }\right\}$ 单调增 (减)；

但,若数列 $\left\{  {f\left( n\right) }\right\}$ 单调增 (减),不一定函数 $f\left( x\right) \left( {x > 0}\right)$ 单调增 (减).

 \ding{174} 数列的单调性与数列的项的最值密切相关.

3. 求数列最大项或最小项的方法:

用不等式 ${a}_{n} \leq  {a}_{n + 1}$ (或 ${a}_{n} \geq  {a}_{n + 1}$ ) 寻找使得 $\left\{  {a}_{n}\right\}$ 的单调递增和单调递减的 $n$ 的范围确定数列的最大项或最小项.

注意: \ding{172}  通常当 ${a}_{n}$ 由等差数列和等比数列相应项的和(或差)构成时利用 ${a}_{n + 1} - {a}_{n} \geq  0$ ；

 \ding{173}  当 ${a}_{n}$ 由等差数列和等比数列相应项的积(或商)构成时利用 $\frac{{a}_{n + 1}}{{a}_{n}} \geq  1$ (注意各项为正)；

 \ding{174}  当数列 $\left\{  {a}_{n}\right\}$ 的各项中有正有负时,常先考虑其各项的绝对值的最大项,再还原.

\section*{二、等差数列与等比数列}

\begin{center}
\begin{tabular}{|c|c|c|}
\hline
\phantom{X} & 等差数列 & 等比数列 \\
\cline{1-3}
定义 & $\left\{  \begin{array}{l} {a}_{1}, \\  {a}_{n + 1} - {a}_{n} = d \end{array}\right.$ (常数) & $\left\{  \begin{array}{l} {a}_{1} \neq  0, \\  \frac{{a}_{n + 1}}{{a}_{n}} = q\text{ (非零常数) } \end{array}\right.$ \\
\cline{1-3}
通项公式 & ${a}_{n} = {a}_{1} + \left( {n - 1}\right) d = {a}_{k} + \left( {n - k}\right) d$ & ${a}_{n} = {a}_{1}{q}^{n - 1} = {a}_{k}{q}^{n - k}$ \\
\cline{1-3}
中项 & $A = \frac{a + b}{2}\left( {a,b \in  \mathbf{R},A\text{ 唯一 }}\right)$ & ${G}^{2} = {ab}\left( {a\text{ 、 }b \in  \mathbf{R},{ab} > 0,G\text{ 不唯一) }}\right)$ \\
\cline{1-3}
前 $n$ 项和 & ${S}_{n} = \frac{n\left( {{a}_{1} + {a}_{n}}\right) }{2} = \frac{n\left( {{a}_{k} + {a}_{n - k + 1}}\right) }{2}  = {a}_{1}n + \frac{n\left( {n - 1}\right) d}{2} = {a}_{n}n - \frac{n\left( {n - 1}\right) d}{2}$ & ${S}_{n} = \left\{  \begin{array}{l} n{a}_{1},q = 1, \\  \frac{{a}_{1}\left( {1 - {q}^{n}}\right) }{1 - q} = \frac{{a}_{1} - {a}_{n}q}{1 - q},q \neq  1 \end{array}\right.$ \\
\cline{1-3}
\phantom{X} & \phantom{X} & 各项和: $\mathop{\lim }\limits_{{n \rightarrow   + \infty }}{S}_{n} = \mathop{\sum }\limits_{{i = 1}}^{{+\infty }}{a}_{i} = \frac{{a}_{1}}{1 - q}\left( {0 < \left| q\right|  < 1}\right)$ \\
\cline{1-3}
\hline
\end{tabular}
\end{center}

\HRule

* 注:本章节所属知识模块是 “函数与分析”.

\HRule

1. 若 $m + n = k + l\left( {m\text{ 、 }n\text{ 、 }k\text{ 、 }l \in  {\mathbf{N}}^{ * }}\right)$ ,则在等差数列中: ${a}_{m} + {a}_{n} = {a}_{k} + {a}_{l}$ ; 在等比数列中: ${a}_{m}{a}_{n} = {a}_{k}{a}_{l}$ . 注意:以上两个性质可推广到若干个，但等式两端的项数必须相等，脚标之和必须相等.

2. $A = \frac{a + b}{2}$ 是 $a,A,b$ 成等差数列的充要条件.

3. ${G}^{2} = {ab}$ 是 $a,G,b$ 成等比数列的必要非充分条件.

4. 用等比数列求和公式求数列的和时,勿忘 $q \neq  1$ .

5. 分析等差数列的前 $n$ 项和,有 ${S}_{{2k} - 1} = \left( {{2k} - 1}\right) {a}_{k}$ .

6. 等差数列 $\left\{  {a}_{n}\right\}$ 的前 $n$ 项和,次 $n$ 项和,再 $n$ 项和, $\cdots$ ,第 $k$ 个 $n$ 项和成等差数列,公差为 ${n}^{2}d$ .

7. 等比数列 $\left\{  {a}_{n}\right\}$ 的前 $n$ 项和,次 $n$ 项和,再 $n$ 项和, $\cdots$ ,第 $k$ 个 $n$ 项和成等比数列,公比 ${q}^{n}$ .

8. 已知等差数列的前 $n$ 项和为 ${S}_{n}$ ,则 $\frac{{S}_{n}}{n} = {dn} + \left( {{a}_{1} - \frac{d}{2}}\right)$ ,表明 $\left\{  \frac{{S}_{n}}{n}\right\}$ 也是等差数列; 即点 $\left( {n,\frac{{S}_{n}}{n}}\right) \left( {n \in  {\mathbf{N}}^{ * }}\right)$ 共线 (也就是当 $n$ 取不同值时,任意两点连线的斜率相等).

9. 项数有限的等差数列奇数项的和 ${S}_{\text{ 奇 }}$ ,偶数项的和 ${S}_{\text{ 偶 }}$ 满足 ${S}_{\text{ 奇 }} + {S}_{\text{ 偶 }} = {S}_{n}$ ; 当 $n$ 为偶数时, ${S}_{\text{ 奇 }} - {S}_{\text{ 偶 }} =  - \frac{nd}{2}$ ; 当 $n$ 为奇数时, ${S}_{\text{ 奇 }} - {S}_{\text{ 偶 }} = \frac{{S}_{n}}{n} = {a}_{\frac{n + 1}{2}}$ .

10. 运算性质:

 \ding{172}  若数列 $\left\{  {a}_{n}\right\}$ 、 $\left\{  {b}_{n}\right\}$ 都是等差数列，则 $\left\{  {A \cdot  {a}_{n} - B \cdot  {b}_{n}}\right\}$ ( $A$ 、 $B$ 为常数)也是等差数列；

 \ding{173} 若数列 $\left\{  {a}_{n}\right\}$ 、 $\left\{  {b}_{n}\right\}$ 都是等比数列，则 $\left\{  {{a}_{n}^{A} \cdot  {b}_{n}^{B}}\right\}$ ( $A$ 、 $B$ 为常数) 也是等比数列；

 \ding{174} 若数列 $\left\{  {a}_{n}\right\}$ 是等差数列， ${b}_{n} = {a}^{{a}_{n}}\left( {a > 0,a \neq  1}\right)$ ，则 $\left\{  {b}_{n}\right\}$ 是等比数列；

 \ding{175} 若数列 $\left\{  {a}_{n}\right\}  \left( {{a}_{n} > 0}\right)$ 是等比数列， ${b}_{n} = {\log }_{a}{a}_{n}\left( {a > 0,a \neq  1}\right)$ ，则 $\left\{  {b}_{n}\right\}$ 是等差数列.

\begin{center}
\begin{tabular}{|c|c|}
\hline
证明数列 $\left\{  {a}_{n}\right\}$ 为等差数列的方法 & 证明数列 $\left\{  {a}_{n}\right\}$ 为等比数列的方法 (注意 ${a}_{1} \neq  0$ ) \\
\cline{1-2}
\makecell{ 据定义: ${a}_{n + 1} - {a}_{n} = d$ ( $d$ 为常数)； \\ 根据通项公式: ${a}_{n} = {a}_{1} + \left( {n - 1}\right) d$ ; \\ 根据性质: $2{a}_{n + 1} = {a}_{n} + {a}_{n + 2}$ } & \makecell{  \ding{172} 据定义: ${a}_{n + 1} = q{a}_{n}$ ( $q$ 为非零常数)； \\  \ding{173}  根据通项公式: ${a}_{n} = {a}_{1}{q}^{n - 1}$ ； \\  \ding{174}  根据性质: ${a}_{n + 1}^{2} = {a}_{n}{a}_{n + 2}$ } \\
\cline{1-2}
\hline
\end{tabular}
\end{center}

11. 设 ${S}_{n}$ 是数列 $\left\{  {a}_{n}\right\}$ 的前 $n$ 项和, $\left\{  {a}_{n}\right\}$ 为等差数列的充要条件是 ${S}_{n} = a{n}^{2} + {bn}\left( {a,b\text{ 为常数 }}\right)$ ,其公差是 ${2a};\left\{  {a}_{n}\right\}$ 为等比数列的充要条件是 ${S}_{n} = a{q}^{n} - a\left( {a\text{ 、 }q\left( {q \neq  1}\right) }\right.$ 为非零常数 $)$ ,其公比是 $q$ .

\section*{三、数列的通项公式}

1. 已知前 $n$ 项和 ${S}_{n}$ ,则 ${a}_{n} = \left\{  \begin{array}{l} {S}_{1},n = 1, \\  {S}_{n} - {S}_{n - 1},n \geq  2, \end{array}\right.$ 只有当 $n = 1$ 适合 $n \geq  2$ 的表达式时才能只用第二个式子表示.

2. 一阶线性递推数列的通项公式: $\left\{  \begin{array}{l} {a}_{1}, \\  {a}_{n + 1} = p{a}_{n} + q\text{ ( }p,q\text{ 为满足条件的常数). } \end{array}\right.$

求法一(迭代法):

${a}_{n} = p{a}_{n - 1} + q = p\left( {p{a}_{n - 2} + q}\right)  + q$

$= {p}^{2}{a}_{n - 2} + {qp} + q = {p}^{2}\left( {p{a}_{n - 3} + q}\right)  + {qp} + q$

$= {p}^{3}{a}_{n - 3} + q{p}^{2} + {qp} + q = \cdots$

$= {p}^{n - 1}{a}_{1} + q{p}^{n - 2} + \cdots  + q{p}^{2} + {qp} + q$

$= {p}^{n - 1}{a}_{1} + \frac{q\left( {1 - {p}^{n - 1}}\right) }{1 - p}$

求法二(构造法 \ding{172} ):

当 $n \geq  2$ 时,由 $\left\{  \begin{array}{l} {a}_{n - 1} = p{a}_{n} + q, \\  {a}_{n} = p{a}_{n - 1} + q \end{array}\right.$ 得, ${a}_{n + 1} - {a}_{n} = p\left( {{a}_{n} - {a}_{n - 1}}\right) ,{a}_{2} = p{a}_{1} + q$ , $\therefore \left\{  {{a}_{n + 1} - {a}_{n}}\right\}$ 是以 ${a}_{2} - {a}_{1} = \left( {p - 1}\right) {a}_{1} + q$ 为首项, $p$ 为公比的等比数列,

$\therefore {a}_{n + 1} - {a}_{n} = \left\lbrack  {\left( {p - 1}\right) {a}_{1} + q}\right\rbrack  {p}^{n - 1}$ ,于是

$\left. \begin{array}{l} {a}_{n} - {a}_{n - 1} = \left\lbrack  {\left( {p - 1}\right) {a}_{1} + q}\right\rbrack  {p}^{n - 2}, \\  {a}_{n - 1} - {a}_{n - 2} = \left\lbrack  {\left( {p - 1}\right) {a}_{1} + q}\right\rbrack  {p}^{n - 3}, \\  \cdots \\  {a}_{2} - {a}_{1} = \left( {p - 1}\right) {a}_{1} + q, \\  {a}_{1} = {a}_{1}, \end{array}\right\}$ 相加得 ${a}_{n} = {a}_{1} + \frac{\left\lbrack  {\left( {p - 1}\right) {a}_{1} + q}\right\rbrack  \left( {1 - {p}^{n - 1}}\right\rbrack  }{1 - p}$ .

求法三(构造法 \ding{173} )

设 ${a}_{n + 1} + x = p\left( {{a}_{n} + x}\right)$ ,则 ${a}_{n + 1} = p{a}_{n} + \left( {p - 1}\right) x$ ,于是 $x = \frac{q}{p - 1}$ .

$\therefore {a}_{n + 1} + \frac{q}{p - 1} = p\left( {{a}_{n} + \frac{q}{p - 1}}\right)$ ,

即数列 $\left\{  {{a}_{n} + \frac{q}{p - 1}}\right\}$ 是以 ${a}_{1} + \frac{q}{p - 1}$ 为首项, $p$ 为公比的等比数列,

$\therefore {a}_{n} + \frac{q}{p - 1} = \left( {{a}_{1} + \frac{q}{p - 1}}\right) {p}^{n - 1}$ ,即 ${a}_{n} = \left( {{a}_{1} + \frac{q}{p - 1}}\right) {p}^{n - 1} - \frac{q}{p - 1}$ .

3. 一阶线性递推数列的特殊情形:

 \ding{172} 当 $p = 1$ 时，即为等差数列；

 \ding{173} 当 $p \neq  0,q = 0$ ，且 $p \neq  1$ 时，即为等比数列.

\section*{四、数列求和}

1. (类似于)等差数列前 $n$ 项和公式的推导方法中的“倒序相加法”，适用于与函数图像的对称性有关的若干项求和问题.

2. (类似于) 等比数列前 $n$ 项和公式的推导方法中的“错位相减法”,适用于等差数列与等比数列相应项相乘构成的数列求和问题.

3. “若 ${c}_{n} = {a}_{n}{b}_{n}$ ,其中 $\left\{  {a}_{n}\right\}$ 是等差数列, $\left\{  {b}_{n}\right\}$ 是等比数列,求 $\left\{  {c}_{n}\right\}$ 的前 $n$ 项”用“错位相减”法.

4. “若 ${b}_{n} = \frac{1}{{a}_{n}{a}_{n + 1}}$ ，其中 $\left\{  {a}_{n}\right\}$ 是等差数列，求 $\left\{  {b}_{n}\right\}$ 的前 $n$ 项的和”用“裂项相消法”(其中常用待定系数法进行裂项,方法如下: 设 ${b}_{n} = \frac{A}{{a}_{n}} + \frac{B}{{a}_{n + 1}}$ ,则 ${b}_{n} = \frac{A{a}_{n + 1} + B{a}_{n}}{{a}_{n}{a}_{n + 1}} = \frac{\left( {A + B}\right) {a}_{n} + {Ad}}{{a}_{n}{a}_{n + 1}}$ ,于是 $\left\{  {\begin{array}{l} A + B = 0, \\  {Ad} = 1, \end{array}A =  - B = }\right.  \frac{1}{d}$ ,即 ${b}_{n} = \frac{1}{d}\left( {\frac{1}{{a}_{n}} - \frac{1}{{a}_{n + 1}}}\right)$ .

注意: \ding{172}  “若 ${b}_{n} = \frac{1}{{a}_{n}{a}_{n + 1}{a}_{n + 2}}$ ，其中 $\left\{  {a}_{n}\right\}$ 是等差数列，求 $\left\{  {b}_{n}\right\}$ 的前 $n$ 项的和”也可用 “裂项相消法”；  \ding{173}  并非只有这种类型的求和问题才能用 “裂项相消法”，要根据具体题目的特征选用，如: $n \cdot  n! = \left( {n + 1}\right) !  - n$ ! 等; 又如 ${a}_{n} = \frac{{2}^{n}}{\left( {{2}^{n - 1} + 1}\right) \left( {{2}^{n} + 1}\right) } = 2\left( {\frac{1}{{2}^{n - 1} + 1} - \frac{1}{{2}^{n} + 1}}\right)$ 等.

5. ”若 ${c}_{n} = A{a}_{n} + B{b}_{n}\left( {A\text{ 、 }B\text{ 为常数 }}\right)$ ，其中 $\left\{  {a}_{n}\right\}$ 、 $\left\{  {b}_{n}\right\}$ 的前 $n$ 项和分别已知或可求，求 $\left\{  {c}_{n}\right\}$ 的前 $n$ 项的和”用“分组求和法”.

6. 分段数列求和要注意分段的要求,分类讨论要到位.

\section*{五、数列的极限与数学归纳法}

1. 数列 $\left\{  {a}_{n}\right\}$ 的极限是 $A$ ,说明随着 $n$ 的无限增大, ${a}_{n}$ 与 $A$ 越来越近,而且无限接近,该极限与数列的前有限项无关, 只与后面的所有项有关.

2.几个常用的数列极限:

 \ding{172}  $\mathop{\lim }\limits_{{n \rightarrow  \infty }}C = C\left( {C\text{ 为常数 }}\right)$ ;  \ding{173}  $\mathop{\lim }\limits_{{n \rightarrow  \infty }}\frac{1}{n} = 0$ ;  \ding{174}  $\mathop{\lim }\limits_{{n \rightarrow  \infty }}{q}^{n} = \left\{  \begin{array}{l} 0,\left| q\right|  < 1, \\  1,q = 1, \\  \text{ 不存在, }q =  - 1\text{ 或 }\left| q\right|  > 1 \end{array}\right.$

3. 若无穷等比数列的和为 $a$ ,由 $\frac{{a}_{1}}{1 - q} = a$ 和 $0 < \left| q\right|  < 1$ 可求 ${a}_{1}$ (或 $q$ )的范围,本质是由定义域求值域.

4. 数学归纳法证明与自然数有关问题的步骤:以下斜体字在证题过程中照抄，“······”是证明过程

 \ding{172}  (验证)当 $n = {n}_{0}$ 时， $\cdots \cdots$ 结论成立；

 \ding{173}  假设 $n = k\left( {k \geq  {n}_{0},n \in  \mathbf{N}}\right)$ 时结论成立，即······

则当 $n = k + 1$ 时， $\cdots \cdots$

即当 $n = k + 1$ 时，结论成立

由 \ding{172}  \ding{173} 可知，结论对任意 $n \geq  {n}_{0},n \in  \mathbf{N}$ 成立.

5. 数学归纳法中,  \ding{172} 是归纳基础;  \ding{173} 是归纳递推, 在第二步中必须使用归纳假设, 否则就不是数学归纳法.

6. 数学归纳法可用于证明一部分与正整数 $n$ 有关的等式、不等式、整除问题和几何问题.

\section*{$\text{ § }{7.1}$ 数列的概念与性质}

\section*{典型例题精讲}

【例 1】根据下列条件,求数列 $\left\{  {a}_{n}\right\}$ 的通项公式.

(1) ${a}_{1} = 2,{a}_{n + 1} = {a}_{n} + n + 1$ ；

(2) ${a}_{1} = \frac{2}{3},{a}_{n + 1} = \frac{n}{n + 1}{a}_{n}$ .

【例 2】已知数列 $\left\{  {a}_{n}\right\}$ 满足 ${a}_{1} = 1,{a}_{n} = 3{a}_{n - 1} + 4\left( {n \geq  2}\right)$ .

(1)求证:数列 $\left\{  {a}_{n}\right\}$ 既不是等差数列，也不是等比数列；

(2)求证:数列 $\left\{  {{a}_{n} + 2}\right\}$ 是等比数列,并求数列 $\left\{  {a}_{n}\right\}$ 的通项公式.

【例 3】已知数列 $\left\{  {a}_{n}\right\}$ 满足 ${a}_{1} = 1,2{a}_{n + 1} = {a}_{n}$ ,数列 $\left\{  {b}_{n}\right\}$ 满足 ${b}_{n} = 2 - {\log }_{2}{a}_{n + 1}^{2}$ .

(1)求数列 $\left\{  {a}_{n}\right\}$ 、 $\left\{  {b}_{n}\right\}$ 的通项公式；

(2)设数列 $\left\{  {b}_{n}\right\}$ 的前 $n$ 项和为 ${T}_{n}$ ，求使得 $2{T}_{n} \leq  4{n}^{2} + m$ 对任意正整数 $n$ 都成立的实数 $m$ 的取值范围.

【例 4】已知数列 $\left\{  {a}_{n}\right\}$ 满足 ${a}_{1} + 3{a}_{2} + 5{a}_{3} + \cdots  + \left( {{2n} - 1}\right) {a}_{n} = {3}^{n}$ .

(1)求数列 $\left\{  {a}_{n}\right\}$ 的通项公式；

(2)若对任意的 $n \in  {\mathbf{N}}^{ * }$ ，不等式 ${a}_{n} \geq  {\left( -1\right) }^{n}\lambda$ 恒成立，求实数 $\lambda$ 的取值范围.

【例 5】已知 $\left\{  {a}_{n}\right\}$ 为等差数列, ${b}_{n} = \left\{  \begin{array}{l} {a}_{n} - 6,n\text{ 为奇数,记 }{S}_{n}\text{ 、 }{T}_{n}\text{ 为 }\left\{  {a}_{n}\right\}  \text{ 和 }\left\{  {b}_{n}\right\}  \text{ 的前 }n\text{ 项和, }{S}_{4} = {32},{T}_{3} = {16}. \end{array}\right.$

(1)求 $\left\{  {a}_{n}\right\}$ 的通项公式；

(2)证明:当 $n > 5$ 时， ${T}_{n} > {S}_{n}$ .

\section*{四基达标精练}

\section*{一、填空题}

1. 在数列 $1,2,\sqrt{7},\sqrt{10},\sqrt{13},\cdots$ 中， $2\sqrt{19}$ 是这个数列的第\_\_\_\_\_项.

2. 在数列 $\left\{  {a}_{n}\right\}$ 中， ${a}_{1} = 2$ ，且 ${a}_{n} = {a}_{n - 1} + \lg \frac{n}{n - 1}\left( {n \geq  2}\right)$ ，则 ${a}_{100} =$ \_\_\_\_\_.

3. 已知数列 $\left\{  {a}_{n}\right\}$ 满足: ${a}_{1} = 1,{a}_{n + 1} = {\left( \sqrt{{a}_{n}} + 1\right) }^{2}$ ,则 ${a}_{5} =$ \_\_\_\_\_.

4. 已知数列 $\left\{  {a}_{n}\right\}$ 的通项公式是 ${a}_{n} = {100} - {dn}$ ,若该数列从第 10 项起以后各项均为负值,则实数 $d$ 的取值范围是\_\_\_\_\_.

5. 数列 $\left\{  {a}_{n}\right\}$ 满足 $\frac{1}{3}{a}_{1} + \frac{1}{{3}^{2}}{a}_{2} + \cdots  + \frac{1}{{3}^{n}}{a}_{n} = {3n} + 1\left( {n \in  {\mathbf{N}}^{ * }}\right)$ ，则 ${a}_{n} =$ \_\_\_\_\_.

6. 已知数列 $\left\{  {a}_{n}\right\}$ 满足 ${a}_{1} = {33},{a}_{n + 1} - {a}_{n} = {2n}$ ，则 $\frac{{a}_{n}}{n}$ 的最小值为\_\_\_\_\_.

7. 数列 $\left\{  {a}_{n}\right\}$ 满足:点 $\left( {n,{a}_{n}}\right)$ 都在抛物线 $y = {x}^{2} + {kx} + {156}$ 上，且 ${a}_{n + 1} > {a}_{n}$ 对任意的正整数 $n$ 恒成立，则实数 $k$ 的范围是\_\_\_\_\_.

8. 设数列 $\left\{  {a}_{n}\right\}$ 的通项公式为 ${a}_{n} = {n}^{2} + {bn}$ ，若数列 $\left\{  {a}_{n}\right\}$ 是单调递增数列，则实数 $b$ 的取值范围为\_\_\_\_\_.

9. 已知数列 $\left\{  {a}_{n}\right\}$ 是单调数列,其前 $n$ 项和为 ${S}_{n}$ ,若 ${S}_{n} + {S}_{n + 1} = {n}^{2} + n$ ,则首项 ${a}_{1}$ 的取值范围是\_\_\_\_\_.

\section*{二、选择题}

10. 已知数列 $\left\{  {a}_{n}\right\}$ 共有 10 项, ${a}_{n} = \frac{{2022} - {10}^{n}}{{2023} - {10}^{n}}$ ,则其中最大项、最小项的情况是 ( )

A. 最大项为 ${a}_{1}$ 、最小项为 ${a}_{10}$ . B. 最大项为 ${a}_{10}$ 、最小项为 ${a}_{1}$ .

C. 最大项为 ${a}_{6}$ 、最小项为 ${a}_{5}$ . D. 最大项为 ${a}_{4}$ 、最小项为 ${a}_{3}$ .

11. 已知数列 $\left\{  {a}_{n}\right\}$ 满足 ${a}_{n} = \left\{  {\begin{array}{l}  - {n}^{2} + {2tn},n \leq  5, \\  \left( {t - 1}\right) n,n > 5 \end{array}\left( {n \in  {\mathbf{N}}^{ * }}\right) }\right.$ ,且数列 $\left\{  {a}_{n}\right\}$ 是严格递增数列,则实数 $t$ 的取值范围是 ( )

A. $\left( {\frac{9}{2},\frac{19}{4}}\right)$ . B. $\left( {\frac{9}{2}, + \infty }\right)$ . C. $\left( {5, + \infty }\right)$ . D. $(1,4\rbrack$ .

12. 数列的发展史, 折射出许多有价值的数学思想方法, 对时代的进步起了重要作用, 比如意大利数学家列昂纳多·斐波拉契以兔子繁殖为例，引入“兔子数列”:即 $1,1,2,3,5,8,{13},{21},{34},{55},{89},{144},{233},\cdots$ ,也即 $F\left( 1\right)  = F\left( 2\right)  = 1,F\left( n\right)  = F\left( {n - 1}\right)  + F\left( {n - 2}\right) \left( {n \geq  3,n \in  \mathbf{N}}\right)$ ,若此数列被 4 整除后的余数构成一个新的数列 $\left\{  {b}_{n}\right\}$ ,则 ${b}_{1} + {b}_{2} + {b}_{3} + \cdots  + {b}_{2023}$ 等于 ( )

A. 2699. B. 3539. C. 2027. D. 2023.

\section*{三、解答题}

13. 已知数列 $\left\{  {a}_{n}\right\}$ 满足 ${a}_{1} = 1,{a}_{2} = \frac{1}{3}$ ,且 $2{a}_{n + 1}{a}_{n}{a}_{n - 1} = {a}_{n + 1}{a}_{n} + {a}_{n}{a}_{n - 1} - 2{a}_{n + 1}{a}_{n - 1}\left( {n \geq  2}\right)$ .

(1)求证: $\left\{  {\frac{1}{{a}_{n + 1}} - \frac{1}{{a}_{n}}}\right\}$ 是等差数列；

(2)求数列 $\left\{  {a}_{n}\right\}$ 的通项公式.

14. 已知数列 $\left\{  {a}_{n}\right\}$ 满足 ${a}_{1} = \frac{5}{6}$ ,若对于任意 $n \in  \mathbf{N},n \geq  2$ ,二次方程 ${a}_{n - 1}{x}^{2} - {a}_{n}x + 1 = 0$ 都有根 $\alpha$ 、 $\beta$ ,且满足 ${3\alpha } - {\alpha \beta } + {3\beta } = 1$ .

(1)求证: $\left\{  {{a}_{n} - \frac{1}{2}}\right\}$ 是等比数列；

(2)求数列 $\left\{  {a}_{n}\right\}$ 的前 $n$ 项和 ${S}_{n}$ .

15. 已知正项等比数列 $\left\{  {a}_{n}\right\}$ 满足 ${a}_{3} = 4,{a}_{1} + {a}_{5} = {17}$ ,公比 $q > 1$ .

(1)求数列 $\left\{  {a}_{n}\right\}$ 的通项公式；

(2)若 ${c}_{n} = \frac{{n}^{2} \cdot  {a}_{n + 1}}{{3}^{n}}$ ，判断数列 $\left\{  {c}_{n}\right\}$ 有没有最大项? 若有，求出第几项为最大项；若没有，请说明理由.

16. 流行性感冒是由流感病毒引起的急性呼吸道传染病. 某市去年 11 月曾发生流感，据统计，11 月 1 日该市的新感染者有 30 人，以后每天的新感染者比前一天的新感染者增加 50 人. 由于该市医疗部门采取措施，使该种病毒的传播得到控制，从 11 月 $k + 1\left( {9 \leq  k \leq  {29},k \in  \mathbf{N}}\right)$ 日起每天的新感染者比前一天的新感染者减少 20 人.

(1)若 $k = 9$ ，求 11 月 1 日至 11 月 10 日新感染者总人数；

(2)若到 11 月 30 日止，该市在这 30 天内的新感染者总人数为 11940 人，问 11 月中哪一天，该市新感染者人数最多? 并求这一天的新感染者人数.

\section*{§ 7.2 等差数列和等比数列一}

\section*{典型例题精讲}

【例 1】等差数列 $\left\{  {a}_{n}\right\}$ 满足 ${a}_{3} + {a}_{5} = {10},{a}_{2} = 9$ .

(1)求 $\left\{  {a}_{n}\right\}$ 的通项公式；

(2)若 $\left\{  {a}_{n}\right\}$ 的前 $n$ 项和为 ${S}_{n}$ ，求 ${S}_{n}$ 的最大值.

【例 2】记 ${S}_{n}$ 为数列 $\left\{  {a}_{n}\right\}$ 的前 $n$ 项和, ${b}_{n}$ 为数列 $\left\{  {S}_{n}\right\}$ 的前 $n$ 项积,已知 $\frac{2}{{S}_{n}} + \frac{1}{{b}_{n}} = 2$ .

(1)证明:数列 $\left\{  {b}_{n}\right\}$ 是等差数列；

(2)求 $\left\{  {a}_{n}\right\}$ 的通项公式.

【例 3】已知函数 $f\left( x\right)  = \frac{1}{2}\left( {{a}_{n + 2} - {a}_{n + 1}}\right) {x}^{2} - \left( {3{a}_{n + 1} - 4{a}_{n}}\right) x\left( {n \in  {\mathbf{N}}^{ * }}\right)$ 的对称轴是 $x = 1$ ,数列 $\left\{  {a}_{n}\right\}$ 满足, ${a}_{1} = 2,{a}_{2} = 8$ .

(1)证明数列 $\left\{  {{a}_{n + 1} - 2{a}_{n}}\right\}$ 是等比数列，并求数列 $\left\{  {a}_{n}\right\}$ 的通项公式；

(2)设 ${3}^{n}{b}_{n} = {\left( -1\right) }^{n}{a}_{n}$ ，且 $\left| {b}_{1}\right|  + \left| {b}_{2}\right|  + \cdots  + \left| {b}_{n}\right|  < m - {3n}{\left( \frac{2}{3}\right) }^{n + 1}$ 对于 $n \in  {\mathbf{N}}^{ * }$ 恒成立，求实数 $m$ 的取值范围.

【例 4】已知数列 $\left\{  {a}_{n}\right\}$ 的前 $n$ 项和为 ${S}_{n}$ ,若 ${a}_{2} = 4,\frac{{S}_{n} - n}{n} = \frac{1}{2}{a}_{n}$ .

(1)求证:数列 $\left\{  {a}_{n}\right\}$ 是等差数列；

(2)若 ${b}_{n} = \left| {{a}_{n} - {11}}\right|$ ，求数列 $\left\{  {b}_{n}\right\}$ 的前 $n$ 项的和 ${T}_{n}$ .

【例 5】已知正项数列 $\left\{  {a}_{n}\right\}$ 的前 $n$ 项和为 ${S}_{n}$ ,数列 $\left\{  {b}_{n}\right\}$ 为等比数列,满足 $4{S}_{n} = {a}_{n + 1}^{2} - {4n} - 4$ ,且 ${a}_{1} =  {b}_{1} + 1 = 2,{a}_{4} = {b}_{4}.$

(1)求证:数列 $\left\{  {a}_{n}\right\}$ 为等差数列；

(2)若从数列 $\left\{  {a}_{n}\right\}$ 中去掉数列 $\left\{  {b}_{n}\right\}$ 的项后余下的项按原来的顺序组成数列 $\left\{  {c}_{n}\right\}$ ，求 ${c}_{1} + {c}_{2} + {c}_{3} + \cdots  + {c}_{100}$ .

\section*{四基达标精练}

\section*{一、填空题}

1. 已知等差数列 $\left\{  {a}_{n}\right\}$ 的前 $n$ 项和为 ${S}_{n}$ ，且 ${a}_{2} = 3,{S}_{5} = {25}$ ，则公差 $d =$ \_\_\_\_\_.

2. 已知各项均为正数的等比数列 $\left\{  {a}_{n}\right\}$ 的前 4 项和为 15，并且 $4{a}_{1},2{a}_{3},{a}_{5}$ 成等差数列，则 ${a}_{1}  =$

3. 记 ${S}_{n}$ 为等差数列 $\left\{  {a}_{n}\right\}$ 的前 $n$ 项和,已知 ${S}_{5} = 0,{a}_{6} = 6$ ,则 ${S}_{n} =$ \_\_\_\_\_.

4. 从等差数列 ${84},{80},{76},\cdots$ 的第\_\_\_\_\_项开始,以后各项均为负值.

5. 已知 $\left\{  {a}_{n}\right\}$ 为等比数列， ${S}_{n}$ 为其前 $n$ 项和，若 $2{S}_{3} = {a}_{2} + {a}_{3} + {a}_{4}$ ，则公比 $q =$ \_\_\_\_\_.

6. 已知等比数列 $\left\{  {a}_{n}\right\}$ 中， ${a}_{2} = 2$ ， ${a}_{5} = \frac{1}{4}$ ，则满足 ${a}_{1}{a}_{2} + {a}_{2}{a}_{3} + \cdots  + {a}_{n}{a}_{n + 1} \leq  \frac{21}{2}$ 成立的最大正整数 $n$ 的值为\_\_\_\_\_.

7. 设 $f\left( x\right)  = \frac{{4}^{x}}{{4}^{x} + 2}$ ,则 $f\left( \frac{1}{11}\right)  + f\left( \frac{2}{11}\right)  + f\left( \frac{3}{11}\right)  + \cdots  + f\left( \frac{10}{11}\right)  =$ \_\_\_\_\_.

8. 等差数列 $\left\{  {a}_{n}\right\}$ 的公差 $d \neq  0$ ，其前 $n$ 项和为 ${S}_{n}$ . 若 ${S}_{10} = 0$ ，则 ${S}_{i}\left( {i = 1,2,\cdots ,{2003}}\right)$ 中不同的数值有\_\_\_\_\_个.

9. 已知数列 $\left\{  {a}_{n}\right\}$ 满足, ${a}_{1} = 1$ ,且对于任意正整数 $m\text{ 、 }n$ 都有 ${a}_{m + n} = {a}_{m} + {a}_{n} + {mn}$ 成立,则 $\frac{1}{{a}_{1}} + \frac{1}{{a}_{2}} + \cdots  +  \frac{1}{{a}_{2023}} =$ \_\_\_\_\_.

\section*{二、选择题}

10. 已知等差数列 $\left\{  {a}_{n}\right\}$ 的公差为 $d,{S}_{20} = {10M}$ ,则下列正确的是 ( )

A. $M = 2{a}_{5} + {a}_{10}$ . B. $M = {a}_{1} + 2{a}_{10}$ . C. $M = {a}_{5} + {a}_{15}$ . D. $M = 2{a}_{10} + d$ .

11. 某林厂现在的森林木材存量是 1800 万 ${{m}}^{3}$ ，木材以每年 25\% 的增长率生长，而每年要砍伐固定的木材量为 $x$ 万 ${{m}}^{3}$ ，为达到经两次砍伐后木材存量增加 50\% 的目标，则 $x$ 的值是 ( )

A. 40. B. 45. C. 50 . D. 55.

12. 数列 $\left\{  {a}_{n}\right\}$ 的前 $n$ 项和为 ${S}_{n},{a}_{1} = m$ ,且对任意的 $n \in  {\mathbf{N}}^{ * }$ 都有 ${a}_{n} + {a}_{n + 1} = {2n} + 1$ ,则下列三个命题中, 所有真命题的序号是 ( )

 \ding{172}  存在实数 $m$ ，使得 $\left\{  {a}_{n}\right\}$ 为等差数列； \ding{173}  存在实数 $m$ ，使得 $\left\{  {a}_{n}\right\}$ 为等比数列；

 \ding{174} 若存在 $k \in  {\mathbf{N}}^{ * }$ ，使得 ${S}_{k} = {S}_{k + 1} = {55}$ ，则实数 $m$ 唯一.

A.  \ding{172} . B.  \ding{172}  \ding{173} . C.  \ding{172}  \ding{174} . D.  \ding{172}  \ding{173}  \ding{174} .

\section*{三、解答题}

13. 已知数列 $\left\{  {a}_{n}\right\}$ 和 $\left\{  {b}_{n}\right\}$ ,其中 ${b}_{n} = {2}^{{a}_{n}}\left( {n \in  {\mathbf{N}}^{ * }}\right)$ ,数列 $\left\{  {{a}_{n} + {b}_{n}}\right\}$ 的前 $n$ 项和为 ${S}_{n}$ .

(1)若 ${a}_{n} = {2n}$ ，求 ${S}_{n}$ ；

(2)若 $\left\{  {b}_{n}\right\}$ 是各项为正的等比数列， ${S}_{n} = {3n}$ ，求数列 $\left\{  {a}_{n}\right\}$ 和 $\left\{  {b}_{n}\right\}$ 的通项公式.

14. 已知数列 $\left\{  {{\log }_{2}\left( {{a}_{n} - 1}\right) }\right\}$ 为等差数列,且 ${a}_{1} = 3,{a}_{3} = 9$ .

(1)求数列 $\left\{  {a}_{n}\right\}$ 的通项公式；

(2)证明 $\frac{1}{{a}_{2} - {a}_{1}} + \frac{1}{{a}_{3} - {a}_{2}} + \cdots  + \frac{1}{{a}_{n + 1} - {a}_{n}} < 1$ .

15. 已知数列 $\left\{  {a}_{n}\right\}$ 的前 $n$ 项和 ${S}_{n}$ 满足 $2{S}_{n} - n{a}_{n} = n\left( {n \in  \mathbf{N}}\right)$ ,且 ${a}_{2} = 3$ .

(1)求证:数列 $\left\{  \frac{{a}_{n} - 1}{n - 1}\right\}  \left( {n \geq  2}\right)$ 是常数列；

(2)若数列 $\left\{  {b}_{n}\right\}$ 通项公式 ${b}_{n} = {3n} - 2$ ，将数列 $\left\{  {a}_{n}\right\}$ 与 $\left\{  {b}_{n}\right\}$ 的公共项按从小到大的顺序排列得到数列 $\left\{  {c}_{n}\right\}$ ,求 $\left\{  {c}_{n}\right\}$ 的前 $n$ 项和.

16. 近两年, 直播带货逐渐成为一种新兴的营销模式, 带来电商行业的新增长点. 某直播平台第 1 年初的启动资金为 500 万元，由于一些知名主播加入，平台资金的年平均增长率可达 40\%，每年年底扣除运营成本 $a$ 万元，再将剩余资金继续投入直播平台.

(1)若 $a = {100}$ ，在第 3 年年底扣除运营成本后，直播平台的资金有多少万元？

(2)每年的运营成本最多控制在多少万元，才能使得直播平台在第 6 年年底扣除运营成本后资金达到 3000 万元?(结果精确到 0.1 万元)

\section*{$\text{ § }{7.3}$ 等差数列和等比数列二}

\section*{典型例题精讲}

【例 1】已知 ${S}_{n}$ 是数列 $\left\{  {a}_{n}\right\}$ 的前 $n$ 项和,且 ${a}_{1} = 1,{a}_{n} + {a}_{n + 1} = {2n} + 1$ . (1)求数列 $\left\{  {a}_{n}\right\}$ 的通项公式；

(2)记 ${b}_{n} = \left\{  \begin{array}{l} {2}^{{a}_{n}},n\text{ 为奇数,求数列 }\left\{  {b}_{n}\right\}  \text{ 的前 }{2n}\text{ 项和 }{T}_{n}. \end{array}\right.$

【例 2】已知数列 $\left\{  {a}_{n}\right\}$ 前 $n$ 项积为 ${T}_{n}$ ,且 ${a}_{n} + {T}_{n} = 1$ .

(1)求证:数列 $\left\{  \frac{1}{1 - {a}_{n}}\right\}$ 为等差数列；

(2)设 ${S}_{n} = {T}_{1}^{2} + {T}_{2}^{2} + \cdots  + {T}_{n}^{2}$ ，求证: ${S}_{n} > {a}_{n + 1} - \frac{1}{2}$ .

【例 3】已知正项等差数列 $\left\{  {a}_{n}\right\}$ 满足: ${a}_{3n} = 3{a}_{n}\left( {n \in  {\mathbf{N}}^{ * }}\right)$ ,且 $2{a}_{1},{a}_{3} + 1,{a}_{8}$ 成等比数列.

(1)求 $\left\{  {a}_{n}\right\}$ 的通项公式；

(2)设 ${c}_{n} = \frac{{2}^{{a}_{n + 1}}}{\left( {1 + {2}^{{a}_{n}}}\right) \left( {1 + {2}^{{a}_{n + 1}}}\right) }$ ， ${R}_{n}$ 是数列 $\left\{  {c}_{n}\right\}$ 的前 $n$ 项和，若对任意正整数 $n$ 均有 ${R}_{n} < \lambda$ 恒成立，求 $\lambda$ 的最小值.

【例 4】已知数列 $\left\{  {a}_{n}\right\}$ 的前 $n$ 项和为 ${S}_{n}$ ,且 ${S}_{n} = 2{a}_{n} - 2\left( {n \in  {\mathbf{N}}^{ * }}\right)$ .

(1)求 $\left\{  {a}_{n}\right\}$ 的通项公式；

(2)设 ${b}_{n + 1} = 2{b}_{n} - {2}^{n + 1},{b}_{1} = 8,{T}_{n}$ 是数列 $\left\{  {b}_{n}\right\}$ 的前 $n$ 项和，求正整数 $k$ ，使得对任意 $n \in  {\mathbf{N}}^{ * }$ 均有 ${T}_{k} \geq  {T}_{n}$ 恒成立.

【例 5】小张计划贷款 60 万元用于购买一套商品住房, 银行给小张提供了两种贷款方式:  \ding{172}  等额本金: 每月的还款额呈递减趋势，且从第二个还款月开始，每月还款额与上月还款额的差均相同； \ding{173}  等额本息:每个月的还款额均相同. 银行规定，在贷款到账日的次月当天开始首次还款(若 2023 年 7 月 7 日贷款到账，则 2023 年 8 月 7 日首次还款).

已知小张该笔贷款年限为 20 年, 月利率为 0.004 .

(1)若小张采取等额本金的还款方式，现已得知第一个还款月应还 4900 元，最后一个还款月应还 2510 元，试计算小张该笔贷款的总利息；

(2)若小张采取等额本息的还款方式，银行规定，每月还款额不得超过家庭平均月收入的一半，已知小张家庭平均月收入为 1 万元，判断小张该笔贷款是否能够获批(不考虑其他因素)；

(3)对比两种还款方式，从经济利益的角度来考虑，小张应选择哪种还款方式.

\section*{四基达标精练}

\section*{一、填空题}

1. 已知数列 $\left\{  {a}_{n}\right\}$ 为等差数列,若 ${a}_{2} + {a}_{8} = \frac{2\pi }{3}$ ,则 $\tan \left( {{a}_{3} + {a}_{7}}\right)  =$ \_\_\_\_\_.

2. 已知等差数列 $\left\{  {a}_{n}\right\}$ 的前 $n$ 项和是 ${S}_{n}$ ，且 ${a}_{9} = 8$ ，则 ${S}_{17} =$ \_\_\_\_\_.

3. 设等差数列 $\left\{  {a}_{n}\right\}$ 的前 $n$ 项和为 ${S}_{n}$ ,若 $\frac{{S}_{2023}}{2023} - \frac{{S}_{2022}}{2022} = 2$ ,则数列 $\left\{  {a}_{n}\right\}$ 的公差为\_\_\_\_\_.

4. 已知等差数列 $\left\{  {a}_{n}\right\}$ 的项数为奇数,其中所有奇数项之和为 319,所有偶数项之和为 290,则该数列的中间项为\_\_\_\_\_.

5. 已知正数等比数列 $\left\{  {a}_{n}\right\}$ 中，公比 $q = 2,{a}_{1}{a}_{2}{a}_{3}\cdots {a}_{30} = {2}^{45}$ ,则 ${a}_{1}{a}_{4}{a}_{7}\cdots {a}_{28} =$ \_\_\_\_\_.

6. 已知实数 $x,{a}_{1},{a}_{2},y$ 成等差数列， $x,{b}_{1},{b}_{2},y$ 成等比数列，则 $\frac{{\left( {a}_{1} + {a}_{2}\right) }^{2}}{{b}_{1}{b}_{2}}$ 的取值范围是\_\_\_\_\_.

7. 数列 $\left\{  {a}_{n}\right\}$ 中， ${a}_{1} = \frac{1}{5},{a}_{n} + {a}_{n + 1} = \frac{6}{{5}^{n + 1}}\left( {n \in  {\mathbf{N}}^{ * }}\right)$ ，则 $\mathop{\sum }\limits_{{i = 1}}^{{+\infty }}{a}_{i} =$ \_\_\_\_\_.

8. 数列 $\left\{  {a}_{n}\right\}$ 满足 ${a}_{n + 2} + {\left( -1\right) }^{n}{a}_{n} = {3n} - 1$ ，前 16 项和为 540，则 ${a}_{1} =$ \_\_\_\_\_.

9. 已知数列 $\left\{  {b}_{n}\right\}$ ， ${b}_{n + 1} - {b}_{n} = 2 = {b}_{1}\left( {n \in  {\mathbf{N}}^{ * }}\right)$ ，等比数列 $\left\{  {a}_{n}\right\}$ 中， ${a}_{1} = {b}_{1},{a}_{4} = {b}_{8}$ ，若数列 $\left\{  {b}_{n}\right\}$ 中去掉与数列 $\left\{  {a}_{n}\right\}$ 相同的项后余下的项按原顺序组成数列 $\left\{  {c}_{n}\right\}$ ,则 $\left\{  {c}_{n}\right\}$ 前 200 项的和为\_\_\_\_\_.

\section*{二、选择题}

10. 已知数列 $\left\{  {a}_{n}\right\}$ 是公差不为零的等差数列, $\left\{  {b}_{n}\right\}$ 是正项等比数列,若 ${a}_{1} = {b}_{1},{a}_{7} = {b}_{7}$ ,则 ( )

A. ${a}_{4} = {b}_{4}$ . B. ${a}_{5} < {b}_{5}$ . C. ${a}_{8} > {b}_{8}$ . D. ${a}_{9} < {b}_{9}$ .

11. 已知函数 $f\left( x\right)  = \frac{1}{3}{x}^{3} + {4x}$ ,记等差数列 $\left\{  {a}_{n}\right\}$ 的前 $n$ 项和为 ${S}_{n}$ ,若 $f\left( {{a}_{1} + 2}\right)  = {100},f\left( {{a}_{2022} + 2}\right)  =$ -100,则 ${S}_{2022}$ 等于 ( )

A. -4044 . B. -2022. C. 2022. D.4044.

12. 各项均为正数的数列 $\left\{  {a}_{n}\right\}$ 满足 ${a}_{n} = 2\sqrt{{S}_{n}} - 1$ ，其中 ${S}_{n}$ 是数列 $\left\{  {a}_{n}\right\}$ 的前 $n$ 项和，若对任意 $n \in  \mathbf{N}$ ， $n \geq  2$ ,总有 $\frac{1}{{S}_{2} - 1} + \frac{1}{{S}_{3} - 1}\cdots  + \frac{1}{{S}_{n} - 1} \leq  \lambda$ 恒成立,则实数 $\lambda$ 的最小值为 ( )

A.1. B. $\frac{3}{4}$ . C. $\frac{2}{3}$ . D. $\frac{1}{3}$ .

\section*{三、解答题}

13. 设数列 $\left\{  {a}_{n}\right\}$ 的前 $n$ 项和为 ${S}_{n}$ ,满足 ${S}_{n} = 1 - n{a}_{n}$ .

(1)求数列 $\left\{  {a}_{n}\right\}$ 的通项公式；

(2)设数列 $\left\{  \frac{{\left( -1\right) }^{n}}{{a}_{n}}\right\}$ 的前 $n$ 项和为 ${T}_{n}$ ，求 ${T}_{2n}$ 的表达式.

14. 赵先生准备通过某银行贷款 5000 元, 然后通过分期付款的方式还款. 银行与赵先生约定: 每个月还款一次，分12次还清所有欠款，且每个月还款的金额都相等，贷款的月利率为 0.5\%，则赵先生每个月所要还款的金额是多少元. (精确到 0.01 元)

15. 已知 ${S}_{n}$ 是公差不为零的等差数列 $\left\{  {a}_{n}\right\}$ 的前 $n$ 项和, ${S}_{3} = 6,{a}_{3}^{2} = {a}_{1} \cdot  {a}_{9}$ .

(1)求数列 $\left\{  {a}_{n}\right\}$ 的通项公式；

(2)设 ${b}_{n} = {\left( -1\right) }^{n}\frac{4{a}_{n}}{4{n}^{2} - 1}$ ，数列 $\left\{  {b}_{n}\right\}$ 的前 ${2n}$ 项和为 ${T}_{2n}$ ，若 $\left| {{T}_{2n} + 1}\right|  < \frac{1}{2022}$ ，求正整数 $n$ 的最小值.

16. 已知数列 $\left\{  {b}_{n}\right\}$ 的前 $n$ 项和 ${S}_{n}$ 满足 $\frac{1}{{S}_{1}} + \frac{1}{{S}_{2}} + \cdots  + \frac{1}{{S}_{n}} = \frac{2n}{n + 1}$ .

(1)求数列 $\left\{  {b}_{n}\right\}$ 的通项公式；

(2)设数列 $\left\{  {c}_{n}\right\}$ 满足 ${c}_{n} = \frac{{b}_{n}}{{2}^{n}}$ ，其前 $n$ 项和为 ${T}_{n}$ ，是否存在正整数 $n$ ，使得 ${T}_{n} = 4 - n$ 成立？若存在，求出所有 $n$ 的值；若不存在，请说明理由.

\section*{$\text{ § }{7.4}$ 数列综合}

\section*{典型例题精讲}

【例 1】已知数列 $\left\{  {a}_{n}\right\}$ 的前 $n$ 项和为 ${S}_{n}$ ,且 ${a}_{1} = 1,{a}_{n} \neq  0,4{S}_{n} = {a}_{n} \cdot  {a}_{n + 1} + 1\left( {n \in  {\mathbf{N}}^{ * }}\right)$ .

(1)求证: $\left\{  {a}_{{2k} - 1}\right\}$ 、 $\left\{  {a}_{2k}\right\}  \left( {k \in  {\mathbf{N}}^{ * }}\right)$ 都是等差数列；

(2)求数列 $\left\{  {a}_{n}\right\}$ 的通项公式.

【例 2】已知数列 $\left\{  {a}_{n}\right\}$ 的前 $n$ 项和 ${S}_{n}$ ,且 ${a}_{n} \in  \mathbf{Z},{a}_{1} = 1$ ,记 $D\left( {a}_{n}\right)  = \frac{1}{n - 1}\left( {\left| {{a}_{2} - {a}_{1}}\right|  + \left| {{a}_{3} - {a}_{2}}\right|  + \cdots  + }\right.  \left. {\mid {a}_{n} - {a}_{n - 1} \mid  }\right) \left( {n \geq  2,n \in  \mathbf{N}}\right)$ .

(1)若 ${a}_{n} = {2n} - 1$ ，求 $D\left( {a}_{3}\right)  + D\left( {S}_{3}\right)$ ；

(2)若对任意正整数 $n$ ， ${3D}\left( {S}_{n}\right)  = {2D}\left( {a}_{n}\right)$ ，求 ${a}_{n}$ .

【例 3】观察下列等式:

\[
1 - \frac{1}{2} = \frac{1}{2},1 - \frac{1}{2} + \frac{1}{3} - \frac{1}{4} = \frac{1}{3} + \frac{1}{4},1 - \frac{1}{2} + \frac{1}{3} - \frac{1}{4} + \frac{1}{5} - \frac{1}{6} = \frac{1}{4} + \frac{1}{5} + \frac{1}{6},\cdots
\]

据此规律,请你猜想出第 $n$ 个等式并证明你的结论.

【例 4】设数列 $\left\{  {a}_{n}\right\}  \text{ 、 }\left\{  {b}_{n}\right\}$ 满足 ${a}_{1} = {b}_{1} = 6,{a}_{2} = {b}_{2} = 4,{a}_{3} = {b}_{3} = 3$ ,且数列 $\left\{  {{a}_{n + 1} - {a}_{n}}\right\}$ 是等差数列,数列 $\left\{  {{b}_{n} - 2}\right\}$ 是等比数列.

(1)求数列 $\left\{  {a}_{n}\right\}$ 、 $\left\{  {b}_{n}\right\}$ 的通项公式；

(2)是否存在正整数 $k$ ，使 ${a}_{k} - {b}_{k} \in  \left( {0,\frac{1}{2}}\right)$ ？若存在，求出 $k$ ；若不存在，说明理由.

【例 5】在平面直角坐标系上,有一点列 ${P}_{0},{P}_{1},{P}_{2},{P}_{3},\cdots ,{P}_{n - 1},{P}_{n}$ ,设点 ${P}_{k}$ 的坐标 $\left( {{x}_{k},{y}_{k}}\right) \left( {k \in  \mathbf{N}}\right)$ ,其中 ${x}_{k} \in  \mathbf{Z}$ 、 ${y}_{k} \in  \mathbf{Z}$ . 记 $\Delta {x}_{k} = {x}_{k} - {x}_{k - 1},\Delta {y}_{k} = {y}_{k} - {y}_{k - 1}$ ,且满足 $\left| {\Delta {x}_{k}}\right|  \cdot  \left| {\Delta {y}_{k}}\right|  = 2\left( {k = 1,2,3,\cdots ,n}\right)$ .

(1)点 ${P}_{0}\left( {0,1}\right)$ ，点 ${P}_{1}$ 坐标满足 ${x}_{0} < {x}_{1},{y}_{0} < {y}_{1}$ ，且点 ${P}_{k}\left( {k = 0,1}\right)$ 落在直线 $l$ 上，求直线 $l$ 的方程；

(2)若 ${P}_{0}\left( {0,0}\right)$ ， ${y}_{n} = {100}$ ，且 $n \geq  {2020}$ ，求 ${T}_{n} = {x}_{0} + {x}_{1} + {x}_{2} + \cdots  + {x}_{n}$ 的最大值.

\section*{四基达标精练}

\section*{一、填空题}

1. 若首项为 1、公比为 $\frac{1}{3}$ 的无穷等比数列的各项和为 $S$ ，则 $S =$ \_\_\_\_\_.

2. 在等差数列 $\left\{  {a}_{n}\right\}$ 中,已知公差 $d \neq  0,{a}_{2}^{2} = {a}_{1}{a}_{4}$ ,若 ${a}_{1},{a}_{3},{a}_{{k}_{1}},{a}_{{k}_{2}},\cdots ,{a}_{{k}_{n}},\cdots$ 成等比数列,则 ${k}_{n} =$ \_\_\_\_\_.

3. 设数列 $\left\{  {a}_{n}\right\}   : {2}^{1},{2}^{2},{2}^{2},{2}^{3},{2}^{3},{2}^{3},{2}^{4},{2}^{4},{2}^{4},{2}^{4},{2}^{5},{2}^{5},{2}^{5},{2}^{5},{2}^{5},{2}^{5},\cdots$ . 则 ${a}_{50} =$ \_\_\_\_\_.

4. 用数学归纳法证明 “ $1 + \frac{1}{2} + \frac{1}{3} + \cdots  + \frac{1}{{2}^{n} - 1} < n\left( {n \geq  2}\right)$ ” 时,由 $n = k$ 的假设证明 $n = k + 1$ 时,不等式左边需增加的项数为\_\_\_\_\_.

5. 已知数列 $\left\{  {a}_{n}\right\}$ 是以 1 为首项、2 为公差的等差数列,且 $\left\{  {{\left( -1\right) }^{n + 1}{a}_{n}^{2}}\right\}$ 的前 $n$ 项和为 ${S}_{n}$ ,则数列 $\left\{  \frac{{S}_{2n}}{n}\right\}$ 的前 10 项和为\_\_\_\_\_.

6. 若 $a$ 、 $b$ 分别是正数 $p$ 、 $q$ 的算术平均数和几何平均数，且 $a$ 、 $b$ 、-2这三个数可适当排序后成等差数列，也可适当排序后成等比数列，则 $p + q + {pq}$ 的值构成的集合是\_\_\_\_\_.

7. 若数列 $\left\{  {a}_{n}\right\}$ 满足:对任意正整数 $n$ ， $\left\{  {{a}_{n + 1} - {a}_{n}}\right\}$ 为递减数列，则称数列 $\left\{  {a}_{n}\right\}$ 为“差递减数列”. 给出下列数列 $\left\{  {a}_{n}\right\}   :$  \ding{172}  ${a}_{n} = {3n}$ ；  \ding{173}  ${a}_{n} = {n}^{2} + 1$ ； \ding{174}  ${a}_{n} = \sqrt{n}$ ； \ding{175}  ${a}_{n} = \frac{n}{n + 1}$ ，其中是“差递减数列”的有\_\_\_\_\_.

8. 若数列 $\left\{  {a}_{n}\right\}$ 满足 ${a}_{1} = {10},{a}_{n + 1} = {a}_{n} + {18n} + {10}\left( {n \in  \mathbf{N},n \geq  1}\right)$ ，记 $\left\lbrack  x\right\rbrack$ 表示不超过实数 $x$ 的最大整数，则 $\mathop{\lim }\limits_{{n \rightarrow  \infty }}\left( {\sqrt{{a}_{n}} - \left\lbrack  \sqrt{{a}_{n}}\right\rbrack  }\right)  =$ \_\_\_\_\_.

9. 已知 $n \in  \mathbf{N},n \geq  2$ ，函数 $y = \frac{n}{{n}^{2} + 3}x + \frac{3n}{n + 3}$ 的图像与 $y$ 轴相交于点 ${A}_{n}$ ，与 $y = {\log }_{\frac{1}{n}}\left( {x - 4}\right)$ 的图像相交于点 ${B}_{n}$ ， $\bigtriangleup  O{A}_{n}{B}_{n}$ 的面积为 ${S}_{n}$ ( $O$ 为坐标原点)，则 $\mathop{\lim }\limits_{{n \rightarrow  \infty }}{S}_{n} =$ \_\_\_\_\_.

\section*{二、选择题}

10. 已知 ${a}_{n} = {n}^{2} + {\lambda n}$ ,则 “ ${a}_{1} < {a}_{2}$ ” 是 “ $\left\{  {a}_{n}\right\}$ 是严格递增数列” 的 ( )

A. 充分非必要条件. B. 必要非充分条件. C. 充要条件. D. 既非充分也非必要条件.

11. 定义在 $\lbrack 0, + \infty )$ 上的函数 $f\left( x\right)$ 满足: 当 $0 \leq  x < 2$ 时, $f\left( x\right)  = {2x} - {x}^{2}$ ; 当 $x \geq  2$ 时, $f\left( x\right)  = {3f}\left( {x - 2}\right)$ . 记函数 $f\left( x\right)$ 的极大值点从小到大依次为 ${a}_{1},{a}_{2},\cdots ,{a}_{n},\cdots$ ,并记相应的极大值依次为 ${b}_{1},{b}_{2},\cdots ,{b}_{n},\cdots$ ,则 ${S}_{20} = {a}_{1}{b}_{1} + {a}_{2}{b}_{2} + \cdots  + {a}_{20}{b}_{20}$ 的值为 ( )

A. ${19} \times  {3}^{20} + 1$ . B. ${19} \times  {3}^{19} + 1$ . C. ${20} \times  {3}^{19} + 1$ . D. ${20} \times  {3}^{20} + {1711}$ .

12. 已知 $\overrightarrow{e}\text{ 、 }\overrightarrow{f}$ 是互相垂直的单位向量,向量 $\overrightarrow{{a}_{n}}$ 满足: $\overrightarrow{e} \cdot  \overrightarrow{{a}_{n}} = n,\overrightarrow{f} \cdot  \overrightarrow{{a}_{n}} = {2n} + 1,{b}_{n}$ 是向量 $\overrightarrow{f}$ 与 $\overrightarrow{{a}_{n}}$ 夹角的正切值,则数列 $\left\{  {b}_{n}\right\}$ 是 ( )

A. 单调递增数列且 $\mathop{\lim }\limits_{{n \rightarrow  \infty }}{b}_{n} = \frac{1}{2}$ . B. 单调递减数列且 $\mathop{\lim }\limits_{{n \rightarrow  \infty }}{b}_{n} = \frac{1}{2}$ .

C. 单调递增数列且 $\mathop{\lim }\limits_{{n \rightarrow  \infty }}{b}_{n} = 2$ . D. 单调递减数列且 $\mathop{\lim }\limits_{{n \rightarrow  \infty }}{b}_{n} = 2$ .

\section*{三、解答题}

13. 设 $\left\{  {a}_{n}\right\}$ 是等差数列,公差为 $d$ ,前 $n$ 项和为 ${S}_{n}$ .

(1)设 ${a}_{1} = {40},{a}_{6} = {38}$ ，求 ${S}_{n}$ 的最大值；

(2)设 ${a}_{1} = 1,{b}_{n} = {2}^{{a}_{n}}\left( {n \in  {\mathbf{N}}^{ * }}\right)$ ，数列 $\left\{  {b}_{n}\right\}$ 的前 $n$ 项和为 ${T}_{n}$ ，且对任意的 $n \in  {\mathbf{N}}^{ * }$ ，都有 ${T}_{n} \leq  {20}$ ，求 $d$ 的取值范围.

14. 已知定理: 设数列 $\left\{  {a}_{n}\right\}$ 、 $\left\{  {b}_{n}\right\}$ 的前 $n$ 项和分别为 ${A}_{n}$ 、 ${B}_{n}$ ,则 ${A}_{n} = {B}_{n} + p$ ( $p$ 为非零常数) 的充要条件是 ${a}_{1} = {b}_{1} + p,{a}_{n} = {b}_{n}\left( {n \geq  2}\right)$ .

(1)若 ${a}_{1} = {10}$ ， ${a}_{n} = {2}^{n}\left( {n \geq  2}\right)$ ，试运用上述定理求数列 $\left\{  {a}_{n}\right\}$ 的前 $n$ 项和 ${A}_{n}$ ；

(2)已知数列 $\left\{  {a}_{n}\right\}$ 的前 $n$ 项和减去数列 $\left\{  {b}_{n}\right\}$ 的前 $n$ 项和的差为 4，数列 $\left\{  {b}_{n}\right\}$ 是首项为 2 的等差数列，且数列 $\left\{  {b}_{n}\right\}$ 的前 5 项和为 40,试运用上述定理求 ${a}_{1}$ 和 ${a}_{5}$ .

15. 定义在 $\mathbf{R}$ 上的偶函数 $f\left( x\right)$ 的图像关于 $x = 1$ 对称,且 $f\left( {{x}_{1} + {x}_{2}}\right)  = f\left( {x}_{1}\right)  \cdot  f\left( {x}_{2}\right)$ 对任意 ${x}_{1}\text{ 、 }{x}_{2} \in  \left\lbrack  {0,\frac{1}{2}}\right\rbrack$ 都成立.

(1)若 $f\left( 1\right)  = a\left( {a > 0}\right)$ ， 探求 $f\left( \frac{1}{2}\right)$ ， $f\left( \frac{1}{4}\right)$ ， $\cdots$ ， $f\left( \frac{1}{{2}^{n}}\right)$ 的值；

(2)若 $f\left( 1\right)  = a\left( {a > 0}\right) ,{a}_{n} = f\left( {{2n} + \frac{1}{2n}}\right)$ ，求 $\mathop{\lim }\limits_{{n \rightarrow  \infty }}{a}_{n}$ 的值.

16. 如果存在常数 $a$ ,使得数列 $\left\{  {a}_{n}\right\}$ 满足: 若 $x$ 是数列 $\left\{  {a}_{n}\right\}$ 中的一项,则 $\frac{a}{x}$ 也是数列 $\left\{  {a}_{n}\right\}$ 中的一项,称数列 $\left\{  {a}_{n}\right\}$ 具有倒置性质 $P$ ,常数 $a$ 是它的 “倒置系数”.

(1)若数列: $1,2,4,m\left( {m > 4}\right)$ 具有倒置性质 $P$ ，求 $m$ 和倒置系数 $a$ ；

(2)已知各项为正的有穷等比数列 $\left\{  {b}_{n}\right\}$ 的项数是 ${n}_{0}\left( {{n}_{0} \geq  3}\right)$ ，公比为 $q$ ，所有项之积是 $B$ ，求证:数列 $\left\{  {b}_{n}\right\}$ 具有倒置性质 $P$ ,并用 ${n}_{0}$ 和 $B$ 表示它的倒置系数 $a$ .

\section*{第 8 章 平面向量与复数(必修 2)}

\section*{知识梳理*}

\section*{一、向量}

1. 向量 (平面向量和空间向量): 既有大小,又有方向的量.

2. 向量的运算:加法、减法(平行四边形、三角形法则).

实数与向量的乘积 $k\overrightarrow{a} : k > 0,k\overrightarrow{a}$ 与 $\overrightarrow{a}$ 同向; $k < 0,k\overrightarrow{a}$ 与 $\overrightarrow{a}$ 反向; $k = 0$ ,零向量; $k\overrightarrow{a}$ 的模为 $\left| k\right| \left| \overrightarrow{a}\right|$

向量的数量积: $\overrightarrow{a} \cdot  \overrightarrow{b} = \left| \overrightarrow{a}\right| \left| \overrightarrow{b}\right| \cos \theta$ ( $\theta$ 为 $\overrightarrow{a}$ 和 $\overrightarrow{b}$ 的夹角); 变式: $\cos \theta  = \frac{\overrightarrow{a} \cdot  \overrightarrow{b}}{\left| \overrightarrow{a}\right|  \cdot  \left| \overrightarrow{b}\right| }$ .

数量积的几何意义: $\left| \overrightarrow{a}\right|$ 与 $\overrightarrow{b}$ 在 $\overrightarrow{a}$ 方向上的数量投影的乘积 (该投影可正可负可为零).

3. 向量的坐标运算: $\overrightarrow{a} = \left( {{x}_{1},{y}_{1}}\right) ,\overrightarrow{b} = \left( {{x}_{2},{y}_{2}}\right)$ ,则 $\overrightarrow{a} \pm  \overrightarrow{b} = \left( {{x}_{1} \pm  {x}_{2},{y}_{1} \pm  {y}_{2}}\right) ;k\overrightarrow{a} = \left( {k{x}_{1},k{y}_{1}}\right) ;\overrightarrow{a} \cdot  \overrightarrow{b} =  {x}_{1} \cdot  {x}_{2} + {y}_{1} \cdot  {y}_{2}$

4. $\overrightarrow{a}//\overrightarrow{b} \Leftrightarrow$ 存在 $k \in  \mathbf{R}$ ,使 $\overrightarrow{a} = k\overrightarrow{b}\left( {\overrightarrow{b} \neq  \overrightarrow{0}}\right)  \Leftrightarrow  {x}_{1}{y}_{2} = {x}_{2}{y}_{1} \Leftrightarrow  {\left( \overrightarrow{a} \cdot  \overrightarrow{b}\right) }^{2} = {\left| \overrightarrow{a}\right| }^{2}{\left| \overrightarrow{b}\right| }^{2}$ .

5. $\overrightarrow{a} \bot  \overrightarrow{b} \Leftrightarrow  \overrightarrow{a} \cdot  \overrightarrow{b} = 0 \Rightarrow  \left| {\overrightarrow{a} + \overrightarrow{b}}\right|  = \left| {\overrightarrow{a} - \overrightarrow{b}}\right|  \Leftrightarrow  {x}_{1} \cdot  {x}_{2} + {y}_{1} \cdot  {y}_{2} = 0$ (空间向量类似).

6. $\overrightarrow{a},\overrightarrow{b}$ 的夹角为锐角 $\Leftrightarrow  \left\{  {\begin{array}{l} \overrightarrow{a} \cdot  \overrightarrow{b} > 0, \\  \overrightarrow{a} \times  \overrightarrow{b} \end{array};\overrightarrow{a}\text{ 、 }\overrightarrow{b}}\right.$ 的夹角为钝角 $\Leftrightarrow  \left\{  \begin{array}{l} \overrightarrow{a} \cdot  \overrightarrow{b} < 0, \\  \overrightarrow{a} \times  \overrightarrow{b}. \end{array}\right.$

7. 定比分点公式: 若 $M\left( {{x}_{1},{y}_{1}}\right) ,N\left( {{x}_{2},{y}_{2}}\right)$ 满足 $\overrightarrow{MP} = \lambda \overrightarrow{PN}$ ,则有 $P\left( {\frac{{x}_{1} + \lambda {x}_{2}}{1 + \lambda },\frac{{y}_{1} + \lambda {y}_{2}}{1 + \lambda }}\right)$ .

8. 常用结论: \ding{172}  ${\overset{ \rightarrow  }{a}}^{2} = {\left| \overset{ \rightarrow  }{a}\right| }^{2}$ ;  \ding{173}  $\overset{ \rightarrow  }{a} \cdot  \overrightarrow{b} \leq  \left| \overset{ \rightarrow  }{a}\right| \left| \overset{ \rightarrow  }{b}\right|$ ;  \ding{174}  ${\left( \overrightarrow{a} \pm  \overrightarrow{b}\right) }^{2} = {\left| \overrightarrow{a} \pm  \overrightarrow{b}\right| }^{2} = {\overrightarrow{a}}^{2} \pm  2\overrightarrow{a} \cdot  \overrightarrow{b} + {\overrightarrow{b}}^{2}$ ；

 \ding{175}  $\left| \right| \overrightarrow{a}\left| -\right| \overrightarrow{b}\left| \right|  \leq  \left| {\overrightarrow{a} \pm  \overrightarrow{b}}\right|  \leq  \left| \overrightarrow{a}\right|  + \left| \overrightarrow{b}\right|$ ;  \circled{5}  ${\left| \overrightarrow{a} + \overrightarrow{b}\right| }^{2} + {\left| \overrightarrow{a} - \overrightarrow{b}\right| }^{2} = 2\left( {{\left| \overrightarrow{a}\right| }^{2} + {\left| \overrightarrow{b}\right| }^{2}}\right)$ ;

 \circled{6} 非零向量 $\overrightarrow{a}$ 、 $\overrightarrow{b}$ 、 $\overrightarrow{c}$ 共面的充要条件是存在 $\lambda$ 、 $u \in  \mathbf{R}$ 使得 $\overrightarrow{a} = \lambda \overrightarrow{b} + u\overrightarrow{c}$ .

9. 向量的投影和数量投影:

$\left| \overrightarrow{a}\right| \cos \langle \overrightarrow{a},\overrightarrow{b}\rangle  \cdot  \frac{\overrightarrow{b}}{\left| \overrightarrow{b}\right| } = \frac{\overrightarrow{a} \cdot  \overrightarrow{b}}{\left| \overrightarrow{b}\right| } \cdot  \frac{\overrightarrow{b}}{\left| \overrightarrow{b}\right| }$ 表示向量 $\overrightarrow{a}$ 在向量 $\overrightarrow{b}$ 方向上的投影是向量.

其中系数 $\left| \overrightarrow{a}\right| \cos \langle \overrightarrow{a},\overrightarrow{b}\rangle  = \frac{\overrightarrow{a} \cdot  \overrightarrow{b}}{\left| \overrightarrow{b}\right| }$ 表示向量 $\overrightarrow{a}$ 在向量 $\overrightarrow{b}$ 方向上的数量投影.

10. 平面向量基本定理: 平面内任意一个向量都可以表示为该平面内两个不平行的向量的线性组合.

\section*{二、复数}

1. 引进虚数单位 ${i}$ 的原则: ${\left( {D}{i}\right) }^{2} =  - 1$ ;  \ding{173}  ${i}$ 可以和实数进行四则运算，且不改变原有的运算律.

2. 对于复数 $z = a + b{i}\left( {a\text{ 、 }b \in  \mathbf{R}}\right) ,a$ 叫实部,记为 $\operatorname{Re}z;b$ 叫虚部,记为 $\operatorname{Im}z$ .

3. 复数相等: $a + b{i} = c + d{i}\left( {a\text{ 、 }b\text{ 、 }c\text{ 、 }d \in  \mathbf{R}}\right)  \Leftrightarrow  \left\{  \begin{array}{l} a = b, \\  c = d; \end{array}\right.$ 不全是实数的两个复数不能比较大小.

4. $z \in  \mathbf{R} \Leftrightarrow  \operatorname{Im}z = 0 \Leftrightarrow  z = \bar{z} \Leftrightarrow  {z}^{2} \geq  0;z$ 为纯虚数 $\Leftrightarrow  \left\{  {\begin{array}{l} \operatorname{Re}z = 0, \\  \operatorname{Im}z \neq  0 \end{array} \Leftrightarrow  \left\{  {\begin{array}{l} z + \bar{z} = 0, \\  z \neq  0 \end{array} \Leftrightarrow  {z}^{2} < 0}\right. }\right.$ .

\HRule

*注:本章节所属知识模块是 “图形与几何”.

\HRule

5. 复数 $z = a + b{i}\left( {a\text{ 、 }b \in  \mathbf{R}}\right)$ 所对应的点即直角坐标平面内的点 $\left( {a,b}\right)$ ; 也对应于以 $\left( {a,b}\right)$ 为坐标的位置向量.

6.* 复数的三角形式:

$z = r\left( {\cos \theta  + {i}\sin \theta }\right) ,\theta  \in  \mathbf{R},\theta$ 为辐角,记作 $\operatorname{Arg}z,r$ 是复数的模,若取 $\theta  \in  \lbrack 0,{2\pi })$ ,则称 $\theta$ 为辐角主值,记作 $\arg z$ !

7. 复数的运算法则(以下 $a$ 、 $b$ 、 $c$ 、 $d \in  \mathbf{R},m$ 、 $n \in  \mathbf{N},m$ 、 $n \geq  2$ ):

 \ding{172}  $\left( {a + b{i}}\right)  \pm  \left( {c + d{i}}\right)  = \left( {a \pm  c}\right)  + \left( {b \pm  d}\right) {i}$ ；

 \ding{173}  $\left( {a + b{i}}\right)  \div  \left( {c + d{i}}\right)  = \frac{{ac} + {bd}}{{c}^{2} + {d}^{2}} + \frac{{bc} - {ad}}{{c}^{2} + {d}^{2}}{i}\left( {c + d{i} \neq  0}\right)$ ；

 \ding{174}  ${z}^{n} = \underset{n}{\underbrace{z \cdot  z \cdot  z\cdots z}}\left( {n \in  {\mathbf{N}}^{ * }}\right)$ ；

 \ding{175}  $\left( {a + b{i}}\right) \left( {c + d{i}}\right)  = \left( {{ac} - {bd}}\right)  + \left( {{bc} + {ad}}\right) {i}$ ；

 \circled{5}  ${z}^{m}{z}^{n} = {z}^{m + n},{\left( {z}^{m}\right) }^{n} = {z}^{mn},{\left( {z}_{1}{z}_{2}\right) }^{n} = {z}_{1}^{n}{z}_{2}^{n}$ .

本质:复数的加、减、乘、乘方运算仿实数的多项式运算(合并同类项)，并随时将 ${{i}}^{2}$ 换成 -1，复数的除法运算的本质是将分母实数化 (利用 $z\bar{z} = {\left| z\right| }^{2}$ ).

8. 常用的运算结论:

${{i}}^{2} =  - 1,{{i}}^{{4n} + 1} = {i},{{i}}^{{4n} + 2} =  - 1,{{i}}^{{4n} + 3} =  - {i},{{i}}^{4n} = 1,{{i}}^{n} + {{i}}^{n + 1} + {{i}}^{n + 2} + {{i}}^{n + 3} = 0\left( {n \in  \mathbf{Z}}\right)$

${\left( 1 \pm  {i}\right) }^{2} =  \pm  2{i},\frac{1 + {i}}{1 - {i}} = {i},\frac{1 - {i}}{1 + {i}} =  - {i}.$

9. 共轭复数的运算性质: $\bar{z} = z;\overline{{z}_{1} \pm  {z}_{2}} = \overline{{z}_{1}} \pm  \overline{{z}_{2}};\overline{{z}_{1} \cdot  {z}_{2}} = \overline{{z}_{1}} \cdot  \overline{{z}_{2}};\overline{\left( \frac{{z}_{1}}{{z}_{2}}\right) } = \frac{\overline{{z}_{1}}}{{z}_{2}}\left( {{z}_{2} \neq  0}\right) ;\overline{{z}^{n}} = {\left( \bar{z}\right) }^{n}$ .

10. 复数模的性质:

$z \cdot  \bar{z} = {\left| z\right| }^{2} = {\left| \bar{z}\right| }^{2} = \left| {z}^{2}\right| ;\left| {{z}_{1} \cdot  {z}_{2}}\right|  = \left| {z}_{1}\right| \left| {z}_{2}\right| ;\left| \frac{{z}_{1}}{{z}_{2}}\right|  = \frac{\left| {z}_{1}\right| }{\left| {z}_{2}\right| }\left( {\left| {z}_{2}\right|  \neq  0}\right) ;\left| {z}^{n}\right|  = {\left| z\right| }^{n}\left( {n \in  \mathbf{N},n \geq  1}\right) ;$

$\left| \right| {z}_{1}\left| -\right| {z}_{2}\left| \right|  \leq  \left| {{z}_{1} \pm  {z}_{2}}\right|  \leq  \left| {z}_{1}\right|  + \left| {z}_{2}\right| ;{\left| {z}_{1} + {z}_{2}\right| }^{2} + {\left| {z}_{1} - {z}_{2}\right| }^{2} = 2\left( {{\left| {z}_{1}\right| }^{2} + {\left| {z}_{2}\right| }^{2}}\right) .$

11. 解复数问题的两个转化:代数化，几何化. 设出一个复数后要注明实部和虚部都是实数.

12. 解复数方程往往是设 $x = a + b{i}\left( {a\text{ 、 }b \in  \mathbf{R}}\right)$ ,利用复数相等转化为实数方程组求解,其本质是两个复数相等的概念.

13. 对于实系数一元二次方程 $a{x}^{2} + {bx} + c = 0$ ,韦达定理一定成立,

 \ding{172}  当 $\Delta  = {b}^{2} - {4ac} > 0$ 时,方程有两个不相等的实根 ${x}_{1}\text{ 、 }{x}_{2} = \frac{-b \pm  \sqrt{{b}^{2} - {4ac}}}{2a}$ ;

此时 $\left| {{x}_{1} - {x}_{2}}\right|  = \sqrt{{\left( {x}_{1} - {x}_{2}\right) }^{2}} = \sqrt{{x}_{1}^{2} + {x}_{2}^{2} - 2{x}_{1}{x}_{2}} = \sqrt{{\left( {x}_{1} + {x}_{2}\right) }^{2} - 4{x}_{1}{x}_{2}} = \frac{\sqrt{{b}^{2} - {4ac}}}{\left| a\right| }$ ;

$\left| {x}_{1}\right|  + \left| {x}_{2}\right|  = \sqrt{{x}_{1}^{2} + {x}_{2}^{2} + 2\left| {{x}_{1}{x}_{2}}\right| } = \sqrt{{\left( {x}_{1} + {x}_{2}\right) }^{2} - 2{x}_{1}{x}_{2} + 2\left| {{x}_{1}{x}_{2}}\right| } = \frac{\sqrt{{b}^{2} - {2ac} + 2\left| {ac}\right| }}{\left| a\right| };$

 \ding{173}  当 $\Delta  = {b}^{2} - {4ac} = 0$ 时,方程有两个相等的实根 ${x}_{1} = {x}_{2} = \frac{-b}{2a}$ ;

此时 $\left| {{x}_{1} - {x}_{2}}\right|  = 0;\left| {x}_{1}\right|  + \left| {x}_{2}\right|  = \frac{\left| b\right| }{\left| a\right| }$ ;

 \ding{174}  当 $\Delta  = {b}^{2} - {4ac} < 0$ 时,方程有两个共轭虚根 ${x}_{1}\text{ 、 }{x}_{2} = \frac{-b \pm  {i}\sqrt{{4ac} - {b}^{2}}}{2a}$ ;

此时 $\left| {{x}_{1} - {x}_{2}}\right|  = \sqrt{{\left| {x}_{1} - {x}_{2}\right| }^{2}} = \sqrt{\left( {{x}_{1} - {x}_{2}}\right) \left( {{\bar{x}}_{1} - {\bar{x}}_{2}}\right) } = \sqrt{-{\left( {x}_{1} - {x}_{2}\right) }^{2}} = \frac{\sqrt{{4ac} - {b}^{2}}}{\left| a\right| }$ ;

$\left| {x}_{1}\right|  + \left| {x}_{2}\right|  = 2\left| {x}_{1}\right|  = 2\sqrt{{\left| {x}_{1}\right| }^{2}} = 2\sqrt{{x}_{1}{\bar{x}}_{1}} = 2\sqrt{{x}_{1}{x}_{2}} = 2\sqrt{\frac{c}{a}}.$

\section*{§ 8.1 平面向量及其运算}

\section*{典型例题精讲}

【例 1】在三角形 ${ABC}$ 中, $\overrightarrow{AB} \bot  \overrightarrow{AC}$ ,点 $F$ 为边 ${AC}$ 中点,点 $E$ 在边 ${AB}$ 上,且 $\overrightarrow{AE} = 2\overrightarrow{EB},{CE}$ 与 ${BF}$ 相交于点 $P$ .

(1)将向量 $\overrightarrow{AP}$ 用向量 $\overrightarrow{AB}$ 、 $\overrightarrow{AC}$ 表示；

(2)若 ${AB} = 3,{AC} = 4$ ，求 $\overrightarrow{AP} \cdot  \overrightarrow{BC}$ .

\begin{center}
\includegraphics[width=0.2\textwidth]{images/bo_d7tk7oc91nqc73875fu0_143_1207_653_233_273_0.jpg}
\end{center}
 

(例 1 图)

【例 2】已知非零向量 $\overrightarrow{a}$ 、 $\overrightarrow{b}$ 满足 $\left| \overrightarrow{a}\right|  = \left| \overrightarrow{b}\right|  = 2\left| {\overrightarrow{a} + \overrightarrow{b}}\right|$ .

(1)求 $\overrightarrow{a}$ 、 $\overrightarrow{b}$ 的夹角的余弦值；

(2)若 $\left| {2\overrightarrow{a} + \overrightarrow{b}}\right|  = 2\sqrt{6}$ ，求 $\left| \overrightarrow{a}\right|$ .

【例 3】如图所示,在 $\bigtriangleup {ABO}$ 中, $\overrightarrow{OC} = \frac{1}{4}\overrightarrow{OA},\overrightarrow{OD} = \frac{1}{2}\overrightarrow{OB},{AD}$ 与 ${BC}$ 相交于点 $M$ ,设 $\overrightarrow{OA} = \overrightarrow{a},\overrightarrow{OB} = \overrightarrow{b}$ .

(1)试用向量 $\overrightarrow{a}$ 、 $\overrightarrow{b}$ 表示 $\overrightarrow{OM}$ ；

(2)过点 $M$ 作直线 ${EF}$ 分别交线段 ${AC}$ 、 ${BD}$ 于点 $E$ 、 $F$ ，记 $\overrightarrow{OE} = \lambda \overrightarrow{OA},\overrightarrow{OF} = \mu \overrightarrow{OB}$ ，求证:不论点 $E$ 、 $F$ 在线段 ${AC}\text{ 、 }{BD}$ 上如何移动, $\frac{1}{\lambda } + \frac{3}{\mu }$ 为定值.

\begin{center}
\includegraphics[width=0.2\textwidth]{images/bo_d7tk7oc91nqc73875fu0_144_1215_400_267_269_0.jpg}
\end{center}
 

(例 3 图)

【例 4】已知向量 $\overrightarrow{a}$ 、 $\overrightarrow{b}$ 的夹角为 ${135}^{ \circ  }$ ，且 $\left| \overrightarrow{a}\right|  = 2$ ， $\left| \overrightarrow{b}\right|  = 2\sqrt{2}$ .

(1)求 $\left| {\overrightarrow{a} + 2\overrightarrow{b}}\right|$ ；

(2) $\overrightarrow{c} = \overrightarrow{a} + x\overrightarrow{b}$ (其中 $x \in  \mathbf{R}$ )，当 $\left| \overrightarrow{c}\right|$ 取最小值时，求 $\overrightarrow{c}$ 与 $\overrightarrow{b}$ 的夹角的大小.

【例 5】已知在平面直角坐标系中,有点 $A\left( {a,0}\right)$ 、点 $B\left( {0,b}\right)$ (其中 $a\text{ 、 }b$ 为常数,且 ${ab} \neq  0$ ),点 $O$ 为坐标原点.

(1)设点 $P$ 为线段 ${AB}$ 靠近点 $A$ 的三等分点， $\overrightarrow{OP} = \lambda \overrightarrow{OA} + \left( {1 - \lambda }\right) \overrightarrow{OB}$ ，求实数 $\lambda$ 的值；

(2)如图，设点 ${P}_{1}$ 、 ${P}_{2}$ 、 $\cdots$ 、 ${P}_{n - 1}$ 是线段 ${AB}$ 的 $n$ 等分点， $\overrightarrow{O{P}_{k}} = \mu \overrightarrow{OA} + \left( {1 - \mu }\right) \overrightarrow{OB}$ ，其中 $1 \leq  k \leq  n - 1$ ( $n$ 、 $k \in  {\mathbf{N}}^{ * },n \geq  2)$ ,当 $n = {2023}$ 时,求 $\left| {\overrightarrow{OA} + \overrightarrow{O{P}_{1}} + \overrightarrow{O{P}_{2}} + \cdots  + \overrightarrow{O{P}_{n - 1}} + \overrightarrow{OB}}\right|$ 的值(用含 $a$ 、 $b$ 的式子表示).

\begin{center}
\includegraphics[width=0.3\textwidth]{images/bo_d7tk7oc91nqc73875fu0_145_975_387_467_253_0.jpg}
\end{center}
 

(例 5 图)

\section*{四基达标精练}

\section*{一、填空题}

1. 已知向量 $\overrightarrow{a} = \left( {2,1}\right) ,\overrightarrow{b} = \left( {-2,4}\right)$ ，则 $\left| {\overrightarrow{a} - \overrightarrow{b}}\right|  =$ \_\_\_\_\_.

2. 已知向量 $\overrightarrow{a} = \left( {\lambda , - 3}\right) ,\overrightarrow{b} = \left( {1,\lambda }\right)$ ，若 $\left( {\overrightarrow{a} + \overrightarrow{b}}\right)  \bot  \overrightarrow{b}$ ，则 $\lambda  =$ \_\_\_\_\_.

3. 已知向量 $\overrightarrow{a}$ 、 $\overrightarrow{b}$ 不平行，若向量 $\overrightarrow{p} = \overrightarrow{a} + \frac{4}{3}m\overrightarrow{b}$ 与向量 $\overrightarrow{q} = \overrightarrow{b} + {3m}\overrightarrow{a}$ 平行，则 $m =$ \_\_\_\_\_.

4. 若空间向量 $\overrightarrow{a} = \left( {1,0,1}\right) ,\overrightarrow{b} = \left( {0,1, - 1}\right)$ ，则向量 $\overrightarrow{a}$ 与 $\overrightarrow{b}$ 的夹角为\_\_\_\_\_.

5. 在 $\bigtriangleup {ABC}$ 中, ${AB} = 2,{AC} = 1,D$ 是 ${BC}$ 边上的中点,则 $\overrightarrow{AD} \cdot  \overrightarrow{BC} =$ \_\_\_\_\_.

6. 若点 $M$ 是 $\bigtriangleup  {ABC}$ 所在平面内一点，且满足 $\overrightarrow{AM} = \frac{3}{5}\overrightarrow{AB} + \frac{2}{5}\overrightarrow{AC}$ . 则 $\bigtriangleup  {ABM}$ 与 $\bigtriangleup  {ABC}$ 的面积之比为

7. 已知不共线向量 $\overrightarrow{OA}\text{ 、 }\overrightarrow{OB}$ 夹角为 $\alpha ,\left| \overrightarrow{OA}\right|  = 1,\left| \overrightarrow{OB}\right|  = 2,\overrightarrow{OP} = \left( {1 - t}\right) \overrightarrow{OA},\overrightarrow{OQ} = t\overrightarrow{OB}$ ，其中 $0 \leq  t \leq  1$ ， $\left| \overrightarrow{PQ}\right|$ 在 $t = {t}_{0}$ 处取小值，则当 $0 < {t}_{0} < \frac{1}{5}$ 时， $\alpha$ 的取值范围是\_\_\_\_\_.

8. 已知 $P\left( {0,\sqrt{3}}\right)$ ,将函数 $f\left( x\right)  = 4\cos \frac{\pi }{2}x$ 和直线 $g\left( x\right)  = x - 1$ 的所有交点从左到右依次记为 ${A}_{1}$ 、 ${A}_{2}\text{ 、 }\cdots \text{ 、 }{A}_{n}$ ,则 $\left| {\overrightarrow{P{A}_{1}} + \overrightarrow{P{A}_{2}} + \cdots  + \overrightarrow{P{A}_{n}}}\right|  =$ \_\_\_\_\_.

9. 在平面直角坐标系 ${xOy}$ 中,角 $\alpha$ 的顶点在原点,始边与 $x$ 轴非负半轴重合,终边与单位圆交于点 $A$ , 角 $\alpha  + \beta$ 的终边交单位圆于点 $B$ ，且 $\alpha$ 、 $\beta  \in  \left( {0,\frac{\pi }{2}}\right)$ . 记 $A\left( {{x}_{1},{y}_{1}}\right)$ ， $B\left( {{x}_{2},{y}_{2}}\right)$ . 若 ${x}_{2} =  - \frac{1}{4}$ ，且 $\overrightarrow{OA} \cdot  \overrightarrow{OB} = \frac{1}{2}$ ， 那么 $\sqrt{3}{x}_{1} + {y}_{1} =$ \_\_\_\_\_.

\section*{二、选择题}

10. 若 $\overrightarrow{a}\text{ 、 }\overrightarrow{b}$ 为不平行的非零向量, $\overrightarrow{AB} = \overrightarrow{a} + 5\overrightarrow{b},\overrightarrow{BC} =  - 2\overrightarrow{a} + 8\overrightarrow{b},\overrightarrow{CD} = 3\overrightarrow{a} - 3\overrightarrow{b}$ ,则 ( )

A. $A\text{ 、 }B\text{ 、 }C$ 三点共线. B. $A\text{ 、 }B\text{ 、 }D$ 三点共线. C. $B\text{ 、 }C\text{ 、 }D$ 三点共线. D. $A\text{ 、 }C\text{ 、 }D$ 三点共线.

11. 已知椭圆 $\frac{{x}^{2}}{16} + \frac{{y}^{2}}{12} = 1$ 的左顶点为 $A$ ,右焦点为 $F,M$ 是椭圆上任意一点,则 $\overrightarrow{MA} \cdot  \overrightarrow{MF}$ 的取值范围为 ( )

A. $\left\lbrack  {-{16},0}\right\rbrack$ . B. $\left\lbrack  {-8,0}\right\rbrack$ . C. $\left\lbrack  {0,8}\right\rbrack$ . D. $\left\lbrack  {0,{16}}\right\rbrack$ .

12. 设向量 $\overrightarrow{a}$ 与 $\overrightarrow{b}$ 的夹角为 $\theta$ ，定义 $\overrightarrow{a} \oplus  \overrightarrow{b} = \left| {\overrightarrow{a}\sin \theta  - \overrightarrow{b}\cos \theta }\right|$ ，已知 $\left| \overrightarrow{b}\right|  = \left| {\overrightarrow{a} - \overrightarrow{b}}\right|  = 1$ ， $\left| \overrightarrow{a}\right|  = \sqrt{2}$ ，则 $\overrightarrow{a} \oplus  \overrightarrow{b} =$ ( )

A. $\frac{\sqrt{2}}{2}$ . B. $\sqrt{2}$ . C. $\frac{\sqrt{3}}{2}$ . D. $\sqrt{3}$ .

\section*{三、解答题}

13. 已知 $\left| \overrightarrow{a}\right|  = 2,\left| \overrightarrow{b}\right|  = 3,\left( {2\overrightarrow{a} + \overrightarrow{b}}\right)  \cdot  \left( {\overrightarrow{a} - 2\overrightarrow{b}}\right)  =  - {19}$ .

(1)求向量 $\overrightarrow{a}$ 与 $\overrightarrow{b}$ 的夹角；

(2)求 $\left| {\overrightarrow{a} + 2\overrightarrow{b}}\right|$ .

14. 在平面直角坐标系中, $O$ 为坐标原点, $A\left( {1,1}\right) ,B\left( {2, - 3}\right)$ .

(1)若 $\overrightarrow{{OA} \bot  }$ 上( $\overrightarrow{OA} + \lambda \overrightarrow{AB}$ )，求实数 $\lambda$ 的值；

(2)设 $C\left( {-6,k}\right)$ ，若 $\overrightarrow{AB}$ 、 $\overrightarrow{BC}$ 的夹角为钝角，求实数 $k$ 的取值范围.

15. 在平面直角坐标系中,已知 $\overrightarrow{a} = \left( {1, - 2}\right) ,\overrightarrow{b} = \left( {3,4}\right)$ .

(1)若 $\left( {3\overrightarrow{a} - \overrightarrow{b}}\right) //\left( {\overrightarrow{a} + k\overrightarrow{b}}\right)$ ，求实数 $k$ 的值；

(2) 求 $3\overrightarrow{a} - \overrightarrow{b}$ 与 $\overrightarrow{a} + 3\overrightarrow{b}$ 的夹角.

16. 在 $\bigtriangleup {ABC}$ 中,设 $\overrightarrow{AC} = \overrightarrow{a},\overrightarrow{AB} = \overrightarrow{b}$ ,若 $\left| \overrightarrow{a}\right|  = 2,\left| \overrightarrow{b}\right|  = 3,\angle {BAC} = {60}^{ \circ  },\overrightarrow{AD} = 2\overrightarrow{DB},E$ 为 ${BC}$ 的中点， ${CD}$ 与 ${AE}$ 交于点 $O$ .

(1)求 $\overrightarrow{AE} \cdot  \overrightarrow{CD}$ 的值；

(2)若 $\overrightarrow{OD} = \lambda \overrightarrow{CD}$ ，求 $\lambda$ 的值.

\section*{§ 8.2 平面向量及其应用一}

\section*{典型例题精讲}

【例 1】已知点 $A\left( {1,4}\right) ,B\left( {5,2}\right) ,C\left( {2,2}\right)$ ,点 $P$ 是直线 ${OC}$ 上的动点 $(O$ 为坐标原点), $\overrightarrow{PA}//\overrightarrow{PB}$ . (1)求 $\overrightarrow{OP}$ 的坐标；

(2)求 $\overrightarrow{OP}$ 在 $\overrightarrow{AB}$ 方向上的向量投影.

【例 2】在 $\bigtriangleup {ABC}$ 中, ${AB} = 2,{AC} = 4,\cos \angle {BAC} = \frac{1}{4},D$ 是线段 ${BC}$ 的靠近点 $B$ 的三等分点.

(1)用 $\overrightarrow{AB}$ 、 $\overrightarrow{AC}$ 表示 $\overrightarrow{AD}$ ；

(2)求 ${AD}$ 的长度.

\begin{center}
\includegraphics[width=0.3\textwidth]{images/bo_d7tk7oc91nqc73875fu0_148_1100_1297_387_166_0.jpg}
\end{center}
 

(例 2 图)

【例 3】已知 $\bigtriangleup {ABC}$ 中, $\angle {BAC} = {90}^{ \circ  }$ . 以 ${AB}$ 为一边向外作等边三角形 ${ABD}$ (如图所示),且 $\overrightarrow{CD} =  \lambda \overrightarrow{AB} + \mu \overrightarrow{AC}\left( {\lambda ,\mu  \in  \mathbf{R}}\right)$ .

(1)当 ${AB} = {AC}$ 时，求 $\lambda  + \mu$ 的值；

(2)当 $\angle {BCD} = 2\angle {ACD}$ 时，求 $\lambda  - \mu$ 的值.

\begin{center}
\includegraphics[width=0.2\textwidth]{images/bo_d7tk7oc91nqc73875fu0_149_1114_325_338_231_0.jpg}
\end{center}
 

(例 3 图)

【例 4】已知向量 $\overrightarrow{a} = \left( {2\sin \theta ,\sin \theta  + \cos \theta }\right) ,\overrightarrow{b} = \left( {\cos \theta ,2 - m}\right)$ ,函数 $f\left( \theta \right)  = \overrightarrow{a} \cdot  \overrightarrow{b}$ .

(1)当 $m = 0$ 时，求 $f\left( \frac{\pi }{6}\right)$ 的值；

(2)若不等式 $f\left( \theta \right)  + \frac{4}{\sin \theta  + \cos \theta } > {2m} - 3$ 对所有 $\theta  \in  \left\lbrack  {0,\frac{\pi }{2}}\right\rbrack$ 恒成立. 求实数 $m$ 的范围.

【例 5】在 $\bigtriangleup {ABC}$ 中,点 $P$ 为 $\bigtriangleup {ABC}$ 内一点.

(1)记 $\bigtriangleup {PBC}$ 、 $\bigtriangleup {PAC}$ 、 $\bigtriangleup {PAB}$ 的面积分别为 ${S}_{A}$ 、 ${S}_{B}$ 、 ${S}_{C}$ ，求证: ${S}_{A}\overrightarrow{PA} + {S}_{B}\overrightarrow{PB} + {S}_{C}\overrightarrow{PC} = \overrightarrow{0}$ ；

(2)若点 $P$ 为 $\bigtriangleup {ABC}$ 的垂心，且 $\overrightarrow{PA} + 2\overrightarrow{PB} + 3\overrightarrow{PC} = \overrightarrow{0}$ ，求 $\cos \angle {APB}$ .

\begin{center}
\includegraphics[width=0.2\textwidth]{images/bo_d7tk7oc91nqc73875fu0_150_1191_290_289_237_0.jpg}
\end{center}
 

(例 5 图)

\section*{四基达标精练}

\section*{一、填空题}

1. 如图，在 $\bigtriangleup  {ABC}$ 中， $D$ 为 ${AB}$ 的中点， $E$ 为 ${CD}$ 的中点，设 $\overrightarrow{AB} = \overrightarrow{a}$ ， $\overrightarrow{AC} = \overrightarrow{b}$ ，以向量 $\overrightarrow{a}$ 、 $\overrightarrow{b}$ 为基底，则向量 $\overrightarrow{AE} =$

\begin{center}
\includegraphics[width=0.2\textwidth]{images/bo_d7tk7oc91nqc73875fu0_150_185_1185_327_239_0.jpg}
\end{center}
 

(第 1 题图)

\begin{center}
\includegraphics[width=0.3\textwidth]{images/bo_d7tk7oc91nqc73875fu0_150_636_1223_352_197_0.jpg}
\end{center}
 

(第 2 题图)

\begin{center}
\includegraphics[width=0.2\textwidth]{images/bo_d7tk7oc91nqc73875fu0_150_1113_1232_316_182_0.jpg}
\end{center}
 

(第 3 题图)

2. 如图，在 $\bigtriangleup  {ABC}$ 中， ${BD} = {3DC}$ ，若 $\overrightarrow{AD} = \lambda \overrightarrow{AB} + \mu \overrightarrow{AC}$ ，则 $\frac{\lambda }{\mu } = {\_\_\_\_\_ }$ \_\_\_\_\_.

3. 如图，在平行四边形 ${ABCD}$ 中， ${AB} = 2,\cos \angle {BAD} = \frac{1}{2}$ ， $E$ 是边 ${BC}$ 的中点， $F$ 是 ${CD}$ 上靠近 $D$ 的三等分点,若 $\overrightarrow{AE} \cdot  \overrightarrow{BF} = 8$ ,则 $\left| \overrightarrow{AD}\right|  =$ \_\_\_\_\_.

4. 已知向量 $\overrightarrow{a}$ 、 $\overrightarrow{b}$ 的夹角为 ${120}^{ \circ  }$ ， $\left| \overrightarrow{a}\right|  = 4$ ， $\left| \overrightarrow{b}\right|  = 2$ ，则向量 $\overrightarrow{b}$ 在向量 $\overrightarrow{a}$ 方向上的数量投影为\_\_\_\_\_.

5. 已知 $\left| \overrightarrow{OM}\right|  = 2,\left| \overrightarrow{ON}\right|  = 2$ ,且 $\overrightarrow{OM}$ 、 $\overrightarrow{ON}$ 的夹角为 $\frac{\pi }{3}$ ,点 $P$ 在以 $O$ 为圆心的圆弧 $\overset{\text{ ⏜ }}{MN}$ 上运动,若 $\overrightarrow{OP} =  x\overrightarrow{OM} + y\overrightarrow{ON}\left( {x\text{ 、 }y \geq  0}\right)$ ,则 $x + y$ 的取值范围为\_\_\_\_\_.

6. 已知 $O$ 是 $\bigtriangleup {ABC}$ 内一点, $2\overrightarrow{OA} + 3\overrightarrow{OB} + m\overrightarrow{OC} = \overrightarrow{0}$ ,若 $\bigtriangleup {AOB}$ 与 $\bigtriangleup {ABC}$ 的面积之比为 $\frac{4}{7}$ ,则实数 $m$ 的值为\_\_\_\_\_.

7. 在 $\bigtriangleup {ABC}$ 中,内角 $A\text{ 、 }B\text{ 、 }C$ 的对边分别是 $a\text{ 、 }b\text{ 、 }c$ . 若 ${20a}\overrightarrow{BC} + {15b}\overrightarrow{CA} + {12c}\overrightarrow{AB} = \overrightarrow{0}$ ,则 $\bigtriangleup {ABC}$ 中最小角的余弦值等于\_\_\_\_\_.

8. 如图，四个边长为 1 的等边三角形有一条边在同一条直线上，边 ${B}_{4}{C}_{4}$ 上有 10 个不同的点 ${P}_{1}$ 、 ${P}_{2}$ 、...、 ${P}_{10}$ ，记 ${m}_{i} = \overrightarrow{A{B}_{2}} \cdot  \overrightarrow{A{P}_{i}}\left( {i = 1,2,\cdots ,{10}}\right)$ ，则 ${m}_{1} + {m}_{2} + \cdots  + {m}_{10} =$ \_\_\_\_\_.

\begin{center}
\includegraphics[width=0.5\textwidth]{images/bo_d7tk7oc91nqc73875fu0_151_220_281_603_191_0.jpg}
\end{center}
 

(第 8 题图)

\begin{center}
\includegraphics[width=0.3\textwidth]{images/bo_d7tk7oc91nqc73875fu0_151_950_246_372_221_0.jpg}
\end{center}
 

(第 13 题图)

9. 窗花是贴在窗纸或窗户玻璃上的剪纸，它是中国古老的传统民间艺术之一. 在 2022 年虎年新春来临之际，人们设计了一种由外围四个大小相等的半圆和中间正方形所构成的剪纸窗花(如图 1). 已知正方形 ${ABCD}$ 的边长为 2,中心为 $O$ ,四个半圆的圆心均为正方形 ${ABCD}$ 各边的中点 (如图 2),若 $P$ 在 $\overset{\text{ ⏜ }}{BC}$ 的中点, 则 $\left( {\overrightarrow{PA} + \overrightarrow{PB}}\right)  \cdot  \overrightarrow{PO} =$ \_\_\_\_\_.

\section*{二、选择题}

10. 已知平面向量 $\overrightarrow{a} = \left( {1,3}\right) ,\overrightarrow{b} = \left( {-2,4}\right)$ ,则 $\overrightarrow{a}$ 在 $\overrightarrow{b}$ 上的投影为 ( )

A. $\left( {1,3}\right)$ . B. $\left( {-1,2}\right)$ . C. $\sqrt{10}$ . D. $\sqrt{5}$ .

11. 已知单位向量 $\overrightarrow{a}$ 、 $\overrightarrow{b}$ ，若对任意实数 $x$ ， $\left| {x\overrightarrow{a} + \overrightarrow{b}}\right|  \geq  \frac{\sqrt{3}}{2}$ 恒成立，则向量 $\overrightarrow{a}$ 、 $\overrightarrow{b}$ 的夹角的取值范围为\_\_\_\_\_ \ding{115} \_\_\_\_\_.

A. $\left\lbrack  {\frac{\pi }{4},\frac{3\pi }{4}}\right\rbrack$ . B. $\left\lbrack  {\frac{\pi }{3},\frac{2\pi }{3}}\right\rbrack$ . C. $\left\lbrack  {\frac{\pi }{4},\frac{\pi }{2}}\right\rbrack$ . D. $\left\lbrack  {\frac{\pi }{3},\frac{\pi }{2}}\right\rbrack$ .

\begin{center}
\includegraphics[width=0.2\textwidth]{images/bo_d7tk7oc91nqc73875fu0_151_1190_1010_257_271_0.jpg}
\end{center}
 

(第 12 题图)

12. 如图，在 $\bigtriangleup  {ABC}$ 中， $M$ 为线段 ${BC}$ 的中点， $G$ 为线段 ${AM}$ 上一点，且 $\overrightarrow{AG} =  2\overrightarrow{GM}$ ,过点 $G$ 的直线分别交直线 ${AB}\text{ 、 }{AC}$ 于 $P\text{ 、 }Q$ 两点, $\overrightarrow{AB} = x\overrightarrow{AP}\left( {x > 0}\right) ,\overrightarrow{AC} =  y\overrightarrow{AQ}\left( {y > 0}\right)$ ,则 $\frac{1}{x} + \frac{1}{y + 1}$ 的最小值为 ( )

A. $\frac{3}{4}$ . B.1. C. $\frac{4}{3}$ . D. 4.

\section*{三、解答题}

13. 在直角梯形 ${ABCD}$ 中,已知 ${AB}//{CD},\angle {DAB} = {90}^{ \circ  },{AB} = {2AD} = {2CD} = 2$ ,点 $F$ 是 ${BC}$ 边上的中点,点 $E$ 是 ${CD}$ 边上一个动点.

(1)若 $\overrightarrow{DE} = \frac{1}{2}\overrightarrow{DC}$ ，求 $\overrightarrow{AC} \cdot  \overrightarrow{EF}$ 的值；

(2)求 $\overrightarrow{EA}$ 、 $\overrightarrow{EF}$ 的取值范围.

\begin{center}
\includegraphics[width=0.2\textwidth]{images/bo_d7tk7oc91nqc73875fu0_151_1138_1529_315_205_0.jpg}
\end{center}
 

(第 13 题图)

14. 如图,在平面四边形 ${ABCD}$ 中, ${BC}//{AD},{AB} = {BC} = 2,{AD} = 4,\angle {BAD} = {120}^{ \circ  },E\text{ 、 }F$ 分别是 ${AD}$ 、 ${DC}$ 的中点, $G$ 为线段 ${BC}$ 上一点,且 $\overrightarrow{BG} = \lambda \overrightarrow{BC}$ . 设 $\overrightarrow{AB} = \overrightarrow{a},\overrightarrow{AD} = \overrightarrow{b}$ .

(1)若 $\lambda  = \frac{1}{3}$ ，以 $\overrightarrow{a}$ 、 $\overrightarrow{b}$ 为基底表示向量 $\overrightarrow{AF}$ 与 $\overrightarrow{EG}$ ；

( 2 )若 $\lambda  \in  \left( {0,1}\right)$ ，求 $\overrightarrow{AF}$ 。 $\overrightarrow{EG}$ 的取值范围.

\begin{center}
\includegraphics[width=0.3\textwidth]{images/bo_d7tk7oc91nqc73875fu0_152_1083_359_405_175_0.jpg}
\end{center}
 

(第 14 题图)

15. 已知 $\bigtriangleup  {ABC}$ 中，角 $A$ 、 $B$ 、 $C$ 所对的边分别是 $a$ 、 $b$ 、 $c$ ，向量 $\overrightarrow{m} = \left( {b + a + c\text{ ， }\sqrt{3}b}\right)$ ， $\overrightarrow{n} = \left( {\sqrt{3}c,b - a + c}\right)$ ， 且 $\overrightarrow{m}//\overrightarrow{n}$ .

(1)若 $\overrightarrow{BD} = 3\overrightarrow{DC},{AD} = 1$ ，求 $\bigtriangleup  {ABC}$ 面积的最大值；

(2)若 $\bigtriangleup  {ABC}$ 为锐角三角形,且 $a = 3$ ，求 $\bigtriangleup  {ABC}$ 周长的取值范围.

16. 如图，在 $\bigtriangleup  {ABC}$ 中，已知 ${AC} = 1,{AB} = 3,\angle {BAC} = {60}^{ \circ  }$ ，且 $\overrightarrow{PA} + \overrightarrow{PB} + \overrightarrow{PC} = \overrightarrow{0}$ .

(1)求 ${PA}$ 的长；

(2)求 $\cos \angle {APC}$ .

\begin{center}
\includegraphics[width=0.3\textwidth]{images/bo_d7tk7oc91nqc73875fu0_152_1111_1622_377_151_0.jpg}
\end{center}
 

(第 16 题图)

$\text{ § }{8.3}$ 平面向量及其应用二

\section*{典型例题精讲}

【例 1】已知向量 $\overrightarrow{OA} = \left( {0, - 6}\right) ,\overrightarrow{OB} = \left( {8,0}\right) ,\overrightarrow{OC} = \left( {x,y}\right)$ .

(1)若点 $A$ 、 $B$ 、 $C$ 三点共线，求实数 $x$ 、 $y$ 满足的关系；

(2)若 $x = 1$ ，且 $\angle {ACB}$ 为钝角，求实数 $y$ 的取值范围.

【例 2】如图,已知 $\bigtriangleup {ABC}$ 的外接圆 $O$ 的半径为 $4,\overrightarrow{AO} = 2\overrightarrow{AC} + \overrightarrow{AB}$ .

(1)求 $\bigtriangleup  {ABC}$ 中 ${AC}$ 边的长；

(2)求 $\overrightarrow{OC} \cdot  \overrightarrow{AB}$ .

\begin{center}
\includegraphics[width=0.2\textwidth]{images/bo_d7tk7oc91nqc73875fu0_153_1158_1256_296_295_0.jpg}
\end{center}
 

(例 2 图)

【例 3】如图,一条河两岸平行,河的宽度 ${AC} = \sqrt{3}{{km}}$ ,一艘船从河边的点 $A$ 出发到达对岸的点 $B$ ,船只在河内行驶的路程 ${AB} = 2{{km}}$ ，行驶时间为 0.2 h. 已知船在静水中的速度 ${v}_{1}$ 的大小为 $\left| {v}_{1}\right|$ ，水流的速度 ${v}_{2}$ 的大小为 $\left| {v}_{2}\right|  = 2{{km}}/{h}$ . 求:

(1) $\left| {v}_{1}\right|$ ；

(2)船在静水中速度 ${v}_{1}$ 与水流速度 ${v}_{2}$ 夹角的余弦值.

\begin{center}
\includegraphics[width=0.2\textwidth]{images/bo_d7tk7oc91nqc73875fu0_154_1163_365_328_207_0.jpg}
\end{center}
 

(例 3 图)

【例 4】在平面直角坐标系 ${xOy}$ 中,已知点 $A\left( {-1,0}\right) ,B\left( {0,\sqrt{3}}\right) ,C\left( {\cos \theta ,\sin \theta }\right)$ ,其中 $\theta  \in  \left\lbrack  {0,\frac{\pi }{2}}\right\rbrack$ . (1)求 $\overrightarrow{AC} \cdot  \overrightarrow{BC}$ 的最大值；

(2)是否存在 $\theta  \in  \left\lbrack  {0,\frac{\pi }{2}}\right\rbrack$ ，使得 $\bigtriangleup  {ABC}$ 为钝角三角形？若存在，求出 $\theta$ 的取值范围；若不存在，说明理由.

【例 5】已知在平面直角坐标系中， $O$ 为坐标原点，定义非零向量 $\overrightarrow{OM} = \left( {a,b}\right)$ 的“相伴函数”为 $y =  a\sin x + b\cos x\left( {x \in  \mathbf{R}}\right)$ ,向量 $\overrightarrow{OM} = \left( {a,b}\right)$ 称为函数 $y = a\sin x + b\cos x\left( {x \in  \mathbf{R}}\right)$ 的“相伴向量”; 记平面内所有向量的“相伴函数”构成的集合为 $S$ .

(1)已知 $\alpha  \in  \mathbf{R}$ ， $h\left( x\right)  = \cos \left( {x + \alpha }\right)  + 2\cos x$ ，若函数 $h\left( x\right)$ 为集合 $S$ 中的元素，求其 “相伴向量” 的模的取值范围;

(2)当向量 $\overrightarrow{OM} = \left( {\sqrt{3},1}\right)$ 时，“相伴函数”为 $f\left( x\right)$ ，若 $x \in  \left\lbrack  {0,\frac{11\pi }{6}}\right\rbrack$ ，方程 ${f}^{2}\left( x\right)  + \left( {2 - a}\right) f\left( x\right)  + a - 3 = 0$ 存在 4 个不相等的实数根,求实数 $a$ 的取值范围.

\section*{四基达标精练}

\section*{一、填空题}

1. 向量 $\overrightarrow{a} = \left( {3,4}\right)$ 在向量 $\overrightarrow{b} = \left( {1,0}\right)$ 方向上的数量投影为\_\_\_\_\_.

2. 已知 $G$ 为 $\bigtriangleup  {ABC}$ 的重心,用向量 $\overrightarrow{AB}$ 、 $\overrightarrow{BC}$ 表示向量 $\overrightarrow{AG} =$ \_\_\_\_\_.

3. 已知平面单位向量 $\overrightarrow{a}\text{ 、 }\overrightarrow{b}\text{ 、 }\overrightarrow{c}$ 满足 $\langle \overrightarrow{a},\overrightarrow{b}\rangle  = \langle \overrightarrow{b},\overrightarrow{c}\rangle  = \langle \overrightarrow{c},\overrightarrow{a}\rangle  = \frac{2\pi }{3}$ ，则 $\left| {3\overrightarrow{a} + 2\overrightarrow{b} + \overrightarrow{c}}\right|  =$ \_\_\_\_\_.

4. 如图所示，两块斜边长均等于 $\sqrt{2}$ 的直角三角板拼在一起，则 $\overrightarrow{OD}$ . $\overrightarrow{AB} =$ \_\_\_\_\_.

\begin{center}
\includegraphics[width=0.2\textwidth]{images/bo_d7tk7oc91nqc73875fu0_155_259_604_343_177_0.jpg}
\end{center}
 

(第 4 题图)

\begin{center}
\includegraphics[width=0.3\textwidth]{images/bo_d7tk7oc91nqc73875fu0_155_815_549_447_226_0.jpg}
\end{center}
 

(第 8 题图)

5. 已知非零向量 $\overrightarrow{a}$ 、 $\overrightarrow{b}$ 、 $\overrightarrow{c}$ 两两不平行，且 $\overrightarrow{a}//\left( {\overrightarrow{b} + \overrightarrow{c}}\right)$ ，方 $//\left( {\overrightarrow{a} + \overrightarrow{c}}\right)$ ，设 $\overrightarrow{c} = {x\overrightarrow{a}} + {y\overrightarrow{b}}\left( {x\text{ 、 }y \in  \mathbf{R}}\right)$ ，则 $x + {2y}  =$ \_\_\_\_\_.

6. 一条东西方向的河流两岸平行，河宽 ${250}\sqrt{3}{m}$ ，河水的速度为向东 $2{{km}}/{h}$ . 一艘小货船准备从河南岸的码头 $A$ 处出发,航行到位于河对岸 $B$ ( ${AB}$ 与河的方向垂直)的正西方向并且与 $B$ 相距 ${250}{m}$ 的码头 $C$ 处卸货. 若流水的速度与小货船航行的速度的合速度为 $6{{km}}/{h}$ ，则当小货船的航程最短时，小货船航行的速度大小是\_\_\_\_\_km/h.

7. 已知点 $A$ 、 $B$ 在圆 $O : {x}^{2} + {y}^{2} = {16}$ 上,且 $\left| {AB}\right|  = 4,P$ 为圆 $O$ 上任意一点,则 $\overrightarrow{AB} \cdot  \overrightarrow{BP}$ 的最小值为\_\_\_\_\_.

8. 青花瓷，又称白地青花瓷，常简称青花，是中国瓷器的知名品种之一. 如图 1，这是一个青花瓷圆盘. 该圆盘中的两个圆的圆心重合,如图 2,其中大圆半径 $R = 3$ ,小圆半径 $r = 2$ ,点 $P$ 在大圆上,过点 $P$ 作小圆的切线，切点分别是 $E$ 、 $F$ ，则 $\overrightarrow{PE}$ . $\overrightarrow{PF} =$ \_\_\_\_\_.

9. 已知平面向量 $\overrightarrow{a}$ 、 $\overrightarrow{b}$ 、 $\overrightarrow{c}$ 满足 $\left| \overrightarrow{a}\right|  = \left| \overrightarrow{b}\right|  = \overrightarrow{a} \cdot  \overrightarrow{b} = 2$ ，且 $\left( {\overrightarrow{b} - \overrightarrow{c}}\right)  \cdot  \left( {2\overrightarrow{b} - \overrightarrow{c}}\right)  = 0$ ，则 $\left| {\overrightarrow{a} - 2\overrightarrow{c}}\right|$ 的最大值为\_\_\_\_\_.

\section*{二、选择题}

10. 设非零向量 $\overrightarrow{a}$ 、 $\overrightarrow{b}$ 的夹角为 $\theta$ ，定义运算 $\overrightarrow{a} * \overrightarrow{b} = \left| \overrightarrow{a}\right|  \cdot  \left| \overrightarrow{b}\right|  \cdot  \sin \theta$ . 下列叙述正确的是 ( )

A. 若 $\overrightarrow{a} * \overrightarrow{b} = 0$ ，则 $\overrightarrow{a} \bot  \overrightarrow{b}$ .

B. 若 $\left| \overrightarrow{a}\right|  \cdot  \left| \overrightarrow{b}\right|  = 1$ ，则 ${\left( \overrightarrow{a} * \overrightarrow{b}\right) }_{\min } =  - 1$ .

C. 设在 $\bigtriangleup {ABC}$ 中, $\overrightarrow{AB} = \overrightarrow{a},\overrightarrow{AC} = \overrightarrow{b}$ ,则 $2{S}_{\bigtriangleup {ABC}} = \overrightarrow{a} * \overrightarrow{b}$ .

D. $\overrightarrow{a} * \left( {\overrightarrow{b} + \overrightarrow{c}}\right)  = \overrightarrow{a} * \overrightarrow{b} + \overrightarrow{a} * \overrightarrow{c}\left( \overrightarrow{c}\right)$ 为任意非零向量).

11. 已知点 ${A}_{1},{A}_{2},\cdots ,{A}_{n},\cdots$ 和数列 $\left\{  {a}_{n}\right\}$ 、 $\left\{  {b}_{n}\right\}$ 满足 $\overrightarrow{{A}_{n}{A}_{n + 1}} = \left( {\cos \frac{2n\pi }{3},\sin \frac{2n\pi }{3}}\right) \left( {n \in  {\mathbf{N}}^{ * }}\right)$ ,若 ${a}_{1} = 1$ , ${S}_{n}\text{ 、 }{T}_{n}$ 分别为数列 $\left\{  {a}_{n}\right\}  \text{ 、 }\left\{  {b}_{n}\right\}$ 的前 $n$ 项和,则 ${S}_{60} + 2{T}_{60}$ 等于 ( )

A. -20 . B. ${24}\sqrt{3}$ . C. ${48}\sqrt{3} - {20}$ . D.O.

12. 如图， ${Ox}$ 、 ${Oy}$ 是平面内相交成 $\theta$ ( $\theta  \neq  \frac{\pi }{2}$ )角的两条数轴， $\overrightarrow{{e}_{1}}$ 、 $\overrightarrow{{e}_{2}}$ 分别是与 $x\text{ 、 }y$ 轴正方向同向的单位向量，则称 ${xOy}$ 为 $\theta$ 斜坐标系，若 $\overrightarrow{OM} = x\overrightarrow{{e}_{1}} + y\overrightarrow{{e}_{2}}$ ，则把有序数对 $\left( {x,y}\right)$ 叫做向量 $\overrightarrow{OM}$ 的斜坐标,记为 $\overrightarrow{OM} = \left( {x,y}\right)$ . 在 $\frac{\pi }{4}$ 斜坐标系中, $\overrightarrow{a} = \left( {\frac{1}{2},\frac{\sqrt{3}}{2}}\right) ,\overrightarrow{b} = \left( {\sqrt{3}, - 1}\right)$ . 则下列结论中,正确的是 ( )

\begin{center}
\includegraphics[width=0.2\textwidth]{images/bo_d7tk7oc91nqc73875fu0_156_1177_141_310_243_0.jpg}
\end{center}
 

(第 12 题图)

A. $\overrightarrow{a} - \overrightarrow{b} = \left( {\frac{1}{2} - \sqrt{3},\frac{\sqrt{3}}{2} + 1}\right)$ . B. $\left| \overrightarrow{a}\right|  = 1$ .

C. $\overrightarrow{a} \bot  \overrightarrow{b}$ . D. $\overrightarrow{b}$ 在 $\overrightarrow{a}$ 上的投影向量为 $\left( {\frac{2\sqrt{2} + \sqrt{3}}{5},\frac{2\sqrt{6} + 2}{5}}\right)$ .

\section*{三、解答题}

13. 在 $\bigtriangleup {ABC}$ 中,已知 $\overrightarrow{AB} \bot  \overrightarrow{AC},M$ 是 ${BC}$ 的中点.

(1)若 $\left| \overrightarrow{AB}\right|  = \left| \overrightarrow{AC}\right|$ ，求向量 $\overrightarrow{AB} + 2\overrightarrow{AC}$ 与向量 $2\overrightarrow{AB} + \overrightarrow{AC}$ 的夹角的余弦值；

(2)若 $O$ 是线段 ${AM}$ 上任意一点，且 $\left| \overrightarrow{AB}\right|  = \left| \overrightarrow{AC}\right|  = \sqrt{2}$ ，求 $\overrightarrow{OA} \cdot  \overrightarrow{OB} + \overrightarrow{OC} \cdot  \overrightarrow{OA}$ 的最小值.

14. 已知 $A$ 、 $B$ 、 $C$ 为 $\bigtriangleup {ABC}$ 的三个内角,向量 $\overrightarrow{m} = \left( {2 - 2\sin A,\sin A + \cos A}\right)$ 平行于 $\overrightarrow{n} = (\sin A - \cos A$ ， $1 + \sin A)$ ,且 $\overrightarrow{AB} \cdot  \overrightarrow{AC} > 0$ .

(1)求角 $A$ ；

(2)求函数 $y = 2{\sin }^{2}\frac{B}{2} + \cos \frac{C - B}{2}$ 的值域.


15. 假设在一个以 ${m}$ 为单位的空间直角坐标系 $O - {xyz}$ 中,平面 ${xOy}$ 内有一跟踪和控制飞行机器人 $T$ 的控制台 $A,A$ 的位置为 $\left( {{170},{200},0}\right)$ ,上午 10 时 07 分测得飞行机器人 $T$ 在 $P\left( {{150},{80},{120}}\right)$ 处,并对飞行机器人 $T$ 发出指令:以速度 ${v}_{1} = {13}{m}/{s}$ 沿单位向量 ${\overrightarrow{d}}_{1} = \left( {\frac{3}{13},\frac{12}{13}, - \frac{4}{13}}\right)$ 作匀速直线飞行 (飞行中无障碍物)， ${10}{s}$ 后到达点 $Q$ ，再发出指令让机器人在点 $Q$ 原地盘旋 $2{s}$ ，在原地盘旋过程中逐步减速并降速到 $8{m}/{s}$ ，然后保持 $8{m}/{s}$ ，再沿单位向量 ${\overrightarrow{d}}_{2} = \left( {\frac{1}{2}, - \frac{\sqrt{2}}{2}, - \frac{1}{2}}\right)$ 作匀速直线飞行 (飞行中无障碍物),当飞行机器人 $T$ 最终落在平面 ${xOy}$ 内发出指令让它停止运动,机器人 $T$ 近似看成一个点.

(1)求从 $P$ 点开始出发 ${20}{s}$ 后飞行机器人 $T$ 的位置；

(2)求在整个飞行过程中飞行机器人 $T$ 与控制台 $A$ 的最近距离(精确到米).

\begin{center}
\includegraphics[width=0.2\textwidth]{images/bo_d7tk7oc91nqc73875fu0_157_1184_572_259_260_0.jpg}
\end{center}
 

(第 15 题图)

16. 如图所示，在 $\bigtriangleup  {ABO}$ 中， $\overrightarrow{OC} = \frac{1}{3}\overrightarrow{OA}$ ， $\overrightarrow{OD} = \frac{1}{2}\overrightarrow{OB}$ ， ${AD}$ 与 ${BC}$ 相交于点 $M$ . 设 $\overrightarrow{OA} = \overrightarrow{a}$ ， $\overrightarrow{OB} = \overrightarrow{b}$ .

(1)试用向量 $\overrightarrow{a}$ 、 $\overrightarrow{b}$ 表示 $\overrightarrow{OM}$ ；

(2)在线段 ${AC}$ 上取点 $E$ ，在线段 ${BD}$ 上取点 $F$ ，使 ${EF}$ ， 过点 $M$ . 设 $\overrightarrow{OE} = \lambda \overrightarrow{OA},\overrightarrow{OF} = \mu \overrightarrow{OB}$ ，其中 $\lambda$ 、 $\mu  \in  \mathbf{R}$ .

当 ${EF}$ 与 ${AD}$ 重合时, $\lambda  = 1,\mu  = \frac{1}{2}$ ,此时 $\frac{1}{\lambda } + \frac{2}{\mu } = 5$ ;

当 ${EF}$ 与 ${BC}$ 重合时, $\lambda  = \frac{1}{3},\mu  = 1$ ,此时 $\frac{1}{\lambda } + \frac{2}{\mu } = 5$ ;

能否由此得出一般结论: 不论 $E\text{ 、 }F$ 在线段 ${AC}\text{ 、 }{BD}$ 上如何变动,等式 $\frac{1}{\lambda } + \frac{2}{\mu } = 5$ 恒成立,请说明理由.

\begin{center}
\includegraphics[width=0.3\textwidth]{images/bo_d7tk7oc91nqc73875fu0_157_1087_1563_363_261_0.jpg}
\end{center}
 

(第 16 题图)

\section*{$\text{ § }{8.4}$ 复数及其运算}

(本小节若无特别说明, i 表示虚数单位)

\section*{典型例题精讲}

【例 1】已知 $m \in  \mathbf{R}$ ,复数 $z = \frac{m\left( {m + 2}\right) }{m - 1} + \left( {{m}^{2} + {2m} - 3}\right) {i}$ ,当 $m$ 为何值时,

(1) $z \in  \mathbf{R}$ ；(2) $z$ 是虚数；(3) $z$ 是纯虚数；(4) $\bar{z} = \frac{1}{2} + {4i}$ .

【例 2】在复数范围内,下列命题中为真命题的序号是\_\_\_\_\_.

 \ding{172}  ${\left| \bar{z}\right| }^{2} = \left| {z}^{2}\right|$ ； \ding{173}  若 ${z}_{1} - {z}_{2} > 0$ ，则 ${z}_{1} > {z}_{2}$ ； \ding{174}  若 ${\left( {z}_{1} - {z}_{2}\right) }^{2} + {\left( {z}_{2} - {z}_{3}\right) }^{2} = 0$ ，则 ${z}_{1} = {z}_{2} = {z}_{3}$ ；

 \ding{175}  $\left| {{z}_{1} - {z}_{2}}\right|  = \sqrt{{\left( {z}_{1} + {z}_{2}\right) }^{2} - 4{z}_{1}{z}_{2}}$ ;  \circled{5}  ${z}_{1}^{2} = {z}_{2}^{2}$ ，则 ${z}_{1} \cdot  {\bar{z}}_{1} = {z}_{2} \cdot  {\bar{z}}_{2}$ ;  \circled{6}  $2{z}_{1}{z}_{2} \leq  {z}_{1}^{2} + {z}_{2}^{2}$ ；

 \circled{7} 两个共轭复数的差是纯虚数； \circled{8} 若 $\left| {z + {i}}\right|  = \left| {z - {i}}\right|$ ，则 $z$ 必为实数.

【例 3】设复数 $z = a + b{i}\left( {a\text{ 、 }b \in  \mathbf{R}}\right) ,{z}_{1} = z + k{i},{z}_{2} = \bar{z} \cdot  k{i}$ (其中 $k \in  \mathbf{R}$ ).

(1)设 $a = b = \frac{1}{2}$ ，若 $\left| {z}_{1}\right|  = \left| {z}_{2}\right|$ ，求出实数 $k$ 的值；

(2)若复数 $z$ 满足条件:存在实数 $k$ ，使得 ${z}_{1}$ 与 ${z}_{2}$ 是某个实系数一元二次方程的两个虚数根，求符合条件的复数 $z$ 的模的取值范围.

【例 4】设 $\alpha \text{ 、 }\beta$ 为方程 ${x}^{2} + {2x} + t = 0\left( {t \in  \mathbf{R}}\right)$ 的两个根, $f\left( t\right)  = \left| \alpha \right|  + \left| \beta \right|$ .

(1)求 $f\left( t\right)$ 的解析式；

(2)证明关于 $t$ 的方程 $f\left( t\right)  = m$ ，当 $m > 2$ 时恰有两个不等的根，且两根之和为定值.

【例 5】已知 $z = m + 3 + 3\sqrt{3}{i}$ ,其中 $m \in  \mathbf{C}$ ,且 $\frac{m + 3}{m - 3}$ 为纯虚数.

(1)求 $m$ 对应点的轨迹；(2)求 $\left| z\right|$ 的最大值、最小值.

\section*{四基达标精练}

\section*{一、填空题}

1. 已知 $\frac{z - 3}{1 + {2z}} = {i}$ ,则 $z =$ \_\_\_\_\_.

2. 若 $1 + \sqrt{2}{i}$ 是关于 $x$ 的实系数方程 ${x}^{2} + {bx} + c = 0$ 的一个虚根,则 $c =$ \_\_\_\_\_.

3. 已知复数 $a + b{i} = \frac{3 + {i}}{1 - {i}}\left( {a\text{ 、 }b \in  \mathbf{R}}\right)$ ,则 $\left| {a + b{i}}\right|  =$ \_\_\_\_\_.

4. 在复数范围内分解因式: $4{x}^{3} - 8{x}^{2} + {10x} =$ \_\_\_\_\_.

5. 若复数 $z$ 满足 $\left| {z - 1 + \sqrt{3}{i}}\right|  = 3$ ，则 $\left| z\right|$ 的最大值为\_\_\_\_\_.

6. 已知关于 $x$ 的实系数方程 ${x}^{2} + \left( {1 - k}\right) x + {k}^{2} - 1 = 0$ 有两个虚根 ${x}_{1}\text{ 、 }{x}_{2}$ ,且 $\left| {x}_{1}\right|  + \left| {x}_{2}\right|  = 2$ ,则满足条件的实数 $k =$ \_\_\_\_\_.

7. 复数 ${z}_{1}$ 和 ${z}_{2}$ 满足 $\left| {{z}_{1} + 1}\right|  = \left| {{z}_{1} - 1}\right|$ ，且 $\left| {{z}_{2} - 2 - {3{i}}}\right|  = 1$ ，则 $\left| {{z}_{1} - {z}_{2}}\right|$ 的最小值是\_\_\_\_\_.

8. 已知关于 $x$ 的方程 ${x}^{2} + {zx} + 4 + {3{i}} = 0$ 有实数根，则复数 $z$ 的模的最小值为\_\_\_\_\_.

9. 已知 $k + 2$ 个两两互不相等的复数 ${z}_{1}\text{ 、 }{z}_{2}\text{ 、 }\cdots \text{ 、 }{z}_{k}\text{ 、 }{\omega }_{1}\text{ 、 }{\omega }_{2}$ ,满足,且 $\left| {{\omega }_{m} - {z}_{n}}\right|  \in  \{ 1,3\}$ (其中 $m = 1,2$ , $n = 1,2,\cdots ,k)$ ，则 $k$ 的最大值为\_\_\_\_\_.

\section*{二、选择题}

10. 若复数 $z = \left( {m + 2}\right)  + \left( {m - 4}\right) {i}$ 是虚数,则实数 $m$ 取值的集合是 ( )

A. $\{ m \mid  m > 4\}$ . B. $\{ m \mid  m < 4\}$ . C. $\{ m \mid  m \neq  4,m \in  \mathbf{R}\}$ . D.R.

D. $- 3 + {i}$ .

11. 在复平面内,平行四边形 ${ABCD}$ 的三个顶点, $A\text{ 、 }B\text{ 、 }C$ 对应的复数分别为 $- 1 + 2{i}\text{ 、 }3 - {i}\text{ 、 }1 + 2{i}$ ,则点 $D$ 对应的复数为 ( )

A. $- 3 + 5{i}$ . B. $1 - {i}$ . C. $1 + 3{i}$ .


12. 数学家欧拉发现了复指数函数和三角函数的关系,并给出公式 ${{e}}^{{i}x} = \cos x + {i}\sin x$ ,(其中 ${e}$ 是自然对数的底数, $x \in  \mathbf{R}$ ),这个公式在复变论中有非常重要的地位,被称为 “数学中的天桥”,根据此公式,有下列四个结论, 其中正确的是 ( )

A. ${{e}}^{{i}\pi } - 1 = 0$ . B. $2\cos x = {{e}}^{-{i}x} + {{e}}^{{i}x}$ . C. $2\sin x = {{e}}^{{i}x} - {{e}}^{-{i}x}$ . D. ${\left( \frac{\sqrt{2}}{2} + \frac{\sqrt{2}}{2}{i}\right) }^{2022} =  - 1$ .

\section*{三、解答题}

13. 已知 ${i}$ 为虚数单位.

(1)计算: $\frac{-2\sqrt{3} + {i}}{1 + 2\sqrt{3}{i}} - {\left( \frac{1 + {i}}{\sqrt{2}}\right) }^{2022}$ ;

(2)若 ${\left| z\right| }^{2} + \left( {z + \bar{z}}\right) {i} = 1 - {i}$ ，求复数 $z$ .

14. 已知复数 $z = \left( {{m}^{2} + m - 6}\right)  + \left( {{m}^{2} - {3m} + 2}\right) {i}$ ，其中 $m \in  \mathbf{R}$ .

(1)当 $m$ 取何值时， $z$ 为纯虚数；

(2)如果复数 $z$ 在复平面上对应的点位于第二象限，求实数 $m$ 的取值范围.

15. 已知复数 ${z}_{1} = a + 2 + \left( {{a}^{2} - 3}\right) {i},{z}_{2} = 2 - \left( {{3a} + 1}\right) {i}\left( {a \in  \mathbf{R}}\right)$ .

(1)若复数 ${z}_{1} - \overline{{z}_{2}}$ 在复平面上对应点落在第一象限，求实数 $a$ 的取值范围；

(2)若虚数 ${z}_{1}$ 是实系数一元二次方程 ${x}^{2} - {6x} + m = 0$ 的根，求实数 $m$ 的值.

16. 对于一组复数 ${z}_{1}\text{ 、 }{z}_{2}\text{ 、 }\cdots \text{ 、 }{z}_{n}\left( {n \in  \mathbf{N},n \geq  3}\right)$ ,令 ${S}_{n} = {z}_{1} + {z}_{2} + {z}_{3} + \cdots  + {z}_{n}$ ,如果存在 ${z}_{p}(p \in  \{ 1,2$ , $3,\cdots ,n\} )$ ,使得 $\left| {z}_{p}\right|  \geq  \left| {{S}_{n} - {z}_{p}}\right|$ ,那么称 ${z}_{p}$ 是该复数组的 “ $M$ 复数”.

(1)已知 ${z}_{1} = {i},{z}_{2} = 1 + {i}$ ，是否存在复数 ${z}_{3}$ 使得 ${z}_{1}$ 、 ${z}_{2}$ 、 ${z}_{3}$ 均是复数组 ${z}_{1}$ 、 ${z}_{2}$ 、 ${z}_{3}$ 的“ $M$ 复数”? 若存在， 求出所有的 ${z}_{3}$ ; 若不存在,说明理由;

(2)若 ${z}_{n} = {\left( \frac{5}{9}\right) }^{n - 1} + {i} \cdot  {\left( -1\right) }^{n}\left( {n \in  {\mathbf{N}}^{ * }}\right)$ ，复数组 ${z}_{1}$ 、 ${z}_{2}$ 、 $\cdots$ 、 ${z}_{n}\left( {n \in  \mathbf{N},n \geq  3}\right)$ 是否存在 “ $M$ 复数”? 给出你的结论, 并说明理由.

\section*{第 9 章 立体几何(必修 3 \& 选择性必修 1)}

\section*{知识梳理*}

第一部分:空间直线和平面

(以下 $a\text{ 、 }b\text{ 、 }l\text{ 、 }\cdots$ 表示直线， $\alpha \text{ 、 }\beta \text{ 、 }\gamma \text{ 、 }\cdots$ 表示平面)

1. 平面及其基本性质:

公理一:如果一条直线有两个点在一个平面内，则这条直线上任意一点都在该平面内.

作用:判定一条直线在平面内的依据；结论:平面过直线的直线的取

公理二:经过不在同一条直线上的三点有且仅有一个平面.

推论 \ding{172} :经过一条直线和这条直线外一点有且仅有一个平面.

推论 \ding{173} :经过两条相交直线有且仅有一个平面.

推论 \ding{174} :经过两条平行直线有且仅有一个平面.

作用:确定平面的依据.

公理三:如果两个平面有一个公共点，则这两个平面有且仅有一条过这点的公共直线。

作用:判定两个平面相交的依据.

2. 空间直线和平面:

\begin{center}
\begin{tabular}{|c|c|c|c|c|c|c|c|c|}
\hline
\phantom{X} & \multicolumn{3}{|c|}{空间两直线} & \multicolumn{3}{|c|}{直线与平面} & \multicolumn{2}{|c|}{平面与平面} \\
\cline{1-9}
位置关系 & 相交 & 平行 & 异面 & 在平面内 & 平行 & 相交 & 平行 & 相交 \\
\cline{1-9}
记号 & $a \cap  b = P$ & $a//b$ & $a\text{ 、 }b$ 异面 & $a \subset  \alpha$ & $a//\alpha$ & $a \cap  \alpha  = P$ & $\alpha //\beta$ & $\alpha  \cap  \beta  = l$ \\
\cline{1-9}
公共点 & 唯一 & \multicolumn{2}{|c|}{无} & 无穷多个 & 无 & 唯一 & 无 & 无穷多个 \\
\cline{1-9}
\hline
\end{tabular}
\end{center}

3. 直线与直线的位置关系:

公理四:平行于同一直线的两直线平行. $\left( {a//b,b//c \Rightarrow  a//c}\right)$

> 异面直线判定定理:过平面外一点与平面上一点的直线，和此平面上不经过该点的任何一条直线都是异面直线.

4. 直线与平面的位置关系:

$>$ 线面平行:

(1)直线与平面平行的判定定理:如果不在平面上的一条直线与这个平面上的一条直线平行,那么该直线与这个平面平行. $\left( {a \text{ ⊄ } \alpha ,b \subset  \alpha ,a//b \Rightarrow  a//\alpha }\right)$

(2)直线与平面平行的性质定理:如果一条直线与一个平面平行，过这条直线的一个平面与此平面相交，那么其交线必与该直线平行. $\left( {a//\alpha ,a \subset  \beta ,\alpha  \cap  \beta  = b \Rightarrow  a//b}\right)$

$>$ 线面垂直:

(1)定义:如果一条直线与平面上的任意一条直线都垂直，就说这条直线与这个平面互相垂直.

(2)判定定理:如果一条直线与一个平面上的两条相交直线都垂直，那么此直线与该平面垂直. $\left( {l \bot  a}\right.$ , $l \bot  b,a \cap  b = P,a \subset  \alpha ,b \subset  \alpha  \Rightarrow  l \bot  \alpha )$

(3)性质定理:垂直于同一个平面的两条直线互相平行. $\left( {a \bot  a,b \bot  \alpha  \Rightarrow  a//b}\right)$

推论 \ding{172} :过一点有且只有一个平面与给定的直线垂直.

\HRule

*注:本章节所属知识模块是“图形与几何”.

\HRule

推论 \ding{173} :过一点有且只有一条直线与给定的平面垂直.

(4)三垂线定理:平面上的一条直线和这个平面的一条斜线垂直的充要条件是它和这条斜线在平面上的投影垂直.

5. 平面与平面的位置关系:

> 面面平行:

(1)定义:如果两平面没有公共点，则称这两平面平行.

(2)判定定理:如果一个平面上的两条相交直线与另一个平面平行，那么这两个平面平行. $\left( {a \subset  \alpha ,b \subset  \alpha }\right.$ , $a \cap  b = P,a//\beta ,b//\beta  \Rightarrow  \alpha //\beta )$

(3)性质定理:两平行平面和第三平面相交，则交线平行. $\left( {\alpha //\beta ,\gamma  \cap  \alpha  = a,\gamma  \cap  \beta  = b \Rightarrow  a//b}\right)$

> 面面垂直:

(1)二面角的平面角:从一条直线出发的两个半平面所组成的图形叫做二面角，一个二面角的大小等于它的平面角的大小.

面面垂直的定义:两平面所成的二面角是直角，则这两个平面垂直.

(2)判定定理:若一个平面经过另一个平面的一条垂线，则两平面垂直. $\left( {a \subset  \alpha ,a\bot \beta  \Rightarrow  \alpha \bot \beta }\right)$

(3)性质定理:如果两个平面垂直，那么其中一个平面上垂直于两平面交线的直线与另一个平面垂直 $\left( {\alpha  \bot  \beta ,\alpha  \cap  \beta  = l,a \subset  \alpha ,a \bot  l \Rightarrow  a \bot  \beta }\right)$

 \ding{72} 空间位置关系的证明可以借助空间向量完成，具体如下:

 \ding{172} 设直线 ${l}_{1}$ 和 ${l}_{2}$ 的方向向量分别为 ${\overrightarrow{d}}_{1}$ 和 ${\overrightarrow{d}}_{2}$ ，若 ${\overrightarrow{d}}_{1}//{\overrightarrow{d}}_{2}$ ，则 ${l}_{1}//{l}_{2}$ ；若 ${\overrightarrow{d}}_{1}\bot {\overrightarrow{d}}_{2}$ ，则 ${l}_{1}\bot {l}_{2}$ ；

 \ding{173} 设直线 $l$ 的方向向量为 $\overrightarrow{d}$ ，平面 $\alpha$ 的法向量为 $\overrightarrow{n}$ ，则 $\overrightarrow{d}//\overrightarrow{n} \Leftrightarrow  l\bot \alpha$ ；

 \ding{174} 设平面 $\alpha$ 和 $\beta$ 的法向量分别为 ${\overrightarrow{n}}_{1}$ 和 ${\overrightarrow{n}}_{2}$ ，若 ${\overrightarrow{n}}_{1}//{\overrightarrow{n}}_{2}$ ，则 $\alpha //\beta$ ；若 ${\overrightarrow{n}}_{1}\bot {\overrightarrow{n}}_{2}$ ，则 $\alpha \bot \beta$ .

6. 空间中的角(关键是构造三角形或利用空间向量的数量积):

异面直线所成的角: 转化为相交直线所成的锐角或直角 (构造三角形)；范围 $\left( {0,\frac{\pi }{2}}\right\rbrack$ .

 \ding{72} 异面直线所成角的向量求法:设异面直线 ${l}_{1}$ 和 ${l}_{2}$ 的方向向量分别为 ${\overrightarrow{d}}_{1}$ 和 ${\overrightarrow{d}}_{2},{\overrightarrow{d}}_{1}$ 和 ${\overrightarrow{d}}_{2}$ 的夹角为 $\varphi$ , 异面直线 ${l}_{1}$ 和 ${l}_{2}$ 所成的角为 $\theta$ ,则 $\cos \theta  = \left| {\cos \varphi }\right|  = \frac{\left| {\overrightarrow{d}}_{1} \cdot  {\overrightarrow{d}}_{2}\right| }{\left| {\overrightarrow{d}}_{1}\right|  \cdot  \left| {\overrightarrow{d}}_{2}\right| }$ .

线面角:斜线与它在平面内的射影所成的锐角，直线垂直于平面时为直角，在平面内或平行于平面时为零角,取值范围是 $\left\lbrack  {0,\frac{\pi }{2}}\right\rbrack$ . (解决的关键是找到直线在平面内的射影)

 \ding{72}  线面角的向量求法:设直线 $l$ 的方向向量为 $\overrightarrow{d}$ ，平面 $\alpha$ 的法向量为 $\overrightarrow{n},\overrightarrow{d}$ 和 $\overrightarrow{n}$ 的夹角为 $\varphi$ ，直线 $l$ 与平面 $\alpha$ 所成的角为 $\theta ,\sin \theta  = \left| {\cos \varphi }\right|  = \frac{\left| \overrightarrow{d} \cdot  \overrightarrow{n}\right| }{\left| \overrightarrow{d}\right|  \cdot  \left| \overrightarrow{n}\right| }$ .

\begin{itemize}\relax
\item\relax 二面角:从一条直线出发的两个半平面组成的几何图形, 用其平面角度量.
\end{itemize}

平面角的作法: \ding{172} 定义:在棱上任意取一点，过这点分别在两个面内作棱的垂线；

 \ding{173} 在棱上任意取一点，过这点作棱的垂面，得两条交线(射线)所成的最小正角；

 \ding{174}  当已知一个半平面的垂线时，可用三垂线定理作二面角的平面角.

 \ding{72} 二面角 $\alpha  - l - \beta$ 的向量求法:设平面 $\alpha$ 和 $\beta$ 的法向量分别为 ${\overrightarrow{n}}_{1}$ 和 ${\overrightarrow{n}}_{2}$ ，向量 ${\overrightarrow{n}}_{1}$ 和 ${\overrightarrow{n}}_{2}$ 的夹角为 $\theta$ ，则 $\cos \theta  = \frac{{\overrightarrow{n}}_{1} \cdot  {\overrightarrow{n}}_{2}}{\left| {\overrightarrow{n}}_{1}\right|  \cdot  \left| {\overrightarrow{n}}_{2}\right| }$ ,若二面角 $\alpha  - l - \beta$ 是锐角,则其大小为 $\arccos \frac{\left| {\overrightarrow{n}}_{1} \cdot  {\overrightarrow{n}}_{2}\right| }{\left| {\overrightarrow{n}}_{1}\right|  \cdot  \left| {\overrightarrow{n}}_{2}\right| }$ ; 若二面角 $\alpha  - l - \beta$ 是钝角,则

其大小为 $\pi  - \arccos \frac{\left| {\overrightarrow{n}}_{1} \cdot  {\overrightarrow{n}}_{2}\right| }{\left| {\overrightarrow{n}}_{1}\right|  \cdot  \left| {\overrightarrow{n}}_{2}\right| }$ . 7. 空间中的距离(转化为两点间的距离构造三角形或利用空间向量的投影或模):

 \ding{72} 点到平面的距离 (也可用等积法): 设平面 $\alpha$ 的斜线段是 ${AB}\left( {B \in  \alpha }\right)$ ,平面 $\alpha$ 的法向量为 $\overrightarrow{n}$ ,点 $A$ 到平面 $\alpha$ 的距离为 $d$ ,则 $d = \frac{\left| \overrightarrow{AB} \cdot  \overrightarrow{n}\right| }{\left| \overrightarrow{n}\right| }$ .

第二部分:多面体与旋转体

1. 棱柱的概念和主要性质:

\begin{center}
\begin{tabular}{|c|c|c|c|}
\hline
名称 & 棱柱 & 直棱柱。 & 正棱柱 \\
\cline{1-4}
图形 & \includegraphics[width=0.2\textwidth]{images/bo_d7tk7oc91nqc73875fu0_165_389_532_221_176_0.jpg} & \includegraphics[width=0.2\textwidth]{images/bo_d7tk7oc91nqc73875fu0_165_741_526_184_189_0.jpg} & \includegraphics[width=0.2\textwidth]{images/bo_d7tk7oc91nqc73875fu0_165_1074_527_189_189_0.jpg} \\
\cline{1-4}
定义 & 有一对互相平行的面, 且这两个面是两个全等的三角形或平面多边形; 不在这两个面上的棱都相互平行. 这样的多面体叫做棱柱 & 侧棱垂直于底面的棱柱 & 底面是正多边形的直棱柱 \\
\cline{1-4}
侧棱 & 平行且相等 & 平行且相等 & 平行且相等 \\
\cline{1-4}
侧面 & 平行四边形 & 矩形 & 全等的矩形 \\
\cline{1-4}
对角面 & 平行四边形 & 矩形 & 矩形 \\
\cline{1-4}
平行于底面的截面 & 与底面全等的多边形 & 与底面全等的多边形 & 与底面全等的正多边形 \\
\cline{1-4}
\hline
\end{tabular}
\end{center}

2. 棱锥的基本概念和主要性质:

\begin{center}
\begin{tabular}{|c|c|c|}
\hline
名称 & 棱锥 & 正棱锥 \\
\cline{1-3}
图形 & \includegraphics[width=0.2\textwidth]{images/bo_d7tk7oc91nqc73875fu0_165_490_1307_197_149_0.jpg} & \includegraphics[width=0.2\textwidth]{images/bo_d7tk7oc91nqc73875fu0_165_986_1308_197_144_0.jpg} \\
\cline{1-3}
定义 & 有一个面是多边形, 其余各面是有一个公共顶点的三角形的多面体 & 底面是正多边形, 且顶点在底面的射影是底面正多边形的中心 \\
\cline{1-3}
侧棱 & 相交于一点但不一定相等 & 相交于一点且相等 \\
\cline{1-3}
侧面 & 三角形 & 全等的等腰三角形 \\
\cline{1-3}
对角面 & 三角形 & 等腰三角形 \\
\cline{1-3}
平行于底的截面 & 与底面相似的多边形 & 与底面相似的正多边形 \\
\cline{1-3}
其他性质 & \phantom{X} & 高过底面中心; 侧棱与底面、侧面与底面、相邻两侧面所成角分别相等 \\
\cline{1-3}
\hline
\end{tabular}
\end{center}

3. 几种特殊四棱柱的特殊性质:

\begin{center}
\begin{tabular}{|c|c|c|c|}
\hline
名称 & 定义 & \multicolumn{2}{|c|}{特殊性质} \\
\cline{1-4}
平行六面体 & 底面是平行四边形的棱柱 & \multirow{4}{*}{对角线交于一点且被这点平分} & 底面和侧面都是平行四边形; 四条对角线交于一点, 且 被该点平分 \\
\cline{1-2}
\cline{4-4}
直平行六面体 & 侧棱垂直于底面的平行六面体 &  & 侧棱垂直于底面, 各侧面都是矩形; 四条对角线交于一点, 且被该点平分 \\
\cline{1-2}
\cline{4-4}
长方体 & 底面是矩形的直平行六面体 &  & 底面和侧面都是矩形; 四条对角线相等, 交于一点,且被该点平分 \\
\cline{1-2}
\cline{4-4}
正方体 & 底面边长和侧棱长相等的长方体 &  & 棱长都相等，各面都是正方形四条对角线相等， 交于一点,且被该点平分 \\
\cline{1-4}
\hline
\end{tabular}
\end{center}

4. 几种特殊的四棱柱的关系:

\{正方体\} $\subseteq  \{$ 正四棱柱 $\}  \subseteq  \{$ 长方体 $\}  \subseteq  \{$ 直平行六面体 $\}  \subseteq  \{$ 平行六面体 $\}  \subseteq  \{$ 四棱柱 $\}$ .

5. 圆柱、圆锥、球的概念和基本性质:

\begin{center}
\begin{tabular}{|c|c|c|c|}
\hline
名称 & 圆柱 & 圆锥 & 球 \\
\cline{1-4}
定义 & 矩形面绕其一边旋转一周所形成的几何体 & 直角三角形面绕其一直角边旋转一周所形成的几何体 & 半圆面绕其直径旋转一周所形成的几何体 \\
\cline{1-4}
过轴的截面 & 全等的矩形 & 全等的等腰三角形 & 以球的半径为半径的圆 \\
\cline{1-4}
垂直于轴的截面 & 与底面相等的圆 & 半径小于底面半径的圆 & 半径不大于球的半径的圆 \\
\cline{1-4}
重要图形 & 母线、高、相应的半径构成的矩形 & 母线、高、相应的半径构成的直角三角形 & 小圆和大圆圆心连线、和相应的截面半径 \\
\cline{1-4}
\hline
\end{tabular}
\end{center}

6．棱柱、正棱锥、圆柱、圆锥、球的表面积和体积:[表中字母: $S$ (面积), $V$ (体积), $c$ (底面周长), ${c}^{\prime }$ (直截面周长)， $h$ (高)， ${h}^{\prime }$ (斜高)， $l$ (侧棱(母线)长)， $r$ (底面半径).]

\begin{center}
\begin{tabular}{|c|c|c|c|c|c|c|}
\hline
名称 & \multicolumn{2}{|c|}{棱柱} & 正棱锥 & 圆锥 & 圆柱 & 球 (半径 $R$ ) \\
\cline{1-7}
\multirow{2}{*}{展开图} & ${c}^{\prime }$ 斜棱柱 & 直棱柱 & \multirow{2}{*}{\includegraphics[width=0.2\textwidth]{images/bo_d7tk7oc91nqc73875fu0_166_711_1297_165_171_0.jpg}} & \multirow{2}{*}{\includegraphics[width=0.2\textwidth]{images/bo_d7tk7oc91nqc73875fu0_166_883_1298_165_156_0.jpg}} & \multirow{2}{*}{\includegraphics[width=0.2\textwidth]{images/bo_d7tk7oc91nqc73875fu0_166_1064_1305_171_137_0.jpg}} & \multirow{2}{*}{\includegraphics[width=0.2\textwidth]{images/bo_d7tk7oc91nqc73875fu0_166_1252_1288_162_168_0.jpg}} \\
\cline{2-3}
 & \includegraphics[width=0.2\textwidth]{images/bo_d7tk7oc91nqc73875fu0_166_334_1350_177_125_0.jpg} & \includegraphics[width=0.2\textwidth]{images/bo_d7tk7oc91nqc73875fu0_166_524_1342_175_138_0.jpg} &  &  &  &  \\
\cline{1-7}
侧面积 ( ${S}_{\text{ 侧 }}$ ) & ${S}_{\text{ 侧 }} = {c}^{\prime }l$ & ${S}_{\text{ 侧 }} = {cl}$ & ${S}_{\text{ 侧 }} = \frac{1}{2}c{h}^{\prime }$ & ${S}_{\text{ 侧 }} = \frac{1}{2}{cl}$ & ${S}_{\text{ 侧 }} = {cl}$ & \multirow{2}{*}{${S}_{\text{ 表 }} = {4\pi }{R}^{2}$} \\
\cline{1-6}
全面积 ( ${S}_{\text{ 全 }}$ ) & \multicolumn{2}{|c|}{${S}_{\text{ 全 }} = {S}_{\text{ 侧 }} + 2{S}_{\text{ 底 }}$} & \multicolumn{2}{|c|}{${S}_{\text{ 全 }} = {S}_{\text{ 侧 }} + {S}_{\text{ 底 }}$} & ${S}_{\text{ 全 }} = {S}_{\text{ 侧 }} + 2{S}_{\text{ 底 }}$ &  \\
\cline{1-7}
体积 (V) & \multicolumn{2}{|c|}{$V = {Sh}$} & $V = \frac{1}{3}{Sh}$ & $V = \frac{1}{3}\pi {r}^{2}h$ & $V = {Sh}$ & $V = \frac{4}{3}\pi {R}^{3}$ \\
\cline{1-7}
\hline
\end{tabular}
\end{center}

7. “斜二侧”画法 (横等长,纵折半,纵横相交 ${45}^{ \circ  }$ 或 ${135}^{ \circ  }$ ) 的两条重要性质:

(1)平行直线的斜二测图仍是平行直线；

(2)线段及线段上定比分点的斜二测图保持原比例不变.

8. 空间中的面积和体积的计算:关键要记住公式，特别是公式中的系数，找到相应的元素.


提醒:立体几何的求解问题分为“作、证、算、结”四个部分，任意一个环节都不能省略.

9. 祖暅原理:夹在两个平行平面间的两个几何体，如果被平行于这两个平面的任意平面截得的两个截面都有相等的面积, 那么这两个几何体的体积必相等.

\section*{$\text{ § }{9.1}$ 空间直线与平面}

\section*{典型例题精讲}

【例 1】已知 $a$ 、 $b$ 、 $c$ 、 $d$ 是空间的 4 条直线.

(1)若 $a//b//c$ ，且 $d$ 与 $a$ 、 $b$ 、 $c$ 都相交，求证:直线 $a$ 、 $b$ 、 $c$ 、 $d$ 在同一平面上；

(2)若 $a$ 、 $b$ 、 $c$ 、 $d$ 两两相交且不过同一点，求证:直线 $a$ 、 $b$ 、 $c$ 、 $d$ 在同一平面上.

【例 2】已知异面直线 $a$ 、 $b$ 所成的角为 $\theta$ ，过空间中一点的直线 $l$ 与 $a$ 、 $b$ 所成的角都是 $\alpha$ .

(1)若 $\theta  = {60}^{ \circ  }$ ， $\alpha  = {75}^{ \circ  }$ ，求直线 $l$ 的条数；

(2)试根据 $\theta$ 和 $\alpha$ 的大小讨论直线 $l$ 的条数.

【例 3】在空间中,过点 $A$ 作平面 $\pi$ 的垂线,垂足为 $B$ ,记 $B = {f}_{\pi }\left( A\right)$ . 设 $\alpha \text{ 、 }\beta$ 是两个不同的平面,对空间任意一点 $P,{Q}_{1} = {f}_{\beta }\left\lbrack  {{f}_{\alpha }\left( P\right) }\right\rbrack  ,{Q}_{2} = {f}_{\alpha }\left\lbrack  {{f}_{\beta }\left( P\right) }\right\rbrack$ ,恒有 $P{Q}_{1} = P{Q}_{2}$ ,则 ( )

A. 平面 $\alpha$ 与平面 $\beta$ 垂直. B. 平面 $\alpha$ 与平面 $\beta$ 所成的 (锐) 二面角为 ${45}^{ \circ  }$ .

C. 平面 $\alpha$ 与平面 $\beta$ 平行. D. 平面 $\alpha$ 与平面 $\beta$ 所成的 (锐) 二面角为 ${60}^{ \circ  }$ .

【例 4】如图,已知二面角 $\alpha  - {MN} - \beta$ 的大小为 ${60}^{ \circ  }$ ,菱形 ${ABCD}$ 在面 $\beta$ 内, $A\text{ 、 }B$ 两点在棱 ${MN}$ 上, $\angle {BAD} = {60}^{ \circ  },E$ 是 ${AB}$ 的中点, ${DO} \bot$ 面 $\alpha$ ,垂足为 $O$ .

(1)证明: ${AB}\bot$ 平面 ${ODE}$ ;

(2)求异面直线 ${BC}$ 与 ${OD}$ 所成角的余弦值.

\begin{center}
\includegraphics[width=0.2\textwidth]{images/bo_d7tk7oc91nqc73875fu0_168_1148_1099_344_192_0.jpg}
\end{center}
 

(例 4 图)


【例 5】如图,在四棱锥 $P - {ABCD}$ 中, ${AB}//{CD},{AB} \bot  {AD},{CD} = {2AB}$ ,平面 ${PAD} \bot$ 底面 ${ABCD},{PA}  \bot  {AD}.E$ 和 $F$ 分别是 ${CD}$ 和 ${PC}$ 的中点,求证:

(1) ${PA} \bot$ 底面 ${ABCD}$ ；

(2) ${BE}//$ 平面 ${PAD}$ ；

(3)平面 ${BEF} \bot$ 平面 ${PCD}$ .

\begin{center}
\includegraphics[width=0.2\textwidth]{images/bo_d7tk7oc91nqc73875fu0_169_1149_367_296_307_0.jpg}
\end{center}
 

(例 5 图)

\section*{四基达标精练}

\section*{一、填空题}

1. 三个平面两两相交得到三条交线，则这三条交线的位置关系是\_\_\_\_\_.

2. 垂直于同一直线的两直线位置关系是\_\_\_\_\_.

3. 设 $m\text{ 、 }n$ 是两条不同的直线, $\alpha \text{ 、 }\beta$ 是两个不同的平面,若 $m \bot  \beta ,n \bot  \beta ,n \bot  \alpha$ ,则 $m$ 与平面 $\alpha$ 的位置关系的最精确的描述是\_\_\_\_\_.

4. 已知三棱锥 $V - {ABC}$ 中， ${VA} = {VB} = {VC}$ ，则 $V$ 在平面 ${ABC}$ 内的射影是 $\bigtriangleup  {ABC}$ 的\_\_\_\_\_，心.

5. 垂直于同一平面的两直线的位置关系是\_\_\_\_\_.

6. 已知直线 $l\bot$ 平面 $\alpha$ ,直线 $m \subset$ 平面 $\beta$ ,有下列命题: \ding{172}  $\alpha //\beta  \Rightarrow  l\bot m$ ;  \ding{173}  $\alpha \bot \beta  \Rightarrow  l//m$ ;  \ding{174}  $l//m \Rightarrow  \alpha \bot \beta$ ；  \ding{175}  $l\bot m \Rightarrow  \alpha //\beta$ . 其中正确命题的序号是\_\_\_\_\_.

\begin{center}
\includegraphics[width=0.2\textwidth]{images/bo_d7tk7oc91nqc73875fu0_169_1117_1679_330_279_0.jpg}
\end{center}
 

(第 7 题图)

7. 如图，三棱锥 $V - {ABC}$ 中， $\angle {AVB} = \angle {AVC} = {45}^{ \circ  }$ ， $\angle {BVC} = {60}^{ \circ  }$ ， ${VA} = {VB}  = {VC} = 2$ ，过点 $A$ 作的截面分别交 ${VB}$ 、 ${VC}$ 于点 $E$ 、 $F$ ，则截面周长的最小值为\_\_\_\_\_.

8. 设 $E\text{ 、 }F$ 是空间四边形 ${ABCD}$ 的边 ${AD}\text{ 、 }{BC}$ 的中点,则 ${EF}$ 与 $\frac{1}{2}({AB} +  {CD})$ 的大小关系是\_\_\_\_\_.

9. 已知正三棱柱 ${ABC} - {A}_{1}{B}_{1}{C}_{1}$ 的侧棱长与底面边长相等,则 $A{B}_{1}$ 与侧面 ${AC}{C}_{1}{A}_{1}$ 所成角的正弦值等于\_\_\_\_\_.

\section*{二、选择题}

\begin{center}
\includegraphics[width=0.2\textwidth]{images/bo_d7tk7oc91nqc73875fu0_170_1179_202_301_272_0.jpg}
\end{center}
 

(第 10 题图)

10. 如图,在棱长为 10 的正方体 ${ABCD} - {A}_{1}{B}_{1}{C}_{1}{D}_{1}$ 中, $P$ 为左侧面 ${AD}{D}_{1}{A}_{1}$ 上一点,已知点 $P$ 到 ${A}_{1}{D}_{1}$ 的距离为 $3,P$ 到 $A{A}_{1}$ 的距离为 2,则过点 $P$ 且与 ${A}_{1}C$ 平行的直线相交的面是 ( )

A. $A{A}_{1}{B}_{1}B$ . B. $B{B}_{1}{C}_{1}C$ .

C. $C{C}_{1}{D}_{1}D$ . D.ABCD.

11. 已知直线 $a$ 与 $b$ 所成的角和直线 $c$ 与 $b$ 所成的角的大小都是 ${60}^{ \circ  }$ ,则直线 $a$ 与 $c$ 的位置关系是 ( )

A. 平行. B. 相交.

C. 异面. D. 平行、相交、异面都可能.

\begin{center}
\includegraphics[width=0.2\textwidth]{images/bo_d7tk7oc91nqc73875fu0_170_1183_595_301_270_0.jpg}
\end{center}
 

(第 14 题图)

12. 如图，正方体 ${ABCD} - {A}_{1}{B}_{1}{C}_{1}{D}_{1}$ 中， $E$ 、 $F$ 分别为 ${A}_{1}{B}_{1}$ 、 ${CD}$ 的中点，则异面直线 ${D}_{1}E$ 与 ${A}_{1}F$ 所成的角的余弦值为 ( )

A. $\frac{\sqrt{5}}{5}$ . B. $\frac{\sqrt{5}}{6}$ .

C. $\frac{\sqrt{3}}{3}$ . D. $\frac{\sqrt{3}}{6}$ .

\section*{三、解答题}

13. 已知 $\bigtriangleup {ABC}$ 在平面 $\alpha$ 外, ${AB} \cap  \alpha  = P,{BC} \cap  \alpha  = Q,{AC} \cap  \alpha  = R$ .

(1)证明 $\bigtriangleup  {ABC}$ 是平面图形；

(2)证明 $P\text{ 、 }Q\text{ 、 }R$ 在同一直线上.

14. 如图所示，在四棱锥 $P - {ABCD}$ 中， ${PA}\bot$ 底面 ${ABCD},{AB}\bot {AD},{AC}\bot {CD},\angle {ABC} = {60}^{ \circ  },{PA} = {AB}  = {BC},E$ 是 ${PC}$ 的中点. 证明:

\begin{center}
\includegraphics[width=0.2\textwidth]{images/bo_d7tk7oc91nqc73875fu0_170_1232_1613_266_354_0.jpg}
\end{center}
 

(第 14 题图)

(1) ${CD} \bot  {AE}$ ；

(2) ${PD} \bot$ 平面 ${ABE}$ .

15. 已知 $E$ 、 $F$ 、 $G$ 、 $H$ 分别是空间四边形 ${ABCD}$ 的边 ${AB}$ 、 ${BC}$ 、 ${CD}$ 、 ${DA}$ 上的点,若 $\frac{AE}{EB} = \frac{CF}{FB} = \frac{DG}{GC} = \frac{DH}{HA} = \frac{1}{2}$ . 求证:(1) ${EF}$ 、 ${CD}$ 是异面直线；

(2)四边形 ${EFGH}$ 是梯形.

\begin{center}
\includegraphics[width=0.2\textwidth]{images/bo_d7tk7oc91nqc73875fu0_171_1139_319_303_252_0.jpg}
\end{center}
 

(第 15 题图)

16. 如图， ${AB}$ 是圆 $O$ 的直径， ${PA}\bot$ 圆 $O$ 所在的平面， $C$ 是圆 $O$ 上的点.

(1)求证: ${BC}\bot$ 平面 ${PAC}$ ；

(2)若Q为 ${PA}$ 的中点， $G$ 为 $\bigtriangleup  {AOC}$ 的重心，求证: ${QG}//$ 平面 ${PBC}$ .

\begin{center}
\includegraphics[width=0.2\textwidth]{images/bo_d7tk7oc91nqc73875fu0_171_1097_1147_346_257_0.jpg}
\end{center}
 

(第 16 题图)

\section*{§ 9.2 多面体与旋转体}

\section*{典型例题精讲}

【例 1】如图,将圆心角为 ${288}^{ \circ  }$ 的扇形卷成一个底面半径为 4 的圆锥的侧面.

(1)求圆锥母线的长.

(2)求过圆锥顶点 $O$ 所作截面的最大面积.

\begin{center}
\includegraphics[width=0.4\textwidth]{images/bo_d7tk7oc91nqc73875fu0_172_919_584_565_222_0.jpg}
\end{center}
 

(例 1 图)

【例 2】如图,在棱长为 $a$ 的正方体 ${ABCD} - {A}_{1}{B}_{1}{C}_{1}{D}_{1}$ 中, $E\text{ 、 }F$ 分别为棱 ${A}_{1}{B}_{1}\text{ 、 }{B}_{1}{C}_{1}$ 的中点.

(1)求二面角 ${B}_{1} - {BF} - E$ 的大小；

(2)求点 $D$ 到平面 ${BEF}$ 的距离；

(3)能否在棱 ${B}_{1}B$ 上找到一点 $M$ ，使 ${DM}\bot$ 面 ${BEF}$ ？若能，请确定点 $M$ 的位置；若不能，请说明理由.

\begin{center}
\includegraphics[width=0.2\textwidth]{images/bo_d7tk7oc91nqc73875fu0_172_1213_1501_278_242_0.jpg}
\end{center}
 

(例 2 图)

【例 3】如图,在底面为等腰直角三角形的直三棱柱 ${ABC} - {A}_{1}{B}_{1}{C}_{1}$ 中, ${AC} = {BC} = A{A}_{1} = a,\angle {ACB} = {90}^{ \circ  }$ .

(1)求 $A{B}_{1}$ 与 $B{C}_{1}$ 所成角的大小；

(2)求 ${A}_{1}C$ 与平面 ${AB}{B}_{1}{A}_{1}$ 所成角的大小；

(3)求二面角 ${A}_{1} - A{B}_{1} - {C}_{1}$ 的大小.

\begin{center}
\includegraphics[width=0.2\textwidth]{images/bo_d7tk7oc91nqc73875fu0_173_1259_334_228_256_0.jpg}
\end{center}
 

(例 3 图)

【例 4】在棱长为 $a$ 的正方体 ${ABCD} - {A}_{1}{B}_{1}{C}_{1}{D}_{1}$ 中, $E\text{ 、 }F$ 分别为 ${AB}$ 与 ${C}_{1}{D}_{1}$ 的中点.

(1)求证:四边形 ${A}_{1}{ECF}$ 是菱形；

(2)求证: ${EF}\bot$ 平面 ${A}_{1}{B}_{1}C$ ；

(3)求 ${A}_{1}{B}_{1}$ 与平面 ${A}_{1}{ECF}$ 所成角的正切值.

\begin{center}
\includegraphics[width=0.2\textwidth]{images/bo_d7tk7oc91nqc73875fu0_173_1206_1249_278_259_0.jpg}
\end{center}
 

(例 4 图)

【例 5】春节期间，佳怡准备去探望奶奶，她到商店买了一盒点心. 为了美观起见，售货员对点心盒做了一个捆扎[如图()所示]，并在角上配了一个花结. 售货员说，这样的捆扎不仅漂亮，而且比一般的十字捆扎 [如图(2)所示]包装更节省彩绳. 你同意这种说法吗? 请给出你的理由. (注:长方体点心盒的高小于长、宽)

\begin{center}
\includegraphics[width=0.4\textwidth]{images/bo_d7tk7oc91nqc73875fu0_174_906_290_562_181_0.jpg}
\end{center}
 

(例 5 图)

\section*{四基达标精练}

\section*{一、填空题}

1. 正方体 ${ABCD} - {A}_{1}{B}_{1}{C}_{1}{D}_{1}$ 中，异面直线 $C{D}_{1}$ 和 $B{C}_{1}$ 所成的角是\_\_\_\_\_.

2. 一个棱柱有 10 个顶点, 其所有的侧棱长的和为 60, 则每条侧棱长为\_\_\_\_\_.

3. 正四面体的中心到底面的距离与这个四面体的高的比值是\_\_\_\_\_.

4. 用与球心距离为 2 的平面截球 $O$ ，得到半径为 $\sqrt{5}$ 的圆面，则球 $O$ 的半径是\_\_\_\_\_.

5．在半径为 30 米的圆形广场中央设置一个照明光源，它射向地面的光呈圆锥形，其轴截面顶角为 120 米， 若要光源恰好照亮整个广场，则光源的高度应为\_\_\_\_\_米。

6. 将半径为 1 的圆分割成面积之比为 $1 : 2 : 3$ 的三个扇形作为三个圆锥的侧面,设这三个圆锥底面半径依次为 ${r}_{1}\text{ 、 }{r}_{2}\text{ 、 }{r}_{3}$ ，那么 ${r}_{1} + {r}_{2} + {r}_{3} =$ \_\_\_\_\_.

7. 用一个平行于圆锥底面的平面截这个圆锥, 截得的圆台上、下底面的面积之比为 1 : 16，若截去的圆锥的母线长是 3 ，则圆台的母线长是\_\_\_\_\_.

8. 在棱长为 1 的正方体 ${ABCD} - {A}_{1}{B}_{1}{C}_{1}{D}_{1}$ 中,点 $P$ 在侧面 ${BC}{C}_{1}{B}_{1}$ 上运动. 若点 $P$ 到点 $A$ 的距离为 $\frac{2\sqrt{3}}{3}$ ，则动点 $P$ 的轨迹在正方形 ${BC}{C}_{1}{B}_{1}$ 内的一段曲线长为\_\_\_\_\_.

9. 有四根长度都为 2 的直铁条,若再选两根长度都为 $a$ 的直铁条,使这六根铁条端点处相连能够焊接成一个三棱锥形的铁架,则 $a$ 的取值范围是\_\_\_\_\_.

\section*{二、选择题}

\begin{center}
\includegraphics[width=0.2\textwidth]{images/bo_d7tk7oc91nqc73875fu0_174_1154_1744_343_259_0.jpg}
\end{center}
 

(第 10 题图)

10. 纸制的正方体的六个面根据其方位分别标记为上、下、东、南、西、北. 现在沿该正方体的一些棱将正方体剪开、外面朝上展平, 得到如图的平面图形, 则标“ $\bigtriangleup$ ”的面的方位是(从上往下看“东、南、西、北”顺时针排列)(

A. 南. B.北.

C. 西. D. 下.

\begin{center}
\includegraphics[width=0.3\textwidth]{images/bo_d7tk7oc91nqc73875fu0_175_1075_156_374_224_0.jpg}
\end{center}
 

(第 11 题图)

11. 如图,三棱柱 ${ABC} - {A}_{1}{B}_{1}{C}_{1}$ 被平面 ${DE}{E}_{1}{D}_{1}$ 截成两个几何体 ${I}$ 、 II,且平面 ${DE}{E}_{1}{D}_{1}//$ 平面 ${AB}{B}_{1}{A}_{1}$ ,则 ( )

A. I 是棱柱, II 是棱台.

B. I 是棱锥, II 是棱台.

C. I 是棱锥, II 是棱柱.

D. I 是棱柱, II 是棱柱.

12. 如图所示的几何体是由一个圆柱挖去一个以圆柱上底面为底面、下底面圆心为顶点的圆锥而得到的组合体, 现用一个竖直的平面去截这个组合体, 则截面图形可能是 ( )

\begin{center}
\includegraphics[width=0.2\textwidth]{images/bo_d7tk7oc91nqc73875fu0_175_149_524_184_202_0.jpg}
\end{center}
 

A.

\begin{center}
\includegraphics[width=0.2\textwidth]{images/bo_d7tk7oc91nqc73875fu0_175_403_529_215_199_0.jpg}
\end{center}
 

B.

\begin{center}
\includegraphics[width=0.2\textwidth]{images/bo_d7tk7oc91nqc73875fu0_175_683_535_203_193_0.jpg}
\end{center}
 

C.

\begin{center}
\includegraphics[width=0.2\textwidth]{images/bo_d7tk7oc91nqc73875fu0_175_954_536_205_193_0.jpg}
\end{center}
 

D.

\begin{center}
\includegraphics[width=0.2\textwidth]{images/bo_d7tk7oc91nqc73875fu0_175_1229_510_166_222_0.jpg}
\end{center}
 

\section*{三、解答题}

13. 如图,圆锥的顶点为 $P$ ,底面圆为 $O$ ,底面的一条直径为 ${AB},C$ 为半圆弧 $\overset{\text{ ⏜ }}{AB}$ 的中点, $E$ 为劣弧 $\overset{\text{ ⏜ }}{CB}$ 的中点,已知 ${PO} = 2,{OA} = 1$ .

\begin{center}
\includegraphics[width=0.2\textwidth]{images/bo_d7tk7oc91nqc73875fu0_175_1149_912_296_316_0.jpg}
\end{center}
 

(第 13 题图)

(1)求证: ${AO}\bot {PC}$ ；

(2)求异面直线 ${PA}$ 和 ${OE}$ 所成角的大小.

14. 如图,在正三棱柱 ${ABC} - {A}_{1}{B}_{1}{C}_{1}$ 中,已知 $D$ 是 ${AC}$ 的中点.

(1)求证: $A{B}_{1}//$ 平面 ${DB}{C}_{1}$ ；

(2)若 $A{B}_{1} \bot  B{C}_{1}$ ，求二面角 $D - B{C}_{1} - C$ 的大小.

\begin{center}
\includegraphics[width=0.3\textwidth]{images/bo_d7tk7oc91nqc73875fu0_175_1044_1571_401_180_0.jpg}
\end{center}
 

(第 14 题图)

15. 如图,现有一块边长为 $a{{cm}}$ 的正方形铁皮，将它的四个角上的阴影部分去掉，然后焊成一个底面边长为 $x{{cm}}$ 的正四棱锥状的无盖容器.

(1)当 $a = 2\sqrt{5},x = 2$ 时，求该容器的高 $h$ ；

(2)求该容器的高 $h$ 关于 $x$ 的函数 $h = f\left( x\right)$ ，指出定义域.

\begin{center}
\includegraphics[width=0.5\textwidth]{images/bo_d7tk7oc91nqc73875fu0_176_878_338_606_248_0.jpg}
\end{center}
 

(第 15 题图)

16. 在棱长为 1 的正方体 ${ABCD} - {A}_{1}{B}_{1}{C}_{1}{D}_{1}$ 中， $M$ 和 $N$ 分别为 ${A}_{1}{B}_{1}$ . 和 $B{B}_{1}$ 的中点.

(1)求直线 ${AM}$ 和 ${CN}$ 所成角；

(2)求直线 ${CN}$ 与平面 ${AC}{C}_{1}{A}_{1}$ 所成的角.

\begin{center}
\includegraphics[width=0.2\textwidth]{images/bo_d7tk7oc91nqc73875fu0_176_1218_1228_281_243_0.jpg}
\end{center}
 

(第 16 题图)

\section*{§ 9.3 多面体与旋转体的表面积与体积}

\section*{典型例题精讲}

【例 1】如图,在四棱锥 $P - {ABCD}$ 中,底面 ${ABCD}$ 是矩形, ${PA}$ 为点 $P$ 到平面 ${ABCD}$ 的距离, ${AB} = 4$ , ${AD} = 3,{PA} = 3$ ,点 $E\text{ 、 }M$ 分别在线段 ${AB}\text{ 、 }{PC}$ 上,其中 $E$ 是 ${AB}$ 的中点, $\frac{PM}{MC} = \lambda$ ,连接 ${ME}$ .

(1)当 $\lambda  = 1$ 时,证明: 直线 ${ME}//$ 平面 ${PAD}$ ;

(2)当 $\lambda  = 2$ 时,求三棱锥 $M - {BCD}$ 的体积.

\begin{center}
\includegraphics[width=0.2\textwidth]{images/bo_d7tk7oc91nqc73875fu0_177_1113_649_332_230_0.jpg}
\end{center}
 

(例 1 图)

【例 2】边长为 $a$ 的正三角形,要拼接成一个正三棱柱且不剩料.

(1)应如何设计？(在图中用虚线画出)

(2)求拼得的三棱柱的体积.

\begin{center}
\includegraphics[width=0.2\textwidth]{images/bo_d7tk7oc91nqc73875fu0_177_1186_1390_259_256_0.jpg}
\end{center}
 

(例 2 图)

【例 3】如图,三棱柱 ${ABC} - {A}_{1}{B}_{1}{C}_{1}$ 中, $A{A}_{1} \bot  {BC},{A}_{1}B \bot  B{B}_{1}$ .

(1)求证: ${A}_{1}C \bot  C{C}_{1}$ ；

(2)若 ${AB} = 2,{AC} = \sqrt{3},{BC} = \sqrt{7}$ ，问 $A{A}_{1}$ 为何值时，三棱柱 ${ABC} - {A}_{1}{B}_{1}{C}_{1}$ 体积最大，并求此最大值.

\begin{center}
\includegraphics[width=0.3\textwidth]{images/bo_d7tk7oc91nqc73875fu0_178_1086_289_398_189_0.jpg}
\end{center}
 

(例 3 图)

【例 4】正三棱锥的高为 1,底面边长为 $2\sqrt{6}$ ,内有一个球与它的四个面都相切. 求:

(1)棱锥的表面积；

(2)内切球的半径.

\begin{center}
\includegraphics[width=0.3\textwidth]{images/bo_d7tk7oc91nqc73875fu0_178_1121_1203_381_298_0.jpg}
\end{center}
 

(例 4 图)

【例 5】如图,某地质队自水平地面 $A\text{ 、 }B\text{ 、 }C$ 三处垂直向地下钻探,自点 $A$ 向下钻到 ${A}_{1}$ 处发现矿藏,再继续下钻到 ${A}_{2}$ 处后下面已无矿,从而得到在 $A$ 处正下方的矿层厚度为 ${A}_{1}{A}_{2} = {d}_{1}$ . 同样可得在 $B\text{ 、 }C$ 处正下方的矿层厚度分别为 ${B}_{1}{B}_{2} = {d}_{2},{C}_{1}{C}_{2} = {d}_{3}$ ,且 ${d}_{1} < {d}_{2} < {d}_{3}$ . 过 ${AB}\text{ 、 }{AC}$ 的中点 $M\text{ 、 }N$ 且与直线 $A{A}_{2}$ 平行的平面截多面体 ${A}_{1}{B}_{1}{C}_{1} - {A}_{2}{B}_{2}{C}_{2}$ 所得的截面 ${DEFG}$ 为该多面体的一个中截面,其面积记为 ${S}_{\text{ 中 }}$ .

(1)证明:中截面 ${DEFG}$ 是梯形；

(2)在 $\bigtriangleup  {ABC}$ 中，记 ${BC} = a,{BC}$ 边上的高为 $h$ ，面积为 $S$ . 在估测三角形 ${ABC}$ 区域内正下方的矿藏储量 (即多面体 ${A}_{1}{B}_{1}{C}_{1} - {A}_{2}{B}_{2}{C}_{2}$ 的体积 $V$ )时,可用近似公式 ${V}_{\text{ 估 }} = {S}_{\text{ 中 }} \cdot  h$ 来估算. 已知 $V = \frac{1}{3}\left( {{d}_{1} + {d}_{2} + {d}_{3}}\right)  S$ ,试判断 ${V}_{\text{ 估 }}$ 与 $V$ 的大小关系,并加以证明.

\begin{center}
\includegraphics[width=0.2\textwidth]{images/bo_d7tk7oc91nqc73875fu0_179_1116_541_317_322_0.jpg}
\end{center}
 

(例 5 图)

\section*{四基达标精练}

\section*{一、填空题}

1. 已知长方体全部棱长的和为 36 ，表面积为 52 ，则其体对角线的长为\_\_\_\_\_.

\begin{center}
\includegraphics[width=0.2\textwidth]{images/bo_d7tk7oc91nqc73875fu0_179_1158_1241_272_244_0.jpg}
\end{center}
 

(第 2 题图)

2. 如图，在正方体 ${ABCD} - {A}_{1}{B}_{1}{C}_{1}{D}_{1}$ 中， $P$ 为对角线 $B{D}_{1}$ 的三等分点，点 $P$ 到各顶点的距离的不同取值的个数是\_\_\_\_\_.

3. 两相邻边长分别为 3 和 4 的矩形, 以一边所在的直线为轴旋转所成的圆柱的底面积为\_\_\_\_\_.

4. 棱长为 $a$ 的正方体中,连接相邻面的中心,以这些线段为棱的八面体的体积为\_\_\_\_\_.

5. 若圆柱的底面半径是 1 ，其侧面展开是一个正方形，则这个圆柱的侧面积是\_\_\_\_\_.

6. 一个密闭且透明的正方体容器中装有部分液体，已知该正方体的棱长为 1 ，如果任意转动该正方体， 液面的形状都不可能是三角形，那么液体体积的取值范围是\_\_\_\_\_.

7. 有一棱长为 $a$ 的正方体框架,其内放置一气球,使其充气且尽可能地膨胀 (仍保持为球的形状),则气球表面积的最大值为\_\_\_\_\_.

8. 已知三边长分别为 $4\text{ 、 }5\text{ 、 }6$ 的 $\bigtriangleup {ABC}$ 的外接圆恰好是球 $O$ 的一个大圆, $P$ 为球面上一点,若点 $P$ 到 $\bigtriangleup  {ABC}$ 的三个顶点的距离相等，则三棱锥 $P - {ABC}$ 的体积为\_\_\_\_\_.

9. 三棱锥一条棱长为 $x$ ，其余各条棱长都等于 1，则体积 $F\left( x\right)$ 关于 $x$ 的函数解析式是\_\_\_\_\_.

\section*{二、选择题}

10. 将棱长为 $a$ 的正方体切成 27 个全等的小正方体，则这些小正方体的表面积之和是原正方体表面积的 ( )

A. 2 倍. B. 3 倍. C. 4 倍. D. 5 倍.

(第 11 题图)

\begin{center}
\includegraphics[width=0.3\textwidth]{images/bo_d7tk7oc91nqc73875fu0_180_1135_143_353_244_0.jpg}
\end{center}
 

11. 如图,已知正四棱台 ${ABCD} - {A}_{1}{B}_{1}{C}_{1}{D}_{1}$ 由正四棱锥 $S - {ABCD}$ 截得,且 ${AB} = 6,{A}_{1}{B}_{1} = 4$ ,正四棱台的高为 $\sqrt{2},E$ 是 ${A}_{1}{B}_{1}$ 的中点,有下列说法:  \ding{172}  正四棱锥 $S - {ABCD}$ 的体积为 ${24}\sqrt{2}$ ;  \ding{173}  平面 $B{C}_{1}D \bot$ 平面 $A{A}_{1}{C}_{1}C$ ;  \ding{174}  ${AE}//$ 平面 $B{C}_{1}D$ ； \ding{175}  正四棱台 ${ABCD} - {A}_{1}{B}_{1}{C}_{1}{D}_{1}$ 的表面积为 ${{20}\sqrt{3}} + {52}$ . 其中正确的个数是 ( )

A.1. B.2. C.3. D. 4.

12. 已知正方体 ${ABCD} - {A}_{1}{B}_{1}{C}_{1}{D}_{1}$ 的棱长为 $1,\overrightarrow{AP} = \lambda \overrightarrow{AB} + \mu \overrightarrow{AD}$ ,其中 $\lambda \text{ 、 }\mu  \in  \left\lbrack  {0,1}\right\rbrack$ ,点 $E$ 为线段 $A{A}_{1}$ 的中点,则下列选项错误的是 ( )

A. $\mu  = \frac{1}{2}$ 时， ${A}_{1}P \bot  E{D}_{1}$ .

B. $\lambda  = \mu$ 时,三棱锥 ${D}_{1} - {EBP}$ 的体积为定值.

C. $\lambda  + \mu  = 1$ 时,直线 ${A}_{1}P$ 与面 ${B}_{1}{D}_{1}E$ 的交点轨迹长度为 $\frac{\sqrt{2}}{3}$ .

D. 当点 $P$ 落在以 ${C}_{1}$ 为球心、 $\sqrt{2}$ 为半径的球面上时, $\lambda  + \mu$ 的最小值为 1 .

\section*{三、解答题}

13. 圆锥的母线长是它的高的 2 倍, 过顶点最大的截面交底面得一条弦, 底面圆心到弦的距离为 1 , 求:

(1)圆锥的侧面展开图的圆心角的大小；

(2)圆锥的高和侧面积.

\begin{center}
\includegraphics[width=0.2\textwidth]{images/bo_d7tk7oc91nqc73875fu0_180_1166_969_331_226_0.jpg}
\end{center}
 

(第 13 题图)

14. 已知 ${ABCD}$ 是边长为 1 的正方形,正方形 ${ABCD}$ 绕 ${AB}$ 旋转形成一个圆柱.

(1)求该圆柱的表面积；

(2)正方形 ${ABCD}$ 绕 ${AB}$ 逆时针旋转 $\frac{\pi }{2}$ 至 ${AB}{C}_{1}{D}_{1}$ ，求线段 $C{D}_{1}$ 与平面 ${ABCD}$ 所成的角.

\begin{center}
\includegraphics[width=0.2\textwidth]{images/bo_d7tk7oc91nqc73875fu0_180_1211_1619_298_242_0.jpg}
\end{center}
 

(第 14 题图)

15. 如图,已知正三棱柱 ${ABC} - {A}_{1}{B}_{1}{C}_{1}$ 的底面边长是 $2,D\text{ 、 }E$ 是 $C{C}_{1}\text{ 、 }{BC}$ 的中点, ${AE} = {DE}$ . 求:

(1)正三棱柱 ${ABC} - {A}_{1}{B}_{1}{C}_{1}$ 的侧棱长;

(2)正三棱柱 ${ABC} - {A}_{1}{B}_{1}{C}_{1}$ 的表面积.

\begin{center}
\includegraphics[width=0.2\textwidth]{images/bo_d7tk7oc91nqc73875fu0_181_1111_281_326_279_0.jpg}
\end{center}
 

(第 15 题图)

16. 已知四棱锥 $P - {ABCD}$ 的底面 ${ABCD}$ 为菱形， ${PD}\bot$ 平面 ${ABCD},{PD} = {AD} = {CD} = 2,\angle {BAD} = \frac{\pi }{3}$ ， $E$ 为 ${PC}$ 上一点.

(1)设平面 ${PAD} \cap$ 平面 ${PBC} = l$ ，证明: ${BC}//l$ ；

(2)当直线 ${BE}$ 与平面 ${BCD}$ 的夹角为 $\frac{\pi }{6}$ 时，求三棱锥 $P - {BDE}$ 体积.

\begin{center}
\includegraphics[width=0.2\textwidth]{images/bo_d7tk7oc91nqc73875fu0_181_1131_1270_302_298_0.jpg}
\end{center}
 

(第 16 题图)

\section*{§ 9.4 空间向量及其运算}

\section*{典型例题精讲}

【例 1】已知点 $A\left( {-2,3, - 3}\right) ,B\left( {4,5,9}\right)$ .

(1)设平面 $\alpha$ 经过线段 ${AB}$ 的中点，且与直线 ${AB}$ 垂直， $M\left( {x,y,z}\right)$ 是平面 $\alpha$ 内任意一点，求 $x$ 、 $y$ 、 $z$ 满足的关系式;

(2)求到 $A$ 、 $B$ 两点距离相等的点 $P\left( {x,y,z}\right)$ 的坐标 $x$ 、 $y$ 、 $z$ 满足的关系式；

(3)比较 (1) (2) 的结论, 你发现了什么结论?

【例 2】如图,已知平行六面体 ${ABCD} - {A}_{1}{B}_{1}{C}_{1}{D}_{1}$ 中,对角线 $B{D}_{1}$ 与平面 $A{B}_{1}C$ 交于点 $F$ ,连接 ${AF}$ 交 ${B}_{1}C$ 于点 $E$ . 证明:

(1)点 $E$ 为 ${B}_{1}C$ 的中点；

(2) ${AF} = \frac{2}{3}{AE}$ ，且 ${BF} = \frac{1}{3}{B{D}_{1}}$ .

\begin{center}
\includegraphics[width=0.3\textwidth]{images/bo_d7tk7oc91nqc73875fu0_182_1116_1425_390_261_0.jpg}
\end{center}
 

(例 2 图)

【例 3】已知四棱锥 $P - {ABCD}$ 中,底面 ${ABCD}$ 是边长为 $2\sqrt{2}$ 的正方形,平面 ${PAC} \bot$ 底面 ${ABCD},{PA}  = {PC} = 2\sqrt{2}$ .

(1)求证: ${PB} = {PD}$ ；

(2)若点 $M$ 、 $N$ 分别是棱 ${PA}$ 、 ${PC}$ 的中点，平面 ${DMN}$ 与棱 ${PB}$ 的交点 $Q$ ，在线段 ${BC}$ 上是否存在点 $H$ ， 使得 ${DQ} \bot  {PH}$ ，若存在，求 ${BH}$ 的长，否则请说明理由.

\begin{center}
\includegraphics[width=0.3\textwidth]{images/bo_d7tk7oc91nqc73875fu0_183_1050_391_390_262_0.jpg}
\end{center}
 

(例 3 图)

【例 4】已知空间向量 $\overrightarrow{a} = \left( {{a}_{1},{a}_{2},{a}_{3}}\right) ,\overrightarrow{b} = \left( {{b}_{1},{b}_{2},{b}_{3}}\right)$ ,定义两个空间向量 $- 4) = \left( {k - 1}\right) \left( {k + 2}\right)  + {k}^{2} -  8 = 0$ 与 $2{k}^{2} + k - {10} = 0$ 之间的距离为 $d\left( {\overrightarrow{a},\overrightarrow{b}}\right)  = \left| {{b}_{1} - {a}_{1}}\right|  + \left| {{b}_{2} - {a}_{2}}\right|  + \left| {{b}_{3} - {a}_{3}}\right|$ .

(1)若 $\overrightarrow{a} = \left( {1,2,3}\right) ,\overrightarrow{b} = \left( {4,1,1}\right) ,\overrightarrow{c} = \left( {\frac{11}{2},\frac{1}{2},0}\right)$ ,证明: $d\left( {\overrightarrow{a},\overrightarrow{b}}\right)  + d\left( {\overrightarrow{b},\overrightarrow{c}}\right)  = d\left( {\overrightarrow{a},\overrightarrow{c}}\right)$ ；

(2)已知 $\overrightarrow{c} = \left( {{c}_{1},{c}_{2},{c}_{3}}\right)$ .

 \ding{172}  证明:若存在 $\lambda  > 0$ ，使 $\overrightarrow{b} - \overrightarrow{a} = \lambda \left( {\overrightarrow{c} - \overrightarrow{b}}\right)$ ，则 $d\left( {\overrightarrow{a},\overrightarrow{b}}\right)  + d\left( {\overrightarrow{b},\overrightarrow{c}}\right)  = d\left( {\overrightarrow{a},\overrightarrow{c}}\right)$ ；

 \ding{173}  若 $d\left( {\overrightarrow{a},\overrightarrow{b}}\right)  + d\left( {\overrightarrow{b},\overrightarrow{c}}\right)  = d\left( {\overrightarrow{a},\overrightarrow{c}}\right)$ ，是否一定存在 $\lambda  > 0$ ，使 $\overrightarrow{b} - \overrightarrow{a} = \lambda \left( {\overrightarrow{c} - \overrightarrow{b}}\right)$ ？请说明理由.

【例 5】如图,在直角梯形 ${ABCD}$ 中, ${AB}//{DC},\angle {ABC} = {90}^{ \circ  },{AB} = {2DC} = {2BC},E$ 为 ${AB}$ 的中点,沿 ${DE}$ 将 $\bigtriangleup {ADE}$ 折起,使得点 $A$ 到点 $P$ 位置,且 ${PE} \bot  {EB},M$ 为 ${PB}$ 的中点, $N$ 是 ${BC}$ 上的动点 (与点 $B\text{ 、 }C$ 不重合).

(1)求证:平面 ${EMN} \bot$ 平面 ${PBC}$ ；

(2)是否存在点 $N$ ，使得二面角 $B - {EN} - M$ 的余弦值为 $\frac{\sqrt{6}}{6}$ ？若存在，确定 $N$ 点位置；若不存在，说明理由.

\begin{center}
\includegraphics[width=0.5\textwidth]{images/bo_d7tk7oc91nqc73875fu0_184_823_409_664_220_0.jpg}
\end{center}
 

(例 5 图)

\section*{四基达标精练}

\section*{一、填空题}

1. 已知向量 $\overrightarrow{a} = \left( {1,0,2}\right) ,\overrightarrow{b} = \left( {2,1,0}\right)$ ，则 $\overrightarrow{a}$ 与 $\overrightarrow{b}$ 的夹角为\_\_\_\_\_.

2. 在空间直角坐标系中，设点 $A\left( {1, - 3,\sqrt{6}}\right)$ 关于原点的对称点是 $B$ ，则 $\left| {AB}\right|  =$ \_\_\_\_\_.

3. 已知正方体 ${ABCD} - {A}_{1}{B}_{1}{C}_{1}{D}_{1}$ 的棱长为 $a$ ，则 $\overrightarrow{AB} \cdot  \overrightarrow{{A}_{1}{C}_{1}} =$ \_\_\_\_\_.

4. 已知平面 $\alpha$ 的一个法向量是 $\overrightarrow{n} = \left( {1, - 1,2}\right)$ ，且点 $A\left( {0,3,1}\right)$ 在平面 $\alpha$ 上，若 $P\left( {x,y,z}\right)$ 是平面 $\alpha$ 上任意一点,则点 $P$ 的坐标满足\_\_\_\_\_.

5. 已知 $\overrightarrow{a} = \left( {1 - t,{2t} - 1,0}\right) ,\overrightarrow{b} = \left( {3,t,t}\right)$ ,则 $\left| {\overrightarrow{b} - \overrightarrow{a}}\right|$ 的最小值为\_\_\_\_\_.

6. 在空间直角坐标系 $O - {xyz}$ 中, $\overrightarrow{a} = \left( {x - 2,{2y}, - 2}\right) ,\overrightarrow{b} = \left( {{3x},{2y}, - {3x}}\right)$ ,且 $\overrightarrow{a} \cdot  \overrightarrow{b} = {12}$ ,则 $m = {x}^{2} + {y}^{2}  + {2x}$ 的最大值是\_\_\_\_\_.

7. 平行六面体 ${ABCD} - {A}_{1}{B}_{1}{C}_{1}{D}_{1}$ 中, $\angle {BA}{A}_{1} = \angle {BAD} = \angle {DA}{A}_{1} = {60}^{ \circ  },{AB} = 3,{AD} = 4,A{A}_{1} = 5$ , 则 $A{C}_{1} =$ \_\_\_\_\_.

8. 在空间直角坐标系 $O - {xyz}$ 中,点 $A\left( {3,4,5}\right)$ 在平面 ${xOy}\text{ 、 }{xOz}$ 内的射影分别为点 $B\text{ 、 }C$ ,则 $\left| \overrightarrow{BC}\right|  =$

9. 已知棱长为 1 的正方体 ${ABCD} - {A}_{1}{B}_{1}{C}_{1}{D}_{1}$ 中， $E$ 为侧面 $B{B}_{1}{C}_{1}C$ 中心， $F$ 在棱 ${AD}$ 上运动，正方体表面上有一点 $P$ 满足 $\overrightarrow{{D}_{1}P} = x\overrightarrow{{D}_{1}F} + y\overrightarrow{{D}_{1}E}\left( {x \geq  0,y \geq  0}\right)$ ,则所有满足条件的点 $P$ 构成图形的面积为\_\_\_\_\_.

\section*{二、选择题}

D. $\overrightarrow{PA} \cdot  \overrightarrow{PC}$ .

( ) 10. 在四棱锥 $P - {ABCD}$ 中, ${PA} \bot$ 平面 ${ABCD}$ ,底面 ${ABCD}$ 是正方形, ${PA} = {AB}$ ,则下列数量积最大的是

A. $\overrightarrow{BD} \cdot  \overrightarrow{PC}$ . B. $\overrightarrow{PB} \cdot  \overrightarrow{PC}$ . C. $\overrightarrow{BC} \cdot  \overrightarrow{PC}$ .

11. 如图, $M\text{ 、 }N$ 分别是四面体 ${OABC}$ 的边 ${OA}\text{ 、 }{BC}$ 的中点, $P$ 是 ${MN}$ 的中点,设 $\overrightarrow{OA} = \overrightarrow{a},\overrightarrow{OB} = \overrightarrow{b},\overrightarrow{OC} =  \overrightarrow{c}$ ,用 $\overrightarrow{a}\text{ 、 }\overrightarrow{b}\text{ 、 }\overrightarrow{c}$ 表示 $\overrightarrow{OP}$ ,则 ( )

A. $\overrightarrow{OP} = \frac{1}{2}\overrightarrow{a} + \frac{1}{3}\overrightarrow{b} + \frac{1}{4}\overrightarrow{c}$ . B. $\overrightarrow{OP} = \frac{1}{2}\overrightarrow{a} + \frac{1}{4}\overrightarrow{b} + \frac{1}{4}\overrightarrow{c}$ .

C. $\overrightarrow{OP} = \frac{1}{3}\overrightarrow{a} + \frac{1}{2}\overrightarrow{b} + \frac{1}{4}\overrightarrow{c}$ . D. $\overrightarrow{OP} = \frac{1}{4}\overrightarrow{a} + \frac{1}{4}\overrightarrow{b} + \frac{1}{4}\overrightarrow{c}$ .

\begin{center}
\includegraphics[width=0.2\textwidth]{images/bo_d7tk7oc91nqc73875fu0_185_264_422_290_281_0.jpg}
\end{center}
 

(第 11 题图)

\begin{center}
\includegraphics[width=0.3\textwidth]{images/bo_d7tk7oc91nqc73875fu0_185_884_470_378_232_0.jpg}
\end{center}
 

(第 12 题图)

12. 如图,四个棱长为 1 的正方体排成一个正四棱柱, ${AB}$ 是一条侧棱, ${P}_{i}\left( {i = 1,2,\cdots ,8}\right)$ 是上底面上其余的八个点,则 $\overrightarrow{AB} \cdot  \overrightarrow{A{P}_{i}}\left( {i = 1,2,\cdots ,8}\right)$ 的不同值的个数为 ( )

A.1. B.2. C.3. D.4.

\section*{三、解答题}

13. 已知空间向量 $\overrightarrow{a} = \left( {x,4,1}\right) ,\overrightarrow{b} = \left( {-2,y, - 1}\right) ,\overrightarrow{c} = \left( {3, - 2,z}\right) ,\overrightarrow{a}//\overrightarrow{b},\overrightarrow{b} \bot  \overrightarrow{c}$ .

(1)求 $x$ 、 $y$ 、 $z$ 的值；

(2) $\left( {\overrightarrow{a} + \overrightarrow{c}}\right)$ 与 $\left( {\overrightarrow{b} + \overrightarrow{c}}\right)$ 所成的角.

14. 已知空间三点 $A\left( {-2,0,2}\right)$ 、 $B\left( {-1,1,2}\right)$ 、 $C\left( {-3,0,4}\right)$ ，设 $\overrightarrow{a} = \overrightarrow{AB}$ ， $\overrightarrow{b} = \overrightarrow{AC}$ .

(1)求 $\overrightarrow{a}$ 和 $\overrightarrow{b}$ 的夹角的余弦值；

(2)若向量 $k\overrightarrow{a} + \overrightarrow{b}$ 与 $k\overrightarrow{a} - 2\overrightarrow{b}$ 互相垂直，求实数 $k$ 的值.

15. 已知空间三点 $A\left( {0,2,3}\right)$ 、 $B\left( {-2,1,6}\right)$ 、 $C\left( {1, - 1,5}\right)$ .

(1)求 $\bigtriangleup  {ABC}$ 的面积；

(2)若向量 $\overrightarrow{CD}//\overrightarrow{AB}$ ，且 $\left| \overrightarrow{CD}\right|  = \sqrt{21}$ ，求向量 $\overrightarrow{CD}$ 的坐标.

16. 一个几何体由一个直三棱柱 ${ABC} - {A}_{1}{B}_{1}{C}_{1}$ 切割而成，如图，已知 ${AB}\bot {BC},{AB} = {BC} = 4,B{B}_{1} =  {2AN} = 8$ .

(1)求证: ${B}_{1}N \bot  {CN}$ ；

(2)设 $M$ 为 ${AB}$ 的中点，在棱 ${BC}$ 上是否存在一点 $P$ ，使 ${MP}//$ 平面 ${B}_{1}{CN}$ ？若存在，求 $\frac{BP}{PC}$ 的值；若不存在, 请说明理由.

\begin{center}
\includegraphics[width=0.2\textwidth]{images/bo_d7tk7oc91nqc73875fu0_186_1183_1180_309_222_0.jpg}
\end{center}
 

(第 16 题图)

\section*{$\text{ § }{9.5}$ 空间角和距离的度量}

\section*{典型例题精讲}

【例 1】在三棱锥 $A - {BCD}$ 中, ${AB} = {AD} = {BD} = 2,{BC} = {DC} = \sqrt{2},{AC} = 2$ .

(1)求证: ${BD}\bot {AC}$ ；

(2)若 $P$ 为 ${AC}$ 上一点，且 ${AP} = \frac{3}{4}{AC}$ ，求直线 ${BP}$ 与平面 ${ACD}$ 所成角的正弦值.

\begin{center}
\includegraphics[width=0.2\textwidth]{images/bo_d7tk7oc91nqc73875fu0_187_1177_608_273_326_0.jpg}
\end{center}
 

(例 1 图)

【例 2】如图,平行六面体 ${ABCD} - {A}_{1}{B}_{1}{C}_{1}{D}_{1}$ 中, $\angle {BAD} = \angle {DA}{A}_{1} = \angle {BA}{A}_{1} = {60}^{ \circ  },A{A}_{1} = 2$ , ${AB} = {AD} = 1$ . 求:

(1) $A{C}_{1}$ 的长；

(2)异面直线 $B{D}_{1}$ 与 ${AC}$ 所成的角的大小.

\begin{center}
\includegraphics[width=0.3\textwidth]{images/bo_d7tk7oc91nqc73875fu0_187_1023_1397_423_292_0.jpg}
\end{center}
 

(例 2 图)

【例 3】如图,由直三棱柱 ${ABC} - {A}_{1}{B}_{1}{C}_{1}$ 和四棱锥 $D - B{B}_{1}{C}_{1}C$ 构成的几何体中, $\angle {BAC} = {90}^{ \circ  }$ , ${AB} = 1,{BC} = B{B}_{1} = 2,{C}_{1}D = {CD} = \sqrt{5}$ ,平面 $C{C}_{1}D \bot$ 平面 ${AC}{C}_{1}{A}_{1}$ .

(1)求证: ${AC}\bot {DC}{\text{ ， }}_{1}$ ；

(2)若 $M$ 为 $D{C}_{1}$ 的中点，求证: ${AM}//$ 平面 ${DB}{B}_{1}$ .

\begin{center}
\includegraphics[width=0.2\textwidth]{images/bo_d7tk7oc91nqc73875fu0_188_1189_318_306_270_0.jpg}
\end{center}
 

(例 3 图)

【例 4】如图,在四棱锥 $S - {ABCD}$ 中底面 ${ABCD}$ 是直角梯形,侧棱 ${SA} \bot$ 底面 ${ABCD},{AB}$ 垂直于 ${AD}$ 和 ${BC},M$ 为棱 ${SB}$ 上的点, ${SM} = {2MB},{SA} = {AB} = {BC} = 2,{AD} = 1$ .

(1)求平面 ${AMC}$ 与平面 ${SAB}$ 所成的锐二面角的余弦值；

(2)设点 $N$ 是线段 ${CD}$ 上的动点， ${MN}$ 与平面 ${SAB}$ 所成的角为 $\theta$ ，求当 $\sin \theta$ 取最大值时点 $N$ 的位置.

\begin{center}
\includegraphics[width=0.2\textwidth]{images/bo_d7tk7oc91nqc73875fu0_188_1213_1204_295_308_0.jpg}
\end{center}
 

(例 4 图)

【例 5】如图所示,在直三棱柱 ${ABC} - {A}_{1}{B}_{1}{C}_{1}$ 中, ${AB} = {BC} = a,{A}_{1}A = {AC} = \sqrt{2}a,D$ 为 $B{B}_{1}$ 的中点.

(1)证明:平面 ${AD}{C}_{1} \bot$ 平面 ${AC}{C}_{1}{A}_{1}$ ；

(2)求点 $B$ 到平面 ${AD}{C}_{1}$ 的距离；

(3)求平面 ${AD}{C}_{1}$ 与平面 ${ABC}$ 所成的二面角大小.

\begin{center}
\includegraphics[width=0.2\textwidth]{images/bo_d7tk7oc91nqc73875fu0_189_1151_326_286_304_0.jpg}
\end{center}
 

(例 5 图)

\section*{四基达标精练}

\section*{一、填空题}

1. 一条长度为 8 的线段 ${AB}$ 的两个端点到平面 $\alpha$ 的距离分别为 2 和 8,则该线段的中点到这个平面的距离是\_\_\_\_\_.

2. 若球心到截面距离为 $d$ ,球的半径为 $R$ ,截面半径为 $r$ ,则它们的关系是\_\_\_\_\_.

3. 长方体 ${ABCD} - {A}_{1}{B}_{1}{C}_{1}{D}_{1}$ 的一条对角线 $A{C}_{1}$ 与从 $A$ 出发的三条棱所成的锐角分别是 $\alpha$ 、 $\beta$ 、 $\gamma$ ,则 ${\sin }^{2}\alpha  + {\sin }^{2}\beta  + {\sin }^{2}\gamma  =$ \_\_\_\_\_.

4. 在空间四边形 ${ABCD}$ 中， $E$ 、 $H$ 分别是 ${AB}$ 、 ${AD}$ 的中点， $F$ 、 $G$ 分别是 ${CB}$ 、 ${CD}$ 上的点，且 $\frac{CF}{CB} = \frac{CG}{CD} =  \frac{2}{3}$ ，若 ${BD} = 6$ ，梯形 ${EFGH}$ 的面积是 28，则平行线 ${EH}$ 与 ${FG}$ 间的距离是\_\_\_\_\_.

5. 已知正六边形 ${ABCDEF}$ 边长为 $a,{PA}\bot$ 平面 ${AC},{PA} = a$ ，则点 $P$ 到直线 ${CD}$ 的距离是\_\_\_\_\_.

6. 已知长方体 ${ABCD} - {A}_{1}{B}_{1}{C}_{1}{D}_{1}$ 中，棱 $A{A}_{1} = 3,{AB} = 4$ ，那么直线 ${BC}$ 与平面 ${A}_{1}{B}_{1}{C}_{1}{D}_{1}$ 的距离是\_\_\_\_\_.

7.过正方形 ${ABCD}$ 的顶点 $A$ 作 ${PA} \bot$ 平面 ${ABCD}$ ,若 ${AB} = {PA}$ ,则平面 ${PAB}$ 与平面 ${PCD}$ 所成的角为\_\_\_\_\_.

8. ${{Rt}} \bigtriangleup  {ABC}$ 的斜边 ${AB}$ 在平面 $\alpha$ 内, ${AC}\text{ 、 }{BC}$ 与平面 $\alpha$ 分别成 ${45}^{ \circ  }$ 和 ${30}^{ \circ  }$ 角,则平面 ${ABC}$ 和平面 $\alpha$ 所成的角等于\_\_\_\_\_.

9. 若 $\angle {BAC} = {120}^{ \circ  },{BC} \subset  \alpha ,A \notin  \alpha$ ，则 $\angle {BAC}$ 在 $\alpha$ 内的影角 $\angle {B{A}^{\prime }}C$ 的取值范围是\_\_\_\_\_.

\section*{二、选择题}

10. 在正方体 ${ABCD} - {A}_{1}{B}_{1}{C}_{1}{D}_{1}$ 中, $M$ 为 $D{D}_{1}$ 的中点, $O$ 为底面 ${ABCD}$ 的中心, $P$ 为棱 ${A}_{1}{B}_{1}$ 上任意一点,则直线 ${OP}$ 与直线 ${AM}$ 所成的角是 ( )

A. $\frac{\pi }{6}$ . B. $\frac{\pi }{4}$ . C. $\frac{\pi }{3}$ . D. $\frac{\pi }{2}$ .

11. 如果空间四边形的对角线相等, 那么顺次连接它的各边中点所成的四边形一定是 ( )

A. 梯形. B.长方形. C. 菱形. D. 正方形.

12. 若直角三角形的斜边与平面 $M$ 平行,两直角边所在的直线与平面 $M$ 所成的角分别是 $\alpha$ 和 $\beta$ ,则

( )

A. ${\sin }^{2}\alpha  + {\sin }^{2}\beta  \geq  1$ . B. ${\sin }^{2}\alpha  + {\sin }^{2}\beta  \leq  1$ . C. ${\sin }^{2}\alpha  + {\sin }^{2}\beta  > 1$ . D. ${\sin }^{2}\alpha  + {\sin }^{2}\beta  < 1$ .

\section*{三、解答题}

13. 已知 ${OA}$ 、 ${OB}$ 、 ${OC}$ 两两垂直， ${OA} = {OC} = 3,{OB} = 2,M$ 为 ${OB}$ 的中点，点 $N$ 在 ${AC}$ 上， ${AN} = {2NC}$ .

(1)求 ${MN}$ 的长；

(2)若点 $P$ 在线段 ${BC}$ 上，设 $\frac{BP}{PC} = \lambda$ ，当 ${AP}\bot {MN}$ 时，求实数 $\lambda$ 的值.

\begin{center}
\includegraphics[width=0.2\textwidth]{images/bo_d7tk7oc91nqc73875fu0_190_1197_818_301_305_0.jpg}
\end{center}
 

(第 13 题图)

14. 如图，在正方体 ${ABCD} - {A}_{1}{B}_{1}{C}_{1}{D}_{1}$ 中，点 $E$ 在 ${BD}$ 上，且 ${BE} = \frac{1}{3}{BD}$ ，点 $F$ 在 $C{B}_{1}$ 上，且 ${CF} = \frac{1}{3}{C{B}_{1}}$ . 求证:

(1) ${EF}\bot {BD}$ ；

(2) ${EF}\bot C{B}_{1}$ .

\begin{center}
\includegraphics[width=0.2\textwidth]{images/bo_d7tk7oc91nqc73875fu0_190_1227_1599_281_242_0.jpg}
\end{center}
 

(第 14 题图)


15. 如图,在棱长为 1 的正方体 ${ABCD} - {A}_{1}{B}_{1}{C}_{1}{D}_{1}$ 中, $P$ 是侧棱 $C{C}_{1}$ 上一点, ${CP} = m$ . 试确定 $m$ 的值, 使得直线 ${AP}$ 与平面 ${BD}{D}_{1}{B}_{1}$ 所成的角为 ${60}^{ \circ  }$ .

\begin{center}
\includegraphics[width=0.2\textwidth]{images/bo_d7tk7oc91nqc73875fu0_191_1137_244_310_289_0.jpg}
\end{center}
 

(第 15 题图)

16. 如图,直三棱柱 ${ABC} - {A}_{1}{B}_{1}{C}_{1}$ 中, ${AB} \bot  {AC},{AB} = {AC} = 2,A{A}_{1} = 4$ ,点 $D$ 是 ${BC}$ 的中点.

(1)求异面直线 ${A}_{1}B$ 与 ${C}_{1}D$ 所成角的大小；

(2)求点 ${B}_{1}$ 到平面 ${AD}{C}_{1}$ 的距离.

\begin{center}
\includegraphics[width=0.2\textwidth]{images/bo_d7tk7oc91nqc73875fu0_191_1163_1073_284_290_0.jpg}
\end{center}
 

(第 16 题图)

\section*{第 10 章 解析几何(选择性必修 1)}

\section*{知识梳理*}

\begin{center}
\includegraphics[width=0.2\textwidth]{images/bo_d7tk7oc91nqc73875fu0_192_1260_442_232_266_0.jpg}
\end{center}
 

\section*{一、直线的方程与直线的位置关系}

1. 直线的斜率 $k$ 、倾斜角 $\alpha$ 之间的对应关系是:

若 $\alpha  \neq  \frac{\pi }{2}$ ,则 $k = \tan \alpha$ ; 若 $\alpha  = \frac{\pi }{2}$ ,则 $k$ 不存在.

若 $k \geq  0$ ,则 $\alpha  = \arctan k$ ; 若 $k < 0$ ,则 $\alpha  = \pi  - \arctan \left( {-k}\right)$ .

设直线的方向向量为 $\overrightarrow{d} = \left( {u,v}\right)$ ,法向量是 $\overrightarrow{n} = \left( {a,b}\right)$ .

\begin{center}
\begin{tabular}{|c|c|c|c|}
\hline
\phantom{X} & 直线方程的形式 & 条件 & 适用范围 \\
\cline{1-4}
点斜式 & $y - {y}_{0} = k\left( {x - {x}_{0}}\right)$ & $k$ 和点 $P\left( {{x}_{0},{y}_{0}}\right)$ & 不垂直于 $x$ 轴 \\
\cline{1-4}
两点式 & $\frac{x - {x}_{1}}{{x}_{2} - {x}_{1}} = \frac{y - {y}_{1}}{{y}_{2} - {y}_{1}}$ & 点 $P\left( {{x}_{1},{y}_{1}}\right)$ 和 $Q\left( {{x}_{2},{y}_{2}}\right)$ & 不垂直于坐标轴 \\
\cline{1-4}
斜截式 & $y = {kx} + b$ & $k$ 和在 $y$ 轴上的截距 & 不垂直于 $x$ 轴 \\
\cline{1-4}
截距式 & $\frac{x}{a} + \frac{y}{b} = 1$ & 在 $x\text{ 、 }y$ 轴上的截距 $a\text{ 、 }b$ & $a$ 、 $b$ 存在，且都不为 0 \\
\cline{1-4}
* 点方向式 & $\frac{x - {x}_{0}}{u} = \frac{y - {y}_{0}}{v}$ & 方向向量 $\overrightarrow{d}$ 和点 $P\left( {{x}_{0},{y}_{0}}\right)$ & ${uv} \neq  0$ \\
\cline{1-4}
* 点法向式 & $a\left( {x - {x}_{0}}\right)  + b\left( {y - {y}_{0}}\right)  = 0$ & 法向量 $\overrightarrow{n}$ 和点 $P\left( {{x}_{0},{y}_{0}}\right)$ & 任意情形 \\
\cline{1-4}
一般式 & ${ax} + {by} + c = 0\left( {{a}^{2} + {b}^{2} \neq  0}\right)$ & 系数 $a$ 、 $b$ 、 $c$ & 任意情形 \\
\cline{1-4}
\hline
\end{tabular}
\end{center}

2. 直线 $l : y = {kx} + b$ 上两点 ${P}_{1}\left( {{x}_{1},{y}_{1}}\right)$ 和 ${P}_{2}\left( {{x}_{2},{y}_{2}}\right)$ 间的距离:

$\left| {{P}_{1}{P}_{2}}\right|  = \sqrt{{\left( {x}_{2} - {x}_{1}\right) }^{2} + {\left( {y}_{2} - {y}_{1}\right) }^{2}} = \sqrt{1 + {k}^{2}}\left| {{x}_{1} - {x}_{2}}\right|  = \sqrt{1 + {\left( \frac{1}{k}\right) }^{2}}\left| {{y}_{1} - {y}_{2}}\right| .$

3. 设直角坐标平面内直线 $l : {ax} + {by} + c = 0$ 和直线外一点 $P\left( {{x}_{0},{y}_{0}}\right)$ ，则点 $P$ 到直线 $l$ 的距离 $d = \frac{\left| a{x}_{0} + b{y}_{0} + c\right| }{\sqrt{{a}^{2} + {b}^{2}}}.$

4. 设直角坐标平面内直线 ${l}_{1} : {Ax} + {By} + {C}_{1} = 0$ 和直线 ${l}_{2} : {Ax} + {By} + {C}_{2} = 0\left( {{C}_{1} \neq  {C}_{2}}\right)$ ，且 ${l}_{1}//{l}_{2}$ ，则 ${l}_{1}$ 和 ${l}_{2}$ 的距离 $d = \frac{\left| {C}_{1} - {C}_{2}\right| }{\sqrt{{A}^{2} + {B}^{2}}}$ .

5. 当点 $P\left( {{x}_{0},{y}_{0}}\right)$ 位于直线 $l : {ax} + {by} + c = 0$ 的法向量 $\overrightarrow{n} = \left( {a,b}\right)$ 指向的一侧时， $\delta  = \frac{a{x}_{0} + b{y}_{0} + c}{\sqrt{{a}^{2} + {b}^{2}}} =  d > 0$ ; 位于 $\overrightarrow{n} = \left( {a,b}\right)$ 指向的另一侧时, $\delta  = \frac{a{x}_{0} + b{y}_{0} + c}{\sqrt{{a}^{2} + {b}^{2}}} =  - d < 0$ ; 位于直线上时, $\delta  = \frac{a{x}_{0} + b{y}_{0} + c}{\sqrt{{a}^{2} + {b}^{2}}} = 0$ .

6. 两直线的夹角:在直角坐标平面内直线 ${l}_{1} : {a}_{1}x + {b}_{1}y + {c}_{1} = 0\left( {{a}_{1}^{2} + {b}_{1}^{2} \neq  0}\right)$ 和直线 ${l}_{2} : {a}_{2}x + {b}_{2}y + {c}_{2} =  0\left( {{a}_{2}^{2} + {b}_{2}^{2} \neq  0}\right)$ 的夹角为 $\alpha$ ,则 $\cos \alpha  = \frac{\left| {a}_{1}{a}_{2} + {b}_{1}{b}_{2}\right| }{\sqrt{{a}_{1}^{2} + {b}_{1}^{2}}\sqrt{{a}_{2}^{2} + {b}_{2}^{2}}}$ . 若直线 ${l}_{1}$ 和 ${l}_{2}$ 的斜率分别是 ${k}_{1}$ 、 ${k}_{2}$ ，则 ${l}_{1}$ 和 ${l}_{2}$ 的夹角 $\alpha$ 可由 $\tan \alpha  = \left| \frac{{k}_{2} - {k}_{1}}{1 + {k}_{2}{k}_{1}}\right|$ 确定；

若直线 ${l}_{1}$ 和 ${l}_{2}$ 的方向向量分别是 ${\overrightarrow{d}}_{1}\text{ 、 }{\overrightarrow{d}}_{2}$ ,则 ${l}_{1}$ 和 ${l}_{2}$ 的夹角 $\alpha$ 可由 $\cos \alpha  = \frac{\left| {\overrightarrow{d}}_{1} \cdot  {\overrightarrow{d}}_{2}\right| }{\left| {\overrightarrow{d}}_{1}\right|  \cdot  \left| {\overrightarrow{d}}_{2}\right| }$ 确定;

若直线 ${l}_{1}$ 和 ${l}_{2}$ 的法向量分别是 ${\overrightarrow{n}}_{1},{\overrightarrow{n}}_{2}$ ,则 ${l}_{1}$ 和 ${l}_{2}$ 的夹角 $\alpha$ 可由 $\cos \alpha  = \frac{\left| {\overrightarrow{n}}_{1} \cdot  {\overrightarrow{n}}_{2}\right| }{\left| {\overrightarrow{n}}_{1}\right|  \cdot  \left| {\overrightarrow{n}}_{2}\right| }$ 确定;

直线 ${l}_{1}$ 和 ${l}_{2}$ 的夹角 $\alpha$ 的取值范围是 $\left\lbrack  {0,\frac{\pi }{2}}\right\rbrack$ .

7. 过两直线 ${l}_{1} : {A}_{1}x + {B}_{1}y + {C}_{1} = 0,{l}_{2} : {A}_{2}x + {B}_{2}y + {C}_{2} = 0$ 交点的直线系方程: ${A}_{1}x + {B}_{1}y + {C}_{1} +  \lambda \left( {{A}_{2}x + {B}_{2}y + {C}_{2}}\right)  = 0\left( {\text{ 除去 }{l}_{2}}\right)$ .

\section*{二、圆的方程、直线与圆、圆与圆的位置关系}

8. 圆的标准方程: ${x}^{2} + {y}^{2} = {r}^{2}$ (圆心(0，0)，半径 $r$ )；

\[
{\left( x - a\right) }^{2} + {\left( y - b\right) }^{2} = {r}^{2}\text{ (圆心 }\left( {a,b}\right) \text{ ,半径 }r\text{ ). }
\]

9. 圆的一般方程: ${x}^{2} + {y}^{2} + {Dx} + {Ey} + F = 0\left( {{D}^{2} + {E}^{2} - {4F} > 0}\right)$ .

10. 直线与圆(圆与圆)的位置关系:

设圆 $C : {\left( x - a\right) }^{2} + {\left( y - b\right) }^{2} = {r}^{2}$ 的圆心 $C$ 到直线 ${Ax} + {By} + C = 0$ 的距离为 $d$ ；

圆 ${C}_{1} : {\left( x - {a}_{1}\right) }^{2} + {\left( y - {b}_{1}\right) }^{2} = {r}_{1}^{2}$ 与圆 ${C}_{2} : {\left( x - {a}_{2}\right) }^{2} + {\left( y - {b}_{2}\right) }^{2} = {r}_{2}^{2}$ 的圆心间的距离为 ${d}^{\prime }$ ,则

\begin{center}
\begin{tabular}{|c|c|c|}
\hline
\multicolumn{3}{|c|}{直线与圆的位置关系} \\
\cline{1-3}
关系 & 特征 & 表达式 \\
\cline{1-3}
相离 & 无公共点 & $d > r$ \\
\cline{1-3}
相切 & 唯一公共点 & $d = r$ \\
\cline{1-3}
相交 & 两个不同公共点 & $d < r$ \\
\cline{1-3}
\multicolumn{3}{|c|}{也可以利用联立方程组解的个数判断} \\
\cline{1-3}
\hline
\end{tabular}
\end{center}

\begin{center}
\begin{tabular}{|c|c|c|}
\hline
\multicolumn{3}{|c|}{圆与圆的位置关系} \\
\cline{1-3}
关系 & 特征 & 表达式 \\
\cline{1-3}
外离 & 无公共点 & ${d}^{\prime } > {r}_{1} + {r}_{2}$ \\
\cline{1-3}
外切 & 唯一公共点 & ${d}^{\prime } = {r}_{1} + {r}_{2}$ \\
\cline{1-3}
相交 & 两个不同公共点 & $\left| {{r}_{1} - {r}_{2}}\right|  < {d}^{\prime } < {r}_{1} + {r}_{2}$ \\
\cline{1-3}
内切 ${r}_{1} \neq  {r}_{2}$ & 唯一公共点 & ${d}^{\prime } = \left| {{r}_{1} - {r}_{2}}\right|$ \\
\cline{1-3}
内含 ${r}_{1} \neq  {r}_{2}$ & 无公共点 & ${d}^{\prime } < \left| {{r}_{1} - {r}_{2}}\right|$ \\
\cline{1-3}
\hline
\end{tabular}
\end{center}

11. 过两圆 ${C}_{1} : {x}^{2} + {y}^{2} + {D}_{1}x + {E}_{1}y + {F}_{1} = 0,{C}_{2} : {x}^{2} + {y}^{2} + {D}_{2}x + {E}_{2}y + {F}_{2} = 0$ 交点的圆系方程: ${x}^{2} +  {y}^{2} + {D}_{1}x + {E}_{1}y + {F}_{1} + \lambda \left( {{x}^{2} + {y}^{2} + {D}_{2}x + {E}_{2}y + {F}_{2}}\right)  = 0$ (除去 ${C}_{2}$ ).

12. 过圆 ${x}^{2} + {y}^{2} = {r}^{2}$ 上一点 $P\left( {{x}_{0},{y}_{0}}\right)$ 的圆的切线方程为 ${x}_{0}x + {y}_{0}y = {r}^{2}$ .

13. 以 $A\left( {{x}_{1},{y}_{1}}\right)$ 、 $B\left( {{x}_{2},{y}_{2}}\right)$ 为直径端点的圆的方程 $\left( {x - {x}_{1}}\right) \left( {x - {x}_{2}}\right)  + \left( {y - {y}_{1}}\right) \left( {y - {y}_{2}}\right)  = 0$ .

14. 用圆的切线到圆心的距离等于圆的半径可求与切线或圆有关的参数 (如斜率、半径等).

15. 在圆中, 注意利用半径、半弦长, 以及弦心距构成的直角三角形.

\section*{三、圆锥曲线、直线与圆锥曲线的位置关系}

16. 圆锥曲线及其性质:

\begin{center}
\begin{tabular}{|c|c|c|}
\hline
椭圆定义 & \multicolumn{2}{|c|}{$\left| {M{F}_{1}}\right|  + \left| {M{F}_{2}}\right|  = {2a}\left( {{2a} > \left| {{F}_{1}{F}_{2}}\right| }\right)$} \\
\cline{1-3}
椭圆标准方程 & $\frac{{x}^{2}}{{a}^{2}} + \frac{{y}^{2}}{{b}^{2}} = 1\left( {a > b > 0}\right)$ & $\frac{{y}^{2}}{{a}^{2}} + \frac{{x}^{2}}{{b}^{2}} = 1\left( {a > b > 0}\right)$ \\
\cline{1-3}
参数 & ${a}^{2} = {b}^{2} + {c}^{2}$ & ${a}^{2} = {b}^{2} + {c}^{2}$ \\
\cline{1-3}
范围 & $\left| x\right|  \leq  a,\left| y\right|  \leq  b$ & $\left| y\right|  \leq  a,\left| x\right|  \leq  b$ \\
\cline{1-3}
对称性 & 轴: $x = 0\text{ 、 }y = 0$ ; 中心: $\left( {0,0}\right)$ & 轴: $x = 0\text{ 、 }y = 0$ ; 中心: $\left( {0,0}\right)$ \\
\cline{1-3}
顶点 & ${A}_{1,2}\left( {\pm a,0}\right)$ 和 ${B}_{1,2}\left( {0, \pm  b}\right)$ & ${A}_{1,2}\left( {0, \pm  a}\right)$ 和 ${B}_{1,2}\left( {\pm b,0}\right)$ \\
\cline{1-3}
离心率 $e \in  \left( {0,1}\right)$ & $\frac{c}{a}$ & $\frac{c}{a}$ \\
\cline{1-3}
\hline
\end{tabular}
\end{center}

\begin{center}
\begin{tabular}{|c|c|c|}
\hline
双曲线定义 & \multicolumn{2}{|c|}{$\left. {\left. \right| M{F}_{1}\left| -\right| M{F}_{2}\left| \right|  = {2a}\left( {a > 0,{2a} < }\right. \left| {{F}_{1}{F}_{2}}\right| }\right)$} \\
\cline{1-3}
双曲线标准方程 & $\frac{{x}^{2}}{{a}^{2}} - \frac{{y}^{2}}{{b}^{2}} = 1\left( {a,b > 0}\right)$ & $\frac{{y}^{2}}{{a}^{2}} - \frac{{x}^{2}}{{b}^{2}} = 1\left( {a\text{ 、 }b > 0}\right)$ \\
\cline{1-3}
参数 & ${c}^{2} = {a}^{2} + {b}^{2}$ & ${c}^{2} = {a}^{2} + {b}^{2}$ \\
\cline{1-3}
范围 & $\left| x\right|  \geq  a$ & $\left| y\right|  \geq  a$ \\
\cline{1-3}
对称性 & 轴: $x = 0\text{ 、 }y = 0$ ; 中心: $\left( {0,0}\right)$ & 轴: $x = 0\text{ 、 }y = 0$ ; 中心: $\left( {0,0}\right)$ \\
\cline{1-3}
顶点 & ${A}_{1.2}\left( {\pm a,0}\right)$ & ${A}_{1,2}\left( {0, \pm  a}\right)$ \\
\cline{1-3}
渐近线 & $\frac{{x}^{2}}{{a}^{2}} - \frac{{y}^{2}}{{b}^{2}} = 0$ & $\frac{{y}^{2}}{{a}^{2}} - \frac{{x}^{2}}{{b}^{2}} = 0$ \\
\cline{1-3}
离心率 $e \in  \left( {1, + \infty }\right)$ & $\frac{c}{a}$ & $\frac{c}{a}$ \\
\cline{1-3}
\hline
\end{tabular}
\end{center}

 \ding{72} 椭圆的定义中距离和要大于两定点之间的距离，双曲线定义中距离之差的绝对值要小于两定点之间的距离，解有关题目要能联想到这两个定义；

 \ding{72}  椭圆和双曲线中 $a$ 、 $b$ 、 $c$ 、抛物线中 $p$ 的几何意义在解题过程中有重要作用.

\begin{center}
\begin{tabular}{|c|c|c|c|c|}
\hline
抛物线定义 & \multicolumn{4}{|c|}{$\left| {MF}\right|  = d$ ( $d$ 为 $M$ 到定直线的距离)} \\
\cline{1-5}
抛物线标准方程 & ${y}^{2} = {2px}$ & ${y}^{2} =  - {2px}$ & ${x}^{2} = {2py}$ & ${x}^{2} =  - {2py}$ \\
\cline{1-5}
焦参数 & \multicolumn{4}{|c|}{$p\left( {p > 0}\right)$ 为定点到定直线的距离} \\
\cline{1-5}
范围 & $x \geq  0$ & $x \leq  0$ & $y \geq  0$ & $y \leq  0$ \\
\cline{1-5}
对称性 & 轴: $y = 0$ & 轴: $y = 0$ & 轴: $x = 0$ & 轴: $x = 0$ \\
\cline{1-5}
顶点 & $O\left( {0,0}\right)$ & $O\left( {0,0}\right)$ & $O\left( {0,0}\right)$ & $O\left( {0,0}\right)$ \\
\cline{1-5}
\hline
\end{tabular}
\end{center}

17. 圆锥曲线中各几何元素间的相对位置关系:如椭圆中，注意焦点、中心、短轴端点所组成的直角三角形，双曲线的中心、顶点和虚轴端点构成的直角三角形。

18. 在用圆锥曲线与直线联立求解时,消元后得到的方程中要注意:二次项的系数是否为零，判别式 $\Delta  > 0$ 或 $\Delta  \geq  0$ 的限制(求交点，弦长，中点，斜率，对称，存在性问题都在 $\Delta  > 0$ 下进行).

19. 解析几何中的取值范围问题解法有二:  \ding{172} 建立某个变量的函数关系,求函数的值域;  \ding{173} 建立有关的等式 (组) 并解之, 有时还需要一些等量关系; 判别式常用于建立不等式, 韦达定理是常用的工具.

20. “点差法” 很有用，但要注意在解题中的局限性.

21. 椭圆的通径 (过焦点垂直于长轴的弦),长为 $\frac{2{b}^{2}}{a}$ ,它是焦点弦中长度最短的弦;

双曲线的通径 (过焦点垂直于实轴的弦),长为 $\frac{2{b}^{2}}{a}$ ,它和实轴长之一是焦点弦中长度最短的弦;

抛物线的通径 (过焦点垂直于对称轴的弦),长为 ${2p}$ 是抛物线焦点弦中长度最短的弦.

22. 点 $P$ 在椭圆上， $\bigtriangleup  P{F}_{1}{F}_{2}$ 的面积 ${b}^{2}\tan \frac{\alpha }{2}$ ；点 $P$ 在双曲线上， $\bigtriangleup  P{F}_{1}{F}_{2}$ 的面积 ${b}^{2}\cot \frac{\alpha }{2}$ ，(其中 ${F}_{1}$ 、 ${F}_{2}$ 是焦点， $\left. {\angle {F}_{1}P{F}_{2} = \alpha }\right)$ .

23. 求圆锥曲线的标准方程时, 要注意焦点所在的位置.

24. 抛物线 ${y}^{2} = {2px}\left( {p > 0}\right)$ 的焦点弦长 $\left| {AB}\right|  = \frac{2p}{{\sin }^{2}\theta } = {x}_{1} + {x}_{2} + p$ .

25. 曲线上 $A\text{ 、 }B$ 与原点的连线垂直 $\Leftrightarrow  {x}_{1}{x}_{2} + {y}_{1}{y}_{2} = 0 \Leftrightarrow$ 以线段 ${AB}$ 为直径的圆过原点.

26. 抛物线上找一点使这点到抛物线内一点和到焦点的距离之和最小, 只要作抛物线对称轴的平行线抛物线上找一点使这点到抛物线外一点和到焦点 (或准线) 的距离之和最小, 只要连结这点和焦点.

27. 求弦的中点轨迹要注意含于曲线的内部.

28. 一般弦长公式 $\left| {AB}\right|  = \sqrt{1 + {k}^{2}}\left| {{x}_{1} - {x}_{2}}\right|  = \sqrt{1 + {k}^{2}}\frac{\sqrt{\Delta }}{\left| a\right| }$ (其中 $a$ 是化简后所得一元二次方程的二次项系数; ${x}_{1}\text{ 、 }{x}_{2}$ 分别是 $A\text{ 、 }B$ 两点的横坐标; $k$ 是直线 ${AB}$ 的斜率).

29. 双曲线上任意一点到两渐近线的距离之积为常数.

30. 求抛物线的焦点坐标时要先将抛物线方程标准化.

31. ”圆的参数方程 $\left\{  {\begin{array}{l} x = a + r\cos \theta , \\  y = b + r\sin \theta  \end{array}\text{ ( }\theta }\right.$ 为参数),其中 $\theta$ 的几何意义是旋转角, $\left( {a,b}\right)$ 是圆心, $r$ 是半径.

\section*{§ 10.1 直线及其位置关系}

\section*{典型例题精讲}

【例 1】过点 $P\left( {1,2}\right)$ ,作直线 $l$ 与 $x$ 轴、 $y$ 轴正半轴分别交于 $A\text{ 、 }B$ 两点, $O$ 为坐标原点.

(1)求 $\bigtriangleup  {AOB}$ 面积的最小值及此时直线 $l$ 的方程；

(2)求直线 $l$ 在两坐标轴上截距之和的最小值及此时直线 $l$ 的方程.

【例 2】已知实数 $a$ 满足 $0 < a < 2$ ,直线 ${l}_{1} : {ax} - {2y} - {2a} + 4 = 0$ 和 ${l}_{2} : {2x} + {a}^{2}y - 2{a}^{2} - 4 = 0$ 与两坐标轴围成一个四边形.

(1)求证:无论实数 $a$ 如何变化，直线 ${l}_{1}$ 、 ${l}_{2}$ 必过定点.

(2)求实数 $a$ 取何值时，所围成的四边形面积最小?

\begin{center}
\includegraphics[width=0.2\textwidth]{images/bo_d7tk7oc91nqc73875fu0_195_1147_1591_296_239_0.jpg}
\end{center}
 

(例 2 图)

【例 3】已知三条直线 ${l}_{1} : {2x} - y + a = 0\left( {a > 0}\right)$ ，直线 ${l}_{2} :  - {4x} + {2y} + 1 = 0$ 和直线 ${l}_{3} : x + y - 1 = 0$ ，且 ${l}_{1}$ 与 ${l}_{2}$ 的距离是 $\frac{7}{10}\sqrt{5}$ .

(1) 求 $a$ 的值； $P$ 到 ${l}_{2}$ 的距离的 $\frac{1}{2}$ ； \ding{174}  点 $P$ 到 ${l}_{1}$ 的距离与点 $P$ 到 ${l}_{3}$ 的距离之比是 $\sqrt{2}$ ； $\sqrt{5}$ . 若能，求点 $P$ 坐标； 明理由.

【例 4】已知直线 $l : \left( {{2m} + 1}\right) x + \left( {m - 1}\right) y - {5m} - 1 = 0$ ,且与坐标轴形成的三角形面积为 $S$ . 求:

(1)求证:不论 $m$ 为何实数，直线 $l$ 过定点 $P$ ；

(2)分别求 $S = 3$ 和 $S = 5$ 时，所对应的直线条数；

(3)针对 $S$ 的不同取值，讨论集合 $\{ l \mid$ 直线经过 $P$ 且与坐标轴围成的三角形面积为 $S$ \}中的元系个数.

【例 5】在平面直角坐标系 ${xOy}$ 中,对于直线 $l : {ax} + {by} + c = 0$ 和点 ${P}_{1}\left( {{x}_{1},{y}_{1}}\right)$ 、 ${P}_{2}\left( {{x}_{2},{y}_{2}}\right)$ ，记 $\eta  =  \left( {a{x}_{1} + b{y}_{1} + c}\right) \left( {a{x}_{2} + b{y}_{2} + c}\right)$ ,若 $\eta  < 0$ ,则称点 ${P}_{1}\text{ 、 }{P}_{2}$ 被直线 $l$ 分隔,若曲线 $C$ 与直线 $l$ 没有公共点,且曲线 $C$ 上存在点 ${P}_{1}\text{ 、 }{P}_{2}$ 被直线 $l$ 分隔,则称直线 $l$ 为曲线 $C$ 的一条分隔线.

(1)求证:点 $A\left( {1,2}\right)$ 、 $B\left( {-1,0}\right)$ 被直线 $x + y - 1 = 0$ 分隔；

(2)若直线 $y = {kx}$ 是曲线 ${x}^{2} - 4{y}^{2} = 1$ 的分隔线，求实数 $k$ 的取值范围；

(3)动点 $M$ 到点 $Q\left( {0,2}\right)$ 的距离与到 $y$ 轴的距离之积为 1，设点 $M$ 的轨迹为曲线 $E$ ，求证:通过原点的直线中,有且仅有一条直线是 $E$ 的分隔线.

\section*{四基达标精练}

\section*{一、填空题}

1. 过点 $P\left( {0,1}\right)$ 且与两点 $A\left( {-2,0}\right) ,B\left( {2, - 3}\right)$ 距离相等的直线方程是 .

2. 已知点 $A\left( {-2, - 2}\right) \text{ 、 }B\left( {5,3}\right)$ 和 $P\left( {-4,2}\right)$ ,若过点 $P$ 的直线 $l$ 与线段 ${AB}$ 有公共点,则 $l$ 的倾斜角的取值范围是\_\_\_\_\_.

3. 已知三角形 ${ABC}$ 的三个顶点为 $A\left( {2,1}\right) \text{ 、 }B\left( {6,1}\right) \text{ 、 }C\left( {5,5}\right)$ ,则 $\angle A$ 平分线所在直线的方程为\_\_\_\_\_.

4. 若直线 $\sqrt{3}x - y - 1 = 0$ 与直线 $x - {ay} = 0$ 的夹角为 $\frac{\pi }{6}$ ，则实数 $a =$ \_\_\_\_\_.

5. 在直角坐标平面内,向量 $\overrightarrow{OA} = \left( {1,3}\right) ,\overrightarrow{OB} = \left( {-3,1}\right)$ 在直线 $l$ 上的射影长度相等,则直线 $l$ 的斜率为\_\_\_\_\_.

6. 平面直角坐标系中,给出点 $A\left( {1,0}\right) \text{ 、 }B\left( {4,0}\right)$ ,若直线 $x + {my} - 1 = 0$ 上存在点 $P$ ,使得 $\left| {PA}\right|  = 2\left| {PB}\right|$ ,则实数 $m$ 的取值范围是\_\_\_\_\_.

7. 已知 $A\left( {4,0}\right)$ 、 $B\left( {0,4}\right)$ ，从点 $P\left( {2,0}\right)$ 射出的光线经直线 ${AB}$ 反射后再射到直线 ${OB}$ 上，最后经直线 ${OB}$ 反射后又回到点 $P$ ，则光线所经过的路程为\_\_\_\_\_.

8. 在平面直角坐标系中，已知点 $P\left( {-2,2}\right)$ ，对于任意不全为零的实数 $a$ 、 $b$ ，直线 $l : {ax} + {by} + {2b} - a = 0$ ， 若点 $P$ 到直线 $l$ 的距离为 $d$ ，则 $d$ 的取值范围是\_\_\_\_\_.

9. 在平面直角坐标系中,若 $x$ 与 $y$ 都是整数,就称 $\left( {x,y}\right)$ 为整点,下列命题中:

 \ding{172} 存在这样的直线，既不与坐标轴平行又不经过任何整点；

 \ding{173} 如果 $k$ 与 $b$ 都是无理数，则直线 $y = {kx} + b$ 不经过任何整点；

 \ding{174} 直线 $l$ 经过无穷多个整点，当且仅当 $l$ 经过两个不同的整点；

 \ding{175} 直线 $y = {kx} + b$ 经过无穷多个整点的充分必要条件是: $k$ 与 $b$ 都是有理数；

 \circled{5} 存在恰经过一个整点的直线. 其中真命题的序号为\_\_\_\_\_.

\section*{二、选择题}

10. $A\left( {3,0}\right) ,B\left( {0,4}\right)$ ,动点 $P\left( {x,y}\right)$ 在线段 ${AB}$ 上移动,则 ${xy}$ 的最大值为 ( )

A. $\frac{12}{5}$ . B. $\frac{144}{49}$ . C. 3. D. 4.

11. 已知直线 ${l}_{1} : y = x\sin \alpha \left( {\alpha  \in  \mathbf{R}}\right)$ 和直线 ${l}_{2} : y = {2x} + c$ ,则下述关于直线 ${l}_{1}\text{ 、 }{l}_{2}$ 关系的判断正确的是 ( )

A. 通过平移可以重合. B. 不可能垂直.

C. 可能与 $x$ 轴围成等腰直角三角形. D. 通过绕 ${l}_{1}$ 上某点旋转可以重合.

12. 已知点 $M\left( {a,b}\right)$ 与点 $N\left( {0, - 1}\right)$ 在直线 ${3x} - {4y} + 5 = 0$ 的异侧，给出以下结论: \ding{172}  ${3a} - {4b} + 5 > 0$ ；  \ding{173} 当 $a > 0$ 时， $a + b$ 有最小值，无最大值； \ding{174}  ${a}^{2} + {b}^{2} > 1$ ； \ding{175}  当 $a > 0$ 且 $a \neq  1$ 时， $\frac{b + 1}{a - 1}$ 的取值范围是 $( - \infty$ ， $\left. {-\frac{9}{4}}\right)  \cup  \left( {\frac{3}{4}, + \infty }\right)$ . 其中正确命题的个数是 ( )

A.1. B.2. C. 3. D. 4.

\section*{三、解答题}

13. 已知直线 $l : x + {2y} - 2 = 0$ . 试求:

(1)点 $P\left( {-2, - 1}\right)$ 关于直线 $l$ 的对称点坐标；

(2)直线 $l$ 关于点 $\left( {1,1}\right)$ 对称的直线方程.

14. 直线 ${l}_{1} : y = {mx} + 1,{l}_{2} : x =  - {my} + 1$ 相交于点 $P$ ,其中 $\left| m\right|  \leq  1$ .

(1)求证: ${l}_{1}$ 、 ${l}_{2}$ 分别过定点 $A$ 、 $B$ ，并求点 $A$ 、 $B$ 的坐标；

(2)当 $m$ 为何值时， $\bigtriangleup  {ABP}$ 的面积 $S$ 取得最大值，并求出最大值.

15. 已知点 $A\left( {0,1}\right) \text{ 、 }B\left( {1,1}\right)$ ,设过点 $P\left( {0, - 1}\right)$ 的直线 $l$ 与 $\bigtriangleup {AOB}$ 的边 ${AB}$ 交于点 $M$ (其中点 $M$ 异于 $A\text{ 、 }B$ 两点),与边 ${OB}$ 交于 $N$ (其中点 $N$ 异于 $O\text{ 、 }B$ 两点),若设直线 $l$ 的斜率为 $k\left( {k > 2}\right)$ .

(1)试用 $k$ 来表示点 $M$ 和 $N$ 的坐标；

(2)当 $k$ 为何值时， $\bigtriangleup  {OMN}$ 的面积 $S$ 取得最大值？并求此最大值.

16. 如图，Rt $\bigtriangleup  {OAB}$ 的直角边 ${OA}$ 在 $x$ 轴上，顶点 $B$ 的坐标为 $\left( {6,8}\right)$ ，直线 ${CD}$ 交 ${AB}$ 于点 $D\left( {6,3}\right)$ ，交 $x$ 轴于点 $C\left( {{12},0}\right)$ .

(1)求直线 ${CD}$ 的方程；

(2)动点 $P$ 在 $x$ 轴上从点 $\left( {-{10},0}\right)$ 出发，向 $x$ 轴正方向运动，点 $P$ 在运动过程中，是否存在某个位置， 使得 $\angle {PDA} = \angle B$ ？若存在，请求出点 $P$ 的坐标；若不存在，请说明理由.

\begin{center}
\includegraphics[width=0.2\textwidth]{images/bo_d7tk7oc91nqc73875fu0_199_1180_1160_293_219_0.jpg}
\end{center}
 

(第 16 题图)

\section*{§ 10.2 圆的方程与性质}

\section*{典型例题精讲}

【例 1】已知 $\odot  O$ 的圆心为坐标原点， $\odot  O$ 上的点到直线 $l : x + y - 2\sqrt{2} = 0$ 的距离的最小值为 1 .

(1)求 $\odot  O$ 的方程；

(2)过点 $P\left( {4,2}\right)$ 作 $\odot  O$ 的两条切线，切点分别为 $A$ 、 $B$ ，求直线 ${AB}$ 的方程.

【例 2】已知直线 $l : y = {kx}\left( {k \neq  0}\right)$ 与圆 $C : {x}^{2} + {y}^{2} - {2x} - 3 = 0$ 相交于 $A\text{ 、 }B$ 两点.

(1)若 $\left| {AB}\right|  = \sqrt{13}$ ，求 $k$ ；

(2)在 $x$ 轴上是否存在点 $M$ ，使得当 $k$ 变化时，总有直线 ${MA}$ 、 ${MB}$ 的斜率之和为 0 ？若存在，求出点 $M$ 的坐标; 若不存在, 说明理由.


【例 3】已知矩形 ${ABCD}$ 的对角线交于点 $P\left( {2,0}\right)$ ,边 ${AB}$ 所在直线的方程为 $x - {3y} - 6 = 0$ ,点 $\left( {-1,1}\right)$ 在边 ${AD}$ 所在的直线上,

(1)求矩形 ${ABCD}$ 的外接圆的方程；

(2)已知直线 $l : \left( {1 - {2k}}\right) x + \left( {1 + k}\right) y - 5 + {4k} = 0\left( {k \in  \mathbf{R}}\right)$ ，求证:直线 $l$ 与矩形 ${ABCD}$ 的外接圆恒相交， 并求出相交的弦长最短时的直线 $l$ 的方程.

【例 4】已知过点 $A\left( {-1,0}\right)$ 的直线 $l$ 与圆 $C : {x}^{2} + {\left( y - 3\right) }^{2} = 4$ 相交于 $P\text{ 、 }Q$ 两点, $M$ 是弦 ${PQ}$ 的中点; 且直线 $l$ 与直线 $m : x + {3y} + 6 = 0$ 相交于点 $N$ .

(1)当直线 $l$ 与直线 $m$ 垂直时,求证:直线 $l$ 经过圆心 $C$ ；

(2)设 $t = \overrightarrow{AM} \cdot  \overrightarrow{AN}$ ，试问 $t$ 是否为定值？若为定值，请求出 $t$ 的值；若不为定值，请说明理由.

\begin{center}
\includegraphics[width=0.2\textwidth]{images/bo_d7tk7oc91nqc73875fu0_201_1147_1293_303_268_0.jpg}
\end{center}
 

(例 4 图)

【例 5】已知圆 $C$ 过点 $P\left( {1,1}\right)$ ,且与圆 $M : {\left( x + 2\right) }^{2} + {\left( y + 2\right) }^{2} = {r}^{2}\left( {r > 0}\right)$ 关于直线 $x + y + 2 = 0$ 对称.

(1)求圆 $C$ 的方程；

(2)直线 $l$ 过点 $Q\left( {1,\frac{1}{2}}\right)$ ，截圆 $C$ 所得的弦长为 2，求直线 $l$ 的方程；

(3)过点 $P$ 作两条相异直线分别与圆 $C$ 相交于 $A$ 、 $B$ ，且直线 ${PA}$ 和直线 ${PB}$ 的倾斜角互补， $O$ 为坐标原点,试判断直线 ${OP}$ 和 ${AB}$ 是否平行? 请说明理由.

\section*{四基达标精练}

\section*{一、填空题}

1. 若不过第四象限的直线 $l$ 平分圆 ${x}^{2} + {y}^{2} - {2x} - {4y} = 0$ ，则 $l$ 斜率的取值范围是\_\_\_\_\_.

2.过点 $P\left( {2,3}\right)$ 且与圆 ${\left( x + 1\right) }^{2} + {\left( y + 1\right) }^{2} = 9$ 相切的直线方程为\_\_\_\_\_.

3. 过点 $\left( {1,\sqrt{2}}\right)$ 的直线 $l$ 将圆 ${\left( x - 2\right) }^{2} + {y}^{2} = 4$ 分成两段弧,当劣弧所对的圆心角最小时,直线 $l$ 的斜率 $k =$ \_\_\_\_\_.

4. 若 $y = x + b$ 与曲线 $x = \sqrt{1 - {y}^{2}}$ 恰有一个公共点,则实数 $b$ 的取值范围是\_\_\_\_\_.

5. 若函数 $y = {ax} + {2a} - \sqrt{1 - {x}^{2}}$ 存在零点,则实数 $a$ 的取值范围是\_\_\_\_\_.

6. 若圆 ${x}^{2} + {y}^{2} - {4x} - {4y} - {10} = 0$ 上至少有三个不同的点到直线 $l : {ax} + {by} = 0$ 的距离为 $2\sqrt{2}$ ，则直线 $l$ 倾斜角的取值范围是\_\_\_\_\_.

7. 已知圆 $O$ 的半径为 $1,{PA}\text{ 、 }{PB}$ 为该圆的两条切线， $A$ 、 $B$ 为两切点，那么 $\overrightarrow{PA}$ . $\overrightarrow{PB}$ 的最小值为\_\_\_\_\_.

8. 已知圆 ${x}^{2} + {y}^{2} = 1$ 与圆 ${x}^{2} + {2mx} + {y}^{2} + {ny} = 0\left( {m\text{ 、 }n}\right.$ 是正实数) 相交于 $A\text{ 、 }B$ 两点， $O$ 为坐标原点,当 $\bigtriangleup {AOB}$ 的面积最大时,则 $\frac{\left( {4{m}^{2} + 1}\right) \left( {{n}^{2} + 1}\right) }{mn}$ 的最小值为\_\_\_\_\_.

9. 已知正方形 ${ABCD}$ 边长为 $8,\overrightarrow{BE} = \overrightarrow{EC},\overrightarrow{DF} = 3\overrightarrow{FA}$ ，若在正方形边上恰有 6 个不同的点 $P$ ，使 $\overrightarrow{PE}$ . $\overrightarrow{PF} = \lambda$ ，则实数 $\lambda$ 的取值范围是\_\_\_\_\_.

\section*{二、选择题}

10. 若圆 ${C}_{1} : {x}^{2} + {y}^{2} = 1$ 和圆 ${C}_{2} : {x}^{2} + {y}^{2} - {6x} - {8y} - k = 0$ 没有公共点,则实数 $k$ 的取值范围是( )

A. $\left( {-9,{11}}\right)$ . B. $\left( {-{25}, - 9}\right)$ .

C. $\left( {-\infty , - 9}\right)  \cup  \left( {{11}, + \infty }\right)$ . D. $\left( {-{25}, - 9}\right)  \cup  \left( {{11}, + \infty }\right)$ .

11. 已知 $x\text{ 、 }y$ 满足 ${x}^{2} + {y}^{2} = {4y} - 3$ ,则 $\frac{\sqrt{3}x + y}{\sqrt{{x}^{2} + {y}^{2}}}$ 的最大值为 ( )

A.1. B.2. C. $\sqrt{3}$ . D. $\sqrt{5}$ .

12. 已知圆 $M : {\left( x + \cos \theta \right) }^{2} + {\left( y - \sin \theta \right) }^{2} = 1$ ,直线 $l : y = {kx}$ ,下面四个命题:

 \ding{172} 对任意实数 $k$ 与 $\theta$ ，直线 $l$ 和圆 $M$ 相切；

 \ding{173} 对任意实数 $k$ 与 $\theta$ ，直线 $l$ 和圆 $M$ 有公共点；

 \ding{174} 对任意实数 $\theta$ ，必存在实数 $k$ ，使得直线 $l$ 与圆 $M$ 相切；

 \ding{175} 对任意实数 $k$ ，必存在实数 $\theta$ ，使得直线 $l$ 与圆 $M$ 相切.

其中真命题为 ( )

A.  \ding{172}  \ding{173} . B.  \ding{173}  \ding{174} . C.  \ding{173}  \ding{175} . D.  \ding{172}  \ding{175} .

\section*{三、解答题}

13. 已知 $P\left( {x,y}\right)$ 为圆 $C : {\left( x + 3\right) }^{2} + {\left( y - 4\right) }^{2} = 1$ 上任意一点.

(1)求 $\frac{y - 6}{x}$ 的最值;

(2)求 $x - {2y}$ 的最值.

14. 已知点 $A\left( {1,2}\right)$ ,圆 $C : {x}^{2} + {y}^{2} + {2mx} + {2y} + 2 = 0$ .

(1)若过点 $A$ 可以做两条圆的切线，求 $m$ 的取值范围；

(2)当 $m =  - 2$ 时,过直线 ${2x} - y + 3 = 0$ 上一点 $P$ 作圆的两条切线 ${PM}$ 、 ${PN}$ ，求四边形 ${PMCN}$ 面积的最小值.

15. 已知半径为 $\frac{8}{3}$ 的圆 $C$ 的圆心在 $y$ 轴的正半轴上,且直线 $l : {2x} - {9y} - 1 = 0$ 与圆 $C$ 相切.

(1)求圆 $C$ 的标准方程；

(2)已知 $A\left( {0, - 1}\right)$ ， $P$ 为圆 $C$ 上任意一点，试问在 $y$ 轴上是否存在定点 $B$ (异于点 $A$ ，使得 $\frac{\left| PB\right| }{\left| PA\right| }$ 为定值? 若存在,求点 $B$ 的坐标; 若不存在,请说明理由.

16. 点 $P$ 到 $A\left( {-2,0}\right)$ 的距离是点 $P$ 到 $B\left( {1,0}\right)$ 的距离的 2 倍.

(1)求点 $P$ 的轨迹方程；

(2)点 $P$ 与点 $Q$ 关于点 $\left( {2,1}\right)$ 对称，点 $C\left( {3,0}\right)$ ，求 ${\left| QA\right| }^{2} + {\left| QC\right| }^{2}$ 的最大值和最小值.

\section*{§ 10.3 椭圆及其性质}

\section*{典型例题精讲}

【例 1】设椭圆 $\frac{{x}^{2}}{{a}^{2}} + \frac{{y}^{2}}{{b}^{2}} = 1\left( {a > b > 0}\right)$ 的左焦点为 $F$ ，左顶点为 $A$ ，上顶点为 $B$ . 已知、 $\sqrt{3}\left| {OA}\right|  = 2\left| {OB}\right|$ ( $O$ 为原点).

(1)求椭圆的离心率；

(2)设经过点 $F$ 且斜率为 $\frac{3}{4}$ 的直线 $l$ 与椭圆在 $x$ 轴上方的交点为 $P$ ，圆 $C$ 同时与 $x$ 轴和直线 $l$ 相切，圆心 $C$ 在直线 $x = 4$ 上,且 ${OC}//{AP}$ . 求椭圆的方程.

【例 2】设常数 $m > 0$ 且 $m \neq  1$ ,椭圆 $\Gamma  : \frac{{x}^{2}}{{m}^{2}} + {y}^{2} = 1$ ,点 $P$ 是 $\Gamma$ 上的动点.

(1)若点 $P$ 的坐标为 $\left( {2,0}\right)$ ，求 $\Gamma$ 的焦点坐标；

(2)设 $m = \frac{1}{2}$ ，若 $\Gamma$ 上的另一动点 $Q$ 满足 ${OP} \bot  {OQ}$ ( $O$ 为坐标原点)，求证: $O$ 到直线 ${PQ}$ 的距离是定值.

【例 3】已知曲线 $E : \frac{{x}^{2}}{2} + {y}^{2} = 1\left( {y \neq  0}\right)$ 的左右焦点为 ${F}_{1}\text{ 、 }{F}_{2},P$ 是曲线 $E$ 上一动点.

(1)过 ${F}_{2}$ 的直线与曲线 $E$ 交于 ${AB}$ 两点，且 $\overrightarrow{A{F}_{2}} = 2\overrightarrow{{F}_{2}B}$ ，求直线 ${AB}$ 的斜率；

(2)若存在过点 $H\left( {0,h}\right) \left( {h > 0}\right)$ 的两条直线 ${l}_{1}$ 和 ${l}_{2}$ 与曲线 $E$ 都只有一个公共点,且 ${l}_{1} \bot  {l}_{2}$ ,求 $h$ 的值.

【例 4】如图所示,由半椭圆 ${C}_{1} : \frac{{x}^{2}}{4} + \frac{{y}^{2}}{{b}^{2}} = 1\left( {y < 0}\right)$ 和两个半圆 ${C}_{2} : {\left( x + 1\right) }^{2} + {y}^{2} = 1\left( {y \geq  0}\right) \text{ 、 }{C}_{3}$ : ${\left( x - 1\right) }^{2} + {y}^{2} = 1\left( {y \geq  0}\right)$ 组成曲线 $C : F\left( {x,y}\right)  = 0$ ,其中点 ${A}_{1}\text{ 、 }{A}_{2}$ 依次为 ${C}_{1}$ 的左、右顶点,点 $B$ 为 ${C}_{1}$ 的下顶点,点 ${F}_{1}\text{ 、 }{F}_{2}$ 依次为 ${C}_{1}$ 的左、右焦点. 若点 ${F}_{1}\text{ 、 }{F}_{2}$ 分别为曲线 ${C}_{2}\text{ 、 }{C}_{3}$ 的圆心.

(1)求 ${C}_{1}$ 的方程；

(2)若点 $P$ 、 $Q$ 分别在 ${C}_{2}$ 、 ${C}_{3}$ 上运动，求 $\left| {BP}\right|  + \left| {BQ}\right|$ 的最大值，并求出此时点 $P$ 、 $Q$ 的坐标；

(3)若点 $M$ 在曲线 $C : F\left( {x,y}\right)  = 0$ 上运动，点 $N\left( {0, - 1}\right)$ ，求 $\left| {NM}\right|$ 的取值范围.

\begin{center}
\includegraphics[width=0.2\textwidth]{images/bo_d7tk7oc91nqc73875fu0_206_1176_1279_321_233_0.jpg}
\end{center}
 

(例 4 图)

【例 5】已知 ${F}_{1}\text{ 、 }{F}_{2}$ 分别为椭圆 $C : \frac{{x}^{2}}{{a}^{2}} + \frac{{y}^{2}}{{b}^{2}} = 1\left( {a > b > 0}\right)$ 的左、右焦点， $P\left( {\sqrt{3},\frac{1}{2}}\right)$ 在椭圆 $C$ 上，且 $P{F}_{2} \bot  x$ 轴.

(1)求椭圆 $C$ 的方程；

(2)设 $Q$ 是椭圆 $C$ 上一点， $T\left( {\frac{\sqrt{3}}{2},\frac{\sqrt{3}}{2}}\right)$ ，求 $\left| {QT}\right|  + \left| {Q{F}_{1}}\right|$ 的最大值；

(3)已知点 $A\left( {0,1}\right)$ 及椭圆 $C$ 上的点 $M$ 、 $N$ 满足 ${AM}\bot {AN}$ ,过点 ${F}_{2}$ 作直线 ${MN}$ 的垂线,垂足为 $H$ ,求点 $H$ 的轨迹方程.

\section*{四基达标精练}

\section*{一、填空题}

1. 方程 $5{x}^{2} + k{y}^{2} = {5k}$ 表示焦点在 $y$ 轴上的椭圆，则 $k$ 的取值范围是\_\_\_\_\_.

2. ${F}_{1}$ 是 $\frac{{x}^{2}}{9} + \frac{{y}^{2}}{5} = 1$ 的左焦点， $P$ 是椭圆上的动点， $A\left( {1,1}\right)$ 为定点，则 $\left| {PA}\right|  + \left| {P{F}_{1}}\right|$ 的最小值是\_\_\_\_\_.

3. 椭圆 $\frac{{x}^{2}}{9} + \frac{{y}^{2}}{4} = 1$ 的焦点为 ${F}_{1}\text{ 、 }{F}_{2}$ ,点 $P$ 为其上的动点,当 $\angle {F}_{1}P{F}_{2}$ 为锐角时,点 $P$ 横坐标的取值范围为\_\_\_\_\_.

4. 设 $A$ 是椭圆 $\frac{{x}^{2}}{{a}^{2} + {b}^{2} - {1}^{2}} = 1\left( {a > 0}\right)$ 上的动点,点 $F$ 的坐标为 $\left( {-2,0}\right)$ ,若满足 $\left| {AF}\right|  = {10}$ 的点 $A$ 有且仅有两个，则实数 $a$ 的取值范围为\_\_\_\_\_.

5. 已知 $c$ 是椭圆 $\frac{{x}^{2}}{{a}^{2}} + \frac{{y}^{2}}{{b}^{2}} = 1\left( {a > b > 0}\right)$ 的半焦距,则 $\frac{b + c}{a}$ 的取值范围是\_\_\_\_\_.

6. 已知椭圆 $\frac{{x}^{2}}{4} + \frac{{y}^{2}}{3} = 1$ 上有不同的两点关于直线 $y = {4x} + m$ 对称，则实数 $m$ 的取值范围是\_\_\_\_\_.

7. 已知椭圆 $\frac{{x}^{2}}{16} + \frac{{y}^{2}}{9} = 1$ 的左、右焦点分别为 ${F}_{1}$ 、 ${F}_{2}$ ，过焦点 ${F}_{1}$ 的直线交椭圆于 $A$ 、 $B$ 两点， $\bigtriangleup  {AB}{F}_{2}$ 的内切圆的面积为 $\pi$ ,若 $A\text{ 、 }B$ 两点的坐标分别为 $\left( {{x}_{1},{y}_{1}}\right)$ 和 $\left( {{x}_{2},{y}_{2}}\right)$ ,则 $\left| {{y}_{2} - {y}_{1}}\right|$ 的值为\_\_\_\_\_.

\begin{center}
\includegraphics[width=0.2\textwidth]{images/bo_d7tk7oc91nqc73875fu0_207_1134_1771_311_258_0.jpg}
\end{center}
 

(第 9 题图)

8. 点 ${F}_{1}$ 、 ${F}_{2}$ 分别是椭圆 $C : \frac{{x}^{2}}{2} + {y}^{2} = 1$ 的左、右两焦点,点 $N$ 为椭圆 $C$ 的上顶点, 若动点 $M$ 满足: ${\left| \overrightarrow{MN}\right| }^{2} = 2\overrightarrow{M{F}_{1}} \cdot  \overrightarrow{M{F}_{2}}$ ,则 $\left| {\overrightarrow{M{F}_{1}} + 2\overrightarrow{M{F}_{2}}}\right|$ 的最大值为\_\_\_\_\_.

9. 如图,椭圆 $C : \frac{{x}^{2}}{9} + \frac{{y}^{2}}{4} = 1$ 的上、下顶点分别为 $A\text{ 、 }C$ ,点 $B$ 在椭圆上运动且不与 $A\text{ 、 }C$ 重合,平面四边形 ${ABCD}$ 满足 $\angle {BAD} = \angle {BCD} = {90}^{ \circ  }$ ,则 $\frac{{S}_{\bigtriangleup {ABC}}}{{S}_{\bigtriangleup {ADC}}} =$ \_\_\_\_\_.

\section*{二、选择题}

10. 设 $\frac{{x}^{2}}{\sin \alpha } + \frac{{y}^{2}}{\cos \alpha } = 1\left( {0 < \alpha  < \frac{\pi }{2}}\right)$ 表示焦点在 $x$ 轴上的椭圆,则 $\alpha$ 的取值范围是 ( )

A. $\left( {0,\frac{\pi }{4}}\right\rbrack$ . B. $\left( {\frac{\pi }{4},\frac{\pi }{2}}\right)$ . C. $\left( {0,\frac{\pi }{4}}\right)$ . D. $\left\lbrack  {\frac{\pi }{4},\frac{\pi }{2}}\right)$ .

11. 已知 ${F}_{1}\text{ 、 }{F}_{2}$ 是椭圆的两个焦点,满足 $\overrightarrow{M{F}_{1}} \cdot  \overrightarrow{M{F}_{2}} = 0$ 的点 $M$ 总在椭圆内部,则椭圆离心率的取值范围是 ( )

A. $\left( {0,1}\right)$ . B. $\left( {0,\frac{1}{2}}\right\rbrack$ . C. $\left( {0,\frac{\sqrt{2}}{2}}\right)$ . D. $\left\lbrack  {\frac{\sqrt{2}}{2},1}\right)$ .

12. 设 ${F}_{1}$ 、 ${F}_{2}$ 是椭圆 $\frac{{x}^{2}}{9} + \frac{{y}^{2}}{4} = 1$ 的两焦点， $A$ 与 $B$ 分别是该椭圆的右顶点与上顶点， $P$ 是该椭圆上的一个动点, $O$ 是坐标原点,记 $s = 2\overrightarrow{O{P}^{2}} - \overrightarrow{{F}_{1}P} \cdot  \overrightarrow{{F}_{2}P}$ ,在动点 $P$ 在第一象限内从 $A$ 沿椭圆向左上方运动到 $B$ 的过程中, $s$ 的大小的变化情况为 ( )

A. 逐渐变大. B. 逐渐变小. C. 先变大后变小. D. 先变小后变大.

\section*{三、解答题}

13. 设椭圆 $\frac{{x}^{2}}{{a}^{2}} + \frac{{y}^{2}}{{b}^{2}} = 1\left( {a > b > 0}\right)$ 的左、右顶点分别为 ${A}_{1}$ 、 ${A}_{2}$ ，右焦点为 $F$ ，已知 $\left| {{A}_{1}F}\right|  = 3$ ， $\left| {{A}_{2}F}\right|  = 1$ .

(1)求椭圆方程及其离心率；

(2)已知点 $P$ 是椭圆上一动点 (不与顶点重合),直线 ${A}_{2}P$ 交 $y$ 轴于点 $Q$ ，若 $\bigtriangleup {A}_{1}{PQ}$ 的面积是 $\bigtriangleup {A}_{2}{FP}$ 面积的两倍,求直线 ${A}_{2}P$ 的方程.

\begin{center}
\includegraphics[width=0.2\textwidth]{images/bo_d7tk7oc91nqc73875fu0_208_1155_1026_334_273_0.jpg}
\end{center}
 

(第 13 题图)

14. 已知椭圆 $C : \frac{{x}^{2}}{{a}^{2}} + \frac{{y}^{2}}{{b}^{2}} = 1\left( {a > b > 0}\right)$ 经过点 $A\left( {-2,0}\right)$ 且两个焦点及短轴两顶点围成四边形的面积为 4 .

(1)求椭圆 $C$ 的方程和离心率；

(2)设 $P\text{ 、 }Q$ 为椭圆 $C$ 上不同的两个点，直线 ${AP}$ 与 $y$ 轴交于点 $E$ ，直线 ${AQ}$ 与 $y$ 轴交于点 $F$ ，且 $P\text{ 、 }O\text{ 、 }Q$ 三点共线. 其中 $O$ 为坐标原点. 问: $x$ 轴上是否存在点 $M$ ,使得 $\angle {AME} = \angle {EFM}$ ? 若存在,求点 $M$ 的坐标; 若不存在, 说明理由.

\begin{center}
\includegraphics[width=0.2\textwidth]{images/bo_d7tk7oc91nqc73875fu0_208_1162_1720_335_295_0.jpg}
\end{center}
 

(第 14 题图)

15. 已知椭圆 $E : \frac{{x}^{2}}{{a}^{2}} + \frac{{y}^{2}}{{b}^{2}} = 1\left( {a > b > 0}\right)$ 经过点 $A\left( {0, - 1}\right)$ ，且离心率为 $e = \frac{\sqrt{2}}{2}$ .

(1)求椭圆 $E$ 的方程和焦点 ${F}_{1}$ 、 ${F}_{2}$ 的坐标；

(2)设点 $P$ 为椭圆 $E$ 上的任一点(不在坐标轴上)，直线 $P{F}_{1}$ 与椭圆 $E$ 交于另一点 $M$ ， $P{F}_{2}$ 与 $E$ 交于另一点 $N,O$ 为坐标原点,证明: 直线 ${OP}$ 与 ${MN}$ 的斜率之积为定值.

\begin{center}
\includegraphics[width=0.2\textwidth]{images/bo_d7tk7oc91nqc73875fu0_209_1107_376_336_260_0.jpg}
\end{center}
 

(第 15 题图)

16. 已知椭圆 $C : \frac{{x}^{2}}{{a}^{2}} + \frac{{y}^{2}}{{b}^{2}} = 1\left( {a > b > 0}\right)$ 的一个焦点为 ${F}_{1}\left( {-1,0}\right)$ ，其左顶点为 $A$ ，上顶点为 $B$ ，且 ${F}_{1}$ 到直线 ${AB}$ 的距离为 $\frac{\sqrt{7}}{7}\left| {OB}\right|$ ( $O$ 为坐标原点).

(1)求 $C$ 的方程；

(2)若椭圆 $E : \frac{{x}^{2}}{{a}^{2}} + \frac{{y}^{2}}{{b}^{2}} = \lambda \left( {\lambda  > 0\text{ 且 }\lambda  \neq  1}\right)$ ，则称椭圆 $E$ 为椭圆 $C$ 的 $\lambda$ 倍相似椭圆. 已知椭圆 $E$ 的 3 倍相似椭圆,直线 $l : y = {kx} + m$ 与椭圆 $C$ 、椭圆 $E$ 交于四点 (依次为 $M\text{ 、 }N\text{ 、 }P\text{ 、 }Q$ ,如图),且 $\overrightarrow{MQ} + \overrightarrow{PQ} =  2\overrightarrow{NQ}$ ，证明:点 $T\left( {k,m}\right)$ 在定曲线上.

\begin{center}
\includegraphics[width=0.3\textwidth]{images/bo_d7tk7oc91nqc73875fu0_209_1080_1323_366_271_0.jpg}
\end{center}
 

(第 16 题图)

\section*{§ 10.4 双曲线及其性质}

\section*{典型例题精讲}

【例 1】双曲线 $\Gamma  : \frac{{x}^{2}}{16} - \frac{{y}^{2}}{9} = 1$ 的左、右焦点分别为 ${F}_{1}\text{ 、 }{F}_{2}$ ,直线 $l$ 经过 ${F}_{2}$ 且与 $\Gamma$ 的两条渐近线中的一条平行,与另一条相交且交点在第一象限.

(1)设 $P$ 为 $\Gamma$ 右支上的任意一点，求 $\left| {P{F}_{1}}\right|$ 的最小值；

(2)设 $O$ 为坐标原点，求 $O$ 到 $l$ 的距离，并求 $l$ 与 $\Gamma$ 的交点坐标.

【例 2】已知 ${F}_{1}\text{ 、 }{F}_{2}$ 分别是双曲线 $C : \frac{{x}^{2}}{{a}^{2}} - \frac{{y}^{2}}{{b}^{2}} = 1\left( {a > 0,b > 0}\right)$ 的左、右焦点,其中一条渐近线的倾斜角为 ${150}^{ \circ  }$ ,点 $M$ 在双曲线上,且 $\left| {\overrightarrow{M{F}_{1}} - \overrightarrow{M{F}_{2}}}\right|  = 4$ .

(1)求双曲线 $C$ 的标准方程；

(2)设椭圆 $M$ 以双曲线 $C$ 的顶点为焦点,焦点为顶点,直线 $l : y = {kx} + m\left( {0 < m < 1}\right)$ 交 $M$ 于 $A$ 、 $B$ 两点 (均不在坐标轴上),若 $\bigtriangleup {AOB}$ 的面积为 1,求 $2{k}^{2} - {m}^{2}$ 的值.

\begin{center}
\includegraphics[width=0.3\textwidth]{images/bo_d7tk7oc91nqc73875fu0_210_1082_1481_419_264_0.jpg}
\end{center}
 

(例 2 图)

【例 3】已知双曲线 $C : \frac{{x}^{2}}{{a}^{2}} - \frac{{y}^{2}}{{b}^{2}} = 1\left( {a > 0,b > 0}\right)$ 的离心率为 $\frac{5}{3}$ ，点 $A\left( {a,0}\right)$ 到一条渐近线的距离为 $\frac{12}{5}$ .

(1)求双曲线的方程；

(2)若直线 $l$ 与双曲线 $C$ 交于 $M$ 、 $N$ 两点，且 $\overrightarrow{AM} \cdot  \overrightarrow{AN} = 0$ ，试探究直线 $l$ 是否过定点？若过定点，求出该定点坐标; 若不过定点,说明理由.

【例 4】已知双曲线 $C$ 中心为坐标原点,左焦点为 $\left( {-2\sqrt{5},0}\right)$ ,离心率为 $\sqrt{5}$ .

(1)求 $C$ 的方程；

(2)记 $C$ 的左、右顶点分别为 ${A}_{1}$ 、 ${A}_{2}$ ，过点 $\left( {-4\text{ ， }0}\right)$ 的直线与 $C$ 的左支交于 $M$ 、 $N$ 两点， $M$ 在第二象限， 直线 $M{A}_{1}$ 与 $N{A}_{2}$ 交于 $P$ ,证明 $P$ 在定直线上.

【例 5】数学家将几何学原理运用到折纸中, 并且利用折纸来研究几何学, 很好地把折纸艺术与数学相结合. 将一张纸片折叠一次，纸片上会留下一条折痕，如果在纸片上按照一定的规律折出很多折痕后，纸上能显现出一条漂亮曲线的轮廓.

如图，一张圆形纸片的圆心为点 $D$ ， $A$ 是圆外的一个定点， $P$ 是圆 $D$ 上任意一点，把纸片折叠使得点 $A$ 与 $P$ 重合,然后展平纸片,折痕与直线 ${DP}$ 相交于点 $Q$ ,当点 $P$ 在圆上运动时,得到点 $Q$ 的轨迹.

(1)证明:点 $Q$ 的轨迹是双曲线；

(2)设定点 $A$ 坐标为 $\left( {2,0}\right)$ ，纸片圆的边界方程为 ${\left( x + 2\right) }^{2} + {y}^{2} = {r}^{2}$ . 若点 $M\left( {2,3}\right)$ 位于 (1) 中所描述的双曲线上,过点 $M$ 的直线 $l$ 交该双曲线的渐近线于 $E\text{ 、 }F$ 两点,且点 $E\text{ 、 }F$ 位于 $y$ 轴右侧, $O$ 为坐标原点,求 $\bigtriangleup {EOF}$ 面积的最小值.

\begin{center}
\includegraphics[width=0.2\textwidth]{images/bo_d7tk7oc91nqc73875fu0_212_1180_538_304_264_0.jpg}
\end{center}
 

(例 5 图)

\section*{四基达标精练}

\section*{一、填空题}

1. 双曲线 $\frac{{x}^{2}}{9} - \frac{{y}^{2}}{4} = 1$ 上一点 $P$ 到左焦点的距离为 2，则 $P$ 到右焦点的距离为\_\_\_\_\_.

2. 已知双曲线 $\frac{{x}^{2}}{{a}^{2}} - \frac{{y}^{2}}{{b}^{2}} = 1\left( {a > 0,b > 0}\right)$ 的一条渐近线平行于直线 $l : y = {2x} + {10}$ ，双曲线的一个焦点在直线 $l$ 上，则双曲线的方程为\_\_\_\_\_.

3. 已知 $F$ 是双曲线 $\frac{{x}^{2}}{4} - \frac{{y}^{2}}{12} = 1$ 的左焦点， $A\left( {1,4}\right)$ ，点 $P$ 是双曲线右支上的动点，则 $\left| {PF}\right|  + \left| {PA}\right|$ 的最小值为\_\_\_\_\_.

4. 已知双曲线 ${x}^{2} - m{y}^{2} = 1\left( {m > 0}\right)$ 的右顶点为 $A$ ,而 $B\text{ 、 }C$ 是双曲线右支上两点,若三角形 ${ABC}$ 为等边三角形,则 $m$ 的取值范围是\_\_\_\_\_.

5. 已知双曲线 $\frac{{x}^{2}}{9} - \frac{{y}^{2}}{16} = 1$ ，其右焦点为 $F$ ， $P$ 为双曲线上一点，点 $M$ 满足 $\left| \overrightarrow{MF}\right|  = 1$ ， $\overrightarrow{MF} \cdot  \overrightarrow{MP} = 0$ ，则 $\overrightarrow{MP} \mid$ 的最小值为

\begin{center}
\includegraphics[width=0.2\textwidth]{images/bo_d7tk7oc91nqc73875fu0_213_1170_285_267_256_0.jpg}
\end{center}
 

(第 6 题图)

6. $P$ 是双曲线 $\frac{{x}^{2}}{4} - \frac{{y}^{2}}{5} = 1$ 右支在第一象限内一点， ${F}_{1}$ 、 ${F}_{2}$ 分别为其左、右焦点， $A$ 为右顶点,如图圆 $C$ 是 $\bigtriangleup P{F}_{1}{F}_{2}$ 的内切圆,设圆与 $P{F}_{1}\text{ 、 }P{F}_{2}$ 分别切于点 $D\text{ 、 }E$ ,当圆 $C$ 的面积为 ${4\pi }$ 时,直线 $P{F}_{2}$ 的斜率为\_\_\_\_\_.

7. 已知 ${F}_{1}$ 、 ${F}_{2}$ 为双曲线 $\frac{{x}^{2}}{2} - {y}^{2} = 1$ 的左、右焦点,过 ${F}_{2}$ 且斜率为 $k\left( {k <  - 1}\right)$ 的直线 $l$ 交双曲线右支于 $P\text{ 、 }Q$ 两点,若 $\bigtriangleup {F}_{1}{PQ}$ 的面积为 4,则直线 $l$ 与 $y$ 轴的交点坐标为\_\_\_\_\_.

8. 过双曲线 $C : \frac{{x}^{2}}{{a}^{2}} - \frac{{y}^{2}}{{b}^{2}} = 1\left( {a > 0,b > 0}\right)$ 右焦点 $F$ 作直线 $l$ ,且直线 $l$ 与双曲线 $C$ 的一条渐近线垂直,垂足为 $A$ ,直线 $l$ 与另一条渐近线交于点 $B$ . 已知 $O$ 为坐标原点,若 $\bigtriangleup {OAB}$ 的内切圆的半径为 $\frac{\sqrt{3} - 1}{2}a$ ,则双曲线 $C$ 的离心率为\_\_\_\_\_.

9. 已知 $P$ 为双曲线 $C : \frac{{x}^{2}}{{a}^{2}} - \frac{{y}^{2}}{{b}^{2}} = 1\left( {a > 0,b > 0}\right)$ 右支上的任意一点,经过点 $P$ 的直线与双曲线 $C$ 的两条渐近线分别相交于 $A\text{ 、 }B$ 两点. 若点 $A\text{ 、 }B$ 分别位于第一、第四象限, $O$ 为坐标原点,当 $\overrightarrow{AP} = \frac{1}{2}\overrightarrow{PB}$ 时, $\bigtriangleup {AOB}$ 的面积为 ${2b}$ ，则双曲线 $C$ 的实轴长为\_\_\_\_\_.

\section*{二、选择题}

10. 过双曲线 $\frac{{x}^{2}}{{a}^{2}} - \frac{{y}^{2}}{{b}^{2}} = 1\left( {a > 0,b > 0}\right)$ 的右焦点 $F$ 作倾斜角为 ${60}^{ \circ  }$ 的直线交双曲线右支于 $A$ 、 $B$ 两点,若 $\overrightarrow{AF} = 5\overrightarrow{FB}$ ,则双曲线的离心率为 ( )

A. $\frac{6}{5}$ . B. $\sqrt{3}$ . C. 2.

D. $\frac{4}{3}$ .

11. 已知 $O$ 为坐标原点,双曲线 $C : \frac{{x}^{2}}{{a}^{2}} - \frac{{y}^{2}}{{b}^{2}} = 1\left( {a > 0,b > 0}\right)$ 的左、右焦点分别为 ${F}_{1}\text{ 、 }{F}_{2}$ ,离心率为 $\frac{\sqrt{6}}{2}$ ,点 $P\left( {{x}_{1},{y}_{1}}\right)$ 是 $C$ 的右支上异于顶点的一点,过 ${F}_{2}$ 作 $\angle {F}_{1}P{F}_{2}$ 的平分线的垂线,垂足是 $M,\left| {MO}\right|  = \sqrt{2}$ ,则 $b$ 等于 ( )

A. $\sqrt{2}$ . B. $\sqrt{3}$ . C.1. D.2.

\begin{center}
\includegraphics[width=0.2\textwidth]{images/bo_d7tk7oc91nqc73875fu0_213_1128_1527_318_273_0.jpg}
\end{center}
 

(第 12 题图)

12. 如图， ${F}_{1}$ 、 ${F}_{2}$ 为双曲线的左、右焦点，过 ${F}_{2}$ 的直线交双曲线于 $B$ 、 $D$ 两点， 且 $\overrightarrow{{F}_{2}D} = 3\overrightarrow{{F}_{2}B},E$ 为线段 $D{F}_{1}$ 的中点,若对于线段 $D{F}_{1}$ 上的任意点 $P$ ，都有 $\overrightarrow{P{F}_{1}} \cdot  \overrightarrow{PB} \geq  \overrightarrow{E{F}_{1}} \cdot  \overrightarrow{EB}$ 成立，则双曲线的离心率是 ( )

A. $\sqrt{2}$ . B. $\sqrt{3}$ .

C. 2. D. $\sqrt{5}$

\section*{三、解答题}

13. 已知双曲线 $C$ 的渐近线方程为 $x \pm  \sqrt{3}y = 0$ ,且过点 $\left( {3,\sqrt{2}}\right)$ .

(1)求双曲线 $C$ 的标准方程；

(2)若点 $Q\left( {\frac{3}{2},0}\right)$ ，过右焦点 $F$ 且与坐标轴都不垂直的直线 $l$ 与 $C$ 交于 $A$ 、 $B$ 两点，求证: $\angle {AQF}  = \angle {BQF}.$

14. 已知双曲线 $C : \frac{{x}^{2}}{{a}^{2}} - \frac{{y}^{2}}{{b}^{2}} = 1\left( {a > 0,b > 0}\right)$ 经过点 $P\left( {-2,1}\right)$ 且 $C$ 的右顶点到一条渐近线的距离为 $\frac{\sqrt{6}}{3}$ .

(1)求双曲线 $C$ 的方程；

(2)过点 $P$ 分别作两条直线 ${l}_{1}$ 、 ${l}_{2}$ 与 $C$ 交于 $A$ 、 $B$ 两点( $A$ 、 $B$ 两点均不与点 $P$ 重合)，设直线 ${l}_{1}$ 、 ${l}_{2}$ 的斜率分别为 ${k}_{1}$ 、 ${k}_{2}$ ，若 ${k}_{1} + {k}_{2} = 1$ ，试问直线 ${AB}$ 是否经过定点？若经过定点，求出定点坐标；若不经过定点，请说明理由.

15. 已知双曲线 $C : \frac{{x}^{2}}{{a}^{2}} - \frac{{y}^{2}}{{b}^{2}} = 1\left( {a > 0,b > 0}\right)$ 的焦距为 $2\sqrt{5}$ ,坐标原点 $O$ 到直线 ${BD}$ 的距离是 $\frac{3\sqrt{13}}{13}$ ,其中 $B\text{ 、 }D$ 的坐标分别为 $\left( {0,b}\right) \text{ 、 }\left( {-\frac{3}{2},0}\right)$ .

(1)求双曲线 $C$ 的方程；

(2)是否存在过点 $D$ 的直线 $l$ 与双曲线 $C$ 交于 $M\text{ 、 }N$ 两点，使得 $\bigtriangleup  {BMN}$ 构成以 $B$ 为顶点的等腰三角形? 若存在,求出所有直线 $l$ 的方程; 若不存在,请说明理由.

16. 已知双曲线 ${C}_{1} : \frac{{x}^{2}}{4} - \frac{{y}^{2}}{{b}^{2}} = 1\left( {b > 0}\right)$ 与圆 ${C}_{2} : {x}^{2} + {y}^{2} = 4 + {b}^{2}\left( {b > 0}\right)$ 在第一象限的交点为 $A\left( {{x}_{A},{y}_{A}}\right)$ , 圆 ${C}_{2}$ 与 $x$ 轴交于 ${F}_{1}\text{ 、 }{F}_{2}$ ,曲线 $\Gamma$ 满足 $\left\{  \begin{array}{l} \frac{{x}^{2}}{4} - \frac{{y}^{2}}{{b}^{2}} = 1,\left| x\right|  > {x}_{A}, \\  {x}^{2} + {y}^{2} = 4 + {b}^{2},\left| x\right|  \leq  {x}_{A}. \end{array}\right.$

(1)当 $b = \sqrt{5}$ ，第一象限一点 $P$ 在曲线 ${C}_{1}$ 上，且 $\left| {P{F}_{1}}\right|  = 8$ ，求 $\angle {F}_{1}P{F}_{2}$ ；

(2)过点 $D\left( {0,2 + \frac{{b}^{2}}{2}}\right)$ 且斜率为 $- \frac{b}{2}$ 的直线 $l$ 与曲线 $\Gamma$ 有且只有两个交点 $M$ 、 $N$ ，用含 $b$ 的代数式表示 $\overrightarrow{OM} \cdot  \overrightarrow{ON}$ ，并求 $\overrightarrow{OM} \cdot  \overrightarrow{ON}$ 的取值范围.

\section*{§ 10.5 抛物线及其性质}

\section*{典型例题精讲}

【例 1】设抛物线 $C : {y}^{2} = {2px}\left( {p > 0}\right)$ 的焦点为 $F$ ，点 $P\left( {4,m}\right)$ 是抛物线 $C$ 上一点，且 $\left| {PF}\right|  = 5$ .

(1)求抛物线 $C$ 的方程；

(2)设直线 $l : y = {kx} + m$ 与抛物线 $C$ 交于 $A$ 、 $B$ 两点,若 $\left| {AF}\right|  + \left| {BF}\right|  = 6$ ,求证: 线段 ${AB}$ 的垂直平分线过定点.

\begin{center}
\includegraphics[width=0.2\textwidth]{images/bo_d7tk7oc91nqc73875fu0_216_1213_592_275_261_0.jpg}
\end{center}
 

(例 1 图)

【例 2】已知抛物线 $C : {y}^{2} = {2px}\left( {0 < p < 4}\right)$ 的焦点为 $F$ ，直线 $l$ 与抛物线 $C$ 交于不同的两点 $A$ 、 $B$ . 当 $\left| {AF}\right|  = 5$ 时,以线段 ${AF}$ 为直径的圆过点 $\left( {0,2}\right)$ .

(1)求抛物线 $C$ 的方程；

(2)若线段 ${AB}$ 的中点在曲线 ${y}^{2} = {2x} - 4$ 上运动，求 $\bigtriangleup  {AOB}$ (其中 $O$ 为平面直角坐标系的原点)的面积的最小值.


【例 3】已知动点 $M$ 到直线 $x + 2 = 0$ 的距离比到点 $F\left( {1,0}\right)$ 的距离大 1 .

(1)求动点 $M$ 所在的曲线 $C$ 的方程；

(2)已知点 $P\left( {1,2}\right)$ ， $A$ 、 $B$ 是曲线 $C$ 上的两个动点，如果直线 ${PA}$ 的斜率与直线 ${PB}$ 的斜率互为相反数， 证明直线 ${AB}$ 的斜率为定值，并求出这个定值；

(3)已知点 $P\left( {1,2}\right)$ ， $A$ 、 $B$ 是曲线 $C$ 上的两个动点，如果直线 ${PA}$ 的斜率与直线 ${PB}$ 的斜率之和为 2，证明: 直线 ${AB}$ 过定点.

【例 4】已知动圆 $M$ 过点 $F\left( {1,0}\right)$ ,且与直线 $x =  - 1$ 相切,圆心 $M$ 的轨迹是曲线 $C$ .

(1)求曲线 $C$ 的方程；

(2)过 $F$ 作弦 ${PQ}$ 、 ${RS}$ ，设 ${PQ}$ 、 ${RS}$ 的中点分别为 $A$ 、 $B$ ，若 $\overrightarrow{PQ} \cdot  \overrightarrow{RS} = 0$ ，求 $\left| \overrightarrow{AB}\right|$ 最小时，弦 ${PQ}$ 、 ${RS}$ 所在直线的方程;

(3)在(2)条件下，是否存在一定点 $T$ ，使得 $\overrightarrow{AF} = \lambda \overrightarrow{TB} - \overrightarrow{FT}$ ？若存在，求出 $T$ 的坐标；若不存在，试说明理由.

【例 5】已知抛物线 $C : {x}^{2} = {2py}\left( {p > 0}\right)$ 的焦点为 $F$ ， $M$ 是抛物线 $C$ 上位于第一象限内的任意一点， $O$ 为坐标原点,记经过 $M\text{ 、 }F\text{ 、 }O$ 三点的圆的圆心为 $Q$ ,且点 $Q$ 到抛物线 $C$ 的准线的距离为 $\frac{3}{2}$ .

(1)求点 $Q$ 的纵坐标 (用 $p$ 表示)；

(2)求抛物线 $C$ 的方程；

(3)设直线 $l : y = {kx} + \frac{1}{2}$ 与抛物线 $C$ 有两个不同的交点 $A$ 、 $B$ . 若点 $M$ 的横坐标为 2，且 $\bigtriangleup  {QAB}$ 的面积为 $2\sqrt{5}$ ，求直线 $l$ 的方程.

\section*{四基达标精练}

\section*{一、填空题}

1. 经过点 $\left( {2,4}\right)$ 的抛物线 $y = a{x}^{2}$ 焦点坐标是\_\_\_\_\_.

2. 抛物线 ${y}^{2} = {2x}$ 上的动点 $P$ 到直线 $y = {2x} + 1$ 的距离最小值是\_\_\_\_\_.

3. 过抛物线 ${y}^{2} = {8x}$ 的焦点 $F$ ，斜率为 $- \sqrt{3}$ 的直线 $l$ 与抛物线相交于点 $A$ 、 $B$ ，则弦 ${AB}$ 的长度为\_\_\_\_\_.

4. 已知点 $P$ 是抛物线 ${y}^{2} = {4x}$ 上的动点,点 $P$ 在 $y$ 轴上的射影是 $M$ ,点 $A$ 的坐标是 $\left( {4,a}\right)$ ,则当 $\left| a\right|  > 4$ 时, $\left| {PA}\right|  + \left| {PM}\right|$ 的最小值是\_\_\_\_\_.

5. 设抛物线 $y = {x}^{2} + m$ 过椭圆 $\frac{{x}^{2}}{{a}^{2}} + \frac{{y}^{2}}{{b}^{2}} = 1\left( {a > b > 0}\right)$ 的两个焦点,且和椭圆有三个交点,以这三个交点为顶点的三角形面积为 1,则 $m =$ \_\_\_\_\_.

6. 已知 $F$ 是抛物线 ${y}^{2} = {4x}$ 的焦点,其准线与 $x$ 轴交于 $P$ 点,过 $P$ 的直线 $l$ 与抛物线交于 $A\text{ 、 }B$ 两点,设 ${FA}\text{ 、 }{FB}$ 的斜率分别为 $m\text{ 、 }n$ ,则 $\frac{m}{n} =$ \_\_\_\_\_.

7. 已知不过原点的直线 $l$ 与抛物线 $C : {y}^{2} = {2px}\left( {p > 0}\right)$ 交于 $A\text{ 、 }B$ 两点,若 $\left| {AF}\right|  = 2\left| {BF}\right|$ 且 $\angle {AFB} =  {90}^{ \circ  }$ ,则直线 $l$ 的斜率为\_\_\_\_\_.

8. 已知 $F$ 是抛物线 ${y}^{2} = {4x}$ 的焦点,其准线与 $x$ 轴交于 $P$ 点,过 $P$ 的直线 $l$ 与抛物线交于 $A\text{ 、 }B$ 两点,若线段 ${AB}$ 上有一点 $M\left( {x,y}\right)$ ，满足 $\left| \overrightarrow{AM}\right|$ ， $\left| \overrightarrow{PB}\right|  = \left| \overrightarrow{BM}\right|$ ， $\left| \overrightarrow{PA}\right|$ ，则 $M$ 点的轨迹方程是\_\_\_\_\_.

9. 已知斜率为 $k$ 的直线 $l$ 过抛物线 $E : {y}^{2} = 4\sqrt{3}x$ 的焦点 $F$ ，分别与抛物线 $E$ 和 $\odot  {O}_{1} : {\left( x - \sqrt{3}\right) }^{2} + {y}^{2} = 3$ 从上到下依次交于为 $A\text{ 、 }B\text{ 、 }C\text{ 、 }D$ 四点,则 $\left| {AB}\right|  \cdot  \left| {CD}\right|  =$ \_\_\_\_\_.

\section*{二、选择题}

10. 设 $O$ 为坐标原点, $F$ 为抛物线 $C : {x}^{2} = {2py}\left( {p > 0}\right)$ 的焦点,直线 $y = 1$ 与抛物线 $C$ 交于 $A\text{ 、 }B$ 两点,若 $\angle {AFB} = {120}^{ \circ  }$ ,则抛物线 $C$ 的准线方程为 ( )

A. $y =  - \frac{2}{3}$ . B. $y =  - 3$ .

C. $y =  - \frac{1}{3}$ 或 $y =  - 3$ . D. $y =  - \frac{2}{3}$ 或 $y =  - 6$ .

11. 设抛物线焦点为 $F$ ,准线与对称轴交于点 $E$ ,过 $F$ 的直线交抛物线于 $A\text{ 、 }B$ 两点,对称轴上一点 $C$ 满足 $\overrightarrow{CA} = 3\overrightarrow{BE}$ ,若 $\bigtriangleup {ACF}$ 的面积为 $\frac{9\sqrt{3}}{4}$ ,则 $F$ 到抛物线准线的距离为 ( )

A. $\sqrt{6}$ .

B. $\frac{\sqrt{6}}{2}$ . C. $\frac{2\sqrt{3}}{3}$ . D. $\frac{\sqrt{3}}{3}$ .

12. 设 $A\text{ 、 }B$ 是抛物线 $C : {y}^{2} = {4x}$ 上两个不同的点, $O$ 为坐标原点,若直线 ${OA}$ 与 ${OB}$ 的斜率之积为 -4, 有下列结论:  \ding{172}  $\left| {AB}\right|  \geq  4$ ;  \ding{173}  $\left| {OA}\right|  + \left| {OB}\right|  > 8$ ;  \ding{174}  直线 ${AB}$ 过抛物线 $C$ 的焦点;  \ding{175}  $\bigtriangleup {OAB}$ 面积的最小值是 2.其中正确的有 ( )

A.  \ding{172}  \ding{174}  \ding{175} . B.  \ding{172}  \ding{173}  \ding{175} . C.  \ding{173}  \ding{174}  \ding{175} . D.  \ding{172}  \ding{173}  \ding{174}  \ding{175} .

\section*{三、解答题}

13. 设抛物线 $\Gamma  : {y}^{2} = {4x}$ 的焦点为 $F$ ，直线 $l : x - {my} - n = 0$ 经过 $F$ 且与 $\Gamma$ 交于 $A$ 、 $B$ 两点.

(1)若 $\left| {AB}\right|  = 8$ ，求 $m$ 的值；

(2)设 $O$ 为坐标原点，直线 ${AO}$ 与 $\Gamma$ 的准线交于点 $C$ ,求证: 直线 ${BC}$ 平行于 $x$ 轴.

14. 已知抛物线 ${y}^{2} = {4x}$ 的焦点为 $F$ ,直线 $l$ 交抛物线于不同的 $A$ 、 $B$ 两点.

(1)若直线 $l$ 经过点 $P\left( {-1,0}\right)$ ，点 $A$ 关于 $x$ 轴的对称点为 ${A}^{\prime }$ ，求证: ${A}^{\prime }$ 、 $F$ 、 $B$ 三点共线；

(2)若直线 $l$ 经过点 $M\left( {8, - 4}\right)$ ，抛物线上是否存在定点 $N$ ，使得以线段 ${AB}$ 为直径的圆恒过点 $N$ ？若存在,求出点 $N$ 的坐标; 若不存在,请说明理由.

15. 已知过点 $P\left( {1,0}\right)$ 的直线 $l$ 与抛物线 $C : {x}^{2} = {2py}\left( {p > 0}\right)$ 相交于 $A$ 、 $B$ 两点,当直线 $l$ 过抛物线 $C$ 的焦点时, $\left| {AB}\right|  = 8$ .

(1)求抛物线 $C$ 的方程；

(2)若点 $Q\left( {0, - 2}\right)$ ，连接 ${QA}\text{ 、 }{QB}$ 分别交抛物线 $C$ 于点 $E\text{ 、 }F$ ，且 $\bigtriangleup  {QAB}$ 与 $\bigtriangleup  {QEF}$ 的面积之比为 1:2， 求直线 ${AB}$ 的方程.

16. 设抛物线 $C : {y}^{2} = {2px}\left( {p > 0}\right)$ ，直线 $x - {2y} + 1 = 0$ 与 $C$ 交于 $A$ 、 $B$ 两点，且 $\left| {AB}\right|  = 4\sqrt{15}$ .

(1)求 $p$ 的值；

(2) $F$ 为 ${y}^{2} = {2px}$ 的焦点， $M$ 、 $N$ 为抛物线上的两点，且 $\overrightarrow{MF} \cdot  \overrightarrow{NF} = 0$ ，求 $\bigtriangleup  {MNF}$ 面积的最小值.

\section*{§ 10.6 圆锥曲线综合一}

\section*{典型例题精讲}

【例 1】已知: 以点 $C\left( {t,\frac{2}{t}}\right) \left( {t \in  \mathbf{R},t \neq  0}\right)$ 为圆心的圆与 $x$ 轴交于点 $O\text{ 、 }A$ ,与 $y$ 轴交于点 $O\text{ 、 }B$ ,其中 $O$ 为坐标原点.

(1)求证: $\bigtriangleup  {OAB}$ 的面积为定值；

(2)设直线 $y =  - {2x} + 4$ 与圆 $C$ 交于点 $M$ 、 $N$ ，若 $\left| {OM}\right|  = \left| {ON}\right|$ ，求圆 $C$ 的方程.

【例 2】在平面直角坐标系 ${xOy}$ 中,抛物线 $C$ 的顶点在原点,经过点 $A\left( {2,2}\right)$ ,其焦点 $F$ 在 $x$ 轴上.

(1)求抛物线 $C$ 的标准方程；

(2)求过点 $F$ ，且与直线 ${OA}$ 垂直的直线方程；

(3)设过点 $M\left( {m,0}\right)$ 的直线交抛物线 $C$ 于 $D$ 、 $E$ 两点， $\left| {ME}\right|  = 2\left| {DM}\right|$ ，记 $D$ 和 $E$ 两点间的距离为 $f\left( m\right)$ ,求 $f\left( m\right)$ 关于 $m$ 的表达式.

\begin{center}
\includegraphics[width=0.2\textwidth]{images/bo_d7tk7oc91nqc73875fu0_221_1228_1440_211_239_0.jpg}
\end{center}
 

(例 2 图)

【例 3】已知两定点 ${F}_{1}\left( {-\sqrt{2},0}\right) ,{F}_{2}\left( {\sqrt{2},0}\right)$ ,满足 $\left| {\overrightarrow{PF}}_{2}\right|  - \left| {\overrightarrow{PF}}_{1}\right|  = 2$ 的点 $P$ 的轨迹是曲线 $E$ ,直线 $y = {kx} - 1$ 与曲线 $E$ 交于 $A\text{ 、 }B$ 两点.

(1)求 $k$ 的取值范围；

(2)如果 $\left| \overrightarrow{AB}\right|  = 6\sqrt{3}$ ，且曲线 $E$ 上存在点 $C$ ，使 $\overrightarrow{OA} + \overrightarrow{OB} = m\overrightarrow{OC}$ ，求 $m$ 的值和 $\bigtriangleup  {ABC}$ 的面积 $S$ .

【例 4】已知抛物线 $C : {y}^{2} = x$ 上有两点 $A\text{ 、 }B$ ，且 $A$ 在 $x$ 轴上方， $O$ 为坐标原点， $F$ 为抛物线焦点，若 $\angle {AOB} = {90}^{ \circ  }$ ,求 ${S}_{\bigtriangleup {AOB}} + {S}_{\bigtriangleup {AOF}}$ 的最小值.

【例 5】若椭圆 ${E}_{1} : \frac{{x}^{2}}{{a}_{1}^{2}} + \frac{{y}^{2}}{{b}_{1}^{2}} = 1$ 和椭圆 ${E}_{2} : \frac{{x}^{2}}{{a}_{2}^{2}} + \frac{{y}^{2}}{{b}_{2}^{2}} = 1$ 满足 $\frac{{a}_{2}}{{a}_{1}} = \frac{{b}_{2}}{{b}_{1}} = m\left( {m > 0}\right)$ ,则称这两个椭圆相似, $m$ 称为其相似比.

(1)求经过点 $\left( {2,\sqrt{6}}\right)$ ，且与椭圆 $\frac{{x}^{2}}{4} + \frac{{y}^{2}}{2} = 1$ 相似的椭圆方程；

(2)设过原点的一条射线 $l$ 分别与(1)中的两个椭圆交于 $A$ 、 $B$ 两点(其中点 $A$ 在线段 ${OB}$ 上)，求 $\left| {OA}\right|  \cdot  \left| {OB}\right|$ 的最大值和最小值;

(3)对于真命题“过原点的一条射线分别与相似比为 2 的两个椭圆 ${C}_{1} : \frac{{x}^{2}}{{2}^{2}} + \frac{{y}^{2}}{{\left( \sqrt{2}\right) }^{2}} = 1$ 和 ${C}_{2} : \frac{{x}^{2}}{{4}^{2}} +  \frac{{y}^{2}}{{\left( 2\sqrt{2}\right) }^{2}} = 1$ 交于 $A\text{ 、 }B$ 两点, $P$ 为线段 ${AB}$ 上的一点,若 $\left| {OA}\right| \text{ 、 }\left| {OP}\right| \text{ 、 }\left| {OB}\right|$ 成等比数列,则点 $P$ 的轨迹方程为 $\frac{{x}^{2}}{{\left( 2\sqrt{2}\right) }^{2}} + \frac{{y}^{2}}{{2}^{2}} = 1$ ”. 请用推广或类比的方法提出类似的一个真命题，并给予证明.

\section*{四基达标精练}

\section*{一、填空题}

1. 已知椭圆中心在原点，一个焦点为 $F\left( {\sqrt{3},0}\right)$ ，且长轴长是短轴长的 2 倍，则该椭圆的标准方程是\_\_\_\_\_.

2. 已知椭圆 $\frac{{x}^{2}}{16} + \frac{{y}^{2}}{7} = 1$ 的右焦点为 $F$ ，过原点 $O$ 作直线交椭圆于 $A$ 、 $B$ 两点，点 $A$ 在 $x$ 轴的上方，若三角形 ${ABF}$ 的面积为 2，则点 $A\left( {p,q}\right)$ 的纵坐标 $q =$ \_\_\_\_\_.

3. 若实数 $x$ 、 $y$ 满足 $y = 2\sqrt{1 - {x}^{2}}$ ，则 $y - {2x}$ 的取值范围为\_\_\_\_\_.

4. 已知双曲线 $\Gamma  : \frac{{x}^{2}}{3} - \frac{{y}^{2}}{2} = 1,{F}_{1}\text{ 、 }{F}_{2}$ 是 $\Gamma$ 的左、右焦点，点 $P$ 是 $\Gamma$ 上的点,若 $\left| {P{F}_{1}}\right|  \cdot  \left| {P{F}_{2}}\right|  = 4$ ,则 $\overrightarrow{P{F}_{1}} \cdot  \overrightarrow{P{F}_{2}}$ 的值为\_\_\_\_\_.

5. 设直线 $l$ 与抛物线 ${y}^{2} = {4x}$ 相交于 $A$ 、 $B$ 两点,与圆 ${\left( x - 5\right) }^{2} + {y}^{2} = {r}^{2}\left( {r > 0}\right)$ 相切于点 $M$ ，且 $M$ 为线段 ${AB}$ 的中点，若这样的直线 $l$ 恰有 4 条，则 $r$ 的取值范围是\_\_\_\_\_.

6. 设 $P$ 是双曲线 ${x}^{2} - \frac{{y}^{2}}{15} = 1$ 的右支上一点，过点 $P$ 分别作圆 ${\left( x + 4\right) }^{2} + {y}^{2} = 4$ 和 ${\left( x - 4\right) }^{2} + {y}^{2} = 1$ 的切线，切点分别为 $M$ 、 $N$ ，则 ${\left| PM\right| }^{2} - {\left| PN\right| }^{2}$ 的最小值为\_\_\_\_\_.

7. 已知椭圆 $\frac{{x}^{2}}{{a}^{2}} + \frac{{y}^{2}}{{b}^{2}} = 1\left( {a > b > 0}\right)$ 的左、右焦点分别为 ${F}_{1}\left( {-c,0}\right)$ 、 ${F}_{2}\left( {c,0}\right)$ ，若椭圆上存在一点 $P$ 使 $\frac{a}{\sin \angle P{F}_{1}{F}_{2}} = \frac{c}{\sin \angle P{F}_{2}{F}_{1}}$ ，则椭圆离心率 $e$ 的取值范围是\_\_\_\_\_.

8. 已知 ${F}_{1}$ 、 ${F}_{2}$ 分别为双曲线 ${x}^{2} - \frac{{y}^{2}}{3} = 1$ 的左、右焦点,过 ${F}_{2}$ 的直线与双曲线的右支交于 $A$ 、 $B$ 两点,记 $\bigtriangleup A{F}_{1}{F}_{2}$ 的内切圆 ${O}_{1}$ 的半径为 ${r}_{1},\bigtriangleup B{F}_{1}{F}_{2}$ 的内切圆 ${O}_{2}$ 的半径为 ${r}_{2}$ ,圆 ${O}_{1}\text{ 、 }{O}_{2}$ 的面积分别为 ${S}_{1}\text{ 、 }{S}_{2}$ , 则 ${S}_{1} + {S}_{2}$ 的取值范围是\_\_\_\_\_.

9. 以下四个关于圆锥曲线的命题中: \ding{172}  当 $1 < k < 4$ 时，曲线 $C : \frac{{x}^{2}}{4 - k} + \frac{{y}^{2}}{k - 1} = 1$ 表示椭圆； \ding{173}  双曲线 $\frac{{x}^{2}}{25} -  \frac{{y}^{2}}{9} = 1$ 和椭圆 $\frac{{x}^{2}}{35} + {y}^{2} = 1$ 有相同的焦点;  \ding{174} 设 $A\text{ 、 }B$ 为两个定点, $k$ 为非零常数, $\left| \overrightarrow{PA}\right|  - \left| \overrightarrow{PB}\right|  = k$ ,则动点 $P$ 的轨迹为双曲线;  \ding{175} 过定圆 $C$ 上一定点 $A$ 和动点 $B$ 作圆的弦 ${AB},O$ 为坐标原点,若 $\overrightarrow{OA} + \overrightarrow{OB} = 2\overrightarrow{OP}$ 则动点 $P$ 的轨迹为圆; 其中真命题的序号为\_\_\_\_\_(写出所有真命题的序号).

\section*{二、选择题}

10. 已知 $A\text{ 、 }B$ 是双曲线 $\Gamma  : \frac{{x}^{2}}{4} - \frac{{y}^{2}}{3} = 1$ 的左、右顶点,动点 $P$ 在 $\Gamma$ 上且 $P$ 在第一象限,若 ${PA}\text{ 、 }{PB}$ 的斜率分别为 ${k}_{1}\text{ 、 }{k}_{2}$ ,则以下总为定值的是 ( )

A. ${k}_{1} + {k}_{2}$ . B. $\left| {{k}_{1} - {k}_{2}}\right|$ . C. ${k}_{1} \cdot  {k}_{2}$ . D. ${k}_{1}^{2} + {k}_{2}^{2}$ .

11. 已知椭圆 ${C}_{1} : \frac{{x}^{2}}{{a}^{2}} + \frac{{y}^{2}}{{b}^{2}} = 1\left( {a > b > 0}\right)$ 与圆 ${C}_{2} : {x}^{2} + {y}^{2} = \frac{4{b}^{2}}{5}$ ,若在椭圆 ${C}_{1}$ 上不存在点 $P$ ,使得由点 $P$ 所作的圆 ${C}_{2}$ 的两条切线互相垂直,则椭圆 ${C}_{1}$ 的离心率的取值范围是 ( )

A. $\left( {0,\frac{\sqrt{3}}{3}}\right)$ . B. $\left( {0,\frac{\sqrt{6}}{4}}\right)$ . C. $\left\lbrack  {\frac{\sqrt{3}}{3},1}\right)$ . D. $\left\lbrack  {\frac{\sqrt{6}}{4},1}\right)$ .

12. 在圆锥 ${PO}$ 中,已知高 ${PO} = 2$ ,底面圆的半径为 $4,M$ 为母线 ${PB}$ 的中点; 根据圆锥曲线的定义,下列四个图中的截面边界曲线分别为圆、椭圆、双曲线及抛物线,下面四个命题:  \ding{172}  圆的面积为 ${4\pi }$ ;  \ding{173}  椭圆的长轴为 $\sqrt{37}$ ;  \ding{174}  双曲线两渐近线的夹角正切值为 $- \frac{3}{4}$ ;  \ding{175} 抛物线中焦点到准线的距离为 $\frac{4\sqrt{5}}{5}$ . 其中正确的个数为 ( )

\begin{center}
\includegraphics[width=0.2\textwidth]{images/bo_d7tk7oc91nqc73875fu0_224_1073_1626_312_180_0.jpg}
\end{center}
 

\begin{center}
\includegraphics[width=0.2\textwidth]{images/bo_d7tk7oc91nqc73875fu0_224_248_1636_288_174_0.jpg}
\end{center}
 

$P$

$M  M$

O B $O \cdot$

(第 12 题图)

A. 1 个. B. 2 个. C. 3 个. D. 4 个.

\section*{三、解答题}

13. 设椭圆 $\frac{{x}^{2}}{{a}^{2}} + \frac{{y}^{2}}{{b}^{2}} = 1\left( {a > b > 0}\right)$ 的左、右顶点分别为是 $A\text{ 、 }B$ ,点 $P$ 在椭圆上且异于 $A\text{ 、 }B$ 两点, $O$ 为坐标原点.

(1)若直线 ${AP}$ 与 ${BP}$ 的斜率之积为 $- \frac{1}{2}$ ，求椭圆的离心率；

(2)若 $\left| {AP}\right|  = \left| {OA}\right|$ ，证明直线 ${OP}$ 的斜率 $k$ 满足 $\left| k\right|  > \sqrt{3}$ .

14. 已知抛物线 $\Gamma  : {y}^{2} = {2px}$ 的焦点为 $F\left( {2,0}\right)$ ,点 $P$ 在抛物线 $\Gamma$ 上.

(1)求抛物线 $\Gamma$ 的方程；

(2)过点 $T\left( {t,0}\right) \left( {t > 0}\right)$ 作两条互相垂直的直线分别交抛物线 $\Gamma$ 于 $A$ 、 $B$ 、 $C$ 、 $D$ 四点，且点 $M$ 、 $N$ 分别为线段 ${AB}$ 、 ${CD}$ 的中点,求 $\bigtriangleup {TMN}$ 的面积的最小值.

\begin{center}
\includegraphics[width=0.2\textwidth]{images/bo_d7tk7oc91nqc73875fu0_225_1172_1266_273_309_0.jpg}
\end{center}
 

(第 14 题图)

15. 已知双曲线 ${x}^{2} - {y}^{2} = 2$ 的左、右焦点分别为 ${F}_{1}$ 、 ${F}_{2}$ ，过点 ${F}_{2}$ 的动直线与双曲线相交于 $A$ 、 $B$ 两点.

(1)若动点 $M$ 满足 $\overrightarrow{{F}_{1}M} = \overrightarrow{{F}_{1}A} + \overrightarrow{{F}_{1}B} + \overrightarrow{{F}_{1}O}$ (其中 $O$ 为坐标原点)，求点 $M$ 的轨迹方程；

(2)在 $x$ 轴上是否存在定点 $C$ ，使 $\overrightarrow{CA}$ - $\overrightarrow{CB}$ 为常数？若存在，求出点 $C$ 的坐标；若不存在，请说明理由.

16. 已知椭圆 $C : \frac{{x}^{2}}{{a}^{2}} + \frac{{y}^{2}}{{b}^{2}} = 1\left( {a > b > 0}\right)$ 的右焦点为 $F\left( {1,0}\right)$ ，且点 $P\left( {1,\frac{3}{2}}\right)$ 在椭圆 $C$ 上.

(1)求椭圆 $C$ 的标准方程；

(2)过椭圆 ${C}_{1} : \frac{{x}^{2}}{{a}^{2}} + \frac{{y}^{2}}{{b}^{2} - \frac{5}{3}} = 1$ 上异于其顶点的任意一点 $Q$ 作圆 $O : {x}^{2} + {y}^{2} = \frac{4}{3}$ 的两条切线，切点分别为 $M\text{ 、 }N\left( {M\text{ 、 }N}\right.$ 不在坐标轴上 $)$ ,若直线 ${MN}$ 在 $x$ 轴、 $y$ 轴上的截距分别为 $m\text{ 、 }n$ ,证明: $\frac{1}{3{m}^{2}} + \frac{1}{{n}^{2}}$ 为定值.

\section*{§ 10.7 圆锥曲线综合二}

\section*{典型例题精讲}

【例 1】已知椭圆 $C : \frac{{x}^{2}}{{a}^{2}} + \frac{{y}^{2}}{{b}^{2}} = 1\left( {a > b > 0}\right)$ 的离心率为 $\frac{1}{2}$ ，左、右焦点分别为 ${F}_{1}$ 、 ${F}_{2}$ ， $O$ 为坐标原点，点 $P$ 在椭圆 $C$ 上，且满足 $\left| {P{F}_{1}}\right|  + \left| {P{F}_{2}}\right|  = 8$ .

(1)求椭圆 $C$ 的方程；

(2)已知过点 $\left( {2,0}\right)$ 且不与 $x$ 轴重合的直线 $l$ 与椭圆 $C$ 交于 $M$ ， $N$ 两点，在 $x$ 轴上是否存在定点 $Q$ ，使得 $\angle {MQO} = \angle {NQO}$ ? 若存在,求出点 $Q$ 的坐标; 若不存在,说明理由.

【例 2】已知双曲线 $C : \frac{{x}^{2}}{{a}^{2}} - \frac{{y}^{2}}{{b}^{2}} = 1\left( {a > 0,b > 0}\right)$ 的一个焦点为 $F\left( {\sqrt{2},0}\right)$ ,且经过点 $T\left( {\frac{\sqrt{5}}{2},\frac{1}{2}}\right)$ .

(1)求双曲线 $C$ 的标准力程；

(2)已知点 $A$ 是 $C$ 上一定点，过点 $B\left( {0,1}\right)$ 的动直线与双曲线 $C$ 交于 $P$ 、 $Q$ 两点，若 ${k}_{AP} + {k}_{AQ}$ 为定值 $\lambda$ ， 求点 $A$ 的坐标及实数 $\lambda$ 的值.

【例 3】已知斜率为 $k$ 的直线 $l$ 经过抛物线 $C : {y}^{2} = {4x}$ 的焦点 $F$ ，且与抛物线 $C$ 交于不同的两点 $A\left( {{x}_{1},{y}_{1}}\right)$ 、 $B\left( {{x}_{2},{y}_{2}}\right)$ .

(1)若点 $A$ 和 $B$ 到抛物线准线的距离分别为 $\frac{3}{2}$ 和3，求 $\left| {AB}\right|$ ；

(2)若 $\left| {AF}\right|  + \left| {AB}\right|  = 2\left| {BF}\right|$ ，求 $k$ 的值.

【例 4】在平面直角坐标系 ${xOy}$ 中,已知椭圆 $\Gamma  : \frac{{x}^{2}}{2} + {y}^{2} = 1$ ,过右焦点 $F$ 做两条互相垂直的弦 ${AB}$ 、 ${CD}$ ,设 ${AB}\text{ 、 }{CD}$ 的中点分别为 $M\text{ 、 }N$ .

(1)写出椭圆右焦点 $F$ 的坐标及椭圆的离心率；

(2)证明:直线 ${MN}$ 必过定点,并求出该定点坐标.

【例 5】在 ${xOy}$ 坐标平面内,已知椭圆 $\Gamma  : \frac{{x}^{2}}{9} + \frac{{y}^{2}}{5} = 1$ 的左、右焦点分别为 ${F}_{1}\text{ 、 }{F}_{2}$ ,直线 $y = {k}_{1}x\left( {{k}_{1} \neq  0}\right)$ 与 $\Gamma$ 相交于 $A\text{ 、 }B$ 两点.

(1)记 $d$ 为 $A$ 到直线 ${2x} + 9 = 0$ 的距离，当 ${k}_{1}$ 变化时，求证: $\frac{\left| A{F}_{1}\right| }{d}$ 为定值；

(2)当 $\angle {A{F}_{2}B} = {120}^{ \circ  }$ 时，求 $\left| {A{F}_{2}}\right|  \cdot  \left| {B{F}_{2}}\right|$ 的值；

(3)过 $B$ 作 ${BM} \bot  x$ 轴,垂足为 $M,{OM}$ 的中点为 $N$ ,延长 ${AN}$ 交 $\Gamma$ 于另一点 $P$ ,记直线 ${PB}$ 的斜率为 ${k}_{2}$ ,当 ${k}_{1}$ 取何值时, $\left| {{k}_{1} - {k}_{2}}\right|$ 有最小值? 并求出此最小值.

\begin{center}
\includegraphics[width=0.2\textwidth]{images/bo_d7tk7oc91nqc73875fu0_229_1131_504_306_258_0.jpg}
\end{center}
 

(例 5 图)

\section*{四基达标精练}

\section*{一、填空题}

1. 抛物线 ${y}^{2} = {4x}$ 的焦点为 $F$ ，过 $F$ 的直线交该抛物线于 $A$ 、 $B$ 两点，则 $\left| {AF}\right|  + 4\left| {BF}\right|$ 的最小值为\_\_\_\_\_.

2. 已知 $M\text{ 、 }N$ 为双曲线 $\frac{{x}^{2}}{4} - {y}^{2} = 1$ 上关于坐标原点 $O$ 对称的两点, $P$ 为双曲线上异于 $M\text{ 、 }N$ 的点,若直线 ${PM}$ 的斜率的取值范围是 $\left\lbrack  {\frac{1}{2},2}\right\rbrack$ ，则直线 ${PN}$ 的斜率的取值范围是\_\_\_\_\_.

3. 从一块短轴长为 ${2b}$ 的椭圆形玻璃镜中划出一块面积最大的矩形,其面积的取值范围是 $\left\lbrack  {3{b}^{2},4{b}^{2}}\right\rbrack$ ,则这个椭圆的离心率的取值范围是\_\_\_\_\_.

4. 设 $O$ 是坐标原点, $F$ 是抛物线 ${y}^{2} = {2px}\left( {p > 0}\right)$ 的焦点, $A$ 是抛物线上的一点, $\overrightarrow{FA}$ 与 $x$ 轴正向的夹角为 ${60}^{ \circ  }$ ,则 $\left| \overrightarrow{OA}\right|  =$ \_\_\_\_\_.

5. 如图, $A\text{ 、 }B$ 是直线 $l$ 上的两点,且 ${AB} = 2$ . 两个半径相等的动圆分别与 $l$ 相切于 $A\text{ 、 }B$ 点, $C$ 是这两个圆的公共点,则圆弧 $\overset{\text{ ⏜ }}{AC}\text{ 、 }\overset{\text{ ⏜ }}{CB}$ 与线段 ${AB}$ 围成图形面积 $S$ 的取值范围是\_\_\_\_\_.

\begin{center}
\includegraphics[width=0.2\textwidth]{images/bo_d7tk7oc91nqc73875fu0_229_1113_1491_326_146_0.jpg}
\end{center}
 

(第 5 题图)

6. 设 $F$ 是双曲线 $\frac{{x}^{2}}{4} - \frac{{y}^{2}}{12} = 1$ 的左焦点, $A\left( {1,4}\right) ,P$ 是双曲线右支上的动点, 则 $\left| {PA}\right|  + \left| {PF}\right|$ 的最小值为\_\_\_\_\_.

7. 过抛物线 ${y}^{2} = {4x}$ 的焦点 $F$ 的直线交抛物线于 $A$ 、 $B$ 两点，则 $\frac{1}{\left| AF\right| } + \frac{1}{\left| BF\right| } =$ \_\_\_\_\_.

8. 已知 $m\text{ 、 }n$ 是正常数， $s\text{ 、 }t \in  {\mathbf{R}}^{ + },m + n = 2$ ， $\frac{m}{s} + \frac{n}{t} = 9$ ，其中 $m\text{ 、 }n$ 是常数，且 $s + t$ 的最小值是 $\frac{4}{9}$ ，满足条件的点 $\left( {m,n}\right)$ 是圆 ${\left( x - 2\right) }^{2} + {\left( y - 2\right) }^{2} = 4$ 中一条弦的中点,则此弦所在的直线方程为\_\_\_\_\_.

9. 在直线 $l : y =  - 2$ 上取一点 $D$ 作抛物线 $C : {x}^{2} = {4y}$ 的切线，切点分别为 $A$ 、 $B$ ，直线 ${AB}$ 与圆 $E$ : ${x}^{2} +  {y}^{2} - {2x} - {2020} = 0$ 交于 $M$ 、 $N$ 两点，当 $\left| {MN}\right|$ 最小时，点 $D$ 的横坐标是\_\_\_\_\_.

\section*{二、选择题}

\begin{center}
\includegraphics[width=0.2\textwidth]{images/bo_d7tk7oc91nqc73875fu0_230_1245_213_249_274_0.jpg}
\end{center}
 

(第 10 题图)

( )

10. 从双曲线 $\frac{{x}^{2}}{{a}^{2}} - \frac{{y}^{2}}{{b}^{2}} = 1\left( {a > 0,b > 0}\right)$ 的左焦点 $F$ 引圆 ${x}^{2} + {y}^{2} = {a}^{2}$ 的切线，切点为 $T$ ,延长 ${FT}$ 交双曲线右支于 $P$ 点,若 $M$ 为线段 ${FP}$ 的中点, $O$ 为坐标原点,则 $\left| {MO}\right|  -  \left| {MT}\right|$ 与 $b - a$ 的大小关系为 ( )

A. $\left| {MO}\right|  - \left| {MT}\right|  > b - a$ . B. $\left| {MO}\right|  - \left| {MT}\right|  = b - a$ .

C. $\left| {MO}\right|  - \left| {MT}\right|  < b - a$ . D. 不确定.

11. 点 $P$ 在直线 $l : y = x - 3$ 上,若存在过 $P$ 的直线交抛物线 ${x}^{2} = {4y}$ 于 $A\text{ 、 }B$ 两点, 且 $\left| {PA}\right|  = \left| {AB}\right|$ ,则称点 $P$ 是“倍起点”,那么下列结论中正确的是

A. 直线 $l$ 上仅有有限个点是 “倍起点”. B. 直线 $l$ 上不存在 “倍起点”.

C. 直线 $l$ 上有无穷多个点 (不是所有点) 是 “倍起点”. D. 直线 $l$ 上每个点都是 “倍起点”.

12. 设双曲线 $\frac{{x}^{2}}{{a}^{2}} - \frac{{y}^{2}}{{b}^{2}} = 1\left( {a > 0,b > 0}\right)$ 的左、右焦点为 ${F}_{1}\text{ 、 }{F}_{2}$ ,若双曲线右支上存在点 $P$ ,使得 $\left| {P{F}_{2}}\right|$ , $\left| {P{F}_{1}}\right| ,\left| {{F}_{1}{F}_{2}}\right|$ 依次成等差数列,则该双曲线的离心率的取值范围为 ( )

A. $\lbrack 3, + \infty )$ . B. $\left( {1,3}\right)$ . C. $\left( {3, + \infty }\right)$ . D. $\left( {1,3}\right)$ .

\section*{三、解答题}

13. 已知双曲线 $\Gamma  : \frac{{x}^{2}}{2} - \frac{{y}^{2}}{4} = 1$ 的右顶点为 $A$ ,点 $B$ 的坐标为 $\left( {1,\sqrt{2}}\right)$ .

(1)设点 $D$ 是双曲线 $\Gamma$ 上的动点，若点 $N$ 满足 $\overrightarrow{BN} = \overrightarrow{ND}$ ，求点 $N$ 的轨迹方程；

(2)过点 $B$ 的动直线 $l$ 交双曲线 $\Gamma$ 于 $P$ 、 $Q$ 两个不同的点， $M$ 为线段 ${PQ}$ 的中点，求直线 ${AM}$ 斜率的取值范围.

14. 已知抛物线 $\Gamma  : {y}^{2} = {4x}$ .

(1)求抛物线 $\Gamma$ 的焦点 $F$ 的坐标和准线 $l$ 的方程；

(2)过焦点 $F$ 且斜率为 $\frac{1}{2}$ 的直线与抛物线 $\Gamma$ 交于不同两点 $A$ 、 $B$ ，求线段 ${AB}$ 的长；

(3)已知点 $P\left( {1,2}\right)$ ，是否存在定点 $Q$ ，使得过点 $Q$ 的直线与抛物线 $\Gamma$ 交于两个不同的点 $M$ 、 $N$ (均不与点 $P$ 重合),且以线段 ${MN}$ 为直径的圆恒过点 $P$ ? 若存在,求出点 $Q$ 的坐标; 若不存在,请说明理由.


15. 已知 $a > b > 0$ ,如图,曲线 $\Gamma$ 由曲线 ${C}_{1} : \frac{{x}^{2}}{{a}^{2}} + \frac{{y}^{2}}{{b}^{2}} = 1\left( {y \leq  0}\right)$ 和曲线 ${C}_{2} : \frac{{x}^{2}}{{a}^{2}} - \frac{{y}^{2}}{{b}^{2}} = 1\left( {y > 0}\right)$ 组成,其中点 ${F}_{1}\text{ 、 }{F}_{2}$ 为曲线 ${C}_{1}$ 所在圆锥曲线的焦点,点 ${F}_{3}\text{ 、 }{F}_{4}$ 为曲线 ${C}_{2}$ 所在圆锥曲线的焦点.

(1)若 ${F}_{2}\left( {2,0}\right)$ ， ${F}_{3}\left( {-6,0}\right)$ ，求曲线 $\Gamma$ 的方程；

(2)对于(1)中的曲线 $\Gamma$ ，若直线 ${l}_{1}$ 过点 ${F}_{4}$ 交曲线 ${C}_{1}$ 于点 $C$ 、 $D$ ，求 $\bigtriangleup  {CD}{F}_{1}$ 面积的最大值.

\begin{center}
\includegraphics[width=0.3\textwidth]{images/bo_d7tk7oc91nqc73875fu0_231_1072_363_374_274_0.jpg}
\end{center}
 

(第 15 题图)

16. 已知双曲线 $C$ 的中心在坐标原点,左焦点 ${F}_{1}$ 与右焦点 ${F}_{2}$ 都在 $x$ 轴上,离心率为 3,过点 ${F}_{2}$ 的动直线 $l$ 与双曲线 $C$ 交于点 $A\text{ 、 }B$ ,设 $\frac{\left| {A{F}_{2}}\right|  \cdot  \left| {B{F}_{2}}\right| }{{\left| AB\right| }^{2}} = \lambda$ .

(1)求双曲线 $C$ 的渐近线方程；

(2)若点 $A$ 、 $B$ 都在双曲线 $C$ 的右支上，求 $\lambda$ 的最大值以及 $\lambda$ 取最大值时 $\angle {A{F}_{1}}B$ 的正切值；(关于求 $\lambda$ 的最值,某学习小组提出了如下的思路可供参考: \ding{172}  利用基本不等式求最值;  \ding{173}  设 $\frac{\left| A{F}_{2}\right| }{\left| AB\right| }$ 为 $\mu$ ，建立相应数量关系并利用它求最值;  \ding{174} 设直线 $l$ 的斜率为 $k$ ，建立相应数量关系并利用它求最值).

\begin{center}
\includegraphics[width=0.2\textwidth]{images/bo_d7tk7oc91nqc73875fu0_231_1186_1361_260_320_0.jpg}
\end{center}
 

(第 16 题图)

\section*{第 11 章 排列组合、二项式定理(选择性必修 2)}

\section*{知识梳理*}

\section*{一、排列、组合}

1. 加法原理(分类计数原理):做一件事，完成它有 $n$ 类不同的办法. 第一类办法中有 ${a}_{1}$ 种方法，第二类办法中有 ${a}_{2}$ 种方法 $\cdots \cdots$ ,第 $n$ 类办法中有 ${a}_{n}$ 种方法,则完成这件事共有: $N = {a}_{1} + {a}_{2} + \cdots  + {a}_{n}$ 种不同的方法.

2. 乘法原理(分步计数原理):做一件事，需要依次完成 $n$ 个步骤. 其中完成第一步有 ${a}_{1}$ 种方法，完成第二步有 ${a}_{2}$ 种方法, $\cdots \cdots$ ,完成第 $n$ 步有 ${a}_{n}$ 种方法,则完成这件事共有: $N = {a}_{1}{a}_{2}\cdots {a}_{n}$ 种方法.

注意:处理实际问题时，要善于区分是用分类计数原理还是分步计数原理，这两个原理的标志是“分类” 还是“分步骤”.

3. 加法原理和乘法原理的共同点在于都是研究完成一件事情, 共有多少种不同的方法, 其差别在于完成一件事情的方式不同，加法原理“分类完成”，即用不同类型的办法去完成一件事情，每类办法中的任何一类方法，都可以独立完成这件事；而乘法原理是“分步完成”，即完成一件事情而分成连续和完整的不同步骤， 而且只有每一步都完成了, 才能完成这件事.

4. 排列:从 $n$ 个不同的元素中取 $m\left( {m \leq  n}\right)$ 个元素，按照一定的顺序排成一排，叫做从 $n$ 个不同的元素中取 $m$ 个元素的排列;

其个数记作: ${{P}}_{n}^{m}$ ,公式: ${{P}}_{n}^{m} = n\left( {n - 1}\right) \left( {n - 2}\right) \cdots \left( {n - m + 1}\right)  = \frac{n!}{\left( {n - m}\right) !}$ (规定 $0! = 1$ ).

5. 组合: 从 $n$ 个不同的元素中,任取 $m\left( {m \leq  n}\right)$ 个元素并成一组,叫做从 $n$ 个不同的元素中取 $m$ 个元素的组合;

其个数记作: ${{C}}_{n}^{m}$ ,公式: ${{C}}_{n}^{m} = \frac{{{P}}_{n}^{m}}{{{P}}_{m}^{m}} = \frac{n\left( {n - 1}\right) \left( {n - 2}\right) \cdots \left( {n - m + 1}\right) }{m\left( {m - 1}\right) \left( {m - 2}\right) \cdots 3 \times  2 \times  1} = \frac{n!}{m!\left( {n - m}\right) !}$ .

6. 组合数的两个性质: ${{C}}_{n}^{m} = {{C}}_{n}^{n - m};{{C}}_{n}^{m} = {{C}}_{n - 1}^{m} + {{C}}_{n - 1}^{m - 1}$ .

7. 解排列组合问题的依据是:分类相加，分步相乘，有序排列，无序组合.

8. 枚举法是解排列组合和概率问题的基本方法，千万别忘了.

9. 解排列组合问题的规律是:相邻问题捆绑法；不邻问题插空法；多排问题单排法；定位问题优先法；定序问题倍缩法；多元问题分类法；有序分配问题法；选取问题先选后排法；至多至少问题间接法.

10. 求解排列组合应用问题的思考程序:

(1)分析:找出条件和结论；明确元素和位置；

(2)分辨:弄清是否有序(排列)或无序(组合)；明确元素与位置有无控制；

(3)分类:元素分类法；位置分类法. 用加法原理关联类与类；

(4)分步:问题分解；厘清步骤. 用乘法原理关联步与步；

(5) 择法:直接法(元素为主法；位置为主法)；间接法(用 $n\left( A\right)  = n\left( I\right)  - n\left( \overline{A}\right)$ ).

\section*{二、二项式定理}

1. 二项式定理: ${\left( a + b\right) }^{n} = {{C}}_{n}^{0}{a}^{n} + {{C}}_{n}^{1}{a}^{n - 1}b + {{C}}_{n}^{2}{a}^{n - 2}{b}^{2} + \cdots  + {{C}}_{n}^{k}{a}^{n - k}{b}^{k} + \cdots  + {{C}}_{n}^{n}{b}^{n}$ .

\HRule

* 注:本章节所属知识模块是“数据与统计”.

\HRule

2. 二项展开式的通项公式: ${T}_{k + 1} = {{C}}_{n}^{k}{a}^{n - k}{b}^{k}\left( {k = 0,1,2,\cdots ,n}\right)$ ,其中 ${{C}}_{n}^{k}$ 叫二项式系数.

3. 杨辉三角: 第 $n$ 行系数为 ${\left( a + b\right) }^{n}$ 的展开式的二项式系数 ${{C}}_{n}^{0},{{C}}_{n}^{1},{{C}}_{n}^{2},\cdots ,{{C}}_{n}^{n}$ .

\begin{center}
\includegraphics[width=0.8\textwidth]{images/bo_d7tk7oc91nqc73875fu0_233_242_243_1011_362_0.jpg}
\end{center}
 

4.二项式系数的性质:

(1)二项式系数具有对称性:与首末两端等距离的两项的二项式系数相等; 当 $n$ 为奇数时,中间两项 ${T}_{\frac{n - 1}{2} + 1}$ 和 ${T}_{\frac{n + 1}{2} + 1}$ 的二项式系数最大,为 ${{C}}_{n}^{\frac{n - 1}{2}}$ 和 ${{C}}_{n}^{\frac{n + 1}{2}}$ ; 当 $n$ 为偶数时,中间一项 ${T}_{\frac{n}{2} + 1}$ 的二项式系数最大,为 ${{C}}_{n}^{\frac{n}{2}}$ ;

(2) ${{C}}_{n}^{m} = \frac{n}{m}{{C}}_{n - 1}^{m - 1} = \frac{m + 1}{n + 1}{{C}}_{n + 1}^{m + 1}$ ；

(3) ${{C}}_{n}^{0} + {{C}}_{n}^{1} + {{C}}_{n}^{2} + \cdots  + {{C}}_{n}^{n} = {2}^{n}$ ； ${{C}}_{n}^{0} + {{C}}_{n}^{2} + {{C}}_{n}^{4} + \cdots  = {{C}}_{n}^{1} + {{C}}_{n}^{3} + {{C}}_{n}^{5} + \cdots  = {2}^{n - 1}$ .

\section*{$\text{ § }$ 11.1 排列与组合}

\section*{典型例题精讲}

【例 1】有甲、乙、丙、丁、戊 5 名同学, 求:

(1)5 名同学站成一排，甲、戊不在两端有多少种不同的排法?

(2)5 名同学站成一排，要求甲乙必须相邻，丙丁不能相邻，有多少种不同的排法?

【例 2】用 0、1、2、3、4 五个数字组成没有重复数字的五位数，并把它们从小到大排列.

(1)问 23140 是第几个数?

(2)问第 70 个数是什么?

【例 3】平面上有相异的 11 个点, 每两点连成一条直线, 共得 43 条不同的直线.

(1)这 11 个点中有无三点或三个以上的点共线? 若有共线, 情形怎样?

(2)这 11 个点构成多少个三角形？

【例 4】设集合 $P = \{ m \mid  m = {2k},k \in  \mathbf{N},1 \leq  k \leq  6\} ,Q = \{ n \mid  n = {6k},k \in  \mathbf{N},1 \leq  k \leq  6\}$ .

(1)若集合 $A$ 含有三个元素，且 $A \subseteq  P$ ，这样的集合 $A$ 有多少个？所有集合 $A$ 中各元素之和是多少?

(2)若集合 $A$ 、 $B$ 各含有三个元素，且 $A \subseteq  P$ ， $B \subseteq  Q$ ， $A \cap  B = \varnothing$ ，这样的集合 $A$ 、 $B$ 有多少种配对方式？

【例 5】设 $n$ 是正整数,对于 $1,2,\cdots ,n$ 的一个排列 ${i}_{1}{i}_{2}\cdots {i}_{n}$ ,如果当 $s < t$ 时,有 ${i}_{s} > {i}_{t}$ ,则称 $\left( {{i}_{s},{i}_{t}}\right)$ 是排列 ${i}_{1}{i}_{2}\cdots {i}_{n}$ 的一个逆序,排列 ${i}_{1}{i}_{2}\cdots {i}_{n}$ 的所有逆序的总个数称为其逆序数. 例如: 对1,2,3的一个排列 231,只有两个逆序 $\left( {2,1}\right) \text{ 、 }\left( {3,1}\right)$ ,则排列 231 的逆序数为 2 .记 ${f}_{n}\left( k\right)$ 为 $1,2,\cdots ,n$ 的所有排列中逆序数为 $k$ 的全部排列的个数.

(1)求 ${f}_{3}$ (2) 和 ${f}_{4}$ (2)的值；

(2)求 ${f}_{n}\left( 2\right) \left( {n \geq  5}\right)$ 的表达式(用 $n$ 表示).

\section*{四基达标精练}

\section*{一、填空题}

1. 正整数 $n$ 满足 ${{C}}_{25}^{2n} = {{C}}_{25}^{{18} - n}$ ，则 $n =$ \_\_\_\_\_.

2. 从甲、乙等 5 名同学中随机选 3 名参加社区服务工作，则甲、乙都入选的概率为\_\_\_\_\_.

3. 从 6 个人挑选 4 个人去值班, 每人值班一天, 第一天安排 1 个人, 第二天安排 1 个人, 第三天安排 2 个人，则共有\_\_\_\_\_种安排方案.

4.参观团由 13 人组成，要选出 2 男 1 女分别担任团长、政委、联络员，所选女性不为团长，有 2 名男性不在候选人之列，已知选出且确定职务的方法有 300 种，则此团中有\_\_\_\_\_，名男团员.

5. 已知花博会有四个不同的场馆 $A$ 、 $B$ 、 $C$ 、 $D$ ，甲、乙两人每人选 2 个去参观，则他们的选择中，恰有一个馆相同的概率为\_\_\_\_\_.

6. 有 7 个高矮不等的人排成一行，其中最高的恰好在中间，而从中间向两边看正好由高到矮，这样的排法共有\_\_\_\_\_种.

7. 有 5 男 3 女共 8 个人列成一队，若女的顺序一定，则共有\_\_\_\_\_种不同的排法.

8. 正方体的八个顶点所确定的直线中有\_\_\_\_\_对异面直线。

9. 设集合 $A = \{ 1,2,3,\cdots ,{2023}\}$ ，选择 $A$ 的两个非空子集 $B$ 和 $C$ ，要使 $C$ 中最小的数大于 $B$ 中的最大数，则不同的选择方法有\_\_\_\_\_.

\section*{二、选择题}

10. 从 2 至 8 的 7 个整数中随机取 2 个不同的数, 则这 2 个数互质的概率为 ( )

A. $\frac{1}{6}$ . B. $\frac{1}{3}$ . C. $\frac{1}{2}$ . D. $\frac{2}{3}$ .

11. 将 4 个 1 和 2 个 0 随机排成一行，则 2 个 0 不相邻的概率为 ( )

A. $\frac{1}{3}$ . B. $\frac{2}{5}$ . C. $\frac{2}{3}$ . D. $\frac{4}{5}$ .

12. 方程 ${ay} = {b}^{2}{x}^{2} + c$ 中的 $a$ 、 $b$ 、 $c \in  \{  - 2,0,1,2,3\}$ ，且 $a$ 、 $b$ 、 $c$ 互不相同，在所有这些方程所表示的曲线中, 不同的抛物线共有 ( )

A. 28 条. B. 32 条. C. 36 条. D. 48 条.

\section*{三、解答题}

13. 已知1、2、3、4、7、9 六个数.

(1)可以组成多少没有重复数字的五位数；

(2)其中有多少个是偶数；

(3)其中有多少个是 3 的倍数.

14. 袋中有 $m$ 个白球和 $n$ 个黑球, $m \geq  n \geq  1$ .

(1)若 $m = 6,n = 5$ ，一次随机抽取两个球，求两个球颜色相同的概率；

(2)有放回地抽取两次，每次随机抽取一个球，若两次取出的球的颜色相同的概率为 $\frac{5}{8}$ ，求 $m : n$ .

15. 从一批苹果中, 随机抽取 50 个, 其重量 (单位: 克) 的频数分布表如下:

\begin{center}
\begin{tabular}{|c|c|c|c|c|}
\hline
分组(重量) & [80,85) & [85, 90) & [90,95) & [95,100) \\
\cline{1-5}
频数(个) & 5 & 10 & 20 & 15 \\
\cline{1-5}
\hline
\end{tabular}
\end{center}

(1)根据频数分布表计算苹果的重量在 $\lbrack {90},{95})$ 的频率;

(2)用分层抽样的方法从重量在 $\lbrack {80},{85})$ 和 $\lbrack {95},{100})$ 的苹果中共抽取 4 个，任取 2 个，求重量在 $\lbrack {80},{85})$ 和 $\lbrack {95},{100})$ 中各有 1 个的概率.

16. 已知 $A = \left\{  {x\left| {1 < {\log }_{2}x < 3}\right. }\right. ,x$ 是正整数 $\} ,B = \{ x\left| \right| x - 6 \mid   < 3,x$ 是正整数 $\}$ .

(1)从集合 $A$ 、 $B$ 中各取一个元素构成直角坐标系中点的坐标，共可以得到多少个不同的点；

(2)从 $A \cup  B$ 中取出 3 个不同元素组成三位数，从左到右的数字要逐渐增大，这样的三位数有多少个;

(3)从集合 $A$ 中取一个元素，从集合 $B$ 中取 3 个元素，可以组成多少个无重复数字且比 4000 大的自然数.

\section*{$\text{ § }$ 11.2 二项式定理}

\section*{典型例题精讲}

【例 1】已知 ${\left( 1 + x\right) }^{n}$ 的展开式中前三项的系数之和等于 79,求 ${\left( \frac{1}{2} + 2x\right) }^{n}$ 的展开式中系数最大的项.

【例 2】设 ${\left( 2 - \sqrt{3}x\right) }^{100} = {a}_{0} + {a}_{1}x + {a}_{2}{x}^{2} + \cdots  + {a}_{100}{x}^{100}$ ,

(1) 求 ${a}_{0}$ ；

(2) $A = {\left( {a}_{0} + {a}_{2} + {a}_{4} + \cdots  + {a}_{100}\right) }^{2},B = {\left( {a}_{1} + {a}_{3} + {a}_{5} + \cdots  + {a}_{99}\right) }^{2}$ ,求 $A - B$ .

【例 3】 ${\left( x + 3y - 4z\right) }^{12}$ 的展开式中,

(1)求 ${x}^{2}{y}^{7}{z}^{3}$ 的系数；

(2)求各项的系数和.

\section*{高分钥匙高考数学复习一轮通}

【例 4】设 ${\left( 1 + x\right) }^{n} = {a}_{0} + {a}_{1}x + {a}_{2}{x}^{2} + \cdots  + {a}_{n}{x}^{n},n \geq  4,n \in  \mathbf{N}$ . 已知 ${a}_{3}^{2} = 2{a}_{2}{a}_{4}$ .

(1)求 $n$ 的值；

(2)设 ${\left( 1 + \sqrt{3}\right) }^{n} = a + b\sqrt{3}$ ，其中 $a$ 、 $b \in  \mathbf{N}$ ，求 ${a}^{2} - 3{b}^{2}$ 的值.

【例 5】( 1 )求 $7{{C}}_{6}^{3} - 4{{C}}_{7}^{4}$ 的值；

(2)设 $m$ 、 $n$ 是正整数， $n \geq  m$ ，求证: $\left( {m + 1}\right) {{C}}_{m}^{m} + \left( {m + 2}\right) {{C}}_{m + 1}^{m} + \left( {m + 3}\right) {{C}}_{m + 2}^{m} + \cdots  + n{{C}}_{n - 1}^{m} + \left( {n + 1}\right) {{C}}_{n}^{m}  = \left( {m + 1}\right) {{C}}_{n + 2}^{m + 2}$ .

\section*{四基达标精练}

\section*{一、填空题}

1. 已知二项式 ${\left( x + a\right) }^{5}$ 展开式中， ${x}^{2}$ 的系数为 80 ，则 $a =$ \_\_\_\_\_.

2. 在 ${\left( {x}^{3} + \frac{1}{x}\right) }^{12}$ 的展开式中，则含 $\frac{1}{{x}^{4}}$ 项的系数为\_\_\_\_\_.

3. ${\left( \sqrt{x} + \frac{3}{{x}^{2}}\right) }^{5}$ 的展开式中的常数项为\_\_\_\_\_.

4. 二项式 ${\left( 3 + x\right) }^{n}$ 的展开式中， ${x}^{2}$ 项的系数是常数项的 5 倍，则 $n =$ \_\_\_\_\_.

5. $\left( {1 - \frac{y}{x}}\right) {\left( x + y\right) }^{8}$ 的展开式中， ${x}^{2}{y}^{6}$ 的系数为\_\_\_\_\_.

6．已知 $\left( {x + 2}\right) {\left( x - 1\right) }^{4} = {a}_{0} + {a}_{1}x + {a}_{2}{x}^{2} + {a}_{3}{x}^{3} + {a}_{4}{x}^{4} + {a}_{5}{x}^{5}$ ，则 ${a}_{1} + {a}_{2} + {a}_{3} + {a}_{4} + {a}_{5} =$ \_\_\_\_\_.

${7.2}^{{6n} - 3} + {3}^{{2n} - 1}$ ( $n$ 是正整数) 除以 11 所得的余数是\_\_\_\_\_.

8. ${\left( \left| x\right|  + \frac{1}{\left| x\right| } - 2\right) }^{3}$ 的展开式中的常数项为\_\_\_\_\_.

9. 对给定函数 $f\left( x\right)$ ，定义一个多项式 ${f}_{n}\left( x\right)$ 如下: ${f}_{n}\left( x\right)  = \mathop{\sum }\limits_{{k = 0}}^{n}{{C}}_{n}^{k}f\left( \frac{k}{n}\right) {x}^{k}{\left( 1 - x\right) }^{n - k}$ ，其中 $n = 1,2$ ， $3,\cdots$ ,若 $f\left( x\right)  = 1\left( {0 \leq  x \leq  1}\right)$ ,则 ${f}_{n}\left( x\right)  =$ \_\_\_\_\_.

\section*{二、选择题}

10. 在二项式 ${\left( \sqrt{2} + x\right) }^{10}$ 展开式中,系数为有理数的项数是 ( )

A.4. B.5. C. 6. D. 7.

11. 已知 ${\left( 1 + x\right) }^{n}$ 的展开式中,有且仅有 ${x}^{3}$ 的系数最大,则 ${\left( 1 + x\right) }^{n}$ 的系数和为 ( )

A. 16. B.32. C. 64. D.128.

12. 已知 ${x}^{2023} = {a}_{0} + {a}_{1}\left( {x + 1}\right)  + {a}_{2}{\left( x + 1\right) }^{2} + \cdots  + {a}_{2023}{\left( x + 1\right) }^{2023}$ ,则 ${a}_{2022} =$ ( )

A.2022. B. -2022. C. 2023. D. -2023.

\section*{三、解答题}

13. 已知二项式 ${\left( \sqrt[3]{x} + \frac{1}{x}\right) }^{n}$ 的展开式中各项系数的和为 256 .

(1)求 $n$ ；

(2)求展开式中的常数项.

14. 已知 ${\left( 2x + 1\right) }^{4} = {a}_{0} + {a}_{1}x + {a}_{2}{x}^{2} + {a}_{3}{x}^{3} + {a}_{4}{x}^{4}$ .

(1)求 $\frac{{a}_{1}}{2} + \frac{{a}_{2}}{4} + \frac{{a}_{3}}{8} + \frac{{a}_{4}}{16}$ 的值；

(2) 求 ${a}_{1} + 2{a}_{2} + 3{a}_{3} + 4{a}_{4}$ 的值.

15. 已知二项式 ${\left( a{x}^{2} + \frac{1}{\sqrt{2x}}\right) }^{5}\left( {a > 0}\right)$ 的展开式中第二项与第四项的系数相同.

(1)求 $a$ 的值与展开式中的常数项；

(2)求展开式中系数最大的项.

16. 记 ${\left( 1 - 2x\right) }^{{2n} + 1} = {a}_{0} + {a}_{1}x + {a}_{2}{x}^{2} + \cdots  + {a}_{{2n} + 1}{x}^{{2n} + 1},n$ 是正整数,

(1)求 $\left| {a}_{0}\right|  + \left| {a}_{1}\right|  + \cdots  + \left| {a}_{{2n} + 1}\right|$ ；

(2)设 ${a}_{k} = {\left( -2\right) }^{k}{b}_{k}$ ，求和: ${b}_{0} + 2{b}_{1} + 3{b}_{2} + \cdots  + \left( {k + 1}\right) {b}_{k} + \cdots  + \left( {{2n} + 2}\right) {b}_{{2n} + 1}$ .

\section*{第 12 章 概率与统计(必修 3 \& 选择性必修 2)}

\section*{知识梳理}

\section*{一、随机现象与样本空间}

1. 随机现象:具不确定性的现象称为随机现象，或者说具有随机性.

2. 必然现象和不可能现象

必然现象:必然会出现的现象；不可能现象:必然不会出现的现象.

3. 随机试验: 对可随意重复的随机现象的观察、记录、试验统称为随机试验.

4. 基本事件:随机试验的所有可能的结果都叫做试验结果或基本事件。

5. 样本空间与基本事件:一个随机现象中依某个角度观察其所有可能出现(发生)的结果所组成的集合称为一个样本空间,用 $\Omega$ 表示,其中的元素称为基本事件或者样本点. 记作 $\Omega  = \left\{  {{\omega }_{1},{\omega }_{2},{\omega }_{2},\cdots ,{\omega }_{n}}\right\}$ .

注:样本空间中的基本事件必须是互斥的，即任何两个均不会同时发生.

(1)“样本空间”，这里的“空间”不是生活中所指物理意义上的三维空间，而是随机现象可能的结果所组成的集合.

(2)基本事件(互斥)是样本空间中的元素，“事件”是样本空间的子集.

\section*{二、古典概率}

1. 等可能性:若一个随机试验的所有结果出现的可能性都一样，则称之为具有等可能性.

2.古典概率模型:如果一个随机试验满足下面两个条件:

(1)包含有限个可能出现的结果(基本事件);

(2)这些结果出现是等可能的，那么这样的随机试验就称为古典概率模型.

3. 古典概型中概率的计算公式:

由于一般的随机事件 $A$ 是基本事件的某个集合，即样本空间的一个子集，用符号 $P\left( A\right)$ 表示事件 $A$ 发生的概率,那么事件 $A$ 发生的概率等于事件 $A$ 中的基本事件个数除以样本空间中的基本事件总数,即 $P\left( A\right)  = \frac{\left| A\right| }{\left| \Omega \right| }$ . 其中 $\left| A\right|$ 表示事件中的基本事件个数,而 $\left| \Omega \right|$ 表示样本空间中的基本事件个数.

4. 怎么得到一个等可能的样本空间?

通常要将一个随机试验依次分解为若干个等可能的随机试验来处理,方法如下:

设一个随机试验分两步完成,第一步有 $m$ 个等可能的结果,记作 ${a}_{1},{a}_{2},\cdots ,{a}_{m}$ ; 而对第一步得到的每个结果,第二步总有 $n$ 个等可能的结果,记作 ${b}_{1},{b}_{2},\cdots ,{b}_{n}$ . 那么,该随机试验的样本空间就是 $\Omega  = \left\{  {\left( {{a}_{i},{b}_{j}}\right)  \mid  i = 1}\right.$ , $2,\cdots ,m;j = 1,2,\cdots ,n\}$ ,它是等可能的,共有 ${mn}$ 个元素. 对多步的等可能随机试验可以类似地构造等可能的样本空间.

5. 事件的关系与运算.

(1)事件的关系:设事件A对应于子集 $A$ ，事件B对应于子集 $B$ . 如果 $A$ 的基本事件都在 $B$ 中，那么 $A$ 发生必然 $B$ 发生. 此时，称 $B$ 包含 $A$ 或者 $A$ 包含于 $B$ ，即 $A \subseteq  B$ . 或等价地说， $B$ 不发生则 $A$ 也不发生.

(2)事件的运算:事件 $A$ 、 $B$ 至少有一个发生,其对应的子集是 $A \cup  B$ ;

事件 $A\text{ 、 }B$ 同时发生,其对应的子集是 $A \cap  B$ ;

注:事件的关系就是样本空间中相应子集的关系，事件的运算就是相应子集的运算，两者是互相对应的.

\HRule

*注，本章节所属知识模块是“数据与统计”.

\HRule

(1)互斥事件:如果 $A$ 与 $B$ 没有共同的基本事件，即两个子集不相交: $A \cap  B = \varnothing$ ，那么这两个事件不可能同时发生，或者说互斥.

(2)对立事件:“事件 $A$ 发生”的否定就是“事件 $A$ 不发生”，它也是一个事件，称为事件 $A$ 的对立事件， 简称为 “非 $A$ ”. 对应的子集是不属于 $A$ 的基本事件全体,从而是 $A$ 在样本空间 $\Omega$ 中的补集 $\bar{A}$ . 即 $A \cap  \bar{A} =  \varnothing ,A \cup  \bar{A} = \Omega$ .

6. 古典概率模型的概率性质.

概率性质 1 : 必然事件的概率是 1,不可能事件的概率是 0,即 $P\left( \Omega \right)  = 1,P\left( \varnothing \right)  = 0$ .

概率性质 2 :设 $A$ 是一个事件，那么 $0 \leq  P\left( A\right)  \leq  1$ .

概率性质 3 (概率的可加性): 两个不可能同时发生的事件至少有一个发生的概率是这两个事件的概率之和. 换言之, $A \cap  B = \varnothing$ ,那么 $P\left( {A \cup  B}\right)  = P\left( A\right)  + P\left( B\right)$ .

概率性质 4: 对任一给定事件,其发生的概率与不发生的概率的和总是 1. 换言之,即 $P\left( A\right)  = 1 - P\left( \bar{A}\right)$ .

7. 所有概率模型中事件的概率性质.

(1)任意事件 $A$ 的概率值介于 0、1 之间，即 $0 \leq  P\left( A\right)  \leq  1$ ；

(2)规范性:必然事件的概率为 1，即 $P\left( \Omega \right)  = 1$ ；

(3)可加性:如果事件 $A$ 、 $B$ 互斥，那么 $P\left( {A \cup  B}\right)  = P\left( A\right)  + P\left( B\right)$ .

注:以上 3 条构成了概率的三条基本性质. 不仅仅是古典概率，所有概率模型中事件的概率都满足以上三条基本性质. 其中,从两个事件的可加性可以推出任意多个事件的可加性: 如果 ${A}_{1}\text{ 、 }{A}_{2}\text{ 、 }\cdots \text{ 、 }{A}_{n}$ 是 $n$ 个两两互斥的事件，那么

\[
P\left( {{A}_{1} \cup  {A}_{2} \cup  \cdots  \cup  {A}_{n}}\right)  = P\left( {A}_{1}\right)  + P\left( {A}_{2}\right)  + \cdots  + P\left( {A}_{n}\right) .
\]

\section*{三、频率与概率}

1. 伯努利试验: 如果一个随机试验只有两种可能的结果,一个是“成功”,概率是 $p\left( {0 < p < 1}\right)$ ,一个是“失败”,概率是 $1 - p$ ,那么这个随机试验称为伯努利试验.

2. 频率: 假设我们可以独立地重复一个伯努利试验 $n$ 次,其中成功的次数记作 ${S}_{n}$ ,那么 $\frac{{S}_{n}}{n}$ 就被称为 (n 次试验中)成功的频率.频率是一个数，依赖于试验次数 $n$ . 它不是一个确定的数，而是一个随机的数.

3. 伯努利大数定律: 独立地重复一个伯努利试验 $n$ 次,当 $n$ 很大时,频率 $\frac{{S}_{n}}{n}$ 逼近概率.

4.经验概率:频率在大数次的试验中稳定于某一常数(概率)这个事实非常重要. 正因为如此，在实际中可以把频率作为概率 (的估计值) 来应用. 频率也称为经验概率,计算它通常是为了估计概率 $P$ . 为了区别于概率,经验概率用 $\widehat{P}$ 来表示.

注:这个结论之所以被称为定律，或许是因为它直观地像是一个自然定律. 该定律成立的前提是:独立地重复一个伯努利试验，而且试验次数要很大. 如果试验次数 $n$ 较小，往往所求出的频率的变化幅度就可能很大，这样就无法合理地利用频率来估计概率. 通过不断增加试验次数 $n$ ，计算这一过程中的频率大小、观察其变化, 就会发现频率会逐步稳定于某个常数, 而这个常数就是该随机事件的概率. 伯努利大数定律的成立有一个条件，即“假设我们可以独立地重复一个伯努利试验”，这个条件非常关键. 抛掷硬币、掷骰子这类随机试验可以独立地重复，然而许多随机现象是不可以独立地重复的. 例如，某人今年会不会得流感是随机的，每个人的高考成绩也是随机的，但这些现象都不能独立地重复。可以独立地重复的才称得上是一个试验. 虽然人们对于不能独立地重复的随机现象也谈论概率, 但那是主观的概率, 并不能检验.

\section*{四、随机事件的独立性}

独立随机事件定义:两个事件 $A$ 与 $B$ (相互)独立是指它们同时发生的概率等于它们各自发生概率的乘积,即 $P\left( {A \cap  B}\right)  = P\left( A\right)  \cdot  P\left( B\right)$ .

\section*{五、条件概率及其相关公式}

1. 条件概率.

事件 $A$ 发生之后事件 $B$ 发生的概率,称为事件 $B$ 基于条件 $A$ 的概率,或在事件 $A$ 发生的条件下,事件 $B$ 发生的概率,或已知事件 $A$ 发生,事件 $B$ 发生的概率,记为 $P\left( {B \mid  A}\right)$ ,且 $P\left( {B \mid  A}\right)  = \frac{P\left( {A \cap  B}\right) }{P\left( A\right) }$ .

在古典概率模型中,事件 $A$ 发生之后,随机现象的结果就剩下事件 $A$ 中的基本事件,所以事件 $A$ 变成了由这些基本事件所构成的新的样本空间. 这个样本空间仍然是等可能的，这时 $P\left( {B \mid  A}\right)  = \frac{\left| A \cap  B\right| }{\left| A\right| }$ .

2. 乘法公式.

事件 $A$ 与事件 $B$ 同时发生的概率: $P\left( {A \cap  B}\right)  = P\left( A\right) P\left( {B \mid  A}\right)$ ,

或者 $P\left( {A \cap  B}\right)  = P\left( B\right) P\left( {A \mid  B}\right)$ .

3. 样本空间 $\Omega$ 的一个划分.

设样本空间 $\Omega$ 可分成 $n$ 个两两互斥的事件 ${\Omega }_{1}\text{ 、 }{\Omega }_{2}\text{ 、 }\cdots \text{ 、 }{\Omega }_{n}$ ,即

(1) ${\Omega }_{1} \cup  {\Omega }_{2} \cup  \cdots  \cup  {\Omega }_{n} = \Omega$ ；

(2) ${\Omega }_{i} \cap  {\Omega j} = \varnothing \left( {i \neq  j,1 \leq  i,j \leq  n}\right)$ ，

则称 ${\Omega }_{1},{\Omega }_{2},\cdots ,{\Omega }_{n}$ 构成了样本空间 $\Omega$ 的一个划分.

注: 一个事件 $A$ 发生可以看作是在不同情况下分别发生,即 $A = \left( {A \cap  {\Omega }_{1}}\right)  \cup  \left( {A \cap  {\Omega }_{2}}\right)  \cup  \ldots  \cup  \left( {A \cap  {\Omega }_{n}}\right)$ .

4. 全概率公式.

设 ${\Omega }_{1}\text{ 、 }{\Omega }_{2}\text{ 、 }\cdots \text{ 、 }{\Omega }_{n}$ 构成了样本空间 $\Omega$ 的一个划分,且,则对 $\Omega$ 中的任一事件 $A$ 都有:

$P\left( A\right)  = \mathop{\sum }\limits_{{k = 1}}^{n}P\left( {A \mid  {\Omega }_{k}}\right) P\left( {\Omega }_{k}\right) .$

上式称为全概率公式.

5.* 贝叶斯公式.

若事件 ${\Omega }_{1}\text{ 、 }{\Omega }_{2}\text{ 、 }\cdots \text{ 、 }{\Omega }_{n}$ 构成了样本空间 $\Omega$ 的一个划分,且 $P\left( {\Omega }_{i}\right)  > 0\left( {i = 1,2,\cdots ,n}\right) ,A$ 为 $\Omega$ 中的任一事件, $P\left( A\right)  > 0$ ,则 $P\left( {{\Omega }_{i} \mid  A}\right)  = \frac{P\left( {A \cap  {\Omega }_{i}}\right) }{P\left( A\right) } = \frac{P\left( {\Omega }_{i}\right) P\left( {A \mid  {\Omega }_{i}}\right) }{\mathop{\sum }\limits_{{k = 1}}^{n}P\left( {\Omega }_{k}\right) P\left( {A \mid  {\Omega }_{k}}\right) }$ .

这个公式称为贝叶斯公式, 也称为后验公式.

\section*{六、随机变量}

1. 以样本空间 (是实数) 作为定义域的一个函数 $X$ 称为一个随机变量,即对样本空间 $\Omega$ 中任意给定的元素 $\omega$ ,都有唯一的实数 $X\left( \omega \right)$ 与之对应.

2. 随机变量所有可能的取值及相应的概率, 称为随机变量的分布, 一般来说, 对于取值有限的离散型随机变量，我们可以通过如下的方式表示随机变量的分布:

\[
\left( \begin{array}{llll} {x}_{1} & {x}_{2} & \cdots & {x}_{n} \\  {p}_{1} & {p}_{2} & \cdots & {p}_{n} \end{array}\right)
\]

其中, ${x}_{1}{x}_{2}\cdots {x}_{n}$ 为随机变量所有可能的取值, ${p}_{k}$ 为 ${x}_{k}$ 对应的概率,即 ${p}_{k} = P\left( {X = {x}_{k}}\right)$ .

因此自然要满足 $0 \leq  {p}_{k} \leq  1$ 和 $\mathop{\sum }\limits_{{k = 1}}^{n}{p}_{k} = 1$ .

分布中 ${p}_{k}$ 的值一般为正数,因为概率为 0 这一列总可以删去. 分布直观上表示了总值为 1 的量是如何分布在随机变量 $X$ 的取值上的.

两种简单的分布:

 \ding{172} 如果随机变量取所有值时，对应的概率相等，则称这个随机变量满足均匀分布，它的分布为:

\[
\left( \begin{matrix} {x}_{1} & {x}_{2} & \cdots & {x}_{n} \\  \frac{1}{n} & \frac{1}{n} & \cdots & \frac{1}{n} \end{matrix}\right)
\]

 \ding{173} 如果随机变量的取值只有两个，则称为伯努利型随机变量，对应的分布称为伯努利分布，即: $\left( \begin{matrix} {x}_{1} & {x}_{2} \\  p & 1 - p \end{matrix}\right)$ ；

通常一个只有成功或者失败的随机实验就可以以成功记为 1 , 失败记为 0 构成一个伯努利型随机变量. 七、期望与方差

1. 定义:对于一个取值有限项的离散型随机变量 $X$ ，如果它的分布为:

$\left( \begin{array}{llll} {x}_{1} & {x}_{2} & \cdots & {x}_{n} \\  {p}_{1} & {p}_{2} & \cdots & {p}_{n} \end{array}\right)$ ,那么可以定义随机变量的期望为 $E\left\lbrack  X\right\rbrack   = \mathop{\sum }\limits_{{k = 1}}^{n}{x}_{k} \cdot  {p}_{k}$ ,

方差为 $D\left\lbrack  X\right\rbrack   = E\left\lbrack  {\left( X - E\left\lbrack  X\right\rbrack  \right) }^{2}\right\rbrack   = E\left\lbrack  {X}^{2}\right\rbrack   - {\left( E\left\lbrack  X\right\rbrack  \right) }^{2}$ ,前者为方差的定义式,后者是推导的结果,有时更便于计算.

2.期望与方差的意义:

 \ding{172} 期望实际上就是随机变量的加权平均数，反映了随机变量平均取值的大小. 是一个确定的实数.

 \ding{173} 对于无限项的离散型随机变量，期望的定义依然是随机变量取值的加权平均，只不过具体计算时，通常不用定义来计算，具体参考课本中关于二项分布的期望的计算.

 \ding{174} 对于连续型随机变量(高中阶段只涉及正态分布)，其期望是由积分定义的，但其意义依然是表示了随机变量取值的平均.

 \ding{175}  方差表示的是随机变量的分散程度，方差越大，随机变量的分布越分散、越不稳定，反之则越集中、越稳定.

 \circled{5}  $D\left\lbrack  X\right\rbrack   \geq  0$ ，等号当且仅当随机变量 $X$ 以百分百的概率取一定值， $\sqrt{D\left\lbrack  X\right\rbrack  }$ 称为随机变量的标准差.

3.期望与方差的性质:

如果 $X\text{ 、 }Y$ 是随机变量, $a\text{ 、 }b$ 是实数,则:

 \ding{172}  $E\left\lbrack  {{aX} + b}\right\rbrack   = {aE}\left\lbrack  X\right\rbrack   + b$ ；

 \ding{173}  $E\left\lbrack  {X + Y}\right\rbrack   = E\left\lbrack  X\right\rbrack   + E\left\lbrack  Y\right\rbrack$ ；

 \ding{174}  $D\left\lbrack  {{aX} + b}\right\rbrack   = {a}^{2}D\left\lbrack  X\right\rbrack$ ；

 \ding{175}  当 $X\text{ 、 }Y$ 是独立的随机变量时, $D\left\lbrack  {X + Y}\right\rbrack   = D\left\lbrack  X\right\rbrack   + D\left\lbrack  Y\right\rbrack$ .

\section*{八、常用分布}

1. 二项分布:独立地重复一个成功概率为 $p$ ，失败概率为 $1 - p$ 的伯努利实验 $n$ 次，记这 $n$ 次中成功的次数为 $X$ ,则称 $X$ 是一个服从参数为 $n\text{ 、 }p$ 的二项分布,记为 $X \sim  B\left( {n,p}\right)$ .

性质: 如果 $X \sim  B\left( {n,p}\right)$ ,则:

 \ding{172}  $X$ 的分布为 $P\left( {X = k}\right)  = {{C}}_{n}^{k}{p}^{k}{\left( 1 - p\right) }^{n - k}\left( {k = 0,1,2,\cdots ,n}\right)$ ;

 \ding{173}  $E\left\lbrack  X\right\rbrack   = {np},D\left\lbrack  X\right\rbrack   = {np}\left( {1 - p}\right)$ .

2. 超几何分布: 从装有大小与质地相同的 $a$ 个白球和 $b$ 个黑球的袋中不放回地随机取 $n$ 个球 $(n \leq  a +  b)$ ,记 $X$ 为取到的白球的个数,则随机变量 $X$ 的分布称为超几何分布.

性质:

 \ding{172} 如果约定当 $k < 0$ 或者 $k > n$ 时，有 ${{C}}_{n}^{k} = 0$ ，则 $X$ 的分布由 $P\left( {X = k}\right)  = \frac{{{C}}_{a}^{k} \cdot  {{C}}_{b}^{n - k}}{{{C}}_{a + b}^{n}}$ 给出，注意: $k \leq  n$ ， $k \leq  a,n - k \leq  b,n \leq  a + b;$

 \ding{173}  $E\left\lbrack  X\right\rbrack   = \frac{na}{a + b}$ ；

 \ding{174} 二项分布和超几何分布的区别,实际上就是摸球模型中放回摸球和不放回摸球的区别,当 $a + b$ 远大于 $n$ 时,两者可以近似.

3.正态分布:

(1)定义:数学中正态分布是指由下面的函数所表达的分布:

\[
{\varphi }_{\mu ,{\sigma }^{2}}\left( x\right)  = \frac{1}{\sqrt{{2\pi }{\sigma }^{2}}}{{e}}^{-\frac{\left( {x - \mu }\right) 2}{2{\sigma }^{2}}}, - \infty  < x <  + \infty ,
\]

这个图形(下图)形似钟型，因此称为钟形曲线。

\begin{center}
\includegraphics[width=0.6\textwidth]{images/bo_d7tk7oc91nqc73875fu0_246_426_217_744_393_0.jpg}
\end{center}
 

此时称随机变量 $X$ 是服从参数为 $\mu$ 、 ${\sigma }^{2}$ 的正态分布，记为 $X \sim  N\left( {\mu ,{\sigma }^{2}}\right)$

其中, ${\varphi }_{\mu ,{\sigma }^{2}}\left( x\right)$ 是 $X$ 的概率密度函数,且 $\mu$ 和 ${\sigma }^{2}$ 分别为随机变量的期望和方差 (总是假设 $\sigma  > 0$ )

(2)钟形曲线的性质

 \ding{172}  钟形曲线关于 $x = \mu$ 对称，总在 $x$ 轴上方，以 $x$ 轴为渐近线.

 \ding{173}  ${\sigma }^{2}$ 越小，钟形曲线越“瘦高”， ${\sigma }^{2}$ 越大，钟形曲线越“矮胖”

 \ding{174} 如下图， $P\left( {b < X < a}\right)$ 等于由钟形曲线、 $x$ 轴、 $x = a$ 、 $x = b$ 围成的面积.

\begin{center}
\includegraphics[width=0.3\textwidth]{images/bo_d7tk7oc91nqc73875fu0_246_588_881_435_267_0.jpg}
\end{center}
 

 \ding{175} 钟形曲线与 $x$ 轴围成的总面积为 1 .

(3)标准正态分布

 \ding{172}  设 $X \sim  N\left( {\mu ,{\sigma }^{2}}\right)$ ，如果 $\mu  = 0$ 、 ${\sigma }^{2} = 1$ ，则称 $X$ 服从标准正态分布，记为 $X \sim  N\left( {0,1}\right)$ .

此时对应的概率密度函数可以简记为 $\varphi \left( x\right)$ .

 \ding{173} 如果 $X \sim  N\left( {\mu ,{\sigma }^{2}}\right)$ ,令 $Y = \frac{X - \mu }{\sigma }$ ,则 $Y \sim  N\left( {0,1}\right)$ .

 \ding{174} 如果 $X \sim  N\left( {0,1}\right)$ ，则用 $\Phi \left( x\right)$ 表示标准正态分布的概率密度函数 $\varphi \left( x\right)$ 从 $- \infty$ 到 $x$ 的累积面积，即 $\Phi \left( x\right)  = P\left( {X < x}\right) .\Phi \left( x\right)$ 的值可以通过查表获得.

标准正态分布的意义在于可以将一般的正态分布先标准化, 再查表, 进而计算一般的正态分布的相关信息.

\section*{九、统计初步}

1. 总体与样本

(1)总体与个体

在统计问题中，我们把研究对象的全体叫做总体，总体中的每一个对象叫做个体.

(2)总体容量与样本容量

总体中所含个体的数量, 称为总体的容量. 从总体中抽取的一部分个体叫做这个总体的一个样本, 样本所含个体的数量称为样本量, 也称样本容量.

2. 数据的获取

(1)观测数据与实验数据

观测数据是指通过调查或观测而收集到的数据，是在没有对有关事物进行人为控制的条件下得到的。以全国人口普查为例，人口普查员通过请社区住户填写人口普查表来调查人口信息，但并不进行干预，这样得到的数据就是观测数据. 实验数据是指在实验中控制实验对象而收集到的数据. 例如，我们想知道每天服用一定剂量的维生素C能否预防感冒。为了完成实验，研究人员让一部分人每天服用一定剂量的维生素C，另一部分人不服用维生素C，然后收集两组对象的感冒发病率并进行比较.

(2)普查与抽查

对总体的每个个体分别进行调查，我们称之为普查. 从总体中抽取样本的过程称为抽样，通过抽样进行调查研究的方法叫做抽样调查.

3. 抽样方法

(1)简单随机抽样

在抽样的过程中通过逐个抽取的方法抽取样本, 且总体的每一个个体都有同样的可能性被选入样本, 这种抽样方法叫做简单随机抽样.

 \ding{172} 抽签法的定义

一般地，抽签法就是把总体中的 $N$ 个个体编号，把号码写在号签上，将号签放在一个容器中，搅拌均匀后,每次从中抽取一个号签,连续抽取 $n$ 次,就得到一个容量为 $n$ 的样本.

 \ding{173} 随机数表法

随机数表 (随机数表不唯一,例如见课本附录) 是一种现成的号签,它是由随机数 (通常为 $0,1,2,\cdots ,9$ ) 组成的数表,表中的数是由一连串的 0-9 之间的数字构成的,且满足以下两个条件:

(i) 对表中任何一个给定位置的数字,其为 0-9 中任何一个数字的概率都相等;

(ii) 不同位置数字的取值是相互独立的. 下列是一个随机数表的一部分.

我们先把总体中的个体逐个编号, 然后按照某种事先确定的规则直接从随机数表 (或其他工具产生的随机数) 中取一定量的随机数, 这些随机数所对应的个体即组成一个样本. 我们称这种抽样方法为随机数法.

(2)分层(随机)抽样

一般地，当总体由差异明显的几个部分组成时，先把总体分成若干个部分，然后从不同的部分中独立、 随机地抽取样本，这种抽样的方法称为分层随机抽样，简称分层抽样.

4.统计图表

(1)频率分布表

 \ding{172}  求极差. 绘制频率分布表最重要的就是分组. 要对数据进行分组，首先要找出这一组数据的最大值和最小值, 最大值与最小值的差称为极差, 又称全距.

 \ding{173} 确定组距与组数. 组距是指每个小组的区间端点之间的距离，组距的选取决定了组数的多少:组数 = 极差. 如果组数太少, 所有数据就会落在少数几个组里, 导致这几个组的柱形太高; 如果组数太多, 大部分的组中只有少数的数据或者没有数据，均无法有效呈现数据的分布. 一般地，第一组的下限应低于最小值，最后一组的上限应高于最大值.

 \ding{174} 统计每组的频数及频率.

将数据分好组以后，每个小组内的数据个数称为频数，频数与样本容量的比值叫做这一小组的频率. 统计每个小组的频数, 再计算各组的频率.

 \ding{175} 绘制频率分布表. 将分组、频数及频率填入频率分布表.

(2)频率分布直方图

依据频率分布表, 制作频率分布直方图. 为此, 我们在直角坐标系中把横轴分成若干段, 每一段对应一个组距, 然后以组距形成的线段为底作一矩形, 矩形的高为小组的组距离, 就可以得到频率分布直方图.

(3)频率分布折线图

在频率分布直方图中, 按照分组原则, 再在左边和右边各加一个区间, 这两个组的频率取值为 0 . 然后从左边第一个区间的中点 (称为组中值) 开始, 从左至右依次连结个矩形上底边的中点, 直至右边所加区间的中点, 再将矩形的边去除, 就可以得到一条折线, 称为频率分布折线图, 简称频率折线图.

(4)总体分布密度曲线

如果将样本容量取得足够大，分组的组距取得足够小，那么相应的频率折线图将趋于一条光滑曲线，我们称这条光滑曲线为总体分布的密度曲线.

(5)茎叶图

顾名思义，茎是指中间的一列数，叶就是从茎的旁边生长出来的数. 如，当数据是两位有效数字时，用中间的数字表示十位数，两边的数字表示个位数.

(6) 散点图

在考虑两组数据时, 为了对两组数据之间的关系形成大致的了解, 通常将这两组数据所对应的点描在平面直角坐标系内, 这些点组成的统计图称为散点图.

5.统计估计

(1)估计总体的分布

总体的分布指的是总体中不同范围或类型的个体所占的比例.

(2)估计总体的数字特征

我们在初中阶段学习过的平均数、中位数和众数等，就是用来刻画一组数据集中趋势的数字特征，而方差和标准差则是用来刻画一组数据离散程度的数字特征.

已知一组数据的样本 ${x}_{1}\text{ 、 }{x}_{2}\text{ 、 }{x}_{3}\text{ 、 }\cdots \text{ 、 }{x}_{n}$ ( $n$ 为正整数),则这组数据的样本数据的平均数:

\[
\bar{x} = \frac{{x}_{1} + {x}_{2} + {x}_{3} + \cdots  + {x}_{n}}{n};
\]

样本数据的方差:

\[
{s}^{2} = \frac{1}{n}\left\lbrack  {{\left( {x}_{1} - \bar{x}\right) }^{2} + {\left( {x}_{2} - \bar{x}\right) }^{2} + {\left( {x}_{3} - \bar{x}\right) }^{2} + \cdots  + {\left( {x}_{n} - \bar{x}\right) }^{2}}\right\rbrack   = \frac{1}{n}\mathop{\sum }\limits_{{i = 1}}^{n}{\left( {x}_{i} - \bar{x}\right) }^{2}
\]

样本数据的标准差:

\[
{s}^{2} = \sqrt{\frac{{\left( {x}_{1} - \bar{x}\right) }^{2} + {\left( {x}_{2} - \bar{x}\right) }^{2} + {\left( {x}_{3} - \bar{x}\right) }^{2} + \cdots  + {\left( {x}_{n} - \bar{x}\right) }^{2}}{n}} = \sqrt{\frac{1}{n}\mathop{\sum }\limits_{{i = 1}}^{n}{\left( {x}_{i} - \bar{x}\right) }^{2}}
\]

(3)估计百分位数

 \ding{172}  百分位数定义:第 $k$ 百分位数， $k$ 为 1 到 100 之间的整数，记作 ${P}_{k}$ ) 即是将一组数据从小到大排列后， 将数据分成两部分:小于或等于第 $k$ 百分位数的数据占 $k\%$ ，大于或等于第 $k$ 百分位数的数据占 $\left( {{100} - k}\right) \%$ .

 \ding{173} 计算一组 $n$ 个数据的第 $k$ 百分位数的一般步骤如下:

第 1 步，按从小到大排列原始数据；

第 2 步,计算 $i = n \cdot  k\%$ ;

第 3 步,若 $i$ 是整数,则第 $k$ 百分位数为第 $i$ 项与第 $\left( {i + 1}\right)$ 项的数据的平均数; 若 $i$ 不是整数,则将 $i$ 向上取整,则第 $k$ 百分位数为第 $i + 1$ 项数据.

(4)统计活动的主要过程

确定调查方案、确定调查对象、收集数据、数据整理与分析、初步结论

\section*{十、成对数据的相关分析}

我们把这样来自同一对象的两组数据称为成对数据，研究成对数据相关性的方法称为相关分析.

1. 相关系数

散点图可反映两个变量是否存在相互性或相关方向, 但不能定量分析. 于是, 著名统计学家卡尔·皮尔逊设计了统计指标——相关系数. 相关系数是用以反映变量之间相关关系密切程度的统计指标.

\begin{center}
\begin{tabular}{|c|c|c|c|c|c|}
\hline
变量 $x$ & ${x}_{1}$ & ${x}_{2}$ & ${x}_{3}$ & ... & ${x}_{n}$ \\
\cline{1-6}
变量 $y$ & ${y}_{1}$ & ${y}_{2}$ & ${y}_{3}$ & ... & ${y}_{n}$ \\
\cline{1-6}
\hline
\end{tabular}
\end{center}

可以将 $\left( {{x}_{1},{x}_{2},\cdots ,{x}_{n}}\right)$ 和 $\left( {{y}_{1},{y}_{2},\cdots ,{y}_{n}}\right)$ 看成两个高维向量,对于上述数据,规定变量 $x$ 和 $y$ 的线性相关系数为:

\[
r = \frac{\mathop{\sum }\limits_{{i = 1}}^{n}\left( {{x}_{i} - \bar{x}}\right) \left( {{y}_{i} - \bar{y}}\right) }{\sqrt{\mathop{\sum }\limits_{{i = 1}}^{n}{\left( {x}_{i} - \bar{x}\right) }^{2}}\sqrt{\mathop{\sum }\limits_{{i = 1}}^{n}{\left( {y}_{i} - \bar{y}\right) }^{2}}} = \frac{\mathop{\sum }\limits_{{i = 1}}^{n}{x}_{i}{y}_{i} - n\bar{x}\bar{y}}{\sqrt{\mathop{\sum }\limits_{{i = 1}}^{n}{x}_{i}^{2} - n{\bar{x}}^{2}}\sqrt{\mathop{\sum }\limits_{{i = 1}}^{n}{y}_{i}^{2} - n{\bar{y}}^{2}}}
\]

2. 相关系数的性质

可以证明,相关系数 $r$ 的值满足 $\left| r\right|  \leq  1.\left| r\right|$ 越接近 1,两个变量的线性相关程度越高; $\left| r\right|$ 越接近 0,两个变量的线性相关程度越低. $r > 0$ 时,当 $x$ 的值由小变大, $y$ 的值具有由小变大的变化趋势,称这种相关为正相关; $r < 0$ 时,当 $x$ 的值由小变大, $y$ 的值具有由大变小的变化趋势,称这种相关为负相关.

相关系数 $r$ 描述的是两个变量之间线性关系的方向与程度,是一种定量分析的方法. 相关系数具有以下特点:

(1)相关系数的计算公式关于 $x$ 和 $y$ 这两个变量是对称的. 画散点图时，不论以哪个变量作为横轴 (纵轴), 所得的相关系数都一样.

(2)两个变量的相关系数与这两个变量的单位无关. 例如，在计算身高与体重的相关系数时，身高单位不管取“米”还是“厘米”，相关系数的结果都一样.

(3)与平均数和标准差一样，相关系数不仅会受到数据量多少的影响，也会受到少数异常值较大的影响.

\section*{十一、一元线性回归分析}

\section*{1. 离差和拟合误差}

一般地,设给定一组有线性相关关系的成对数据 $\left( {{x}_{1},{y}_{1}}\right) \text{ 、 }\left( {{x}_{2}{y}_{2}}\right) \text{ 、 }\cdots \text{ 、 }\left( {{x}_{n},{y}_{n}}\right)$ 和一个线性方程 (或称线性模型) $y = {ax} + b$ .

当变量 $x$ 取值 ${x}_{i}\left( {i = 1,2,\cdots ,n}\right)$ 时,令 ${\widehat{y}}_{i} = a{x}_{i} + b$ ,它是变量 $y$ 与 ${x}_{i}$ 对应的理想值.

但数据中的 ${y}_{i}$ 与 ${\widehat{y}}_{i}$ 不一定相同,它们的差 ${y}_{i} - {\widehat{y}}_{i}$ 称为在 ${x}_{i}$ 处的离差,当

${y}_{i} - {\widehat{y}}_{i} \geq  0$ 时称为正离差,而当 ${y}_{i} - {\widehat{y}}_{i} < 0$ 时称为负离差.

显然, 离差直观地描述了单对数据与线性方程 \ding{172} 的贴近度.

由于离差可正可负, 考虑数据整体与线性方程 \ding{172} 的贴近度时, 不能简单地用离差的代数和作为指标. 我们可以像计算方差那样,用离差的平方和 $Q = \mathop{\sum }\limits_{{i = 1}}^{n}{\left( {y}_{i} - {\widehat{y}}_{i}\right) }^{2}$ 来刻画直线与点之间的拟合程度. $Q$ 称为拟合误差, 它是一个很好的描述数据与函数贴合程度的指标.

2. 线性回归方程

我们把拟合误差取得最小值时得到的线性方程 (线性模型) 记为 $y = \widehat{a}x + \widehat{b}$  \ding{173} 

并称之为变量 $y$ 随 $x$ 波动的回归方程或回归模型,其中自变量 $x$ 称为解释变量,因变量 $y$ 称为反应变量. 回归方程所定义的直线称为回归直线，回归方程的系数 (或称回归模型的参数) $\widehat{a}$ 与 $\widehat{b}$ 称为回归系数. 由一组有某种线性关系的成对数据求其回归方程的方法称为一元线性回归分析. 回归系数 $\widehat{a}$ 与 $\widehat{b}$ 的计算公式如下:

\[
\left\{  \begin{array}{l} \widehat{a} = \frac{\mathop{\sum }\limits_{{i = 1}}^{n}\left( {{x}_{i} - \bar{x}}\right) \left( {{y}_{i} - \bar{y}}\right) }{\mathop{\sum }\limits_{{i = 1}}^{n}{\left( {x}_{i} - \bar{x}\right) }^{2}} = \frac{\mathop{\sum }\limits_{{i = 1}}^{n}{x}_{i}{y}_{i} - n\bar{x}\bar{y}}{\mathop{\sum }\limits_{{i = 1}}^{n}{x}_{i}^{2} - n{\bar{x}}^{2}}, \\  \widehat{b} = \bar{y} - \widehat{a}\bar{x} \end{array}\right.
\]

其中 $\bar{x}$ 与 $\bar{y}$ 分别是数据 ${x}_{i}$ 与 ${y}_{i}\left( {i = 1,2,\cdots ,n}\right)$ 的算术平均值,数对 $\left( {\bar{x},\bar{y}}\right)$ 称为样本点的中心,从回归系数公式 \ding{174} 可以看出，回归直线经过样本点的中心 $\left( {\bar{x},\bar{y}}\right)$ ，也就是散点图中数据点的中心.

回归系数 $\widehat{a}$ 与 $\widehat{b}$ 的值只与给定的数据有关，而与回归分析过程中初始选择的线性模型无关. 有了回归系数数 $\widehat{a}$ 与 $\widehat{b}$ ,代入方程 \ding{173} 就得到了这一组成对数据的回归方程这种回归分析的方法称为最小二乘法,由最小二乘法导出的估计量称为最小二乘估计量,所得到的回归系数 $\left( {\widehat{a}\text{ 与 }\widehat{b}}\right.$ 又称为模型参数 $)a$ 与 $b$ 的最小二乘估计.

\section*{十二、2×2 列联表}

1. 分类变量与 2 × 2 列联表

\begin{center}
\begin{tabular}{|c|c|c|c|}
\hline
\phantom{X} & $A$ 组 & $B$ 组 & 总计 \\
\cline{1-4}
0 & $a$ & $b$ & $a + b$ \\
\cline{1-4}
1 & $c$ & $d$ & $c + d$ \\
\cline{1-4}
总计 & $a + c$ & $b + d$ & $n = a + b + c + d$ \\
\cline{1-4}
\hline
\end{tabular}
\end{center}

其中， $a$ 、 $b$ 、 $c$ 、 $d$ 为实际观察值. 这种变量的不同“值”表示对象所属的不同类别，这类变量称为分类变量. 两个分类变量分别占两行和两列, 形成 4 个格子, 每个格子中的数据是同时满足所在行列对应类别的个体的频数. 这些数据都是通过实际调查得到的, 称为观察值. 这些观察值形成的 2 行、 2 列的频数表格，称为 2 行 2 列列联表,简称 $2 \times  2$ 列联表,也称为四格表.

要检验两个随机变量是否有关, 统计上一般先假设它们没有关系, 即相互独立, 再进行统计检验. 这种假设称为原假设，也称为零假设，习惯上用 ${H}_{0}$ 表示。以上述问题为例，我们提出的原假设是:

${H}_{0}$ : 两个随机变量没有关系,即它们相互独立.

要检验上述假设,我们需要对 $2 \times  2$ 列联表中的观察值与预期值进行比较. 预期值是当原假设 ${H}_{0}$ 成立的成时的预期结果. 例如,由 $2 \times  2$ 联表可知,总计 $n$ 位样本对象中有 $a + b$ 位状态为 0,其占比为 $\frac{\left( a + b\right) }{n}$ . 假设两个随机变量没有关系,那么 $a + c$ 位 $A$ 组对象中应该有 $\frac{\left( {a + b}\right) \left( {a + c}\right) }{n}$ 位状态为 0,这里的 $\frac{\left( {a + b}\right) \left( {a + c}\right) }{n}$ 就是原假设 ${H}_{0}$ 成立时计算得到的预期值.

为了描述观察值与预期值之间的总体偏差，我们引入统计量 ${\chi }^{2} : {\chi }^{2} = \sum \frac{{\left( \text{ 观察值一预期值 }\right) }^{2}}{\text{ 预期值 }}$ ，经过变形可得 ${\chi }^{2}$ 的一般计算公式:

${\chi }^{2} = \frac{n{\left( ad - bc\right) }^{2}}{\left( {a + b}\right) \left( {c + d}\right) \left( {a + c}\right) \left( {b + d}\right) }$ . 其中, $n = a + b + c + d$ .

${\chi }^{2}$ 的值越大，说明 $2 \times  2$ 列联表中观察值与预期值的总体偏差越大，原假设成立的可能性就越小. 那么究竟 ${\chi }^{2}$ 多大时，我们才可以拒绝原假设呢？这涉及 ${\chi }^{2}$ 分布. 通过查阅 ${\chi }^{2}$ 分布概率表，可以得到 ${\chi }^{2}$ 值超过某些界限的概率. 例如， $P\left( {{\chi }^{2} \geq  {6.635}}\right)  \approx  {0.01}$ ， $P\left( {{\chi }^{2} \geq  {5.024}}\right)  \approx  {0.025}$ ， $P\left( {{\chi }^{2} \geq  {3.841}}\right)  \approx  {0.05}$ ， $P\left( {{\chi }^{2} \geq  {2.706}}\right)  \approx  {0.1}.$

${0.05}\text{ 、 }{0.01}$ 等小概率值在统计上称为显著性水平,记作 $\alpha$ .

以 $P\left( {{\chi }^{2} \geq  {3.841}}\right)  \approx  {0.05}$ 为例,其含义是: 如果原假设成立,那么 ${\chi }^{2} \geq  {3.841}$ 成立的概率约为 0.05 . 这是一个小概率事件，不太可能发生. 因此我们可以推断原假设 “两个随机变量没有关系”成立的可能性小于 5\%. 或者说， 我们有 95\%的把握认为两个随机变量有关. 虽然小概率事件发生的概率很小，但在统计上由于小概率事件发生而作出拒绝原假设的推断还是有风险的. 显著性水平为 0.05 , 可以理解为错误拒绝的概率不超过 0.05 .

以上所用的 ${\chi }^{2}$ 检验方法在统计学中称为 $2 \times  2$ 列联表独立性检验.

${2.2} \times  2$ 列联表独立性检验步骤

(1)提出两个随机变量没有关系的原假设 ${H}_{0}$ : 患慢性气管炎与吸烟没有关系.

(2)确定显著性水平 $\alpha$ ,教材中规定 $\alpha  = {0.05}$ ,也即 $P\left( {{\chi }^{2} \geq  {3.841}}\right)  \approx  {0.05}$ .

(3)计算统计量 ${\chi }^{2}$ 的值.

(4)统计决断:比较上述 ${\chi }^{2}$ 值与 3.841 的大小，若 ${\chi }^{2}$ 值 $\geq  {3.841}$ ，则拒绝(或否定) ${H}_{0}$ ；若 ${\chi }^{2}$ 值 $< {3.841}$ ， 则不能拒绝 (或否定) ${H}_{0}$ ，即接受 ${H}_{0}$ . 根据上述推断作出结论.

\section*{§ 12.1 古典概率}

\section*{典型例题精讲}

【例 1】某商场举行有奖促销活动, 顾客购买一定金额的商品后即可抽奖, 抽奖方法是: 从装有 2 个红球 ${A}_{1}\text{ 、 }{A}_{2}$ 和 1 个白球 $B$ 的甲箱与装有 2 个红球 ${a}_{1}\text{ 、 }{a}_{2}$ 和 2 个白球 ${b}_{1}\text{ 、 }{b}_{2}$ 的乙箱中各随机摸出 1 个球,若摸出的 2 个球都是红球则中奖, 否则不中奖.

(1)用球的标号列出所有可能的摸出结果的样本空间 $\Omega$ ;

(2)有人认为:两个箱子中的红球比白球多，所以中奖的概率大于不中奖的概率，你认为正确吗？请说明理由.

【例 2】一个盒子里装有三张卡片,分别标记有数字 $1\text{ 、 }2\text{ 、 }3$ ,这三张卡片除标记的数字外完全相同. 随机有放回地抽取 3 次，每次抽取 1 张，将抽取的卡片上的数字依次记为 $a$ 、 $b$ 、 $c$ .

(1)求“抽取的卡片上的数字满足 $a + b = c$ ”的概率；

(2)求“抽取的卡片上的数字 $a$ 、 $b$ 、 $c$ 不完全相同”的概率.

【例 3】两个人比赛, 对于弱者(赢的概率较小者)来说, 一局定胜负和三局两胜比较, 哪个更有利？(这里所说的“三局两胜”是常见的比赛模式，指先赢得两局者为胜，最多三局结束)

【例 4】甲、乙各抛掷若干枚硬币,甲抛掷的硬币总数恰比乙少一枚. 求甲得到的正面数比乙得到的正面数少的概率.

【例 5】甲、乙、丙三位同学进行羽毛球比赛，约定赛制如下:累计负两场者被淘汰:比赛前抽签决定首先比赛的两人，另一人轮空；每场比赛的胜者与轮空者进行下一场比赛，负者下一场轮空，直至有一人被淘汰；当一人被淘汰后，剩余的两人继续比赛，直至其中一人被淘汰，另一人最终获胜，比赛结束. 经抽签，甲、乙首先比赛,丙轮空. 设每场比赛双方获胜的概率都为 $\frac{1}{2}$ .

(1)求甲连胜四场的概率；

(2)求需要进行第五场比赛的概率；

(3)求丙最终获胜的概率.

\section*{四基达标精练}

\section*{一、填空题}

1. 从 1、2、3、6 这 4 个数中一次随机抽取 2 个数，则所取 2 个数的乘积为 6 的倍数的概率是\_\_\_\_\_.

2. 我国数学家陈景润在哥德巴赫猜想的研究中取得了世界领先的成果. 哥德巴赫猜想是“每个大于 2 的偶数可以表示为两个素数的和”,如 ${30} = 7 + {23}$ . 在不超过 30 的素数中,随机选取两个不同的数,其和等于 30 的概率是\_\_\_\_\_.

3. 从正方体的 8 个顶点中任选 4 个，则这 4 个点在同一个平面的概率为\_\_\_\_\_.

4. 为强化安全意识，某商场拟在未来的连续 10 天中随机选择 3 天进行紧急疏散演练，则选择的 3 天恰好为连续 3 天的概率是\_\_\_\_\_.

5. 从 2、3、8、9 中任取两个不同的数字，分别记为 $a$ 、 $b$ ，则 ${\log }_{a}b$ 为整数的概率是\_\_\_\_\_.

6.3 位男同学与 3 位女同学随机排成一行, 其中两端都不是女同学的概率为\_\_\_\_\_.

7. 为了检测学生的身体素质指标，从游泳类 1 项、球类 3 项、田径类 4 项共 8 个项目中随机抽取 4 项进行检测，则每一类都被抽到的概率为\_\_\_\_\_.

8. 某三位数密码, 每位数字可在 0-9 这 10 个数字中任选一个, 则该三位数密码中, 恰有两位数字相同的概率是\_\_\_\_\_.

9. 甲、乙两队进行篮球决赛, 采取七场四胜制 (当一队赢得四场胜利时, 该队获胜, 决赛结束). 根据前期比赛成绩, 甲队的主客场安排依次为“主主客客主客主”. 设甲队主场取胜的概率为 0.6 , 客场取胜的概率为 0.5 ，且各场比赛结果相互独立，则甲队以 4:1 获胜的概率是\_\_\_\_\_.

\section*{二、选择题}

10. 从分别写有1、2、3、4、5、6的 6 张卡片中无放回随机抽取 2 张，则抽到的 2 张卡片上的数字之积是 4 的倍数的概率为 ( )

A. $\frac{1}{5}$ . B. $\frac{1}{3}$ . C. $\frac{2}{5}$ . D. $\frac{2}{3}$ .

11. 有 6 个相同的球,分别标有数字1、2、3、4、5、6，从中有放回的随机取两次，每次取 1 个球. 甲表示事件 “第一次取出的球的数字是 1 ”，乙表示事件“第二次取出的球的数字是 2 ”，丙表示事件“两次取出的球的数字之和是 8 ”，丁表示事件“两次取出的球的数字之和是 7 ”，则 ( )

A. 甲与丙相互独立. B. 甲与丁相互独立. C. 乙与丙相互独立. D. 丙与丁相互独立.

12. 某棋手与甲、乙、丙三位棋手各比赛一盘, 各盘比赛结果相互独立. 已知该棋手与甲、乙、丙比赛获胜的概率分别为 ${p}_{1}\text{ 、 }{p}_{2}\text{ 、 }{p}_{3}$ ,且 ${p}_{3} > {p}_{2} > {p}_{1} > 0$ . 记该棋手连胜两盘的概率为 $p$ ,则 ( )

A. $p$ 与该棋手和甲、乙、丙的比赛次序无关. B. 该棋手在第二盘与甲比赛, $p$ 最大.

C. 该棋手在第二盘与乙比赛， $p$ 最大. D. 该棋手在第二盘与丙比赛, $p$ 最大.

\section*{三、解答题}

13. 设平面向量 ${\overrightarrow{a}}_{m} = \left( {m,1}\right) ,{\overrightarrow{b}}_{n} = \left( {2,n}\right)$ ，其中 $m$ 、 $n \in  \{ 1,2,3,4\}$ .

(1)请写出有序数组 $\left( {m,n}\right)$ 的组成的样本空间 $\Omega$ ;

(2)记“使得 ${\overrightarrow{a}}_{m} \bot  \left( {{\overrightarrow{a}}_{m} - {\overrightarrow{b}}_{n}}\right)$ 成立的 $\left( {m,n}\right)$ ”为事件 $A$ ，求事件 $A$ 发生的概率.

14. 甲、乙、丙各自独立投篮一次. 已知乙投中的概率是 $\frac{2}{3}$ ，甲投中并且丙投中的概率是 $\frac{3}{8}$ ，乙投不中并且丙投不中的概率是 $\frac{1}{6}$ .

(1)求甲投中的概率；

(2)求甲、乙、丙 3 人中恰有 2 人投中的概率.

15. 判断下列事件 $A$ 与 $B$ 是否相互独立.

(1)从一副去掉大小王的 52 张扑克牌中随机抽取一张牌,事件 $A$ 表示 “取得的牌面数是 10”,事件 $B$ 表示 “取得的牌的花色是红桃”；

(2)掷一颗骰子，事件 $A$ 表示 “结果是偶数”，事件 $B$ 表示 “结果是奇数”；

16. 设每个工作日甲、乙、丙、丁 4 人需使用某种设备的概率分别为0.6、0.5、0.5、0.4、6各人是否需使用设备相互独立.

(1)求同一工作日至少 3 人需使用设备的概率;

(2)实验室计划购买 $k$ 台设备供甲、乙、丙、丁使用，若要求“同一工作日需使用设备的人数大于 $k$ ”的概率小于 0.1 , 求 $k$ 的最小值.

\section*{$\text{ § }{12.2}$ 条件概率及相关公式}

\section*{典型例题精讲}

【例 1】1 号箱中有 2 个白球和 4 个红球, 2 号箱中有 5 个白球和 3 个红球, 现随机地从 1 号箱中取出一球放入 2 号箱，然后从 2 号箱随机取出一球，问:

(1)从 1 号箱中取出的是红球的条件下，从 2 号箱取出红球的概率是多少？

(2)从 2 号箱取出红球的概率是多少？

【例 2】设某工厂有甲、乙、丙三个车间，它们生产同一种工件，每个车间的产量占该厂总产量的百分比依次为 ${25}\% ,{35}\% ,{40}\%$ ,它们的次品率依次为 $5\% ,4\% ,2\%$ . 现从这批工件中任取一件.

(1)求取到次品的概率；

(2)已知取到的是次品,求它是甲车间生产的概率(精确到 0.01).

【例 3】设在 $n$ 张彩票中有一张奖券, $n$ 个人依次随机地摸奖券.

(1)分别求第一、第二、第三个人摸到奖券的概率；

(2)试证明中奖(摸到奖券)和摸奖顺序无关.

【例 4】设有两个罐子, $A$ 罐中放有 2 个白球、 1 个黑球, $B$ 罐中放有 3 个白球,这些球的大小与质地相同. 现从两个罐子中各摸 1 个球并交换,

(1)求这样交换 2 次后,黑球还在 $A$ 罐中的概率；

(2)记这样交换 $n$ 次后,黑球还在 $A$ 罐中的概率为 $P\left( {A}_{n}\right)$ ，

 \ding{172}  求 $\left\{  {P\left( {A}_{n}\right) }\right\}$ 的递推公式； \ding{173}  求 $\left\{  {P\left( {A}_{n}\right) }\right\}$ 的通项公式.

【例 5】某汽车模型公司共有 25 个汽车模型, 其外观与内饰颜色数量如下表所示.

\begin{center}
\begin{tabular}{|c|c|c|}
\hline
\phantom{X} & 红色外观 & 蓝色外观 \\
\cline{1-3}
棕色内饰 & 8 & 12 \\
\cline{1-3}
米色内饰 & 2 & 3 \\
\cline{1-3}
\hline
\end{tabular}
\end{center}

(1)若小明从这些模型中随机拿一个模型，记事件 $A$ 为小明取到的模型为红色外观，事件 $B$ 取到模型有棕色内饰. 求 $P\left( B\right)$ 、 $P\left( {B \mid  A}\right)$ ，并据此判断事件 $A$ 和事件 $B$ 是否独立；

(2)该公司举行了一个抽奖活动，规定在一次抽奖中，每人可以一次性从这些模型中拿两个汽车模型， 给出以下假设:

假设 1 :拿到的两个模型会出现三种结果，即外观和内饰均为同色、外观内饰都异色，以及仅外观或仅内饰同色;

假设 2: 按结果的可能性大小,概率越小奖项越高;

假设 3: 奖金额为一等奖 600 元,二等奖 300 元,三等奖 150 元.

请你分析奖项对应的结果，设 $X$ 为奖金额，写出 $X$ 的分布和数学期望.

\section*{四基达标精练}

\section*{一、填空题}

1. 从 1、2、3、4、5 中任取 2 个不同的数，事件 $A$ : “取到的 2 个数之和为偶数”，事件 $B$ : “取到的 2 个数均为偶数”，则 $P\left( {B \mid  A}\right)  =$ \_\_\_\_\_.

2. 从标有1、2、3、4、5、6 的六张卡片中，依次不放回地抽出两张，则在第一次抽到奇数的情况下，第二次抽到奇数的概率为\_\_\_\_\_.

3. 计算机中心有三台打字机 $A$ 、 $B$ 、 $C$ ，某打字员使用各台打字机打字的概率依次为 0.6、0.3 和 0.1，打字机发生故障的概率依次为 ${0.01}\text{ 、 }{0.05}\text{ 、 }{0.04}$ . 该打字员因打字机发生故障而耽误了工作进度的概率为\_\_\_\_\_.

4. 假设有两箱零件，第一箱内装有 10 件，其中有 2 件次品；第二箱内装有 20 件，其中有 3 件次品，现从两箱中随意挑选一箱，然后从该箱中随机取 1 个零件，则取出的零件是次品的概率为\_\_\_\_\_.

5. 甲、乙两位同学高考后打算去旅游，准备分别从北京、成都、长沙、武汉和拉萨 5 个城市中随机选择其中一个,记事件 $A$ : 甲和乙选择的城市不同,事件 $B$ : 甲和乙恰好一人选择北京,则条件概率 $P\left( {B \mid  A}\right)  =$

6. 现从 3 个男生 2 个女生共 5 人中任意选出 3 人参加某校高三年级的百日誓师大会，若选出的 3 人中， 在有 1 人是女生的条件下，另 2 人是男生的概率为\_\_\_\_\_。

7.需要从某医院某科室的 5 名男医生(含一名主任医师)、4 名女医生(含一名主任医师)中分别选派 3 名男医生和2 名女医生组成医疗队，则在男主任医师被选派的条件下，两名主任医师都被选派的概率为\_\_\_\_\_.

8. 甲箱中有 5 个红球、 2 个白球和 3 个黑球，乙箱中有 4 个红球、 3 个白球和 3 个黑球(球除颜色外，大小质地均相同). 先从甲箱中随机取出一球放入乙箱,分别以 ${A}_{1}\text{ 、 }{A}_{2}$ 和 ${A}_{3}$ 表示由甲箱中取出的球是红球,白球和黑球的事件；再从乙箱中随机取出一球，以 $B$ 表示由乙箱中取出的球是红球的事件，下列说法正确的序号是\_\_\_\_\_.

 \ding{172} 事件 ${A}_{1}\text{ 、 }{A}_{2}$ 相互独立;  \ding{173}  $P\left( {A}_{3}\right)  = \frac{1}{5}$ ； \ding{174}  $P\left( B\right)  = \frac{9}{22}$ ； \ding{175}  $P\left( {B \mid  {A}_{2}}\right)  = \frac{4}{11}$ ； \circled{5}  $P\left( {{A}_{1} \mid  B}\right)  = \frac{5}{9}$ .

9.某种电子玩具按下按钮后, 会出现红球或绿球. 已知按钮第一次按下后, 出现红球与绿球的概率都是 $\frac{1}{2}$ ，从按钮第二次按下起，若前一次出现红球，则下一次出现红球、绿球的概率分别为 $\frac{1}{3}$ 和 $\frac{2}{3}$ ，若前一次出现绿球,则下一次出现红球、绿球的概率分别为 $\frac{3}{5}$ 和 $\frac{2}{5}$ ,记第 $n\left( {n \in  {\mathbf{N}}^{ * }}\right)$ 次按下按钮后出现红球的概率为 ${P}_{n}$ , 则 $\left\{  {P}_{n}\right\}$ 的通项公式为 ${P}_{n} =$ \_\_\_\_\_.

\section*{二、选择题}

10.某地空气质量监测资料表明,一天的空气质量为优良的概率是 0.75 , 连续两天为优良的概率是 0.6 , 已知某天的空气质量为优良, 则随后一天的空气质量为优良的概率是 ( )

A.0.8. B.0.75. C.0.6. D.0.45.

11. 现有 12 道四选一的单选题,学生甲对其中的 9 道题有思路,3 道题完全没有思路,有思路的题做对的概率为 0.9 , 没有思路的题只好任意猜一个答案, 猜对答案的概率为 0.25 . 学生甲从这 12 道题中随机选择 1 题, 他做对的概率是 ( )

A.0.225. B.0.4125. C.0.7375. D.0.8325.

12. 已知 5 只动物中有 1 只患有某种疾病，需要通过血液化验来确定患病的动物，血液化验结果呈阳性的为患病动物，下面是两种化验方案: ( )

方案甲:将各动物的血液逐个化验，直到查出患病动物为止；

方案乙:先取 3 只动物的血液进行混合，然后检查，若呈阳性，对这 3 只动物的血液再逐个化验，直到查出患病动物; 若不成阳性, 则检查剩下的两只动物中 1 只动物的血液.

A. 若利用方案甲，化验次数为 4 次的概率为 0.2 .

B. 若利用方案甲，平均检查次数为 2.8 .

C. 若利用方案乙，最多需要检查次数为 4 次.

D. 若利用方案乙，化验次数为 2 次的概率为 0.4 .

\section*{三、解答题}

13. 已知 $P\left( \bar{A}\right)  = \frac{2}{3},P\left( {\bar{B} \mid  A}\right)  = \frac{3}{4},P\left( {\bar{B} \mid  \bar{A}}\right)  = \frac{1}{3}$ ,计算:

(1) $P\left( {B \mid  A}\right)$ ；

(2) $P\left( B\right)$ .

14. 研究表明: 人体内某部位的半径约 $5{{mm}}$ 的结节约有 0.2\% 的可能性会在 1 年内发生病变. 某医院引进一台检测设备，可以通过对血液检测，估计患者体内半径约为 $5{{mm}}$ 的结节是否会在 1 年内发生病变，若检测结果为阳性, 则提示该结节会在 1 年内发生病变, 若检测结果为阴性, 则提示该结节不会在 1 年内发生病变.这种检测的准确率为 90\%，即一个会在 1 年内发生病变的患者有 90\%的可能性被检出阳性，一个不会在 1 年内发生病变的患者有 90\% 的可能性被检出阴性. 患者甲被检查出体内长了一个半径约为 $5{{mm}}$ 的结节，他做了该项血液检测.

(1)求患者甲检查结果为阳性的概率；

(2)若患者甲的检查结果为阳性，求他的这个结节在 1 年内发生病变的概率(结果保留 4 位小数).

15. 某机构对某品牌机电产品进行了质量调查, 下面是消费者关于质量投诉的数据:

\begin{center}
\begin{tabular}{|c|c|c|c|c|}
\hline
\phantom{X} & 擦伤(\%) & 凹痕(\%) & 外观(\%) & 合计(\%) \\
\cline{1-5}
保质期内 & 18 & 13 & 32 & 63 \\
\cline{1-5}
保质期后 & 12 & 22 & 3 & 37 \\
\cline{1-5}
合计 & 30 & 35 & 35 & 100 \\
\cline{1-5}
\hline
\end{tabular}
\end{center}

(1)已知投诉发生在保质期后，投诉的原因是产品外观的概率是多少？(结果保留两位有效数字)

(2)若事件 $A$ :投诉的原因是产品外观，事件 $B$ :投诉发生在保质期内，则 $A$ 和 $B$ 是独立事件吗？

16. 某保险的基本保费为 $a$ (单位:元)，继续购买该保险的投保人成为续保人，续保人本年度的保费与其上年度出险次数的关联如下:

\begin{center}
\begin{tabular}{|c|c|c|c|c|c|c|}
\hline
上年度出险次数 & 0 & 1 & 2 & 3 & 4 & ≥5 \\
\cline{1-7}
保费 & ${0.85a}$ & $a$ & ${1.25a}$ & 1.5a & 1.75a & ${2a}$ \\
\cline{1-7}
\hline
\end{tabular}
\end{center}

设该险种一续保人一年内出险次数与相应概率如下:

\begin{center}
\begin{tabular}{|c|c|c|c|c|c|c|}
\hline
一年内出险次数 & 0 & 1 & 2 & 3 & 4 & ≥5 \\
\cline{1-7}
概率 & 0.30 & 0.15 & 0.20 & 0.20 & 0.10 & 0.05 \\
\cline{1-7}
\hline
\end{tabular}
\end{center}

(1)求一续保人本年度的保费高于基本保费的概率；

(2)若一续保人本年度的保费高于基本保费，求其保费比基本保费高出 60\% 的概率；

(3)求续保人本年度的平均保费与基本保费的比值.

\section*{§ 12.3 随机变量的分布与特征}

\section*{典型例题精讲}

【例 1】从集合 $\{ 1,2,3,4,5\}$ 的所有非空子集中,等可能地取出一个.

(1)记性质 $r$ :集合中的所有元素之和为 10，求所取出的非空子集满足性质 $r$ 的概率；

(2)记所取出的非空子集的元素个数为 $\xi$ ，求 $\xi$ 的分布和数学期望 $E\left\lbrack  \xi \right\rbrack$ .

【例 2】甲、乙两个学校进行体育比赛，比赛共设三个项目，每个项目胜方得 10 分，负方得 0 分，没有平局. 三个项目比赛结束后, 总得分高的学校获得冠军. 已知甲学校在三个项目中获胜的概率分别为 0.5、0.4、 0.8 , 各项目的比赛结果相互独立.

(1)求甲学校获得冠军的概率；

(2)用 $X$ 表示乙学校的总得分，求 $X$ 的分布与期望和方差.

【例 3】某商场周年庆, 准备提供一笔资金, 对消费满一定金额的顾客以参与活动的方式进行奖励, 顾客从一个装有大小相同的 2 个红球和 4 个黄球的袋中按指定规则取出 2 个球, 根据取到的红球数确定奖励金额,具体金额设置如下表:

\begin{center}
\begin{tabular}{|c|c|c|c|}
\hline
取到的红球数 & 0 & 1 & 2 \\
\cline{1-4}
奖励(单位:元) & 5 & 10 & 50 \\
\cline{1-4}
\hline
\end{tabular}
\end{center}

现有两种取球规则的方案:

方案一:一次性随机取出 2 个球；

方案二:依次有放回取出 2 个球.

(1)比较两种方案下，一次抽奖获得 50 元奖金概率的大小；

(2)作为公司负责人，你会选择哪种方案？请说明理由.

【例 4】随着“五一”黄金周的到来，各大旅游景点热闹非凡，为了解 $A$ 、 $B$ 两个旅游景点游客的满意度， 某研究性学习小组采用随机抽样的方法,获得关于 $A$ 旅游景点的问卷 100 份,关于 $B$ 旅游景点的问卷 80 份. 问卷中, 对景点的满意度等级为: 非常满意、满意、一般、差评, 对应分数分别为: 4 分、3 分、2 分、1 分, 数据统计如下:

\begin{center}
\begin{tabular}{|c|c|c|c|c|}
\hline
\phantom{X} & 非常满意 (4 分) & 满意(3 分) & 一般(2 分) & 差评 (1 分) \\
\cline{1-5}
$A$ 景点 & 50 & 30 & 5 & 15 \\
\cline{1-5}
$B$ 景点 & 35 & 30 & 7 & 8 \\
\cline{1-5}
\hline
\end{tabular}
\end{center}

假设用频率估计概率,且游客对 $A\text{ 、 }B$ 两个旅游景点的满意度评价相互独立.

(1)从所有(人数足够多)在 $A$ 旅游景点的游客中随机抽取 2 人，从所有(人数足够多)在 $B$ 旅游景点的游客中随机抽取 2 人，估计这 4 人中恰有 2 人给出“非常满意”的概率；

(2)根据上述数据，你若旅游，你会选择 $A$ 、 $B$ 哪个旅游景点？说明理由.

【例 5】甲、乙两人投篮，每次由其中一人投篮，规则如下:若命中则此人继续投篮，若未命中则换对方投篮. 无论之前投篮情况如何, 甲每次投篮的命中率均为 0.6 , 乙每次投篮的命中率均为 0.8 . 由抽签确定第 1 次投篮的人选，第 1 次投篮的人是甲、乙的概率各为 0.5 .

(1)求第 2 次投篮的人是乙的概率；

(2)求第 $i$ 次投篮的人是甲的概率；

(3)已知:若随机变量 ${X}_{i}$ 服从两点分布，且 $P\left( {{X}_{i} = 1}\right)  = 1 - P\left( {{X}_{i} = 0}\right)  = {q}_{i},i = 1,2,\cdots ,n$ ,则 $E\left( {\mathop{\sum }\limits_{{i = 1}}^{n}{X}_{i}}\right)  = \mathop{\sum }\limits_{{i = 1}}^{n}{q}_{i}$ . 记前 $n$ 次 (即从第 1 次到第 $n$ 次投篮) 中甲投篮的次数为 $Y$ ,求 $E\left( Y\right)$ .

\section*{四基达标精练}

\section*{一、填空题}

1. 设 $X$ 是一个离散型随机变量,其分布为: $\left( \begin{matrix}  - 1 & 0 & 1 \\  {0.5} & 1 - {2q} & {q}^{2} \end{matrix}\right)$ ,则 $q =$ \_\_\_\_\_.

2. 设随机变量 $X$ 的可能取值为 $1\text{ 、 }2\text{ 、 }\cdots \text{ 、 }n$ ,并且取 $1\text{ 、 }2\text{ 、 }\cdots \text{ 、 }n$ 是等可能的. 若 $P\left( {X < 3}\right)  = {0.2}$ ,则 $n =$ \_\_\_\_\_.

3. 从一批含有 13 只正品,2 只次品的产品中,不放回地抽取 3 次,每次抽取 1 只,设抽得次品数为 $X$ ,则 $E\left\lbrack  {{5X} + 1}\right\rbrack   =$

4. 已知离散型随机变量 $\xi$ 的分布为 $\left( \begin{matrix} 1 & 3 & 5 \\  {0.5} & m & {0.2} \end{matrix}\right)$ ，则均值 $E\left\lbrack  \xi \right\rbrack   =$ \_\_\_\_\_.

5. 已知随机变量 $X$ 的分布为 $\left( \begin{matrix} 0 & 1 \\  \frac{1}{2} & a \end{matrix}\right)$ ，随机变量 $Y$ 的分布为 $\left( \begin{matrix} 1 & 2 \\  b & \frac{1}{3} \end{matrix}\right)$ ，且 $X$ 与 $Y$ 相互独立，则 $D\lbrack X +  {2Y}\rbrack  =$

6. 已知 $\eta$ 的分布为 $\left( \begin{matrix}  - 1 & 0 & 1 \\  \frac{1}{2} & \frac{1}{3} & \frac{1}{6} \end{matrix}\right)$ ,设 $\xi  = {3\eta } - 2$ ,则 $D\left\lbrack  \xi \right\rbrack$ 的值为\_\_\_\_\_.

7. 现有 7 张卡片,分别写上数字 $1\text{ 、 }2\text{ 、 }3\text{ 、 }4\text{ 、 }5\text{ 、 }6$ . 从这 7 张卡片中随机抽取 3 张,记所抽取卡片上数字的最小值为 $\xi$ ,则 $P\left( {\xi  = 2}\right)  =$ \_\_\_\_\_， $E\left\lbrack  \xi \right\rbrack   =$ \_\_\_\_\_.

8. 设随机变量 $\xi$ 的概率满足: $P\left( {X = k}\right)  = {a}_{k}\left( {k = 1,2,\cdots ,{10}}\right)$ ,且数列 $\left\{  {a}_{n}\right\}$ 满足 $P\left( {\xi  \leq  k}\right)  = k{a}_{k}(k = 1$ , $2,\cdots ,{10})$ ,则 $E\left\lbrack  \xi \right\rbrack   =$ \_\_\_\_\_.

9. 设随机变量 $\xi$ 的分布为: $\left( \begin{matrix} 1 & 2 & 3 & \cdots & {2020} & {2021} \\  {a}_{1} & {a}_{2} & {a}_{3} & \cdots & {a}_{2020} & {a}_{2021} \end{matrix}\right)$ ,则

 \ding{172}  若 $\left\{  {a}_{n}\right\}$ 为等差数列，则 ${a}_{2} + {a}_{2020} = \frac{2}{2021}$ ；

 \ding{173} 数列 $\left\{  {a}_{n}\right\}$ 的通项公式可能为 ${a}_{n} = \frac{2022}{{2021n}\left( {n + 1}\right) }$ ;

 \ding{174}  若数列 $\left\{  {a}_{n}\right\}$ 满足 ${a}_{n} = \frac{1}{{2}^{n}}\left( {n = 1,2,\cdots ,{2020}}\right)$ ，则 ${a}_{2021} = \frac{1}{{2}^{2021}}$ ；

 \ding{175} 若数列 $\left\{  {a}_{n}\right\}$ 满足 $P\left( {\xi  \leq  k}\right)  = {k}^{2}{a}_{k}\left( {k = 1,2,\cdots ,{2021}}\right)$ ,则 ${a}_{1} = \frac{1011}{2021}$ .

上列说法正确的是\_\_\_\_\_.

\section*{二、选择题}

10. 某射手射击所得环数 $\xi$ 的分布: $\left( \begin{matrix} 7 & 8 & 9 & {10} \\  x & {0.1} & {0.3} & y \end{matrix}\right)$ ,已知 $\xi$ 的数学期望 $E\left\lbrack  \xi \right\rbrack   = {8.9}$ ,则 $y$ 的值为 ( )

A.0.8. B.0.6. C.0.4. D.0.2.

11. 设随机变量 $X$ 的分布为 $P\left( {X = i}\right)  = \frac{i}{10}\left( {i = 1,2,3,4}\right)$ ,若 $P\left( {1 \leq  X < a}\right)  = \frac{3}{5}$ ,则实数 $a$ 的值可以是 ( )

A.1. B.2. C.3. D.4.

12. 已知 $0 < p < \frac{1}{2}$ ,随机变量 $\xi \text{ 、 }\eta$ 相互独立,随机变量 $\xi$ 的分布为 $\left( \begin{matrix}  - 1 & 1 \\  \frac{2}{3} & \frac{1}{3} \end{matrix}\right) ,\eta$ 的分布为 $\left( \begin{matrix}  - 1 & 1 \\  1 - p & p \end{matrix}\right)$ ,则当 $p$ 在 $\left( {0,\frac{1}{2}}\right)$ 内增大时 ( )

A. $E\left\lbrack  {\xi  + \eta }\right\rbrack$ 减小, $D\left\lbrack  {\xi  + \eta }\right\rbrack$ 增大. B. $E\left\lbrack  {\xi  + \eta }\right\rbrack$ 减小, $D\left\lbrack  {\xi  + \eta }\right\rbrack$ 减小.

C. $E\left\lbrack  {\xi  + \eta }\right\rbrack$ 增大, $D\left\lbrack  {\xi  + \eta }\right\rbrack$ 增大. D. $E\left\lbrack  {\xi  + \eta }\right\rbrack$ 增大, $D\left\lbrack  {\xi  + \eta }\right\rbrack$ 减小.

\section*{三、解答题}

13. 某小组共 10 人，利用假期参加义工活动，已知参加义工活动次数为 1、2、3 的人数分别为 3、3、4 . 现从这 10 人中随机选出 2 人作为该组代表参加座谈会.

(1)设 $A$ 为事件“选出的 2 人参加义工活动次数之和为 4”，求事件 $A$ 发生的概率；

(2)设随机变量 $X$ 为选出的 2 人参加义工活动次数差的绝对值，求 $X$ 的分布和数学期望.

14. 甲、乙两名运动员进行五局三胜制的乒乓球比赛, 先赢得 3 局的运动员获胜, 并结束比赛. 设各局比赛的结果相互独立,每局比赛甲赢的概率为 $\frac{2}{3}$ ,乙赢的概率为 $\frac{1}{3}$ .

(1)求甲获胜的概率；

(2)设 $X$ 为结束比赛所需要的局数,求随机变量 $X$ 的分布及数学期望.

15. 某厂接受了一项加工业务，加工出来的产品(单位:件)按标准分为 $A$ 、 $B$ 、 $C$ 、 $D$ 四个等级. 加工业务约定:对于 $A$ 级品、 $B$ 级品、 $C$ 级品，厂家每件分别收取加工费 90 元、50 元、20 元；对于 $D$ 级品，厂家每件要赔偿原料损失费 50 元. 该厂有甲、乙两个分厂可承接加工业务. 甲分厂加工成本费为 25 元/件, 乙分厂加工成本费为 20 元/件. 厂家为决定由哪个分厂承接加工业务，在两个分厂各试加工了 100 件这种产品，并统计了这些产品的等级, 整理如下:

甲分厂产品等级的频数分布表

\begin{center}
\begin{tabular}{|c|c|c|c|c|}
\hline
等级 & A & B & $C$ & $D$ \\
\cline{1-5}
频数 & 40 & 20 & 20 & 20 \\
\cline{1-5}
\hline
\end{tabular}
\end{center}

乙分厂产品等级的频数分布表

\begin{center}
\begin{tabular}{|c|c|c|c|c|}
\hline
等级 & $A$ & $B$ & $C$ & $D$ \\
\cline{1-5}
频数 & 28 & 17 & 34 & 21 \\
\cline{1-5}
\hline
\end{tabular}
\end{center}

(1)分别估计甲、乙两分厂加工出来的一件产品为 A 级品的概率;

(2)分别求甲、乙两分厂加工出来的产品的平均利润，以平均利润为依据，厂家应选哪个分厂承接加工业务?

16. 在校运动会上，只有甲、乙、丙三名同学参加铅球比赛，比赛成绩达到 ${9.50}{m}$ 以上(含 ${9.50}{m}$ )的同学将获得优秀奖. 为预测获得优秀奖的人数及冠军得主，收集了甲、乙、丙以往的比赛成绩，并整理得到如下数据(单位:m)

甲:9.80,9.70,9.55,9.54,9.48,9.42,9.40,9.35,9.30,9.25;

乙:9.78,9.56,9.51,9.36,9.32,9.23;

丙:9.85,9.65,9.20,9.16.

假设用频率估计概率, 且甲、乙、丙的比赛成绩相互独立.

(1)估计甲在校运动会铅球比赛中获得优秀奖的概率;

(2)设 $X$ 是甲、乙、丙在校运动会铅球比赛中获得优秀奖的总人数，估计 $X$ 的数学期望 $E\left\lbrack  X\right\rbrack$ 和方差 $D\left\lbrack  X\right\rbrack$ ；

(3)在校运动会铅球比赛中，甲、乙、丙谁获得冠军的概率估计值最大？(结论不要求证明).

\section*{§ 12.4 常用分布}

\section*{典型例题精讲}

【例 1】上海市某高中组织知识竞赛，分小组进行问题竞答. 甲、乙两个小组分别从 6 个问题中随机抽取 3 个问题进行回答, 答对题目多者为胜. 已知这 6 个问题中, 甲组能正确回答其中 4 个问题, 而乙组能正确回答每个问题的概率均为 $\frac{2}{3}$ . 甲、乙两个小组的选题以及对每题的回答都是相互独立、互不影响的.

(1)求甲组学生至少答对 2 个问题的概率；

(2)若从甲、乙两个组中选拔一组代表学校参加全市决赛，请分析说明选择哪个小组更好？

【例 2】已知随机变量 $X$ 服从二项分布 $B\left( {{1001},\frac{1}{2}}\right)$ .  \ding{172}  $P\left( {X = k}\right)  = {{C}}_{1001}^{k}{\left( \frac{1}{2}\right) }^{1001}$ ;  \ding{173}  $P\left( {X \leq  {301}}\right)  =  P\left( {X \geq  {701}}\right)$ ;  \ding{174}  $P\left( {X > E\left\lbrack  X\right\rbrack  }\right)  = \frac{1}{2}$ ;  \ding{175}  $P\left( {X = k}\right)$ 最大时 $k = {500}$ 或 501 .

指出上述命题中的真命题, 并说明理由.


【例 3】某校在校庆期间举办羽毛球比赛, 某班派出甲、乙两名单打主力, 为了提高两位主力的能力, 体育老师安排了为期一周的对抗训练，比赛规则如下:甲、乙两人每轮分别与体育老师打 2 局，当两人获胜局数不少于 3 局时,则认为这轮训练过关; 否则不过关.若甲、乙两人每局获胜的概率分别为 ${p}_{1}\text{ 、 }{p}_{2}\left( {0 \leq  {p}_{1}\text{ 、 }{p}_{2} \leq  1}\right)$ ,且满足 ${p}_{1} + {p}_{2} = \frac{3}{2}$ ,每局之间相互独立. 记甲、乙在 $n$ 轮训练中训练过关的轮数为 $X$ ,若 $E\left\lbrack  X\right\rbrack   = {24}$ ,则从期望的角度来看, 甲、乙两人至少需要训练几轮?

【例 4】经过不懈的奋力拼搏，农村建设取得了巨大进步，农民收入也逐年增加，为了更好地制定 2023 年“关于加快提升农民年收入、力争早日奔小康的工作计划”, 该地统计了 2022 年 50 位农民的年收入并制成如下频率分布直方图.

(1)根据频率分布直方图，估计这 50 位农民的年平均收入 $\overline{x}$ (单位:千元)(同一组数据用该组数据区间的中点值表示)；

(2)由频率分布直方图，可以认为该地区农民收入 $X$ 服从正态分布 $N\left( {\mu ,{\sigma }^{2}}\right)$ ，其中 $\mu$ 近似为年平均收入 $\bar{x},{\sigma }^{2}$ 近似为样本方差 ${s}^{2}$ ,经计算得 ${s}^{2} = {6.92}$ . 利用该正态分布,求:

 \ding{172}  在 2022 年脱贫攻坚工作中，若使该地区约有占总农民人数的 84.13\% 的农民的年收入高于最低年收人标准，则最低年收入大约为多少千元?

 \ding{173} 为了调研政策落实情况，有关部门随机走访了 1000 位农民. 若每位农民的年收入相互独立，问:这 1000 位农民中的年收入不少于 12.14 千元的人数最有可能是多少?

参考数据:标准正态分布函数 $\Phi \left( 1\right)  = {0.8413},\Phi \left( 2\right)  = {0.9772},\Phi \left( 3\right)  = {0.9987}$ .

\begin{center}
\includegraphics[width=0.4\textwidth]{images/bo_d7tk7oc91nqc73875fu0_267_881_1504_548_309_0.jpg}
\end{center}
 

【例 5】最新研发的某产品每次试验结果为成功或不成功,且试验成功的概率为 $p\left( {0 < p < 1}\right)$ . 现对该产品进行独立重复试验,若试验成功,试验结束; 若试验不成功,则继续试验,且最多试验 10 次. 记 $X$ 为试验结束时所进行的试验次数.

(1) 求 $P\left( {X = i}\right) \left( {i = 1,2,\cdots ,9,{10}}\right)$ ；

(2)证明: $E\left\lbrack  X\right\rbrack   < \frac{1}{p}$ .

\section*{四基达标精练}

\section*{一、填空题}

1. 已知 $X \sim  B\left( {n,p}\right)$ ，且 ${3E}\left\lbrack  X\right\rbrack   = {10D}\left\lbrack  X\right\rbrack$ ，则 $p =$ \_\_\_\_\_.

2. 已知随机变量 $X$ 服从正态分布 $N\left( {\mu ,{\sigma }^{2}}\right)$ ，若 $P\left( {2 < X < 6}\right)  = {0.6}$ ， $P\left( {X \geq  6}\right)  = {0.2}$ ，则 $\mu  =$ \_\_\_\_\_.

3. 上海环球金融中心的观光厅分布在第 94、第 97 种第 100 层, 若该电梯在底层载有 5 位游客，且每位游客在这三层的每一层下电梯是等可能的，用 $X$ 表示这 5 位乘客中在第 100 层下电梯的人数，则 $E\left\lbrack  X\right\rbrack   =$ \_\_\_\_\_.

4. 有 10 件产品,其中 3 件是次品,从中任取两件,若 $X$ 表示取得次品的件数,则 $P\left( {X < 2}\right)  =$ \_\_\_\_\_.

5. 某次招聘考试共有 50 个人参加, 假设每个人获得通过的概率都为 0.4 , 且各人通过与否相互独立. 设这 50 人中获得通过的人数为 $\xi$ ，则 $D\left\lbrack  {{3\xi } + {2022}}\right\rbrack   =$ \_\_\_\_\_.

6. 某地区为了解高中毕业年级男同学身体发育情况, 从全地区高三年级男同学中随机抽取了 10000 同学为样本,分别测量样本中每名同学的体重 $X$ (单位: ${{kg}}$ ),已知 $X \sim  N\left( {{60},{\sigma }^{2}}\right) ,P\left( {{58} < X < {62}}\right)  = {0.7}$ ,则样本中体重不低于 ${62}{{kg}}$ 的人数约为\_\_\_\_\_.

7.某袋中装有大小相同质地均匀的黑球和白球共 5 个.从袋中随机取出 3 个球,已知恰好全为黑球的概率为 $\frac{1}{10}$ ，若记取出 3 个球中黑球的个数为 $X$ ，则 $D\left\lbrack  X\right\rbrack   =$ \_\_\_\_\_.

8. 在一次抗洪抢险中,准备用射击的方法引爆从桥上游漂流而下的一巨大汽油罐. 已知只有 5 发子弹备用,且首次命中只能使汽油流出,再次命中才能引爆成功,两次命中不一定连续,每次射击命中率都是 $\frac{2}{3}$ ,每次命中与否互相独立. 则在油罐被引爆条件下，第 4 次击中油罐引爆成功的概率为\_\_\_\_\_.

9. 一批产品共 100 件,其中有 3 件不合格品,从中随机 $n\left( {n \in  \mathbf{N}}\right)$ 件,若用 $X$ 表示所抽取的 $n$ 件产品中不合格品的件数,则使 $X = 1$ 的概率取得最大值时, $n =$ \_\_\_\_\_.

\section*{二、选择题}

10. 生物实验室有 5 只兔子, 其中只有 3 只测量过某项指标. 若从这 5 只兔子中随机取出 3 只, 则恰有 2 只测量过该指标的概率为 ( )

A. $\frac{2}{3}$ . B. $\frac{3}{5}$ . C. $\frac{2}{5}$ . D. $\frac{1}{5}$ .

\begin{center}
\includegraphics[width=0.4\textwidth]{images/bo_d7tk7oc91nqc73875fu0_269_940_391_496_292_0.jpg}
\end{center}
 

(第 11 题图)

11. 已知三个正态分布密度函数 ${\varphi }_{i}\left( x\right)  = \frac{1}{\sqrt{2\pi }\sigma }{{e}}^{-\frac{{\left( x - {\mu }_{i}\right) }^{2}}{2{\sigma }^{2}}}(x \in  \mathbf{R}$ , $i = 1\text{ 、 }2\text{ 、 }3)$ 的图像如图所示,则下列结论正确的是 ( )

A. ${\mu }_{1} < {\mu }_{2} = {\mu }_{3},{\sigma }_{1} = {\sigma }_{2} > {\sigma }_{3}$ .

B. ${\mu }_{1} > {\mu }_{2} = {\mu }_{3},{\sigma }_{1} = {\sigma }_{2} < {\sigma }_{3}$ .

C. ${\mu }_{1} = {\mu }_{2} < {\mu }_{3},{\sigma }_{1} < {\sigma }_{2} = {\sigma }_{3}$ .

D. ${\mu }_{1} < {\mu }_{2} = {\mu }_{3},{\sigma }_{1} = {\sigma }_{2} < {\sigma }_{3}$ .

12. 某群体中的每位成员使用移动支付的概率都为 $p$ ,各成员

的支付方式相互独立. 设 $X$ 为该群体的 10 位成员中使用移动支付的人数,已知 $D\left\lbrack  X\right\rbrack   = {2.4},P\left( {X = 4}\right)  <  P\left( {X = 6}\right)$ ,则 $p$ 的值是 ( )

A.0.7. B.0.6. C.0.4. D.0.3.

\section*{三、解答题}

13. 根据以往统计资料, 某地车主购买甲种保险的概率为 0.5 , 购买乙种保险但不购买甲种保险的概率为 0.3 . 假设车主购买甲保险与乙保险相互独立且各车主购买保险相互独立.

(1)求该地 1 位车主至少购买甲、乙两种保险中的 1 种的概率；

(2) $X$ 表示该地的 100 位车主中，甲、乙两种保险都不购买的车主数. 求 $X$ 的期望.

14. 某居民小区有两个相互独立的安全防范系统 (简称系统) $A$ 和 $B$ ,系统 $A$ 和系统 $B$ 在任意时刻发生故障的概率分别为 $\frac{1}{10}$ 和 $p$ .

(1)若在任意时刻至少有一个系统不发生故障的概率为 $\frac{49}{50}$ ，求 $p$ 的值；

(2)求系统 $A$ 在 3 次相互独立的检测中不发生故障的次数大于发生故障的次数的概率.

15. 对某种产品的抽检规则如下:从一批 10 件产品中随机抽取 2 件，逐一检查，如果未发现次品，则该批产品抽检通过, 现有一批 10 件产品.

(1)若其中有 1 件次品，求该批产品通过抽检的概率；

(2)若该批产品通过抽检的概率为 $\frac{1}{3}$ ，其中次品有几件？

16. 某市教育部门规定，高中学生三年在校期间必须参加不少于 80 小时的社区服务，该市教育部门随机抽取了全市 200 位高中学生参加社区服务时间的数据,按时间段 $\lbrack {75},{80}),\lbrack {80},{85}),\lbrack {85},{90}),\lbrack {90},{95}),\lbrack {95}$ , 100] (单位:小时)进行统计，其频率分布直方图如图所示.

(1)求抽取的 200 位学生中，参加社区服务时间不少于 90 小时的学生人数，并估计从全市高中学生中任意选取一人，其参加社区服务时间不少于 90 小时的概率；

(2)从全市高中学生中任意选取 3 位学生，记 $X$ 为 3 位学生中参加社区服务时间不少于 90 小时的人数,试求随机变量 $X$ 的分布与期望.

\begin{center}
\includegraphics[width=0.4\textwidth]{images/bo_d7tk7oc91nqc73875fu0_270_903_1245_570_377_0.jpg}
\end{center}
 

(第 16 题图)

\section*{§ 12.5 统计初步}

\section*{典型例题精讲}

【例 1】在一次高三年级统一考试中,数学试卷里有一道满分 10 分的选做题,学生可以从 $A\text{ 、 }B$ 两道题目中任选一题作答. 某校有 800 名高三学生参加了本次考试，为了了解该校学生解答该选做题的得分情况， 计划从 800 名考生的选做题成绩中随机抽取一个容量为 10 的样本，为此将 800 名考生选做题的成绩按照随机顺序依次编号为 000-799.

(1)若采用随机数法抽样，并按照以下随机数表，从第 2 行的第 18 列的数字 9 为起点，从左向右依次读取数据，每次读取三位随机数，一行读数用完之后接下一行左端，写出样本编号；

05269370602235851513920351597759567806835291057074

07971088230998429964617162991506512916935805770951

51268785855487664754733208111244959263162956242948

26996165535837788070421050674232175585749444671694

14655268758759362241267863065513082701501529393943

(2)求出(1)中抽取的样本编号对应的数的极差与中位数.

【例 2】在一次男子 10 米气手枪射击比赛中, 甲、乙运动员的成绩如图所示, 射击队该从哪一位选手中参加比赛呢?

\begin{center}
\includegraphics[width=0.3\textwidth]{images/bo_d7tk7oc91nqc73875fu0_271_1078_1627_390_214_0.jpg}
\end{center}
 

(例 2 图)

【例 3】甲、乙两支田径队体检结果为:甲队的体重的平均数为 ${60}{{kg}}$ ，方差为 200，乙队体重的平均数为 70kg，方差为 300 ，又已知甲、乙两队的队员人数之比为 1 :4，那么甲、乙两队全部队员的平均体重和方差分别是什么？

【例 4】下表是 13-17 岁未成年人的身高的主要百分位数(单位:cm). 小明今年 16 岁，他的身高为 176cm，他所在城市男性同龄人约有 6.4 万人. 可以估计出小明的身高至少高于他所在城市多少男性同龄人.

13-17 岁未成年人的身高的主要百分位数

\begin{center}
\begin{tabular}{|c|c|c|c|c|c|c|c|c|c|c|}
\hline
\multicolumn{2}{|c|}{\phantom{X}} & ${P1}$ & ${P5}$ & P10 & P25 & P50 & P75 & P90 & P95 & P99 \\
\cline{1-11}
\multirow{2}{*}{13-15 岁} & 男 & 141 & 147 & 151 & 157 & 164 & 169 & 174 & 177 & 182 \\
\cline{2-11}
 & 女 & 143 & 147 & 150 & 153 & 157 & 161 & 165 & 167 & 171 \\
\cline{1-11}
\multirow{2}{*}{16-17 岁} & 男 & 155 & 160 & 163 & 167 & 171 & 175 & 179 & 181 & 186 \\
\cline{2-11}
 & 女 & 147 & 150 & 152 & 155 & 159 & 163 & 166 & 169 & 172 \\
\cline{1-11}
\hline
\end{tabular}
\end{center}

数据来源:《中国未成年人人体尺寸》(标准号:GB/T 26158-2010).

【例 5】某研究小组经过研究发现某种疾病的患病者与未患病者的某项医学指标有明显差异, 经过大量调查,得到如下的患病者和未患病者该指标的频率分布直方图:

\begin{center}
\includegraphics[width=0.7\textwidth]{images/bo_d7tk7oc91nqc73875fu0_273_319_246_934_458_0.jpg}
\end{center}
 

利用该指标制定一个检测标准,需要确定临界值 $c$ ,将该指标大于 $c$ 的人判定为阳性,小于或等于 $c$ 的人判定为阴性，此检测标准的漏诊率是将患病者判定为阴性的概率，记为 $p\left( c\right)$ ；误诊率是将未患病者判定为阳性的概率,记为 $q\left( c\right)$ . 假设数据在组内均匀分布,以事件发生的频率作为相应事件发生的概率.

(1)当漏诊率 $p\left( c\right)  = {0.5}\%$ 时，求临界值 $c$ 和误诊率 $q\left( c\right)$ ；

(2)设函数 $f\left( c\right)  = p\left( c\right)  + q\left( c\right)$ . 当 $c \in  \left\lbrack  {{95},{105}}\right\rbrack$ ,求 $f\left( c\right)$ 的解析式，并求 $f\left( c\right)$ 在区间 $\left\lbrack  {{95},{105}}\right\rbrack$ 的最小值.

\section*{四基达标精练}

\section*{一、填空题}

1. 甲、乙两名射击运动员参加某大型运动会的预选赛，他们分别射击了 5 次，成绩如下表 (单位:环):

\begin{center}
\begin{tabular}{|c|c|c|c|c|c|}
\hline
甲 & 10 & 8 & 9 & 9 & 9 \\
\cline{1-6}
乙 & 10 & 10 & 7 & 9 & 9 \\
\cline{1-6}
\hline
\end{tabular}
\end{center}

如果甲、乙两人中只有 1 人入选，则人选的最佳人选应是\_\_\_\_\_.

2. 某工厂在 12 月共生产了 3600 双皮靴，在出厂前要检查这批产品的质量，决定采用分层抽样的方法进行抽取,若从一、二、三车间抽取的产品数分别为 $a\text{ 、 }b\text{ 、 }c$ ,且 ${2b} = a + c$ ,则第二车间生产的产品数为\_\_\_\_\_.

3. 设一组样本数据 ${x}_{1},{x}_{2},\cdots ,{x}_{n}$ 的方差为 0.01,则数据 ${10}{x}_{1} + 1,{10}{x}_{2} + 1,\cdots ,{10}{x}_{n} + 1$ 的方差为\_\_\_\_\_.

4. 已知有四个数 $1\text{ 、 }2\text{ 、 }a\text{ 、 }b$ ,这四个数的中位数是 3,平均数是 4,则 ${ab} =$ \_\_\_\_\_.

5. 为了解某校高三年级男生的体重，从该校高三年级男生中抽取 17 名，测得他们的体重数据如下 (按从小到大的顺序排列,单位: ${{kg}}$ ):

5656575859596163646566686970737483

据此估计该校高三年级男生体重的第 75 百分位数为\_\_\_\_\_ ${{kg}}$ .

6. 已知某地一年四个季度的 GDP(亿元)逐季度增长，且中位数与平均数相同，其中第一季度 GDP 为 232(亿元)，第四季度 GDP 为 241(亿元)，则该地一年的 GDP 为\_\_\_\_\_亿元.

7. 有一组样本数据 ${x}_{1},{x}_{2},\cdots ,{x}_{6}$ ，其中 ${x}_{1}$ 是最小值， ${x}_{6}$ 是最大值，则下列说法正确的是\_\_\_\_\_.

 \ding{172}  ${x}_{2},{x}_{3},{x}_{4},{x}_{5}$ 的平均数等于 ${x}_{1},{x}_{2},\cdots ,{x}_{6}$ 的平均数;

 \ding{173}  ${x}_{2},{x}_{3},{x}_{4},{x}_{5}$ 的中位数等于 ${x}_{1},{x}_{2},\cdots ,{x}_{6}$ 的中位数;

 \ding{174}  ${x}_{2},{x}_{3},{x}_{4},{x}_{5}$ 的标准差不小于 ${x}_{1},{x}_{2},\cdots ,{x}_{6}$ 的标准差;

 \ding{175}  ${x}_{2},{x}_{3},{x}_{4},{x}_{5}$ 的极差不大于 ${x}_{1},{x}_{2},\cdots ,{x}_{6}$ 的极差.

8. 为了解某大学射击社团的射击水平, 分析组用分层抽样的方法抽取了 6 名老学员和 2 名新学员的某次射击成绩进行分析,经测算,6 名老学员的射击成绩样本均值为 8 (单位:环),方差为 $\frac{5}{3}$ (单位:环); 2 名新学员的射击成绩分别为3 环和 5 环，则抽取的这 8 名学员的射击成绩的方差为\_\_\_\_\_环。

9. 已知一组样本数据 ${x}_{1},{x}_{2},\cdots ,{x}_{n}\left( {{x}_{1} < {x}_{2} < \cdots  < {x}_{n}}\right)$ ,现有一组新的数据 $\frac{{x}_{1} + {x}_{2}}{2},\frac{{x}_{2} + {x}_{3}}{2},\cdots$ , $\frac{{x}_{n - 1} + {x}_{n}}{2},\frac{{x}_{n} + {x}_{1}}{2}$ ,则与原样本数据相比,对于新的数据有以下四个判断:  \ding{172} 平均数不变;  \ding{173} 中位数不变;  \ding{174}  极差变小;  \ding{175} 方差变小, 其中所有正确判断的序号是\_\_\_\_\_.

\section*{二、选择题}

10. 下列统计量中,不能度量样本 ${x}_{1},{x}_{2},\cdots ,{x}_{n}$ 的离散程度的有 ( )

A. 样本 ${x}_{1},{x}_{2},\cdots ,{x}_{n}$ 的标准差. B. 样本 ${x}_{1},{x}_{2},\cdots ,{x}_{n}$ 的方差.

C. 样本 ${x}_{1},{x}_{2},\cdots ,{x}_{n}$ 的极差. D. 样本 ${x}_{1},{x}_{2},\cdots ,{x}_{n}$ 的平均数.

11.FRM(Financial Risk Manager) ——金融风险管理师考试，是全球金融风险管理领域的一种资格认证考试. 某研究机构用随机数表法抽取了 2017 年参加 FRM 考试的某市 50 名考生的成绩进行分析, 先将 50 名考生按 01,02,03, $\cdots ,{50}$ 进行编号,然后从随机数表第 8 行第 11 列的数开始向右读,则选出的第 12 个个体是

(注:下面为随机数表的第 8 行和第 9 行)

16 95 55 67 19 9810507175 1286735807 4439523879

第 9 行:33211234297864560782 9966027954

A.12. B.21. C. 29. D.34.

12. 甲、乙两位射击爱好者,各射击 10 次,甲的环数从小到大排列为4,5,5,6,6,7,7,8,8,9，乙的环数从小到大排列为2,5,6,6,7,7,7,8.9,10. 则不正确的是 ( )

A. 甲的环数的 70\% 分位数是 7 .

B. 甲的平均环数比乙的平均环数小.

C. 这 20 个数据的平均值为 6.6 .

D. 甲的方差为 2.25 ，乙的方差为 4.41 ，这 20 个数据的方差为 3.34 .

\section*{三、解答题}

13.某车间 20 名工人年龄数据如下表:

\begin{center}
\begin{tabular}{|c|c|c|c|c|c|c|c|}
\hline
年龄(岁) & 19 & 28 & 29 & 30 & 31 & 32 & 40 \\
\cline{1-8}
工人数 & 1 & 3 & 3 & 5 & 4 & 3 & 1 \\
\cline{1-8}
\hline
\end{tabular}
\end{center}

(1)以十位数为茎，个位数为叶，作出这 20 名工人年龄的茎叶图；

(2)求这 20 名工人年龄的平均数、中位数、众数、极差和方差.

14.从某企业生产的产品中抽取 100 件, 测量这些产品的一项质量指标值, 由测量结果得如下频数分布表:

\begin{center}
\begin{tabular}{|c|c|c|c|c|c|}
\hline
质量指标值分组 & [75,85) & [85,95) & [95,105) & [105,115) & [115,125) \\
\cline{1-6}
频数 & 6 & 26 & 38 & 22 & 8 \\
\cline{1-6}
\hline
\end{tabular}
\end{center}

(1)在表格中作出这些数据的频率分布直方图；

\begin{center}
\includegraphics[width=0.4\textwidth]{images/bo_d7tk7oc91nqc73875fu0_275_519_1019_538_590_0.jpg}
\end{center}
 

(2)估计这种产品质量指标的平均数及方差(同一组中的数据用该组区间的中点值作代表)；

(3)根据以上抽样调查数据，能否认为该企业生产的这种产品符合“质量指标值不低于 95 的产品至少要占全部产品 80\%” 的规定?

15.某厂为比较甲乙两种工艺对橡胶产品伸缩率的处理效应, 进行 10 次配对试验, 每次配对试验选用材质相同的两个橡胶产品, 随机地选其中一个用甲工艺处理, 另一个用乙工艺处理, 测量处理后的橡胶产品的伸缩率,甲、乙两种工艺处理后的橡胶产品的伸缩率分别记为 ${x}_{i}\text{ 、 }{y}_{i}\left( {i = 1,2,\cdots {10}}\right)$ . 试验结果如下:

\begin{center}
\begin{tabular}{|c|c|c|c|c|c|c|c|c|c|c|}
\hline
试验序号 $i$ & 1 & 2 & 3 & 4 & 5 & 6 & 7 & 8 & 9 & 10 \\
\cline{1-11}
伸缩率 ${x}_{i}$ & 545 & 533 & 551 & 522 & 575 & 544 & 541 & 568 & 596 & 548 \\
\cline{1-11}
伸缩率 ${y}_{i}$ & 536 & 527 & 543 & 530 & 560 & 533 & 522 & 550 & 576 & 536 \\
\cline{1-11}
\hline
\end{tabular}
\end{center}

记 ${z}_{i} = {x}_{i} - {y}_{i}\left( {i = 1,2,\cdots ,{10}}\right)$ ,记 ${z}_{1},{z}_{2},\cdots ,{z}_{10}$ 的样本平均数为 $\bar{z}$ ,样本方差为 ${s}^{2}$ .

(1)求 $\frac{}{z},{s}^{2}$ ；

(2)判断甲工艺处理后的橡胶产品的伸缩率较乙工艺处理后的橡胶产品的伸缩率是否有显著提高. (如果 $\bar{z} \geq  2\sqrt{\frac{{s}^{2}}{10}}$ ,则认为甲工艺处理后的橡胶产品的伸缩率较乙工艺处理后的橡胶产品的伸缩率有显著提高, 否则不认为有显著提高)

16. 某厂商调查甲、乙两种不同型号电视机在 8 个卖场的销售量(单位:台)，并根据这 8 个卖场的销售情况，得到如图所示的茎叶图. 为了鼓励卖场，在同型号电视机的销售中，该厂商将销售量高于数据平均数的卖场命名为该型号电视机的“星级卖场”.

(1)当 $a = 1,b = 1$ 时，记甲型号电视机的“星级卖场”数量为 $m$ ，乙型号电视机的“星级卖场”数量为 $n$ ，比较 $m$ 、 $n$ 的大小关系；

(2)在这 8 个卖场中，随机选取 2 个卖场，求这两个卖场都是甲型号电视机的“星级卖场”的概率；

(3)记乙型号电视机销售量的方差为 ${s}^{2}$ ，根据茎叶图推断， $a$ 与 $b$ 分别取何值时， ${s}^{2}$ 达到最小值(只需写出结论).

\begin{center}
\includegraphics[width=0.2\textwidth]{images/bo_d7tk7oc91nqc73875fu0_276_1189_1636_309_129_0.jpg}
\end{center}
 

(第 16 题图)

\section*{§ 12.6 成对数据的统计分析}

\section*{典型例题精讲}

【例 1】某科技公司为制订下一年的研发投入计划，需了解年研发资金投入量X(单位:亿元)对年销售额 $y$ (单位:亿元)的影响,对近 10 年研发资金投入量 ${x}_{i}$ 和销售额 ${y}_{i}\left( {i = 1,2,\cdots ,{10}}\right)$ 数据作了初步处理,得到如图的散点图.

\begin{center}
\includegraphics[width=0.4\textwidth]{images/bo_d7tk7oc91nqc73875fu0_277_981_574_486_458_0.jpg}
\end{center}
 

(例 1 图)

(1)根据散点图判断， $y = a + {bx}$ 与 $y = {{e}}^{c + {dx}}$ 哪一个适宜作为年销售额 $y$ 关于年研发资金投入量 $x$ 的回归方程类型? (给出判断即可, 不必说明理由)

(2) \ding{172} 根据(1)的选择及表中数据，建立 $y$ 关于 $x$ 的回归方程 (精确到 0.001 );

 \ding{173} 若下一年销售额 $y$ 需达到 90 亿元，预测下一年的研发资金投入量 $x$ 约为多少亿元? (精确到 1 )

附:对于一组数据 $\left( {{\mu }_{1},{v}_{1}}\right) ,\left( {{\mu }_{2},{v}_{2}}\right) ,\cdots ,\left( {{\mu }_{n},{v}_{n}}\right)$ ,其回归直线 $v = {a\mu } + b$ 的斜率和截距的最小二乘估计分别为:

$\widehat{a} = \frac{\mathop{\sum }\limits_{{i = 1}}^{n}\left( {{\mu }_{i} - \overline{\mu }}\right) \left( {{v}_{i} - \bar{v}}\right) }{\mathop{\sum }\limits_{{i = 1}}^{n}{\left( {\mu }_{i} - \overline{\mu }\right) }^{2}},\widehat{b} = \bar{v} - \widehat{a}\overline{\mu }$ . 参考数据: 其中 $w = \ln y$ .

\begin{center}
\begin{tabular}{|c|c|c|c|c|c|c|c|}
\hline
$\frac{ - }{x}$ & − & $\bar{w}$ & $\mathop{\sum }\limits_{{i = 1}}^{{10}}{\left( {x}_{i} - \bar{x}\right) }^{2}$ & $\mathop{\sum }\limits_{{i = 1}}^{{10}}{\left( {y}_{i} - \bar{y}\right) }^{2}$ & $\mathop{\sum }\limits_{{i = 1}}^{{10}}{\left( {w}_{i} - \bar{w}\right) }^{2}$ & $\mathop{\sum }\limits_{{i = 1}}^{{10}}\left( {{x}_{i} - \bar{x}}\right) \left( {{y}_{i} - \bar{y}}\right)$ & $\mathop{\sum }\limits_{{i = 1}}^{{10}}\left( {{x}_{i} - \bar{x}}\right) \left( {{w}_{i} - \bar{w}}\right)$ \\
\cline{1-8}
20 & 30 & 3.2 & 900 & 300 & 3.78 & 1600 & 58.21 \\
\cline{1-8}
\hline
\end{tabular}
\end{center}

【例 2】为了监控某种零件的一条生产线的生产过程,检验员每隔 ${30}{{min}}$ 从该生产线上随机抽取一个零件,并测量其尺寸 (单位:cm). 下面是检验员在一天内依次抽取的 16 个零件的尺寸:

\begin{center}
\begin{tabular}{|c|c|c|c|c|c|c|c|c|}
\hline
抽取次序 & 1 & 2 & 3 & 4 & 5 & 6 & 7 & 8 \\
\cline{1-9}
零件尺寸 & 9.95 & 10.12 & 9.96 & 9.96 & 10.01 & 9.92 & 9.98 & 10.04 \\
\cline{1-9}
抽取次序 & 9 & 10 & 11 & 12 & 13 & 14 & 15 & 16 \\
\cline{1-9}
零件尺寸 & 10.26 & 9.91 & 10.13 & 10.02 & 9.22 & 10.04 & 10.05 & 9.95 \\
\cline{1-9}
\hline
\end{tabular}
\end{center}

(1)求 $\left( {{x}_{i},i}\right) \left( {i = 1,2,\cdots ,{16}}\right)$ 的相关系数 $r$ ，并回答是否可以认为这一天生产的零件尺寸不随生产过程的进行而系统地变大或变小 (若 $\left| r\right|  < {0.25}$ ，则可以认为零件的尺寸不随生产过程的进行而系统地变大或变小).

(2)一天内抽检零件中，如果出现了尺寸在 $\left( {\bar{x} - {3s},\bar{x} + {3s}}\right)$ 之外的零件，就认为这条生产线在这一天的生产过程可能出现了异常情况，需对当天的生产过程进行检查.

(1)从这一天抽检的结果看，是否需对当天的生产过程进行检查？

(ii)在( $\bar{x} - {3s},\bar{x} + {3s}$ )之外的数据称为离群值，试剔除离群值，估计这条生产线当天生产的零件尺寸的均值与标准差. (精确到 0.01)

附:样本 $\left( {{x}_{i},{y}_{i}}\right) \left( {i = 1,2,\cdots ,n}\right)$ 的相关系数 $r = \frac{\mathop{\sum }\limits_{{i = 1}}^{n}\left( {{x}_{i} - \bar{x}}\right) \left( {{y}_{i} - \bar{y}}\right) }{\sqrt{\mathop{\sum }\limits_{{i = 1}}^{n}{\left( {x}_{i} - \bar{x}\right) }^{2}}\sqrt{\mathop{\sum }\limits_{{i = 1}}^{n}{\left( {y}_{i} - \bar{y}\right) }^{2}}}$ .

【例 3】“绿色出行, 低碳环保”已成为新的时尚, 近几年国家相继出台了一系列的环保政策, 为新能源汽车行业的发展开辟了广阔的前景. 某公司对 $A$ 充电桩进行生产投资，所获得的利润有如下统计数据:

\begin{center}
\begin{tabular}{|c|c|c|c|c|c|c|}
\hline
$A$ 充电桩投资金额 $x/$ 万元 & 3 & 4 & 6 & 7 & 9 & 10 \\
\cline{1-7}
所获利润 $y/$ 百万元 & 1.5 & 2 & 3 & 4.5 & 6 & 7 \\
\cline{1-7}
\hline
\end{tabular}
\end{center}

(1)已知可用一元线性回归模型拟合 $y$ 与 $x$ 的关系,求其回归方程；

(2)若规定所获利润 $y$ 与投资金额 $x$ 的比值不低于 $\frac{2}{3}$ ，则称对应的投入额为 “优秀投资额”. 记 2 分，所获利润 $y$ 与投资金额 $x$ 的比值低于 $\frac{2}{3}$ 且大于 $\frac{1}{2}$ ，则称对应的投入额为 “良好投资额”，记 1 分，所获利润 $y$ 与投资金额 $x$ 的比值不超过 $\frac{1}{2}$ ，则称对应的投入额为 “不合格投资额”，记 0 分，现从表中 6 个投资金额中任意选 2 个,用 $X$ 表示记分之和,求 $X$ 的分布及数学期望.

附:对于一组数据 $\left( {{x}_{1},{y}_{1}}\right) ,\left( {{x}_{2},{y}_{2}}\right) ,\cdots ,\left( {{x}_{n},{y}_{n}}\right)$ ,其回归直线方程 $\widehat{y} = \widehat{a}x + \widehat{b}$ 的斜率和截距的最小二乘法估计公式分别为: $\widehat{a} = \frac{\mathop{\sum }\limits_{{i = 1}}^{n}\left( {{x}_{i} - \bar{x}}\right) \left( {{y}_{i} - \bar{y}}\right) }{\mathop{\sum }\limits_{{i = 1}}^{n}{\left( {x}_{i} - \bar{x}\right) }^{2}} = \frac{\mathop{\sum }\limits_{{i = 1}}^{n}{x}_{i}{y}_{i} - n\bar{x} \cdot  \bar{y}}{\mathop{\sum }\limits_{{i = 1}}^{n}{x}_{i}^{2} - n{\bar{x}}^{2}},\widehat{b} = \bar{y} - \widehat{a}\bar{x}$ .

【例 4】一医疗团队为研究某地的一种地方性疾病与当地居民的卫生习惯 (卫生习惯分为良好和不够良好两类)的关系，在已患该疾病的病例中随机调查了 100 例(称为病例组)，同时在未患该疾病的人群中随机调查了 100 人(称为对照组)，得到如下数据:

\begin{center}
\begin{tabular}{|c|c|c|}
\hline
\phantom{X} & 不够良好 & 良好 \\
\cline{1-3}
病例组 & 40 & 60 \\
\cline{1-3}
对照组 & 10 & 90 \\
\cline{1-3}
\hline
\end{tabular}
\end{center}

(1)能否有 99\% 的把握认为患该疾病群体与未患该疾病群体的卫生习惯有差异？

(2)从该地的人群中任选一人，A 表示事件“选到的人卫生习惯不够良好”，B 表示事件“选到的人患有该疾病”, $\frac{P\left( {B \mid  A}\right) }{P\left( {\bar{B} \mid  A}\right) }$ 与 $\frac{P\left( {B \mid  \bar{A}}\right) }{P\left( {\bar{B} \mid  \bar{A}}\right) }$ 的比值是卫生习惯不够良好对患该疾病风险程度的一项度量指标,记该指标为 $R$ .

 \ding{172}  证明: $R = \frac{P\left( {A \mid  B}\right) }{P\left( {\overline{A} \mid  B}\right) }$ . $\frac{P\left( {\overline{A} \mid  \overline{B}}\right) }{P\left( {A \mid  \overline{B}}\right) }$ ；

 \ding{173}  利用该调查数据，给出 $P\left( {A \mid  B}\right)$ 、 $P\left( {A \mid  \bar{B}}\right)$ 的估计值，并利用 \ding{172} 的结果给出 $R$ 的估计值.

附: ${K}^{2} = \frac{n{\left( ad - bc\right) }^{2}}{\left( {a + b}\right) \left( {c + d}\right) \left( {a + c}\right) \left( {b + d}\right) }$ .

\begin{center}
\begin{tabular}{|c|c|c|c|}
\hline
$P\left( {{K}^{2} \geq  k}\right)$ & 0.050 & 0.010 & 0.001 \\
\cline{1-4}
$k$ & 3.841 & 6.635 & 10.828 \\
\cline{1-4}
\hline
\end{tabular}
\end{center}

【例 5】某地在种植棉花有着得天独厚的自然条件，土质呈碱性，夏季温差大，阳光充足，光合作用充分，生长时间长，这种环境下种植的棉花绒长、品质好、产量高，但每年五月，该地进入灾害天气高发期，灾害天数对当年棉花产量有着重要影响，根据过去五年的数据统计，得到相关数据如下表:

\begin{center}
\begin{tabular}{|c|c|c|c|c|c|}
\hline
灾害天气天数 $x$ (天) & 2 & 3 & 4 & 5 & 8 \\
\cline{1-6}
棉花产量 $y$ (吨/公顷) & 3.2 & 2.4 & 2 & 1.9 & 1.7 \\
\cline{1-6}
\hline
\end{tabular}
\end{center}

根据以上数据, 技术人员分别借助甲、乙两种不同的回归模型, 得到两个回归方程,

模型甲: ${\widehat{y}}_{1} = \frac{4}{x} + {1.1}$ ,模型乙: ${\widehat{y}}_{2} = \frac{6.4}{{x}^{2}} + {1.6}$ .

(1)分别利用这两个模型, 求受灾害天气为 7 天时, 当年棉花产量的预测值; (精确到 0.01

(2)你认为用哪个模型得到的预测值更可靠？并说明理由.

\section*{四基达标精练}

\section*{一、填空题}

\begin{center}
\includegraphics[width=0.2\textwidth]{images/bo_d7tk7oc91nqc73875fu0_280_1213_1568_257_259_0.jpg}
\end{center}
 

(第 1 题图)

1. 如图所示，有 $A$ 、 $B$ 、 $C$ 、 $D$ 、 $E$ 共 5 组数据，去掉\_\_\_\_\_组数据后，剩下的 4 组数据具有较强的线性相关关系. (请用 $A$ 、 $B$ 、 $C$ 、 $D$ 、 $E$ 作答)

2. 甲、乙、丙、丁四位同学各自对 $x$ 、 $y$ 两变量的线性相关性作试验,求得相关系数 $r$ 如表:

\begin{center}
\begin{tabular}{|c|c|c|c|c|}
\hline
相关系数 & 甲 & 乙 & 丙 & 丁 \\
\cline{1-5}
$r$ & -0.82 & 0.78 & 0.69 & 0.87 \\
\cline{1-5}
\hline
\end{tabular}
\end{center}

则\_\_\_\_\_同学的试验结果体现两变量有更强的线性相关性.

3. 在一组样本数据 $\left( {{x}_{1},{y}_{1}}\right) ,\left( {{x}_{2},{y}_{2}}\right) ,\cdots ,\left( {{x}_{n},{y}_{n}}\right) \left( {n \geq  2}\right.$ ,且 ${x}_{1}\text{ 、 }{x}_{2}\text{ 、 }\cdots \text{ 、 }{x}_{n}$ 不全相等) 的散点图中,若所有样本点 $\left( {{x}_{i},{y}_{i}}\right) \left( {i = 1,2,\cdots ,n}\right)$ 都在直线 $y =  - \frac{1}{2}x + 1$ 上，则这组数据的样本相关系数为\_\_\_\_\_.

4. 调查某种花萼长度和花瓣长度,所得数据如图所示,其中相关系数 $r = {0.8245}$ ,下列说法正确的是\_\_\_\_\_.

\begin{center}
\includegraphics[width=0.5\textwidth]{images/bo_d7tk7oc91nqc73875fu0_281_478_242_642_260_0.jpg}
\end{center}
 

(第 3 题图)

 \ding{172} 花瓣长度和花萼长度没有相关性； \ding{173} 花瓣长度和花萼长度呈现负相关； \ding{174} 花瓣长度和花萼长度呈现正相关;  \ding{175} 若从样本中抽取一部分, 则这部分的相关系数一定是 0.8245 .

5. 为了规定工时定额，需要确定加工零件所花费的时间，为此进行了 5 次试验. 根据收集到的数据(如表),由最小二乘法求得回归方程为 $y = {0.67x} + {54.9}$ .

\begin{center}
\begin{tabular}{|c|c|c|c|c|c|}
\hline
零件数 $x$ 个 & 10 & 20 & 30 & 40 & 50 \\
\cline{1-6}
加工时间 $y/{{min}}$ & 62 & \includegraphics[width=0.2\textwidth]{images/bo_d7tk7oc91nqc73875fu0_281_760_784_118_76_0.jpg} & 75 & 81 & 89 \\
\cline{1-6}
\hline
\end{tabular}
\end{center}

现发现表中有一个数据模糊看不清，则该数据为\_\_\_\_\_.

6. 调查了某地若干户家庭的年收入 $x$ (单位:万元)和年饮食支出 $y$ (单位:万元). 调查显示年收入 $x$ 与年饮食支出 $y$ 具有线性相关关系,并由调查数据得到 $y$ 对 $x$ 的回归方程: $y = {0.254x} + {0.321}$ . 由回归方程可知, 家庭年收入每增加 1 万元，年饮食支出平均增加\_\_\_\_\_万元.

7. 某池塘中水生植物的覆盖水塘面积 $x$ (单位: ${{{dm}}}^{2}$ ) 与水生植物的株数 $y$ (单位:株) 之间的相关关系,收集了 4 组数据,用模型 $y = c{{e}}^{kx}\left( {c > 0}\right)$ 去拟合 $x$ 与 $y$ 的关系,设 $z = \ln y,x$ 与 $z$ 的数据如表格所示:

\begin{center}
\begin{tabular}{|c|c|c|c|c|}
\hline
$x$ & 3 & 4 & 6 & 7 \\
\cline{1-5}
$z$ & 2.5 & 3 & 4 & 5.9 \\
\cline{1-5}
\hline
\end{tabular}
\end{center}

得到 $x$ 与 $z$ 的线性回归方程 $\widehat{z} = {0.7x} + a$ ,则 $c =$ \_\_\_\_\_. (用指数形式表示)

8. 某数学老师身高 ${176}{{cm}}$ ，他爷爷、父亲和儿子的身高分别是 ${173}{{cm}}$ 、 ${170}{{cm}}$ 和 ${182}{{cm}}$ . 因儿子的身高与父亲的身高有关，该老师用线性回归分析的方法预测他孙子的身高为\_\_\_\_\_cm.

9. 足球是一项大众喜爱的运动，为了解喜爱足球是否与性别有关，随机抽取了若干人进行调查，抽取女性人数是男性的 2 倍，男性喜爱足球的人数占男性人数的 $\frac{5}{6}$ ，女性喜爱足球的人数占女性人数的 $\frac{1}{3}$ ，若本次调查得出“在犯错误的概率不超过 0.005 的前提下认为喜爱足球与性别有关”的结论，则被调查的男性至少有\_\_\_\_\_人。

\begin{center}
\begin{tabular}{|c|c|c|c|c|c|}
\hline
$a$ & 0.10 & 0.05 & 0.01 & 0.005 & 0.001 \\
\cline{1-6}
${x}_{a}$ & 2.706 & 3.841 & 5.635 & 7.879 & 10.828 \\
\cline{1-6}
\hline
\end{tabular}
\end{center}

附: ${\chi }^{2} = \frac{n{\left( ad - bc\right) }^{2}}{\left( {a + b}\right) \left( {c + d}\right) \left( {a + c}\right) \left( {b + d}\right) }$

\begin{center}
\includegraphics[width=0.4\textwidth]{images/bo_d7tk7oc91nqc73875fu0_281_938_1816_518_213_0.jpg}
\end{center}
 

(第 10 题图)

\section*{二、选择题}

10. 根据散点图, 下列说法正确的是 ( )

A. 身高越大，体重越大.

B. 身高越大，体重越小.

C. 身高和体重成正相关.

D. 身高和体重成负相关.

11. 在某次独立性检验中, 得到如下列联表:

\begin{center}
\begin{tabular}{|c|c|c|c|}
\hline
\phantom{X} & $A$ & $\bar{A}$ & 总计 \\
\cline{1-4}
$B$ & 200 & 800 & 1000 \\
\cline{1-4}
$\bar{B}$ & 180 & $a$ & ${180} + a$ \\
\cline{1-4}
总计 & 380 & ${800} + a$ & ${1180} + a$ \\
\cline{1-4}
\hline
\end{tabular}
\end{center}

最后发现,两个分类变量没有关联,则 $a$ 的值可能是 ( )

A. 960. B. 720. C. 360. D. 180.

12.某学校对高二学生是否喜欢阅读进行随机调查, 调查的数据如表所示:

\begin{center}
\begin{tabular}{|c|c|c|c|}
\hline
\phantom{X} & 喜欢阅读 & 不喜欢阅读 & 总计 \\
\cline{1-4}
男学生 & 30 & 20 & 50 \\
\cline{1-4}
女学生 & 40 & 10 & 50 \\
\cline{1-4}
总计 & 70 & 30 & 100 \\
\cline{1-4}
\hline
\end{tabular}
\end{center}

\begin{center}
\begin{tabular}{|c|c|c|c|c|c|c|c|}
\hline
$P\left( {{\chi }^{2} \geq  k}\right)$ & 0.25 & 0.15 & 0.10 & 0.05 & 0.025 & 0.010 & 0.001 \\
\cline{1-8}
$k$ & 1.323 & 2.072 & 2.706 & 3.841 & 5.024 & 6.635 & 10.828 \\
\cline{1-8}
\hline
\end{tabular}
\end{center}

根据表中的数据, 下列对该校高二学生的说法正确的是 ( )

A. 没有 95\%以上的把握认为“性别与是否喜欢阅读有关”.

B. 有 99\%以上的把握认为“性别与是否喜欢阅读有关”.

C. 在犯错误的概率不超过 0.025 的前提下认为 “性别与是否喜欢阅读有关”.

D. 在犯错误的概率不超过 0.05 的前提下认为“性别与是否喜欢阅读有关".

\section*{三、解答题}

13. 某地经过多年的环境治理, 已将荒山改造成了绿水青山. 为估计一林区某种树木的总材积量, 随机选取了10棵这种树木，测量每棵树的根部横截面积(单位: ${{m}}^{2}$ )和树积量(单位: ${{m}}^{3}$ )，得到数据，并计算得 $\mathop{\sum }\limits_{{i = 1}}^{{10}}{x}_{i} = {0.6},\mathop{\sum }\limits_{{i = 1}}^{{10}}{y}_{i} = {3.9},\mathop{\sum }\limits_{{i = 1}}^{{10}}{x}_{i}^{2} = {0.038},\mathop{\sum }\limits_{{i = 1}}^{{10}}{y}_{i}^{2} = {1.6158},\mathop{\sum }\limits_{{i = 1}}^{{10}}{x}_{i}{y}_{i} = {0.2474}.$

(1)求该林区这种树木的根部横截面积与材积量的样本相关系数(精确到 0.01)；

(2)现测量了该林区所有这种树木的根部横截面积, 并得到所有这种树木的根部横截面积总和为 ${186}{{m}}^{2}$ . 利用以上数据和回归方程 $y = \widehat{a}x + \widehat{b}$ ,给出该林区这种树木的总材积量的估计值.

附:相关系数 $r = \frac{\mathop{\sum }\limits_{{i = 1}}^{n}{x}_{i}{y}_{i} - n\bar{x} \cdot  \bar{y}}{\sqrt{\left( {\mathop{\sum }\limits_{{i = 1}}^{n}{x}_{i}^{2} - n{\bar{x}}^{2}}\right) \left( {\mathop{\sum }\limits_{{i = 1}}^{n}{y}_{i}^{2} - n{\bar{y}}^{2}}\right) }\text{ ,回归系数 }\widehat{a} = \frac{\mathop{\sum }\limits_{{i = 1}}^{n}{x}_{i}{y}_{i} - n\bar{x} \cdot  \bar{y}}{\mathop{\sum }\limits_{{i = 1}}^{n}{x}_{i}^{2} - n{\bar{x}}^{2}},\widehat{b} = \bar{y} - \widehat{a}\bar{x}}$ .

14. 一项试验旨在研究臭氧效应, 试验方案如下: 选 40 只小白鼠, 随机地将其中 20 只分配到试验组, 另外 20 只分配到对照组, 试验组的小白鼠饲养在高浓度臭氧环境, 对照组的小白鼠饲养在正常环境, 一段时间后统计每只小白鼠体重的增加量(单位:g).试验结果如下:

对照组的小白鼠体重的增加量从小到大排序为:

15.2 18.8 20.2 21.3 22.5 23.2 25.8 26.5 27.5 30.1

32.634.334.835.635.635.836.237.340.543.2

试验组的小白鼠体重的增加量从小到大排序为:

7.8 9.2 11.4 12.4 13.2 15.5 16.5 18.0 18.8 19.2

19.8 20.2 21.6 22.8 23.6 23.9 25.1 28.2 32.3 36.5

(1)计算试验组的样本平均数；

(2) \ding{172}  求 40 只小白鼠体重的增加量的中位数 $m$ ，再分别统计两样本中小于 $m$ 与不小于 $m$ 的数据的个数，完成如下列联表；

\begin{center}
\begin{tabular}{|c|c|c|}
\hline
\phantom{X} & $< m$ & $\geq  m$ \\
\cline{1-3}
对照组 & \phantom{X} & \phantom{X} \\
\cline{1-3}
试验组 & \phantom{X} & \phantom{X} \\
\cline{1-3}
\hline
\end{tabular}
\end{center}

 \ding{173} 根据(i)中的列联表，能否有 95\%的把握认为小白鼠在高浓度臭氧环境中与在正常环境中体重的增加量有差异?

附: ${K}^{2} = \frac{n{\left( ad - bc\right) }^{2}}{\left( {a + b}\right) \left( {c + d}\right) \left( {a + c}\right) \left( {b + d}\right) }$ .

\begin{center}
\begin{tabular}{|c|c|c|c|}
\hline
$P\left( {{K}^{2} \geq  k}\right)$ & 0.100 & 0.050 & 0.010 \\
\cline{1-4}
$k$ & 2.706 & 3.841 & 6.635 \\
\cline{1-4}
\hline
\end{tabular}
\end{center}

\begin{center}
\includegraphics[width=0.3\textwidth]{images/bo_d7tk7oc91nqc73875fu0_284_1095_142_390_315_0.jpg}
\end{center}
 

15.某同学用随机抽样的办法对所在学校老师的饮食习惯进行了一次调查, 饮食指数结果用茎叶图表示如图, 图中饮食指数低于 70 的人是饮食以蔬菜为主: 饮食指数高于 70 的人是饮食以肉类为主.

(1)完成下列 $2 \times  2$ 列联表:

\begin{center}
\begin{tabular}{|c|c|c|c|}
\hline
\phantom{X} & 主食蔬菜 & 主食肉类 & 总计 \\
\cline{1-4}
不超过 45 岁 & 4 & \_\_\_\_\_ & \_\_\_\_\_ \\
\cline{1-4}
45 岁以上 & \_\_\_\_\_ & \_\_\_\_\_ & \_\_\_\_\_ \\
\cline{1-4}
总计 & \_\_\_\_\_ & \_\_\_\_\_ & \_\_\_\_\_ \\
\cline{1-4}
\hline
\end{tabular}
\end{center}

能否有 99\% 的把握认为老师的饮食习惯与年龄有关？

(2)从所在学校任一名老师设“选到 45 岁以上老师为事件 $A$ ”，“饮食指数高于 70 的老师”为事件 $B$ ，用调查的结果估计 $P\left( {B \mid  A}\right)$ 及 $P\left( {B \mid  \bar{A}}\right)$ (用最简分数作答)；

(3)为了给食堂提供老师的饮食信息，根据(1)(2)的结论，能否有更好的抽样方法来估计老师的饮食习惯, 并说明理由.

附:

\begin{center}
\begin{tabular}{|c|c|c|c|}
\hline
$P\left( {{K}^{2} \geq  {k}_{0}}\right)$ & 0.050 & 0.010 & 0.001 \\
\cline{1-4}
${k}_{0}$ & 3.841 & 6.635 & 10.828 \\
\cline{1-4}
\hline
\end{tabular}
\end{center}

${K}^{2} = \frac{n{\left( ad - bc\right) }^{2}}{\left( {a + b}\right) \left( {c + d}\right) \left( {a + c}\right) \left( {b + d}\right) }$

16. 为了选拔培养有志于服务国家重大战略需求且综合素质优秀或基础学科拔尖的学生, 聚焦高端芯片与软件、智能科技、新材料、先进制造和国家安全等关键领域以及国家人才紧缺的人文社会科学领域, 某校在一次模拟考试后,从全体考生中随机抽取 52 名,获取他们本次考试的数学成绩 $\left( x\right)$ 和物理成绩 $\left( y\right)$ ，绘制成如图散点图:根据散点图可以看出 $y$ 与 $x$ 之间有线性相关关系，但图中有两个异常点 $A$ 、 $B$ . 经调查得知， $A$ 考生由于重感冒导致物理考试发挥失常， $B$ 考生因故未能参加物理考试，为了使分析结果更科学准确,剔除这两组数据后,对剩下的数据作处理,得到一些统计的值: $\mathop{\sum }\limits_{{i = 1}}^{{50}}{x}_{i} = {5800},\mathop{\sum }\limits_{{i = 1}}^{{50}}{y}_{i} = {3900}$ , $\mathop{\sum }\limits_{{i = 1}}^{{50}}{x}_{i}{y}_{i} = {462770},\mathop{\sum }\limits_{{i = 1}}^{{50}}{\left( {x}_{i} - \bar{x}\right) }^{2} = {28540},\mathop{\sum }\limits_{{i = 1}}^{{50}}{\left( {y}_{i} - \bar{y}\right) }^{2} = {18930}$ ,其中 ${x}_{i}\text{ 、 }{y}_{i}$ 分别表示这 50 名考生的数学成绩、 物理成绩, $i = 1,2,\cdots ,{50},y$ 与 $x$ 的相关系数 $r \approx  {0.45}$ .

\begin{center}
\includegraphics[width=0.8\textwidth]{images/bo_d7tk7oc91nqc73875fu0_285_335_550_939_284_0.jpg}
\end{center}
 

(1)若不剔除 $A$ 、 $B$ 两名考生的数据，用 52 组数据作回归分析，设此时 $y$ 与 $x$ 的相关系数为 ${r}_{0}$ . 试判断 ${r}_{0}$ 与 $r$ 的大小关系 (不必说明理由);

(2)求 $y$ 关于 $x$ 的线性回归方程(系数精确到 0.01)，并估计如果 $B$ 考生加了这次物理考试，物理成绩是多少? (精确到 1)

参考数据: $r = \frac{\mathop{\sum }\limits_{{i = 1}}^{n}\left( {{x}_{i} - \bar{x}}\right) \left( {{y}_{i} - \bar{y}}\right) }{\sqrt{\mathop{\sum }\limits_{{i = 1}}^{n}{\left( {x}_{i} - \bar{x}\right) }^{2}\mathop{\sum }\limits_{{i = 1}}^{n}{\left( {y}_{i} - \bar{y}\right) }^{2}}},\widehat{a} = \frac{\mathop{\sum }\limits_{{i = 1}}^{n}{x}_{i}{y}_{i} - n\bar{x} \cdot  \bar{y}}{\mathop{\sum }\limits_{{i = 1}}^{n}{x}_{i}^{2} - n{\bar{x}}^{2}},\widehat{b} = \bar{y} - \widehat{a}\bar{x}$ .

\section*{第 13 章 数学建模(必修 4 \& 选择性必修 3)}

\section*{知识梳理}

\section*{1. 数学建模的概念}

\begin{center}
\includegraphics[width=0.2\textwidth]{images/bo_d7tk7oc91nqc73875fu0_286_1161_480_320_402_0.jpg}
\end{center}
 

数学建模是对某个现实问题经过必要的简化、合理的假设得到一个数学问题, 称为数学模型, 再求解所得到的数学问题, 然后根据实际情况验证该数学答案是否合理；在合理性得不到保证时，还要进行反复迭代和修正模型.

数学建模是一个建立数学模型的过程, 而数学模型是数学建模的一个结果. 数学建模的全过程一般包括在实际情境中从数学的视角发现问题、提出问题、分析问题、建立模型、求解模型得到答案、回到实际情境验证答案、如果答案不合理则加以改进等步骤, 如此循环直到实际情境中的问题得到合理的回答. 这个过程统称为数学建模活动.

2. 数学建模的基本过程

数学建模的基本过程包括:

 \ding{172} 在实际情境中从数学的视角发现问题

 \ding{173} 提出问题

 \ding{174}  分析问题、建立模型

 \ding{175}  确定参数、计算求解

 \circled{5} 检验结果、改进模型

 \circled{6} 最终解决实际问题、形成报告

3. 数学建模的水平要求

\begin{center}
\begin{tabular}{|c|c|}
\hline
情境与问题 & \makecell{ 水平一:了解熟悉的数学模型的实际背景及其数学描述，了解数学模型中的参数、结论的实际含义水平二:能够在熟悉的情境中，发现问题并转化为数学问题，知道数学问题的价值与作用 \\ 水平三:能够在综合的情境中，运用数学思维进行分析，发现情境中的数学关系，提出数学问题 } \\
\cline{1-2}
知识与技能 & \makecell{ 水平一:知道数学建模的过程包括:提出问题、建立模型、求解模型、检验结果、完善模型.能够在熟悉的实际情境中, 模仿学过的数学建模过程解决问题 \\ 水平二、能够选择合适的数学模型表达所要解决的数学问题; 理解模型中参数的意义, 知道如何确定参数, 律立模型, 求解模型; 能够根据问题的实际意义检验结果, 完善模型, 解决问题 \\ 水平三:能够运用数学建模的一般方法和相关知识，创造性地建立数学模型，解决问题 } \\
\cline{1-2}
思维与表达 & \makecell{ 水平一:对于学过的数学模型，能够举例说明建模的意义，体会其蕴含的数学思想；感悟数学表达对数学建模的重要性 \\ 水平二:能够在关联的情境中，经历数学建模的过程，理解数学建模的意义；能够运用数学语言，表述数学律模过程中的问题以及解决问题的过程和结果，形成研究报告，展示研究成果 \\ 水平三、能够理解数学建模的意义和作用; 能够运用数学语言, 清晰、准确地表达数学建模的过程和结果 } \\
\cline{1-2}
交流与反思 & \makecell{ 水平一:在交流的过程中，能够借助或引用已有数学建模的结果说明问题 \\ 水平二:在交流的过程中，能够用模型的思想说明问题 \\ 水平三.在交流的过程中, 能够通过数学建模的结论和思想阐释科学规律和社会现象 } \\
\cline{1-2}
\hline
\end{tabular}
\end{center}

\section*{§ 13.1 数学建模的基本步骤}

\section*{典型例题精讲}

【例 1】某晚报曾刊登过一则生活趣事, 市民唐某乘坐出租车时, 在半途中骂骂咧咧要求司机临时停靠，打表计价结账，然后重新计价，继续前行，该市民解释说，根据经验，这样分开支付车费比一次性付费便宜一些, 他的这一说法有道理吗?

确实, 由于出租车运价上调, 有些人出行时会估计一下可能的价格, 再决定是否乘坐出租车.

据了解，该市出租车在 5 时到 23 时之间起租价为 14 元/3 km，超起租里程单价为 2.50 元/km，总里程超过 ${15}{{km}}$ (不含 ${15}{{km}}$ ) 部分按超起租里程单价加 50\%. 此外，相关部门还规定了低速等候费和其他时段的计价办法, 以及适合其他车型的计价办法.

你乘坐过出租车吗? 你会仿效那位市民唐某的做法吗? 为什么?

(1)根据上述情境你能提出什么数学问题？为了解决你的问题，请给出一些合理假设；

(2)你能否根据你的假设建立数学模型, 并回答你所提出的问题.

【例 2】雨天外出虽然有雨伞，却时常免不了淋湿衣袖、裤脚、背包等，小明想通过数学建模的方法研究如何撑伞可以让淋湿的面积尽量小. 为了简化问题小明做出下列假设:

假设 1 :在网上查阅了人均身高和肩宽的数据后，小明把人假设为身高、肩宽分别为 170cm、40cm 的矩形“纸片人”:

假设 2 : 受风的影响,雨滴下落轨迹视为与水平地面所成角为 ${60}^{ \circ  }$ 的直线;

假设 3: 伞柄 ${OT}$ 长为 ${60}{{cm}}$ ,可绕矩形 “纸片人”上点 $O$ 旋转;

假设 4: 伞面为被伞柄 ${OT}$ 垂直平分的线段 ${AB},{AB} = {120}{{cm}}$ .

以如图 1 方式撑伞，矩形“纸片人”将淋湿“裤脚”；以如图 2 方式撑，则矩形“纸片人”将淋湿“头和肩膀”。

\begin{center}
\includegraphics[width=0.6\textwidth]{images/bo_d7tk7oc91nqc73875fu0_288_409_511_791_301_0.jpg}
\end{center}
 

(例 2 图)

(1)如图 3 在矩形“纸片人”上身恰好不被淋湿时，求其“裤脚”被淋湿(阴影)部分的面积(结果精确到 0.1c ${\text{ m }}^{2}$ );

(2)请根据你的生活经验对小明建立的数学模型提出两条改进建议 (无需求解改进后的模型，如果建议超过两条仅对前两条评分).

\begin{center}
\includegraphics[width=0.2\textwidth]{images/bo_d7tk7oc91nqc73875fu0_289_1290_160_169_307_0.jpg}
\end{center}
 

(例 3 图)

【例 3】小明同学最近注意到一条新闻, 易拉罐 (如右图所示) 作为饮品的容器, 每年的用量可达数万亿个. 这让他想到一个用料最优化的问题, 即在易拉罐的体积一定的情况下，如何确定易拉罐的高和半径才能使得用料最省？他研究发现易拉罐的上盖、下底和侧壁的厚度是不同的, 进而结合数学建模知识进行了深入研究.

以下是小明的研究过程, 请你补全缺失的部分.

(1)模型假设: \ding{172} 易拉罐近似看成圆柱体； \ding{173} 上盖、下底、侧壁的厚度处处均匀； \ding{174} 上盖、 下底、侧壁所用金属相同； \ding{175} 易拉罐接口处的所用材料忽略不计.

(2)建立模型:记圆柱体积为 $V$ ，高为 $h$ ，底面半径为 $r$ ，上盖、下底和侧壁的厚度分别为 ${d}_{1}$ 、 ${d}_{2}$ 、 ${d}_{3}$ ，金属用料总量为 $C$ . 由几何知识得到如下数量关系:

$\left\{  \begin{array}{l} V = \pi {r}^{2}h,\text{  \ding{172}  } \\  C = \pi {r}^{2}{d}_{1} + \pi {r}^{2}{d}_{2} + {2\pi rh}{d}_{3}. \end{array}\right.$

由  \ding{172}  得 $h = \frac{V}{\pi {r}^{2}}$ ，代入  \ding{173}  得: $C = \pi \left( {{d}_{1} + {d}_{2}}\right) {r}^{2} + \frac{{2V}{d}_{3}}{r}$

$\because V\text{ 、 }{d}_{1}\text{ 、 }{d}_{2}\text{ 、 }{d}_{3}$ 都是常数,不妨设 $a = \pi \left( {{d}_{1} + {d}_{2}}\right) ,b = {2V}{d}_{3}$ ,

则用料总量的函数简化为 $C = a{r}^{2} + \frac{b}{r}$ .

请写出表格中代入整理这一步的目的是:\_\_\_\_\_【分】

(3)求解模型:

$\therefore$ 当 $r =$ \_\_\_\_\_(用 $a\text{ 、 }b$ 表示)时， $C$ 取得最小值，即在此种情况下用料最省.

(4)检验模型:小明上网查阅到目前 330 毫升可乐易拉罐的数据，得知 ${d}_{1} = {0.028}{{cm}}$ ， ${d}_{2} = {0.021}{{cm}}$ ， ${d}_{3} = {0.011}{{cm}}$ ，代入(3)的模型结果，经计算得 $r \approx  {2.868}{{cm}}$ .

经验算,确认计算无误,但是这与实际罐体半径 ${3.305}{{cm}}$ 差异较大.

实际上, 在经济利益驱动之下, 目前的罐体成本应该已经达最优.

(5)模型评价与改进:

模型计算结果与现实数据存在较大差异的原因可能为:\_\_\_\_\_，

相应改进措施为:\_\_\_\_\_

【例 4】某企业新研发了一种产品, 产品的成本由原料成本及非原料成本组成. 每件产品的非原料成本 $y$ (单位:元) 与生产该产品的数量 $x$ (单位:千件) 有关，经统计得到如下数据:

\begin{center}
\begin{tabular}{|c|c|c|c|c|c|c|c|c|}
\hline
$x$ & 1 & 2 & 3 & 4 & 5 & 6 & 7 & 8 \\
\cline{1-9}
$y$ & 112 & 61 & 44.5 & 35 & 30.5 & 28 & 25 & 24 \\
\cline{1-9}
\hline
\end{tabular}
\end{center}

根据以上数据,绘制了散点图.

\begin{center}
\includegraphics[width=0.3\textwidth]{images/bo_d7tk7oc91nqc73875fu0_290_1107_409_378_472_0.jpg}
\end{center}
 

(例 4 图)

观察散点图, 两个变量不具有线性相关关系, 现考虑用反比例函数模型 $y = a + \frac{b}{x}$ 和指数函数模型 $y = c{{e}}^{ax}$ 分别对两个变量的关系进行拟合. 已求得用指数函数模型拟合的经验回归方程为 $\widehat{y} = {96.54}{{e}}^{-{0.2x}},\ln y$ 与 $x$ 的相关系数 ${r}_{1} =  - {0.94}$ .

(1)用反比例函数模型求 $y$ 关于 $x$ 的经验回归方程；

(2)用相关系数对模型进行分析，问上述两个函数模型哪一个拟合效果更好(精确到 0.01)，并用其估计产量为 10 千件时每件产品的非原料成本；

(3)该企业采取订单生产模式(根据订单数量进行生产，即产品全部售出). 根据市场调研数据, 若该产品单价定为 100 元, 则签订 9 千件订单的概率为 0.8 ，签订 10 千件订单的概率为 0.2 ；若单价定为 90 元，则签订 10 千件订单的概率为 0.3 ，签订 11 千件订单的概率为 0.7 . 已知每件产品的原料成本为 10 元, 根据(2)的结果, 企业要想获得更高利润, 产品单价应选择 100 元还是 90 元, 请说明理由.

参考公式: 数据 $\left( {{u}_{1},{v}_{1}}\right) ,\left( {{u}_{2},{v}_{2}}\right) ,\cdots ,\left( {{u}_{n},{v}_{n}}\right)$ 的经验回归直线 $\widehat{v} = \widehat{\alpha }u + \widehat{\beta }$ 的斜率和截距的最小二乘估计为:

$\widehat{\alpha } = \frac{\mathop{\sum }\limits_{{i = 1}}^{n}{u}_{i}{v}_{i} - n\bar{u}\bar{v}}{\mathop{\sum }\limits_{{i = 1}}^{n}{u}_{i}^{2} - n{\bar{u}}^{2}},\widehat{\beta } = \bar{v} - \widehat{\alpha }\bar{u}$ ,相关系数 $r = \frac{\mathop{\sum }\limits_{{i = 1}}^{n}{u}_{i}{v}_{i} - n\bar{u}\bar{v}}{\sqrt{\left( {\mathop{\sum }\limits_{{i = 1}}^{n}{u}_{i}^{2} - n{\bar{u}}^{2}}\right) \left( {\mathop{\sum }\limits_{{i = 1}}^{n}{v}_{i}^{2} - n{\bar{v}}^{2}}\right) }}$ .

【例 5】一种微生物群体可以经过自身繁殖不断生存下来,设一个这种微生物为第 0 代,经过 $k(k = 0$ , $1,2,\cdots )$ 次繁殖后为第 $k$ 代,该微生物每代繁殖的个数相互独立,且具有相同的分布,设 $X$ 表示 1 个微生物个体繁殖下一代的个数, $P\left( {X = i}\right)  = {p}_{i}\left( {i = 0,1,2,3}\right)$ .

(1)已知 ${p}_{0} = {0.4},{p}_{1} = {0.3},{p}_{2} = {0.2},{p}_{3} = {0.1}$ ,求 $E\left\lbrack  X\right\rbrack$ ;

(2)已知该种微生物经过多代繁殖后临近灭绝的概率 $p$ 是方程 ${p}_{0} + {p}_{1}x + {p}_{2}{x}^{2} + {p}_{3}{x}^{3} = x$ 的一个最小正实根，求证:当 $E\left\lbrack  X\right\rbrack   \leq  1$ 时， $p = 1$ ，当 $E\left\lbrack  X\right\rbrack   > 1$ 时， $p < 1$ ；

(3)利用数学模型解释现实生活，根据你的理解说明(2)问结论的实际含义.

\section*{四基达标精练}

\section*{一、填空题}

1.一般的数学建模包含如下活动过程: \ding{172} 建立模型； \ding{173} 实际情境； \ding{174} 提出问题； \ding{175} 求解模型； \circled{5} 实际结果；  \circled{6} 检验结果，请写出正确的序号顺序\_\_\_\_\_.

2. 珠穆朗玛峰高达 8848.86m，但即使你拥有良好的视力，你也无法在上海看到它. 一个观察者距离珠穆朗玛峰多远，才能在地面上看到它呢？为了能够通过几何方法解决这个问题，需要利用简单的几何模型表示这个问题情境，在此过程中，有下列假设: \ding{172}  珠穆朗玛峰的形状为等腰梯形； \ding{173}  地球的形状是一个球体；  \ding{174} 太阳光线沿直线传播； \ding{175} 没有事物可以阻碍人们看到珠穆朗玛峰的视线，你认为最不重要的一个假设是\_\_\_\_\_.

3. 某兴趣小组设计制作一枚不均匀的骰子, 满足:

假设1:这枚骰子有六个面，点数分别为1、2、3、4、5、6；

假设 2 : 当投掷骰子的时候, 偶数点出现的概率比奇数点出现的概率大一倍;

假设 3: 不同偶数点出现的概率相同, 不同奇数点出现的概率也相同.

现将骰子投掷一次，则求点数不小于 4 的概率为\_\_\_\_\_.

4. "学如逆水行舟，不进则退；心似平原跑马，易放难收” (明。《增广贤文》) 是勉励人们专心学习的. 如果每天的“进步”率都是 $1\%$ ，那么一年后是 ${\left( 1 + 1\% \right) }^{365} = {1.01}^{365}$ ；如果每天的“退步”率都是 $1\%$ ，那么一年后是 ${\left( 1 - 1\% \right) }^{365} = {0.99}^{365}$ . 一年后“进步”的是“退步”的 $\frac{{1.01}^{365}}{{0.99}^{365}} = {\left( \frac{1.01}{0.99}\right) }^{365} \approx  {1481}$ 倍. 如果每月的“进步”率和“退步”率都是 20\%，那么大约经过\_\_\_\_\_月后“进步”的是 “退步”的一万倍. (精确到 1 月)

5.2023 年 1 月底, 人工智能研究公司 OpenAI 发布的名为 “ChatGTP”的人工智能聊天程序进入中国，迅速以其极高的智能化水平引起国内关注. 深度学习是人工智能的一种具有代表性的实现方法, 它是以神经网络为出发点的,在神经网络优化中,指数衰减的学习率模型为 $L = {L}_{0}{D}^{\frac{G}{G0}}$ ,其中 $L$ 表示每一轮优化时使用的学习率, ${L}_{0}$ 表示初始学习率, $D$ 表示衰减系数, $G$ 表示训练迭代轮数, ${G}_{0}$ 表示衰减速度. 已知某个指数衰减的学习率模型的初始学习率为 0.8 , 衰减速度为 12 , 且当训练迭代轮数为 12 时, 学习率衰减为 0.5 . 则学习率衰减到 0.2 以下(不含 0.2 )所需的训练迭代轮数至少为\_\_\_\_\_.

6. 某传染病在某地蔓延，当地医疗机构采取了各种针对性的治疗方法，取得了不错的成效，每日治愈的患者人数如表所示: 由表格可得 $y$ 关于 $x$ 的线性回归方程为 $\widehat{y} = \widehat{b}x + 1$ ,若第 6 日实际治愈人数为18人,则此回归模型第 6 日的离差(实际值减去预报值)为\_\_\_\_\_.

\begin{center}
\begin{tabular}{|c|c|c|c|c|c|}
\hline
第 $x$ 日 & 1 & 2 & 3 & 4 & 5 \\
\cline{1-6}
治愈人数 $y$ (人) & 2 & 9 & 10 & 13 & 16 \\
\cline{1-6}
\hline
\end{tabular}
\end{center}

7.小明和他的数学建模小队现有这样一个问题:提高过江大桥的车辆通行能力可改善整个城市的交通状况，那么，怎样才可以提高呢? 他们理想化地建立这样一个关系，在一般情况下，大桥上的车流速度 $v$ (单位: ${k}{m}/{h}$ )是车流密度 $x$ (单位:辆/ ${{km}}$ )的函数，当桥上的车流密度达到 200 辆/ ${{km}}$ 时，造成堵塞，此时车流速度为 0 ; 当车流密度不超过 20 辆/km 时,车流速度为 60 km/h. 研究表明,当 $x \in  \left\lbrack  {{20},{200}}\right\rbrack$ 时,车流速度 $v$ 是车流密度 $x$ 的一次函数. 问:当车流密度为\_\_\_\_\_辆/ ${{km}}$ 时，车流量(单位时间内通过桥上某观测点的车辆数，单位:辆/h)可以达到最大.

8. 某一时段内，从天空降落到地面上的雨水，未经蒸发、渗漏、流失而在水平面上积聚的深度，称为这个时段的降雨量(单位:mm).24h 降雨量的等级划分如下:

\begin{center}
\begin{tabular}{|c|c|}
\hline
等级 & 24 小时降雨量(精确到 0.1) \\
\cline{1-2}
... ... & ... ... \\
\cline{1-2}
小雨 & 0.1\textasciitilde 9.9 \\
\cline{1-2}
中雨 & 10.0\textasciitilde 24.9 \\
\cline{1-2}
大雨 & 25.0\textasciitilde 49.9 \\
\cline{1-2}
暴雨 & 50.0\textasciitilde 99.9 \\
\cline{1-2}
... ... & ... ... \\
\cline{1-2}
\hline
\end{tabular}
\end{center}

\begin{center}
\includegraphics[width=0.2\textwidth]{images/bo_d7tk7oc91nqc73875fu0_292_1007_933_303_285_0.jpg}
\end{center}
 

在综合实践活动中,某小组自制了一个底面直径为 ${200}{{mm}}$ 、高为 ${300}{{mm}}$ 的圆锥形雨量器. 若一次降雨过程中，该雨量器收集的 24h 的雨水高度是 150mm (如图所示)，则这 24h 降雨量的等级是\_\_\_\_\_. (填 “小雨”“中雨”“大雨”或“暴雨”)

9. 某公司通过统计分析发现，工人工作效率 $E$ 与工作年限 $r$ ( $r > 0$ )、劳累程度 $T\left( {0 < T < 1}\right)$ 、劳动动机 $b\left( {1 < b < 5}\right)$ 相关,并建立了数学模型 $E = {10} - {10T} \cdot  {b}^{-{0.14r}}$ . 已知甲、乙为该公司的员工，则下列说法正确的有\_\_\_\_\_.

 \ding{172} 甲与乙工作年限相同，且甲比乙工作效率高，劳动动机低，则甲比乙劳累程度强；

 \ding{173} 甲与乙劳动动机相同，且甲比乙工作效率高，工作年限短，则甲比乙劳累程度弱；

 \ding{174} 甲与乙劳累程度相同，且甲比乙工作年限长，劳动动机高，则甲比乙工作效率高；

 \ding{175} 甲与乙劳动动机相同，且甲比乙工作年限长，劳累程度弱，则甲比乙工作效率高。

\section*{二、选择题}

10. 工厂需定期购买原料并存放在仓库供生产之用，因此必须考虑解决什么才是合理的存贮量问题. 为了建立数学模型解决相关问题，需要分析问题情境，提出合理假设，以便简化实际问题情境，抓住问题核心， 如我们可以提出:“假设1:该工厂对于原料的需求量是恒定的”“假设2:为了保障生产，仓库内的原料不可以缺货”. 那么为了更好地建立模型, 你认为还需要下面哪些假设 ( )

 \ding{172} 该厂每天的产能是个定值，所有产品都能售出； \ding{173} 每件产品所需原料的每日存储费用是个常数；

 \ding{174}  每件产品所需购买原料的价格不变；  \ding{175} 工厂不能保证所生产的每一件产品都是正品.

A.  \ding{172}  \ding{173}  \ding{174}  \ding{175} . B.  \ding{172}  \ding{173} . C.  \ding{172}  \ding{173}  \ding{174} . D.  \ding{172}  \ding{175} .

11. 数学与音乐有着紧密的关联. 声音中也包含正弦函数，声音是由于物体的振动产生的能引起听觉的波，每一个音都是由纯音合成的. 纯音的数学模型是函数 $y = A\sin {\omega x}$ ，我们平时听到的音乐一般不是纯音，而是有多种波叠加而成的复合音. 已知刻画某复合音的函数为 $\sin x + \frac{1}{2}\sin {2x} + \frac{1}{3}\sin {3x}$ ,则其部分图像大致为 ( )

\begin{center}
\includegraphics[width=0.2\textwidth]{images/bo_d7tk7oc91nqc73875fu0_293_1114_379_298_202_0.jpg}
\end{center}
 

\begin{center}
\includegraphics[width=0.2\textwidth]{images/bo_d7tk7oc91nqc73875fu0_293_206_376_304_208_0.jpg}
\end{center}
 

y y

0

A. B. C. D.

(第 11 题图)

12. 为了估计一批产品的不合格品率 $p$ ,现从这批产品中随机抽取一个样本容量为 $n$ 的样本 ${\xi }_{1}\text{ 、 }{\xi }_{2}$ 、 ${\xi }_{3}\text{ 、 }\cdots \text{ 、 }{\xi }_{n}$ ,定义 ${\xi }_{i} = \left\{  \begin{array}{l} 1,\text{ 第 }i\text{ 次不合格,则 }P\left( {{\xi }_{i} = 1}\right)  = p,P\left( {{\xi }_{i} = 0}\right)  = 1 - p,\text{ 其中 }i = 1,2,\cdots ,n,\text{ 记 }L\left( p\right)  =  \end{array}\right.  P\left( {{\xi }_{1} = {x}_{1},{\xi }_{2} = {x}_{2},\cdots ,{\xi }_{n} = {x}_{n}}\right)$ (其中 ${x}_{i} = 0$ 或 $1,i = 1,2,\cdots ,n$ ),称 $L\left( p\right)$ 表示 $p$ 为参数的似然函数. 极大似然估计法是建立在极大似然原理基础上的一个统计方法, 极大似然原理的直观想法是: 一个随机试验如有若干个可能的结果 $A\text{ 、 }B\text{ 、 }C\text{ 、 }\cdots$ ,若在一次试验中,结果 $A$ 出现,则一般认为试验条件对 $A$ 出现有利,也即 $A$ 出现的概率很大. 极大似然估计是一种用给定观察数据来评估模型参数的统计方法, 即 “模型已定, 参数未知", 通过若干次试验, 观察其结果, 利用试验结果得到某个参数值能够使样本出现的概率为最大.根据以上原理, 下面说法不正确的是 ( )

A. 有外形完全相同的两个箱子, 甲箱有 99 个白球 1 个黑球, 乙箱有 1 个白球 99 个黑球. 今随机地抽取一箱, 再从取出的一箱中抽取一球, 结果取得白球, 那么该球一定是从甲箱子中抽出的.

B. 一个池塘里面有鲤鱼和草鱼，打捞了 100 条鱼，其中鲤鱼 80 条、草鱼 20 条，那么推测鲤鱼和草鱼的比例为 $4 : 1$ 时,出现 80 条鲤鱼、 20 条草鱼的概率是最大的.

C. $L\left( p\right)  = {p}^{\mathop{\sum }\limits_{{i = 1}}^{n}{x}_{i}}{\left( 1 - p\right) }^{n - \mathop{\sum }\limits_{{i = 1}}^{n}{x}_{i}}\left( {{x}_{i} = 0\text{ 或 }1,i = 1,2,\cdots ,n}\right)$ .

D. $L\left( p\right)$ 达到极大值时,参数 $p$ 的极大似然估计值为 $\frac{1}{n}\mathop{\sum }\limits_{{i = 1}}^{n}{x}_{i}$ .

\section*{三、解答题}

13. 甲虫是行动较快的昆虫之一，如表记录了某种类型的甲虫的爬行速度:

\begin{center}
\begin{tabular}{|c|c|c|c|c|c|c|c|}
\hline
时间 $t\left( {s}\right)$ & 1 & 2 & 3 & ... & ? & ... & 60 \\
\cline{1-8}
距离 $s\left( {{cm}}\right)$ & 9.8 & 19.6 & 29.4 & ... & 49 & ... & ? \\
\cline{1-8}
\hline
\end{tabular}
\end{center}

(1)你能建立一个等差数列的模型, 表示甲虫的爬行距离和时间之间的关系吗？

(2)利用建立的模型计算，甲虫 $1{{min}}$ 能爬多远？它爬行 ${49}{{cm}}$ 需要多长时间？

14. 学校新建的体育器材室的一面外墙如图 1 所示,它的轮廓由抛物线和矩形 ${ABCD}$ 构成. 数学兴趣小组要为器材室设计一个矩形标牌 ${EFGH}$ ,要求矩形 ${EFGH}$ 的顶点 $E\text{ 、 }H$ 在抛物线上,顶点 $F\text{ 、 }G$ 在矩形 ${ABCD}$ 的边 ${AD}$ 上. 为了设计面积最大的矩形 ${EFGH}$ ,兴趣小组对矩形 ${EFGH}$ 的面积与它的一边 ${FG}$ 的长之间的关系进行研究.

\begin{center}
\includegraphics[width=0.5\textwidth]{images/bo_d7tk7oc91nqc73875fu0_294_490_355_673_340_0.jpg}
\end{center}
 

(第 14 题图)

具体研究过程如下, 请补充完整.

(1)建立模型:

以 ${FG}$ 的中点为坐标原点,建立如图 2 所示的平面直角坐标系 ${xOy}$ ,通过研究发现,抛物线满足函数关系 $y =  - \frac{1}{4}{x}^{2} + 1\left( {-2 \leq  x \leq  2}\right)$ . 设矩形 ${EFGH}$ 的面积为 $S{{m}}^{2},{FG}$ 的长为 $a{m}$ ,则另一边 ${HG}$ 的长为 m(用含 $a$ 的代数式表示)，得到 $S$ 与 $a$ 的关系式为:\_\_\_\_\_ $\left( {0 < a < 4}\right)$ ；

我们的数学模型为: 当 ${FG}$ 为多长时,矩形 ${EFGH}$ 面积最大.

(2)求解模型:

利用数学知识,求解上述数学模型:

(3)模型总结:

当 ${FG}$ 的长约为\_\_\_\_\_ ${m}$ 时,矩形面积最大. (精确到 0.1)


15.  \ding{172}  在菜市上，摊主小王叫卖:“洋芋便宜卖了，两元一斤，三元两斤，四元三斤，五元四斤，六元五斤，快来买啊！”结果一群人都在买六元五斤的. 由此得到如下结论:一次购买的斤数越多，单价越低，请建立一个函数模型,来说明以上结论;

 \ding{173} 小王卖洋芋赚到了钱，想进行某个项目的投资，约定如下: \ding{172} 投资金额固定； \ding{173} 投资年数可自由选择，但最短 3 年，最长不超过 10 年； \ding{174} 投资年数 $x\left( {x \in  {\mathbf{N}}^{ * }}\right)$ 与总回报 $y$ 的关系，可选择下述三种方案中的一种:

方案一:当 $x = 3$ 时， $y = 6$ ，以后 $x$ 每增加 1 时， $y$ 增加 2；

方案二: $y = \frac{1}{3}{x}^{2}$ ;

方案三: $y = {\left( \sqrt[3]{3}\right) }^{x}$ .

请你根据以上材料，结合你的分析，为小王提供一个最佳投资方案.

16. 小张决定加盟某专营店自主创业,该专营店统计了近五年来创收利润数 ${y}_{i}$ (单位: 万元) 与时间 ${t}_{i}$ (单位:年)的数据，列表如下:

\begin{center}
\begin{tabular}{|c|c|c|c|c|c|}
\hline
${t}_{i}$ & 1 & 2 & 3 & 4 & 5 \\
\cline{1-6}
${y}_{i}$ & 2.4 & 2.7 & 4.1 & 6.4 & 7.9 \\
\cline{1-6}
\hline
\end{tabular}
\end{center}

(1)依据表中给出的数据，是否可用线性回归模型拟合 $y$ 与 $t$ 的关系，请计算相关系数 $r$ ，并加以说明 (计算结果精确到 0.01)(若 $\left| r\right|  > {0.75}$ ，则线性相关程度很高，可用线性回归模型拟合)；

(2)该专营店为吸引顾客，特推出两种促销方案.

方案一:每满 500 元可减 50 元；

方案二:每满 500 元可抽奖一次，每次中奖的概率都为 $\frac{2}{5}$ ，中奖就可以获得 100 元现金奖励. 假设顾客每次抽奖的结果相互独立.

 \ding{172} 某位顾客购买了 1050 元的产品，该顾客选择参加两次抽奖，求该顾客获得 100 元现金奖励的概率；

 \ding{173} 某位顾客购买了 1500 元的产品，作为专营店老板，是希望该顾客直接选择返回 150 元现金，还是选择参加三次抽奖? 说明理由.

\section*{§ 13.2 数学建模案例}

\section*{典型例题精讲}

【例 1】如图, “十字歇山”是由两个直三棱柱重叠后的景象, 重叠后的底面为正方形, 直三棱柱的底面是顶角为 ${120}^{ \circ  }$ ,腰为 3 的等腰三角形,求该几何体的体积.

\begin{center}
\includegraphics[width=0.4\textwidth]{images/bo_d7tk7oc91nqc73875fu0_296_938_531_543_266_0.jpg}
\end{center}
 

(例 1 图)

【例 2】某企业因技术升级,决定实行新的绩效方案.方案起草后,为了解员工对新绩效方案是否满意, 决定采取如下“随机化回答技术”进行问卷调查:

一个袋子中装有三个大小相同的小球, 其中 1 个黑球、 2 个白球. 企业所有员工从袋子中有放回的随机摸两次球, 每次摸出一球. 约定 “若两次摸到的球的颜色不同, 则按方式 I 回答问卷, 否则按方式 II 回答问卷”.

方式 I :若第一次摸到的是白球，则在问卷中画“○”，否则画“ \ding{55} ”；

方式 II :若你对新绩效方案满意，则在问卷中画“○”，否则画“ \ding{55} ”.

当所有员工完成问卷调查后, 统计画  \bigcirc , 画 $\times$ 的比例. 用频率估计概率, 由所学概率知识即可求得该企业员工对新绩效方案的满意度的估计值. 其中，满意度 $= \frac{\text{ 企业所有对新绩效方案满意的员工人数 }}{\text{ 企业所有员工人数 }}$

(1)若该企业某部门有 9 名员工，用 $X$ 表示其中按方式 I 回答问卷的人数，求 $X$ 的数学期望；

(2)若该企业的所有调查问卷中，画“○”与画“ \ding{55} ”的比例为 $4 : 5$ ，试估计该企业员工对新绩效方案的满意度.

【例 3】为了节能环保、节约材料，定义建筑物的“体形系数” $S = \frac{{F}_{0}}{{V}_{0}}$ ，其中 ${F}_{0}$ 为建筑物暴露在空气中的面积 (单位: ${{m}}^{2}$ ), ${V}_{0}$ 为建筑物的体积 (单位: ${{m}}^{3}$ ).

(1)若有一个圆柱体建筑的底面半径为 $R$ ，高度为 $H$ ，暴露在空气中的部分为上底面和侧面，试求该建筑体的“体形系数” $S$ ; (结果用含 $R$ 、 $H$ 的代数式表示)

(2)定义建筑物的“形状因子”为 $f = \frac{{L}^{2}}{A}$ ，其中 $A$ 为建筑物底面面积， $L$ 为建筑物底面周长，又定义 $T$ 为总建筑面积,即为每层建筑面积之和 (每层建筑面积为每一层的底面面积). 设 $n$ 为某宿舍楼的层数,层高为 $3{m}$ ,则可以推导出该宿舍楼的 “体形系数”为 $S = \sqrt{\frac{f \cdot  n}{T}} + \frac{1}{3n}$ . 当 $f = {18},T = {10000}$ 时,试求当该宿舍楼的层数 $n$ 为多少时,“体形系数” $S$ 最小.

【例 4】有一条长为 ${120}{m}$ 的步行道 ${OA},A$ 是垃圾投放点 ${\omega }_{1}$ ,若以 $O$ 为原点, ${OA}$ 为 $x$ 轴正半轴建立直角坐标系,设点 $B\left( {x,0}\right)$ ,现要建设另一座垃圾投放点 ${\omega }_{2}\left( {t,0}\right)$ ,函数 ${f}_{t}\left( x\right)$ 表示与点 $B$ 距离最近的垃圾投放点的距离.

(1)若 $t = {60}$ ，求 ${f}_{60}\left( {10}\right)$ 、 ${f}_{60}\left( {80}\right)$ 、 ${f}_{60}\left( {95}\right)$ 的值，并写出 ${f}_{60}\left( x\right)$ 的函数解析式；

(2)若可以通过 ${f}_{t}\left( x\right)$ 与坐标轴围成的面积来测算扔垃圾的便利程度，面积越小越便利. 问:垃圾投放点 ${\omega }_{2}$ 建在何处才能比建在中点时更加便利?

【例 5】某制药公司研制了一款针对某种病毒的新疫苗. 该病毒一般通过病鼠与白鼠之间的接触传染， 现有 $n$ 只白鼠,每只白鼠在接触病鼠后被感染的概率为 $\frac{1}{2}$ ,被感染的白鼠数用随机变量 $X$ 表示,假设每只白鼠是否被感染之间相互独立.

(1)若 $P\left( {X = 5}\right)  = P\left( {X = {95}}\right)$ ，求数学期望 $E\left\lbrack  X\right\rbrack$ ；

(2)接种疫苗后的白鼠被病鼠感染的概率为 $p$ ，现有两个不同的研究团队理论研究发现概率 $p$ 与参数 $\theta \left( {0 < \theta  < 1}\right)$ 的取值有关. 团队 $A$ 提出函数模型为 $p = \ln \left( {1 + \theta }\right)  - \frac{2}{3}{\theta }^{2}$ . 团队 $B$ 提出函数模型为 $p = \frac{1}{2}(1 -  {{e}}^{-\theta }$ ). 现将白鼠分成 10 组，每组 10 只，进行实验，随机变量 ${X}_{i}\left( {i = 1,2,\cdots ,{10}}\right)$ 表示第 $i$ 组被感染的白鼠数， 现将随机变量 ${X}_{i}\left( {i = 1,2,\cdots ,{10}}\right)$ 的实验结果 ${x}_{i}\left( {i = 1,2,\cdots ,{10}}\right)$ 绘制成频数分布图，如图所示.

 \ding{172} 试写出事件“ ${X}_{1} = {x}_{1},{X}_{2} = {x}_{2},\cdots ,{X}_{10} = {x}_{10}$ ”发生的概率表达式(用 $p$ 表示，组合数不必计算)；

 \ding{173} 在统计学中，若参数 $\theta  = {\theta }_{0}$ 时使得概率 $P\left( {{X}_{1} = {x}_{1},{X}_{2} = {x}_{2},\cdots ,{X}_{10} = {x}_{10}}\right)$ 最大，称 ${\theta }_{0}$ 是 $\theta$ 的最大似然估计. 根据这一原理和团队 $A\text{ 、 }B$ 提出的函数模型,判断哪个团队的函数模型可以求出 $\theta$ 的最大似然估计, 并求出最大似然估计.

\begin{center}
\includegraphics[width=0.4\textwidth]{images/bo_d7tk7oc91nqc73875fu0_298_902_745_580_342_0.jpg}
\end{center}
 

\section*{四基达标精练}

\section*{一、填空题}

1. 青少年视力是社会普遍关注的问题, 视力情况可借助视力表测量. 通常用五分记录法和小数记录法记录视力数据,五分记录法的数据 $L$ 和小数记录法的数据 $V$ 满足 $L = 5 + \lg V$ . 已知某同学视力的五分记录法的数据为 4.9 ，则其视力的小数记录法的数据约为\_\_\_\_\_. (精确到 0.1)

2. 从某网络平台推荐的影视作品中抽取 400 部, 统计其评分数据,将所得 400 个评分数据分为 8 组: $\lbrack {66},{70})$ , $\lbrack {70},{74}),\cdots \left\lbrack  {{94},{98}}\right\rbrack$ ,并整理得到如图所示频率分布直方图, 则评分在区间 $\lbrack {82},{86})$ 内的影视作品数量是\_\_\_\_\_.

\begin{center}
\includegraphics[width=0.5\textwidth]{images/bo_d7tk7oc91nqc73875fu0_298_877_1675_621_371_0.jpg}
\end{center}
 

3. 图 1 是中国古代建筑中的举架结构， $A{A}^{\prime }$ 、 $B{B}^{\prime }$ 、 $C{C}^{\prime }$ 、 $D{D}^{\prime }$ 是桁，相邻桁的水平距离称为步，垂直距离称为举. 图 2 是某古代建筑屋顶截面的示意图,其中 $D{D}_{1}$ 、 $C{C}_{1}$ 、 $B{B}_{1}$ 、 $A{A}_{1}$ 是举, $O{D}_{1}$ 、 $D{C}_{1}$ 、 $C{B}_{1}$ 、 $B{A}_{1}$ 是相等的步,相邻桁的举步之比分别为 $\frac{D{D}_{1}}{O{D}_{1}} = {0.5},\frac{C{C}_{1}}{D{C}_{1}} = {k}_{1},\frac{B{B}_{1}}{C{B}_{1}} = {k}_{2},\frac{A{A}_{1}}{B{A}_{1}} = {k}_{3}$ . 已知 ${k}_{1},{k}_{2},{k}_{3}$ 成公是为 0.1 的等差数列,且直线 ${OA}$ 的斜率为 0.725 ,则 ${k}_{3} =$ \_\_\_\_\_.

\begin{center}
\includegraphics[width=0.6\textwidth]{images/bo_d7tk7oc91nqc73875fu0_299_433_380_743_247_0.jpg}
\end{center}
 

(第 3 题图)

4. 在卫星导航系统中,地球静止同步轨道卫星的轨道位于地球赤道所在平面,轨道高度为 ${36000}{{km}}$ (轨道高度是指卫星到地球表面的距离). 将地球看作是一个球心为 $O$ ,半径 $r$ 为 ${6400}{{km}}$ 的球,其上点 $A$ 的纬度是指 ${OA}$ 与赤道平面所成角的度数. 地球表面上能直接观测到的一颗地球静止同步轨道卫星点的纬度最大值为 $\alpha$ ,该卫星信号覆盖地球表面的表面积 $S = {2\pi }{r}^{2}\left( {1 - \cos \alpha }\right)$ (单位: ${{{km}}}^{2}$ ),则 $S$ 占地球表面积的百分比约为\_\_\_\_\_. (精确到 $1\%$ )

5. 某年 $A$ 公司的股价在年初 $\left( {t = 0\text{ 时 }}\right)$ 的股价是每股 5 元人民币,到了年末 $\left( {t = {12}\text{ 时 }}\right)$ 涨到了每股 6 元人民币. 经过建立模型分析发现,在该年第 $t$ 个月的时候, $A$ 公司的股价可以用函数 $A = {A}_{0}{{e}}^{kt}$ 来表示,其中 $k$ 为常数. 假设 $A$ 公司的股价继续按照上述的模型持续增长,则当 $A$ 公司的股价涨到 10 元时, $t$ 的值约为\_\_\_\_\_(结果精确到个位数).

6. 甲烷分子式为 ${{{CH}}}_{4}$ ，其结构抽象成的立体几何模型如图所示，碳原子位于四个氢原子的正中间位置，四个碳氢键长度相等,且任意两个氢原子等距排列. 用 $C$ 表示碳原子的位置,用 ${H}_{1}\text{ 、 }{H}_{2}\text{ 、 }{H}_{3}\text{ 、 }{H}_{4}$ 表示四个氢原子的位置，设 $\alpha  = \left\langle  {\overrightarrow{C{H}_{1}},\overrightarrow{C{H}_{4}}}\right\rangle$ ，则 $\cos {2\alpha } =$ \_\_\_\_\_.

7.2020 年 12 月 8 日,中国和尼泊尔联合公布珠穆朗玛峰最新高程为 8848.86 (单位: ${m}$ ),三角高程测量法是珠峰高程测量方法之一. 如图是三角高程测量法的一个示意图,现有 $A\text{ 、 }B\text{ 、 }C$ 三点在同一水平面上的投影 ${A}^{\prime }\text{ 、 }{B}^{\prime }\text{ 、 }{C}^{\prime }$ ,且 $\angle {A}^{\prime }{C}^{\prime }{B}^{\prime } = {45}^{ \circ  },\angle {A}^{\prime }{B}^{\prime }{C}^{\prime } = {60}^{ \circ  }$ ,由点 $C$ 测得点 $B$ 的仰角为 ${15}^{ \circ  },B{B}^{\prime }$ 与 $C{C}^{\prime }$ 的差为 100 ; 由点 $B$ 测得点 $A$ 的仰角为 ${45}^{ \circ  }$ ，则 $A$ 、 $C$ 两点到水平面 ${A}^{\prime }{B}^{\prime }{C}^{\prime }$ 的高度差 $A{A}^{\prime } - C{C}^{\prime }$ 约为\_\_\_\_\_m.

\begin{center}
\includegraphics[width=0.2\textwidth]{images/bo_d7tk7oc91nqc73875fu0_299_373_1426_263_228_0.jpg}
\end{center}
 

(第 6 题图)

\begin{center}
\includegraphics[width=0.2\textwidth]{images/bo_d7tk7oc91nqc73875fu0_299_927_1389_298_314_0.jpg}
\end{center}
 

(第 7 题图)

8. 某公园欲建设一段斜坡，坡顶是一条直线，斜坡顶点距水平地面的高度为 4m，坡面与水平面所成夹角为 $\theta$ . 行人每沿着斜坡向上走 $1{m}$ 消耗的体力为 $\left( {{1.025} - \cos \theta }\right)$ ,欲使行人走上斜坡所消耗的总体力最小，则 $\theta  =$

9. 上海海关大楼的顶部为逐级收拢的四面钟楼，如图，四个大钟分布在四棱柱的四个侧面, 则每天 0 点至 12 点 (包含 0 点, 不含 12 点) 相邻两钟面上的时针相互垂直的次数为\_\_\_\_\_.

\begin{center}
\includegraphics[width=0.2\textwidth]{images/bo_d7tk7oc91nqc73875fu0_299_1179_1837_278_227_0.jpg}
\end{center}
 

(第 9 题图)

\section*{二、选择题}

\begin{center}
\includegraphics[width=0.4\textwidth]{images/bo_d7tk7oc91nqc73875fu0_300_938_228_543_308_0.jpg}
\end{center}
 

(第 10 题图)

10.某社区通过公益讲座以普及社区居民的垃圾分类知识. 为了解讲座效果，随机抽取 10 位社区居民，让他们在讲座前和讲座后各回答一份垃圾分类知识问卷, 这 10 位社区居民在讲座前和讲座后问卷答题的正确率如图: 则 ( )

A. 讲座前问卷答题的正确率的中位数小于 70\%.

B. 讲座后问卷答题的正确率的平均数大于 85\%.

C. 讲座前问卷答题的正确率的标准差小于讲座后正确率的标准差.

D. 讲座后问卷答题的正确率的极差大于讲座前正确率的极差.

11. 关于外来物种泛滥的有害性经常引起热议. 为了研究某池塘里某种植物生长面积 $S$ (单位: ${{m}}^{2}$ ) 与时间 $t$ (单位: 月)之间的关系,通过观察建立了函数模型 $S\left( t\right)  = k{a}^{t}\left( {t \in  \mathbf{Z},k > 0,a > 0\text{ 且 }a \neq  1}\right)$ . 已知第一个月该植物的生长面积为 $1{{m}}^{2}$ ，第 3 个月该植物的生长面积为 $4{{m}}^{2}$ ，则该植物的生长面积达到 ${100}{{m}}^{2}$ ，至少要经过\_\_\_\_\_ ( )

A. 6 个月. B. 8 个月. C. 9 个月. D. 11 个月.

12. 阻尼器是一种以提供运动的阻力, 从而达到减振效果的专业工程装置. 如图, 某阻尼器模型的运动过程可近似为单摆运动,其离开平衡位置的位移 $s\left( {{cm}}\right)$ 和时间 $t\left( {s}\right)$ 的函数关系式 $\left( {\omega  > 0}\right)$ 为 $s\left( t\right)  = 3\sin \left( {{\omega t} + \varphi }\right)$ ,若该阻尼器模型在摆动过程中连续三次位移为 ${s}_{0}\left( {-3 < {s}_{0} < 3}\right)$ 的时间分别为 ${t}_{1}\text{ 、 }{t}_{2}\text{ 、 }{t}_{3}$ ,且 ${t}_{1} + {t}_{2} = 2,{t}_{2} + {t}_{3} = 4$ ,则 $s\left( t\right)$ 的单调区间是 ( )

\begin{center}
\includegraphics[width=0.3\textwidth]{images/bo_d7tk7oc91nqc73875fu0_300_1131_951_354_243_0.jpg}
\end{center}
 

(第 12 题图)

A. $\left\lbrack  {k,k + 1}\right\rbrack  \left( {k \in  \mathbf{N}}\right)$ .

B. $\left\lbrack  {k + \frac{1}{2},k + \frac{3}{2}}\right\rbrack  \left( {k \in  \mathbf{N}}\right)$ .

C. $\left\lbrack  {{2k} + 1,{2k} + 2}\right\rbrack  \left( {k \in  \mathbf{N}}\right)$ .

D. $\left\lbrack  {{2k} + \frac{1}{2},{2k} + \frac{3}{2}}\right\rbrack  \left( {k \in  \mathbf{N}}\right)$ .

\section*{三、解答题}

13. 甲、乙两地相距 1004km，汽车从甲地匀速驶向乙地，速度不得超过 120km/h，已知汽车每小时的运输成本 (以 1 元为单位) 由可变部分和固定部分组成:可变部分与速度 $v\left( {{{km}}/{h}}\right)$ 的立方成正比，比例系数为 2，固定部分为 $a$ 元 $\left( {a > 0}\right)$ .

(1)把全部运输成本 $y$ 元表示为速度 $v\left( {{{km}}/{h}}\right)$ 的函数，并指出这个函数的定义域；

(2)为了使全部运输成本最小，汽车应以多大速度行驶 利(用 $a$ 表示)？


14. 某学校组织科普知识竞赛, 有 A、B 两类问题. 每位参加比赛的同学先在两类问题中选择一类并从中随机抽取一个问题回答, 若回答错误则该同学比赛结束; 若回答正确则从另一类问题中再随机抽取一个问题回答，无论回答正确与否，该同学比赛结束. A 类问题中的每个问题回答正确得 20 分，否则得 0 分；B 类问题中的每个问题回答正确得 80 分，否则得 0 分.

已知小明能正确回答 A 类问题的概率为 0.8 ，能正确回答 B 类问题的概率为 0.6 ，且能正确回答问题的概率与回答次序无关.

(1)若小明先回答 A 类问题，记 $X$ 为小明的累计得分，求 $X$ 的分布；

(2)为使累计得分的期望最大，小明应选择先回答哪类问题？并说明理由.

15. 摩天轮是一种大型转轮状的机械建筑设施, 游客坐在摩天轮的座舱里慢慢地往上转, 可以从高处俯瞰四周景色，如图，某摩天轮最高点距离地面 100m，最低点距离地面 10m，摩天轮上均匀设置了依次标号为 1\textasciitilde 36 号的 36 个座舱. 开启后摩天轮按照逆时针方向匀速旋转，游客在座舱转到距离地面最近的位置进舱， 开始转动 $t$ 分钟后距离地面的高度为 $H{m}$ ，转一周需要 ${30}{{min}}$ .

(1)求在转动一周的过程中， $H$ 关于 $t$ 的函数解析式；

(2)若甲、乙两人分别坐在 1 号和 7 号座舱里，在转动一周的过程中，求两人距离地面的高度差 $h$ (单位: ${m}$ ) 关于 $t$ 的函数解析式，并求高度差的最大值.

\begin{center}
\includegraphics[width=0.6\textwidth]{images/bo_d7tk7oc91nqc73875fu0_301_716_1432_707_292_0.jpg}
\end{center}
 

(第 15 题图)

16. 近年来短视频平台吸引了大量流量和网络博主的加入. 红人榜的数据推出是体现各平台 KOL 网络博主商业价值的榜单. 每周一期，红人榜能反应最近一周 KOL 网络的综合价值，以粉丝数、集均评论、集均赞，以及集均分享来进行综合衡量. 红人榜单在统计时发现某平台一网络博主的累计粉丝数 $y$ (百万) 与人驻平台周次 $x$ (周)之间的关系如图所示.

设 $\omega  = \ln x$ ,数据经过初步处理得 $\mathop{\sum }\limits_{{i = 1}}^{5}{y}_{i} = {73},\mathop{\sum }\limits_{{i = 1}}^{5}{w}_{i} = 5,\mathop{\sum }\limits_{{i = 1}}^{5}{w}_{i}^{2} = {6.3},\mathop{\sum }\limits_{{i = 1}}^{5}{x}_{i}{y}_{i} = {258},\mathop{\sum }\limits_{{i = 1}}^{5}{w}_{i}{y}_{i} = {86},\mathop{\sum }\limits_{{i = 1}}^{5}\left( {{y}_{i} - }\right.  \bar{y}{)}^{2} = {160}$ ,线性回归模型残差平方和 $\mathop{\sum }\limits_{{i = 1}}^{5}{\left( {y}_{i} - \widehat{y}\right) }^{2} = 9$ . (其中 ${x}_{1}\text{ 、 }{y}_{i}$ 分别为观测数据中的周次和累计粉丝数)

(1)求出 $y$ 关于 $x$ 的线性回归模型 $\widehat{y} = \widehat{b}x + \widehat{a}$ 的相关指数 ${R}_{1}^{2}$ . 若用非线性回归模型 $\widehat{y} = \widehat{c}\ln x + \widehat{d}$ 求得的相关指数 ${R}_{2}^{2} = {0.9998}$ ,试用相关指数 ${R}^{2}$ 判断哪种模型的拟合效果较好? (相关指数越接近于 1,拟合效果越好)

(2)根据(1)中拟合效果较好的模型求出 $y$ 关于 $x$ 的回归方程，并由此预测入驻平台 8 周后，对应的累计粉丝数 $y$ 为多少?

附参考公式: 相关指数 ${R}^{2} = 1 - \frac{\mathop{\sum }\limits_{{i = 1}}^{n}{\left( {y}_{i} - {\widehat{y}}_{i}\right) }^{2}}{\mathop{\sum }\limits_{{i = 1}}^{n}{\left( {y}_{i} - \bar{y}\right) }^{2}}$ ,

$y$ 与 $x$ 的线性回归方程 $\widehat{y} = \widehat{b}x + \widehat{a}$ 中, $\widehat{b} = \frac{\mathop{\sum }\limits_{{i = 1}}^{n}{x}_{i}{y}_{i} - n\overline{xy}}{\mathop{\sum }\limits_{{i = 1}}^{n}{x}_{i}^{2} - n{\bar{x}}^{2}},\widehat{a} = \bar{y} - \widehat{b}\bar{x}$ .

\begin{center}
\includegraphics[width=0.3\textwidth]{images/bo_d7tk7oc91nqc73875fu0_302_1023_915_462_284_0.jpg}
\end{center}